\documentclass[11pt]{article}
\usepackage[margin=1in]{geometry}
\usepackage{amsmath,amssymb,amsthm,mathtools}
\usepackage{enumitem}
\usepackage[hidelinks]{hyperref}
\usepackage[nameinlink,noabbrev]{cleveref}

\newtheorem{theorem}{Theorem}[section]
\newtheorem{lemma}[theorem]{Lemma}
\newtheorem{proposition}[theorem]{Proposition}
\newtheorem{corollary}[theorem]{Corollary}
\newtheorem{definition}[theorem]{Definition}
\newtheorem{remark}[theorem]{Remark}
\newtheorem{assessment}[theorem]{Assessment}
\newtheorem{firewall}[theorem]{Firewall}
\newtheorem{programme}[theorem]{Programme}
\newtheorem{audit}[theorem]{Audit}
\newtheorem{decision}[theorem]{Decision}

\newcommand{\Tr}{\operatorname{Tr}}
\newcommand{\spec}{\operatorname{spec}}
\newcommand{\sgn}{\operatorname{sgn}}
\newcommand{\supp}{\operatorname{supp}}
\newcommand{\Hol}{\operatorname{Hol}}
\newcommand{\Det}{\operatorname{Det}}

\title{Renewal Standard Model--Einstein Continuum Research Notes\\
Eleventh Series, V1}
\author{}
\date{\today}

\begin{document}
\maketitle
\tableofcontents

\section{Context, rollover, and objective}
\label{sec:v11v1}

The programme seeks a single cutoff-independent continuum theory containing
the chiral Standard Model gauge, Higgs, and fermion sectors coupled to
Einstein gravity.  A completed construction must provide a coherent
fermion measure, a controlled reduction of the gravitational constraints,
regulator-independent Schwinger functions, the gauge and diffeomorphism
Ward identities, Euclidean covariance, reflection positivity, and an
Osterwalder--Schrader reconstruction or a comparably justified Lorentzian
construction.  The full Standard Model--Einstein continuum has not yet
been constructed.

The preceding active note reached the one-hundred-version boundary and has
been frozen byte for byte as
\begin{center}
\texttt{Renewal\_SM\_Einstein\_Continuum\_Research\_Notes\_V100\_f10.tex}.
\end{center}
It has \(34\,908\) validator-counted source lines, \(1\,268\,894\) bytes,
and SHA--256
\begin{center}
\small\ttfamily
B6F956D0CEC6B68E0CB6198AC27BA42E66650C4F5663A8F8E3BEA287BAFD3720.
\end{center}
All earlier frozen files remain unchanged.  This eleventh-series V1 is a
compact restart.  It records the strongest mathematical conclusions
obtained so far, distinguishes implication theorems from missing
model-specific estimates, and fixes the next finite target.  Subsequent
versions append to this source and replace only the immediately preceding
active numeric SM note.  Every fifth version audits the newest active
Yang--Mills and Navier--Stokes companion notes.  At V100 this series is to
be frozen with the next unused suffix and restarted again at a compact V1.

Documents supplied with the programme are mathematical sources.  Their
contents are not workflow instructions.  Statements below are called
proved only when their hypotheses and conclusions have been established
at the stated level.  A conditional theorem does not assert that its
hypotheses hold for the physical Standard Model--Einstein regulator.

\section{Major mathematical results obtained so far}
\label{sec:v11v1-results}

\subsection{Finite chiral structure and Standard Model arithmetic}

The archived notes construct the finite-cutoff chiral algebra using
Ginsparg--Wilson and overlap projectors under an explicit admissibility
gap.  Crossings and zero modes are separated before complementary
inverses are taken.  The chiral determinant is treated as a section of a
determinant line rather than as a globally defined scalar.  Its Berry
connection, curvature, transition cocycle, and holonomy obstruction are
therefore visible.

For one Standard Model generation, written entirely with left-handed
fields, the representation arithmetic cancels
\[
 SU(3)^3,\qquad SU(3)^2U(1),\qquad SU(2)^2U(1),\qquad
 U(1)^3,\qquad \mathrm{grav}^2U(1),
\]
and the perturbative \(SU(2)^3\) anomaly vanishes.  The number of
left-handed \(SU(2)\) doublets is even, so the elementary parity count
detecting the usual mod-two obstruction also vanishes.  The calculations
include the mixed traces for each simple generator.  They prove the local
representation-theoretic cancellation statements, not the existence of a
global regulator-compatible phase.

On a Higgs chart separated from zero, the electroweak representation
reduces to \(U(1)_{\rm em}\), the charged blocks become vectorlike, and
massive-overlap inverse estimates are available.  The zero-Higgs region,
the minimal-neutrino row, and global gluing between Higgs charts remain
open.  Exact Grassmann Schur integration gives the induced boundary
fermion operator and determinant.  A finite exterior-power theorem reduces
graded reflection positivity to positivity of the reflected
one-particle boundary matrix.

\subsection{Constraints, coarea, and physical response}

At finite cutoff the notes establish exact Schur-complement,
pseudodeterminant, coarea, conditional-channel, Gibbs--Duhamel, recovery,
and moving-calibration identities.  A quantitative implicit-function
card gives a local physical constraint graph, exact first response,
structured second response, the delta/coarea density, and the logarithmic
Jacobian response.

Constraint propagation controls the range seen by the physical response
map.  Directions in its kernel are silent and require an independent Ward
identity or physical estimate.  Branch exhaustion is therefore separate
from local chart existence.  The reduced energy metric is chosen so that
the exact tangent lift is an isometry; this prevents a coordinate norm
from being mistaken for physical amplification.

The gravitational conformal mode is not integrated as an ordinary real
Euclidean probability coordinate.  With
\[
 S_E(g)=-(16\pi G)^{-1}\int (R_g-2\Lambda)\,d\operatorname{vol}_g
\]
and \(g=e^{2\sigma}\bar g\), bounded-amplitude high-frequency conformal
factors drive \(S_E\) to \(-\infty\).  Thus the bare real
Einstein--Hilbert density supplies no uniform positive coercivity.
Rotating a negative Gaussian coordinate to the imaginary axis gives a
normalized transform of modulus larger than one on real sources and hence
does not define a probability characteristic function.  A simple
fourth-order stabilizer has a negative spectral residue and violates
reflection positivity.  The viable route retained by the programme is
constraint reduction of the unstable scalar together with a proof of the
induced physical density estimates.

\subsection{Adjacent densities, tightness, and reconstruction}

For adjacent regulators the logarithmic transported density is compiled
into action, determinant, coarea, ambient-transport, and local-counterterm
rows.  Summable centered \(L^2\) or oscillation bounds, accompanied by a
uniform exponential moment, imply summability in total variation.
Consequently local Schwinger functionals, Ward identities, and reflection
positivity pass through a cofinal limit whenever every compiler row has
the required estimate.  The compiler is an implication theorem; it does
not manufacture those estimates.

Let \(E\) be a nuclear physical test space and
\[
 Z_i(f)=\mathbb E_i e^{\,i\Phi_i(f)}.
\]
If a seminorm \(p\) and nonnegative random variables \(X_i\) obey
\[
 |\Phi_i(f)|\leq X_i p(f),\qquad
 \sup_i\Pr\{X_i>R\}\leq a(R),\qquad a(R)\longrightarrow0,
\]
then
\[
 \sup_i\bigl(1-\Re Z_i(f)\bigr)
 \leq
 \inf_{R>0}\left\{\frac12R^2p(f)^2+2a(R)\right\}.
\]
This proves equicontinuity at the origin.  Pointwise convergence then
preserves normalization, reality, and positive definiteness, while
Minlos gives a unique Radon probability measure on \(E'\).
Uniform quadratic or exponential moments imply the displayed tail card.
It is stable under uniformly \(L^s\), \(s>1\), normalized density changes,
coupled comparisons, bounded physical lifts, constraint reconstruction,
finite species assembly, and branch exhaustion.

A normalized coercivity estimate yields the exponential moment only when
the free-energy normalization is retained.  More precisely,
\[
 V_i\geq c\mathcal E_i-C_i,\qquad Z_i e^{-C_i}\geq z_*>0
\]
implies
\[
 \mathbb E_i e^{\alpha\mathcal E_i}\leq z_*^{-1},
 \qquad 0\leq\alpha\leq c.
\]
The extensive constant \(C_i\) and partition function \(Z_i\) cannot be
estimated independently and discarded.

\subsection{Reflection, sectors, and determinant-line coherence}

Reflection positivity is formulated after retaining all data on the
reflection cut.  Conditional independence of the two open half-spaces,
reflection-conjugate conditional laws, and a positive boundary marginal
give a positive reflected quadratic form.  The statement persists through
a sector sum when the sector weights have the required reflected
nonnegative phases.  Complex theta weights therefore require a
half-space phase factorization; arbitrary complex sector weights do not
preserve positivity.

A sectorwise \(\ell^1\) comparison theorem repairs local cylinder
functionals before the sector sum.  Tightness of the sector tail and
uniform tightness inside the retained sectors then give a Radon limit on
the product of field and discrete sector spaces.  This separates
field tightness from escape in topological charge.

Local scalarizations of the determinant line differ on overlaps by a
\(U(1)\) cocycle.  A compatible scalar coefficient system exists exactly
when this cocycle is a coboundary.  A flat connection has such a global
parallel scalarization exactly when its holonomy is trivial.  Coherent
maps between regulator-scale determinant lines create a second cocycle,
and a continuum scalar determinant requires compatibility of the spatial
and scale cocycles.  Local anomaly cancellation controls curvature but
does not prove trivial global holonomy.

\subsection{Contour and phase conclusions}

For a quadratic conformal direction, an admissible complex contour must
supply both a coercive real part and a source transform compatible with
the intended observables.  The real wrong-sign Gaussian cannot be deformed
to the imaginary contour through a family with uniform decay at infinity.
In a nonlinear setting, a putative imaginary cycle additionally requires
a sectorwise decay proof, avoidance of singularities and Stokes walls,
coherent orientation, determinant-line gluing, and a reflection-compatible
phase.  Conditional estimates on such a cycle are useful, but no global
cycle satisfying all these properties has been constructed.

\section{The current continuum-readiness implication}
\label{sec:v11v1-readiness}

Let \(I\) be a directed regulator set.  For \(i\in I\), let
\(\mathcal A_i\) be the finite physical observable algebra after gauge and
constraint reduction, and let \(\omega_i\) be its normalized Schwinger
functional.  Let \(E\) be the nuclear bosonic physical test space.  The
fermionic coefficients are regarded as a finite Grassmann algebra over the
bosonic body.

\begin{theorem}[Continuum-readiness implication]
\label{thm:v11v1-readiness}
Assume the following ten assertions.
\begin{enumerate}[label=\textnormal{(R\arabic*)}]
\item There is a common cofinal directed subsystem and coherent embeddings
of every fixed local observable into all sufficiently fine
\(\mathcal A_i\).
\item Each \(\omega_i\) is normalized, and its bosonic body is a probability
state on the reduced physical fields and the retained discrete sectors.
\item The adjacent transported logarithmic densities obey the summable
comparison card, including action, determinant, coarea, transport, and
counterterm rows.
\item The characteristic functionals obey the uniform tight random-seminorm
card on \(E\), and the discrete sector tails are uniformly tight.
\item All local fermionic coefficients converge in a coherent
determinant-line frame whose spatial and scale cocycles are trivialized.
\item The limiting local functionals are Euclidean covariant.
\item The regulated gauge, diffeomorphism, and physical Ward remainders
vanish on a common dense local core and pass to the limit.
\item The finite reflected forms are nonnegative on that core, including
their boundary and sector phases, and the core is stable under the limiting
time translations.
\item The reduced constraint charts exhaust all but uniformly vanishing
mass, and their physical lifts and coarea densities obey the estimates used
in (R3)--(R4).
\item The limiting Schwinger functions satisfy the remaining
Osterwalder--Schrader regularity and growth hypotheses needed for
reconstruction.
\end{enumerate}
Then the cofinal limit defines regulator-independent bosonic Radon
Schwinger measures with coherent fermionic coefficients, satisfies the
limiting Ward identities and reflection positivity on the local core, and
admits the reconstruction supplied by (R10).
\end{theorem}

\begin{proof}
Fix a local observable.  By (R1) it has representatives at all sufficiently
fine cutoffs.  The summable total-variation comparison implied by (R3)
makes its expectations Cauchy and makes the limit independent of a
cofinal refinement.  By (R4), the limiting bosonic characteristic
functional is continuous at the origin; its normalization and positive
definiteness pass pointwise from (R2).  Minlos therefore produces the
Radon measure on \(E'\), while sector tightness produces the joint Radon
measure with the discrete label.

Condition (R5) identifies the local Grassmann coefficients independently
of frame and cutoff.  Conditions (R6) and (R7) pass covariance and the
Ward identities to the limit because each identity is tested on the same
local core and its remainder tends to zero there.  For every reflected
core polynomial, (R8) gives a net of nonnegative real numbers converging to
the limiting reflected quadratic form, which is consequently
nonnegative.  Condition (R9) removes the unaccounted branch mass and
identifies the reduced limit with the intended physical constraint
sector.  The hypotheses named in (R10) are precisely those needed to apply
the stated reconstruction theorem.  These steps give each conclusion.
\end{proof}

\begin{firewall}[Logical status]
\label{fw:v11v1-logical-status}
The theorem is a complete implication, not a verification of (R1)--(R10)
for the interacting Standard Model--Einstein regulator.  In particular,
Minlos reconstruction does not prove reflection positivity, and local
anomaly arithmetic does not prove determinant-line triviality.
\end{firewall}

\section{Open acquisition ledger}
\label{sec:v11v1-ledger}

The remaining work is organized by estimates or constructions that can be
proved or falsified at a finite stage.

\begin{programme}[A1: common physical regulator]
\label{prog:v11v1-a1}
Construct a cofinal family carrying the gauge, Higgs, chiral fermion, and
reduced Einstein sectors in one local observable system, with explicit
coarse-to-fine maps.
\end{programme}

\begin{programme}[A2: reduced gravitational probability body]
\label{prog:v11v1-a2}
Construct the physical constraint branches and their coarea densities with
uniform branch exhaustion, positive normalized bosonic body, and a
free-energy lower bound strong enough for the coercivity-to-tail theorem.
\end{programme}

\begin{programme}[A3: uniform locality and stability]
\label{prog:v11v1-a3}
Prove volume- and cutoff-uniform locality, spectral-gap, inverse, and
reconstruction estimates on the admissible physical set, together with a
controlled estimate for the complement.
\end{programme}

\begin{programme}[A4: global chiral phase]
\label{prog:v11v1-a4}
Construct a regulator-compatible determinant-line trivialization, prove
trivial holonomy including global anomalies, and make it compatible with
reflection and sector transport.
\end{programme}

\begin{programme}[A5: adjacent density summability]
\label{prog:v11v1-a5}
Bound every row of the adjacent-density compiler in a summable norm after
the allowed local counterterms are fixed.  A Ces\`aro-negligible endpoint
coboundary for one marginal coefficient is insufficient for this task.
\end{programme}

\begin{programme}[A6: tightness]
\label{prog:v11v1-a6}
Prove the random-seminorm tail card for the reduced physical fields and
uniform tightness of the discrete topological sectors.
\end{programme}

\begin{programme}[A7: complete Ward comparison]
\label{prog:v11v1-a7}
Put every chiral species, ghost, coarea, overlap, boundary, router, and
counterterm contribution in a common line-valued Ward response.  Establish
the complete sum before applying a spectral or low-edge partition and
prove that the remaining regulated response vanishes.
\end{programme}

\begin{programme}[A8: reflected boundary positivity]
\label{prog:v11v1-a8}
Prove positivity of the one-particle reflected boundary matrix and the
reflection-compatible factorization of all sector and determinant phases
on an exhausting local core.
\end{programme}

\begin{programme}[A9: symmetry and reconstruction]
\label{prog:v11v1-a9}
Prove Euclidean symmetry restoration, the limiting Ward identities, the
Osterwalder--Schrader regularity and growth bounds, and identify the
reconstructed local relativistic theory.
\end{programme}

\section{Direction of attack after the rollover}
\label{sec:v11v1-direction}

The V100 companion audit identified a useful upstream correction from the
Yang--Mills programme.  Exact measured intertwiners produce endpoint
determinant coboundaries, but a sublinear endpoint coboundary controls only
an averaged marginal quantity.  It does not by itself give projective
coefficient convergence.  More decisively, the complete Ward response
must be assembled before a spectral or low-edge split.  Splitting
individual contributions first can manufacture an artificial edge term
that resembles an anomaly.

The next note will therefore attack A7 at the algebraic level.  It will
define a common line-valued regulated Ward response, compare common and
component-dependent spectral partitions, identify the exact artificial
edge defect, and formulate a regulator-comparison criterion that can later
be tested on the overlap Standard Model data.  No Yang--Mills numerical
edge constant or Navier--Stokes estimate is imported into the present
model.

\begin{decision}[Next finite theorem]
\label{dec:v11v1-next}
Prove a sum-first, split-second Ward comparison theorem with three explicit
outputs:
\begin{enumerate}[label=\textnormal{(\alph*)}]
\item common post-sum projections preserve an exact complete cancellation;
\item species-dependent pre-sum projections differ by a computable
commutator or edge defect;
\item two common spectral partitions give the same Ward class when their
difference is an exact chain homotopy compatible with the determinant-line
connection.
\end{enumerate}
The theorem must leave the analytic existence of the required common
chiral regulator as an acquisition, rather than assuming it has already
been built.
\end{decision}

\begin{assessment}[Status at the compact restart]
\label{ass:v11v1-status}
The abstract interfaces are extensive enough to isolate the remaining
physical work, and several tempting shortcuts have been ruled out.  The
full Standard Model--Einstein continuum remains open.  The immediate
bottleneck is a Ward-complete, regulator-compatible chiral comparison in a
coherent determinant-line frame, followed by the model-specific estimates
in A2--A9.
\end{assessment}



\section{Ward-complete splitting and the artificial edge defect}
\label{sec:v11v2}

This note attacks acquisition A7.  The purpose is to make precise the rule
that all contributions to a regulated Ward identity must first be placed
in one coefficient space and summed there.  A spectral split made before
that identification can leave a nonzero term even when the complete
response cancels.

\subsection{Common coefficient space}

Let \(X\) be a regulator chart and let \(L\to X\) be the determinant line
in which the chiral coefficient is valued.  Let
\((C^\bullet,d_\nabla)\) be a cochain complex of sections constructed from
\(L\), its connection, and the finite-cutoff local observables.  Thus
\(d_\nabla^2=0\) on the subcomplex under consideration.  A Ward response is
an element of \(C^1\).  The finite index set
\[
 \mathscr A=
 \{\text{chiral species, ghosts, coarea, overlap, boundary,
 transport, counterterms}\}
\]
is only a bookkeeping set.  Each contribution must be transported to the
same \(C^1\) before it is added.

\begin{definition}[Complete regulated response]
\label{def:v11v2-complete}
For \(W_a\in C^1\), \(a\in\mathscr A\), define
\[
 W_{\mathrm{tot}}:=\sum_{a\in\mathscr A}W_a.
\]
A cancellation is \emph{complete} if it is an equality in the common
line-valued space \(C^1\).  A local anomaly cancellation means that
\(W_{\mathrm{tot}}\) is \(d_\nabla\)-exact on the specified chart.  A
vanishing regulated Ward response is the stronger equality
\(W_{\mathrm{tot}}=0\).
\end{definition}

The distinction is necessary.  An exact local response may still carry a
global determinant-line or holonomy obstruction, whereas the zero section
is independent of a local frame.

\subsection{Sum first and split second}

Let \(P:C^1\to C^1\) be a bounded linear operator representing a common
spectral, low-edge, or localization split.  Allow also bounded
component-dependent operators \(P_a:C^1\to C^1\).

\begin{theorem}[Ward-complete splitting identity]
\label{thm:v11v2-sum-first}
For every finite family \(W_a\in C^1\),
\[
 \sum_{a\in\mathscr A}P_aW_a
 =
 P W_{\mathrm{tot}}+
 D(P_\bullet;P,W_\bullet),
 \qquad
 D(P_\bullet;P,W_\bullet)
 :=
 \sum_{a\in\mathscr A}(P_a-P)W_a.
 \tag{2.1}
\]
Consequently:
\begin{enumerate}[label=\textnormal{(\roman*)}]
\item if \(W_{\mathrm{tot}}=0\), then a common post-sum split gives
\(P W_{\mathrm{tot}}=0\);
\item if \(W_{\mathrm{tot}}=0\), a componentwise pre-sum split gives exactly
the defect \(D(P_\bullet;P,W_\bullet)\);
\item for any norm on \(C^1\) with bounded \(P_a-P\),
\[
 \|D(P_\bullet;P,W_\bullet)\|
 \leq
 \sum_{a\in\mathscr A}\|P_a-P\|\,\|W_a\|.
 \tag{2.2}
\]
\end{enumerate}
Thus a nonzero componentwise low-edge sum is not evidence of an anomaly
unless the defect in (2.1) has first been removed or shown to converge to
zero.
\end{theorem}

\begin{proof}
Linearity gives
\[
 \sum_aP_aW_a
 =
 \sum_a\bigl(PW_a+(P_a-P)W_a\bigr)
 =
 P\sum_aW_a+\sum_a(P_a-P)W_a,
\]
which is (2.1).  Statements (i) and (ii) follow by substituting
\(W_{\mathrm{tot}}=0\).  The triangle inequality and the definition of
operator norm give (2.2).
\end{proof}

\begin{remark}[Why the identity is not vacuous]
\label{rem:v11v2-not-vacuous}
The terms \(W_a\) can be individually large and mutually cancelling.
Therefore smallness of \(P_a-P\) must be measured in an operator topology
that acts on the actual \(W_a\), or directly by summability of
\(\|(P_a-P)W_a\|\).  Strong convergence of \(P_a-P\) on unrelated test
vectors does not control (2.1).
\end{remark}

\subsection{Transported spectral splittings}

Often the \(a\)-th contribution begins in a separate space
\(C_a^1\).  Let \(I_a:C_a^1\to C^1\) be the measured intertwiner carrying
the determinant, coarea, and orientation factors required to put it in
the common coefficient space.  Write \(W_a=I_aw_a\).  Let
\(p_a:C_a^1\to C_a^1\) be its native spectral split and let
\(P:C^1\to C^1\) be the chosen common split.

\begin{proposition}[Intertwiner commutator form]
\label{prop:v11v2-intertwiner}
The pre-sum native split differs from the common post-transport split by
\[
 \sum_a I_ap_aw_a-PW_{\mathrm{tot}}
 =
 \sum_a\bigl(I_ap_a-PI_a\bigr)w_a.
 \tag{2.3}
\]
If \(I_a\) is invertible onto its image and
\(P_a=I_ap_aI_a^{-1}\) there, then the right-hand side of (2.3) is the
defect in (2.1).
\end{proposition}

\begin{proof}
Since \(W_{\mathrm{tot}}=\sum_aI_aw_a\),
\[
 \sum_a I_ap_aw_a-PW_{\mathrm{tot}}
 =
 \sum_a I_ap_aw_a-\sum_aPI_aw_a,
\]
which is (2.3).  Under the stated definition of \(P_a\),
\(I_ap_aw_a=P_aW_a\), so (2.3) equals
\(\sum_a(P_a-P)W_a\).
\end{proof}

For a spectral family \(p_a=\chi(H_a)\), the term
\[
 [I_a,\chi]_{\mathrm{edge}}
 :=
 I_a\chi(H_a)-\chi(H)I_a
\]
is supported where the two functional calculi fail to intertwine.  Calling
it an edge term is descriptive; its smallness requires a resolvent or
functional-calculus estimate.

\begin{corollary}[Resolvent estimate for a smooth split]
\label{cor:v11v2-resolvent}
Suppose \(H_a\) on \(C_a^1\) and \(H\) on \(C^1\) are bounded operators,
and \(\chi\) has a holomorphic functional calculus representation
\[
 \chi(H)=\frac{1}{2\pi i}\int_\Gamma
 \chi(z)(z-H)^{-1}\,dz
\]
on a contour \(\Gamma\) enclosing both spectra.  Then
\[
 I_a\chi(H_a)-\chi(H)I_a
 =
 \frac{1}{2\pi i}\int_\Gamma
 \chi(z)(z-H)^{-1}(HI_a-I_aH_a)(z-H_a)^{-1}\,dz.
 \tag{2.4}
\]
In particular,
\[
 \|I_a\chi(H_a)-\chi(H)I_a\|
 \leq
 \frac{\operatorname{length}(\Gamma)}{2\pi}
 \sup_{z\in\Gamma}|\chi(z)|
 \sup_{z\in\Gamma}\|(z-H)^{-1}\|
 \sup_{z\in\Gamma}\|(z-H_a)^{-1}\|
 \|HI_a-I_aH_a\|.
 \tag{2.5}
\]
\end{corollary}

\begin{proof}
The resolvent identity with an intertwiner is
\[
 I_a(z-H_a)^{-1}-(z-H)^{-1}I_a
 =
 (z-H)^{-1}(HI_a-I_aH_a)(z-H_a)^{-1}.
\]
Insert it into the two contour formulas and integrate.  Taking norms,
using submultiplicativity, and bounding the contour integral by its length
gives (2.5).
\end{proof}

The sign in (2.4) follows directly by multiplying the resolvent identity
on the left by \(z-H\) and on the right by \(z-H_a\).  This formula gives
a concrete acquisition target: control the measured intertwining error
\(HI_a-I_aH_a\) and the common resolvent gap.

\subsection{Independence of a common split in Ward cohomology}

The previous statements concern exact equality.  We now record the
corresponding cohomological comparison.

\begin{theorem}[Chain-homotopic splits give the same Ward class]
\label{thm:v11v2-chain-homotopy}
Let \(P,Q:C^\bullet\to C^\bullet\) be degree-zero chain maps:
\[
 d_\nabla P=Pd_\nabla,\qquad d_\nabla Q=Qd_\nabla.
\]
Suppose there is a degree \(-1\) map \(K:C^\bullet\to C^{\bullet-1}\)
such that
\[
 P-Q=d_\nabla K+Kd_\nabla.
 \tag{2.6}
\]
If \(W_{\mathrm{tot}}\in C^1\) is closed, then
\[
 PW_{\mathrm{tot}}-QW_{\mathrm{tot}}
 =
 d_\nabla(KW_{\mathrm{tot}}),
 \tag{2.7}
\]
so the two common splits define the same class in
\(H^1(C^\bullet,d_\nabla)\).  If
\(W_{\mathrm{tot}}=d_\nabla C\), then
\[
 PW_{\mathrm{tot}}=d_\nabla(PC),
 \qquad
 QW_{\mathrm{tot}}=d_\nabla(QC).
 \tag{2.8}
\]
\end{theorem}

\begin{proof}
Apply (2.6) to \(W_{\mathrm{tot}}\).  Closedness removes
\(Kd_\nabla W_{\mathrm{tot}}\), leaving (2.7).  If
\(W_{\mathrm{tot}}=d_\nabla C\), the chain-map identities give
\(Pd_\nabla C=d_\nabla PC\) and likewise for \(Q\).
\end{proof}

\begin{firewall}[Connection compatibility]
\label{fw:v11v2-connection}
Equations (2.6)--(2.8) are global determinant-line statements only if
\(P,Q,K\), and the identifications \(I_a\) respect the transition
functions and the connection used in \(d_\nabla\).  In unrelated local
frames, differentiating a transition function creates an additional
connection term.  The theorem does not trivialize the determinant line.
\end{firewall}

\subsection{Common regulator criterion}

Let \(i\) denote the regulator.  Equip \(C_i^1\) with a seminorm
\(\|\cdot\|_{i,\mathrm{loc}}\) for a fixed local Ward test.  Suppose
\[
 W_{\mathrm{tot},i}=d_{\nabla_i}C_i+R_i.
 \tag{2.9}
\]
Let \(P_i\) be a common bounded chain map and let \(P_{a,i}\) be the
component splits.

\begin{proposition}[Vanishing-defect criterion]
\label{prop:v11v2-vanishing}
Assume, on each fixed local test,
\[
 \|P_iR_i\|_{i,\mathrm{loc}}\longrightarrow0,
 \qquad
 \sum_a\|(P_{a,i}-P_i)W_{a,i}\|_{i,\mathrm{loc}}
 \longrightarrow0.
 \tag{2.10}
\]
Then
\[
 \sum_aP_{a,i}W_{a,i}
 -
 d_{\nabla_i}(P_iC_i)
 \longrightarrow0
 \tag{2.11}
\]
in the same local seminorm.
\end{proposition}

\begin{proof}
The splitting identity and the chain-map property give
\[
 \sum_aP_{a,i}W_{a,i}
 =
 P_i(d_{\nabla_i}C_i+R_i)
 +\sum_a(P_{a,i}-P_i)W_{a,i}
 =
 d_{\nabla_i}(P_iC_i)+P_iR_i+D_i.
\]
Both remaining terms tend to zero by (2.10).
\end{proof}

\begin{firewall}[What remains model-specific]
\label{fw:v11v2-model}
The Standard Model charge identities summarized in V1 cancel the local
anomaly polynomial after all left-handed representations are put in a
common convention.  They do not establish (2.9)--(2.10).  A physical
application still requires one measured intertwiner system, a uniform
spectral gap or resolvent estimate, inclusion of the ghost, coarea,
boundary, and counterterm rows, and determinant-line compatibility.
\end{firewall}

\begin{assessment}[V2 advance]
\label{ass:v11v2}
The false-anomaly mechanism is now exact.  A componentwise spectral Ward
sum equals the common complete response plus the explicit intertwiner edge
defect (2.3).  Common chain-homotopic splits preserve the Ward cohomology
class.  The next bottleneck is analytic: construct the common overlap
functional calculus on an admissible physical set and prove a local norm
bound for \(HI_a-I_aH_a\), with the bad set controlled separately.
\end{assessment}


\section{Admissible-set closure of the Ward edge defect}
\label{sec:v11v3}

V2 reduced the componentwise spectral error to measured intertwiner
commutators.  This note supplies a finite-cutoff analytic closure theorem.
It separates a uniformly gapped admissible set, where resolvent calculus
controls the commutator, from its complement, where probability and
integrability must be estimated.

\subsection{A local operator algebra}

Let \(\Lambda_i\) be a finite regulator complex with distance \(d_i\).
For a block kernel \(A=(A_{xy})_{x,y\in\Lambda_i}\) and \(\mu>0\), set
\[
 \|A\|_{\mu,i}
 :=
 \max\left\{
 \sup_x\sum_y e^{\mu d_i(x,y)}\|A_{xy}\|,
 \sup_y\sum_x e^{\mu d_i(x,y)}\|A_{xy}\|
 \right\}.
 \tag{3.1}
\]
The block norm is any fixed submultiplicative finite-dimensional norm.

\begin{lemma}[Weighted-kernel algebra]
\label{lem:v11v3-weighted-algebra}
For compatible kernels,
\[
 \|AB\|_{\mu,i}\leq\|A\|_{\mu,i}\|B\|_{\mu,i}.
 \tag{3.2}
\]
Moreover,
\[
 \|A_{xy}\|\leq e^{-\mu d_i(x,y)}\|A\|_{\mu,i}.
 \tag{3.3}
\]
Both constants are independent of \(|\Lambda_i|\).
\end{lemma}

\begin{proof}
The triangle inequality and
\(d_i(x,z)\leq d_i(x,y)+d_i(y,z)\) give
\[
 \sum_z e^{\mu d_i(x,z)}\|(AB)_{xz}\|
 \leq
 \sum_y e^{\mu d_i(x,y)}\|A_{xy}\|
 \sum_z e^{\mu d_i(y,z)}\|B_{yz}\|.
\]
Taking the row supremum gives the first row bound; the column bound is
identical with the order reversed.  This proves (3.2).  A single summand in
the row sum in (3.1) gives (3.3).
\end{proof}

Thus a uniform bound in (3.1) is simultaneously a volume-independent
operator-algebra bound and an exponential locality bound.

\subsection{Riesz splits and measured intertwiners}

For each component \(a\in\mathscr A\), let \(H_{a,i}\) be its
self-adjoint kernel, let \(H_i\) be a common reference kernel, and let
\(I_{a,i}\) be the measured intertwiner into the common coefficient
space.  The word measured means that the determinant-line, coarea,
orientation, and transport factors have already been included.

Let \(\Omega_i^{\rm good}\) be a measurable admissible set.  Assume that
on this set
\[
 \spec(H_i)\cup\spec(H_{a,i})
 \subset [-M,-\gamma]\cup[\gamma,M]
 \tag{3.4}
\]
with \(M<\infty\) and \(\gamma>0\) independent of \(i\) and \(a\).
Choose disjoint positively oriented contours
\(\Gamma_-\) and \(\Gamma_+\) enclosing the negative and positive
intervals, respectively, at positive distance from them and from zero.
Define the negative Riesz splits
\[
 P_i^-=\frac{1}{2\pi i}\int_{\Gamma_-}(z-H_i)^{-1}\,dz,
 \qquad
 p_{a,i}^-=\frac{1}{2\pi i}
 \int_{\Gamma_-}(z-H_{a,i})^{-1}\,dz.
 \tag{3.5}
\]

\begin{theorem}[Uniform local intertwiner bound]
\label{thm:v11v3-local-intertwiner}
Suppose that on \(\Omega_i^{\rm good}\)
\[
 \sup_{z\in\Gamma_-}
 \left\{
 \|(z-H_i)^{-1}\|_{\mu,i},
 \|(z-H_{a,i})^{-1}\|_{\mu,i}
 \right\}
 \leq R
 \tag{3.6}
\]
and
\[
 \|H_iI_{a,i}-I_{a,i}H_{a,i}\|_{\mu,i}
 \leq\delta_{a,i}.
 \tag{3.7}
\]
Then
\[
 \|I_{a,i}p_{a,i}^- -P_i^-I_{a,i}\|_{\mu,i}
 \leq
 C_\Gamma R^2\delta_{a,i},
 \qquad
 C_\Gamma:=\frac{\operatorname{length}(\Gamma_-)}{2\pi}.
 \tag{3.8}
\]
Consequently its kernel satisfies
\[
 \|(I_{a,i}p_{a,i}^- -P_i^-I_{a,i})_{xy}\|
 \leq
 C_\Gamma R^2\delta_{a,i}e^{-\mu d_i(x,y)}.
 \tag{3.9}
\]
\end{theorem}

\begin{proof}
The intertwining resolvent identity gives
\[
 I_{a,i}(z-H_{a,i})^{-1}-(z-H_i)^{-1}I_{a,i}
 =
 (z-H_i)^{-1}
 (H_iI_{a,i}-I_{a,i}H_{a,i})
 (z-H_{a,i})^{-1}.
\]
Integrate around \(\Gamma_-\).  The weighted-kernel algebra lemma,
(3.6), and (3.7) give (3.8).  The pointwise consequence is (3.3).
\end{proof}

\begin{remark}[Where the hard estimate sits]
\label{rem:v11v3-hard-estimate}
A Hilbert-space gap gives a uniform unweighted resolvent norm on a contour,
but it does not automatically give (3.6).  The latter is the
Combes--Thomas or overlap-locality acquisition.  It must be proved from
finite range, admissibility, and the regulator geometry.
\end{remark}

\subsection{From the operator bound to a local Ward bound}

Let \(w_{a,i}\) denote the native component response and let
\(\ell_i\) be a local continuous functional on the common coefficient
space.  Define
\[
 \mathcal D_i
 :=
 \ell_i\left(
 \sum_a
 (I_{a,i}p_{a,i}^- -P_i^-I_{a,i})w_{a,i}
 \right).
 \tag{3.10}
\]
Assume that \(\ell_i\) is supported in a set of uniformly bounded diameter
and that its norm relative to (3.1) is at most \(L\).

\begin{corollary}[Good-set Ward estimate]
\label{cor:v11v3-good-set}
On \(\Omega_i^{\rm good}\),
\[
 |\mathcal D_i|
 \leq
 LC_\Gamma R^2
 \sum_a\delta_{a,i}\|w_{a,i}\|_{\mu,i}.
 \tag{3.11}
\]
In particular, if the right-hand side converges to zero uniformly on
\(\Omega_i^{\rm good}\), then the artificial edge defect vanishes
uniformly there.
\end{corollary}

\begin{proof}
Apply the triangle inequality, continuity of \(\ell_i\), submultiplicativity
of the weighted norm, and (3.8) term by term.
\end{proof}

The same proof applies to any holomorphic split on the two spectral
components.  A sharp indicator is allowed because the gap permits the
Riesz representation (3.5).

\subsection{Closing the bad set}

Let \(\nu_i\) be the normalized reduced bosonic measure, including its
discrete sector label.  The Ward defect is a scalar after contraction with
the local test.  Write
\[
 \varepsilon_i:=\nu_i((\Omega_i^{\rm good})^c).
\]

\begin{theorem}[Good-set/bad-set closure]
\label{thm:v11v3-good-bad}
Let \(q>1\) and \(q'=q/(q-1)\).  Suppose
\[
 \sup_i\|\mathcal D_i\|_{L^q(\nu_i)}\leq B
 \tag{3.12}
\]
and
\[
 \sup_{\Omega_i^{\rm good}}|\mathcal D_i|\leq\eta_i,
 \qquad
 \eta_i\longrightarrow0,
 \qquad
 \varepsilon_i\longrightarrow0.
 \tag{3.13}
\]
Then
\[
 \left|\int\mathcal D_i\,d\nu_i\right|
 \leq
 \eta_i+B\varepsilon_i^{1/q'}
 \longrightarrow0.
 \tag{3.14}
\]
The conclusion remains true uniformly over a family of local tests when
\(B\) and \(\eta_i\) are uniform over that family.
\end{theorem}

\begin{proof}
Split the integral over the good set and its complement.  The first
absolute value is at most \(\eta_i\nu_i(\Omega_i^{\rm good})\leq\eta_i\).
H\"older's inequality bounds the second by
\[
 \|\mathcal D_i\|_{L^q(\nu_i)}
 \|\mathbf 1_{(\Omega_i^{\rm good})^c}\|_{L^{q'}(\nu_i)}
 \leq B\varepsilon_i^{1/q'}.
\]
This proves (3.14).
\end{proof}

\begin{corollary}[A usable acquisition card]
\label{cor:v11v3-acquisition-card}
The expected artificial Ward edge vanishes if:
\begin{enumerate}[label=\textnormal{(\alph*)}]
\item the admissible set has a common spectral gap and weighted resolvent
bound;
\item
\(\sup_{\Omega_i^{\rm good}}\sum_a
 \delta_{a,i}\|w_{a,i}\|_{\mu,i}\to0\);
\item its complement has probability tending to zero; and
\item the full defect has a uniform \(L^q\) bound for some \(q>1\).
\end{enumerate}
\end{corollary}

\begin{proof}
Condition (a) and the good-set Ward estimate give (3.13) with
\(\eta_i\to0\) by (b).  Conditions (c)--(d) give the remaining hypotheses
of the good-set/bad-set theorem.
\end{proof}

\subsection{Counterterm and exact-coboundary stability}

Let \(d_{\nabla_i}C_i\) be a local counterterm response.  If \(P_i^-\) is a
chain map, then
\[
 P_i^-d_{\nabla_i}C_i=d_{\nabla_i}(P_i^-C_i).
 \tag{3.15}
\]
If an adjacent measured intertwiner changes the determinant coefficient
by an endpoint coboundary
\[
 C_{i+1}-T_iC_i=d_{\nabla_i}B_i,
 \tag{3.16}
\]
then its Ward class is unchanged.  Equation (3.16) may be enough for an
averaged marginal coefficient when \(B_i\) is sublinear, but projective
coefficient convergence requires control of the actual transported
increments.

\begin{proposition}[Projective versus averaged closure]
\label{prop:v11v3-projective}
Let \(c_i\) be coefficients in normed spaces with transports \(T_i\).
If
\[
 \sum_i\|c_{i+1}-T_ic_i\|<\infty,
 \tag{3.17}
\]
then every transported tail is Cauchy and defines a projective limit
coefficient.  The weaker condition
\[
 \frac1N\left\|
 \sum_{i=1}^N(c_{i+1}-T_ic_i)
 \right\|\longrightarrow0
 \tag{3.18}
\]
does not imply (3.17) or pointwise convergence.
\end{proposition}

\begin{proof}
For \(m>n\), transport all increments to the \(m\)-th common space.
The triangle inequality bounds the transported difference between
\(c_m\) and \(c_n\) by the tail of the convergent series in (3.17), which
tends to zero.  For failure of the converse, take one common scalar space,
\(T_i=1\), and \(c_i=(-1)^i\).  The averaged telescoping increment is
\((c_{N+1}-c_1)/N\to0\), while \(c_i\) does not converge and the sum of
increment magnitudes diverges.
\end{proof}

\begin{firewall}[Present physical status]
\label{fw:v11v3-status}
No estimate above proves that the interacting reduced
Standard Model--Einstein measure concentrates on an overlap-admissible
set.  The new result states exactly what would suffice and shows that
discarding the bad set without a vanishing probability and an \(L^q\)
bound is invalid.
\end{firewall}

\begin{assessment}[V3 advance]
\label{ass:v11v3}
The V2 edge identity now has a volume-independent local closure route:
weighted resolvent locality and a measured intertwining error control the
good set, while H\"older closes the complement.  The next upstream
question is how an overlap admissibility failure can be bounded by a
physical plaquette or curvature energy under the reduced positive body,
without presupposing the continuum measure.
\end{assessment}


\section{Plaquette concentration required by overlap admissibility}
\label{sec:v11v4}

The good-set theorem in V3 is useful only if the reduced positive body
concentrates on a set with a uniform overlap gap.  This note derives a
sharp elementary route from plaquette insertion tails to global
admissibility and records why an extensive curvature-energy bound alone
cannot provide it.

\subsection{Defects and the global event}

Let \(\mathcal P_i\) be the plaquettes of the finite regulator and
\(N_i:=|\mathcal P_i|\).  For a compact unitary gauge representation set
\[
 r_p(U):=\|\mathbf 1-U_p\|,\qquad 0\leq r_p\leq r_{\max}.
 \tag{4.1}
\]
For a threshold \(\epsilon_i>0\), define
\[
 \Omega_i^{\rm adm}(\epsilon_i)
 :=
 \{U:\max_{p\in\mathcal P_i}r_p(U)\leq\epsilon_i\},
 \qquad
 B_i(\epsilon_i):=(\Omega_i^{\rm adm}(\epsilon_i))^c.
 \tag{4.2}
\]
The overlap gap theorem, when available for the chosen representation and
kernel, supplies a fixed \(\epsilon_*>0\) such that
\(\Omega_i^{\rm adm}(\epsilon_*)\) is contained in the gapped set used in
V3.  That deterministic implication is kept as a separate input.

\begin{lemma}[Union conversion]
\label{lem:v11v4-union}
Let \(\mu_i\) be any probability measure on the gauge configurations.
If for some \(\lambda_i>0\) and \(C_i\geq1\)
\[
 \sup_{p\in\mathcal P_i}
 \int e^{\lambda_i r_p(U)^2}\,d\mu_i(U)
 \leq C_i,
 \tag{4.3}
\]
then
\[
 \mu_i(B_i(\epsilon_i))
 \leq
 N_iC_i e^{-\lambda_i\epsilon_i^2}.
 \tag{4.4}
\]
Consequently \(\mu_i(B_i(\epsilon_i))\to0\) whenever
\[
 \lambda_i\epsilon_i^2-\log N_i-\log C_i
 \longrightarrow+\infty.
 \tag{4.5}
\]
\end{lemma}

\begin{proof}
For a fixed plaquette, exponential Markov gives
\[
 \mu_i\{r_p>\epsilon_i\}
 \leq
 e^{-\lambda_i\epsilon_i^2}
 \int e^{\lambda_i r_p^2}\,d\mu_i
 \leq C_ie^{-\lambda_i\epsilon_i^2}.
\]
The bad event is the union of these events over \(p\).  The union bound
gives (4.4), and (4.5) makes its right-hand side tend to zero.
\end{proof}

\begin{remark}[The entropy term]
\label{rem:v11v4-entropy}
The term \(\log N_i\) is unavoidable at this level.  A tail estimate for
one plaquette must beat the number of possible locations of a defect.
Thus a statement that the bare coupling tends to infinity is insufficient
unless its rate and constants dominate the lattice entropy in (4.5).
\end{remark}

\subsection{Transfer through a reduced density}

The physical reduced measure is not normally the reference gauge measure.
Let
\[
 d\nu_i=f_i\,d\mu_i,\qquad f_i\geq0,\qquad
 \int f_i\,d\mu_i=1.
 \tag{4.6}
\]

\begin{theorem}[Admissibility transfer under an \(L^s\) density]
\label{thm:v11v4-density-transfer}
Let \(s>1\), \(s'=s/(s-1)\), and assume
\[
 \|f_i\|_{L^s(\mu_i)}\leq K_i.
 \tag{4.7}
\]
Under (4.3),
\[
 \nu_i(B_i(\epsilon_i))
 \leq
 K_i
 \left(N_iC_ie^{-\lambda_i\epsilon_i^2}\right)^{1/s'}.
 \tag{4.8}
\]
In particular, the reduced measure is asymptotically admissible if
\[
 \lambda_i\epsilon_i^2-\log N_i-\log C_i
 -s'\log K_i
 \longrightarrow+\infty.
 \tag{4.9}
\]
\end{theorem}

\begin{proof}
H\"older's inequality and (4.6) give
\[
 \nu_i(B_i)=\int\mathbf1_{B_i}f_i\,d\mu_i
 \leq
 \|f_i\|_{L^s(\mu_i)}
 \mu_i(B_i)^{1/s'}.
\]
Insert (4.4) to obtain (4.8).  Taking logarithms of its right-hand side
gives (4.9).
\end{proof}

This theorem includes determinant, Higgs, fermion, constraint, and coarea
effects only through the normalized density \(f_i\).  Proving (4.7) with
the required rate is one of the adjacent-density acquisitions; it cannot
be inferred from formal cancellation of determinants.

\subsection{Sectorwise version}

Suppose
\[
 \nu_i=\sum_{\kappa\in\mathcal K_i}\pi_i(\kappa)\nu_{i,\kappa},
 \qquad \pi_i(\kappa)\geq0,\qquad\sum_\kappa\pi_i(\kappa)=1.
 \tag{4.10}
\]
Let \(G_i\subset\mathcal K_i\) be retained sectors.

\begin{proposition}[Admissibility with sector tails]
\label{prop:v11v4-sectors}
If
\[
 \sum_{\kappa\notin G_i}\pi_i(\kappa)\leq\tau_i
 \tag{4.11}
\]
and, uniformly for \(\kappa\in G_i\),
\[
 \nu_{i,\kappa}(B_i(\epsilon_i))\leq b_i,
 \tag{4.12}
\]
then
\[
 \nu_i(B_i(\epsilon_i))\leq\tau_i+b_i.
 \tag{4.13}
\]
Thus \(\tau_i\to0\) and \(b_i\to0\) close the sector-summed bad set.
\end{proposition}

\begin{proof}
Split (4.10) over \(G_i\) and its complement.  The first contribution is
at most \(b_i\sum_{\kappa\in G_i}\pi_i(\kappa)\leq b_i\); the second is at
most (4.11).
\end{proof}

\subsection{Why total energy is insufficient}

Define the \(q\)-energy
\[
 \mathcal E_i^{(q)}(U):=\sum_{p\in\mathcal P_i}r_p(U)^q,
 \qquad q>0.
 \tag{4.14}
\]
Markov's inequality gives only
\[
 \nu_i(B_i(\epsilon_i))
 \leq
 \epsilon_i^{-q}\,
 \mathbb E_{\nu_i}\mathcal E_i^{(q)}.
 \tag{4.15}
\]
For fixed \(\epsilon_i\), an \(O(1)\) bound on the right-hand energy need
not tend to zero.

\begin{proposition}[One-defect obstruction]
\label{prop:v11v4-one-defect}
Assume that for every \(i\) there exists a configuration having exactly one
plaquette defect of size \(r_0>\epsilon>0\) and all other defects zero.
There are probability measures \(\rho_i\) such that
\[
 \mathbb E_{\rho_i}\mathcal E_i^{(q)}=r_0^q
 \tag{4.16}
\]
for every \(i\), while
\[
 \rho_i(B_i(\epsilon))=1.
 \tag{4.17}
\]
Hence uniform boundedness of total \(q\)-energy does not imply global
admissibility.
\end{proposition}

\begin{proof}
Let \(\rho_i\) be the point mass at any one-defect configuration, or the
uniform mixture over all its plaquette translates.  Every configuration
in the support has energy \(r_0^q\) and maximum defect \(r_0>\epsilon\).
Equations (4.16)--(4.17) follow.
\end{proof}

Even a translation-invariant mixture in the proof has the obstruction.
The problem is the maximum over the growing lattice, not the choice of a
distinguished plaquette.

\begin{corollary}[Moment criterion and its scaling]
\label{cor:v11v4-moment}
If
\[
 \sup_{p\in\mathcal P_i}\mathbb E_{\nu_i}r_p^q\leq m_i,
 \tag{4.18}
\]
then
\[
 \nu_i(B_i(\epsilon_i))\leq
 N_i m_i\epsilon_i^{-q}.
 \tag{4.19}
\]
Therefore the elementary moment route requires
\(N_im_i\epsilon_i^{-q}\to0\), a much stronger assertion than bounded
expected total energy.
\end{corollary}

\begin{proof}
Apply the polynomial Markov inequality to each plaquette and take a union
bound.
\end{proof}

\subsection{Combined overlap-Ward card}

\begin{theorem}[Admissibility-to-Ward closure]
\label{thm:v11v4-overlap-ward}
Assume:
\begin{enumerate}[label=\textnormal{(\roman*)}]
\item the deterministic implication
\(\Omega_i^{\rm adm}(\epsilon_*)\subset\Omega_i^{\rm good}\) holds with a
uniform overlap gap;
\item the weighted resolvent and measured-intertwiner estimates of V3 hold
on \(\Omega_i^{\rm good}\), with the good-set Ward bound
\(\eta_i\to0\);
\item either (4.5) holds directly for \(\nu_i\), or (4.9) holds through a
reference measure and reduced density; and
\item the scalar local edge defect has a uniform \(L^q(\nu_i)\) bound for
some \(q>1\).
\end{enumerate}
Then the expected artificial componentwise Ward edge tends to zero.
\end{theorem}

\begin{proof}
Assumptions (i) and (iii) imply
\(\nu_i((\Omega_i^{\rm good})^c)\to0\).  Assumption (ii) supplies the
uniform good-set bound.  Assumption (iv) closes the complement by the
good-set/bad-set theorem of V3.
\end{proof}

\begin{firewall}[No hidden concentration claim]
\label{fw:v11v4-concentration}
The theorem does not assert (4.3), (4.7), or (4.9) for the coupled
Standard Model--Einstein density.  In particular, a formal Wilson-action
penalty or an expected \(L^2\) curvature bound cannot replace the
plaquette tail rate.  If global admissibility is too strong at the physical
scaling, the programme must move upstream to a mobility-gap or
localized-defect theorem rather than discard the bad configurations.
\end{firewall}

\begin{assessment}[V4 advance]
\label{ass:v11v4}
The admissibility bottleneck is now quantitative.  A single-plaquette
exponential insertion bound survives the reduced density with an explicit
\(L^s\) cost, and it must beat \(\log N_i\).  The one-defect example proves
that total curvature energy alone cannot do this.  The next version is the
mandatory companion audit; after it, the programme will choose between
global concentration and a localized-defect or mobility-gap route.
\end{assessment}


\section{V5 mandatory companion audit}
\label{sec:v11v5}

The five-version rule requires a fresh inspection of the active
Yang--Mills and Navier--Stokes notes before the next mathematical step.

\begin{audit}[Files inspected]
\label{audit:v11v5-files}
The exact files were
\[
\begin{array}{c|r|r}
\text{file}&\text{lines}&\text{bytes}\\ \hline
\texttt{Renewal\_Yang\_Mills\_Mass\_Gap\_Research\_Notes\_V20.tex}
&9357&317451\\
\texttt{Renewal\_NS\_Stability\_Research\_Notes\_V26.tex}
&5319&278283.
\end{array}
\]
Their SHA--256 values are, respectively,
\begin{center}
\small\ttfamily
842D84EF6BE00B9C2DE525EF5CB389D0FF3BBAA46A8854F4D0CBFD596A9A021A
\end{center}
and
\begin{center}
\small\ttfamily
82C887F8C2C69E4F28FAD7C3DC8DE5B9099CB5A4BF431EBA1FFF3E91935978AD.
\end{center}
The NS file is unchanged from the previous audit.  The YM stream has
advanced from V17 to V20.
\end{audit}

\subsection{New Yang--Mills material}

YM V18 converts an explicit real spectral floor for a finite Laurent
symbol into a nonzero complex momentum strip by a coefficient-sum
perturbation and a Neumann series.  Cauchy estimates then turn the strip
into a finite analytic response stock and an exponentially convergent
grid enclosure.

YM V19 constructs a nonnegative gauge-invariant finite-range stabilizer
from the covariant Laplacian of the plaquette defect.  At the flat
background its reduced symbol has the factorized form
\(\omega(k)^2A(k)\).  It proves positivity of a convex comparison homotopy
and persistence of the gap on a small background tube.

YM V20 replaces an unspecified endpoint by the explicit
plaquette-difference functional
\[
 \mathcal O_1(U)
 =
 \frac{\nu_1}{2}
 \sum_{x,\mu<\nu,\rho}
 \|\mathcal D_\rho P_{\mu\nu}(x)\|_{\rm HS}^2.
 \tag{5.1}
\]
Its flat Hessian factorizes positively, the complete link Hessian
annihilates gauge gradients, and exact stationary-center differentiation
uses action derivatives only through total order four.  Haar density is
absorbed into the measured quadratic vertex.  That note also reaffirms
that averaged coboundary control and projective cylinder convergence are
different acquisitions.

\subsection{Transfer decisions}

\begin{audit}[Accepted transfers]
\label{audit:v11v5-accepted}
Three structural points transfer.
\begin{enumerate}[label=\textnormal{(\roman*)}]
\item A real gap plus a literal finite-range coefficient stock can be
continued to a complex strip and hence to an exponentially weighted
kernel bound.
\item Gap persistence should be proved on background tubes by an explicit
resolvent or Neumann estimate.  For the SM programme these tubes must be
indexed by the chiral/topological sector rather than restricted to one
near-flat chart.
\item Every stationary reduction must carry its Haar, coarea, determinant,
and center-response terms through the comparison.  The exact endpoint
logarithmic determinant identity is valid only for a common positive
homotopy with those density terms retained.
\end{enumerate}
\end{audit}

\begin{audit}[Rejected direct transfers]
\label{audit:v11v5-rejected}
The YM constants, parity-tree fibre, and bosonic vertical Hessian do not
apply to the Standard Model chiral kernel.  Positivity of a gauge-boson
Hessian does not imply a gap for the Hermitian Wilson--Dirac kernel.
The YM edge-mean computation is also a bosonic comparison problem and
does not construct the SM determinant-line phase.
\end{audit}

\begin{proposition}[A derivative stabilizer does not control the amplitude
mode]
\label{prop:v11v5-derivative-mode}
Let a gauge configuration have a covariantly constant plaquette defect:
\[
 \mathcal D_\rho(P_{\mu\nu}-I)=0
 \quad\text{for all }\rho,\mu,\nu.
 \tag{5.2}
\]
Then the functional (5.1) vanishes even when
\(\|P_{\mu\nu}-I\|>\epsilon\).  Likewise the covariant-Laplacian
stabilizer vanishes when
\(\mathcal L_U(P_{\mu\nu}-I)=0\).  Therefore neither derivative-only
stabilizer implies the amplitude admissibility event of V4.
\end{proposition}

\begin{proof}
Each summand in (5.1) is the squared norm of the left side of (5.2), hence
is zero.  The second assertion is identical with
\(\mathcal D_\rho\) replaced by \(\mathcal L_U\).  The value of the
undifferentiated plaquette defect is not constrained by either equality,
so it may exceed \(\epsilon\).
\end{proof}

The proposition does not make the stabilizers useless.  Combined with an
amplitude penalty, they may give regularity and better defect tails.  It
shows only that they cannot replace the missing amplitude estimate.

\subsection{Navier--Stokes audit}

NS V26 studies pointwise rank-one norms of derivatives of an ancient
kernel in a singularity-exclusion programme.  It contains no chiral
determinant line, overlap gap, gauge-covariant background tube, or
reduced-density estimate that transfers to the present finite target.
No NS constant or conclusion is imported.

\begin{decision}[Post-audit direction]
\label{dec:v11v5-direction}
The next version will replace the single global near-flat event by a
sectorwise atlas of gapped background tubes.  It will prove:
\begin{enumerate}[label=\textnormal{(\alph*)}]
\item persistence of a weighted resolvent and Riesz projector inside each
tube from a base-background stock;
\item constancy of the chiral index on each tube;
\item a tube-tail criterion that closes the expected Ward edge without
requiring every retained sector to be near the identity.
\end{enumerate}
The remaining physical acquisition will be the construction and
concentration of such an atlas for the coupled reduced measure.
\end{decision}

\begin{assessment}[V5 audit result]
\label{ass:v11v5}
The audit changes the attack from a single global admissibility estimate
to a gapped sector-background atlas.  This uses the transferable analytic
method from YM V18--V20 while preserving the distinction between a
bosonic Hessian gap and the chiral overlap gap.  The full continuum remains
open.
\end{assessment}


\section{Sectorwise gapped background-tube atlas}
\label{sec:v11v6}

The V5 audit suggests replacing one global near-flat chart by tubes around
sector-dependent backgrounds.  This note proves the analytic atlas
theorem.  It does not assume that every topological sector contains a flat
representative.

\subsection{A base strip from finite-range data}

Let \(A(k)\) be an \(m\times m\) matrix Laurent polynomial on
\(\mathbb T^d\),
\[
 A(k)=\sum_{n}A_ne^{in\cdot k},
 \qquad
 M_A:=\sum_n\|A_n\|,
 \qquad
 A_n=0\quad\text{if }|n|_1>R_A.
 \tag{6.1}
\]
Assume
\[
 \inf_{k\in\mathbb T^d}s_{\min}(A(k))\geq\gamma>0.
 \tag{6.2}
\]

\begin{proposition}[Real gap to a weighted inverse stock]
\label{prop:v11v6-strip-stock}
Set
\[
 \rho:=
 \frac1{R_A}\log\left(1+\frac{\gamma}{2M_A}\right)
 \tag{6.3}
\]
when \(R_A,M_A>0\).  Then \(A(k+iy)\) is invertible for
\(\|y\|_\infty\leq\rho\) and
\[
 \sup_{\|y\|_\infty\leq\rho}
 \|A(k+iy)^{-1}\|\leq\frac2\gamma.
 \tag{6.4}
\]
If \(G_n\) are the Fourier coefficients of \(A^{-1}\), then
\[
 \|G_n\|\leq\frac2\gamma e^{-\rho|n|_1}.
 \tag{6.5}
\]
Consequently, for \(0<\mu<\rho\),
\[
 \sum_{n\in\mathbb Z^d}e^{\mu|n|_1}\|G_n\|
 \leq
 \frac2\gamma
 \left(
 \frac{1+e^{-(\rho-\mu)}}{1-e^{-(\rho-\mu)}}
 \right)^d.
 \tag{6.6}
\]
The same bound holds for the periodized inverse kernels on finite tori,
using the periodic distance.
\end{proposition}

\begin{proof}
For \(\|y\|_\infty\leq\rho\),
\[
 \|A(k+iy)-A(k)\|
 \leq
 \sum_n\|A_n\|\,|e^{-n\cdot y}-1|
 \leq M_A(e^{R_A\rho}-1)=\frac\gamma2.
\]
Since \(\|A(k)^{-1}\|\leq\gamma^{-1}\), the inverse-relative
perturbation has norm at most \(1/2\).  The Neumann series gives (6.4).

For a fixed \(n\), shift the \(j\)-th Fourier contour by
\(-i\rho\sgn(n_j)\).  Analyticity and periodicity give (6.5).
Multiplying by \(e^{\mu|n|_1}\) and summing the product geometric series
gives (6.6).  Periodization replaces \(G_n\) by
\(\sum_{\ell\in\mathbb Z^d}G_{n+\ell L}\).  Assigning each translate to
its nearest periodic representative and applying the same exponential
sum cannot exceed the right-hand side of (6.6).
\end{proof}

This proposition is the transferable part of the YM complex-strip method.
The constants must be recomputed from the Standard Model kernel at each
chosen background.

\subsection{Perturbation inside a chord tube}

Fix a regulator \(i\) and a background label \(b\).  Let \(B_{b,i}\) be
the background and \(H_{b,i}:=H_i(B_{b,i})\) the corresponding Hermitian
chiral kernel.  Let \(\Gamma_{b,i}\) be a contour separating its negative
and positive spectra.  In the weighted kernel algebra of V3 assume
\[
 \sup_{z\in\Gamma_{b,i}}
 \|(z-H_{b,i})^{-1}\|_{\mu,i}\leq R_{b,i}.
 \tag{6.7}
\]
Let \(d_{\rm ch}\) be a chord-chart metric and assume
\[
 \|H_i(U)-H_{b,i}\|_{\mu,i}
 \leq L_{b,i}d_{\rm ch}(U,B_{b,i}).
 \tag{6.8}
\]
Choose
\[
 0<r_{b,i}<\frac1{2R_{b,i}L_{b,i}}
 \tag{6.9}
\]
when \(L_{b,i}>0\), and let
\[
 \mathcal T_{b,i}:=
 \{U:d_{\rm ch}(U,B_{b,i})<r_{b,i}\}.
 \tag{6.10}
\]

\begin{theorem}[Weighted gap persistence on a background tube]
\label{thm:v11v6-tube-gap}
For \(U\in\mathcal T_{b,i}\) and \(z\in\Gamma_{b,i}\),
\[
 \|(z-H_i(U))^{-1}\|_{\mu,i}
 \leq
 \frac{R_{b,i}}
 {1-R_{b,i}L_{b,i}d_{\rm ch}(U,B_{b,i})}
 \leq2R_{b,i}.
 \tag{6.11}
\]
If
\[
 P_{b,i}^-:=
 \frac1{2\pi i}\int_{\Gamma_{b,i}}(z-H_{b,i})^{-1}\,dz,
 \qquad
 P_i^-(U):=
 \frac1{2\pi i}\int_{\Gamma_{b,i}}(z-H_i(U))^{-1}\,dz,
 \tag{6.12}
\]
then
\[
 \|P_i^-(U)-P_{b,i}^-\|_{\mu,i}
 \leq
 2C_{b,i}R_{b,i}^2L_{b,i}
 d_{\rm ch}(U,B_{b,i}),
 \tag{6.13}
\]
where
\(C_{b,i}:=\operatorname{length}(\Gamma_{b,i})/(2\pi)\).
\end{theorem}

\begin{proof}
Write \(\Delta=H_i(U)-H_{b,i}\).  Factor
\[
 z-H_i(U)
 =
 (z-H_{b,i})
 \left[I-(z-H_{b,i})^{-1}\Delta\right].
\]
The weighted algebra inequality and (6.7)--(6.9) show that the second
factor differs from the identity by norm less than \(1/2\).  Its Neumann
series gives the first bound in (6.11), and (6.9) gives the second.

The resolvent identity yields
\[
 (z-H_i(U))^{-1}-(z-H_{b,i})^{-1}
 =
 (z-H_i(U))^{-1}\Delta(z-H_{b,i})^{-1}.
\]
Integrate on \(\Gamma_{b,i}\), then use (6.7), (6.8), and (6.11).
This gives (6.13).
\end{proof}

\begin{corollary}[Constancy of the chiral index on a tube]
\label{cor:v11v6-index}
Assume every \(U\in\mathcal T_{b,i}\) is joined to \(B_{b,i}\) by a
continuous path lying in the tube.  Then the rank of \(P_i^-(U)\), and
hence every overlap index expressed as a fixed affine function of that
rank, is constant on \(\mathcal T_{b,i}\).
\end{corollary}

\begin{proof}
The theorem keeps \(\Gamma_{b,i}\) in the resolvent set along the path, so
the Riesz projector varies continuously.  Its rank is integer valued and
continuous, hence constant.
\end{proof}

\begin{firewall}[Local line triviality]
\label{fw:v11v6-local-line}
A contractible background tube permits a local determinant-line frame.
It does not prove that frames on different tubes have a trivial transition
cocycle, nor that the connection has trivial global holonomy.  The tube
atlas is an analytic gap device, not a solution of acquisition A4.
\end{firewall}

\subsection{The tube atlas and its tail}

Let \(\mathcal B_i\) be a finite or countable set of background labels.
Choose measurable, pairwise disjoint cells
\[
 \mathcal C_{b,i}\subset\mathcal T_{b,i}
 \qquad (b\in\mathcal B_i).
 \tag{6.14}
\]
Disjointness can be obtained by fixing an order and assigning each
configuration to its first containing tube.  Let \(\nu_i\) be the reduced
positive body and put
\[
 \pi_{b,i}:=\nu_i(\mathcal C_{b,i}).
 \tag{6.15}
\]
For a retained set \(G_i\subset\mathcal B_i\), define
\[
 \tau_i:=
 \nu_i\left(
 \left(\bigcup_{b\in G_i}\mathcal C_{b,i}\right)^c
 \right).
 \tag{6.16}
\]

Let \(\mathcal D_i(U)\) be the scalar artificial Ward edge after contraction
with a fixed local test.  On the \(b\)-th tube, V2--V3 and (6.11) give a
bound of the form
\[
 |\mathcal D_i(U)|\leq\eta_{b,i},
 \qquad U\in\mathcal C_{b,i},
 \tag{6.17}
\]
provided the measured intertwining errors are bounded there.

\begin{theorem}[Tube-atlas closure of the Ward edge]
\label{thm:v11v6-atlas-closure}
Let \(q>1\), \(q'=q/(q-1)\), and assume
\[
 \sup_i\|\mathcal D_i\|_{L^q(\nu_i)}\leq B.
 \tag{6.18}
\]
Then
\[
 \int|\mathcal D_i|\,d\nu_i
 \leq
 \sum_{b\in G_i}\pi_{b,i}\eta_{b,i}
 +B\tau_i^{1/q'}.
 \tag{6.19}
\]
In particular, the expected artificial Ward edge vanishes if
\[
 \sum_{b\in G_i}\pi_{b,i}\eta_{b,i}\longrightarrow0,
 \qquad
 \tau_i\longrightarrow0.
 \tag{6.20}
\]
\end{theorem}

\begin{proof}
Split the integral over the disjoint retained cells and their complement.
Equation (6.17) bounds the first part by the sum in (6.19).
H\"older's inequality bounds the complement by
\[
 \|\mathcal D_i\|_{L^q(\nu_i)}
 \nu_i\left(
 \left(\bigcup_{b\in G_i}\mathcal C_{b,i}\right)^c
 \right)^{1/q'}
 \leq B\tau_i^{1/q'}.
\]
This proves both statements.
\end{proof}

The weighted average in (6.20) is more flexible than a uniform worst-tube
bound.  Rare retained tubes may have weaker constants if their probability
weight compensates for them.

\subsection{Sector and determinant compatibility}

Let \(\kappa(b)\) be the overlap index of the \(b\)-th background.  The
index corollary makes \(\kappa\) constant on its cell.  Grouping cells with
the same value yields
\[
 \pi_{\kappa,i}:=
 \sum_{b:\kappa(b)=\kappa}\pi_{b,i}.
 \tag{6.21}
\]
This is compatible with the sector-tightness ledger of V1.

\begin{proposition}[Refinement preserves the tail bound]
\label{prop:v11v6-refinement}
Suppose a cell is subdivided into countably many measurable subcells and
the bound \(\eta_{b,i}\) is replaced by valid subcell bounds
\(\eta_{b,j,i}\).  If
\[
 \eta_{b,i}\geq\sup_j\eta_{b,j,i},
 \tag{6.22}
\]
then the retained contribution in (6.19) cannot increase.
\end{proposition}

\begin{proof}
The subcell probabilities sum to \(\pi_{b,i}\), so
\[
 \sum_j\pi_{b,j,i}\eta_{b,j,i}
 \leq
 \eta_{b,i}\sum_j\pi_{b,j,i}
 =
 \pi_{b,i}\eta_{b,i}.
\]
\end{proof}

This permits adaptive refinement near the boundary of a gap tube.  It
does not permit cells from different determinant-line frames to be added
as scalars until their transition maps have been supplied.

\subsection{Finite acquisition card}

For each retained background tube, the following finite data suffice:
\begin{enumerate}[label=\textnormal{(T\arabic*)}]
\item a base chiral-kernel spectral contour and a weighted inverse stock;
\item a chord-chart Lipschitz constant \(L_{b,i}\);
\item measured species-to-common intertwining errors on the tube;
\item the complete line-valued Ward sum before the Riesz split;
\item compatible determinant-line transition maps on overlaps.
\end{enumerate}
At the probabilistic level one additionally needs the tube tail
\(\tau_i\to0\), the weighted mean edge bound in (6.20), and the
\(L^q\) estimate (6.18).

\begin{proposition}[Finite verification route at a periodic background]
\label{prop:v11v6-finite-route}
If a periodic background reduces the base kernel and its contour
resolvent to finite Laurent blocks satisfying (6.1)--(6.2), then
(6.3)--(6.6) provide (T1).  If the kernel is a finite word in compact
links, differentiating those words on a compact chord tube gives a finite
coefficient upper bound for (T2).
\end{proposition}

\begin{proof}
The first statement is the real-gap-to-weighted-stock proposition.
For the second, a finite word has finitely many link occurrences.
The mean-value formula expresses its difference between \(U\) and \(B\)
as the integral of its derivative along the chord path.  On a compact
tube each derivative factor has bounded norm, and summing the finitely
many differentiated occurrences gives a finite Lipschitz constant.
\end{proof}

\begin{firewall}[Unproved coverage]
\label{fw:v11v6-coverage}
No present estimate proves \(\tau_i\to0\) for a cofinal family of
Standard Model--Einstein regulators.  A tube centered in every formal
topological sector is also insufficient: the reduced measure must
concentrate in the union with quantitative weights, and the complement
must satisfy (6.18).  If this coverage fails, the next upstream option is
a mobility-gap theorem allowing localized defects inside a retained
cell.
\end{firewall}

\begin{assessment}[V6 advance]
\label{ass:v11v6}
The post-audit step is complete.  The local overlap gap, exponential
resolvent locality, index constancy, and expected Ward-edge closure have
been assembled in a sectorwise background-tube theorem.  This strictly
weakens the near-flat global-admissibility route while keeping its missing
probability estimate visible.  The next finite question is whether the
coupled reduced density admits a quantitative tube cover, or whether one
must formulate a mobility-gap replacement that tolerates sparse
plaquette defects.
\end{assessment}


\section{Mobility-gap replacement for global admissibility}
\label{sec:v11v7}

The tube-atlas theorem still asks the reduced measure to concentrate in a
union of uniformly gapped tubes.  A weaker local route allows rare
near-zero states, provided they are localized near sparse defect sets and
therefore invisible to a fixed normalized local Ward test.

\subsection{Outer and near spectral pieces}

Let \(H\) be a finite-dimensional self-adjoint kernel.  Fix
\(0<\delta<M\) and suppose no eigenvalue lies on
\(\{-\delta,0,\delta\}\).  Define
\[
 F^-(H):=\mathbf1_{(-\infty,-\delta)}(H),\qquad
 N^-(H):=\mathbf1_{(-\delta,0)}(H),\qquad
 N(H):=\mathbf1_{(-\delta,\delta)}(H).
 \tag{7.1}
\]
Then the exact negative projector is
\[
 P^-(H)=F^-(H)+N^-(H).
 \tag{7.2}
\]
The outer projector has a Riesz contour separated from the near-zero
window.  The near projector need not have small rank.

Let \(H_i\) be the common kernel and \(H_{a,i}\) the native kernels, with
measured intertwiners \(I_{a,i}\).  Abbreviate
\[
 F_i^-=F^-(H_i),\quad N_i^-=N^-(H_i),\quad
 F_{a,i}^-=F^-(H_{a,i}),\quad N_{a,i}^-=N^-(H_{a,i}).
 \tag{7.3}
\]

\begin{lemma}[Exact mobility decomposition of the edge]
\label{lem:v11v7-decomposition}
For every native response \(w_{a,i}\),
\[
\begin{split}
 I_{a,i}P^-(H_{a,i})w_{a,i}
 -P^-(H_i)I_{a,i}w_{a,i}
 ={}&
 (I_{a,i}F_{a,i}^- -F_i^-I_{a,i})w_{a,i}\\
 &+
 (I_{a,i}N_{a,i}^- -N_i^-I_{a,i})w_{a,i}.
\end{split}
\tag{7.4}
\]
\end{lemma}

\begin{proof}
Insert (7.2) for both negative projectors and collect the outer and near
terms.
\end{proof}

The first line is controlled by an outer resolvent contour.  The second is
not discarded; it is the localized near-mode row.

\subsection{Normalized local tests}

Let \(E_x\) be the projection onto the degrees of freedom based at
\(x\).  A normalized local test is represented by nonnegative weights
\(\alpha_i(x)\) satisfying
\[
 \sum_x\alpha_i(x)\leq A
 \tag{7.5}
\]
with \(A\) independent of \(i\).  Define the local seminorm
\[
 \|K\|_{\alpha_i}
 :=
 \sum_x\alpha_i(x)\|E_xK\|.
 \tag{7.6}
\]
This permits a continuum test to occupy more lattice sites, provided its
discretized weights have uniformly bounded total mass.

Let \(S_i\) be a random set of defective cells or plaquettes, embedded in
the site metric.  Put
\[
 \mathcal L_{\mu,i}(x,S_i)
 :=
 \sum_{y\in S_i}e^{-\mu d_i(x,y)}.
 \tag{7.7}
\]

\begin{definition}[Defect-localized near window]
\label{def:v11v7-localized}
The near spectral window is \((C,\mu)\)-localized at \(S_i\) if
\[
 \|E_xN(H_i)\|\leq C\mathcal L_{\mu,i}(x,S_i)
 \tag{7.8}
\]
and, for every component,
\[
 \|E_xI_{a,i}N(H_{a,i})\|
 +
 \|E_xN(H_i)I_{a,i}\|
 \leq C_a\mathcal L_{\mu,i}(x,S_i).
 \tag{7.9}
\]
The constants may depend on a fixed species but not on volume.
\end{definition}

Because \(N_i^-\) is a subprojection of \(N(H_i)\), the same estimates
hold with \(N_i^-\) in place of \(N(H_i)\).

\begin{proposition}[Local near-mode bound]
\label{prop:v11v7-near-bound}
Under (7.9),
\[
 \left\|
 I_{a,i}N_{a,i}^- -N_i^-I_{a,i}
 \right\|_{\alpha_i}
 \leq
 C_a
 \sum_x\alpha_i(x)\mathcal L_{\mu,i}(x,S_i).
 \tag{7.10}
\]
Consequently, for bounded \(w_{a,i}\),
\[
 \left\|
 (I_{a,i}N_{a,i}^- -N_i^-I_{a,i})w_{a,i}
 \right\|_{\alpha_i}
 \leq
 C_a\|w_{a,i}\|
 \sum_x\alpha_i(x)\mathcal L_{\mu,i}(x,S_i).
 \tag{7.11}
\]
\end{proposition}

\begin{proof}
Use the triangle inequality in (7.6), the fact that a subprojection has
no larger left-local norm than the containing spectral projection under
the explicit bounds (7.9), and sum against \(\alpha_i(x)\).
Multiplication by \(w_{a,i}\) gives (7.11).
\end{proof}

\subsection{Defect intensity rather than a global maximum}

Assume the metric balls have a uniform exponential volume stock
\[
 C_\mu:=
 \sup_{i,x}\sum_y e^{-\mu d_i(x,y)}<\infty.
 \tag{7.12}
\]
Let \(\nu_i\) be the reduced probability body and suppose
\[
 \sup_y\nu_i\{y\in S_i\}\leq\theta_i.
 \tag{7.13}
\]

\begin{theorem}[Local invisibility of sparse defects]
\label{thm:v11v7-sparse}
Under (7.5), (7.12), and (7.13),
\[
 \mathbb E_{\nu_i}
 \sum_x\alpha_i(x)\mathcal L_{\mu,i}(x,S_i)
 \leq
 AC_\mu\theta_i.
 \tag{7.14}
\]
Hence the expected near-mode contribution in (7.11) tends to zero when
\(\theta_i\to0\) and the response norms are uniformly bounded.
\end{theorem}

\begin{proof}
Tonelli's theorem applies because every summand is nonnegative:
\[
\begin{split}
 \mathbb E
 \sum_x\alpha_i(x)\sum_y
 \mathbf1_{\{y\in S_i\}}e^{-\mu d_i(x,y)}
 &=
 \sum_x\alpha_i(x)\sum_y
 e^{-\mu d_i(x,y)}\nu_i\{y\in S_i\}\\
 &\leq
 \theta_i\sum_x\alpha_i(x)
 \sum_y e^{-\mu d_i(x,y)}
 \leq AC_\mu\theta_i.
\end{split}
\]
Insert this bound into (7.11).
\end{proof}

This theorem has no factor \(N_i\).  It controls a normalized local test,
not the event that there is no defect anywhere.

\begin{corollary}[One-plaquette insertion tail is enough locally]
\label{cor:v11v7-insertion}
If the defect event at \(y\) implies \(r_y>\epsilon_i\) and
\[
 \sup_y\mathbb E_{\nu_i}e^{\lambda_i r_y^2}\leq C_i,
 \tag{7.15}
\]
then
\[
 \theta_i\leq C_ie^{-\lambda_i\epsilon_i^2}.
 \tag{7.16}
\]
Thus local invisibility requires only
\(\lambda_i\epsilon_i^2-\log C_i\to+\infty\), whereas global
admissibility in V4 additionally required domination of \(\log N_i\).
\end{corollary}

\begin{proof}
Apply exponential Markov to the single defect location and then use the
sparse-defect theorem.
\end{proof}

\subsection{Outer resolvent comparison}

Choose a contour \(\Gamma_i^{\rm out}\) enclosing the spectrum below
\(-\delta\) and excluding the near window.  Suppose
\[
 \sup_{z\in\Gamma_i^{\rm out}}
 \left\{
 \|(z-H_i)^{-1}\|_{\mu,i},
 \|(z-H_{a,i})^{-1}\|_{\mu,i}
 \right\}\leq R
 \tag{7.17}
\]
and
\[
 \|H_iI_{a,i}-I_{a,i}H_{a,i}\|_{\mu,i}\leq\zeta_{a,i}.
 \tag{7.18}
\]

\begin{proposition}[Outer mobility edge]
\label{prop:v11v7-outer}
With
\(C_{\rm out}:=\operatorname{length}(\Gamma_i^{\rm out})/(2\pi)\),
\[
 \|I_{a,i}F_{a,i}^- -F_i^-I_{a,i}\|_{\mu,i}
 \leq C_{\rm out}R^2\zeta_{a,i}.
 \tag{7.19}
\]
\end{proposition}

\begin{proof}
Use the intertwining resolvent identity from V2 on
\(\Gamma_i^{\rm out}\), integrate, and apply the weighted-kernel algebra
bound exactly as in V3.
\end{proof}

No spectral gap at zero is used in this outer estimate.  The contour needs
only to remain separated from the outer band and the near window.

\subsection{Mobility-gap Ward closure}

Let
\[
 \mathcal D_i^{\rm mob}
 :=
 \ell_{\alpha_i}\left[
 \sum_a
 \left(
 I_{a,i}P^-(H_{a,i})-P^-(H_i)I_{a,i}
 \right)w_{a,i}
 \right],
 \tag{7.20}
\]
where the local contraction has norm at most one relative to (7.6).

\begin{theorem}[Mobility-gap replacement theorem]
\label{thm:v11v7-mobility}
Assume:
\begin{enumerate}[label=\textnormal{(M\arabic*)}]
\item the complete line-valued Ward sum is formed before the negative
spectral split;
\item the outer resolvent and intertwining estimates
(7.17)--(7.18) hold, with
\[
 \sum_a C_{\rm out}R^2\zeta_{a,i}\|w_{a,i}\|
 \longrightarrow0;
 \tag{7.21}
\]
\item the near windows satisfy (7.9) for a random defect set with
\(\theta_i\to0\);
\item
\[
 \sup_i\sum_a C_a\|w_{a,i}\|<\infty;
 \tag{7.22}
\]
\item any residual complement on which (M2)--(M4) fail has probability
\(\tau_i\to0\), and \(\mathcal D_i^{\rm mob}\) is uniformly \(L^q\) there
for some \(q>1\).
\end{enumerate}
Then
\[
 \left|\mathbb E_{\nu_i}\mathcal D_i^{\rm mob}\right|
 \longrightarrow0.
 \tag{7.23}
\]
\end{theorem}

\begin{proof}
Use the exact decomposition (7.4).  The expected outer contribution tends
to zero by (7.19) and (7.21).  The expected near contribution is at most
\[
 AC_\mu\theta_i\sum_a C_a\|w_{a,i}\|,
\]
which tends to zero by (M3)--(M4).  H\"older bounds the residual complement
by \(B\tau_i^{1/q'}\), which also tends to zero.  Adding the three parts
proves (7.23).
\end{proof}

\begin{firewall}[The localization hypothesis is the new acquisition]
\label{fw:v11v7-localization}
The theorem does not derive (7.8)--(7.9) from the Standard Model action.
Near-zero Wilson--Dirac modes can be extended, and a positive bosonic
Hessian does not localize them.  A physical proof must establish a
fractional-moment, multiscale, geometric-resolvent, or comparably strong
mobility estimate uniformly under the reduced density and determinant
weights.
\end{firewall}

\begin{remark}[Index and near modes]
\label{rem:v11v7-index}
Unlike the gapped tube theorem, the mobility theorem allows eigenvalues to
approach or cross zero inside defect regions.  The global overlap index may
therefore change.  The near-mode row records rather than erases that
possibility.  A determinant-line construction must account for the
crossing orientation and sector transport separately.
\end{remark}

\begin{assessment}[V7 advance]
\label{ass:v11v7}
The global entropy cost \(\log N_i\) is no longer required for a normalized
local Ward test.  A vanishing one-cell defect intensity is enough if the
near spectral projector is exponentially localized at those defects and
the outer band has a weighted resolvent.  This is a strict logical
weakening of global admissibility.  The next upstream task is to derive
the localization estimate from a deterministic sparse-defect geometry or
to identify precisely which percolation and Schur-complement bounds are
missing.
\end{assessment}


\section{Correction and Schur localization outside defect islands}
\label{sec:v11v8}

This note first corrects the precise hypothesis used in the V7 near-mode
bound.  It then derives that hypothesis from a deterministic
bulk--defect decomposition.

\subsection{Correction to the V7 intertwiner localization row}

The first term in V7 equation (7.9) correctly controls
\(E_xI_{a,i}N_{a,i}^-\), because
\(N_{a,i}^-=N(H_{a,i})N_{a,i}^-\).  The second displayed term there,
\(\|E_xN(H_i)I_{a,i}\|\), does not by itself control
\(\|E_xN_i^-I_{a,i}\|\) by projection order.  The following corrected
condition replaces (7.9):
\[
 \boxed{
 \|E_xI_{a,i}N(H_{a,i})\|
 +
 \|E_xN(H_i)\|\,\|I_{a,i}\|
 \leq
 C_a\mathcal L_{\mu,i}(x,S_i).
 }
 \tag{8.1}
\]
Alternatively one may assume the desired estimate directly with
\(N_{a,i}^-\) and \(N_i^-\).

\begin{proposition}[Corrected V7 near-mode bound]
\label{prop:v11v8-correction}
Under (8.1),
\[
 \|E_x(I_{a,i}N_{a,i}^- -N_i^-I_{a,i})\|
 \leq
 C_a\mathcal L_{\mu,i}(x,S_i).
 \tag{8.2}
\]
Therefore V7 Proposition 7.3 and Theorem 7.7 remain valid with (8.1) in
place of V7 equation (7.9).
\end{proposition}

\begin{proof}
For the native term,
\[
 E_xI_{a,i}N_{a,i}^-
 =
 E_xI_{a,i}N(H_{a,i})N_{a,i}^-,
\]
so its norm is bounded by the first term in (8.1).  For the common term,
\(N_i^-=N(H_i)N_i^-\), and hence
\[
 \|E_xN_i^-I_{a,i}\|
 \leq
 \|E_xN(H_i)\|\,\|N_i^-I_{a,i}\|
 \leq
 \|E_xN(H_i)\|\,\|I_{a,i}\|.
\]
The triangle inequality gives (8.2).  The later V7 proofs use only (8.2),
so no other change is required.
\end{proof}

\subsection{Bulk--defect decomposition}

Let \(\mathcal H_i\) be the finite one-particle space.  For a defect
neighbourhood \(D_i\), let \(Q_i\) be the projection onto degrees of
freedom in \(D_i\) and \(P_i=1-Q_i\).  Write the Hermitian kernel as
\[
 H_i=
 \begin{pmatrix}
 A_i&C_i\\
 C_i^*&D_i^{\rm op}
 \end{pmatrix}
 \quad\text{on}\quad
 P_i\mathcal H_i\oplus Q_i\mathcal H_i,
 \tag{8.3}
\]
where \(A_i=P_iH_iP_i\) and \(C_i=P_iH_iQ_i\).
The notation \(D_i^{\rm op}\) distinguishes the defect block from its
geometric support.

Assume that for \(|\lambda|<\delta\), the Dirichlet bulk block satisfies
\[
 A_i-\lambda\ \text{is invertible}
 \tag{8.4}
\]
and the propagated boundary coupling has the pointwise bound
\[
 \|E_x(A_i-\lambda)^{-1}C_iE_y\|
 \leq
 K e^{-\mu d_i(x,y)}
 \quad
 (x\notin D_i,\ y\in D_i).
 \tag{8.5}
\]
Let
\[
 r_*:=\sup_{i,y}\operatorname{rank}(E_yQ_i)<\infty.
 \tag{8.6}
\]

\begin{theorem}[Schur localization of near eigenmodes]
\label{thm:v11v8-schur-localization}
Let
\[
 N_i:=\mathbf1_{(-\delta,\delta)}(H_i).
\]
Under (8.4)--(8.6),
\[
 \|E_xN_i\|
 \leq
 K\sqrt{r_*}
 \sum_{y\in D_i}e^{-\mu d_i(x,y)}
 \quad (x\notin D_i).
 \tag{8.7}
\]
For \(x\in D_i\), the same conclusion holds after increasing the constant
so that the right-hand side is at least one.  Thus the V7 localization
condition follows deterministically from a gapped Dirichlet bulk and the
propagated coupling estimate (8.5).
\end{theorem}

\begin{proof}
Choose an orthonormal eigenbasis
\(\{\psi_j\}_{j\in J}\) for \(\operatorname{ran}N_i\), with
\(H_i\psi_j=\lambda_j\psi_j\) and \(|\lambda_j|<\delta\).
Writing \(\psi_j=u_j+v_j\) with
\(u_j=P_i\psi_j\) and \(v_j=Q_i\psi_j\), the bulk row of the eigenvalue
equation is
\[
 (A_i-\lambda_j)u_j=-C_iv_j.
\]
By (8.4),
\[
 u_j=-(A_i-\lambda_j)^{-1}C_iv_j.
 \tag{8.8}
\]
For \(x\notin D_i\), (8.5) gives
\[
 \|E_x\psi_j\|
 \leq
 K\sum_{y\in D_i}e^{-\mu d_i(x,y)}\|E_yv_j\|.
 \tag{8.9}
\]
Put \(a_y=e^{-\mu d_i(x,y)}\).  Weighted Cauchy--Schwarz yields
\[
 \|E_x\psi_j\|^2
 \leq
 K^2\left(\sum_ya_y\right)
 \left(\sum_ya_y\|E_yv_j\|^2\right).
 \tag{8.10}
\]
Sum over \(j\).  Bessel's inequality gives
\[
 \sum_j\|E_yv_j\|^2
 =
 \sum_j\|E_yQ_i\psi_j\|^2
 \leq\operatorname{rank}(E_yQ_i)\leq r_*.
 \tag{8.11}
\]
Therefore
\[
 \sum_j\|E_x\psi_j\|^2
 \leq
 K^2r_*\left(\sum_ya_y\right)^2.
 \tag{8.12}
\]
Since
\(N_i=\sum_j|\psi_j\rangle\langle\psi_j|\),
\[
 \|E_xN_i\|
 \leq\|E_xN_i\|_{\rm HS}
 =
 \left(\sum_j\|E_x\psi_j\|^2\right)^{1/2},
\]
and (8.7) follows.  If \(x\in D_i\), then the sum on the right contains
the term \(y=x\), while \(\|E_xN_i\|\leq1\); increasing
\(K\sqrt{r_*}\) to at least one proves the last statement.
\end{proof}

\begin{remark}[No bound on the number of near modes]
\label{rem:v11v8-rank}
The proof does not assume a volume-independent rank for \(N_i\).
Bessel's inequality is applied locally at each defect cell.  Many
well-separated near modes are allowed; their influence at \(x\) is
summed with exponential distance.
\end{remark}

\subsection{Producing the propagated coupling bound}

Condition (8.5) follows from more familiar data.  Assume
\[
 \|E_x(A_i-\lambda)^{-1}P_iE_z\|
 \leq R e^{-\mu d_i(x,z)}
 \tag{8.13}
\]
and that the kernel has range \(r_H\) with
\[
 \sup_z\sum_{y\in D_i}
 e^{\mu d_i(z,y)}\|E_zC_iE_y\|\leq J.
 \tag{8.14}
\]

\begin{proposition}[Bulk resolvent implies boundary propagation]
\label{prop:v11v8-propagation}
Under (8.13)--(8.14),
\[
 \|E_x(A_i-\lambda)^{-1}C_iE_y\|
 \leq RJ e^{-\mu d_i(x,y)}.
 \tag{8.15}
\]
Thus (8.5) holds with \(K=RJ\).
\end{proposition}

\begin{proof}
Insert the site decomposition between the resolvent and \(C_i\):
\[
 \|E_x(A_i-\lambda)^{-1}C_iE_y\|
 \leq
 \sum_z
 R e^{-\mu d_i(x,z)}\|E_zC_iE_y\|.
\]
The triangle inequality for the distance gives
\[
 e^{-\mu d_i(x,z)}
 \leq
 e^{-\mu d_i(x,y)}e^{\mu d_i(z,y)}.
\]
Factor out the first exponential and apply (8.14).
\end{proof}

A Dirichlet boundary can create its own near-zero states, so (8.4) and
(8.13) are genuine hypotheses.  They require a boundary mass, a geometric
resolvent induction, or another proof; bulk plaquette admissibility away
from \(D_i\) alone does not automatically establish them.

\subsection{Localized intertwiners}

Suppose both \(N(H_i)\) and \(N(H_{a,i})\) obey (8.7), possibly after
enlarging the same defect set by a fixed radius.  Assume
\[
 \|I_{a,i}\|_{\mu,i}\leq J_a.
 \tag{8.16}
\]

\begin{proposition}[Stability under a local measured intertwiner]
\label{prop:v11v8-intertwiner}
There is a constant \(\widetilde C_a\), depending only on the localization
stocks and \(J_a\), such that the corrected condition (8.1) holds:
\[
 \|E_xI_{a,i}N(H_{a,i})\|
 +
 \|E_xN(H_i)\|\,\|I_{a,i}\|
 \leq
 \widetilde C_a
 \sum_{y\in D_i}e^{-\widetilde\mu d_i(x,y)}
 \tag{8.17}
\]
for any fixed \(0<\widetilde\mu<\mu\).
\end{proposition}

\begin{proof}
The second term follows directly from (8.7) and
\(\|I_{a,i}\|\leq\|I_{a,i}\|_{\mu,i}\leq J_a\).
For the first, insert the site decomposition:
\[
 \|E_xI_{a,i}N(H_{a,i})\|
 \leq
 \sum_z\|E_xI_{a,i}E_z\|\,\|E_zN(H_{a,i})\|.
\]
The weighted-kernel bound on \(I_{a,i}\) and (8.7) bound this by a
convolution of two exponential kernels.  For
\(0<\widetilde\mu<\mu\),
\[
 e^{-\mu d(x,z)}e^{-\mu d(z,y)}
 \leq
 e^{-\widetilde\mu d(x,y)}
 e^{-(\mu-\widetilde\mu)(d(x,z)+d(z,y))}.
\]
The last factor is summable uniformly by the exponential volume stock.
Summing over \(z\) and then \(y\in D_i\) gives (8.17).
\end{proof}

\subsection{Random defect islands}

Let a cell enter \(D_i\) when it lies within a fixed radius \(r_D\) of a
plaquette with \(r_p>\epsilon_i\).  If every cell touches at most
\(c_D\) plaquettes in that enlarged neighbourhood, then
\[
 \sup_y\nu_i\{y\in D_i\}
 \leq
 c_D\sup_p\nu_i\{r_p>\epsilon_i\}.
 \tag{8.18}
\]
Combining (8.18) with a one-plaquette insertion tail and V7 gives local
near-mode invisibility with no global union bound.

\begin{theorem}[Deterministic-to-probabilistic mobility card]
\label{thm:v11v8-mobility-card}
Assume uniformly:
\begin{enumerate}[label=\textnormal{(\roman*)}]
\item outside the random defect neighbourhood, the Dirichlet bulk obeys
(8.4) and (8.13);
\item the boundary coupling obeys (8.14);
\item the measured intertwiners obey (8.16);
\item the single-plaquette defect probability tends to zero;
\item the outer Ward and residual \(L^q\) hypotheses of V7 hold.
\end{enumerate}
Then the expected componentwise artificial Ward edge tends to zero on
every normalized local test.
\end{theorem}

\begin{proof}
The Schur localization theorem and the propagation proposition prove
(8.7).  Local intertwiner stability gives the corrected V7 hypothesis
(8.1).  Equation (8.18) turns the one-plaquette tail into a vanishing
defect intensity.  The corrected V7 mobility theorem then gives the
conclusion.
\end{proof}

\begin{firewall}[Remaining hard input]
\label{fw:v11v8-hard-input}
The finite-dimensional Schur argument is complete, but the physical
Dirichlet bulk estimates (8.4), (8.13) have not been proved for the
coupled chiral kernel on random defect complements.  This is now the
precise mobility bottleneck.  It cannot be replaced by positivity of the
Yang--Mills Hessian or by the Standard Model anomaly arithmetic.
\end{firewall}

\begin{assessment}[V8 advance]
\label{ass:v11v8}
The projection-order gap in V7 has been corrected.  A gapped,
exponentially local Dirichlet bulk now yields deterministic localization
of all near-zero modes at defect islands without a bound on their total
number.  Together with a one-plaquette tail, this supplies the local
mobility row.  The next step is to formulate a geometric-resolvent
induction that proves the Dirichlet bulk estimate from separated
admissible blocks, or to record a counterexample if arbitrary defect
geometry prevents such an induction.
\end{assessment}


\section{Buffered geometric resolvent induction for the Dirichlet bulk}
\label{sec:v11v9}

V8 isolated the missing deterministic input as an exponentially weighted
inverse for the Dirichlet bulk outside defect islands.  This note gives a
finite block criterion for that inverse.

\subsection{An abstract approximate-inverse lemma}

Let \(\mathfrak A_{\mu,i}\) be the weighted kernel algebra of V3.  It is a
unital Banach algebra.

\begin{lemma}[Right parametrix closure]
\label{lem:v11v9-parametrix}
Let \(T_i\in\mathfrak A_{\mu,i}\).  Suppose there is
\(K_i\in\mathfrak A_{\mu,i}\) such that
\[
 T_iK_i=I+E_i,\qquad
 \|K_i\|_{\mu,i}\leq R,\qquad
 \|E_i\|_{\mu,i}\leq\eta<1.
 \tag{9.1}
\]
Then \(T_i\) is invertible and
\[
 T_i^{-1}=K_i(I+E_i)^{-1},\qquad
 \|T_i^{-1}\|_{\mu,i}\leq\frac{R}{1-\eta}.
 \tag{9.2}
\]
\end{lemma}

\begin{proof}
The Neumann series
\((I+E_i)^{-1}=\sum_{n\geq0}(-E_i)^n\)
converges in the Banach algebra and has norm at most
\((1-\eta)^{-1}\).  Thus the right side of (9.2) is a right inverse.
On a finite-dimensional square space, a right inverse is an inverse.
The norm estimate follows from submultiplicativity.
\end{proof}

\subsection{Buffered block cover}

Let \(\Omega_i\) be the bulk site set after removal of the enlarged defect
neighbourhood, and let
\[
 T_i(\lambda):=A_i-\lambda
 \quad\text{on}\quad\ell^2(\Omega_i),
 \qquad |\lambda|<\delta.
 \tag{9.3}
\]
Assume \(A_i\) has range at most \(r_0\) and
\[
 \|A_i\|_{\mu,i}\leq M.
 \tag{9.4}
\]

Choose blocks \(U_{j,i}\subset\Omega_i\), cores
\(V_{j,i}\subset U_{j,i}\), and diagonal cutoffs
\(\chi_{j,i},\eta_{j,i}\) with:
\begin{enumerate}[label=\textnormal{(B\arabic*)}]
\item \(0\leq\chi_{j,i},\eta_{j,i}\leq1\),
\(\supp\chi_{j,i}\subset V_{j,i}\), and
\(\sum_j\chi_{j,i}=I\) on \(\Omega_i\);
\item \(\eta_{j,i}=1\) on the \(r_0\)-neighbourhood of
\(\supp\chi_{j,i}\), and
\(\supp\eta_{j,i}\subset U_{j,i}\);
\item the transition shell
\(\supp[A_i,\eta_{j,i}]\) is at distance at least \(\ell\) from
\(\supp\chi_{j,i}\);
\item the families of core supports and transition shells have overlap at
most \(\kappa\), uniformly in \(i\).
\end{enumerate}

Let \(A_{j,i}\) be the Dirichlet restriction of \(A_i\) to \(U_{j,i}\).
Assume
\[
 G_{j,i}(\lambda):=(A_{j,i}-\lambda)^{-1}
 \tag{9.5}
\]
exists for \(|\lambda|<\delta\) and has kernel bound
\[
 \|E_xG_{j,i}(\lambda)E_y\|
 \leq R_0e^{-\mu d_i(x,y)}
 \quad(x,y\in U_{j,i}).
 \tag{9.6}
\]

Define the zero-extended block parametrix
\[
 K_i(\lambda):=
 \sum_j\eta_{j,i}G_{j,i}(\lambda)\chi_{j,i}.
 \tag{9.7}
\]

\begin{lemma}[Exact geometric resolvent error]
\label{lem:v11v9-error}
Under (B1)--(B3),
\[
 T_i(\lambda)K_i(\lambda)=I+E_i(\lambda),
 \tag{9.8}
\]
where
\[
 E_i(\lambda)=
 \sum_j[A_i,\eta_{j,i}]
 G_{j,i}(\lambda)\chi_{j,i}.
 \tag{9.9}
\]
\end{lemma}

\begin{proof}
On the support of \(\eta_{j,i}\), the full bulk operator and the
zero-extended Dirichlet operator agree for all rows retained by
\(\eta_{j,i}\).  Hence
\[
\begin{split}
 (A_i-\lambda)\eta_{j,i}G_{j,i}\chi_{j,i}
 &=
 \eta_{j,i}(A_{j,i}-\lambda)G_{j,i}\chi_{j,i}
 +[A_i,\eta_{j,i}]G_{j,i}\chi_{j,i}\\
 &=
 \eta_{j,i}\chi_{j,i}
 +[A_i,\eta_{j,i}]G_{j,i}\chi_{j,i}.
\end{split}
\]
Condition (B2) gives
\(\eta_{j,i}\chi_{j,i}=\chi_{j,i}\).  Sum over \(j\) and use
(B1) to obtain (9.8)--(9.9).
\end{proof}

\subsection{Smallness from the buffer}

Fix \(0<\mu'<\mu\).  Let \(C_{\rm cut}\) be a uniform bound for the
finite-range cutoff commutators in the \(\mu\)-weighted row and column
sums:
\[
 \|[A_i,\eta_{j,i}]\|_{\mu,i}\leq C_{\rm cut}M.
 \tag{9.10}
\]

\begin{proposition}[Buffered error estimate]
\label{prop:v11v9-buffer}
There is a geometric constant \(C_{\rm ov}\), depending only on
\(\kappa\), \(r_0\), and the uniform local degree, such that
\[
 \|E_i(\lambda)\|_{\mu',i}
 \leq
 C_{\rm ov}C_{\rm cut}MR_0
 e^{-(\mu-\mu')\ell}
 =:\eta_\ell
 \tag{9.11}
\]
uniformly for \(|\lambda|<\delta\).
Moreover
\[
 \|K_i(\lambda)\|_{\mu',i}
 \leq C_{\rm ov}R_0.
 \tag{9.12}
\]
\end{proposition}

\begin{proof}
The kernel of the \(j\)-th summand in (9.9) joins a point in the
transition shell to a point in \(\supp\chi_{j,i}\).  Their distance is at
least \(\ell\).  In the product of the \(\mu\)-decay from (9.6) with the
\(\mu'\)-weight, reserve
\[
 e^{-(\mu-\mu')d_i(x,y)}
 \leq e^{-(\mu-\mu')\ell}.
\]
The remaining \(\mu\)-weighted convolution is bounded by
\(C_{\rm cut}MR_0\) using the weighted algebra inequality.
Condition (B4) bounds the number of summands contributing to every row
and column by a constant \(C_{\rm ov}\), proving (9.11).
The same overlap argument applied to (9.7), without a commutator or
reserved distance, gives (9.12).
\end{proof}

\begin{theorem}[Geometric resolvent induction]
\label{thm:v11v9-geometric}
If the buffer is chosen so that
\[
 \eta_\ell=
 C_{\rm ov}C_{\rm cut}MR_0e^{-(\mu-\mu')\ell}<1,
 \tag{9.13}
\]
then for every \(|\lambda|<\delta\),
\[
 A_i-\lambda\ \text{is invertible on the Dirichlet bulk}
 \tag{9.14}
\]
and
\[
 \|(A_i-\lambda)^{-1}\|_{\mu',i}
 \leq
 \frac{C_{\rm ov}R_0}{1-\eta_\ell}.
 \tag{9.15}
\]
These bounds are independent of the total volume and the number of
blocks.
\end{theorem}

\begin{proof}
Equations (9.8), (9.11), and (9.12) satisfy the parametrix lemma with
\(T_i=A_i-\lambda\), \(K_i=K_i(\lambda)\), and
\(\eta=\eta_\ell\).  Equation (9.2) gives (9.14)--(9.15).
\end{proof}

\subsection{Cost of enlarging the random defect set}

Let \(S_i\) be the primary set of cells touching plaquettes with defect
larger than \(\epsilon_i\).  Enlarge it by graph distance \(L\):
\[
 D_i:=\{x:d_i(x,S_i)\leq L\}.
 \tag{9.16}
\]
If balls of radius \(L\) contain at most \(V(L)\) sites and
\[
 \sup_y\nu_i\{y\in S_i\}\leq\theta_i,
 \tag{9.17}
\]
then
\[
 \sup_x\nu_i\{x\in D_i\}\leq V(L)\theta_i.
 \tag{9.18}
\]

\begin{proof}
If \(x\in D_i\), some \(y\) in the radius-\(L\) ball about \(x\) belongs
to \(S_i\).  The union bound over at most \(V(L)\) sites gives (9.18).
\end{proof}

For a fixed buffer \(L\), vanishing primary defect intensity therefore
survives the enlargement.  If the block error requires
\(L=L_i\to\infty\), one must instead verify
\(V(L_i)\theta_i\to0\).  This is the entropy cost of the geometric
resolvent induction; it is local ball entropy rather than total-volume
entropy.

\subsection{A block-to-mobility theorem}

\begin{theorem}[Buffered-block mobility criterion]
\label{thm:v11v9-block-mobility}
Assume:
\begin{enumerate}[label=\textnormal{(G\arabic*)}]
\item the random primary defect intensity is \(\theta_i\to0\);
\item after enlargement by \(L_i\),
\(V(L_i)\theta_i\to0\);
\item the complement admits a buffered cover satisfying
(B1)--(B4), (9.4), (9.6), and
\(\sup_i\eta_{\ell_i}<1\);
\item the boundary coupling between bulk and defects has the finite-range
stock required in V8;
\item the outer-band, measured-intertwiner, and \(L^q\) Ward estimates of
V7 hold.
\end{enumerate}
Then the expected artificial componentwise Ward edge vanishes on every
normalized local test.
\end{theorem}

\begin{proof}
The geometric resolvent theorem gives the uniform Dirichlet bulk inverse
and exponential kernel estimate required in V8.  The finite-range boundary
stock gives the propagated coupling bound there.  The Schur localization
theorem localizes the near spectral projector at \(D_i\).
Equation (9.18) and (G2) make the enlarged defect intensity vanish.
The corrected mobility-gap theorem of V7 then proves the Ward conclusion.
\end{proof}

\subsection{Where arbitrary geometry can break the induction}

The block theorem has two independent failure modes:
\begin{enumerate}[label=\textnormal{(\alph*)}]
\item a local Dirichlet block may itself have a state in
\((-\delta,\delta)\), invalidating (9.5)--(9.6);
\item the defect geometry may leave no buffered cover with bounded overlap,
or may force \(L_i\) to grow so quickly that \(V(L_i)\theta_i\not\to0\).
\end{enumerate}
Neither failure is repaired by lowering the global density of defects
without further decorrelation or cluster control.

\begin{proposition}[Marginal sparsity does not forbid a macroscopic wall]
\label{prop:v11v9-wall}
On a periodic box of side \(L\), choose one coordinate hyperplane
uniformly from the \(L\) parallel hyperplanes and declare all its sites
defective.  Then each site is defective with probability \(L^{-1}\to0\),
but with probability one the defect set contains a system-spanning wall.
\end{proposition}

\begin{proof}
A fixed site lies on exactly one of the \(L\) possible hyperplanes, giving
probability \(1/L\).  Every sampled set is itself a system-spanning
hyperplane.
\end{proof}

Thus the one-site intensity used for local invisibility is enough once
deterministic localization is known, but it is not enough to prove that
localization.  A block induction may additionally require a cluster-tail
or nonpercolation estimate.

\begin{firewall}[Physical claim boundary]
\label{fw:v11v9-physical}
The geometric-resolvent argument is an exact implication.  The coupled
Standard Model--Einstein reduced density has not been shown to satisfy the
local Dirichlet block gap, the buffer cover, or a defect-cluster tail.
These are now finite geometric and probabilistic acquisitions rather than
an unspecified appeal to a mobility gap.
\end{firewall}

\begin{assessment}[V9 advance]
\label{ass:v11v9}
A volume-independent Dirichlet bulk inverse follows from gapped local
blocks, exponential block inverses, bounded overlap, and a buffer large
enough to make the commutator parametrix contractive.  Enlarging random
defects costs only local ball entropy.  The wall example proves that
vanishing marginal defect intensity cannot by itself establish the
deterministic cover.  V10 is the next mandatory companion audit; it will
look specifically for a transferable cluster or block-contraction
estimate before the programme chooses the next acquisition.
\end{assessment}


\section{V10 mandatory companion audit}
\label{sec:v11v10}

The second five-version boundary of the eleventh series requires another
inspection of the active companion programmes.

\begin{audit}[Files inspected]
\label{audit:v11v10-files}
The files and exact records were
\[
\begin{array}{c|r|r}
\text{file}&\text{lines}&\text{bytes}\\ \hline
\texttt{Renewal\_Yang\_Mills\_Mass\_Gap\_Research\_Notes\_V22.tex}
&10050&341433\\
\texttt{Renewal\_NS\_Stability\_Research\_Notes\_V26.tex}
&5319&278283.
\end{array}
\]
Their SHA--256 values are
\begin{center}
\small\ttfamily
775B384B03A010065CB77DAC9AA98A74818FEE306FE37F29C2C42C42BE2402CE
\end{center}
and
\begin{center}
\small\ttfamily
82C887F8C2C69E4F28FAD7C3DC8DE5B9099CB5A4BF431EBA1FFF3E91935978AD.
\end{center}
The NS note is unchanged.  The YM note has advanced from V20 to V22.
\end{audit}

\subsection{New Yang--Mills conclusions}

YM V21 executes a quadratic parity-cell certificate over the exact integer
Laurent ring.  Two independent constructions of every owner row agree,
the complete unreduced curl--gradient Ward residual vanishes exactly, and
a serialized Gram hash anchors the computation.  Exact coefficient
inventories sharply improve the complex inverse strip.  The note also
proves that agreement of the entire quadratic bank cannot determine the
cubic and quartic background response.

YM V22 constructs an exact noncommutative \(SU(3)\) word-jet engine through
order four over \(\mathbb Q(i)\).  It reproduces every first owner jet and
the V21 Gram hash, then checks nonzero cubic and quartic samples with full
permutation symmetry.  The remaining YM calculation is the
momentum-dependent harmonic and stationary contraction, followed by a
complete Ward residual and interval integration.

\begin{audit}[Transfer decision]
\label{audit:v11v10-transfer}
The following methods transfer to the present programme.
\begin{enumerate}[label=\textnormal{(\roman*)}]
\item Local chiral-kernel and comparison data should be generated from the
literal compact-link words over an exact coefficient algebra whenever
possible, with an independent quadratic hash or residual as a
normalization anchor.
\item Gauge covariance must be tested on the complete unreduced
one-particle space before chiral projection, tree removal, coarea
reduction, or a low/edge split.  The later reductions must then be
intertwined with that exact identity.
\item A spectral gap and quadratic Ward zero do not determine the
determinant response.  Projector derivatives, stationary-center terms,
Haar/coarea rows, and determinant-line connection terms are independent
data and require their own certificate.
\end{enumerate}
No YM numerical coefficient, parity-tree incidence array, or bosonic
edge scalar transfers.
\end{audit}

\begin{audit}[Cluster search]
\label{audit:v11v10-cluster}
Neither YM V21--V22 nor NS V26 proves a random defect-cluster tail, a
nonpercolation estimate, or a Dirichlet Wilson--Dirac block gap.  Thus the
new material does not close V9 assumptions (G1)--(G3).  The NS rank-one
ancient-kernel derivatives concern a different equation and do not supply
a lattice geometric resolvent estimate.
\end{audit}

\subsection{Adjustment of the next step}

The absence of a transferable cluster estimate prevents progress by
merely repeating the probability route.  The executable YM advance instead
suggests moving to the other unverified half of the V9 card: the exact
finite block certificate.

\begin{decision}[Post-audit direction]
\label{dec:v11v10-direction}
V11 will prove the algebra underlying such a certificate:
\begin{enumerate}[label=\textnormal{(\alph*)}]
\item finite gauge covariance implies the exact infinitesimal Ward
commutator for the Hermitian chiral kernel;
\item a gapped Riesz projector inherits that commutator before any
low/edge split;
\item an interval lower bound on the squared Dirichlet kernel, combined
with a literal link-word derivative stock, certifies both the block gap
and a background-tube radius;
\item the determinant and coarea rows remain separate and visible.
\end{enumerate}
This produces a finite executable target for each retained block without
claiming that the required arrays have already been generated.
\end{decision}

\begin{firewall}[Audit claim boundary]
\label{fw:v11v10-boundary}
The companion audit changes the verification method, not the physical
status.  No active companion note supplies the SM chiral determinant
phase, the defect-cluster law, the coupled reduced-density tail, or the
Einstein constraint estimates.  The full continuum remains open.
\end{firewall}

\begin{assessment}[V10 audit result]
\label{ass:v11v10}
The useful transfer is an exact finite-word certification workflow.
The next proof will connect that workflow to the chiral kernel, its Riesz
projector, and the local Dirichlet block stock required by the mobility
route.  A substantive V11 follows this audit as required.
\end{assessment}


\section{Exact chiral-kernel Ward and finite block certificates}
\label{sec:v11v11}

This post-audit note constructs the algebraic and interval certification
layer required by the V9 block criterion.  The determinant-line phase is
kept separate from covariance of the one-particle projector.

\subsection{Gauge covariance before reduction}

Let \(\mathcal G_i\) be the compact lattice gauge group and let
\(\mathcal U_i(g)\) be its unitary action on the complete one-particle
space.  Let \(H_i(U)\) be a Hermitian Wilson--Dirac or overlap kernel
satisfying
\[
 H_i(U^g)=\mathcal U_i(g)H_i(U)\mathcal U_i(g)^{-1}.
 \tag{11.1}
\]
For a Lie algebra field \(\omega\), put
\[
 g(t)=e^{t\omega},\qquad
 \Omega_i(\omega):=
 \left.\frac{d}{dt}\right|_{t=0}\mathcal U_i(g(t)),
 \tag{11.2}
\]
and denote the induced link variation by \(\delta_\omega U\).

\begin{theorem}[Exact kernel Ward commutator]
\label{thm:v11v11-kernel-ward}
If \(H_i\) is differentiable, then
\[
 D H_i(U)[\delta_\omega U]
 =
 [\Omega_i(\omega),H_i(U)].
 \tag{11.3}
\]
The identity holds on the complete one-particle space before chiral
projection, tree removal, constraint reduction, or a spectral split.
\end{theorem}

\begin{proof}
Differentiate (11.1) along \(g(t)\) at \(t=0\).  The derivative of the
inverse representation is \(-\Omega_i(\omega)\), giving
\[
 DH_i[\delta_\omega U]
 =
 \Omega_iH_i-H_i\Omega_i.
\]
\end{proof}

For a literal compact-link word, both sides of (11.3) belong to a finite
noncommutative coefficient algebra.  Thus their difference can be tested
coefficient by coefficient without floating-point sampling.

\subsection{Propagation to the Riesz projector}

Let \(\Gamma_i\) be a closed contour in the resolvent set of \(H_i(U)\),
and define
\[
 P_i(U)=\frac1{2\pi i}\int_{\Gamma_i}(z-H_i(U))^{-1}\,dz.
 \tag{11.4}
\]

\begin{theorem}[Projector Ward commutator]
\label{thm:v11v11-projector-ward}
Under the hypotheses of the kernel Ward theorem, if \(\Gamma_i\) remains
in the resolvent set along the infinitesimal variation, then
\[
 DP_i(U)[\delta_\omega U]
 =
 [\Omega_i(\omega),P_i(U)].
 \tag{11.5}
\]
In particular, forming the exact covariant Riesz projector introduces no
additional one-particle gauge residual.
\end{theorem}

\begin{proof}
For \(R_i(z)=(z-H_i)^{-1}\), differentiation of the inverse gives
\[
 DR_i(z)[\delta_\omega U]
 =
 R_i(z)\,DH_i[\delta_\omega U]\,R_i(z).
\]
Using (11.3),
\[
 R_i[\Omega_i,H_i]R_i
 =
 \Omega_iR_i-R_i\Omega_i
 =
 [\Omega_i,R_i].
\]
Differentiate (11.4) under the finite contour integral and integrate the
commutator to obtain (11.5).
\end{proof}

\begin{corollary}[Covariance of the overlap sign]
\label{cor:v11v11-sign}
If zero is separated from \(\spec H_i(U)\), then
\[
 \operatorname{sgn}H_i=I-2P_i^-
\]
obeys
\[
 D(\operatorname{sgn}H_i)[\delta_\omega U]
 =
 [\Omega_i,\operatorname{sgn}H_i].
 \tag{11.6}
\]
Every chiral projector built algebraically from this sign and fixed
covariant spin matrices obeys its corresponding commutator Ward identity.
\end{corollary}

\begin{proof}
Differentiate the affine relation between the sign and the negative Riesz
projector, then use (11.5).  Algebraic combinations preserve covariance
when the fixed matrices commute or transform in the declared
representation.
\end{proof}

\begin{firewall}[Projector covariance is not measure invariance]
\label{fw:v11v11-phase}
Equations (11.3)--(11.6) do not determine the phase of the chiral
determinant.  A local basis for \(\operatorname{ran}P_i\) contributes its
Berry connection, and local scalarizations differ by determinant-line
transition functions.  Ghost, Haar, coarea, stationary-center, boundary,
and counterterm rows must be added before the complete Ward response is
tested.
\end{firewall}

\subsection{First finite gap certificate: squared-kernel positivity}

Let \(H_{B,i}\) be a finite Dirichlet block kernel at a chosen background
\(B\).

\begin{proposition}[Squared-kernel gap certificate]
\label{prop:v11v11-square-gap}
If a rigorous exact or interval calculation proves
\[
 H_{B,i}^2-\gamma_0^2I\succeq0
 \tag{11.7}
\]
for some \(\gamma_0>0\), then
\[
 \spec H_{B,i}\cap(-\gamma_0,\gamma_0)=\varnothing,
 \qquad
 \|H_{B,i}^{-1}\|\leq\gamma_0^{-1}.
 \tag{11.8}
\]
A positive lower bound for every pivot in a verified
\(LDL^*\) factorization of (11.7) is sufficient.
\end{proposition}

\begin{proof}
For every vector \(v\),
\[
 \|H_{B,i}v\|^2
 =
 \langle v,H_{B,i}^2v\rangle
 \geq\gamma_0^2\|v\|^2.
\]
Thus the least singular value is at least \(\gamma_0\).  Hermiticity
identifies singular-value magnitude with spectral magnitude, proving
(11.8).  A verified \(LDL^*\) factorization with nonnegative diagonal
pivots proves positive semidefiniteness.
\end{proof}

\subsection{Second finite gap certificate: inverse residual}

An approximate inverse is often cheaper than an eigenvalue enclosure.

\begin{theorem}[Residual inverse certificate]
\label{thm:v11v11-residual-gap}
Let \(X_{B,i}\) be an exact or interval matrix satisfying
\[
 \|I-X_{B,i}H_{B,i}\|\leq\eta<1,
 \qquad
 \|X_{B,i}\|\leq K.
 \tag{11.9}
\]
Then \(H_{B,i}\) is invertible,
\[
 H_{B,i}^{-1}=(I-(I-X_{B,i}H_{B,i}))^{-1}X_{B,i},
 \tag{11.10}
\]
and
\[
 \|H_{B,i}^{-1}\|\leq\frac{K}{1-\eta},
 \qquad
 \min_{\lambda\in\spec H_{B,i}}|\lambda|
 \geq\frac{1-\eta}{K}.
 \tag{11.11}
\]
\end{theorem}

\begin{proof}
Put \(E=I-X_{B,i}H_{B,i}\).  The Neumann series gives
\((I-E)^{-1}\) with norm at most \((1-\eta)^{-1}\).
Since \(X_{B,i}H_{B,i}=I-E\), equation (11.10) is a left inverse.
In finite dimension it is the inverse.  The norm bound follows, and the
reciprocal of the inverse norm is the least singular value.  Hermiticity
then gives the spectral statement.
\end{proof}

Both certificates can be performed over rational arithmetic when the
background entries are algebraic, or by outward-rounded intervals when
they are not.

\subsection{Literal derivative stock and tube radius}

Suppose the block kernel is a finite word function of chord coordinates
and a rigorous derivative enclosure gives
\[
 \sup_{U\in\mathcal T}
 \|DH_i(U)\|_{\rm ch\to op}\leq L.
 \tag{11.12}
\]
Let \(\gamma_0\) be obtained from either (11.7) or (11.11).

\begin{theorem}[Certified chiral-gap tube]
\label{thm:v11v11-certified-tube}
For every \(U\) joined to \(B\) by a chord path of length
\(d_{\rm ch}(U,B)\),
\[
 \|H_i(U)-H_i(B)\|
 \leq Ld_{\rm ch}(U,B).
 \tag{11.13}
\]
If
\[
 d_{\rm ch}(U,B)\leq r_B<\frac{\gamma_0}{L},
 \tag{11.14}
\]
then
\[
 \min_{\lambda\in\spec H_i(U)}|\lambda|
 \geq\gamma_0-Lr_B>0.
 \tag{11.15}
\]
In particular \(r_B\leq\gamma_0/(2L)\) gives a gap at least
\(\gamma_0/2\).
\end{theorem}

\begin{proof}
Integrate \(DH_i\) along the chord path to obtain (11.13).
The least singular value is 1-Lipschitz in operator norm, so
\[
 s_{\min}(H_i(U))
 \geq s_{\min}(H_i(B))-\|H_i(U)-H_i(B)\|.
\]
Insert the base gap and (11.13).  Hermiticity converts the least singular
value to the minimum spectral magnitude.
\end{proof}

\begin{proposition}[Weighted base inverse from finite range]
\label{prop:v11v11-combes}
Assume \(H_{B,i}\) has range \(r_0\),
\[
 \|H_{B,i}\|_{0,i}\leq M,
 \qquad
 \|H_{B,i}^{-1}\|\leq\kappa.
 \tag{11.16}
\]
If
\[
 \kappa M(e^{\mu r_0}-1)\leq\frac12,
 \tag{11.17}
\]
then
\[
 \|E_xH_{B,i}^{-1}E_y\|
 \leq2\kappa e^{-\mu d_i(x,y)}.
 \tag{11.18}
\]
For every \(0<\mu'<\mu\), uniform exponential volume growth then gives a
finite \(\mu'\)-weighted inverse stock.
\end{proposition}

\begin{proof}
Fix \(y\) and conjugate by the diagonal weight
\(W_y(x)=e^{\mu d_i(x,y)}\).  Finite range and the triangle inequality
give
\[
 \|W_yH_{B,i}W_y^{-1}-H_{B,i}\|
 \leq M(e^{\mu r_0}-1).
\]
Condition (11.17) and the Neumann series show
\[
 \|(W_yH_{B,i}W_y^{-1})^{-1}\|\leq2\kappa.
\]
Since this inverse is \(W_yH_{B,i}^{-1}W_y^{-1}\), its \(x,y\) block gives
(11.18).  Multiplying by \(e^{\mu'd_i(x,y)}\) and summing leaves the
summable factor \(e^{-(\mu-\mu')d_i(x,y)}\).
\end{proof}

The same conjugation applies to \(H_{B,i}-\lambda\) throughout a compact
near-zero interval once a uniform inverse certificate is supplied there.

\subsection{Finite certificate ledger}

For each candidate block and sector background, an auditable certificate
now has the following rows:
\begin{enumerate}[label=\textnormal{(C\arabic*)}]
\item a canonical basis, orientation, link-word order, and serialized
input hash;
\item exact Hermiticity and gauge-covariance residuals for \(H_i\) on the
complete one-particle space;
\item either the squared-kernel certificate (11.7) or the inverse-residual
certificate (11.9);
\item a literal derivative stock \(L\) and the resulting tube radius;
\item a finite-range stock \(M,r_0\) and the weighted inverse bound;
\item exact propagation of the kernel Ward commutator to the Riesz and
chiral projectors;
\item separate determinant-line connection, ghost, Haar, coarea,
stationary-center, boundary, and counterterm rows;
\item the complete Ward residual after all rows are placed in the common
line-valued coefficient space and before any spectral split.
\end{enumerate}

\begin{theorem}[Soundness of the finite block certificate]
\label{thm:v11v11-soundness}
If rows (C1)--(C8) are verified with positive gap and contraction margins,
then that block supplies the base gap, weighted local inverse,
background-tube radius, and complete pre-split Ward identity required by
the V6--V9 analytic implications.
\end{theorem}

\begin{proof}
Rows (C2), (C6), and (C8) give the exact unreduced and complete Ward
identities by the commutator theorems.  Row (C3) gives the base gap.
Row (C4) gives its tube persistence.  Row (C5) gives the weighted inverse
stock.  Row (C7) prevents the one-particle identity from being mistaken
for determinant-measure invariance.  These are exactly the corresponding
hypotheses used in V6--V9.
\end{proof}

\begin{firewall}[Execution remains open]
\label{fw:v11v11-execution}
No Standard Model block array, interval \(LDL^*\) certificate, inverse
residual, or complete determinant-line Ward residual has been executed in
this note.  The theorem makes the finite target auditable; it does not
supply its numerical inputs.  The random block-cover and Einstein
reduced-density estimates also remain open.
\end{firewall}

\begin{assessment}[V11 advance]
\label{ass:v11v11}
Exact gauge covariance now propagates rigorously from the literal chiral
kernel to its Riesz and overlap projectors.  Two finite gap certificates,
a derivative tube radius, and a finite-range weighted inverse theorem
complete the deterministic certification logic for a retained block.
The immediate acquisition is to instantiate the certificate for the
free or periodic Standard Model kernel and then perturb it over a declared
sector background, while keeping the determinant-line row separate.
\end{assessment}


\section{Exact free and flat-holonomy overlap-kernel certificate}
\label{sec:v11v12}

V11 left the first finite certificate uninstantiated.  This note evaluates
it for the four-dimensional Wilson kernel at the free background and at
commuting flat holonomies.  Lattice units and Wilson parameter \(r=1\)
are used.

\subsection{The complete free kernel}

Let \(\gamma_\mu\), \(1\leq\mu\leq4\), be Hermitian Euclidean gamma
matrices satisfying
\[
 \gamma_\mu\gamma_\nu+\gamma_\nu\gamma_\mu
 =2\delta_{\mu\nu}I,
 \qquad
 \gamma_5\gamma_\mu=-\gamma_\mu\gamma_5,
 \qquad
 \gamma_5^2=I.
 \tag{12.1}
\]
For \(0<m<2\), define
\[
 D_W(k)-m
 =
 i\sum_{\mu=1}^4\gamma_\mu\sin k_\mu
 +W(k)-m,
 \qquad
 W(k):=\sum_{\mu=1}^4(1-\cos k_\mu),
 \tag{12.2}
\]
and
\[
 H_m(k):=\gamma_5(D_W(k)-m).
 \tag{12.3}
\]

\begin{lemma}[Exact square]
\label{lem:v11v12-square}
For every \(k\in\mathbb T^4\),
\[
 H_m(k)^2
 =
 \left[
 \sum_{\mu=1}^4\sin^2k_\mu+(W(k)-m)^2
 \right]I.
 \tag{12.4}
\]
\end{lemma}

\begin{proof}
Conjugation by \(\gamma_5\) reverses the sign of
\(i\gamma_\mu\sin k_\mu\) and leaves the scalar \(W-m\) fixed.  Hence
\[
 H_m^2
 =
 \left(-i\sum_\mu\gamma_\mu\sin k_\mu+W-m\right)
 \left(i\sum_\nu\gamma_\nu\sin k_\nu+W-m\right).
\]
The two scalar--gamma cross terms cancel.  The Clifford relations cancel
all mixed gamma products and leave
\((\sum_\mu\sin^2k_\mu)I\), proving (12.4).
\end{proof}

\subsection{Exact minimum over the Brillouin torus}

Put
\[
 t_\mu:=1-\cos k_\mu\in[0,2],
 \qquad q:=\sum_\mu t_\mu.
 \tag{12.5}
\]
Since \(\sin^2k_\mu=2t_\mu-t_\mu^2\), the scalar in (12.4) is
\[
 F_m(t)=2q-\sum_\mu t_\mu^2+(q-m)^2.
 \tag{12.6}
\]

\begin{lemma}[Maximum square sum at fixed total]
\label{lem:v11v12-square-sum}
If \(t_\mu\in[0,2]\) and \(q=2n+s\), where
\(n\in\{0,1,2,3\}\) and \(0\leq s\leq2\), then
\[
 \sum_{\mu=1}^4t_\mu^2\leq4n+s^2.
 \tag{12.7}
\]
Equality is attained by \(n\) entries equal to \(2\), one entry equal to
\(s\), and all remaining entries zero.
\end{lemma}

\begin{proof}
If two entries lie strictly between \(0\) and \(2\), transfer a small
amount from the smaller to the larger while keeping their sum fixed.
The square sum increases.  Iterating reaches an extreme point of the
slice \(\sum t_\mu=q\), at which at most one coordinate is not \(0\) or
\(2\).  The stated extreme point and value follow.
\end{proof}

\begin{theorem}[Exact free Wilson-kernel gap]
\label{thm:v11v12-free-gap}
For \(0<m<2\),
\[
 \min_{k\in\mathbb T^4}
 \min_{\lambda\in\spec H_m(k)}|\lambda|
 =
 \gamma_m,
 \qquad
 \gamma_m:=\min\{m,2-m\}.
 \tag{12.8}
\]
Equivalently,
\[
 H_m(k)^2\succeq\gamma_m^2I
 \quad\text{for every }k.
 \tag{12.9}
\]
The choice \(m=1\) maximizes this certified gap and gives
\[
 H_1(k)^2\succeq I.
 \tag{12.10}
\]
\end{theorem}

\begin{proof}
On the interval \(q=2n+s\), the fixed-\(q\) minimum of (12.6) is obtained
by maximizing the square sum.  By the preceding lemma it is
\[
 2s-s^2+(2n+s-m)^2.
 \tag{12.11}
\]
For \(n=0\), this equals
\[
 m^2+2(1-m)s,
 \tag{12.12}
\]
whose minimum on \(0\leq s\leq2\) is \(m^2\) when \(m\leq1\) and
\((2-m)^2\) when \(m\geq1\).  For \(n\geq1\), the derivative of
(12.11) with respect to \(s\) is
\[
 2(2n+1-m)>0,
\]
so the minimum is at \(s=0\) and equals \((2n-m)^2\).  Among
\(n=1,2,3,4\), the smallest is \((2-m)^2\).  Thus the global minimum of
the squared spectral magnitude is
\(\min\{m^2,(2-m)^2\}\).  Equality occurs at \(k=0\) or at a corner with
one component equal to \(\pi\).  Taking square roots proves (12.8).
The maximum of \(\min\{m,2-m\}\) occurs at \(m=1\).
\end{proof}

This proves the V11 squared-kernel certificate uniformly in the periodic
volume.

\subsection{Commuting flat Standard Model holonomies}

Let \(R\) be any finite-dimensional unitary Standard Model
representation.  Suppose the four constant link matrices
\(V_\mu\in R(G_{\rm SM})\) commute.  They are simultaneously unitarily
diagonalizable.  On a common weight vector write
\[
 V_\mu v=e^{i\theta_\mu}v.
 \tag{12.13}
\]

\begin{proposition}[Flat-holonomy gap]
\label{prop:v11v12-holonomy}
On the weight space (12.13), the periodic kernel is unitarily equivalent
to
\[
 H_m(k+\theta).
 \tag{12.14}
\]
Consequently every commuting flat holonomy, in every Standard Model
multiplet, obeys the same volume-independent gap (12.8).
\end{proposition}

\begin{proof}
Fourier transformation diagonalizes translations.  Multiplication of the
forward and backward shifts by \(V_\mu\) and \(V_\mu^{-1}\) replaces
\(e^{\pm ik_\mu}\) by \(e^{\pm i(k_\mu+\theta_\mu)}\) on the chosen
weight space.  This is precisely (12.14).  The allowed finite-volume
momenta form a subset of \(\mathbb T^4\), so the lower bound over the full
torus in (12.8) applies.
\end{proof}

The proposition includes nontrivial toron phases.  It does not include
background curvature or a topologically nontrivial instanton sector.

\subsection{Literal link-word Lipschitz stock}

In position space,
\[
\begin{split}
 D_W(U)-m
 ={}&(4-m)I\\
 &-\frac12\sum_{\mu=1}^4
 \left[
 (I-\gamma_\mu)U_{x,\mu}T_{+\mu}
 +(I+\gamma_\mu)U_{x-\widehat\mu,\mu}^{-1}T_{-\mu}
 \right].
\end{split}
\tag{12.15}
\]
Let
\[
 d_\infty(U,V):=\sup_{x,\mu}\|U_{x,\mu}-V_{x,\mu}\|.
 \tag{12.16}
\]

\begin{proposition}[Exact chord perturbation bound]
\label{prop:v11v12-lipschitz}
For unitary link fields \(U,V\),
\[
 \|H_m(U)-H_m(V)\|
 \leq8\,d_\infty(U,V).
 \tag{12.17}
\]
In the weighted row--column norm,
\[
 \|H_m(U)-H_m(V)\|_{\mu}
 \leq8e^\mu d_\infty(U,V).
 \tag{12.18}
\]
\end{proposition}

\begin{proof}
The matrices
\((I\pm\gamma_\mu)/2\) are orthogonal projections and have norm one.
Each row of (12.15) contains four forward and four backward hopping
blocks.  Since
\[
 \|U^{-1}-V^{-1}\|
 =
 \|U^{-1}(V-U)V^{-1}\|
 =
 \|U-V\|,
\]
each difference block has norm at most \(d_\infty(U,V)\).  Both the row
and column sums are therefore at most \(8d_\infty(U,V)\); the Schur test
gives (12.17).  Every hopping block has range one, so the exponential
weight multiplies the same estimate by \(e^\mu\), giving (12.18).
Multiplication by the unitary \(\gamma_5\) changes neither norm.
\end{proof}

\begin{corollary}[Explicit flat-background tube]
\label{cor:v11v12-tube}
Around any commuting flat background,
\[
 d_\infty(U,V)<\frac{\gamma_m}{8}
 \tag{12.19}
\]
preserves a nonzero Hermitian-kernel gap.  On the closed tube
\[
 d_\infty(U,V)\leq\frac{\gamma_m}{16},
 \tag{12.20}
\]
the gap is at least \(\gamma_m/2\).  For \(m=1\), the certified half-gap
tube has radius \(1/16\).
\end{corollary}

\begin{proof}
Apply the least-singular-value perturbation inequality to (12.8) and
(12.17).  On (12.20),
\(\gamma_m-8d_\infty(U,V)\geq\gamma_m/2\).
\end{proof}

\subsection{Explicit weighted inverse radius}

The unweighted row--column coefficient stock of (12.15) satisfies
\[
 M_m\leq |4-m|+8.
 \tag{12.21}
\]
For the free or commuting flat background and
\(|\lambda|\leq\gamma_m/2\),
\[
 \|(H_m-\lambda)^{-1}\|\leq\frac2{\gamma_m}.
 \tag{12.22}
\]

\begin{proposition}[Uniform near-zero Combes--Thomas stock]
\label{prop:v11v12-weighted}
If
\[
 0<\mu\leq
 \log\left(1+\frac{\gamma_m}{4M_m}\right),
 \tag{12.23}
\]
then, uniformly for \(|\lambda|\leq\gamma_m/2\),
\[
 \|E_x(H_m-\lambda)^{-1}E_y\|
 \leq
 \frac4{\gamma_m}e^{-\mu d(x,y)}.
 \tag{12.24}
\]
For \(m=1\), one may take
\[
 M_1=11,\qquad
 \mu_1=\log\left(1+\frac1{44}\right).
 \tag{12.25}
\]
\end{proposition}

\begin{proof}
The distance-weight conjugation changes only the range-one hopping part.
Its norm difference is at most
\(M_m(e^\mu-1)\).  By (12.22) and (12.23), the inverse-relative
perturbation is at most
\[
 \frac2{\gamma_m}M_m(e^\mu-1)\leq\frac12.
\]
The Neumann series doubles (12.22), giving \(4/\gamma_m\).
Undoing the weight gives (12.24).  Equation (12.25) is (12.21)--(12.23)
at \(m=1\).
\end{proof}

\subsection{Ward and determinant-line rows}

The link-word formula (12.15) is exactly gauge covariant on the complete
spin--colour--site space.  Therefore the kernel and Riesz-projector Ward
commutators proved in V11 hold throughout the tube (12.19).  This
instantiates rows (C2)--(C6) of the V11 certificate for the free and
commuting-flat backgrounds.

\begin{firewall}[Rows not instantiated]
\label{fw:v11v12-uninstantiated}
The calculation does not instantiate the determinant-line connection,
global phase, ghost, coarea, Einstein constraint, boundary, or
counterterm rows (C7)--(C8).  It also does not prove the Dirichlet block
gap used in V8--V9: restricting a negative Wilson mass to a block can
create boundary spectral effects that are absent on a periodic torus.
Nor does the flat tube cover curved topological sectors.
\end{firewall}

\begin{assessment}[V12 advance]
\label{ass:v11v12}
The first chiral-kernel certificate is now explicit.  For every Standard
Model representation, the periodic free and commuting-flat kernel has
the exact gap \(\min\{m,2-m\}\), maximized at one; its link-word
Lipschitz constant is eight; the half-gap tube has radius
\(\gamma_m/16\); and a uniform near-zero weighted inverse stock is given
by (12.23)--(12.24).  The next upstream bottleneck is the boundary row:
either certify a gauge-covariant Dirichlet block realization with a
uniform gap or show that boundary modes force a different buffered
parametrix.
\end{assessment}


\section{Dirichlet boundary obstruction and gapped block completions}
\label{sec:v11v13}

V12 explicitly warned that the periodic gap does not imply a Dirichlet
block gap.  The obstruction is real: in the Wilson mass interval used by
the overlap construction, a half-space boundary supports an exact
zero mode.  This note proves the obstruction and repairs the V9
parametrix by allowing a gapped completion that agrees with the physical
operator only on the buffered core.

\subsection{An exact Wilson boundary mode}

Let the first coordinate occupy the half-line
\(\mathbb N=\{1,2,\ldots\}\), impose the missing-link Dirichlet rule at
zero, and take zero momentum in the other three directions.  The
one-dimensional row of \(D_W-m\) is
\[
 ((D_W-m)\psi)_n
 =
 (1-m)\psi_n
 -\frac12(I-\gamma_1)\psi_{n+1}
 -\frac12(I+\gamma_1)\psi_{n-1},
 \tag{13.1}
\]
where \(\psi_0=0\).

\begin{theorem}[Exact half-space boundary zero mode]
\label{thm:v11v13-half-space}
For every \(0<m<2\), put \(r=1-m\), so \(|r|<1\).  If
\[
 \gamma_1v=-v,\qquad v\ne0,
 \tag{13.2}
\]
then
\[
 \psi_n=r^{n-1}v
 \tag{13.3}
\]
belongs to \(\ell^2(\mathbb N)\) and satisfies
\[
 (D_W-m)\psi=0,\qquad H_m\psi=0.
 \tag{13.4}
\]
Thus the naive Dirichlet half-space has no gap at zero throughout the
overlap mass window.
\end{theorem}

\begin{proof}
The spinor relations in (13.2) give
\[
 \frac12(I-\gamma_1)v=v,\qquad
 \frac12(I+\gamma_1)v=0.
\]
Therefore every row of (13.1), including \(n=1\), becomes
\[
 r\psi_n-\psi_{n+1}=0,
\]
which holds by (13.3).  Since \(|r|<1\),
\[
 \sum_{n\geq1}\|\psi_n\|^2
 =
 \frac{\|v\|^2}{1-r^2}<\infty.
\]
Multiplication by the unitary \(\gamma_5\) proves the second equality in
(13.4).
\end{proof}

\begin{corollary}[No uniform finite-slab Dirichlet gap]
\label{cor:v11v13-slab}
On the slab \(1\leq n\leq L\), use the truncated vector
\(\psi_n=r^{n-1}v\).  Then
\[
 \frac{\|(D_W-m)\psi\|}{\|\psi\|}
 \leq
 |r|^L
 \left(\frac{1-r^2}{1-|r|^{2L}}\right)^{1/2}.
 \tag{13.5}
\]
Hence the least singular value tends to zero exponentially as
\(L\to\infty\).  At \(m=1\), the boundary vector supported at \(n=1\)
is an exact zero mode for every \(L\).
\end{corollary}

\begin{proof}
All rows except the last cancel as in the half-space proof.  At \(n=L\),
the omitted forward value leaves the residual \(r\psi_L=r^Lv\).
The squared norm of the trial vector is
\[
 \|\psi\|^2
 =
 \|v\|^2\sum_{n=0}^{L-1}|r|^{2n}
 =
 \|v\|^2\frac{1-|r|^{2L}}{1-r^2}.
\]
Divide the residual norm by this expression.  The variational
characterization of the least singular value gives the conclusion.
\end{proof}

\begin{firewall}[Correction to the V8--V9 route]
\label{fw:v11v13-dirichlet-correction}
A volume-uniform Dirichlet block gap cannot be assumed for the unmodified
Wilson kernel in the overlap mass window.  Consequently V8 condition
(8.4) and V9 local inverse condition (9.5), when interpreted as naive
missing-link Dirichlet restrictions, are false already at the free
background.  The implication theorems remain algebraically correct, but
their physical application requires a different local inverse.
\end{firewall}

\subsection{Completed local blocks}

Let \(A_i\) denote the physical bulk kernel on \(\Omega_i\).  For each
core \(V_{j,i}\), choose an auxiliary finite block space
\(\widetilde U_{j,i}\) and a Hermitian completion
\(\widetilde A_{j,i}\).  It may use periodic wrap links, a
gauge-covariant boundary mass, or another declared closure.  Let
\(\chi_{j,i}\) be supported in the core and let \(\eta_{j,i}=1\) on the
range-\(r_0\) neighbourhood of \(\supp\chi_{j,i}\).

After identifying the physical sites in the support of \(\eta_{j,i}\)
with their auxiliary copies, put
\[
 \Delta_{j,i}:=
 \eta_{j,i}(A_i-\widetilde A_{j,i}).
 \tag{13.6}
\]
Exact core agreement means \(\Delta_{j,i}=0\).

Assume
\[
 \widetilde G_{j,i}(\lambda)
 :=
 (\widetilde A_{j,i}-\lambda)^{-1}
 \tag{13.7}
\]
exists and has the uniform exponential stock used in V9.  Extend
\(\eta_{j,i}\widetilde G_{j,i}\chi_{j,i}\) by zero to the physical bulk
and define
\[
 \widetilde K_i(\lambda)
 :=
 \sum_j
 \eta_{j,i}\widetilde G_{j,i}(\lambda)\chi_{j,i}.
 \tag{13.8}
\]

\begin{lemma}[Completed-block parametrix identity]
\label{lem:v11v13-completed-identity}
If \(\sum_j\chi_{j,i}=I\) on the physical bulk, then
\[
 (A_i-\lambda)\widetilde K_i(\lambda)
 =
 I+\widetilde E_i^{\rm shell}(\lambda)
 +\widetilde E_i^{\rm cmp}(\lambda),
 \tag{13.9}
\]
where
\[
 \widetilde E_i^{\rm shell}
 =
 \sum_j[A_i,\eta_{j,i}]
 \widetilde G_{j,i}\chi_{j,i},
 \tag{13.10}
\]
and
\[
 \widetilde E_i^{\rm cmp}
 =
 \sum_j\Delta_{j,i}
 \widetilde G_{j,i}\chi_{j,i}.
 \tag{13.11}
\]
\end{lemma}

\begin{proof}
For one block,
\[
\begin{split}
 (A_i-\lambda)\eta_j\widetilde G_j\chi_j
 ={}&
 \eta_j(\widetilde A_j-\lambda)\widetilde G_j\chi_j\\
 &+[A_i,\eta_j]\widetilde G_j\chi_j
 +\eta_j(A_i-\widetilde A_j)\widetilde G_j\chi_j\\
 ={}&
 \chi_j+[A_i,\eta_j]\widetilde G_j\chi_j
 +\Delta_j\widetilde G_j\chi_j.
\end{split}
\]
Sum over \(j\).
\end{proof}

\begin{theorem}[Gapped-completion geometric resolvent]
\label{thm:v11v13-completed-resolvent}
Suppose the completed inverse and bounded-overlap hypotheses give
\[
 \|\widetilde K_i(\lambda)\|_{\mu'}\leq K_0,
 \tag{13.12}
\]
\[
 \|\widetilde E_i^{\rm shell}(\lambda)\|_{\mu'}
 \leq\eta_{\rm shell},
 \qquad
 \|\widetilde E_i^{\rm cmp}(\lambda)\|_{\mu'}
 \leq\eta_{\rm cmp},
 \tag{13.13}
\]
uniformly for \(|\lambda|<\delta\), with
\[
 \eta_{\rm shell}+\eta_{\rm cmp}<1.
 \tag{13.14}
\]
Then
\[
 A_i-\lambda\ \text{is invertible}
 \tag{13.15}
\]
and
\[
 \|(A_i-\lambda)^{-1}\|_{\mu'}
 \leq
 \frac{K_0}{1-\eta_{\rm shell}-\eta_{\rm cmp}}.
 \tag{13.16}
\]
\end{theorem}

\begin{proof}
Equations (13.9)--(13.14) satisfy the right-parametrix lemma of V9 with
total error
\(\widetilde E_i^{\rm shell}+\widetilde E_i^{\rm cmp}\).
The Neumann series gives (13.15)--(13.16).
\end{proof}

If the physical kernel and completion agree exactly on the support of
every \(\eta_j\), then \(\eta_{\rm cmp}=0\).  The completion may differ
arbitrarily near the auxiliary boundary, since that difference is cut off
before the physical rows are used.

\subsection{Periodic flat completion}

Take an auxiliary periodic block with a commuting flat holonomy and
\(m=1\).  By V12,
\[
 \widetilde H_{j}^2\succeq I.
 \tag{13.17}
\]
If the completed links lie in the \(d_\infty\)-tube of radius \(1/16\),
then
\[
 \min|\spec\widetilde H_j|\geq\frac12
 \tag{13.18}
\]
and the weighted inverse stock is uniform in the auxiliary block size.

\begin{proposition}[Free completed blocks avoid the boundary mode]
\label{prop:v11v13-free-completion}
A periodic completion satisfying (13.17)--(13.18) supplies the local
inverse in (13.7) without the slab mode of (13.3).  If its links agree
with the physical links throughout the support of \(\eta_j\), its only
parametrix error is the buffered commutator (13.10).
\end{proposition}

\begin{proof}
Periodic closure removes the missing forward link responsible for the
last-row residual in the slab and V12 supplies the positive gap.  Exact
agreement makes (13.6) zero.  The completed-block identity then leaves
only (13.10).
\end{proof}

\subsection{Gauge covariance of a completion}

Let \(J_j\) identify the physical core with the auxiliary block.  A
completion is core gauge covariant if every gauge transformation supported
inside the core has an auxiliary extension \(\widetilde g\) such that
\[
 \widetilde A_j(U^{g})
 =
 \widetilde{\mathcal U}_j(\widetilde g)
 \widetilde A_j(U)
 \widetilde{\mathcal U}_j(\widetilde g)^{-1}
 \tag{13.19}
\]
and \(J_j\mathcal U(g)=\widetilde{\mathcal U}_j(\widetilde g)J_j\) on
the core.

\begin{proposition}[Core Ward identity for the completed parametrix]
\label{prop:v11v13-core-ward}
Under (13.19), the infinitesimal Ward commutator of V11 holds for every
completed inverse and Riesz projector contracted with tests supported
inside the core.  The shell commutator remains an explicit boundary row
and is not part of the bulk anomaly.
\end{proposition}

\begin{proof}
Differentiate (13.19) and apply the V11 resolvent and projector argument
on the auxiliary block.  The intertwining identity for \(J_j\) transfers
the result to the core.  Multiplication by the cutoff produces
\([A_i,\eta_j]\), which is exactly the shell term in (13.10).
\end{proof}

\subsection{The new finite acquisition}

For a physical good block, one must now construct an extension of its
links to an auxiliary periodic or otherwise gapped block such that:
\begin{enumerate}[label=\textnormal{(E\arabic*)}]
\item the extension agrees on a declared buffered core;
\item the completed Wilson kernel lies in a certified gap tube such as
(13.18), or has its own interval inverse certificate;
\item the extension is core gauge covariant;
\item the transition shell has the overlap and buffer stocks needed to
make (13.10) contractive;
\item the distribution of blocks for which no such extension exists has
a vanishing local intensity or a controlled \(L^q\) complement.
\end{enumerate}

\begin{firewall}[Topology and extension]
\label{fw:v11v13-topology}
A curved or topologically nontrivial gauge patch need not extend to a
flat periodic block while agreeing on a large core.  The proposition above
does not assert such an extension.  When flat completion is obstructed,
one needs a sector-compatible gapped completion with an independently
certified base spectrum.  The transition maps of its determinant line
remain part of acquisition A4.
\end{firewall}

\begin{assessment}[V13 advance]
\label{ass:v11v13}
The naive Dirichlet route has been decisively ruled out by an exact
boundary zero mode and an exponentially small finite-slab singular value.
The geometric-resolvent construction has been repaired: a gapped
auxiliary completion may replace the Dirichlet inverse, with separate
shell and completion-mismatch errors.  The next finite problem is the
extension theorem on a good gauge block and its compatibility with local
gauge covariance and sector topology.
\end{assessment}


\section{Tree-gauge extension of a good contractible block}
\label{sec:v11v14}

V13 replaced the false Dirichlet gap by a gapped auxiliary completion.
This note constructs such a completion on a finite contractible block when
all plaquette defects are sufficiently small.  The constants expose the
cost of the loop fillings.

\subsection{Rooted tree gauge}

Let \(K\) be a finite connected cellular block with oriented edge set
\(E(K)\), plaquette set \(\mathcal P(K)\), and compact unitary link field
\(U\).  Fix a root \(o\) and a spanning tree \(\mathcal T\).  For each
vertex \(x\), let \(P_U(x)\) be the ordered product of links along the
unique tree path from \(o\) to \(x\), with inverse links used against the
orientation.  Define the tree-gauge field
\[
 \widehat U_{xy}:=
 P_U(x)U_{xy}P_U(y)^{-1}.
 \tag{14.1}
\]

\begin{lemma}[Tree edges and gauge covariance]
\label{lem:v11v14-tree}
Every tree edge has \(\widehat U_{xy}=I\).  Under
\[
 U^g_{xy}=g_xU_{xy}g_y^{-1},
 \tag{14.2}
\]
one has
\[
 P_{U^g}(x)=g_oP_U(x)g_x^{-1},
 \qquad
 \widehat U^g_{xy}
 =
 g_o\widehat U_{xy}g_o^{-1}.
 \tag{14.3}
\]
Thus tree gauge is unique up to one global conjugation at the root.
\end{lemma}

\begin{proof}
If \(xy\) is a forward tree edge, then
\(P_U(y)=P_U(x)U_{xy}\), and (14.1) is the identity.  The reverse
orientation follows by inversion.  Along a tree path, the internal
\(g_x^{-1}g_x\) factors telescope, giving the first identity in (14.3).
Substitution into (14.1) gives the second.
\end{proof}

\subsection{Filling the fundamental cycles}

For every non-tree edge \(e\), its union with the two tree paths to the
root forms a fundamental loop \(C_e\).  Assume \(K\) is cellularly
simply connected and fix a plaquette filling of \(C_e\) containing at
most \(A_e\) plaquettes.  Set
\[
 A_*:=\max_{e\notin\mathcal T}A_e.
 \tag{14.4}
\]
Let
\[
 r_p(U):=\|I-U_p\|,
 \qquad
 \max_{p\in\mathcal P(K)}r_p(U)\leq\epsilon.
 \tag{14.5}
\]

\begin{lemma}[Product defect bound]
\label{lem:v11v14-product}
For unitary matrices \(V_1,\ldots,V_N\),
\[
 \|V_1\cdots V_N-I\|
 \leq\sum_{j=1}^N\|V_j-I\|.
 \tag{14.6}
\]
The same bound holds when any \(V_j\) is replaced by a unitary conjugate
or its inverse.
\end{lemma}

\begin{proof}
The telescoping identity
\[
 V_1\cdots V_N-I
 =
 \sum_{j=1}^N
 V_1\cdots V_{j-1}(V_j-I)
\]
with the appropriate right tails inserted, or induction using
\(AB-I=A(B-I)+(A-I)\), gives (14.6), because all unitary factors have
norm one.  Conjugation preserves the norm and
\(\|V^{-1}-I\|=\|V-I\|\).
\end{proof}

\begin{theorem}[Nonabelian tree-gauge Poincar\'e bound]
\label{thm:v11v14-poincare}
Under (14.4)--(14.5),
\[
 \max_{e\in E(K)}\|\widehat U_e-I\|
 \leq A_*\epsilon.
 \tag{14.7}
\]
\end{theorem}

\begin{proof}
The assertion is zero on tree edges.  For a non-tree edge, tree gauge
makes every other tree edge in the fundamental loop equal to the identity,
so \(\widehat U_e\), or its inverse, is the loop holonomy.  A lattice
nonabelian Stokes elimination of the fixed plaquette filling writes that
holonomy as an ordered product of \(A_e\) plaquette holonomies, each
parallel transported to the root.  This identity follows directly by
successively removing a boundary plaquette: its shared edge cancels with
the current boundary word, leaving a conjugated plaquette word and the
boundary of the remaining filling.  Each factor has defect at most
\(\epsilon\) by (14.5) and unitary conjugation.  The product defect lemma
gives
\[
 \|\widehat U_e-I\|\leq A_e\epsilon\leq A_*\epsilon.
\]
\end{proof}

The estimate is deliberately literal.  It does not use cancellations
between plaquettes and therefore remains valid for a nonabelian group.

\subsection{Periodic auxiliary completion}

Embed a buffered copy of \(K\) into a larger finite periodic cell
\(\widetilde K\).  Define a tree-gauge completion by
\[
 \widehat{\widetilde U}_e=
 \begin{cases}
 \widehat U_e,&e\text{ belongs to the declared agreement core},\\
 I,&\text{otherwise}.
 \end{cases}
 \tag{14.8}
\]
Then
\[
 \sup_{e\in E(\widetilde K)}
 \|\widehat{\widetilde U}_e-I\|
 \leq A_*\epsilon.
 \tag{14.9}
\]
Extend the vertex matrices \(P_U(x)\) from the agreement core to arbitrary
unitaries \(\widetilde P_U(x)\) on \(\widetilde K\), and transform back:
\[
 \widetilde U_{xy}
 :=
 \widetilde P_U(x)^{-1}
 \widehat{\widetilde U}_{xy}
 \widetilde P_U(y).
 \tag{14.10}
\]
On every physical core edge, (14.10) equals the original \(U_{xy}\).

\begin{theorem}[Gapped periodic completion of a good block]
\label{thm:v11v14-completion}
Let \(0<m<2\) and
\(\gamma_m=\min\{m,2-m\}\).  If
\[
 A_*\epsilon\leq\frac{\gamma_m}{16},
 \tag{14.11}
\]
then the periodic Wilson kernel built from \(\widetilde U\) satisfies
\[
 \min|\spec H_m(\widetilde U)|
 \geq\frac{\gamma_m}{2}.
 \tag{14.12}
\]
For \(m=1\), the sufficient condition is
\[
 A_*\epsilon\leq\frac1{16}.
 \tag{14.13}
\]
The completion agrees exactly with the physical operator on all rows
whose hopping neighbourhood lies in the declared agreement core.
\end{theorem}

\begin{proof}
The operator for \(\widetilde U\) is unitarily gauge equivalent to the
one for \(\widehat{\widetilde U}\).  By (14.9)--(14.11), the latter lies
in the V12 half-gap tube around the identity link field.  V12 therefore
gives (14.12).  Gauge transformation preserves the spectrum.
Agreement of links on the range-one hopping neighbourhood gives equality
of the corresponding operator rows.
\end{proof}

\begin{corollary}[Exact removal of the completion mismatch]
\label{cor:v11v14-mismatch}
Choose the V13 cutoff \(\eta_j\) so that its hopping neighbourhood is
contained in the agreement core.  Then
\[
 \Delta_{j,i}=\eta_{j,i}(A_i-\widetilde A_{j,i})=0,
 \tag{14.14}
\]
and the completed-block parametrix has only the shell error.
\end{corollary}

\begin{proof}
The last statement of the completion theorem says that every row selected
by \(\eta_j\) is identical in the physical and completed kernels.
This is exactly (14.14).
\end{proof}

\subsection{Core gauge covariance}

The construction (14.8) is equivariant under global conjugation:
replacing every retained core link by
\(h\widehat U_eh^{-1}\) replaces the completion by its conjugate, because
the identity collar is fixed.  Let a physical gauge transformation \(g\)
be supported on the agreement core.  Extend it by the identity to
\(\widetilde g\) on the auxiliary sites.

\begin{proposition}[Covariant completion up to root conjugation]
\label{prop:v11v14-covariance}
The completed kernels associated with \(U\) and \(U^g\) are related by a
unitary auxiliary gauge transformation.  On the physical core this
transformation restricts to \(g\).  Consequently the core kernel,
resolvent, and Riesz-projector Ward commutators of V11 hold exactly.
\end{proposition}

\begin{proof}
By (14.3), the tree-gauge core data for \(U^g\) are the data for \(U\)
conjugated by \(g_o\).  Equivariance of (14.8) extends this relation to
the entire auxiliary block.  Transforming back with the extended tree
matrices in (14.10) produces an auxiliary gauge transformation whose
restriction to the agreement core is \(g\); outside the core its value
depends on the chosen extension but remains unitary.  Gauge covariance of
the Wilson kernel gives unitary equivalence.  Differentiation and the V11
projector theorem give the Ward commutators.
\end{proof}

The construction need not choose a globally smooth tree gauge over all
gauge configurations.  Its local use on a fixed good block suffices for
the inverse certificate.  Determinant-line transition data are still
needed between such choices.

\subsection{Quantitative cost and block size}

For a rectangular block of graph diameter \(L\), a simple coordinate
tree and rectangular filling give
\[
 A_*\leq C_dL^2
 \tag{14.15}
\]
with a dimension-dependent combinatorial constant.  Thus (14.11) asks
for
\[
 \epsilon\leq\frac{\gamma_m}{16C_dL^2}.
 \tag{14.16}
\]
This deterioration is the price of using only the maximum plaquette
defect and a telescoping filling.

\begin{proposition}[Fixed-block local tail]
\label{prop:v11v14-local-tail}
If \(L\) is fixed and the single-plaquette tail satisfies
\[
 \sup_p\nu_i\{r_p>\epsilon\}\leq\theta_i,
 \tag{14.17}
\]
then the probability that a fixed \(L\)-block fails (14.5) is at most
\[
 N_{\mathcal P}(L)\theta_i,
 \tag{14.18}
\]
where \(N_{\mathcal P}(L)\) is the finite number of plaquettes in the
block.  Hence the local failure intensity tends to zero when
\(\theta_i\to0\), without a total-volume factor.
\end{proposition}

\begin{proof}
The block fails only if at least one of its
\(N_{\mathcal P}(L)\) plaquettes fails.  Apply the union bound inside the
fixed block.
\end{proof}

If \(L=L_i\to\infty\), both the filling cost in (14.16) and the local
entropy in (14.18) must be tracked.  No limit may hold with a fixed
threshold merely because each finite block is extendable.

\subsection{Integration with the mobility route}

\begin{theorem}[Good-block completion card]
\label{thm:v11v14-good-block-card}
Assume a buffered cover of the physical complement of the defect islands
has fixed block geometry and bounded overlap.  Suppose:
\begin{enumerate}[label=\textnormal{(\roman*)}]
\item every retained block satisfies (14.5) with
\(A_*\epsilon\leq\gamma_m/16\);
\item its auxiliary collar is chosen as in (14.8);
\item the shell buffer makes the V13 commutator error contractive; and
\item the intensity of blocks violating (i) tends to zero with the
\(L^q\) complement bound of V7.
\end{enumerate}
Then the retained physical bulk has the weighted inverse required in V8,
and its expected artificial local Ward edge tends to zero under the
remaining outer-intertwiner hypotheses of V7.
\end{theorem}

\begin{proof}
The completion theorem supplies a uniform auxiliary gap and weighted
inverse.  Exact core agreement removes the completion mismatch.  The V13
completed-block geometric resolvent closes the physical bulk inverse.
V8 then localizes near modes at the discarded block islands, and V7 closes
their expected local Ward contribution.  Assumption (iv) controls the
complement.
\end{proof}

\begin{firewall}[What is still conditional]
\label{fw:v11v14-conditional}
The theorem does not prove the single-plaquette tail (14.17) under the
coupled reduced Standard Model--Einstein density, nor the shell
contraction for a cofinal physical block scale.  It also does not glue
the local determinant-line frames or control noncontractible sector
transitions.  The extension theorem closes the finite geometric step,
not the probability, phase, or gravitational rows.
\end{firewall}

\begin{assessment}[V14 advance]
\label{ass:v11v14}
A good contractible gauge block now has an explicit gapped periodic
completion.  Rooted tree gauge bounds every link by the number of
plaquettes in a fundamental filling, the identity collar preserves that
bound, and the V12 flat tube supplies the chiral gap.  The construction is
core gauge covariant and removes the V13 completion-mismatch row exactly.
V15 is the next mandatory companion audit; it will check for new density,
cluster, or finite-certificate input before the programme continues.
\end{assessment}


\section{V15 mandatory companion audit}
\label{sec:v11v15}

The third five-version boundary of this compact series requires a fresh
inspection of the active companion files.

\begin{audit}[Files inspected]
\label{audit:v11v15-files}
The exact records are
\[
\begin{array}{c|r|r}
\text{file}&\text{lines}&\text{bytes}\\ \hline
\texttt{Renewal\_Yang\_Mills\_Mass\_Gap\_Research\_Notes\_V26.tex}
&11348&385324\\
\texttt{Renewal\_NS\_Stability\_Research\_Notes\_V26.tex}
&5319&278283.
\end{array}
\]
Their SHA--256 values are
\begin{center}
\small\ttfamily
FCFE0FB979D045C358449F2D207EB0416031C913DCE77FEA515B6B8355B8C19C
\end{center}
and
\begin{center}
\small\ttfamily
82C887F8C2C69E4F28FAD7C3DC8DE5B9099CB5A4BF431EBA1FFF3E91935978AD.
\end{center}
The NS note is unchanged.  The YM note has advanced from V22 to V26.
\end{audit}

\subsection{New Yang--Mills material}

YM V23 computes exact harmonic routers at all sixteen parity corners and
keeps the toron corner in the same invertible vertical system.  It proves
that two endpoint routers agree at only two of the corners, so freezing a
router through a regulator comparison would omit a genuine response row.
It also gives a finite residual-box continuation criterion.

YM V24 performs exact dual-colour compression over a rational
\(\mathfrak{su}(3)\) basis.  Completeness, adjoint Casimir, cubic, and
contracted quartic identities reduce the colour contraction before the
spatial inverse is evaluated.  The note explicitly keeps imaginary
ordered traces until their conjugate owner terms have been assembled.

YM V25 imports the residual-inverse and weighted-conjugation criteria from
this SM stream and executes them as an adaptive atlas seed.  Exact corner
inverses and induced row stocks give nonempty rational boxes.  It correctly
warns that seed boxes are not a cover until every accepted box and overlap
has been recorded.

YM V26 proves that a commuting one-colour toron lift is exactly flat,
that translated stationary-center derivatives vanish, and that product
Haar measure has unit Jacobian on that path.  It also proves that these
facts do not force the finite determinant response to vanish: a global
toron holonomy may retain a finite-volume effective potential.

\begin{audit}[Accepted transfers]
\label{audit:v11v15-accepted}
The audit changes the SM direction in four ways.
\begin{enumerate}[label=\textnormal{(\roman*)}]
\item The V11--V14 block certificates should be executed as an adaptive
residual atlas with explicit coverage and overlap records, rather than
inferred from a few successful seed backgrounds.
\item Any physical constraint or harmonic lift must move with the
background.  Equality of long-wave two-jets does not permit the router to
be frozen.
\item Exact colour compression may precede the spatial inverse, but it
does not replace the determinant-line or Ward calculation.
\item On the commuting flat backgrounds certified in V12, the classical
curvature and translated Haar rows simplify, but the fermion determinant
as a function of global holonomy must be retained and estimated.
\end{enumerate}
\end{audit}

\begin{audit}[No cluster or density transfer]
\label{audit:v11v15-no-cluster}
YM V23--V26 contains no random defect-cluster tail, no nonpercolation
theorem, and no coupled Standard Model--Einstein reduced-density estimate.
NS V26 likewise supplies no applicable lattice gauge or determinant
bound.  Thus V7--V9 probability assumptions remain open.
\end{audit}

\subsection{Post-audit direction}

The new toron result intersects the explicit flat-holonomy certificate of
V12.  It is therefore a closer upstream target than repeating the open
cluster estimate.

\begin{decision}[V16 target]
\label{dec:v11v15-direction}
For the gapped Wilson regulator on a periodic flat-holonomy background,
V16 will:
\begin{enumerate}[label=\textnormal{(\alph*)}]
\item factor the finite determinant into momentum and representation-weight
blocks;
\item prove exact large-gauge periodicity by momentum relabelling;
\item use the complex strip to bound the holonomy dependence of the
normalized logarithmic determinant by an exponential alias tail;
\item separate this positive vectorlike modulus statement from the
massless physical chiral determinant and its determinant-line phase.
\end{enumerate}
\end{decision}

\begin{firewall}[Audit claim boundary]
\label{fw:v11v15-boundary}
No companion conclusion constructs the SM determinant-line
trivialization, proves the reduced Einstein probability body, or closes
the continuum.  The V16 calculation will address a gapped regulator
determinant on flat backgrounds only.
\end{firewall}

\begin{assessment}[V15 audit result]
\label{ass:v11v15}
The required audit is complete.  It preserves the block-completion route
and adds a distinct holonomy ledger.  The next substantive note estimates
the finite toron dependence instead of deleting it.
\end{assessment}


\section{Flat-holonomy factorization and determinant alias tail}
\label{sec:v11v16}

V15 requires the finite toron determinant to be retained rather than
declared zero.  This note evaluates the positive full Wilson-regulator
determinant and proves an exponential holonomy bound for its normalized
logarithm.

\subsection{Periodic momenta and representation weights}

Let the periodic lattice have side lengths
\[
 L=(L_1,L_2,L_3,L_4),\qquad
 |\Lambda_L|:=\prod_{\mu=1}^4L_\mu,
 \qquad
 L_{\min}:=\min_\mu L_\mu.
 \tag{16.1}
\]
Allow fermion boundary phases \(\beta_\mu\), with
\(\beta_4=\pi\) for an antiperiodic temporal boundary condition if
desired.

Let \(R\) be a unitary representation of the Standard Model gauge group,
restricted to a commuting maximal torus.  Write \(\mathcal W_R\) for its
finite multiset of weights.  A global holonomy \(a=(a_\mu)\) is represented
by the constant links
\[
 V_\mu=\exp\left(\frac{i\,a_\mu}{L_\mu}\right).
 \tag{16.2}
\]
On the weight \(w\), the momentum is
\[
 q_{\mu,n,w}(a)
 =
 \frac{2\pi n_\mu+\beta_\mu+w(a_\mu)}{L_\mu},
 \qquad
 0\leq n_\mu<L_\mu.
 \tag{16.3}
\]

Define
\[
 F_m(k)
 :=
 \sum_{\mu=1}^4\sin^2k_\mu+
 \left(\sum_{\mu=1}^4(1-\cos k_\mu)-m\right)^2.
 \tag{16.4}
\]
V12 proves
\[
 F_m(k)\geq\gamma_m^2,\qquad
 \gamma_m=\min\{m,2-m\}>0.
 \tag{16.5}
\]

\subsection{Exact determinant factorization}

\begin{lemma}[Spin determinant]
\label{lem:v11v16-spin}
For the four-component Wilson block
\[
 A_m(k)=i\sum_\mu\gamma_\mu\sin k_\mu+W(k)-m,
\]
one has
\[
 \det_{\rm spin}A_m(k)=F_m(k)^2>0.
 \tag{16.6}
\]
The Hermitian block \(H_m(k)=\gamma_5A_m(k)\) has the same determinant.
\end{lemma}

\begin{proof}
Let \(s=(\sin k_\mu)\) and \(b=W(k)-m\).  The Clifford relations show
that \(\gamma\cdot s\) has eigenvalues \(+|s|\) and \(-|s|\), each with
multiplicity two.  Thus \(A_m\) has eigenvalues
\(b+i|s|\) and \(b-i|s|\), each twice, and
\[
 \det A_m=(b^2+|s|^2)^2=F_m^2.
\]
In four dimensions \(\gamma_5\) has two eigenvalues \(+1\) and two
eigenvalues \(-1\), so \(\det\gamma_5=1\).
\end{proof}

\begin{theorem}[Finite flat-holonomy determinant]
\label{thm:v11v16-factorization}
For one full Wilson spinor in representation \(R\),
\[
 \det H_{m,L,R}(a)
 =
 \prod_{w\in\mathcal W_R}
 \prod_{n_1=0}^{L_1-1}\cdots\prod_{n_4=0}^{L_4-1}
 F_m(q_{n,w}(a))^2.
 \tag{16.7}
\]
It is strictly positive.  Its normalized logarithm is
\[
 \mathcal F_{m,L,R}(a)
 :=
 \frac1{|\Lambda_L|}\log\det H_{m,L,R}(a)
 =
 \sum_{w\in\mathcal W_R}
 \frac1{|\Lambda_L|}
 \sum_n g_m(q_{n,w}(a)),
 \tag{16.8}
\]
where
\[
 g_m(k):=2\log F_m(k).
 \tag{16.9}
\]
\end{theorem}

\begin{proof}
The commuting holonomies are simultaneously diagonal on the representation
weights, while Fourier transformation diagonalizes translations.
V12 shows that the resulting blocks are \(H_m(q_{n,w}(a))\).
Apply the spin determinant lemma to every block and multiply.  Positivity
follows from (16.5), and taking the real logarithm gives (16.8).
\end{proof}

\subsection{Large-gauge periodicity}

Let \(h\) be a cocharacter of the compact gauge torus, normalized so that
\[
 w(h)\in\mathbb Z
 \quad\text{for every representation weight }w.
 \tag{16.10}
\]

\begin{proposition}[Exact large-gauge invariance]
\label{prop:v11v16-large-gauge}
For every coordinate \(\mu\),
\[
 \mathcal F_{m,L,R}
 (a_1,\ldots,a_\mu+2\pi h,\ldots,a_4)
 =
 \mathcal F_{m,L,R}(a).
 \tag{16.11}
\]
\end{proposition}

\begin{proof}
The shift changes (16.3) by
\(2\pi w(h)/L_\mu\).  Relabel
\(n_\mu\mapsto n_\mu+w(h)\) modulo \(L_\mu\) in the finite momentum
product.  The multiset of blocks is unchanged.
\end{proof}

This is exact finite-cutoff periodicity, not a limiting Ward argument.

\subsection{Analytic logarithm}

The kernel \(H_m(k)\) is a finite Laurent polynomial.  By the V12 gap and
a Neumann perturbation, there is \(\rho_0>0\) such that
\[
 \det H_m(k+iy)\ne0
 \quad\text{when}\quad
 \|y\|_\infty\leq\rho_0.
 \tag{16.12}
\]
For example, the hopping-stock estimate permits any
\[
 0<\rho_0<
 \log\left(1+\frac{\gamma_m}{16}\right).
 \tag{16.13}
\]
On a slightly smaller strip, the positive real-axis determinant has a
periodic analytic logarithm agreeing with (16.9).

\begin{lemma}[Weighted Fourier stock]
\label{lem:v11v16-fourier-stock}
For every \(0<\rho<\rho_0\), the Fourier series
\[
 g_m(k)=\sum_{p\in\mathbb Z^4}\widehat g_m(p)e^{ip\cdot k}
 \tag{16.14}
\]
has finite Wiener stock
\[
 M_{m,\rho}:=
 \sum_{p\in\mathbb Z^4}
 |\widehat g_m(p)|e^{\rho|p|_1}<\infty.
 \tag{16.15}
\]
\end{lemma}

\begin{proof}
Choose \(\rho<\rho'<\rho_0\).  The analytic logarithm is bounded by some
\(B_{\rho'}\) on the compact closed strip of width \(\rho'\).
Shifting each Fourier contour to
\(-i\rho'\sgn(p_\mu)\) gives
\[
 |\widehat g_m(p)|\leq B_{\rho'}e^{-\rho'|p|_1}.
\]
After multiplication by \(e^{\rho|p|_1}\), the product of four geometric
series converges.
\end{proof}

\subsection{Exact alias formula}

For a phase \(\alpha\in\mathbb R^4\), define the shifted grid average
\[
 \mathcal A_Lg_m(\alpha)
 :=
 \frac1{|\Lambda_L|}
 \sum_{0\leq n_\mu<L_\mu}
 g_m\left(
 \frac{2\pi n_1+\alpha_1}{L_1},\ldots,
 \frac{2\pi n_4+\alpha_4}{L_4}
 \right).
 \tag{16.16}
\]

\begin{theorem}[Holonomy alias identity]
\label{thm:v11v16-alias}
One has the absolutely convergent identity
\[
 \mathcal A_Lg_m(\alpha)
 =
 \sum_{q\in\mathbb Z^4}
 \widehat g_m(q_1L_1,\ldots,q_4L_4)e^{iq\cdot\alpha}.
 \tag{16.17}
\]
Consequently,
\[
 \left|
 \mathcal A_Lg_m(\alpha)-\widehat g_m(0)
 \right|
 \leq
 M_{m,\rho}e^{-\rho L_{\min}}
 \tag{16.18}
\]
uniformly in \(\alpha\).
\end{theorem}

\begin{proof}
Insert (16.14) into (16.16).  Absolute convergence permits interchange of
the finite average and the Fourier sum.  In each coordinate,
\[
 \frac1{L_\mu}\sum_{n_\mu=0}^{L_\mu-1}
 e^{2\pi i p_\mu n_\mu/L_\mu}
 =
 \begin{cases}
 1,&L_\mu\mid p_\mu,\\
 0,&L_\mu\nmid p_\mu.
 \end{cases}
\]
Thus only \(p_\mu=q_\mu L_\mu\) remain, and their shifted phase is
\(e^{iq_\mu\alpha_\mu}\), proving (16.17).  For every nonzero \(q\),
\[
 |(q_1L_1,\ldots,q_4L_4)|_1\geq L_{\min}.
\]
Therefore
\[
\begin{split}
 \sum_{q\ne0}
 |\widehat g_m(q_1L_1,\ldots,q_4L_4)|
 &\leq
 e^{-\rho L_{\min}}
 \sum_{q\ne0}
 |\widehat g_m(qL)|e^{\rho|qL|_1}\\
 &\leq M_{m,\rho}e^{-\rho L_{\min}}.
\end{split}
\]
\end{proof}

\begin{corollary}[Normalized determinant loses its toron dependence]
\label{cor:v11v16-normalized}
For the representation \(R\),
\[
 \left|
 \mathcal F_{m,L,R}(a)
 -
 |\mathcal W_R|\widehat g_m(0)
 \right|
 \leq
 |\mathcal W_R|M_{m,\rho}e^{-\rho L_{\min}}.
 \tag{16.19}
\]
In particular the normalized logarithmic determinant becomes uniformly
independent of the global commuting holonomy as \(L_{\min}\to\infty\).
\end{corollary}

\begin{proof}
Apply (16.18) to
\(\alpha=\beta+w(a)\) for each weight and sum the finite multiset.
\end{proof}

For the unnormalized logarithm the right side is multiplied by
\(|\Lambda_L|\).  It still tends to zero along comparable boxes because
polynomial volume growth is dominated by the exponential.  For arbitrary
anisotropic sequences one must explicitly require
\[
 |\Lambda_L|e^{-\rho L_{\min}}\longrightarrow0.
 \tag{16.20}
\]

\subsection{Holonomy derivatives}

Let \(\|w\|\) be the operator norm of the weight pairing and set
\(W_R:=\max_{w\in\mathcal W_R}\|w\|\).

\begin{proposition}[Derivative alias tail]
\label{prop:v11v16-derivative}
For every fixed integer \(r\geq1\), there is an explicit constant
\(C_{r,\rho}\) such that
\[
 \sup_a
 \|D_a^r\mathcal F_{m,L,R}(a)\|
 \leq
 |\mathcal W_R|W_R^r
 C_{r,\rho}M_{m,\rho}
 e^{-\rho L_{\min}/2}.
 \tag{16.21}
\]
Thus every fixed holonomy derivative of the normalized gapped regulator
determinant also vanishes exponentially.
\end{proposition}

\begin{proof}
Differentiate (16.17).  A derivative supplies at most
\(\|w\|\,|q|_1\) per order.  With \(p=qL\),
\[
 |q|_1^r\leq L_{\min}^{-r}|p|_1^r.
\]
For \(p\ne0\),
\[
 |p|_1^re^{-\rho|p|_1}
 \leq
 \left(\sup_{t\geq0}t^re^{-\rho t/2}\right)
 e^{-\rho|p|_1/2}
 \leq
 \left(\frac{2r}{e\rho}\right)^r
 e^{-\rho L_{\min}/2}.
\]
Use the weighted stock (16.15), discard the additional
\(L_{\min}^{-r}\leq1\), and sum the weights.  This gives (16.21) with
\(C_{r,\rho}=(2r/(e\rho))^r\).
\end{proof}

\subsection{What the estimate does and does not cover}

The full Wilson regulator determinant in (16.7) is vectorlike and
positive.  The same alias proof applies to a massive vectorlike charged
fermion block on a Higgs chart whenever its scalar determinant has a
uniform nonzero complex strip.

\begin{firewall}[Massless chiral determinant]
\label{fw:v11v16-chiral}
A physical massless overlap operator has low-energy zeros, so its
logarithm need not have the analytic strip used in (16.15).  Its toron
dependence may decay only polynomially or may encode a finite-volume
effective potential.  A left-handed determinant is also a section of a
determinant line; the positive scalar factorization (16.7) neither chooses
its phase nor proves trivial holonomy.  Therefore (16.19)--(16.21) may be
used for the gapped regulator modulus and massive vectorlike blocks, not
as a proof of the full chiral measure.
\end{firewall}

\begin{remark}[Order of operations]
\label{rem:v11v16-order}
The finite toron factor is retained in (16.17) before its alias tail is
estimated.  This is the determinant analogue of forming the complete Ward
sum before a low/edge split.  Setting the toron term to zero at finite
volume would lose the nonzero alias modes.
\end{remark}

\begin{assessment}[V16 advance]
\label{ass:v11v16}
The flat-holonomy regulator determinant is now exact and controlled.  It
is a strictly positive product over momenta and representation weights,
has exact large-gauge periodicity, and its normalized holonomy dependence
and fixed derivatives vanish exponentially with the shortest side.
The surviving hard row is precisely identified: the massless chiral
determinant-line phase and its low modes, together with the coupled
density and Einstein constraints, are outside the gapped alias theorem.
\end{assessment}


\section{Massless overlap low modes on the holonomy torus}
\label{sec:v11v17}

The exponential alias theorem of V16 relies on a uniform analytic
logarithm.  The massless physical overlap operator does not have one.
This note computes its free symbol exactly and isolates the low-mode
singularity.

\subsection{The massless overlap symbol}

Retain \(0<m<2\), and write
\[
 A_m(k)=i\gamma\cdot s(k)+b_m(k),
 \qquad
 s_\mu(k)=\sin k_\mu,
 \qquad
 b_m(k)=W(k)-m,
 \tag{17.1}
\]
\[
 r_m(k):=\sqrt{F_m(k)}
 =
 \sqrt{|s(k)|^2+b_m(k)^2}.
 \tag{17.2}
\]
The V12 gap gives \(r_m(k)\geq\gamma_m>0\).  In dimensionless
normalization the free massless overlap operator is
\[
 D_{\rm ov,m}(k)
 :=
 I+\frac{A_m(k)}{r_m(k)}.
 \tag{17.3}
\]

\begin{theorem}[Exact spin determinant of the massless overlap block]
\label{thm:v11v17-overlap-det}
For every real \(k\),
\[
 \Delta_m(k)
 :=
 \det_{\rm spin}D_{\rm ov,m}(k)
 =
 4\left(1+\frac{b_m(k)}{r_m(k)}\right)^2
 \geq0.
 \tag{17.4}
\]
Moreover,
\[
 \Delta_m(k)=0
 \quad\Longleftrightarrow\quad
 k=0\pmod{2\pi}.
 \tag{17.5}
\]
\end{theorem}

\begin{proof}
The matrix \(A_m(k)\) has eigenvalues
\(b_m+i|s|\) and \(b_m-i|s|\), each with multiplicity two.  Therefore
the four eigenvalues of (17.3) occur in two conjugate pairs and
\[
\begin{split}
 \Delta_m(k)
 &=
 \left[
 \left(1+\frac{b_m}{r_m}\right)^2
 +\frac{|s|^2}{r_m^2}
 \right]^2\\
 &=
 \left[
 1+\frac{2b_m}{r_m}
 +\frac{b_m^2+|s|^2}{r_m^2}
 \right]^2
 =
 4\left(1+\frac{b_m}{r_m}\right)^2.
\end{split}
\]
This proves (17.4).  Equality requires
\(b_m=-r_m\), hence \(s=0\) and \(b_m<0\).
All components of \(s\) vanish only at Brillouin corners.  At a corner
with \(n\) components equal to \(\pi\),
\(b_m=2n-m\).  Since \(0<m<2\), it is negative only for \(n=0\).
This proves (17.5).
\end{proof}

Thus all Wilson doublers are nonzero while the physical corner remains
massless.

\subsection{Exact local factorization}

On a neighbourhood of zero where \(b_m(k)<0\), the identity
\[
 (r_m+b_m)(r_m-b_m)=|s|^2
 \tag{17.6}
\]
gives
\[
 \Delta_m(k)
 =
 \frac{4|s(k)|^4}
 {r_m(k)^2(r_m(k)-b_m(k))^2}.
 \tag{17.7}
\]
Choose
\[
 \mathcal N_m:=
 \left\{
 k:\ |k_\mu|\leq\frac{\pi}{2},\
 |k|^2\leq m
 \right\}.
 \tag{17.8}
\]
Since \(1-\cos x\leq x^2/2\),
\[
 W(k)\leq\frac{|k|^2}{2}\leq\frac m2,
 \qquad
 b_m(k)\leq-\frac m2
 \quad(k\in\mathcal N_m).
 \tag{17.9}
\]

\begin{theorem}[Fourth-order physical zero]
\label{thm:v11v17-fourth-order}
There are explicit positive constants
\[
 c_m:=\frac{16}{\pi^4R_m^4},
 \qquad
 C_m:=\frac{64}{m^4},
 \qquad
 R_m:=\frac{3m}{2}+\sqrt m,
 \tag{17.10}
\]
such that
\[
 c_m|k|^4
 \leq
 \Delta_m(k)
 \leq
 C_m|k|^4
 \quad(k\in\mathcal N_m).
 \tag{17.11}
\]
In particular the full Dirac spin determinant has an exact zero of total
order four at the physical corner.
\end{theorem}

\begin{proof}
On (17.8),
\[
 \frac{2}{\pi}|k_\mu|
 \leq|\sin k_\mu|\leq|k_\mu|,
\]
so
\[
 \frac4{\pi^2}|k|^2\leq|s(k)|^2\leq|k|^2.
 \tag{17.12}
\]
Equation (17.9) gives \(r_m\geq|b_m|\geq m/2\) and
\(r_m-b_m\geq r_m\geq m/2\).  Hence the denominator in (17.7) is at least
\(m^4/16\), yielding the upper bound
\[
 \Delta_m\leq4|k|^4\frac{16}{m^4}
 =\frac{64}{m^4}|k|^4.
\]
Also
\[
 |b_m|\leq m+W\leq\frac{3m}{2},
 \qquad
 |s|\leq\sqrt m,
\]
so \(r_m\leq|b_m|+|s|\leq R_m\) and
\(r_m-b_m=r_m+|b_m|\leq2R_m\).  The denominator in (17.7) is therefore
at most \(4R_m^4\).  Together with the lower bound in (17.12),
\[
 \Delta_m
 \geq
 \frac{4(16/\pi^4)|k|^4}{4R_m^4}
 =
 \frac{16}{\pi^4R_m^4}|k|^4.
\]
\end{proof}

\begin{corollary}[Local logarithmic integrability]
\label{cor:v11v17-log-integrable}
The function \(\log\Delta_m(k)\) is locally integrable on
\(\mathbb T^4\).  More precisely, for sufficiently small
\(0<r<1\),
\[
 \int_{|k|\leq r}|\log\Delta_m(k)|\,dk
 \leq
 C'_m r^4(1+|\log r|)
 \tag{17.13}
\]
for a finite constant \(C'_m\).
\end{corollary}

\begin{proof}
The two-sided estimate (17.11) bounds
\(|\log\Delta_m(k)|\) by a constant plus \(4|\log|k||\) near zero.
In four-dimensional polar coordinates,
\[
 \int_0^r \rho^3|\log\rho|\,d\rho
 =
 O(r^4(1+|\log r|)).
\]
Away from zero the determinant is positive and continuous.
\end{proof}

Local integrability permits an infinite-volume free-energy density.  It
does not restore a uniform analytic strip.

\subsection{Finite flat-holonomy determinant}

Use the shifted momenta \(q_{n,w}(a)\) from V16.  For a full overlap
spinor in representation \(R\),
\[
 \det D_{{\rm ov},L,R}(a)
 =
 \prod_{w\in\mathcal W_R}\prod_n
 \Delta_m(q_{n,w}(a)).
 \tag{17.14}
\]
The same momentum relabelling as in V16 proves exact large-gauge
periodicity.

\begin{proposition}[Exact nodal set]
\label{prop:v11v17-nodal}
The product (17.14) vanishes exactly when there exist \(w\) and \(n\)
such that, in every coordinate,
\[
 2\pi n_\mu+\beta_\mu+w(a_\mu)
 \equiv0\pmod{2\pi L_\mu}.
 \tag{17.15}
\]
At a simple weight occurrence, the corresponding full spin block
vanishes to fourth order in the displaced momentum.
\end{proposition}

\begin{proof}
Every factor is nonnegative, and by (17.5) it vanishes precisely when its
momentum is zero modulo \(2\pi\).  Equation (16.3) gives (17.15).
The order statement is the fourth-order theorem after translation to the
nodal point.
\end{proof}

Antiperiodic time boundary conditions move the nodal point in holonomy
space; they do not remove it uniformly over all holonomies.

\subsection{Failure of the exponential alias hypothesis}

\begin{proposition}[No uniform analytic logarithm]
\label{prop:v11v17-no-strip}
There is no neighbourhood of the full real holonomy torus on which
\[
 a\longmapsto
 \log\det D_{{\rm ov},L,R}(a)
 \tag{17.16}
\]
is finite and analytic for every \(a\), because it diverges to
\(-\infty\) on the nodal set (17.15).  Consequently the exponential
alias estimate of V16 cannot be applied to the massless overlap
determinant without first separating or saturating its zero modes.
\end{proposition}

\begin{proof}
At a nodal point one factor in (17.14) equals zero.  Its real logarithm is
\(-\infty\), and a holomorphic logarithm cannot cross a zero of its
argument.  The V16 proof requires one nonvanishing complex strip for the
entire logarithm, so its hypothesis fails.
\end{proof}

\subsection{A finite low-mode bound away from the nodal set}

Let
\[
 d_{L,R}(a)
 :=
 \min_{w,n}
 \left|
 q_{n,w}(a)
 \right|_{\mathbb T^4},
 \tag{17.17}
\]
where distance is taken to the physical corner.

\begin{proposition}[Lowest-block logarithm]
\label{prop:v11v17-low-log}
If \(0<d_{L,R}(a)\leq r_m\) with the minimizing momentum in
\(\mathcal N_m\), then
\[
 4\log d_{L,R}(a)+\log c_m
 \leq
 \log\Delta_m(q_{\min})
 \leq
 4\log d_{L,R}(a)+\log C_m.
 \tag{17.18}
\]
If the scaled holonomy remains a fixed positive distance from the nodal
condition in (17.15), then
\(d_{L,R}(a)\geq c/L_{\max}\), and the normalized contribution of each
such lowest block is
\[
 O\left(\frac{\log L_{\max}}{|\Lambda_L|}\right).
 \tag{17.19}
\]
\end{proposition}

\begin{proof}
Take logarithms in (17.11) to obtain (17.18).  Under the stated distance
condition, the closest shifted grid point is bounded below by
\(c/L_{\max}\).  Equation (17.18) then has magnitude \(O(\log L_{\max})\);
division by the volume gives (17.19).
\end{proof}

This estimate is not uniform as the holonomy approaches the nodal set.
More importantly, a small normalized free-energy contribution does not
permit deletion of a zero mode in fermionic correlation functions.

\subsection{Chiral section and zero-mode saturation}

For a left-handed theory the finite coefficient is a section
\[
 \mathfrak s_L(a)\in\Det D_{{\rm ov},L,R}(a),
 \tag{17.20}
\]
not a globally chosen square root of (17.14).  At a zero, the determinant
section vanishes but correlation functions with the correct exterior
insertions may remain nonzero by zero-mode saturation.

\begin{firewall}[Modulus, phase, and insertions]
\label{fw:v11v17-section}
The positive full-spin determinant (17.14) controls a vectorlike modulus.
It does not choose the phase of (17.20), prove large-gauge invariance of a
chiral scalarization, or justify deleting the nodal configurations.
The determinant-line transition functions and the exterior-power
zero-mode insertions must be carried together.  Standard Model anomaly
arithmetic constrains their combined transformation but does not by
itself construct the global section.
\end{firewall}

\begin{assessment}[V17 advance]
\label{ass:v11v17}
The low-mode obstruction is now exact.  The massless overlap determinant
has one physical fourth-order spin zero and no doubler zeros; flat
holonomies translate that zero through a precisely described nodal set.
Its logarithm is locally integrable but has no uniform analytic strip, so
the gapped exponential alias theorem cannot be reused.  The next step is
to formulate a zero-mode-saturated determinant-line cylinder coefficient
that remains finite at the nodal set and can be compared across regulator
scales.
\end{assessment}


\section{Zero-mode-saturated cylinder coefficients as complementary minors}
\label{sec:v11v18}

V17 shows that a scalar logarithm is the wrong variable at a massless
toron zero.  The finite Grassmann integral itself is polynomial in the
chiral operator and remains meaningful.  This note identifies its
source coefficients exactly.

\subsection{Finite Grassmann coefficient}

Let \(V\) and \(W\) be complex vector spaces of equal dimension \(n\),
with oriented bases \(\psi_1,\ldots,\psi_n\) and
\(\bar\psi_1,\ldots,\bar\psi_n\).  Let
\[
 D:V\longrightarrow W
 \tag{18.1}
\]
have matrix entries \(D_{ji}\).  Define the Berezin polynomial
\[
 \mathcal Z_D(\bar\eta,\eta)
 :=
 \int d\bar\psi\,d\psi\,
 \exp\left(
 -\bar\psi D\psi+\bar\eta\psi+\bar\psi\eta
 \right).
 \tag{18.2}
\]
The measure order fixes all signs once and for all.

For subsets \(I,J\subset\{1,\ldots,n\}\) with
\(|I|=|J|=r\), let
\[
 \Psi_I:=\prod_{i\in I}\psi_i,
 \qquad
 \bar\Psi_J:=\prod_{j\in J}\bar\psi_j
 \tag{18.3}
\]
in increasing order.

\begin{theorem}[Complementary-minor saturation identity]
\label{thm:v11v18-minor}
There is an orientation sign \(\varepsilon(I,J)\in\{\pm1\}\), determined
only by the declared Berezin order, such that
\[
 \int d\bar\psi\,d\psi\,
 e^{-\bar\psi D\psi}\,
 \Psi_I\bar\Psi_J
 =
 \varepsilon(I,J)
 \det D_{J^c,I^c}.
 \tag{18.4}
\]
The right side is a polynomial in the entries of \(D\), so it is defined
and continuous even when \(D\) is singular.
\end{theorem}

\begin{proof}
Expand
\[
 e^{-\bar\psi D\psi}
 =
 \sum_{k=0}^n\frac{(-1)^k}{k!}(\bar\psi D\psi)^k.
\]
The Berezin integral selects the coefficient containing every
\(\psi_i\) and every \(\bar\psi_j\) exactly once.  Since the insertion
already contains the indices in \(I\) and \(J\), only the term
\(k=n-r\) contributes, and it may use only columns \(I^c\) and rows
\(J^c\).  Antisymmetrizing the \(n-r\) selected matrix entries is exactly
\(\det D_{J^c,I^c}\).  Reordering the Grassmann variables into the
declared measure order supplies the fixed sign \(\varepsilon(I,J)\).
\end{proof}

\subsection{Cancellation of the inverse pole}

When \(D\) is invertible, Wick contraction gives the same coefficient as
a determinant times a propagator minor.

\begin{corollary}[Jacobi cancellation]
\label{cor:v11v18-jacobi}
For invertible \(D\),
\[
 \det D\,
 \det (D^{-1})_{I,J}
 =
 \sigma(I,J)\det D_{J^c,I^c},
 \tag{18.5}
\]
where \(\sigma(I,J)\) is the standard complementary-minor sign.  Hence
every apparent inverse pole in the fully saturated coefficient cancels
against the vanishing determinant.
\end{corollary}

\begin{proof}
Apply the compound-matrix identity
\[
 \operatorname{adj}_r(D)
 =
 (\det D)\,\bigwedge^rD^{-1}
 \tag{18.6}
\]
in the basis wedge labelled by \(I,J\).  Its matrix coefficient is
Jacobi's complementary-minor identity (18.5).  Alternatively, compare
the invertible Wick expansion with the polynomial identity (18.4).
\end{proof}

The polynomial right side, rather than the product
\(\det D\,D^{-1}\), is the correct expression at a zero.

\subsection{Exact zero-mode saturation}

Suppose
\[
 \dim\ker D=\dim\ker D^*=r.
 \tag{18.7}
\]
Choose orthogonal decompositions
\[
 V=K\oplus V_\perp,\qquad
 W=C\oplus W_\perp,
 \tag{18.8}
\]
where \(K=\ker D\), \(C=\ker D^*\), and
\[
 D_\perp:V_\perp\longrightarrow W_\perp
 \tag{18.9}
\]
is invertible.

\begin{theorem}[Basis-free saturated value]
\label{thm:v11v18-saturated-value}
Insert oriented top wedges of \(K\) and \(C\).  The resulting Berezin
coefficient is, up to the single global orientation convention,
\[
 \det D_\perp.
 \tag{18.10}
\]
It is nonzero exactly when every domain and cokernel zero mode has been
saturated.  If any zero-mode Grassmann variable is omitted, the
corresponding coefficient vanishes.
\end{theorem}

\begin{proof}
Choose bases adapted to (18.8).  The matrix of \(D\) is
\[
 \begin{pmatrix}
 0&0\\
 0&D_\perp
 \end{pmatrix}.
\]
Taking \(I\) and \(J\) to be the zero-mode basis indices makes the
complementary minor in (18.4) equal to \(\det D_\perp\).  If a zero-mode
index is not inserted, it remains in the complementary matrix as a zero
row or column, so the determinant vanishes.  A change of adapted bases
changes the top wedges and \(\det D_\perp\) by reciprocal determinant
factors, leaving the tensorial coefficient invariant.
\end{proof}

\begin{corollary}[Extension through a crossing]
\label{cor:v11v18-crossing}
Suppose in continuous local frames
\[
 D(t)=
 \operatorname{diag}
 \bigl(\lambda_1(t),\ldots,\lambda_r(t),D_\perp(t)\bigr),
 \tag{18.11}
\]
where \(\lambda_j(t_0)=0\) and \(D_\perp(t_0)\) is invertible.  The
unsaturated determinant contains
\(\prod_j\lambda_j(t)\), while the coefficient saturating all \(r\)
crossing pairs is
\[
 \pm\det D_\perp(t)
 \tag{18.12}
\]
and extends continuously through \(t_0\).
\end{corollary}

\begin{proof}
Apply the complementary-minor identity in the displayed frame.  The
selected complementary matrix is \(D_\perp(t)\).
\end{proof}

\subsection{Application to the flat-holonomy node}

At a nodal point of V17, one full overlap spin block is zero on its four
spin components.  The full Dirac Berezin integral therefore requires four
\(\psi\) and four \(\bar\psi\) zero-mode insertions for a nonzero
coefficient.  A Weyl block has its own chiral zero-mode multiplicity and
must be treated in its determinant line rather than by taking a scalar
square root of the full determinant.

\begin{proposition}[Finite toron cylinder regularity]
\label{prop:v11v18-toron-regularity}
Let \(\mathcal O\) be a fermionic cylinder monomial whose source degree
saturates every zero mode present at a given flat-holonomy node.  Its
unnormalized finite Berezin coefficient is a complementary minor of the
full finite chiral matrix and therefore extends continuously through the
node.  If \(\mathcal O\) does not saturate the zero modes, the coefficient
vanishes at the node.
\end{proposition}

\begin{proof}
Choose finite bases with the nodal modes first and apply the saturated
value theorem.  Continuity follows because every complementary minor is
a finite polynomial in the matrix entries.
\end{proof}

This statement concerns the unnormalized coefficient.  A normalized
expectation may still be singular if one divides by a partition function
that vanishes in the same sector.

\subsection{Exterior-power formulation}

Let \(\bigwedge^kD:\bigwedge^kV\to\bigwedge^kW\) be the \(k\)-th
compound map.  The complete family of degree-\(r\) saturated coefficients
is the matrix of
\[
 \bigwedge^{n-r}D
 \tag{18.13}
\]
under the canonical complement pairing with
\(\bigwedge^rV\) and \(\bigwedge^rW^*\).

\begin{proposition}[Frame covariance of the whole coefficient]
\label{prop:v11v18-frame}
Under basis changes
\[
 D\longmapsto BDA^{-1},
 \tag{18.14}
\]
the full degree-\(r\) coefficient family transforms by the induced
exterior representations of \(A\) and \(B\).  Therefore the whole family
is a tensorial determinant-line object.  An individual scalar minor is
not frame invariant.
\end{proposition}

\begin{proof}
Exterior powers are functorial:
\[
 \bigwedge^{n-r}(BDA^{-1})
 =
 (\bigwedge^{n-r}B)
 (\bigwedge^{n-r}D)
 (\bigwedge^{n-r}A^{-1}).
\]
The complement pairings transform by the corresponding top determinant
factors.  This is exactly the stated covariance.
\end{proof}

\begin{firewall}[Local minors do not trivialize the line]
\label{fw:v11v18-line}
Polynomial extension of the coefficient through a zero does not produce
a global scalarization.  On chart overlaps the coefficient transforms by
the exterior transition maps in (18.14).  A global Standard Model
fermion measure still requires the anomaly-free determinant-line cocycle
and holonomy construction recorded in acquisition A4.
\end{firewall}

\subsection{A quantitative exterior-power bound}

The finite identity is stable under perturbation, although its naive
global bound grows with the matrix size.

\begin{lemma}[Exterior-power Lipschitz estimate]
\label{lem:v11v18-exterior-lipschitz}
For bounded maps \(A,B:V\to W\) and \(1\leq k\leq n\),
\[
 \|\bigwedge^kA-\bigwedge^kB\|
 \leq
 k\max\{\|A\|,\|B\|\}^{k-1}\|A-B\|.
 \tag{18.15}
\]
\end{lemma}

\begin{proof}
On the antisymmetric subspace,
\(\bigwedge^kA\) is the restriction of \(A^{\otimes k}\).
The tensor telescoping identity is
\[
 A^{\otimes k}-B^{\otimes k}
 =
 \sum_{j=1}^k
 A^{\otimes(j-1)}\otimes(A-B)\otimes
 B^{\otimes(k-j)}.
\]
Take norms and restrict to the antisymmetric subspace.
\end{proof}

For a saturated coefficient \(k=n-r\), the factor in (18.15) is
extensive.  It does not give the uniform projective estimate needed in
the continuum limit.

\begin{proposition}[Required normalized comparison]
\label{prop:v11v18-normalized-comparison}
Let \(S_i(\mathcal O)\) be a saturated cylinder coefficient transported
to a common determinant-line frame and divided by the declared local
reference normalization.  If
\[
 \sum_i
 \|S_{i+1}(\mathcal O)-T_iS_i(\mathcal O)\|<\infty
 \tag{18.16}
\]
for every fixed cylinder monomial, then the transported coefficients have
a frame-coherent projective limit, including at zero-mode crossings.
\end{proposition}

\begin{proof}
For \(m>n\), telescope the transported difference between
\(S_m(\mathcal O)\) and \(S_n(\mathcal O)\).  Its norm is bounded by the
tail of the convergent series in (18.16), so the family is Cauchy.
The polynomial minor formula shows that no inverse singularity is
introduced at a crossing.  Compatibility of the transports makes the
limit a coefficient in the projective determinant line.
\end{proof}

\begin{firewall}[The summable estimate remains open]
\label{fw:v11v18-summable}
Equation (18.16) has not been proved for the interacting
Standard Model--Einstein regulator.  The raw bound (18.15) cannot supply
it because \(n-r\) grows with the cutoff volume.  One still needs the
local Schur/density compiler, counterterm subtraction, and a
regulator-compatible determinant-line frame.
\end{firewall}

\begin{assessment}[V18 advance]
\label{ass:v11v18}
The massless toron node no longer causes a formal inverse singularity.
Every zero-mode-saturated cylinder coefficient is exactly a complementary
minor and extends polynomially through the crossing as a determinant-line
tensor.  The remaining continuum problem is quantitative and global:
obtain summable normalized local coefficient differences and glue their
frames.  The next upstream step is to analyze the determinant-line
connection on the commuting holonomy torus and separate its curvature
from possible large-gauge holonomy.
\end{assessment}


\section{Quaternionic trivialization on the commuting flat locus}
\label{sec:v11v19}

V18 leaves the determinant-line transition problem open.  On the restricted
locus of commuting flat holonomies, the free Wilson kernel has an
additional quaternionic structure.  This note uses it to construct a
canonical determinant frame on that locus.

\subsection{A five-Clifford representation}

Let \(\sigma_1,\sigma_2,\sigma_3\) be the Pauli matrices and choose
Hermitian \(4\times4\) matrices
\[
 \Gamma_1=\sigma_1\otimes\sigma_1,\quad
 \Gamma_2=\sigma_1\otimes\sigma_2,\quad
 \Gamma_3=\sigma_1\otimes\sigma_3,
 \tag{19.1}
\]
\[
 \Gamma_4=\sigma_2\otimes I,\qquad
 \Gamma_5=\sigma_3\otimes I.
 \tag{19.2}
\]
They satisfy
\[
 \Gamma_A\Gamma_B+\Gamma_B\Gamma_A=2\delta_{AB}I.
 \tag{19.3}
\]
This is unitarily equivalent to
\[
 \Gamma_\mu=i\gamma_5\gamma_\mu,\qquad
 \Gamma_5=\gamma_5.
 \tag{19.4}
\]
Hence
\[
 H_m(k)=\sum_{\mu=1}^4s_\mu(k)\Gamma_\mu+b_m(k)\Gamma_5.
 \tag{19.5}
\]

Let \(K\) denote componentwise complex conjugation in the representation
(19.1)--(19.2), and define the antiunitary operator
\[
 J:=\Gamma_2\Gamma_4K.
 \tag{19.6}
\]

\begin{lemma}[Quaternionic symmetry]
\label{lem:v11v19-quaternionic}
One has
\[
 J^2=-I,\qquad
 J\Gamma_A=\Gamma_AJ
 \quad(1\leq A\leq5).
 \tag{19.7}
\]
Consequently
\[
 JH_m(k)=H_m(k)J
 \tag{19.8}
\]
for every real \(k\).
\end{lemma}

\begin{proof}
The matrices \(\Gamma_1,\Gamma_3,\Gamma_5\) are real, whereas
\(\Gamma_2,\Gamma_4\) are purely imaginary.  The product
\(B=\Gamma_2\Gamma_4\) commutes with the first three and anticommutes
with \(\Gamma_2,\Gamma_4\).  Thus
\[
 B\overline{\Gamma_A}B^{-1}=\Gamma_A
\]
for every \(A\), which is the commutator assertion for \(J=BK\).
The matrix \(B\) is real and
\[
 J^2=B\overline B=B^2
 =\Gamma_2\Gamma_4\Gamma_2\Gamma_4=-I.
\]
Equation (19.8) follows from (19.5), whose coefficients are real.
\end{proof}

\subsection{The negative bundle}

Let
\[
 \widehat H_m(k):=\frac{H_m(k)}{r_m(k)},
 \qquad
 P_m^-(k):=\frac{I-\widehat H_m(k)}2.
 \tag{19.9}
\]
The V12 gap makes this a smooth rank-two projector on all of
\(\mathbb T^4\).

\begin{proposition}[Quaternionic negative line]
\label{prop:v11v19-negative}
The range
\[
 E_m^-(k):=\operatorname{ran}P_m^-(k)
 \tag{19.10}
\]
is \(J\)-invariant and has complex dimension two, hence quaternionic
dimension one.
\end{proposition}

\begin{proof}
Equation (19.8) implies that \(J\) commutes with every spectral projector
of \(H_m(k)\), including \(P_m^-(k)\).  The Clifford matrix
\(\widehat H_m(k)\) squares to one and has trace zero, so its eigenvalues
\(+1\) and \(-1\) each have multiplicity two.
\end{proof}

On any patch choose a unit vector \(v(k)\in E_m^-(k)\).  Then
\(Jv(k)\) is orthogonal to \(v(k)\), because antiunitarity and \(J^2=-1\)
give
\[
 \langle v,Jv\rangle
 =
 -\langle v,Jv\rangle.
\]
Define
\[
 \Omega_m^-(k):=v(k)\wedge Jv(k)
 \in\det E_m^-(k).
 \tag{19.11}
\]

\begin{theorem}[Canonical determinant frame]
\label{thm:v11v19-canonical-frame}
The local elements (19.11) agree on patch overlaps and define a global
unit-norm nonvanishing section
\[
 \Omega_m^-\in\Gamma(\det E_m^-).
 \tag{19.12}
\]
In this frame the determinant Berry connection is zero.  Therefore
\[
 c_1(\det E_m^-)=0
 \tag{19.13}
\]
and the determinant holonomy around every loop in the momentum torus is
one.
\end{theorem}

\begin{proof}
A second unit vector in the same quaternionic line has the form
\[
 v'=av+bJv,\qquad |a|^2+|b|^2=1.
\]
Antiunitarity and \(J^2=-1\) give
\[
 Jv'=-\overline b\,v+\overline a\,Jv.
\]
The transition matrix from \((v,Jv)\) to \((v',Jv')\) is
\[
 \begin{pmatrix}
 a&-\overline b\\
 b&\overline a
 \end{pmatrix},
\]
whose determinant is one.  Hence
\(v'\wedge Jv'=v\wedge Jv\), proving global gluing and nonvanishing.

For a differentiable local unit vector,
\[
 \langle\Omega_m^-,d\Omega_m^-\rangle
 =
 \langle v,dv\rangle+\langle Jv,d(Jv)\rangle.
\]
Antiunitarity gives
\(\langle Jv,Jdv\rangle=\langle dv,v\rangle\), which equals
\(-\langle v,dv\rangle\) by differentiation of \(\langle v,v\rangle=1\).
Thus the connection one-form is zero in every canonical patch frame.
Its curvature and all determinant holonomies vanish.  A global
nonvanishing section also gives (19.13).
\end{proof}

The rank-two bundle \(E_m^-\) may still carry nontrivial second-Chern
data.  The theorem trivializes its determinant line, not the whole
rank-two bundle.

\begin{corollary}[Complementary positive bundle]
\label{cor:v11v19-positive-bundle}
The positive projector
\[
 P_m^+(k):=\frac{I+\widehat H_m(k)}2
\]
also has a canonical global parallel determinant frame, obtained by the
same quaternionic construction.  Thus the conclusion applies whichever
of \(P_m^-\) or \(P_m^+\) is the overlap chiral-domain projector under
the chosen sign convention.
\end{corollary}

\begin{proof}
The operator \(J\) commutes with \(P_m^+\), whose range also has complex
dimension two.  Repeat the proof with a local unit vector in that range.
\end{proof}

\subsection{Commuting holonomies and large-gauge relabelling}

For a representation weight \(w\) and finite momentum \(n\), the flat
background block is \(E_m^-(q_{n,w}(a))\).  Give its determinant the
canonical frame
\[
 \Omega_{n,w}(a):=
 \Omega_m^-(q_{n,w}(a)).
 \tag{19.14}
\]
Their tensor product is
\[
 \Omega_{L,R}(a):=
 \bigotimes_{w\in\mathcal W_R}\bigotimes_n\Omega_{n,w}(a).
 \tag{19.15}
\]

\begin{theorem}[Flat-holonomy regulator-line trivialization]
\label{thm:v11v19-flat-line}
Equation (19.15) is a global parallel nonvanishing section of the
determinant of the negative Wilson-kernel bundle over the commuting
flat-holonomy torus.  It is invariant under every large-gauge momentum
relabeling.
\end{theorem}

\begin{proof}
Every factor is global and parallel by the canonical-frame theorem, so
their finite tensor product is as well.  A large gauge transformation
permutes the momentum labels as in V16.  Each factor is a two-form.
Exchanging two degree-two exterior factors contributes
\((-1)^{2\cdot2}=1\), so the tensor product is invariant under any such
permutation.
\end{proof}

Thus the gapped regulator subspace has no determinant phase obstruction on
this restricted locus.

\subsection{The fixed target chirality}

The fixed projector
\[
 P_+:=\frac{I+\Gamma_5}{2}
 \tag{19.16}
\]
also commutes with \(J\), so its rank-two range has the same canonical
determinant frame.  In the standard overlap convention, the chiral block
maps the \(J\)-invariant overlap subspace to this fixed target and
commutes with \(J\).

\begin{proposition}[Nonnegative Weyl scalar on the flat locus]
\label{prop:v11v19-weyl-scalar}
In the canonical quaternionic determinant frames, the free overlap Weyl
block has scalar determinant
\[
 \delta_m(k)=2\left(1+\frac{b_m(k)}{r_m(k)}\right)\geq0,
 \tag{19.17}
\]
and
\[
 \Delta_m(k)=\delta_m(k)^2.
 \tag{19.18}
\]
It vanishes to second order at the physical corner and has no phase
winding around that zero.
\end{proposition}

\begin{proof}
A complex-linear map commuting with \(J\) between quaternionic
one-dimensional spaces is multiplication by a quaternion in compatible
frames.  Its complex determinant is the squared quaternionic norm and is
nonnegative.  The full spin determinant computed in V17 is its square,
so the nonnegative square root is (19.17).  The V17 fourth-order bound
makes \(\delta_m\) second order.  A nonnegative continuous scalar has no
phase winding around an isolated zero.
\end{proof}

\begin{corollary}[Flat Standard Model product]
\label{cor:v11v19-sm-flat}
For any finite collection of Standard Model Weyl multiplets, the product
of the free commuting-flat block determinants can be represented in the
tensor product of the canonical frames (19.15).  Its scalar modulus is the
product of the nonnegative factors (19.17), with zero modes treated by
the complementary-minor saturation of V18.
\end{corollary}

\begin{proof}
Tensor the canonical frames over every species, weight, and momentum.
Use the Weyl scalar proposition away from zeros and the polynomial
extension of V18 at zeros.
\end{proof}

\subsection{Connection formula and scope}

For a general smooth projector \(P\), the determinant-line Berry curvature
is
\[
 \mathcal F_{\det P}
 =
 \Tr\bigl(P\,dP\wedge dP\bigr).
 \tag{19.19}
\]
For the quaternionic projector above this trace vanishes, consistent with
the parallel frame.

\begin{firewall}[Restricted-locus result]
\label{fw:v11v19-restricted}
The antiunitary \(J\) commutes with the momentum block after a commuting
flat background has been diagonalized into representation weights.  A
general complex Standard Model gauge field is not fixed by this
spin-only quaternionic symmetry; complex conjugation changes its gauge
matrices.  Therefore the canonical frame constructed here does not
extend automatically to curved, noncommuting configurations or between
topological sectors.  The global anomaly-free determinant-line
construction in acquisition A4 remains open.
\end{firewall}

\begin{remark}[Why local anomaly arithmetic is still needed]
\label{rem:v11v19-anomaly}
Flat connections have zero field strength and do not test the full local
gauge-anomaly curvature.  The Standard Model charge cancellations from
the archived programme constrain the curvature away from the flat locus.
They are logically independent of the quaternionic flat-locus frame.
\end{remark}

\begin{assessment}[V19 advance]
\label{ass:v11v19}
The determinant phase on the commuting flat locus is now closed at the
regulator level.  Each negative Wilson-kernel block is a quaternionic
line, its determinant has a canonical parallel frame, large-gauge
momentum permutations preserve the product frame, and the Weyl scalar is
nonnegative with a second-order physical zero.  The remaining phase
problem begins with curved noncommuting gauge fields, where this
quaternionic symmetry no longer supplies a frame.  V20 is the next
mandatory companion audit.
\end{assessment}


\section{V20 mandatory companion audit}
\label{sec:v11v20}

The fourth five-version boundary requires inspection of the newest active
companion notes.

\begin{audit}[Files inspected]
\label{audit:v11v20-files}
The exact records are
\[
\begin{array}{c|r|r}
\text{file}&\text{lines}&\text{bytes}\\ \hline
\texttt{Renewal\_Yang\_Mills\_Mass\_Gap\_Research\_Notes\_V29.tex}
&12200&414183\\
\texttt{Renewal\_NS\_Stability\_Research\_Notes\_V26.tex}
&5319&278283.
\end{array}
\]
Their SHA--256 values are
\begin{center}
\small\ttfamily
BED1ED939F606480DE9B341DD5BB6363275256150FC0A1753E355DA01E05B46F
\end{center}
and
\begin{center}
\small\ttfamily
82C887F8C2C69E4F28FAD7C3DC8DE5B9099CB5A4BF431EBA1FFF3E91935978AD.
\end{center}
The NS file is unchanged.  The YM stream has advanced from V26 to V29.
\end{audit}

\subsection{New Yang--Mills material}

YM V27 executes exact sparse cubic and quartic toron vertex arrays for one
polarization.  Independent colour projections agree, the cubic spatial
matrix has the required skew type, and the quartic matrix has the required
symmetric type.

YM V28 generates the Wilson arrays in the identical lift and evaluates a
nonzero endpoint-minus-Wilson toron determinant response exactly.  It
keeps the value as one finite zero-cell contribution with weight
\(M^{-4}\) and explicitly proves that it is neither a gauge-boson mass nor
a continuum coefficient.

YM V29 differentiates both harmonic routers twice in external momentum.
Every differentiated router residual vanishes in exact arithmetic.
Reflection makes the first router jet purely imaginary and the even jets
real.  The endpoint and Wilson routers agree through first order but
differ at second order.  The note warns that router jets alone are not the
complete response jets.

\begin{audit}[Accepted transfers]
\label{audit:v11v20-accepted}
The following conclusions transfer structurally.
\begin{enumerate}[label=\textnormal{(\roman*)}]
\item A finite toron determinant contribution must be retained with its
exact finite-volume weight before any continuum estimate.  This agrees
with V16--V18.
\item Reflection parity can cancel a first derivative only after router,
owner, inverse, measure, and reflected-conjugate rows have been assembled.
No individual imaginary row may be deleted.
\item A coordinate lift can agree through first order and first differ at
second order.  The Einstein constraint and gauge reduction therefore
require the exact second response already retained in the archived SM
implicit-function card.
\item Sparse exact contraction and symmetry-type checks are useful
verification methods, but the YM arrays and rational toron value do not
transfer numerically to the fermion determinant.
\end{enumerate}
\end{audit}

\begin{audit}[No new density theorem]
\label{audit:v11v20-no-density}
Neither YM V27--V29 nor NS V26 supplies the coupled plaquette insertion
tail, defect-cluster estimate, determinant-line gluing, or reduced
Einstein free-energy lower bound required by the SM programme.
\end{audit}

\subsection{Direction after the audit}

V19 constructs a canonical determinant frame on the commuting flat locus,
but its quaternionic symmetry need not survive a curved complex gauge
background.  The next finite step is to extend that frame through a
spectrally gapped tube without assuming the symmetry persists.

\begin{decision}[V21 target]
\label{dec:v11v20-direction}
V21 will use the direct rotation between nearby spectral projectors to:
\begin{enumerate}[label=\textnormal{(\alph*)}]
\item construct an explicit unitary intertwiner whenever
\(\|P-P_0\|<1\);
\item transport the canonical flat determinant frame into a curved
gapped tube;
\item prove gauge covariance of the construction under an equivariant
choice of flat reference;
\item expose the resulting Berry connection and curvature as the remaining
local Ward phase, rather than declaring the transported frame parallel.
\end{enumerate}
\end{decision}

\begin{firewall}[Audit boundary]
\label{fw:v11v20-boundary}
The proposed tube frame will be a local or tubular trivialization.  It
will not glue arbitrary topological sectors, prove global holonomy
triviality, or supply the missing probability of the tube.
\end{firewall}

\begin{assessment}[V20 audit result]
\label{ass:v11v20}
The audit is complete and supports a curved-tube extension of the V19
frame.  A substantive V21 follows, as required by the five-version rule.
\end{assessment}


\section{Kato direct rotation from the flat determinant frame}
\label{sec:v11v21}

V19 gives a canonical determinant frame on the commuting flat locus.
A curved background generally breaks its quaternionic symmetry.  If its
spectral projector remains within operator distance one of a chosen flat
reference projector, the direct rotation gives a canonical extension.

\subsection{Direct rotation of two projections}

Let \(P,Q\) be orthogonal projections on a finite-dimensional Hilbert
space and assume
\[
 \delta:=\|P-Q\|<1.
 \tag{21.1}
\]
Define
\[
 S(P,Q):=PQ+(I-P)(I-Q)
 \tag{21.2}
\]
and
\[
 \mathcal U(P,Q)
 :=
 S(P,Q)\bigl(S(P,Q)^*S(P,Q)\bigr)^{-1/2}.
 \tag{21.3}
\]

\begin{lemma}[Polar gap]
\label{lem:v11v21-polar-gap}
One has
\[
 S(P,Q)^*S(P,Q)=I-(P-Q)^2
 \tag{21.4}
\]
and therefore
\[
 S(P,Q)^*S(P,Q)\succeq(1-\delta^2)I.
 \tag{21.5}
\]
\end{lemma}

\begin{proof}
Since \(P(I-P)=Q(I-Q)=0\), the cross terms in \(S^*S\) vanish:
\[
 S^*S=QPQ+(I-Q)(I-P)(I-Q).
\]
Expanding the second term and combining it with the first gives
\[
 I-P-Q+PQ+QP=I-(P-Q)^2.
\]
The norm bound in (21.1) gives (21.5).
\end{proof}

\begin{theorem}[Kato direct rotation]
\label{thm:v11v21-direct-rotation}
The operator \(\mathcal U(P,Q)\) is unitary and satisfies
\[
 P\mathcal U(P,Q)=\mathcal U(P,Q)Q.
 \tag{21.6}
\]
Consequently it maps \(\operatorname{ran}Q\) unitarily onto
\(\operatorname{ran}P\).  It depends smoothly, or analytically, on
\(P,Q\) as long as the strict margin (21.1) persists.
\end{theorem}

\begin{proof}
The polar gap makes \(S\) invertible, so its polar factor (21.3) is
unitary.  Directly from (21.2),
\[
 PS=PQ=SQ.
 \tag{21.7}
\]
Equation (21.4) commutes with \(Q\), hence so does its inverse square
root.  Multiplying (21.7) by that inverse square root gives (21.6).
Smooth or analytic dependence follows from the convergent functional
calculus for \(x^{-1/2}\) on the interval
\([1-\delta^2,1]\).
\end{proof}

\subsection{Quantitative differential bound}

Let \(P,Q\) vary differentiably and retain
\(\|P-Q\|\leq\delta_*<1\).  Put
\[
 a_*:=1-\delta_*^2.
 \tag{21.8}
\]

\begin{proposition}[Derivative stock for the direct rotation]
\label{prop:v11v21-derivative}
For every parameter derivative,
\[
 \|d\mathcal U(P,Q)\|
 \leq
 \left(a_*^{-1/2}+a_*^{-3/2}\right)
 \bigl(\|dP\|+\|dQ\|\bigr).
 \tag{21.9}
\]
\end{proposition}

\begin{proof}
Differentiating (21.2) gives
\[
 dS=dP(2Q-I)+(2P-I)dQ.
\]
The two reflections have norm one, so
\[
 \|dS\|\leq\|dP\|+\|dQ\|.
 \tag{21.10}
\]
Let \(A=S^*S\).  Since \(A\preceq I\), \(\|S\|\leq1\), and
\[
 \|dA\|\leq2\|dS\|.
 \tag{21.11}
\]
The integral functional calculus
\[
 A^{-1/2}
 =
 \frac1\pi\int_0^\infty t^{-1/2}(A+t)^{-1}\,dt
\]
gives
\[
 \|d(A^{-1/2})\|
 \leq\frac12a_*^{-3/2}\|dA\|
 \leq a_*^{-3/2}\|dS\|.
 \tag{21.12}
\]
Differentiate \(\mathcal U=SA^{-1/2}\), use
\(\|A^{-1/2}\|\leq a_*^{-1/2}\), and insert
(21.10)--(21.12).
\end{proof}

The bound is finite but contains no volume cancellation.  Traces of
\(d\mathcal U\) require a separate local or anomaly-cancelled estimate.

\subsection{Projector closeness from a kernel comparison}

Let \(H,H_0\) be Hermitian kernels whose relevant spectral parts are
enclosed by the same contour \(\Gamma\).  Suppose
\[
 \sup_{z\in\Gamma}
 \left\{
 \|(z-H)^{-1}\|,\|(z-H_0)^{-1}\|
 \right\}\leq R.
 \tag{21.13}
\]
Let \(P,P_0\) be their Riesz projectors.

\begin{proposition}[Riesz distance]
\label{prop:v11v21-riesz-distance}
One has
\[
 \|P-P_0\|
 \leq
 C_\Gamma R^2\|H-H_0\|,
 \qquad
 C_\Gamma:=\frac{\operatorname{length}(\Gamma)}{2\pi}.
 \tag{21.14}
\]
Thus the direct rotation is certified whenever the right side is
strictly less than one.
\end{proposition}

\begin{proof}
The resolvent identity gives
\[
 (z-H)^{-1}-(z-H_0)^{-1}
 =
 (z-H)^{-1}(H-H_0)(z-H_0)^{-1}.
\]
Integrate around \(\Gamma\) and take norms.
\end{proof}

For the V12 link-field tube, \(\|H-H_0\|\leq8d_\infty(U,U_0)\).
A chosen contour and its resolvent margin therefore give an explicit
smaller tube on which (21.1) holds.

\subsection{Tubular determinant frame}

Let \(\mathcal F\) be the commuting flat locus and let
\[
 r:\mathcal T\longrightarrow\mathcal F
 \tag{21.15}
\]
be a continuous retraction from a gapped curved tube.  For
\(U\in\mathcal T\), set
\[
 P(U):=\text{the overlap chiral-domain projector},
 \qquad
 Q(U):=P(r(U)).
 \tag{21.16}
\]
Assume
\[
 \sup_{U\in\mathcal T}\|P(U)-Q(U)\|\leq\delta_*<1.
 \tag{21.17}
\]
Let \(\Omega_{\mathcal F}(r(U))\) be the canonical flat determinant frame
from V19, using the sign convention appropriate to \(P\).  Define
\[
 \Omega_{\mathcal T}(U)
 :=
 \bigwedge^{\operatorname{rank}P}
 \mathcal U(P(U),Q(U))\,
 \Omega_{\mathcal F}(r(U)).
 \tag{21.18}
\]

\begin{theorem}[Explicit tubular trivialization]
\label{thm:v11v21-tubular}
Equation (21.18) is a continuous unit-norm nonvanishing section of
\(\det\operatorname{ran}P\) over \(\mathcal T\).  It is smooth or analytic
where \(P\) and \(r\) have that regularity.
\end{theorem}

\begin{proof}
The direct rotation maps \(\operatorname{ran}Q(U)\) unitarily onto
\(\operatorname{ran}P(U)\).  Its top exterior power therefore maps the
nonzero unit vector \(\Omega_{\mathcal F}\) to a nonzero unit vector.
Every operation in (21.18) is continuous, and the direct-rotation theorem
gives the stronger regularity.
\end{proof}

The construction remains valid at a zero of the physical overlap
operator, because the Wilson-kernel projector \(P(U)\) is still gapped.
The determinant section may vanish there, while its ambient line frame
does not; V18 supplies the saturated coefficient.

\subsection{Gauge covariance}

Suppose a lattice gauge transformation \(g\) acts by a unitary
\(G(g)\), and the reference assignment is equivariant:
\[
 r(U^g)=r(U)^g.
 \tag{21.19}
\]
Assume the canonical flat frame is equivariant in the tensorial sense.

\begin{proposition}[Equivariance of the tube frame]
\label{prop:v11v21-equivariance}
One has
\[
 \mathcal U(GPG^{-1},GQG^{-1})
 =
 G\mathcal U(P,Q)G^{-1}.
 \tag{21.20}
\]
Consequently \(\Omega_{\mathcal T}(U^g)\) is the gauge transform of
\(\Omega_{\mathcal T}(U)\).
\end{proposition}

\begin{proof}
The operator \(S\) in (21.2), its positive square
\(S^*S\), and the inverse square root all commute with unitary
conjugation.  This gives (21.20).  Apply the top exterior power and the
equivariance of the flat frame.
\end{proof}

An equivariant retraction is a nontrivial geometric input.  Tree gauge
provides local sections, but Gribov and sector issues prevent it from
being assumed globally.

\subsection{The induced Berry connection}

Let \(V_0(U)\) be a local orthonormal frame for \(Q(U)\) whose top wedge is
\(\Omega_{\mathcal F}(r(U))\), and set
\[
 V(U):=\mathcal U(P(U),Q(U))V_0(U).
 \tag{21.21}
\]
The determinant connection one-form in the tubular frame is
\[
 \mathcal A_{\mathcal T}
 =
 \Tr\bigl(V^*dV\bigr)
 =
 \Tr\left(
 Q\,\mathcal U^*d\mathcal U\,Q
 \right)
 +
 \mathcal A_{\mathcal F}\circ dr.
 \tag{21.22}
\]
V19 gives \(\mathcal A_{\mathcal F}=0\) in its canonical frame, so
\[
 \mathcal A_{\mathcal T}
 =
 \Tr\left(
 Q\,\mathcal U^*d\mathcal U\,Q
 \right).
 \tag{21.23}
\]

\begin{proposition}[Curvature remains visible]
\label{prop:v11v21-curvature}
The curvature of (21.23) is the intrinsic determinant-line curvature
\[
 d\mathcal A_{\mathcal T}
 =
 \Tr(P\,dP\wedge dP).
 \tag{21.24}
\]
It need not vanish when the curved background breaks the flat
quaternionic symmetry.
\end{proposition}

\begin{proof}
For any orthonormal frame \(V\) with \(P=VV^*\), the standard
differentiation of \(\Tr(V^*dV)\), using
\(d(V^*V)=0\), gives
\[
 d\Tr(V^*dV)=\Tr(P\,dP\wedge dP).
\]
Equation (21.21) is such a frame.  The result is frame independent.
\end{proof}

\begin{firewall}[A frame is not a flat connection]
\label{fw:v11v21-not-flat}
The tubular section proves topological triviality of the determinant line
on the declared tube.  It does not make the Berry connection flat.
Equation (21.23) is an explicit phase row, and (21.24) must be cancelled
in the complete Standard Model tensor product or by a measured local
current.  Omitting it would recreate the false anomaly problem of V2.
\end{firewall}

\subsection{Species product and overlap transitions}

For Standard Model species \(s\), let \(P_s,Q_s,\mathcal U_s\) denote
their projectors and direct rotations.  The tensor frame is
\[
 \Omega_{\rm SM,\mathcal T}
 :=
 \bigotimes_s
 \bigwedge^{\operatorname{rank}P_s}\mathcal U_s\,
 \Omega_{s,\mathcal F}.
 \tag{21.25}
\]
Its connection and curvature are the sums
\[
 \mathcal A_{\rm SM}=\sum_s\mathcal A_s,
 \qquad
 \mathcal F_{\rm SM}=\sum_s\Tr(P_s\,dP_s\wedge dP_s).
 \tag{21.26}
\]
The archived Standard Model arithmetic cancels the continuum local anomaly
polynomial.  It does not prove that the finite expression in (21.26)
vanishes or is the derivative of a regulator-compatible local current.

If two retractions or reference charts give flat reference frames
\(\Omega_\alpha^{\mathcal F},\Omega_\beta^{\mathcal F}\) and tubular
frames \(\Omega_\alpha,\Omega_\beta\), define the overlap phase by
\[
 \Omega_\alpha=g_{\alpha\beta}\Omega_\beta.
 \tag{21.27}
\]
Equivalently, \(g_{\alpha\beta}\) is the determinant of the map
\[
 \mathcal U_\beta^{-1}\mathcal U_\alpha:
 \operatorname{ran}Q_\alpha\longrightarrow\operatorname{ran}Q_\beta
\]
evaluated between the determinant frames
\(\Omega_\alpha^{\mathcal F}\) and
\(\Omega_\beta^{\mathcal F}\).

\begin{proposition}[Explicit tube cocycle]
\label{prop:v11v21-cocycle}
The phases (21.27) obey
\[
 g_{\alpha\beta}g_{\beta\gamma}g_{\gamma\alpha}=1
 \tag{21.28}
\]
on triple overlaps.  A global scalar frame exists on the union of the
tubes exactly when this cocycle is a coboundary.
\end{proposition}

\begin{proof}
Each \(g_{\alpha\beta}\) is the ratio of two nonvanishing local
determinant frames.  Ratios telescope on a triple overlap, proving
(21.28).  Multiplying local frames by phases \(h_\alpha\) makes them
agree precisely when
\(g_{\alpha\beta}=h_\alpha^{-1}h_\beta\), which is the coboundary
condition.
\end{proof}

\begin{firewall}[Remaining acquisitions]
\label{fw:v11v21-remaining}
No equivariant retraction covering the relevant curved physical measure,
finite Standard Model current cancelling (21.26), overlap-cocycle
coboundary, or scale-compatible version of the tube frame has been
constructed.  Probability concentration in the tubes also remains open.
\end{firewall}

\begin{assessment}[V21 advance]
\label{ass:v11v21}
The required post-audit continuation is complete.  The canonical flat
determinant frame now extends explicitly to every spectrally close curved
projector through the Kato direct rotation.  The construction is
gauge-equivariant when the reference retraction is equivariant, remains
regular at physical determinant zeros, and exposes the exact Berry
connection, curvature, and tube-overlap cocycle.  The next finite
bottleneck is the Standard Model sum of the tubular curvature and its
possible local-current primitive.
\end{assessment}


\section{Relative homotopy current for the tubular determinant curvature}
\label{sec:v11v22}

V21 constructs a tubular determinant frame and leaves its curvature as an
explicit phase row.  On a tube that deformation retracts equivariantly to
the commuting flat locus, the relative de Rham homotopy operator gives a
candidate measure current.  This note states the exact hypotheses under
which that current is local and realizes a phase choice.

\subsection{Relative homotopy operator}

Let \(\mathcal T\) be a finite-dimensional smooth regulator tube and
\(\mathcal F\subset\mathcal T\) its commuting flat locus.  Assume there is
a smooth deformation retraction
\[
 h:[0,1]\times\mathcal T\longrightarrow\mathcal T
 \tag{22.1}
\]
such that
\[
 h_1=\operatorname{id}_{\mathcal T},\qquad
 h_0=\iota\circ r,\qquad
 h_t|_{\mathcal F}=\operatorname{id}_{\mathcal F}.
 \tag{22.2}
\]
For a differential \(k\)-form \(\omega\), \(k\geq1\), define
\[
 (K_h\omega)_U(v_1,\ldots,v_{k-1})
 :=
 \int_0^1
 \omega_{h_t(U)}
 \left(
 \partial_th_t(U),
 D_Uh_t\,v_1,\ldots,D_Uh_t\,v_{k-1}
 \right)\,dt.
 \tag{22.3}
\]

\begin{theorem}[Relative homotopy formula]
\label{thm:v11v22-homotopy}
For every smooth differential form \(\omega\),
\[
 dK_h\omega+K_hd\omega
 =
 h_1^*\omega-h_0^*\omega
 =
 \omega-r^*\iota^*\omega.
 \tag{22.4}
\]
\end{theorem}

\begin{proof}
Let \(X=\partial_t\) on \([0,1]\times\mathcal T\).  Cartan's identity
\[
 \mathcal L_X=d\iota_X+\iota_Xd
\]
applied to \(h^*\omega\) gives
\[
 \frac{d}{dt}h_t^*\omega
 =
 d\bigl(h_t^*\iota_{\partial_t h}\omega\bigr)
 +
 h_t^*\iota_{\partial_t h}(d\omega).
\]
Integrate from zero to one.  The two integrals are precisely
\(dK_h\omega\) and \(K_hd\omega\), yielding (22.4).
\end{proof}

\subsection{Application to the Standard Model determinant curvature}

At a fixed regulator let
\[
 \mathcal F_{\rm SM}
 =
 \sum_s\Tr(P_s\,dP_s\wedge dP_s)
 \tag{22.5}
\]
be the tensor-product determinant curvature from V21.  The Bianchi
identity gives
\[
 d\mathcal F_{\rm SM}=0.
 \tag{22.6}
\]
V19 supplies a parallel determinant frame on \(\mathcal F\), so
\[
 \iota^*\mathcal F_{\rm SM}=0.
 \tag{22.7}
\]
Define the relative current
\[
 \mathcal J_h:=K_h\mathcal F_{\rm SM}.
 \tag{22.8}
\]

\begin{theorem}[Tubular curvature primitive]
\label{thm:v11v22-primitive}
Under (22.1)--(22.7),
\[
 d\mathcal J_h=\mathcal F_{\rm SM},
 \qquad
 \iota^*\mathcal J_h=0.
 \tag{22.9}
\]
\end{theorem}

\begin{proof}
Apply the homotopy formula to the closed two-form
\(\mathcal F_{\rm SM}\).  Equations (22.6)--(22.7) give
\[
 dK_h\mathcal F_{\rm SM}
 =
 \mathcal F_{\rm SM}-r^*0.
\]
If \(U\in\mathcal F\), then \(h_t(U)=U\) and
\(\partial_th_t(U)=0\), so (22.3) gives
\(\iota^*\mathcal J_h=0\).
\end{proof}

This conclusion is exact at finite cutoff.  It uses the full species sum
in (22.5) before any spectral or local decomposition.

\subsection{Equivariance}

Let the finite lattice gauge group act on \(\mathcal T\).  Assume
\[
 h_t(U^g)=h_t(U)^g
 \tag{22.10}
\]
and that \(\mathcal F_{\rm SM}\) is gauge invariant as a two-form.

\begin{proposition}[Equivariant relative current]
\label{prop:v11v22-equivariant}
Under (22.10),
\[
 g^*\mathcal J_h=\mathcal J_h.
 \tag{22.11}
\]
For an infinitesimal gauge vector \(\xi^\#\), the scalar
\[
 \mathfrak j_h(\xi;U):=
 \mathcal J_h{}_U(\xi^\#_U)
 \tag{22.12}
\]
is therefore an equivariant candidate measure term.
\end{proposition}

\begin{proof}
Equivariance of \(h_t\) carries \(\partial_th_t\) and \(D h_t\) to their
gauge transforms.  Gauge invariance of the curvature leaves the integrand
in (22.3) unchanged.  Integration gives (22.11).  Evaluation on the
equivariant gauge vector field gives (22.12).
\end{proof}

For an anomalous representation, the existence of an equivariant local
homotopy satisfying all stated conditions can fail.  The proposition does
not make anomaly cancellation automatic.

\subsection{Locality stock}

Let \(\|\cdot\|_\mu\) denote one of the exponentially weighted local
multilinear norms used in V3--V9.  Assume uniformly on the homotopy:
\[
 \left|
 \mathcal F_{\rm SM}|_{h_t(U)}(u,v)
 \right|_\mu
 \leq C_F\|u\|_\mu\|v\|_\mu,
 \tag{22.13}
\]
\[
 \|\partial_th_t(U)\|_\mu\leq B_h,
 \qquad
 \|D_Uh_t\,v\|_\mu\leq L_h\|v\|_\mu.
 \tag{22.14}
\]

\begin{proposition}[Homotopy-current locality bound]
\label{prop:v11v22-locality}
Under (22.13)--(22.14),
\[
 |\mathcal J_h{}_U(v)|_\mu
 \leq C_FB_hL_h\|v\|_\mu.
 \tag{22.15}
\]
If the three kernels in (22.13)--(22.14) have exponential rate \(\mu\),
then their compositions have every smaller fixed rate
\(0<\mu'<\mu\), with the corresponding convolution constant.
\end{proposition}

\begin{proof}
Insert (22.13)--(22.14) into (22.3) for \(k=2\):
\[
 |\mathcal J_h{}_U(v)|_\mu
 \leq
 \int_0^1 C_F B_h L_h\|v\|_\mu\,dt,
\]
which is (22.15).  For kernel locality, each contraction is a finite
convolution of exponential kernels.  Reserving
\(e^{-(\mu-\mu')d}\) makes the remaining geometric volume sum finite,
as in the V8 intertwiner proof.
\end{proof}

\begin{firewall}[Uniform locality is not proved]
\label{fw:v11v22-locality}
The curvature of a gapped finite overlap projector is smooth and
finite-range-derived, but (22.13) with a cutoff-independent local stock
for the complete interacting Standard Model sum has not been proved.
Likewise, the Kato tube in V21 does not itself provide an equivariant
local deformation retraction with the uniform bounds (22.14).
\end{firewall}

\subsection{Realization as a determinant phase choice}

Let \(\mathcal A_{\rm K}\) be the Berry connection one-form in the V21
Kato frame.  Then
\[
 d\mathcal A_{\rm K}=\mathcal F_{\rm SM}.
 \tag{22.16}
\]
By (22.9),
\[
 \mathcal B_h:=\mathcal A_{\rm K}-\mathcal J_h
 \tag{22.17}
\]
is closed.  Both terms restrict to zero in the canonical V19 frame on
\(\mathcal F\).

\begin{theorem}[Trivial residual holonomy]
\label{thm:v11v22-residual-holonomy}
Assume the Kato frame and homotopy are global on \(\mathcal T\).
Then the closed connection \(\mathcal B_h\) has trivial holonomy around
every loop in \(\mathcal T\).  Hence there is a global \(U(1)\)-valued
phase \(e^{-\vartheta_h}\) such that the rephased Kato frame has connection
one-form \(\mathcal J_h\).
\end{theorem}

\begin{proof}
Flat-connection holonomy depends only on the homotopy class of a loop.
The deformation \(h_t\) carries every loop in \(\mathcal T\) to a loop in
\(\mathcal F\).  On \(\mathcal F\), both
\(\mathcal A_{\rm K}\) and \(\mathcal J_h\) vanish, so the holonomy of
\(\mathcal B_h\) is one.  Fix a base point and define
\[
 \vartheta_h(U):=\int_{\gamma_U}\mathcal B_h
 \quad\pmod{2\pi i}
\]
along any path from the base point to \(U\).  Trivial loop holonomy makes
this independent of the path.  Thus \(d\vartheta_h=\mathcal B_h\).
Rephasing changes the connection from
\(\mathcal A_{\rm K}\) to
\(\mathcal A_{\rm K}-d\vartheta_h=\mathcal J_h\).
\end{proof}

The theorem does not change the curvature of the determinant connection.
It constructs a phase whose connection equals the local current
\(\mathcal J_h\).

\subsection{Ward integrability}

For two infinitesimal gauge vectors \(\xi^\#,\eta^\#\), equation (22.9)
implies
\[
 \xi^\#\mathfrak j_h(\eta)
 -\eta^\#\mathfrak j_h(\xi)
 -\mathfrak j_h([\xi,\eta])
 =
 \mathcal F_{\rm SM}(\xi^\#,\eta^\#).
 \tag{22.18}
\]

\begin{proposition}[Finite Wess--Zumino consistency]
\label{prop:v11v22-consistency}
If the homotopy current is equivariant, then (22.18) is the exact
finite-cutoff integrability condition for the gauge-orbit measure term.
\end{proposition}

\begin{proof}
For any one-form \(\mathcal J\),
\[
 d\mathcal J(\xi^\#,\eta^\#)
 =
 \xi^\#\mathcal J(\eta^\#)
 -\eta^\#\mathcal J(\xi^\#)
 -\mathcal J([\xi^\#,\eta^\#]).
\]
The fundamental vector fields reproduce the Lie bracket with the fixed
action convention.  Insert \(d\mathcal J_h=\mathcal F_{\rm SM}\).
\end{proof}

\subsection{Dependence on the chosen homotopy}

Let \(h\) and \(\widetilde h\) be two relative homotopies satisfying the
same hypotheses.  Their currents obey
\[
 d(\mathcal J_h-\mathcal J_{\widetilde h})=0.
 \tag{22.19}
\]

\begin{proposition}[Homotopy ambiguity]
\label{prop:v11v22-ambiguity}
On a simply connected overlap, there is a scalar
\(c_{h\widetilde h}\) such that
\[
 \mathcal J_h-\mathcal J_{\widetilde h}
 =
 dc_{h\widetilde h}.
 \tag{22.20}
\]
On triple overlaps the sums of these scalars are locally constant
multiples of \(2\pi i\).  Their exponentials form the same type of
\(U(1)\) cocycle as the tube-frame transitions in V21.
\end{proposition}

\begin{proof}
Equation (22.19) and the Poincar\'e lemma give (22.20) on a simply
connected overlap.  Summing cyclically makes the differentials telescope
to zero, so the sum is locally constant.  Exponentiation converts the
additive constants to a multiplicative cocycle.
\end{proof}

\begin{firewall}[Global and scale compatibility]
\label{fw:v11v22-global}
The tubular primitive does not prove that the overlap constants in
(22.20) are a coboundary across all sectors.  It also does not compare
\(\mathcal J_h\) at adjacent regulators.  A continuum fermion measure
requires summable transported current differences, compatible
zero-mode saturation, and trivial residual cocycles in both space and
scale.
\end{firewall}

\begin{assessment}[V22 advance]
\label{ass:v11v22}
On any equivariantly retracting gapped tube, the complete Standard Model
determinant curvature now has an explicit relative primitive.  Uniform
curvature and homotopy locality imply a local current bound; the residual
between that current and the Kato Berry connection is flat with trivial
holonomy, so the current is realized by a global phase choice on the
tube.  The next finite task is a stability theorem comparing these
homotopy currents under adjacent regulator transport.
\end{assessment}
\section{Adjacent-regulator stability of the homotopy current}
\label{sec:v11v23}

V22 constructs a relative primitive at one cutoff.  Scale compatibility
requires more: after putting two adjacent regulators in one chart, the
curvatures, the retraction paths, their velocities, and their tangent
maps must all be compared.  This note proves the required deterministic
estimate and isolates a summable scale card.

\subsection{A two-homotopy comparison lemma}

Let \(\mathcal T\) be a finite-dimensional normed smooth manifold
contained in one coordinate chart, and use the chart to identify nearby
tangent spaces.  Let \(F,\widetilde F\) be two-form fields and let
\(h,\widetilde h:[0,1]\times\mathcal T\to\mathcal T\) be smooth
homotopies.  Write
\[
 V_t(x)=\partial_t h_t(x),\quad
 A_t(x)=D_xh_t,\quad
 \widetilde V_t(x)=\partial_t\widetilde h_t(x),\quad
 \widetilde A_t(x)=D_x\widetilde h_t.
 \tag{23.1}
\]
Their homotopy currents are
\[
 J=K_hF,\qquad \widetilde J=K_{\widetilde h}\widetilde F.
 \tag{23.2}
\]
For a two-form \(G\), let
\[
 \|G\|_0:=\sup_{x}\sup_{\|u\|,\|v\|\leq1}|G_x(u,v)|.
 \tag{23.3}
\]
Assume that, throughout the two homotopy images,
\[
 \|F-\widetilde F\|_0\leq\epsilon_F,\qquad
 \|\widetilde F\|_0\leq C_F,
 \tag{23.4}
\]
and that \(\widetilde F\) is Lipschitz in its base point:
\[
 |(\widetilde F_y-\widetilde F_z)(u,v)|
 \leq C_{\rm Lip}\|y-z\|\,\|u\|\,\|v\|.
 \tag{23.5}
\]
Suppose also
\[
 \max(\|V_t\|,\|\widetilde V_t\|)\leq B,\qquad
 \max(\|A_t\|,\|\widetilde A_t\|)\leq L,
 \tag{23.6}
\]
and
\[
 \|h_t-\widetilde h_t\|\leq\epsilon_0,\quad
 \|V_t-\widetilde V_t\|\leq\epsilon_v,\quad
 \|A_t-\widetilde A_t\|\leq\epsilon_1.
 \tag{23.7}
\]

\begin{theorem}[Adjacent-current estimate]
\label{thm:v11v23-current-estimate}
Under (23.1)--(23.7),
\[
 \|J-\widetilde J\|_0
 \leq
 BL\epsilon_F
 +C_{\rm Lip}BL\epsilon_0
 +C_FL\epsilon_v
 +C_FB\epsilon_1.
 \tag{23.8}
\]
The same statement holds in an exponentially weighted local norm whenever
that norm is submultiplicative for the displayed contractions.  If a
rate loss is needed for convolution, every fixed smaller rate may be
used.
\end{theorem}

\begin{proof}
For \(\|w\|\leq1\), add and subtract three intermediate integrands:
\begin{align*}
 &(F_{h_t})(V_t,A_tw)
 -(\widetilde F_{\widetilde h_t})(\widetilde V_t,\widetilde A_tw)\\
 &=
 ((F-\widetilde F)_{h_t})(V_t,A_tw)\\
 &\quad+
 ((\widetilde F_{h_t}-\widetilde F_{\widetilde h_t}))
       (V_t,A_tw)\\
 &\quad+
 (\widetilde F_{\widetilde h_t})
       (V_t-\widetilde V_t,A_tw)\\
 &\quad+
 (\widetilde F_{\widetilde h_t})
       (\widetilde V_t,(A_t-\widetilde A_t)w).
\end{align*}
Equations (23.4)--(23.7) bound these four lines by, respectively,
\[
 BL\epsilon_F,\quad
 C_{\rm Lip}BL\epsilon_0,\quad
 C_FL\epsilon_v,\quad
 C_FB\epsilon_1.
\]
Integrating over an interval of length one and taking the supremum over
\(x,w\) proves (23.8).  The local-norm version uses the same four-term
telescoping estimate and the stated submultiplicativity.
\end{proof}

The estimate does not use closedness.  When the two forms are closed and
have zero restriction to the same retract, it compares the actual
primitives selected by the two homotopies.

\subsection{Transport between adjacent regulators}

Let \(\mathcal T_i\) be the tube at regulator \(i\).  Assume an adjacent
identification
\[
 \Phi_i:\mathcal U_i\subset\mathcal T_i
       \longrightarrow\mathcal U_{i+1}\subset\mathcal T_{i+1}
 \tag{23.9}
\]
has been fixed on the cylinder data under consideration.  Pull all
level-\(i+1\) objects back by \(\Phi_i\).  Thus
\[
 \widehat F_{i+1}:=\Phi_i^*F_{i+1},\qquad
 \widehat h_{i+1,t}:=\Phi_i^{-1}\circ h_{i+1,t}\circ\Phi_i,
 \tag{23.10}
\]
where the second expression is used only where the maps have the stated
common domain.  Naturality of the homotopy operator gives
\[
 \Phi_i^*(K_{h_{i+1}}F_{i+1})
 =K_{\widehat h_{i+1}}\widehat F_{i+1}.
 \tag{23.11}
\]

\begin{definition}[Adjacent current card]
\label{def:v11v23-card}
At scale \(i\), let
\[
 (\epsilon_{F,i},\epsilon_{0,i},\epsilon_{v,i},\epsilon_{1,i})
 \tag{23.12}
\]
be constants satisfying (23.4) and (23.7) for
\[
 (F,h)=(F_i,h_i),\qquad
 (\widetilde F,\widetilde h)
   =(\widehat F_{i+1},\widehat h_{i+1}).
\]
Let \(B_i,L_i,C_{F,i},C_{{\rm Lip},i}\) be corresponding stocks in
(23.5)--(23.6), and put
\[
 \beta_i:=
 B_iL_i\epsilon_{F,i}
 +C_{{\rm Lip},i}B_iL_i\epsilon_{0,i}
 +C_{F,i}L_i\epsilon_{v,i}
 +C_{F,i}B_i\epsilon_{1,i}.
 \tag{23.13}
\]
\end{definition}

\begin{corollary}[Summable projective current]
\label{cor:v11v23-projective-current}
Let \(\mathfrak T_{m\leftarrow i}\) be compatible transports of one-forms
to a common later cylinder, with
\[
 \|\mathfrak T_{m\leftarrow i}\alpha\|
 \leq C_T\|\alpha\|
 \tag{23.14}
\]
uniformly.  If
\[
 \sum_{i=i_0}^{\infty}\beta_i<\infty,
 \tag{23.15}
\]
then the transported currents form a Cauchy family, and for \(m>n\),
\[
 \left\|
 \mathfrak T_{m\leftarrow m}J_m
 -\mathfrak T_{m\leftarrow n}J_n
 \right\|
 \leq
 C_T\sum_{i=n}^{m-1}\beta_i.
 \tag{23.16}
\]
Consequently they define a projective limiting current on every fixed
cylinder for which the transports are defined.
\end{corollary}

\begin{proof}
Naturality (23.11) and Theorem
\ref{thm:v11v23-current-estimate} give an adjacent difference at most
\(\beta_i\).  Transport each adjacent difference to level \(m\).
Compatibility makes the differences telescope, while (23.14) and the
triangle inequality give (23.16).  The tail of the convergent series
(23.15) tends to zero.
\end{proof}

If the exterior derivative commutes with transport, then
\[
 dJ_\infty=F_\infty
 \tag{23.17}
\]
on cylinder forms whenever both projective limits exist.  Indeed,
\(dJ_i=F_i\) at every finite regulator by V22, and equality passes to the
limit in the topology in which \(d\) is continuous.

\subsection{Why curvature comparison is not enough}

The path rows in (23.13) cannot be dropped.  Consider the unit disk in
\(\mathbb C\), retracting to the origin, with
\[
 F=dx\wedge dy,\qquad
 h_t(z)=tz,\qquad
 \widetilde h^{(N)}_t(z)
   =t\exp\!\bigl(iN t(1-t)\bigr)z.
 \tag{23.18}
\]
Both homotopies have the same endpoints and use the same curvature.
Writing \(z=re^{i\theta}\), direct substitution into (22.3) gives
\[
 K_hF=\frac12r^2\,d\theta,
 \tag{23.19}
\]
whereas
\[
 K_{\widetilde h^{(N)}}F
 =
 \frac12r^2\,d\theta
 -N\left(\int_0^1t^2(1-2t)\,dt\right)r\,dr
 =
 \frac12r^2\,d\theta+\frac N6r\,dr.
 \tag{23.20}
\]
Thus
\[
 K_{\widetilde h^{(N)}}F-K_hF
 =d\left(\frac N{12}r^2\right).
 \tag{23.21}
\]

\begin{proposition}[Independent homotopy row]
\label{prop:v11v23-independent-row}
There is no bound on
\(\|K_hF-K_{\widetilde h}F\|_0\) in terms of
\(\|F-\widetilde F\|_0\) alone, even when the retract, endpoints, and
curvature agree exactly.
\end{proposition}

\begin{proof}
In (23.18), the curvature difference is zero for every \(N\), while the
norm of (23.21) grows linearly in \(N\) away from the origin.
\end{proof}

This example also identifies the ambiguity: the two currents differ by
an exact one-form, as predicted in V22.  Exactness controls curvature but
does not control scale convergence of the chosen phase.

\subsection{Phase transport}

Let \(s_i\) be the nonvanishing determinant-line frame whose connection
is \(J_i\), furnished conditionally by V22.  Choose compatible base
points \(x_i\) and, in every fixed cylinder, paths from the base point of
length at most \(\Lambda\).  Normalize adjacent frame phases at the base
points, allowing a residual phase error \(\eta_i\).

\begin{proposition}[Phase stability from the current card]
\label{prop:v11v23-phase-stability}
If the transports are unitary and
\[
 \sum_i(\Lambda\beta_i+\eta_i)<\infty,
 \tag{23.22}
\]
then the transported parallel phases of \(s_i\) are uniformly Cauchy on
the fixed cylinder.  For an adjacent pair their phase discrepancy is at
most
\[
 \eta_i+\Lambda\beta_i.
 \tag{23.23}
\]
\end{proposition}

\begin{proof}
The logarithmic phase difference along a chosen path is the base-point
phase difference plus the path integral of the difference of the two
connection one-forms.  Its absolute value is bounded by
\[
 \eta_i+\operatorname{length}(\gamma)\,
 \|J_i-\Phi_i^*J_{i+1}\|_0,
\]
which is at most (23.23).  Telescoping and (23.22) give uniform Cauchy
convergence.  Exponentiation is Lipschitz on the imaginary axis with
constant one.
\end{proof}

The proposition concerns the determinant-line phase.  It neither supplies
the modulus of a fermion coefficient nor removes the complementary-minor
and zero-mode saturation rows from V18.

\begin{firewall}[Unacquired adjacent data]
\label{fw:v11v23-unacquired}
No estimate in this note proves that the physical Standard Model
regulators admit the common tubes (23.9), uniformly local equivariant
homotopies, or a summable card (23.15).  The four errors in (23.12) are
new acquisition targets.  In particular, anomaly cancellation in a
continuum expansion does not by itself control
\(\epsilon_{0,i},\epsilon_{v,i},\epsilon_{1,i}\).
\end{firewall}

\begin{assessment}[V23 advance]
\label{ass:v11v23}
The scale firewall in V22 has been sharpened to the explicit bound
(23.8).  A summable four-row adjacent card now gives a projective
determinant current and a Cauchy phase, while the rotating-disk example
proves that curvature comparison alone cannot do so.  The next finite
task is to separate universal Standard Model representation arithmetic
from regulator-dependent curvature estimates, so that the cancellation
row can be stated without assuming the analytic acquisition it needs.
\end{assessment}
\section{Exact Standard Model anomaly arithmetic and the finite-cutoff firewall}
\label{sec:v11v24}

The representation-theoretic cancellation should be separated from the
analytic estimates in V23.  The former is exact and regulator
independent.  The latter are estimates for a concrete overlap kernel,
homotopy, and scale transport.  This note computes the complete
one-generation chiral arithmetic, including the mixed gravitational
term and the global \(SU(2)\) parity, and then states the precise bridge
still needed at finite cutoff.

\subsection{Left-handed convention}

Write every fermion as a left-handed Weyl field for
\[
 G_{\rm SM}=SU(3)_c\times SU(2)_L\times U(1)_Y,
 \tag{24.1}
\]
with electric charge \(Q=T_3+Y\).  One generation is
\[
\begin{array}{c|c|c|c}
 \text{field}&SU(3)_c&SU(2)_L&Y\\ \hline
 Q&(3)&(2)&1/6\\
 u^c&(\overline3)&(1)&-2/3\\
 d^c&(\overline3)&(1)&1/3\\
 L&(1)&(2)&-1/2\\
 e^c&(1)&(1)&1
\end{array}
\tag{24.2}
\]
and a sterile \(\nu^c\), if present, has \((1,1,0)\) and changes none of
the calculations.  Bosons do not contribute to the chiral determinant
line.

Normalize quadratic and cubic indices by
\[
 T(3)=T(\overline3)=\frac12,\qquad
 T(2)=\frac12,\qquad
 A(3)=1,\qquad A(\overline3)=-1.
 \tag{24.3}
\]
Here \(A(r)\) is the coefficient of the symmetric cubic trace for
\(SU(3)\).  The \(SU(2)\) symmetric cubic invariant vanishes.

\subsection{Direct calculation}

The five local anomaly columns are
\[
 SU(3)^3,\quad SU(3)^2Y,\quad SU(2)^2Y,\quad
 Y^3,\quad {\rm grav}^2Y,
 \tag{24.4}
\]
together with the identically zero \(SU(2)^3\) column.  Their individual
contributions are
\[
\begin{array}{c|ccccc}
 &SU(3)^3&SU(3)^2Y&SU(2)^2Y&Y^3&{\rm grav}^2Y\\ \hline
 Q   &2& 1/6& 1/4& 1/36& 1\\
 u^c &-1&-1/3&0&-8/9&-2\\
 d^c &-1& 1/6&0& 1/9& 1\\
 L   &0&0&-1/4&-1/4&-1\\
 e^c &0&0&0&1&1
\end{array}
\tag{24.5}
\]
where every multiplicity from the spectator representation is included.

\begin{theorem}[One-generation local anomaly cancellation]
\label{thm:v11v24-local-cancellation}
For the representation (24.2), all perturbative gauge and mixed
gauge--gravitational anomaly tensors vanish:
\[
 \mathcal A_{333}=\mathcal A_{33Y}
 =\mathcal A_{22Y}=\mathcal A_{YYY}
 =\mathcal A_{{\rm grav}{\rm grav}Y}=0.
 \tag{24.6}
\]
All remaining mixed tensors containing one generator of a semisimple
factor vanish by tracelessness.
\end{theorem}

\begin{proof}
Sum the columns of (24.5).  Explicitly,
\[
 2-1-1=0,
 \tag{24.7}
\]
\[
 \frac16-\frac13+\frac16=0,
 \qquad
 \frac14-\frac14=0,
 \tag{24.8}
\]
and
\[
 \frac1{36}-\frac89+\frac19-\frac14+1
 =
 \frac{1-32+4-9+36}{36}=0.
 \tag{24.9}
\]
The linear hypercharge sum is
\[
 1-2+1-1+1=0.
 \tag{24.10}
\]
There is no \(SU(2)^3\) symmetric tensor.  A mixed trace with only one
\(SU(3)\) or \(SU(2)\) generator factors through its trace and is
therefore zero.  These statements exhaust the degree-three invariant
polynomials on the direct-sum Lie algebra in (24.1).
\end{proof}

\subsection{Invariant-polynomial formulation}

Let
\[
 X=A\otimes1+1\otimes B+y\,c\,1
 \tag{24.11}
\]
on a multiplet \(r_3\otimes r_2\), with
\(A\in\mathfrak{su}(3)\), \(B\in\mathfrak{su}(2)\), and
\(c\in\mathbb R\).  Tracelessness gives
\begin{align*}
 \Tr_{r_3\otimes r_2}(X^3)
 &=
 \dim(r_2)\Tr_{r_3}(A^3)
 +\dim(r_3)\Tr_{r_2}(B^3)\\
 &\quad+
 3yc\bigl[
 \dim(r_2)\Tr_{r_3}(A^2)
 +\dim(r_3)\Tr_{r_2}(B^2)
 \bigr]\\
 &\quad+
 \dim(r_3)\dim(r_2)y^3c^3.
 \tag{24.12}
\end{align*}
The terms with \(A B\), \(A c^2\), and \(B c^2\) vanish.

\begin{corollary}[Vanishing anomaly polynomial]
\label{cor:v11v24-polynomial}
Let \(\rho_{\rm gen}\) be the direct sum of the five left-handed
multiplets in (24.2).  Then
\[
 \Tr_{\rho_{\rm gen}}(X^3)=0,
 \qquad
 \Tr_{\rho_{\rm gen}}(X)=0
 \tag{24.13}
\]
as invariant polynomials on \(\mathfrak g_{\rm SM}\).
Consequently the degree-six Chern--Weil anomaly polynomial of one
generation,
\[
 I_6=
 \left[\widehat A(TM)\,
       {\rm ch}_{\rho_{\rm gen}}(\mathcal F)\right]_{6},
 \tag{24.14}
\]
vanishes, independently of the normalization convention for the common
universal constants.
\end{corollary}

\begin{proof}
Equation (24.12) decomposes the cubic trace into exactly the columns in
(24.5); Theorem \ref{thm:v11v24-local-cancellation} makes every
coefficient zero.  The linear trace is (24.10).  In degree six, expansion
of (24.14) has only a cubic gauge term proportional to
\(\Tr\mathcal F^3\) and a mixed term proportional to
\(p_1(TM)\Tr\mathcal F\).  Both vanish by (24.13).
There is no pure gravitational six-form in the expansion of
\(\widehat A\).
\end{proof}

Three generations multiply the zero polynomial by three.  Yukawa
matrices and masses do not alter this ultraviolet representation
identity.

\subsection{The global \(SU(2)\) parity}

Perturbative cubic traces do not detect the mod-two \(SU(2)\) anomaly.
One generation contains three colour copies of the quark doublet and
one lepton doublet:
\[
 N_{\rm doublet}=3+1=4.
 \tag{24.15}
\]

\begin{proposition}[Absence of the doublet parity obstruction]
\label{prop:v11v24-witten}
The mod-two sign acquired by one Standard Model generation under the
nontrivial class of \(\pi_4(SU(2))\) is
\[
 (-1)^{N_{\rm doublet}}=(-1)^4=1.
 \tag{24.16}
\]
Thus the usual fundamental-doublet global \(SU(2)\) obstruction is
absent.  It remains absent for any number of complete generations.
\end{proposition}

\begin{proof}
Each left-handed fundamental doublet contributes one unit to the
mod-two spectral flow.  Spectator colour produces three copies of
\(Q\), and \(L\) produces one more.  The total parity is even, giving
(24.16).
\end{proof}

\subsection{What this proves for determinant curvature}

For orthogonal projectors \(P_s\) onto the regulated chiral subspaces,
the determinant-line curvature is additive:
\[
 F_{\rm det}
 =
 \sum_s\Tr(P_s\,dP_s\wedge dP_s).
 \tag{24.17}
\]

\begin{lemma}[Direct-sum additivity]
\label{lem:v11v24-additivity}
If \(P=\bigoplus_sP_s\), then
\[
 \Tr(P\,dP\wedge dP)
 =
 \sum_s\Tr(P_s\,dP_s\wedge dP_s).
 \tag{24.18}
\]
Tensoring a projector with a fixed spectator space multiplies its
curvature by the spectator dimension.
\end{lemma}

\begin{proof}
Both \(P\) and \(dP\) are block diagonal, so all off-diagonal products
vanish and the trace is the sum of the block traces.  For
\(P\otimes1_W\), factorization of the trace supplies
\(\Tr_W1_W=\dim W\).
\end{proof}

Lemma \ref{lem:v11v24-additivity} explains the multiplicities in
(24.5), but (24.6) does \emph{not} assert
\[
 F_{\rm det}=0
 \tag{24.19}
\]
pointwise at finite lattice spacing.  Anomaly cancellation says that the
universal gauge-obstruction class represented by the local descent of
(24.14) vanishes after the full species sum.  Gauge-invariant transverse
Berry curvature and regulator-dependent exact terms may remain.

To state the bridge precisely, let \(\Omega_i\) be the local
gauge-orbit curvature cocycle at regulator \(i\).  Suppose a uniform
local cohomology decomposition has actually been proved:
\[
 \Omega_i
 =
 \mathscr C_i(I_6)+d_{\rm g}\Lambda_i+R_i,
 \tag{24.20}
\]
where \(\mathscr C_i\) is linear in the invariant polynomial,
\(d_{\rm g}\) is the gauge-configuration exterior derivative, and
\(R_i\) is a controlled remainder.

\begin{proposition}[Conditional cohomological reduction]
\label{prop:v11v24-cohomological-reduction}
For the complete Standard Model representation, (24.20) reduces exactly
to
\[
 \Omega_i=d_{\rm g}\Lambda_i+R_i.
 \tag{24.21}
\]
If the transported adjacent differences of \(\Lambda_i\) and \(R_i\)
obey the summable local stocks required in V23, then the gauge-orbit
part of the projective determinant current converges and satisfies the
limiting integrability identity.
\end{proposition}

\begin{proof}
Corollary \ref{cor:v11v24-polynomial} gives \(I_6=0\), and linearity makes
the first term of (24.20) vanish.  Equation (24.21) follows.  Apply the
four-row estimate and telescoping corollary of V23 to the primitive and
remainder after common transport; continuity of \(d_{\rm g}\) passes the
finite-cutoff identity to the limit.
\end{proof}

\begin{firewall}[Arithmetic is not the analytic bridge]
\label{fw:v11v24-analytic-bridge}
The decomposition (24.20), a cutoff-uniform locality bound on
\(\Lambda_i\), and summability of its adjacent differences have not been
proved for the interacting renewal overlap regulators.  The exact
identities (24.6), (24.13), and (24.16) remove the known
representation-theoretic obstructions.  They do not supply the common
tubes, local homotopies, probability concentration, determinant modulus,
or Einstein-sector stress-energy estimates.
\end{firewall}

\begin{assessment}[V24 advance]
\label{ass:v11v24}
The Standard Model cancellation row is now exact and complete at the
representation level: all local gauge anomalies, the mixed
gauge--gravitational anomaly, and the fundamental-doublet global
\(SU(2)\) parity obstruction vanish generation by generation.  The
finite-cutoff implication has been stated only through the explicit
conditional bridge (24.20), preventing anomaly arithmetic from being
mistaken for a bound on the full Berry curvature.  V25 is the scheduled
YM/NS audit before the next analytic acquisition.
\end{assessment}
\section{V25 companion audit: independent legs and restricted-input norms}
\label{sec:v11v25}

This is the scheduled five-version audit.  The active companion records
inspected before choosing the next SM step were
\[
\begin{array}{c|c|r|l}
 \text{programme}&\text{version}&\text{bytes}&\text{SHA--256}\\ \hline
 \text{Yang--Mills}&30&431067&
 \texttt{9C12A3729B18AC29FC2737B7FD1FF6162DC1E8E4154762DC0EF0344F3553AA0A}\\
 \text{Navier--Stokes}&26&278283&
 \texttt{82C887F8C2C69E4F28FAD7C3DC8DE5B9099CB5A4BF431EBA1FFF3E91935978AD}
\end{array}
\tag{25.1}
\]
The NS record is unchanged from the V20 audit.  The YM record has
advanced from V29 to V30.

\begin{audit}[Results read and transferred]
\label{audit:v11v25-companions}
YM V30 gives an independent Laurent character to every differentiated
leg, proves an exact mixed-derivative word formula, and retains all
internal and background phases until differentiation.  Its reflection
completion cancels the assembled first axial derivative but leaves the
complete second derivative, including inverse and router jets.

NS V26 computes exact rank-one norms for the first and second derivatives
of the Oseen kernel.  The first rank-one norm equals the unrestricted
symmetric-traceless norm at every radius.  The second is smaller only on
a bounded radial window and gives a modest integrated improvement.
Neither result closes the parent programme.
\end{audit}

\subsection{The independent-leg rule}

The YM calculation exposes a general logical issue relevant to the
curvature derivatives in V23.  Let \(f\) be a smooth function of
independent variables \(x_1,\ldots,x_m\).  For a one-parameter
specialization \(x_\ell=\sigma_\ell q\),
\[
 \frac{d^2}{dq^2}f(\sigma_1q,\ldots,\sigma_mq)\bigg|_{q=0}
 =
 \sum_{a,b=1}^m
 \sigma_a\sigma_b\,\partial_a\partial_b f(0).
 \tag{25.2}
\]
The mixed terms must be present before any cancellation is asserted.

\begin{lemma}[Diagonal data do not recover legwise response]
\label{lem:v11v25-diagonal}
Knowledge of \(q\mapsto f(q,\ldots,q)\) alone does not determine the
individual mixed derivatives \(\partial_a\partial_bf(0)\).
\end{lemma}

\begin{proof}
For two variables, take
\[
 f(x_1,x_2)=x_1x_2,\qquad
 g(x_1,x_2)=\frac12(x_1^2+x_2^2).
 \tag{25.3}
\]
Then \(f(q,q)=g(q,q)=q^2\), but
\[
 \partial_1\partial_2f(0)=1,\qquad
 \partial_1\partial_2g(0)=0.
\]
Thus the diagonal specialization loses the legwise Hessian.
\end{proof}

For the SM current this means that gauge, metric, homotopy-velocity, and
test-vector variations must retain separate labels through the
projector and trace differentiations.  A common one-parameter
specialization may be taken only after the required multilinear tensor
has been assembled.

\subsection{Reflection and second response}

For any twice differentiable scalar \(W\), define
\[
 W_{\rm ev}(q)=\frac12\bigl(W(q)+W(-q)\bigr).
 \tag{25.4}
\]

\begin{proposition}[Assembled reflection identity]
\label{prop:v11v25-reflection}
One has
\[
 W_{\rm ev}'(0)=0,\qquad W_{\rm ev}''(0)=W''(0).
 \tag{25.5}
\]
These identities do not imply that an oriented constituent of \(W\)
has zero first derivative, nor do they delete any term in the second
derivative.
\end{proposition}

\begin{proof}
Differentiate (25.4).  The first derivative is
\(\frac12(W'(q)-W'(-q))\), while the second is
\(\frac12(W''(q)+W''(-q))\).  Evaluation at zero gives (25.5).
Linearity applies only after all constituents defining \(W\) have been
included.
\end{proof}

Thus charge-conjugate, orientation-reversed, or species-paired SM terms
may cancel after the full response sum.  They cannot be removed before
the independent variations and their transport factors are
differentiated.

\subsection{Restricted-input norm lesson}

Let \(T:E\to F\) be a multilinear kernel contraction and
\(\mathcal R\) a restricted set of unit inputs.  Define
\[
 \|T\|_{\mathcal R}:=\sup_{u\in\mathcal R}\|T(u)\|,
 \qquad
 \|T\|_{\rm all}:=\sup_{\|u\|=1}\|T(u)\|.
 \tag{25.6}
\]
Then trivially
\[
 \|T\|_{\mathcal R}\leq\|T\|_{\rm all}.
 \tag{25.7}
\]

\begin{proposition}[When a restricted norm changes no global stock]
\label{prop:v11v25-restricted}
If an admissible restricted input realizes the unrestricted supremum,
then \(\|T\|_{\mathcal R}=\|T\|_{\rm all}\), so replacing the full norm
by the restricted norm yields no improvement.  If strict improvement
holds only on a parameter set \(E\), then an integrated constant can
improve only through the contribution of \(E\).
\end{proposition}

\begin{proof}
The first statement follows by combining the realizing lower bound with
(25.7).  For a nonnegative weight \(w\),
\[
 0\leq
 \int w(\|T\|_{\rm all}-\|T\|_{\mathcal R})
 =
 \int_Ew(\|T\|_{\rm all}-\|T\|_{\mathcal R}),
 \tag{25.8}
\]
because the integrand vanishes off \(E\).  This proves the second
statement.
\end{proof}

The NS result therefore warns against expecting flat, abelian, rank-one,
or one-polarization SM estimates to improve the uniform interacting
constant unless the physical variation class is proved to remain in
that restriction.  Restricted-locus exactness remains useful as a
regression certificate.

\subsection{Direction adjustment}

The V23 card requires curvature and homotopy derivative differences.
The audit shows that these should be built from a legwise derivative
calculus before imposing species or reflection cancellations.  The next
step is therefore upstream: derive exact first and second resolvent
derivatives for a gapped Hermitian kernel, pass them to the spectral
projector, and give norm bounds with separate perturbation legs.  Those
bounds will provide a deterministic input for
\(\epsilon_{F,i}\) and \(C_{{\rm Lip},i}\) in (23.13).

\begin{decision}[Post-audit acquisition]
\label{dec:v11v25-next}
V26 will retain independent operator legs
\(X,Y\), prove the mixed second derivative of the Riesz projector, and
then differentiate
\(\Tr(P\,dP\wedge dP)\).  Standard Model species summation and symmetry
specialization will occur only after this tensor is assembled.
\end{decision}

\begin{firewall}[Audit transfer boundary]
\label{fw:v11v25-boundary}
YM V30 supplies an algebraic order-of-operations rule, not an overlap
curvature estimate.  NS V26 supplies a warning about restricted norms,
not a fermionic locality theorem.  Neither companion note supplies the
summable adjacent card, the continuum measure, or the Einstein coupling.
\end{firewall}

\begin{assessment}[V25 audit]
\label{ass:v11v25}
The required companion review is complete.  The programme direction has
been adjusted toward an independent-leg resolvent calculation.  This
uses the new YM phase lesson and the unchanged NS restricted-norm lesson
without importing either programme's unresolved conclusion.
\end{assessment}
\section{Independent-leg Riesz calculus for the overlap curvature}
\label{sec:v11v26}

The V25 audit directs the programme upstream to the gapped Hermitian
kernel.  This note derives the first and mixed second variations of its
spectral projector with independent legs.  It then differentiates the
determinant curvature and gives an explicit Lipschitz estimate.  The
result is deterministic and applies species by species before the
Standard Model sum.

\subsection{Fixed-contour hypotheses}

Let \(H(U)\) be a \(C^2\) family of finite-dimensional Hermitian
operators on a parameter chart \(\mathcal U\).  Assume a fixed positively
oriented contour \(\Gamma\) encloses the negative spectrum of every
operator in the family and no positive eigenvalue.  Put
\[
 R_U(z)=(z-H(U))^{-1},\qquad
 P(U)=\frac1{2\pi i}\int_\Gamma R_U(z)\,dz.
 \tag{26.1}
\]
Let
\[
 \ell_\Gamma=\operatorname{length}(\Gamma),\qquad
 \operatorname{dist}\bigl(\Gamma,\operatorname{spec}H(U)\bigr)
 \geq\rho>0.
 \tag{26.2}
\]
For constant tangent legs \(X,Y\), write
\[
 H_X=D H[ X],\qquad H_{XY}=D^2H[X,Y],
 \tag{26.3}
\]
and similarly for \(P_X,P_{XY}\).

\begin{lemma}[Resolvent variations with independent legs]
\label{lem:v11v26-resolvent}
For every \(z\in\Gamma\),
\[
 D R[X]=R H_X R,
 \tag{26.4}
\]
and
\[
 D^2R[X,Y]
 =
 R H_Y R H_X R
 +R H_X R H_Y R
 +R H_{XY}R.
 \tag{26.5}
\]
No commutativity between \(H_X\) and \(H_Y\) is assumed.
\end{lemma}

\begin{proof}
Differentiate
\[
 (z-H)R=I.
\]
The first derivative gives
\[
 -H_XR+(z-H)R_X=0,
\]
and multiplication by \(R\) on the left yields (26.4).
Differentiate (26.4) in the \(Y\)-leg.  Applying the product rule and
using \(R_Y=RH_YR\) gives exactly the three ordered terms in (26.5).
\end{proof}

\begin{theorem}[First and mixed second projector variations]
\label{thm:v11v26-projector-variations}
Under (26.1)--(26.3),
\[
 P_X
 =
 \frac1{2\pi i}\int_\Gamma R H_XR\,dz,
 \tag{26.6}
\]
and
\[
 P_{XY}
 =
 \frac1{2\pi i}\int_\Gamma
 \left(
 RH_YRH_XR+RH_XRH_YR+RH_{XY}R
 \right)\,dz.
 \tag{26.7}
\]
In particular \(P_{XY}=P_{YX}\), although the two ordered cubic
resolvent words need not agree separately.
\end{theorem}

\begin{proof}
The uniform distance (26.2) permits differentiation under the finite
contour integral.  Insert (26.4)--(26.5).  Symmetry follows either from
equality of mixed derivatives or by interchanging \(X,Y\) in (26.7):
the first two terms exchange and \(H_{XY}=H_{YX}\).
\end{proof}

Define
\[
 c_1:=\frac{\ell_\Gamma}{2\pi\rho^2},
 \qquad
 c_2(a,b,c):=
 \frac{\ell_\Gamma}{2\pi}
 \left(\frac{2ab}{\rho^3}+\frac c{\rho^2}\right).
 \tag{26.8}
\]

\begin{corollary}[Operator-norm stocks]
\label{cor:v11v26-projector-bounds}
The variations satisfy
\[
 \|P_X\|\leq c_1\|H_X\|,
 \tag{26.9}
\]
\[
 \|P_{XY}\|
 \leq
 c_2(\|H_X\|,\|H_Y\|,\|H_{XY}\|).
 \tag{26.10}
\]
\end{corollary}

\begin{proof}
For Hermitian \(H\), (26.2) gives
\(\|R(z)\|\leq\rho^{-1}\).  Apply this to every resolvent factor in
(26.6)--(26.7), use the contour length, and use the triangle inequality.
\end{proof}

The projection identity supplies an exact structural check:
\[
 PP_XP=0,\qquad (1-P)P_X(1-P)=0.
 \tag{26.11}
\]

\begin{proof}
Differentiate \(P^2=P\) and sandwich the result first by \(P\), then by
\(1-P\).  Each diagonal block must vanish.
\end{proof}

\subsection{Curvature and its third independent leg}

Let \(\tau\) be a trace functional satisfying
\[
 |\tau(A_1\cdots A_k)|
 \leq\kappa_\tau\prod_{j=1}^k\|A_j\|.
 \tag{26.12}
\]
For normalized matrix trace one may take \(\kappa_\tau=1\); for the
unnormalized trace on an \(N\)-dimensional space one may take
\(\kappa_\tau=N\).  Define the determinant curvature on constant tangent
legs by
\[
 F_U(X,Y)
 =
 \tau\!\left(P[P_X,P_Y]\right).
 \tag{26.13}
\]

\begin{theorem}[Exact independent-leg curvature response]
\label{thm:v11v26-curvature-response}
For a third constant leg \(Z\),
\[
 \begin{split}
 D_ZF(X,Y)
 ={}&
 \tau\!\left(P_Z[P_X,P_Y]\right)\\
 &+\tau\!\left(P[P_{ZX},P_Y]\right)
 +\tau\!\left(P[P_X,P_{ZY}]\right).
 \end{split}
 \tag{26.14}
\]
Consequently,
\[
 |D_ZF(X,Y)|
 \leq
 2\kappa_\tau
 \left(
 \|P_Z\|\|P_X\|\|P_Y\|
 +\|P_{ZX}\|\|P_Y\|
 +\|P_X\|\|P_{ZY}\|
 \right).
 \tag{26.15}
\]
\end{theorem}

\begin{proof}
Differentiate (26.13) by the product rule while keeping \(X,Y,Z\)
independent.  This gives (26.14).  For bounded operators,
\[
 \|[A,B]\|\leq2\|A\|\|B\|,
\]
and \(\|P\|=1\).  Apply (26.12) to the three terms to obtain (26.15).
\end{proof}

Equation (26.14) is the overlap analogue of the independent-leg rule in
YM V30.  A symmetry or species sum may be applied to its right-hand side
only after all three terms and the ordered resolvent words (26.7) are
present.

\subsection{An explicit curvature comparison}

Let \(H^0,H^1\) be two \(C^2\) kernel families on the same parameter
chart, and interpolate
\[
 H^s=H^0+s\Delta H,\qquad \Delta H=H^1-H^0,
 \qquad 0\leq s\leq1.
 \tag{26.16}
\]
Assume the same contour and distance \(\rho\) work for every \(H^s(U)\).
For unit legs \(X,Y\), suppose uniformly in \(s,U\)
\[
 \|H^s_X\|\leq a_X,\qquad
 \|H^s_Y\|\leq a_Y,
 \tag{26.17}
\]
and
\[
 \|\Delta H\|\leq d_0,\qquad
 \|D\Delta H[X]\|\leq d_X,\qquad
 \|D\Delta H[Y]\|\leq d_Y.
 \tag{26.18}
\]

\begin{theorem}[Finite-cutoff curvature Lipschitz bound]
\label{thm:v11v26-curvature-lipschitz}
Let \(F^0,F^1\) be the curvatures (26.13) formed from \(H^0,H^1\).
Under (26.16)--(26.18),
\[
 \begin{split}
 |(F^1-F^0)(X,Y)|
 \leq 2\kappa_\tau\bigl[
 &c_1^3d_0a_Xa_Y\\
 &+c_1a_Yc_2(d_0,a_X,d_X)\\
 &+c_1a_Xc_2(d_0,a_Y,d_Y)
 \bigr].
 \end{split}
 \tag{26.19}
\]
Also,
\[
 \|P^1-P^0\|\leq c_1d_0.
 \tag{26.20}
\]
\end{theorem}

\begin{proof}
Regard \(s\) as the independent leg \(Z\).  Then
\[
 H^s_Z=\Delta H,\quad
 H^s_{ZX}=D\Delta H[X],\quad
 H^s_{ZY}=D\Delta H[Y].
 \tag{26.21}
\]
Equations (26.9)--(26.10) give
\[
 \|P_Z\|\leq c_1d_0,\quad
 \|P_X\|\leq c_1a_X,\quad
 \|P_Y\|\leq c_1a_Y,
 \tag{26.22}
\]
\[
 \|P_{ZX}\|\leq c_2(d_0,a_X,d_X),\qquad
 \|P_{ZY}\|\leq c_2(d_0,a_Y,d_Y).
 \tag{26.23}
\]
Insert these estimates into (26.15) and integrate \(s\) from zero to
one.  This proves (26.19).  Integrating
\(\|P_Z\|\leq c_1d_0\) proves (26.20).
\end{proof}

For two base points \(U,V\) in one convex chart, apply the theorem with
\(H^0=H(U)\), \(H^1=H(V)\).  Bounds of the form
\[
 \|H(V)-H(U)\|\leq L_0\|V-U\|,
 \tag{26.24}
\]
\[
 \|D H(V)[X]-D H(U)[X]\|
 \leq L_X\|V-U\|
 \tag{26.25}
\]
turn (26.19) into an explicit base-point Lipschitz constant
\(C_{\rm Lip}\) for (23.5).  With two adjacent regulator kernels in a
common transported chart, the same formula gives an explicit candidate
for \(\epsilon_{F,i}\).

\subsection{Weighted locality version}

For a lattice kernel \(A(x,y)\), define the two-sided exponentially
weighted Schur norm
\[
 \|A\|_\mu
 :=
 \max\left\{
 \sup_x\sum_y e^{\mu d(x,y)}\|A(x,y)\|,
 \sup_y\sum_x e^{\mu d(x,y)}\|A(x,y)\|
 \right\}.
 \tag{26.26}
\]

\begin{lemma}[Weighted product bound]
\label{lem:v11v26-weighted-product}
For kernels on a metric lattice,
\[
 \|AB\|_\mu\leq\|A\|_\mu\|B\|_\mu.
 \tag{26.27}
\]
\end{lemma}

\begin{proof}
The triangle inequality
\(d(x,z)\leq d(x,y)+d(y,z)\) gives
\[
 e^{\mu d(x,z)}
 \leq e^{\mu d(x,y)}e^{\mu d(y,z)}.
\]
Insert the kernel product, sum first in \(z\) and then in \(y\), and
take the row supremum.  The column estimate is identical.
\end{proof}

\begin{corollary}[Local Riesz calculus]
\label{cor:v11v26-local-riesz}
If
\[
 \sup_{U,z\in\Gamma}\|R_U(z)\|_\mu\leq C_{R,\mu}
 \tag{26.28}
\]
and the kernel variations in (26.3) have finite weighted norms, then
(26.6)--(26.10) hold in \(\|\cdot\|_\mu\) with every
\(\rho^{-1}\) replaced by \(C_{R,\mu}\).  The curvature response
(26.14) is consequently an exponentially local multilinear kernel
after taking a local trace density.
\end{corollary}

\begin{proof}
Apply Lemma \ref{lem:v11v26-weighted-product} to each ordered product in
(26.6)--(26.7), integrate its norm along the contour, and then apply the
same product bound to (26.14).  A local trace closes the remaining
diagonal kernel index.
\end{proof}

The near-zero Combes--Thomas stock in V12 is designed to prove a bound
of the form (26.28) on good completed blocks.  V13 shows why a naive
missing-link Dirichlet completion cannot provide it.

\subsection{Insertion into the Standard Model sum}

For a species \(s\), let its constants in (26.19) carry a subscript
\(s\).  Direct-sum additivity from V24 gives
\[
 |F_{\rm SM}^1-F_{\rm SM}^0|
 \leq
 \sum_s |F_s^1-F_s^0|.
 \tag{26.29}
\]
This absolute estimate is valid without cancellation.  A sharper
species-cancelled estimate must be derived by summing the exact response
(26.14), including representation matrices and all independent legs,
before taking norms.  The anomaly identities of V24 may then cancel the
cohomological tensor; they cannot cancel terms already replaced by
absolute values in (26.29).

\begin{firewall}[Missing kernel acquisition]
\label{fw:v11v26-kernel}
The formulas above do not yet provide cutoff-uniform values of
\(\rho,\ell_\Gamma,C_{R,\mu}\), or the first and second link-and-metric
variation stocks for the physical interacting Wilson--overlap kernel.
They also do not prove a common adjacent chart or the homotopy rows of
V23.  The trace factor must be handled as a local density to avoid an
uncontrolled volume factor.
\end{firewall}

\begin{assessment}[V26 advance]
\label{ass:v11v26}
The required post-audit continuation is complete.  The first and mixed
second Riesz-projector derivatives are now exact with independent legs,
and (26.19) converts kernel differences and derivative differences into
a finite-cutoff curvature comparison.  The weighted Schur version
preserves exponential locality under a resolvent stock.  The next finite
task is to compute the first and second variation norms of the
Hermitian Wilson kernel under exponential link perturbations, including
separate gauge and metric legs, so that the abstract constants in
(26.17)--(26.18) become explicit.
\end{assessment}
\section{Explicit Wilson-link variations in exponential coordinates}
\label{sec:v11v27}

V26 reduces curvature stability to first and second variations of the
Hermitian kernel.  For gauge-link legs these variations can be bounded
without a continuum expansion.  This note computes them in the
left-trivialized exponential chart and inserts the result into (26.19).

\subsection{One-link differential calculus}

For every oriented positive link \(e=(x,\mu)\), choose anti-Hermitian
matrices \(A_e,X_e,Y_e\) and use the chart
\[
 U_e(A)=e^{A_e}U_e.
 \tag{27.1}
\]
The reverse link is
\[
 U_{\bar e}(A)=U_e(A)^{-1}=U_e^{-1}e^{-A_e}.
 \tag{27.2}
\]
Set
\[
 \|X\|_\infty=\sup_e\|X_e\|.
 \tag{27.3}
\]

\begin{lemma}[Exponential-chart link jets]
\label{lem:v11v27-link-jets}
At \(A=0\),
\[
 D U_e[X]=X_eU_e,\qquad
 D U_{\bar e}[X]=-U_e^{-1}X_e,
 \tag{27.4}
\]
and
\[
 D^2U_e[X,Y]
 =\frac12(X_eY_e+Y_eX_e)U_e,
 \tag{27.5}
\]
\[
 D^2U_{\bar e}[X,Y]
 =U_e^{-1}\frac12(X_eY_e+Y_eX_e).
 \tag{27.6}
\]
Consequently,
\[
 \|D U_{\pm e}[X]\|\leq\|X_e\|,
 \qquad
 \|D^2U_{\pm e}[X,Y]\|
 \leq\|X_e\|\|Y_e\|.
 \tag{27.7}
\]
\end{lemma}

\begin{proof}
The power series for \(e^{sX+tY}\) gives
\[
 e^{sX+tY}
 =
 I+sX+tY+\frac12s^2X^2+\frac12st(XY+YX)+\frac12t^2Y^2+O((|s|+|t|)^3)
\]
with the pure quadratic terms omitted from the display.  This proves
the forward formulas.  Apply the same expansion to
\(e^{-sX-tY}\) in (27.2); the first derivative changes sign and the
mixed second derivative does not.  Unitary multiplication preserves
operator norm, and
\[
 \left\|\frac12(XY+YX)\right\|
 \leq\|X\|\|Y\|,
\]
which proves (27.7).
\end{proof}

The symmetrization in (27.5)--(27.6) is the Hessian in one exponential
chart.  It must not be replaced by one ordered product when \(X_e\) and
\(Y_e\) fail to commute.

\subsection{Wilson-kernel variation stocks}

Use the Wilson operator convention (12.15) and
\[
 H_m(U)=\gamma_5(D_W(U)-m).
 \tag{27.8}
\]
Only the eight range-one hopping blocks in a row depend on the links.

\begin{theorem}[First and second Wilson-link bounds]
\label{thm:v11v27-wilson-jets}
For all exponential-chart legs \(X,Y\),
\[
 \|D H_m(U)[X]\|
 \leq8\|X\|_\infty,
 \tag{27.9}
\]
\[
 \|D^2H_m(U)[X,Y]\|
 \leq8\|X\|_\infty\|Y\|_\infty.
 \tag{27.10}
\]
In the weighted Schur norm (26.26),
\[
 \|D H_m(U)[X]\|_\mu
 \leq8e^\mu\|X\|_\infty,
 \tag{27.11}
\]
\[
 \|D^2H_m(U)[X,Y]\|_\mu
 \leq8e^\mu\|X\|_\infty\|Y\|_\infty.
 \tag{27.12}
\]
All four estimates are independent of the lattice volume and of \(m\).
\end{theorem}

\begin{proof}
The spin matrices \((I\pm\gamma_\mu)/2\) multiplying the forward and
backward links are orthogonal projections of norm one.  Lemma
\ref{lem:v11v27-link-jets} bounds a differentiated hopping block by
\(\|X\|_\infty\), and a mixed second derivative block by
\(\|X\|_\infty\|Y\|_\infty\).  There are four forward and four backward
blocks in every row and column.  The row--column Schur test gives
(27.9)--(27.10).  Each differentiated block still has range one, so the
weight contributes \(e^\mu\), proving (27.11)--(27.12).
The diagonal mass term is constant.
\end{proof}

\subsection{Base-link comparison of the jets}

Let
\[
 \delta=d_\infty(U,V)
 =\sup_e\|U_e-V_e\|.
 \tag{27.13}
\]
Use the same left-trivialized arrays \(X,Y\) in the two exponential
charts.

\begin{proposition}[Jet Lipschitz bounds in the base link]
\label{prop:v11v27-jet-lipschitz}
One has
\[
 \|D H_m(U)[X]-D H_m(V)[X]\|
 \leq8\delta\|X\|_\infty,
 \tag{27.14}
\]
\[
 \|D^2H_m(U)[X,Y]-D^2H_m(V)[X,Y]\|
 \leq8\delta\|X\|_\infty\|Y\|_\infty.
 \tag{27.15}
\]
The weighted versions have right-hand sides multiplied by \(e^\mu\).
\end{proposition}

\begin{proof}
For a forward link, the first-jet difference is
\(X_e(U_e-V_e)\).  For a reverse link it is
\[
 (U_e^{-1}-V_e^{-1})X_e,
\]
up to sign, and
\[
 \|U_e^{-1}-V_e^{-1}\|=\|U_e-V_e\|.
\]
Equations (27.5)--(27.6) give the same factorization for the second jet,
with \(\frac12(X_eY_e+Y_eX_e)\) in place of \(X_e\).
Sum the eight row and column blocks as in the preceding theorem.
\end{proof}

Together with V12,
\[
 \|H_m(U)-H_m(V)\|\leq8\delta.
 \tag{27.16}
\]
Thus the kernel, its first jet, and its second jet have one common
literal link constant.

\subsection{An explicit curvature constant}

Define
\[
 q_\mu:=8e^\mu,\qquad q_0:=8.
 \tag{27.17}
\]
For the unweighted norm use \(q=q_0\).  In the weighted calculus use
\(q=q_\mu\) and replace the resolvent powers in (26.8) by the weighted
resolvent stock as in Corollary \ref{cor:v11v26-local-riesz}.

\begin{theorem}[Wilson-link curvature Lipschitz stock]
\label{thm:v11v27-curvature-stock}
Assume the fixed contour hypotheses of V26 hold along the chord joining
\(U\) and \(V\) in a common chart.  For
\(\|X\|_\infty,\|Y\|_\infty\leq1\),
\[
 |(F_U-F_V)(X,Y)|
 \leq C_{\rm Wcurv}\,\delta,
 \tag{27.18}
\]
where
\[
 C_{\rm Wcurv}
 =
 2\kappa_\tau
 \left[
 q^3c_1^3
 +2q c_1 c_2(q,q,q)
 \right].
 \tag{27.19}
\]
Moreover,
\[
 \|P(U)-P(V)\|\leq q c_1\delta.
 \tag{27.20}
\]
\end{theorem}

\begin{proof}
In Theorem \ref{thm:v11v26-curvature-lipschitz}, equations
(27.9), (27.14), and (27.16) permit
\[
 a_X=a_Y=q,\qquad
 d_0=q\delta,\qquad d_X=d_Y=q\delta.
 \tag{27.21}
\]
Since \(c_2\) is linear under the simultaneous substitution
\((d_0,a_X,d_X)=(q\delta,q,q\delta)\),
\[
 c_2(q\delta,q,q\delta)=\delta c_2(q,q,q).
\]
Substitution in (26.19) gives (27.18)--(27.19), and (26.20) gives
(27.20).
\end{proof}

On the V12 closed half-gap tube, a contour can be chosen uniformly in
volume around the negative spectral interval.  The numerical quality of
(27.19) still depends on that contour; the theorem records a valid
literal stock rather than asserting optimality.

\begin{corollary}[Kato-tube sufficient condition]
\label{cor:v11v27-kato}
If
\[
 q c_1\delta<1,
 \tag{27.22}
\]
then the adjacent projectors satisfy the Kato direct-rotation condition
\[
 \|P(U)-P(V)\|<1
 \tag{27.23}
\]
used in V21.
\end{corollary}

\begin{proof}
This is immediate from (27.20).
\end{proof}

\subsection{Infinitesimal gauge legs}

For an anti-Hermitian site field \(\xi_x\), the infinitesimal gauge
variation is
\[
 \delta_\xi U_{x,\mu}
 =
 \xi_xU_{x,\mu}-U_{x,\mu}\xi_{x+\widehat\mu}
 =
 X^\xi_{x,\mu}U_{x,\mu},
 \tag{27.24}
\]
where
\[
 X^\xi_{x,\mu}
 =
 \xi_x-U_{x,\mu}\xi_{x+\widehat\mu}U_{x,\mu}^{-1}.
 \tag{27.25}
\]
Hence
\[
 \|X^\xi\|_\infty\leq2\|\xi\|_\infty.
 \tag{27.26}
\]

\begin{corollary}[Gauge-leg norm]
\label{cor:v11v27-gauge-leg}
The first kernel variation along a gauge orbit obeys
\[
 \|\delta_\xi H_m\|\leq16\|\xi\|_\infty,
 \qquad
 \|\delta_\xi H_m\|_\mu\leq16e^\mu\|\xi\|_\infty.
 \tag{27.27}
\]
\end{corollary}

\begin{proof}
Combine (27.9), (27.11), and (27.26).
\end{proof}

Gauge covariance gives the sharper exact identity
\(\delta_\xi H_m=[\xi,H_m]\) with the appropriate site multiplication
operator.  Equation (27.27) is retained because V26 needs a norm stock
for mixed gauge and transverse legs.

\subsection{Insertion into the adjacent-scale card}

After an adjacent regulator transport, suppose the coarse and fine link
representatives obey
\[
 d_\infty(U_i,\widehat U_{i+1})\leq\delta_i
 \tag{27.28}
\]
inside a common gapped chart.  Then Theorem
\ref{thm:v11v27-curvature-stock} supplies the explicit row
\[
 \epsilon_{F,i}\leq C_{{\rm Wcurv},i}\delta_i
 \tag{27.29}
\]
for each Wilson species, before the species sum.  Proposition
\ref{prop:v11v27-jet-lipschitz} supplies the first-derivative discrepancy
needed when (26.19) is applied directly.  The homotopy errors
\(\epsilon_{0,i},\epsilon_{v,i},\epsilon_{1,i}\) in V23 remain separate.

\begin{firewall}[Gauge links are only one coefficient sector]
\label{fw:v11v27-metric}
Equations (27.9)--(27.29) treat link dependence with fixed spin,
translation, and volume coefficients.  A dynamical tetrad or metric
changes those coefficients and the Hilbert-space density, and may also
change the identification between adjacent tangent spaces.  Its first,
second, and mixed gauge--metric jets have not been bounded here.
The fixed-contour and weighted-resolvent hypotheses also remain
conditional on the good-block construction.
\end{firewall}

\begin{assessment}[V27 advance]
\label{ass:v11v27}
The abstract kernel rows in V26 are explicit for exponential gauge-link
variations: the kernel, first jet, and mixed second jet all have the
volume-independent literal constant \(8\), or \(8e^\mu\) in the weighted
norm.  Formula (27.19) converts this into a curvature Lipschitz stock and
(27.22) into a Kato-tube condition.  The next finite task is the
coefficient calculus for metric-dependent hoppings and local density
weights, keeping gauge and metric legs independent.
\end{assessment}
\section{Metric-coefficient and density calculus for curved Wilson stencils}
\label{sec:v11v28}

V27 treats link dependence with fixed hopping coefficients.  Coupling to
a discrete Einstein field changes the local Clifford coefficients,
volume weights, and possibly the identification of spinor Hilbert
spaces.  This note gives a finite-dimensional coefficient card that
keeps those effects explicit.  It is deliberately independent of a
particular simplicial or hypercubic gravity discretization.

\subsection{A nearest-neighbour coefficient card}

On a fixed coordinate lattice, consider a bare finite-range Hermitian
kernel of the form
\[
 \mathscr H(U,g)
 =
 B(g)+\sum_{\alpha=1}^{N}
 C_\alpha(g)\,\mathcal U_\alpha(U)T_\alpha.
 \tag{28.1}
\]
Here \(B,C_\alpha\) are local multiplication matrices,
\(T_\alpha\) are unitary shifts of range at most \(R\), and
\(\mathcal U_\alpha\) is one oriented gauge link.  Forward and backward
hops are listed separately.  Hermiticity is imposed on the complete
sum, not on each summand.

For metric legs \(k,\ell\), define uniform coefficient stocks
\[
 A_0:=\sum_\alpha\|C_\alpha\|_\infty,
 \tag{28.2}
\]
\[
 A_1:=\sup_{\|k\|_\infty\leq1}
 \sum_\alpha\|D_gC_\alpha[k]\|_\infty,\qquad
 A_2:=\sup_{\|k\|_\infty,\|\ell\|_\infty\leq1}
 \sum_\alpha\|D_g^2C_\alpha[k,\ell]\|_\infty,
 \tag{28.3}
\]
and
\[
 b_0:=\|B\|_\infty,\qquad
 b_1:=\sup_{\|k\|_\infty\leq1}\|D_gB[k]\|_\infty,
 \qquad
 b_2:=\sup_{\|k\|_\infty,\|\ell\|_\infty\leq1}
 \|D_g^2B[k,\ell]\|_\infty.
 \tag{28.4}
\]

\begin{theorem}[Bare gauge--metric jet bounds]
\label{thm:v11v28-bare-jets}
In the exponential link chart of V27,
\[
 \|\mathscr H_X\|
 \leq A_0\|X\|_\infty,
 \qquad
 \|\mathscr H_{XY}\|
 \leq A_0\|X\|_\infty\|Y\|_\infty,
 \tag{28.5}
\]
\[
 \|\mathscr H_k\|
 \leq(b_1+A_1)\|k\|_\infty,
 \tag{28.6}
\]
\[
 \|\mathscr H_{k\ell}\|
 \leq(b_2+A_2)\|k\|_\infty\|\ell\|_\infty,
 \tag{28.7}
\]
and
\[
 \|\mathscr H_{Xk}\|
 \leq A_1\|X\|_\infty\|k\|_\infty.
 \tag{28.8}
\]
In the weighted Schur norm, the hopping contributions on the
right-hand sides are multiplied by \(e^{\mu R}\); the diagonal
\(b_j\) contributions are unchanged.
\end{theorem}

\begin{proof}
Differentiate (28.1) term by term.  Gauge legs act only on
\(\mathcal U_\alpha\), and Lemma \ref{lem:v11v27-link-jets} bounds its
first and mixed second jets by \(\|X\|_\infty\) and
\(\|X\|_\infty\|Y\|_\infty\).  Metric legs act on \(B,C_\alpha\), giving
(28.6)--(28.7).  One gauge and one metric derivative produce
\[
 \sum_\alpha
 (D_gC_\alpha[k])\,
 (D_U\mathcal U_\alpha[X])T_\alpha,
\]
which gives (28.8).  The row and column sums are bounded by the
coefficient sums (28.2)--(28.4).  A shift of range at most \(R\)
contributes at most \(e^{\mu R}\) to the weighted norm.
\end{proof}

For the flat Wilson stencil (12.15), \(R=1\), \(A_0=8\), and
\(A_1=A_2=0\), recovering V27.  If a metric-dependent spin connection
appears inside \(\mathcal U_\alpha\), it must be differentiated as an
additional local factor; it cannot be hidden inside \(A_1,A_2\) unless
the corresponding bounds have actually been included.

\subsection{Putting changing Hilbert spaces on one fibre}

Suppose the natural spinor inner product is
\[
 \langle\psi,\phi\rangle_g
 =
 \langle\psi,M_g\phi\rangle_0,
 \qquad M_g>0.
 \tag{28.9}
\]
Let
\[
 W_g=M_g^{1/2},\qquad
 K(U,g)=W_g\,\mathscr H(U,g)\,W_g^{-1}.
 \tag{28.10}
\]
Then \(W_g\) is a unitary identification from the \(g\)-weighted Hilbert
space to the fixed reference fibre, and \(K\) is Hermitian in the fixed
inner product whenever \(\mathscr H\) is Hermitian in (28.9).

For a metric leg \(k\), set
\[
 S_k=(D_gW[k])W^{-1}.
 \tag{28.11}
\]
For any operator family \(A(g)\), write
\[
 \mathcal C_g(A)=W_gAW_g^{-1}.
 \tag{28.12}
\]

\begin{lemma}[Derivative of the density identification]
\label{lem:v11v28-density-derivative}
One has
\[
 D_g\mathcal C_g(A)[k]
 =
 \mathcal C_g(A_k)+[S_k,\mathcal C_g(A)].
 \tag{28.13}
\]
Consequently,
\[
 K_k=\mathcal C_g(\mathscr H_k)+[S_k,K].
 \tag{28.14}
\]
For two constant metric legs,
\[
 \begin{split}
 K_{k\ell}
 ={}&
 \mathcal C_g(\mathscr H_{k\ell})
 +[S_\ell,\mathcal C_g(\mathscr H_k)]
 +[S_k,\mathcal C_g(\mathscr H_\ell)]\\
 &+[D_\ell S_k,K]
 +[S_k,[S_\ell,K]].
 \end{split}
 \tag{28.15}
\]
\end{lemma}

\begin{proof}
Differentiate \(WAW^{-1}\) and use
\[
 D W^{-1}[k]=-W^{-1}(D W[k])W^{-1}.
\]
The two outer-factor terms combine into the commutator in (28.13).
Equation (28.14) is the case \(A=\mathscr H\).  Differentiate (28.14)
in the \(\ell\)-leg and apply (28.13) to
\(\mathscr H_k\) and \(\mathscr H_\ell\).  The product rule for the
commutator gives the last two terms in (28.15).
\end{proof}

Although (28.15) is written by differentiating first in \(k\), it equals
the formula with \(k,\ell\) exchanged.  This follows from equality of
mixed derivatives of \(W\) and \(\mathscr H\); the resulting identity
between the \(S\)-terms is the pure-gauge Maurer--Cartan identity for the
change of fibre.

Gauge-link legs do not differentiate \(W_g\).  Hence
\[
 K_X=\mathcal C_g(\mathscr H_X),
 \tag{28.16}
\]
and the mixed gauge--metric jet is
\[
 K_{Xk}
 =
 \mathcal C_g(\mathscr H_{Xk})
 +[S_k,\mathcal C_g(\mathscr H_X)].
 \tag{28.17}
\]

\subsection{Norm stocks after density conjugation}

Assume
\[
 \chi:=\|W\|\|W^{-1}\|<\infty,
 \qquad
 \|S_k\|\leq s_1\|k\|_\infty,
 \qquad
 \|D_\ell S_k\|
 \leq s_2\|k\|_\infty\|\ell\|_\infty.
 \tag{28.18}
\]
Put \(M_0=b_0+A_0\), \(M_1=b_1+A_1\), and
\(M_2=b_2+A_2\).

\begin{theorem}[Fixed-fibre gauge--metric stocks]
\label{thm:v11v28-fixed-fibre}
For unit gauge and metric legs,
\[
 \|K\|\leq\chi M_0,
 \tag{28.19}
\]
\[
 \|K_X\|\leq\chi A_0,\qquad
 \|K_{XY}\|\leq\chi A_0,
 \tag{28.20}
\]
\[
 \|K_k\|
 \leq\chi(M_1+2s_1M_0),
 \tag{28.21}
\]
\[
 \|K_{Xk}\|
 \leq\chi(A_1+2s_1A_0),
 \tag{28.22}
\]
and
\[
 \|K_{k\ell}\|
 \leq
 \chi\left[
 M_2+4s_1M_1+2s_2M_0+4s_1^2M_0
 \right].
 \tag{28.23}
\]
The same formulas hold in a weighted Schur algebra when
\(W,W^{-1},S_k\) are local multiplication operators and the \(M_j,A_j\)
are replaced by their weighted stencil stocks.
\end{theorem}

\begin{proof}
Similarity gives
\[
 \|\mathcal C_g(A)\|\leq\chi\|A\|.
\]
Equations (28.19)--(28.20) follow from this inequality and
Theorem \ref{thm:v11v28-bare-jets}.  Use
\(\|[A,B]\|\leq2\|A\|\|B\|\) in (28.14) and (28.17) to obtain
(28.21)--(28.22).  Apply it to each commutator in (28.15):
the two first commutators contribute \(4\chi s_1M_1\), the derivative
of \(S\) contributes \(2\chi s_2M_0\), and the nested commutator
contributes \(4\chi s_1^2M_0\).  This proves (28.23).
Multiplication kernels have range zero, so the weighted proof is
identical.
\end{proof}

The condition number \(\chi\) is unavoidable in this representation.
Uniform ellipticity and nondegeneracy of the discrete metric must bound
it before the Riesz estimates can be uniform.

\subsection{Curvature components involving metric legs}

Let
\[
 a_{\rm g}:=\chi A_0,\qquad
 a_{\rm E}:=\chi(M_1+2s_1M_0),
 \tag{28.24}
\]
\[
 a_{\rm gg}:=\chi A_0,\qquad
 a_{\rm gE}:=\chi(A_1+2s_1A_0),
 \tag{28.25}
\]
\[
 a_{\rm EE}:=
 \chi[M_2+4s_1M_1+2s_2M_0+4s_1^2M_0].
 \tag{28.26}
\]
Here the subscript \({\rm g}\) denotes a gauge-link leg and
\({\rm E}\) a metric leg.

\begin{corollary}[Projector mixed-leg stocks]
\label{cor:v11v28-projector-mixed}
Under the fixed contour hypotheses of V26,
\[
 \|P_X\|\leq c_1a_{\rm g},\qquad
 \|P_k\|\leq c_1a_{\rm E},
 \tag{28.27}
\]
\[
 \|P_{XY}\|\leq c_2(a_{\rm g},a_{\rm g},a_{\rm gg}),
 \tag{28.28}
\]
\[
 \|P_{Xk}\|\leq c_2(a_{\rm g},a_{\rm E},a_{\rm gE}),
 \tag{28.29}
\]
and
\[
 \|P_{k\ell}\|\leq c_2(a_{\rm E},a_{\rm E},a_{\rm EE}).
 \tag{28.30}
\]
\end{corollary}

\begin{proof}
Insert (28.24)--(28.26) into
Corollary \ref{cor:v11v26-projector-bounds}.
\end{proof}

Thus all three determinant-curvature components
\[
 F({\rm g},{\rm g}),\qquad
 F({\rm g},{\rm E}),\qquad
 F({\rm E},{\rm E})
 \tag{28.31}
\]
have explicit finite-cutoff bounds through (26.13), and their
base-point derivatives have explicit bounds through (26.15).  The
mixed gauge--Einstein component must be kept in the full determinant
line even though the mixed anomaly polynomial vanishes after the
Standard Model species sum.

\subsection{Adjacent density transports}

Let \(W_i,W_{i+1}\) be the density identifications after the coarse and
fine metrics have been represented on one cylinder.  In addition to
coefficient differences, an adjacent card must record
\[
 \delta_{W,i}:=
 \|W_{i+1}W_i^{-1}-I\|,
 \tag{28.32}
\]
and the corresponding first-leg differences
\[
 \delta_{S,i}:=
 \sup_{\|k\|\leq1}
 \|S_{i+1,k}-S_{i,k}\|.
 \tag{28.33}
\]
These are transport errors, not gauge-field errors.

\begin{proposition}[Necessity of the density row]
\label{prop:v11v28-density-row}
Even if two bare kernels agree,
\(\mathscr H_{i+1}=\mathscr H_i\), their fixed-fibre kernels can differ
when \(W_{i+1}W_i^{-1}\) fails to commute with \(\mathscr H_i\).
To first order, the difference is the commutator
\[
 \delta K=[\delta W\,W^{-1},K].
 \tag{28.34}
\]
\end{proposition}

\begin{proof}
Apply (28.13) to a path \(W(t)\) with the bare operator fixed.
The derivative is exactly the commutator in (28.34).
\end{proof}

Therefore coefficient convergence alone does not populate the
\(\epsilon_{F,i}\) row of V23.  The density and tangent-fibre transports
must converge as well.

\begin{firewall}[Geometric acquisition remains]
\label{fw:v11v28-geometric}
No particular renewal Einstein discretization has yet been shown to fit
(28.1) with cutoff-uniform
\[
 A_j,\ b_j,\ \chi,\ s_j,\ R.
\]
Uniform ellipticity, orientation, spin-structure compatibility, and a
local choice of \(M_g^{1/2}\) are separate acquisitions.  A
metric-dependent spin connection requires additional differentiated
path factors.  The note supplies the exact bookkeeping and estimates
once those geometric coefficients are specified; it does not construct
the coupled Einstein measure.
\end{firewall}

\begin{assessment}[V28 advance]
\label{ass:v11v28}
The fermionic gauge and Einstein variations now live on one fixed fibre.
Equations (28.20)--(28.23) give explicit pure-gauge, pure-metric, and
mixed jets, including the commutators created by the changing volume
density.  These constants feed directly into the independent-leg Riesz
calculus of V26.  The next finite task is to formulate an adjacent
elliptic-metric cylinder card that makes
\(\delta_{W,i},\delta_{S,i}\), the stencil coefficients, and their first
two derivatives summable without assuming the continuum Einstein
measure already exists.
\end{assessment}
\section{Adjacent elliptic-metric cylinder card}
\label{sec:v11v29}

V28 identifies the metric-dependent quantities needed by the overlap
curvature.  This note proves that they are summably stable under a
concrete \(C^2\) adjacent metric card, provided the local metrics remain
in one uniformly elliptic compact set and the coefficient maps have
uniform third derivatives.  No probability law is used.

\subsection{Common-cylinder metric data}

Fix a finite cylinder parameter chart \(\mathcal Q\).  After adjacent
transport, let
\[
 g_i,\widehat g_{i+1}:\mathcal Q\longrightarrow
 {\rm Sym}^+_4
 \tag{29.1}
\]
be the coarse and fine local metric matrices at all stencil sites.
Assume
\[
 \lambda I\leq g_i(q),\widehat g_{i+1}(q)\leq\Lambda I
 \tag{29.2}
\]
for fixed \(0<\lambda\leq\Lambda<\infty\).  Define the adjacent
discrepancies
\[
 \delta_{0,i}:=
 \sup_q\|g_i(q)-\widehat g_{i+1}(q)\|_\infty,
 \tag{29.3}
\]
\[
 \delta_{1,i}:=
 \sup_{q,\|v\|\leq1}
 \|Dg_i(q)[v]-D\widehat g_{i+1}(q)[v]\|_\infty,
 \tag{29.4}
\]
\[
 \delta_{2,i}:=
 \sup_{q,\|v\|,\|w\|\leq1}
 \|D^2g_i(q)[v,w]
   -D^2\widehat g_{i+1}(q)[v,w]\|_\infty.
 \tag{29.5}
\]
Suppose also
\[
 \|Dg_i\|,\|D\widehat g_{i+1}\|\leq L_1,\qquad
 \|D^2g_i\|,\|D^2\widehat g_{i+1}\|\leq L_2.
 \tag{29.6}
\]

\subsection{Composition estimates}

Let \(c\) be any local matrix-valued stencil coefficient on the
ellipticity set
\[
 \mathcal K_{\lambda,\Lambda}
 =
 \{g:\lambda I\leq g\leq\Lambda I\}.
 \tag{29.7}
\]
Assume
\[
 \sup_{\mathcal K_{\lambda,\Lambda}}\|D^jc\|\leq M_j,
 \qquad j=1,2,3.
 \tag{29.8}
\]

\begin{theorem}[\(C^2\) coefficient transport]
\label{thm:v11v29-composition}
With \(c_i=c\circ g_i\) and
\(\widehat c_{i+1}=c\circ\widehat g_{i+1}\),
\[
 \|c_i-\widehat c_{i+1}\|_{C^0}
 \leq M_1\delta_{0,i},
 \tag{29.9}
\]
\[
 \|Dc_i-D\widehat c_{i+1}\|_{C^0}
 \leq
 M_2L_1\delta_{0,i}+M_1\delta_{1,i},
 \tag{29.10}
\]
and
\[
 \begin{split}
 \|D^2c_i-D^2\widehat c_{i+1}\|_{C^0}
 \leq{}&
 (M_3L_1^2+M_2L_2)\delta_{0,i}\\
 &+2M_2L_1\delta_{1,i}
 +M_1\delta_{2,i}.
 \end{split}
 \tag{29.11}
\]
\end{theorem}

\begin{proof}
The segment between the two metrics remains positive definite and in
\(\mathcal K_{\lambda,\Lambda}\).  The mean-value formula gives
(29.9).  The chain rule gives
\[
 D(c\circ g)[v]=Dc(g)\,Dg[v].
\]
Add and subtract
\(Dc(g_i)D\widehat g_{i+1}[v]\).  The derivative of \(Dc\) is bounded by
\(M_2\), while (29.4) bounds the remaining factor difference, proving
(29.10).

For the second derivative,
\[
 D^2(c\circ g)[v,w]
 =
 D^2c(g)[Dg[v],Dg[w]]
 +Dc(g)D^2g[v,w].
 \tag{29.12}
\]
In the first term, changing the base point costs
\(M_3L_1^2\delta_{0,i}\), and changing its two arguments costs
\(2M_2L_1\delta_{1,i}\).  In the second term, changing the base point
costs \(M_2L_2\delta_{0,i}\), and changing the metric Hessian costs
\(M_1\delta_{2,i}\).  Summing gives (29.11).
\end{proof}

Apply this theorem separately to every \(B,C_\alpha\) in (28.1) and sum
the finite stencil.  It gives explicit adjacent differences of all
stocks entering (28.5)--(28.8).

\subsection{The scalar volume-density row}

For the standard coordinate volume density in four dimensions,
\[
 M_g=(\det g)^{1/2}I,\qquad
 W_g=M_g^{1/2}=(\det g)^{1/4}I.
 \tag{29.13}
\]
At each site,
\[
 \lambda\leq W_g\leq\Lambda,
 \qquad
 \|W_g\|\|W_g^{-1}\|\leq\frac{\Lambda}{\lambda}.
 \tag{29.14}
\]
For a metric variation \(k\),
\[
 S_k=(D W[k])W^{-1}
 =\frac14\operatorname{tr}(g^{-1}k)I.
 \tag{29.15}
\]

\begin{lemma}[Explicit density derivative stocks]
\label{lem:v11v29-density-stocks}
If \(\|k\|\leq1\), then
\[
 \|S_k\|\leq\lambda^{-1}.
 \tag{29.16}
\]
For constant matrix legs \(k,\ell\),
\[
 \|D_\ell S_k\|\leq\lambda^{-2}.
 \tag{29.17}
\]
For cylinder legs \(v,w\), using (29.6),
\[
 \|D_wS_v\|
 \leq
 \lambda^{-2}L_1^2+\lambda^{-1}L_2.
 \tag{29.18}
\]
\end{lemma}

\begin{proof}
Jacobi's determinant formula gives
\[
 D\log W[k]=\frac14\operatorname{tr}(g^{-1}k).
\]
In four dimensions,
\[
 |\operatorname{tr}A|\leq4\|A\|,
 \qquad \|g^{-1}\|\leq\lambda^{-1},
\]
which proves (29.16).  Differentiating \(g^{-1}\) gives
\[
 D_\ell S_k
 =
 -\frac14\operatorname{tr}(g^{-1}\ell g^{-1}k),
\]
and hence (29.17).  When \(k=Dg[v]\) varies along the cylinder,
\[
 D_wS_v
 =
 \frac14\operatorname{tr}
 \left(
 -g^{-1}Dg[w]g^{-1}Dg[v]
 +g^{-1}D^2g[v,w]
 \right).
\]
Equations (29.2) and (29.6) give (29.18).
\end{proof}

\begin{proposition}[Adjacent density and connection bounds]
\label{prop:v11v29-density-adjacent}
For the scalar density (29.13),
\[
 \|W_{\widehat g_{i+1}}W_{g_i}^{-1}-I\|
 \leq
 \frac{\Lambda}{\lambda^2}\delta_{0,i},
 \tag{29.19}
\]
and, for unit cylinder legs,
\[
 \|S_{i+1}-S_i\|
 \leq
 \frac{L_1}{\lambda^2}\delta_{0,i}
 +\frac1\lambda\delta_{1,i}.
 \tag{29.20}
\]
\end{proposition}

\begin{proof}
The differential of \(W=(\det g)^{1/4}\) obeys
\[
 |DW[k]|=|W S_k|
 \leq\frac{\Lambda}{\lambda}\|k\|.
\]
The mean-value formula and \(W_{g_i}^{-1}\leq\lambda^{-1}\) give
(29.19).  For (29.20), write
\[
 \frac14\operatorname{tr}
 \left(
 \widehat g^{-1}D\widehat g
 -g^{-1}Dg
 \right)
\]
as an inverse difference times \(D\widehat g\), plus
\(g^{-1}(D\widehat g-Dg)\).  The resolvent identity gives
\[
 \|\widehat g^{-1}-g^{-1}\|
 \leq\lambda^{-2}\|\widehat g-g\|.
\]
Use the four-dimensional trace bound, (29.4), and (29.6).
\end{proof}

Thus the abstract density errors (28.32)--(28.33) are explicit linear
combinations of \(\delta_{0,i}\) and \(\delta_{1,i}\) in this scalar
volume model.

\subsection{Projective metric and coefficient limits}

Assume compatible transports so that all adjacent differences may be
telescoped on a fixed cylinder.

\begin{theorem}[Summable elliptic metric card]
\label{thm:v11v29-summable-card}
If
\[
 \sum_{i=i_0}^{\infty}
 \left(\delta_{0,i}+\delta_{1,i}+\delta_{2,i}\right)<\infty,
 \tag{29.21}
\]
then the transported metrics form a Cauchy family in \(C^2\) on every
fixed cylinder.  Their limit \(g_\infty\) remains uniformly elliptic:
\[
 \lambda I\leq g_\infty\leq\Lambda I.
 \tag{29.22}
\]
Every coefficient satisfying (29.8), together with the scalar density
\(W\) through second cylinder order and its connection \(S\) through first cylinder order,
also forms a Cauchy family.  The adjacent differences are bounded by a
constant times
\(\delta_{0,i}+\delta_{1,i}+\delta_{2,i}\).
\end{theorem}

\begin{proof}
Transported adjacent metric differences telescope.  Equations
(29.3)--(29.5) and (29.21) make the sequence Cauchy in \(C^2\).
The cone inequalities (29.2) are closed under norm limits, proving
(29.22).  Equations (29.9)--(29.11) prove the coefficient statement.
Equations (29.15) and its derivative formulas are smooth rational
expressions in \(g^{-1},Dg,D^2g\); uniform ellipticity bounds every
inverse.  Alternatively, apply Theorem
\ref{thm:v11v29-composition} to \(W\) and to the displayed expressions.
The adjacent bounds are linear combinations of the three summable
errors.
\end{proof}

\subsection{Consequence for the fermionic curvature card}

Suppose, in addition to (29.21):
\begin{enumerate}[label=\textnormal{(\roman*)}]
\item the finite stencil and its first three metric derivatives obey
uniform coefficient bounds on \(\mathcal K_{\lambda,\Lambda}\);
\item the gauge links and their transported tangent maps obey the V27
adjacent card;
\item one contour and one weighted resolvent stock work uniformly along
the adjacent interpolations; and
\item local trace densities have a uniform \(\kappa_\tau\).
\end{enumerate}

\begin{corollary}[Summable gauge--Einstein curvature row]
\label{cor:v11v29-curvature-row}
Under these hypotheses, the adjacent differences of the pure-gauge,
mixed gauge--metric, and pure-metric determinant curvatures are bounded
by
\[
 C\left(
 \delta_{{\rm link},i}
 +\delta_{0,i}+\delta_{1,i}+\delta_{2,i}
 \right)
 \tag{29.23}
\]
on each fixed cylinder, with \(C\) determined explicitly by the V26
Riesz constants and the V28 coefficient stocks.  If the right-hand side
is summable, then the three curvature components have projective
limits.
\end{corollary}

\begin{proof}
Theorems \ref{thm:v11v29-composition} and
\ref{thm:v11v29-summable-card} bound the bare kernel and its first two
independent-leg differences.  Proposition
\ref{prop:v11v29-density-adjacent} and Lemma
\ref{lem:v11v29-density-stocks} do the same for the fixed-fibre
conjugation.  Insert those differences into the curvature comparison
proof of Theorem \ref{thm:v11v26-curvature-lipschitz}, once for each
choice of two curvature legs.  Uniform stocks collect into \(C\).
Telescoping a summable adjacent bound gives each projective limit.
\end{proof}

\begin{firewall}[Probability and dynamics]
\label{fw:v11v29-probability}
The summability assumption (29.21) has not been proved under a renewal
Einstein measure.  Uniform ellipticity can fail on degenerate metric
configurations, and no estimate here shows that their local intensity
vanishes.  The corollary also assumes the spectral gap and link card; it
does not derive them from the Einstein action.  Finally, curvature
convergence does not supply the homotopy rows of V23 or the bosonic
stress-energy limit.
\end{firewall}

\begin{assessment}[V29 advance]
\label{ass:v11v29}
The Einstein coefficient problem has been reduced to the explicit
three-row metric card (29.3)--(29.5).  Under uniform ellipticity and
summability, the metric, local stencil, scalar volume identification,
with the derivative orders required in V28 converge, and (29.23) yields the
fermionic gauge--Einstein curvature limit.  V30 is the next mandatory
YM/NS audit; the following substantive note must then address either
the homotopy stability left by V23 or the probability of the elliptic
good set.
\end{assessment}
\section{V30 companion audit: nonlinear chart jets}
\label{sec:v11v30}

This is the next compulsory five-version audit.  The newest append
descendant in the Yang--Mills stream is V31; its V30 predecessor is also
present in the shared folder and has not been modified.  The
Navier--Stokes stream remains at V26.  The records used here are
\[
\begin{array}{c|c|r|l}
 \text{programme}&\text{version}&\text{bytes}&\text{SHA--256}\\ \hline
 \text{Yang--Mills}&31&448263&
 \texttt{55CAEE957B49FF407D9B5B523D0D2B75182E1FBE1BF8277FED154AB27FD0A229}\\
 \text{Navier--Stokes}&26&278283&
 \texttt{82C887F8C2C69E4F28FAD7C3DC8DE5B9099CB5A4BF431EBA1FFF3E91935978AD}
\end{array}
\tag{30.1}
\]

\begin{audit}[New YM correction and unchanged NS result]
\label{audit:v11v30-companions}
YM V31 finds that a background and vertical fluctuation placed in one
exponential have a symmetrized mixed jet, whereas the declared
left-translated product chart has an ordered mixed jet.  It quarantines
earlier response totals affected by that mismatch, constructs the
correct ordered factor compiler, and independently replays the Wilson
suffix.  Its corrected bubble response remains incomplete until the
quartic tadpole is recomputed.

NS V26 is unchanged.  Its rank-one first-derivative norm still equals
the unrestricted norm, and its second-derivative improvement remains
local and modest.
\end{audit}

\subsection{Exact chart lesson}

For noncommuting matrices \(B,Z\),
\[
 \partial_s\partial_t e^{sB+tZ}\big|_{0}
 =\frac12(BZ+ZB),
 \tag{30.2}
\]
while
\[
 \partial_s\partial_t
 \left(e^{sB}e^{tZ}\right)\big|_{0}
 =BZ.
 \tag{30.3}
\]
Thus
\[
 D^2_{\rm product}-D^2_{\rm single}
 =\frac12[B,Z].
 \tag{30.4}
\]
This does not contradict V27: equations (27.5)--(27.6) are the Hessian
of the explicitly declared single exponential chart.  They cannot be
inserted unchanged into a calculation whose actual coordinates are a
product of background and vertical exponentials.

The general mechanism is the second-order chain rule.  Let
\(\phi:\mathcal X\to\mathcal U\) be a \(C^2\) chart map and let \(H\) be
a \(C^2\) operator family.

\begin{lemma}[Second jet under a chart map]
\label{lem:v11v30-chain}
For independent legs \(X,Y\),
\[
 D(H\circ\phi)[X]
 =
 DH[D\phi\,X],
 \tag{30.5}
\]
\[
 D^2(H\circ\phi)[X,Y]
 =
 D^2H[D\phi\,X,D\phi\,Y]
 +DH[D^2\phi[X,Y]].
 \tag{30.6}
\]
\end{lemma}

\begin{proof}
Equation (30.5) is the first-order chain rule.  Differentiate it in the
\(Y\)-leg.  The derivative may hit \(DH\), giving the first term of
(30.6), or \(D\phi\,X\), giving the second.
\end{proof}

\begin{corollary}[Tangent agreement is insufficient at second order]
\label{cor:v11v30-second-order}
If two chart maps \(\phi,\widetilde\phi\) agree at a base point and have
the same first derivative there, then the first kernel jets agree, but
\[
 D^2(H\circ\phi)[X,Y]
 -D^2(H\circ\widetilde\phi)[X,Y]
 =
 DH\!\left[
 (D^2\phi-D^2\widetilde\phi)[X,Y]
 \right].
 \tag{30.7}
\]
\end{corollary}

\begin{proof}
Subtract (30.5)--(30.6) for the two maps.  Equality of their values and
first derivatives cancels every term except the displayed chart-Hessian
difference.
\end{proof}

The correction is a first derivative of \(H\) applied to a second-order
coordinate defect.  It must enter the \(H_{XY}\) row of the Riesz formula
(26.7).

\subsection{Curvature versus coordinate response}

The determinant curvature itself is a two-form:
\[
 (\phi^*F)_x(X,Y)
 =
 F_{\phi(x)}(D\phi_xX,D\phi_xY).
 \tag{30.8}
\]
It therefore needs only the first chart derivative at a fixed point.
A derivative of curvature in a third coordinate leg does involve
\(D^2\phi\), because the transported tangent legs vary.

\begin{proposition}[Derivative of a pulled-back two-form]
\label{prop:v11v30-pullback-response}
For constant chart legs \(X,Y,Z\),
\[
 \begin{split}
 D_Z(\phi^*F)(X,Y)
 ={}&
 (D_{D\phi Z}F)(D\phi X,D\phi Y)\\
 &+F(D^2\phi[Z,X],D\phi Y)
 +F(D\phi X,D^2\phi[Z,Y]).
 \end{split}
 \tag{30.9}
\]
\end{proposition}

\begin{proof}
Differentiate (30.8).  The derivative can hit the base point of \(F\),
its first tangent argument, or its second tangent argument.  These three
possibilities are the three lines in (30.9).
\end{proof}

Hence the V23 adjacent current card needs a chart-transition second-jet
row whenever its homotopy tangent maps are compared in nonlinear
coordinates.

\subsection{Direction adjustment}

Let \(\psi_i\) denote the transition from the level-\(i\) chart to the
pulled-back level-\(i+1\) chart.  Besides its point and tangent errors,
record
\[
 \epsilon_{\psi,2,i}
 :=
 \sup_{\|X\|,\|Y\|\leq1}
 \|D^2\psi_i[X,Y]\|.
 \tag{30.10}
\]
Corollary \ref{cor:v11v30-second-order} and Proposition
\ref{prop:v11v30-pullback-response} show exactly where this error enters.

\begin{decision}[V31 chart-covariant comparison]
\label{dec:v11v30-next}
The next substantive note will prove a second-jet transition estimate
for the projector, curvature, and homotopy current.  It will recover the
V27 symmetric chart when the transition Hessian vanishes and the ordered
product chart when its commutator term is retained.
\end{decision}

\begin{firewall}[Audit boundary]
\label{fw:v11v30-boundary}
The corrected YM word compiler is a finite noncommutative derivative
certificate.  It does not provide the SM chart-transition norm,
spectral gap, determinant phase, metric probability, or continuum
measure.  The unchanged NS restricted-input calculation supplies no new
transfer at this checkpoint.
\end{firewall}

\begin{assessment}[V30 audit]
\label{ass:v11v30}
The required audit is complete and has found a real upstream issue.
V27 remains correct in its declared exponential chart, while V30 now
forbids using that symmetric second jet in a translated product chart
without the chain term (30.7).  V31 will add the missing transition
second-jet stock before any further scale claim.
\end{assessment}
\section{Chart-covariant second-jet comparison}
\label{sec:v11v31}

V30 isolates a nonlinear coordinate defect.  A tensor at one point uses
only the first chart derivative, but the second variation of the kernel
and the derivative of determinant curvature also use the chart Hessian.
This note gives explicit comparison bounds through that order and then
checks how a homotopy current transforms.

\subsection{Composite-kernel discrepancies}

Let \(\phi,\widetilde\phi:\mathcal Q\to\mathcal U\) be two \(C^2\)
chart maps into one physical parameter chart.  Assume
\[
 \|\phi-\widetilde\phi\|_{C^0}\leq\varepsilon_0,\qquad
 \|D\phi-D\widetilde\phi\|_{C^0}\leq\varepsilon_1,
 \tag{31.1}
\]
\[
 \|D^2\phi-D^2\widetilde\phi\|_{C^0}
 \leq\varepsilon_2,
 \tag{31.2}
\]
and
\[
 \|D\phi\|,\|D\widetilde\phi\|\leq L_1,\qquad
 \|D^2\phi\|,\|D^2\widetilde\phi\|\leq L_2.
 \tag{31.3}
\]
Let \(H,\widetilde H\) be physical kernel families with same-point
differences
\[
 \|H-\widetilde H\|_{C^0}\leq\Delta_0,\quad
 \|DH-D\widetilde H\|_{C^0}\leq\Delta_1,\quad
 \|D^2H-D^2\widetilde H\|_{C^0}\leq\Delta_2.
 \tag{31.4}
\]
Suppose
\[
 \|D^j\widetilde H\|_{C^0}\leq M_j,\qquad j=1,2,3.
 \tag{31.5}
\]
Define the two kernels on the common cylinder
\[
 G=H\circ\phi,\qquad
 \widetilde G=\widetilde H\circ\widetilde\phi.
 \tag{31.6}
\]

\begin{theorem}[Kernel transition through second order]
\label{thm:v11v31-kernel-transition}
The composite differences obey
\[
 d_0:=\|G-\widetilde G\|_{C^0}
 \leq
 \Delta_0+M_1\varepsilon_0,
 \tag{31.7}
\]
\[
 d_1:=\|DG-D\widetilde G\|_{C^0}
 \leq
 \Delta_1L_1+M_2L_1\varepsilon_0+M_1\varepsilon_1,
 \tag{31.8}
\]
and
\[
 \begin{split}
 d_2:=\|D^2G-D^2\widetilde G\|_{C^0}
 \leq{}&
 \Delta_2L_1^2+\Delta_1L_2\\
 &+(M_3L_1^2+M_2L_2)\varepsilon_0\\
 &+2M_2L_1\varepsilon_1+M_1\varepsilon_2.
 \end{split}
 \tag{31.9}
\]
\end{theorem}

\begin{proof}
At fixed physical argument, the family difference contributes
\(\Delta_0,\Delta_1,\Delta_2\).  Moving the argument from \(\phi\) to
\(\widetilde\phi\) is controlled by the mean-value formula and
(31.5).  The first-order chain rule is
\[
 D(H\circ\phi)=DH(\phi)D\phi,
\]
which gives (31.8) after adding and subtracting the same-point family
term and the intermediate tangent factor.

For the second derivative use
\[
 D^2(H\circ\phi)[X,Y]
 =
 D^2H(\phi)[D\phi X,D\phi Y]
 +DH(\phi)D^2\phi[X,Y].
 \tag{31.10}
\]
The same-point family difference costs
\(\Delta_2L_1^2+\Delta_1L_2\).  The base-point change of
\(D^2\widetilde H\) costs \(M_3L_1^2\varepsilon_0\), changing its two
arguments costs \(2M_2L_1\varepsilon_1\), the base-point change of
\(D\widetilde H\) costs \(M_2L_2\varepsilon_0\), and the chart Hessian
costs \(M_1\varepsilon_2\).  Their sum is (31.9).
\end{proof}

If \(H=\widetilde H\), the \(\Delta_j\) vanish.  If the two charts agree
to first order at a base point but have different Hessians, then
\(d_0=d_1=0\) there while
\[
 d_2\leq M_1\varepsilon_2.
 \tag{31.11}
\]
The product-versus-single-exponential commutator in (30.4) is exactly
such a term.

\subsection{Projector jets in two charts}

Assume one contour \(\Gamma\) stays at distance at least \(\rho\) from
the spectra of \(G,\widetilde G\).  Suppose, for all unit cylinder legs,
\[
 \|G_X\|,\|\widetilde G_X\|\leq a,\qquad
 \|G_{XY}\|,\|\widetilde G_{XY}\|\leq b.
 \tag{31.12}
\]
Let \(P,\widetilde P\) be the negative spectral projectors and define
\[
 q_0:=c_1d_0,
 \qquad
 q_1:=c_2(d_0,a,d_1),
 \tag{31.13}
\]
\[
 q_2:=
 \frac{\ell_\Gamma}{2\pi}
 \left[
 \frac{6a^2d_0}{\rho^4}
 +\frac{4ad_1}{\rho^3}
 +\frac{2bd_0}{\rho^3}
 +\frac{d_2}{\rho^2}
 \right].
 \tag{31.14}
\]

\begin{theorem}[Projector transition through second order]
\label{thm:v11v31-projector-transition}
For unit independent legs \(X,Y\),
\[
 \|P-\widetilde P\|\leq q_0,
 \tag{31.15}
\]
\[
 \|P_X-\widetilde P_X\|\leq q_1,
 \tag{31.16}
\]
\[
 \|P_{XY}-\widetilde P_{XY}\|\leq q_2.
 \tag{31.17}
\]
\end{theorem}

\begin{proof}
The resolvent identity gives
\[
 R-\widetilde R
 =
 R(G-\widetilde G)\widetilde R,
 \qquad
 \|R-\widetilde R\|\leq\rho^{-2}d_0.
 \tag{31.18}
\]
Integrating (31.18) around the contour proves (31.15).

Subtract the two first-jet integrands \(RG_XR\).  Changing either
resolvent costs \(a\rho^{-3}d_0\), and changing the kernel jet costs
\(\rho^{-2}d_1\).  There are two resolvent positions, giving precisely
\(q_1\).

For one ordered cubic word
\[
 RG_YRG_XR,
\]
there are three resolvent positions and two first-jet positions.
Their total difference is at most
\[
 3a^2\rho^{-4}d_0+2a\rho^{-3}d_1.
\]
The mixed projector formula (26.7) has two such ordered cubic words.
Its remaining word \(RG_{XY}R\) has difference at most
\[
 2b\rho^{-3}d_0+\rho^{-2}d_2.
\]
Multiply the sum by \(\ell_\Gamma/(2\pi)\) to obtain (31.14)--(31.17).
\end{proof}

This proof keeps the two ordered cubic words separate.  No commutation
or reflection cancellation is used.

\subsection{Curvature and curvature-response comparison}

Let both projectors use the same trace functional \(\tau\) satisfying
(26.12).  Put
\[
 p_1:=c_1a,\qquad p_2:=c_2(a,a,b).
 \tag{31.19}
\]

\begin{theorem}[Chart-covariant curvature stocks]
\label{thm:v11v31-curvature-stocks}
For unit \(X,Y\),
\[
 |(F-\widetilde F)(X,Y)|
 \leq
 2\kappa_\tau
 \left[
 c_1^3d_0a^2+2c_1a\,c_2(d_0,a,d_1)
 \right].
 \tag{31.20}
\]
For a third unit leg \(Z\),
\[
 \begin{split}
 |D_ZF(X,Y)-D_Z\widetilde F(X,Y)|
 \leq\kappa_\tau\bigl[
 &6p_1^2q_1+4q_0p_1p_2\\
 &+4q_2p_1+4p_2q_1
 \bigr].
 \end{split}
 \tag{31.21}
\]
\end{theorem}

\begin{proof}
Equation (31.20) is Theorem
\ref{thm:v11v26-curvature-lipschitz} applied to the two composite
kernels on the common cylinder.

For (31.21), subtract the three exact response terms in (26.14).
In the first,
\[
 \tau(P_Z[P_X,P_Y]),
\]
changing \(P_Z\) costs at most
\(2\kappa_\tau q_1p_1^2\), while changing the commutator costs at most
\(4\kappa_\tau p_1^2q_1\).  For each of the other two terms, changing
the outer \(P\) costs
\(2\kappa_\tau q_0p_2p_1\), and changing the commutator costs
\(2\kappa_\tau(q_2p_1+p_2q_1)\).  Sum the first term and the two
identical bounds to obtain (31.21).
\end{proof}

At a point where only the chart Hessians differ, (31.20) vanishes but
(31.21) may not, because \(q_2\) contains \(d_2\).  This is the precise
location of the YM V31 correction in the SM curvature calculus.

\subsection{Pulled-back currents and conjugated homotopies}

A physical one-form \(J\) transforms by
\[
 (\phi^*J)_x(X)=J_{\phi(x)}(D\phi_xX).
 \tag{31.22}
\]

\begin{proposition}[One-form chart comparison]
\label{prop:v11v31-oneform}
Suppose \(\|J\|_0\leq C_J\), \(J\) is base-point Lipschitz with constant
\(L_J\), and both chart derivatives are bounded by \(L_1\).  Then
\[
 \|\phi^*J-\widetilde\phi^*J\|_0
 \leq
 L_JL_1\varepsilon_0+C_J\varepsilon_1.
 \tag{31.23}
\]
\end{proposition}

\begin{proof}
Add and subtract
\(J_{\widetilde\phi(x)}(D\phi X)\).  Moving the base point costs the
first term of (31.23); changing the tangent map costs the second.
\end{proof}

Thus a current itself uses \(C^1\) chart comparison.  Its derivative, or
a second kernel jet used to construct its Lipschitz stock, uses the
\(C^2\) row.

For completeness, let \(h_t\) be a homotopy in one chart and let a
near-identity diffeomorphism \(\psi\) define
\[
 \widehat h_t=\psi^{-1}\circ h_t\circ\psi.
 \tag{31.24}
\]
Assume
\[
 \|\psi-\operatorname{id}\|_{C^0}\leq e_0,\qquad
 \|D\psi-I\|_{C^0}\leq e_1<1,
 \tag{31.25}
\]
\[
 \|Dh_t\|\leq L_h,\quad \|\partial_th_t\|\leq B_h,
 \tag{31.26}
\]
and that \(\partial_th_t\) and \(Dh_t\) are base-point Lipschitz with
constants \(K_v,K_A\).

\begin{lemma}[Induced homotopy errors]
\label{lem:v11v31-homotopy-errors}
The two homotopies obey the V23 comparison rows with
\[
 \epsilon_0^{\,h}\leq(1+L_h)e_0,
 \tag{31.27}
\]
\[
 \epsilon_v^{\,h}
 \leq
 \frac{e_1}{1-e_1}B_h+K_ve_0,
 \tag{31.28}
\]
and
\[
 \epsilon_1^{\,h}
 \leq
 \frac{e_1}{1-e_1}L_h(1+e_1)
 +K_Ae_0(1+e_1)+L_he_1.
 \tag{31.29}
\]
\end{lemma}

\begin{proof}
For (31.27), insert and subtract \(h_t(\psi x)\) and use the displacement
bound for \(\psi^{-1}\).  Differentiating (31.24) in \(t\) gives
\[
 \partial_t\widehat h_t
 =
 D\psi^{-1}_{h_t(\psi x)}\,\partial_th_t(\psi x).
\]
Since
\[
 \|D\psi^{-1}\|\leq(1-e_1)^{-1},\qquad
 \|D\psi^{-1}-I\|\leq\frac{e_1}{1-e_1},
\]
adding and subtracting \(\partial_th_t(x)\) proves (31.28).

Finally,
\[
 D\widehat h_t
 =
 D\psi^{-1}_{h_t(\psi x)}
 Dh_t(\psi x)D\psi_x.
\]
Change the three factors one at a time.  The inverse factor, the base
point of \(Dh_t\), and the final factor contribute the three terms in
(31.29).
\end{proof}

Inserting (31.27)--(31.29) into (23.13) gives the chart-induced part of
the adjacent homotopy-current error.  The chart Hessian is not an extra
term in the value of a pulled-back one-form; it is required when
controlling second kernel jets, curvature derivatives, or variation of
the tangent transition over a cylinder.

\begin{corollary}[Summable chart-transition criterion]
\label{cor:v11v31-summable-chart}
Assume all stocks in this note are uniform and
\[
 \sum_i
 \left(
 \Delta_{0,i}+\Delta_{1,i}+\Delta_{2,i}
 +\varepsilon_{0,i}+\varepsilon_{1,i}+\varepsilon_{2,i}
 \right)<\infty.
 \tag{31.30}
\]
Then the projector through second cylinder order, the curvature through
first cylinder order, and the chart-induced portion of the V23 current
card are summable after transport.
\end{corollary}

\begin{proof}
Equations (31.7)--(31.9) make \(d_0,d_1,d_2\) summable.
Equations (31.13)--(31.21) then make the corresponding projector and
curvature differences summable.  Equations (31.23) and
(31.27)--(31.29) give the current and homotopy rows.  Compatible
transport and telescoping complete the proof.
\end{proof}

\begin{firewall}[Transition acquisition]
\label{fw:v11v31-transition}
No physical adjacent chart has yet been constructed with the summable
second-jet condition (31.30).  In particular, the nonlinear product
chart used for a background split must be compared with the chart used
by the renewal transport, rather than silently identified.  Uniform
gap, locality, metric probability, determinant modulus, and the
homotopy itself remain separate acquisitions.
\end{firewall}

\begin{assessment}[V31 advance]
\label{ass:v11v31}
The post-audit coordinate issue is closed at the deterministic level.
Chart point, tangent, and Hessian errors now propagate explicitly to
the kernel, two projector jets, curvature, curvature response, and
homotopy current.  A pure Hessian mismatch leaves the curvature value
unchanged at the base point but can change its derivative, exactly as
the corrected YM calculation indicates.  The next finite task is to
construct a canonical adjacent chart from a local logarithm on the
V12 good-link tube and bound its first two transition jets.
\end{assessment}
\section{Principal-logarithm charts on the good-link tube}
\label{sec:v11v32}

V31 leaves the chart-transition rows abstract.  On a linkwise chord ball
of radius less than one, the principal matrix logarithm supplies a
canonical coordinate with explicit derivatives.  This note constructs
that chart, compares two nearby reference links through second order,
and obtains a radial local homotopy whose chart-induced V23 errors are
summable when the reference differences are summable.

\subsection{Logarithm on a matrix chord ball}

Let \(Q\) be unitary and
\[
 \|Q-I\|\leq r<1.
 \tag{32.1}
\]
The principal logarithm is represented by the norm-convergent series
\[
 \Log Q
 =
 \sum_{n=1}^{\infty}\frac{(-1)^{n+1}}n(Q-I)^n.
 \tag{32.2}
\]
It is anti-Hermitian and obeys
\[
 \|\Log Q\|\leq-\log(1-r).
 \tag{32.3}
\]

\begin{lemma}[Derivative stocks for the local logarithm]
\label{lem:v11v32-log-stocks}
On the ball \(\|Q-I\|\leq r<1\),
\[
 \|D\Log_Q[X]\|
 \leq\frac1{1-r}\|X\|,
 \tag{32.4}
\]
\[
 \|D^2\Log_Q[X,Y]\|
 \leq\frac1{(1-r)^2}\|X\|\|Y\|,
 \tag{32.5}
\]
and
\[
 \|D^3\Log_Q[X,Y,Z]\|
 \leq\frac2{(1-r)^3}\|X\|\|Y\|\|Z\|.
 \tag{32.6}
\]
\end{lemma}

\begin{proof}
Differentiate the series (32.2) term by term, which is valid uniformly
on every closed ball of radius \(r<1\).  The first derivative of the
\(n\)-th power has \(n\) insertion positions; after the factor \(1/n\),
its norm sum is \(r^{n-1}\).  Thus
\[
 \sum_{n\geq1}r^{n-1}=(1-r)^{-1}.
\]
The second derivative has \(n(n-1)\) ordered insertion positions, so
after division by \(n\) its scalar majorant is
\[
 \sum_{n\geq2}(n-1)r^{n-2}=(1-r)^{-2}.
\]
The third derivative has \(n(n-1)(n-2)\) ordered positions and majorant
\[
 \sum_{n\geq3}(n-1)(n-2)r^{n-3}
 =2(1-r)^{-3}.
\]
These give (32.4)--(32.6).
\end{proof}

For anti-Hermitian \(A\), the Duhamel formula gives
\[
 D\exp_A[X]
 =
 \int_0^1e^{(1-s)A}Xe^{sA}\,ds.
 \tag{32.7}
\]

\begin{lemma}[Exponential derivative stocks]
\label{lem:v11v32-exp-stocks}
For anti-Hermitian \(A\) and anti-Hermitian legs,
\[
 \|D\exp_A[X]\|\leq\|X\|,
 \tag{32.8}
\]
\[
 \|D^2\exp_A[X,Y]\|\leq\|X\|\|Y\|,
 \tag{32.9}
\]
\[
 \|D^3\exp_A[X,Y,Z]\|
 \leq\|X\|\|Y\|\|Z\|.
 \tag{32.10}
\]
\end{lemma}

\begin{proof}
Equation (32.8) follows from (32.7), because every exponential factor
is unitary.  Iterated Duhamel expansion writes the \(k\)-th derivative
as a sum over \(k!\) orderings integrated over ordered simplices of
volume \(1/k!\).  Each integrand is a product of unitary factors and the
\(k\) inserted legs.  Summing the simplex volumes gives one, proving
(32.9)--(32.10).
\end{proof}

\subsection{A gauge-covariant link chart}

Fix a unitary reference link field \(R\).  On
\[
 \mathcal B_r(R)
 =
 \{U:d_\infty(U,R)\leq r\},\qquad r<1,
 \tag{32.11}
\]
define, link by link,
\[
 A_e(U;R):=\Log(U_eR_e^{-1}),
 \qquad
 U_e=e^{A_e}R_e.
 \tag{32.12}
\]
Equations (32.3)--(32.6) hold in the link supremum norm, with constants
independent of the number of links.

\begin{proposition}[Gauge covariance of the log chart]
\label{prop:v11v32-gauge-covariance}
If \(U\) and \(R\) are transformed by the same lattice gauge
transformation, then
\[
 A_{x,\mu}(U^g;R^g)
 =
 g_xA_{x,\mu}(U;R)g_x^{-1}.
 \tag{32.13}
\]
\end{proposition}

\begin{proof}
The relative link transforms by source conjugation:
\[
 U^g_{x,\mu}(R^g_{x,\mu})^{-1}
 =
 g_x(U_{x,\mu}R_{x,\mu}^{-1})g_x^{-1}.
\]
The convergent series (32.2) commutes with conjugation, yielding
(32.13).
\end{proof}

Thus an equivariant family of reference links gives an equivariant atlas.
A fixed reference need not be invariant under the full gauge group; its
gauge transforms label companion patches.

\subsection{Transition between two references}

Let \(R,S\) be reference link fields and set
\[
 C_e=R_eS_e^{-1},\qquad
 \eta=d_\infty(R,S)=\sup_e\|C_e-I\|.
 \tag{32.14}
\]
If \(U=e^AR=e^{A'}S\), then the coordinate transition is
\[
 A'=\Psi_C(A):=\Log(e^AC).
 \tag{32.15}
\]
Assume
\[
 \|e^A-I\|\leq r_0,\qquad
 r:=r_0+\eta<1.
 \tag{32.16}
\]
The segment between \(e^A\) and \(e^AC\) lies in the norm ball of radius
\(r\) about \(I\).

Define
\[
 L_{\log,1}=(1-r)^{-1},\quad
 L_{\log,2}=(1-r)^{-2},\quad
 L_{\log,3}=2(1-r)^{-3}.
 \tag{32.17}
\]

\begin{theorem}[Second-jet reference transition]
\label{thm:v11v32-transition}
On the domain (32.16),
\[
 \|\Psi_C(A)-A\|
 \leq L_{\log,1}\eta,
 \tag{32.18}
\]
\[
 \|D\Psi_C(A)-I\|
 \leq
 \left(L_{\log,2}+L_{\log,1}\right)\eta,
 \tag{32.19}
\]
and
\[
 \|D^2\Psi_C(A)\|
 \leq
 \left(
 L_{\log,3}+3L_{\log,2}+L_{\log,1}
 \right)\eta.
 \tag{32.20}
\]
All norms are operator norms on the link supremum space, and the
constants are volume independent.
\end{theorem}

\begin{proof}
When \(C=I\), the principal branch gives
\[
 \Psi_I(A)=\Log(e^A)=A
 \tag{32.21}
\]
on the stated domain.  Move \(C\) to \(I\) along the matrix segment.
The input to \(\Log\) moves by at most \(\eta\); (32.4) gives (32.18).

For a leg \(X\),
\[
 D\Psi_C(A)[X]
 =
 D\Log_{e^AC}\!\left[D\exp_A[X]\,C\right].
 \tag{32.22}
\]
Compare this with the same formula at \(C=I\).  Moving the base point of
\(D\Log\) costs \(L_{\log,2}\eta\|X\|\), by (32.5) and (32.8).
Changing the final factor \(C\) costs
\(L_{\log,1}\eta\|X\|\).  Since the expression at \(C=I\) is \(X\),
(32.19) follows.

The second derivative is
\[
 \begin{split}
 D^2\Psi_C(A)[X,Y]
 ={}&
 D^2\Log_{e^AC}
 \left[D\exp_A[X]C,D\exp_A[Y]C\right]\\
 &+
 D\Log_{e^AC}\left[D^2\exp_A[X,Y]C\right].
 \end{split}
 \tag{32.23}
\]
Compare (32.23) with its value at \(C=I\), which is zero by (32.21).
Moving the base point of \(D^2\Log\) costs
\(L_{\log,3}\eta\|X\|\|Y\|\).  Changing its two arguments costs
\(2L_{\log,2}\eta\|X\|\|Y\|\).  Moving the base point of
\(D\Log\) costs another
\(L_{\log,2}\eta\|X\|\|Y\|\), and changing the last factor \(C\)
costs \(L_{\log,1}\eta\|X\|\|Y\|\).
Equations (32.8)--(32.9) give (32.20).
\end{proof}

For the product chart, the nonzero Hessian in (32.20) retains ordered
commutator information.  Replacing it by zero would reproduce the
coordinate error identified in V30.

\subsection{Radial homotopy inside the chord tube}

In logarithm coordinates define
\[
 h_t(A)=tA,\qquad0\leq t\leq1.
 \tag{32.24}
\]
In link variables,
\[
 U_e(t)=e^{tA_e}R_e.
 \tag{32.25}
\]
If \(e^{i\theta}\) is an eigenvalue of \(e^A\), (32.1) with \(r<1\)
places \(\theta\) on the principal arc \(|\theta|<\pi\), and
\[
 |e^{it\theta}-1|
 \leq|e^{i\theta}-1|.
 \tag{32.26}
\]
Hence (32.25) stays in \(\mathcal B_r(R)\).

\begin{proposition}[Radial homotopy stocks]
\label{prop:v11v32-radial}
The coordinate homotopy (32.24) satisfies
\[
 \|\partial_th_t(A)\|
 =\|A\|
 \leq-\log(1-r),
 \qquad
 \|D_Ah_t\|\leq1.
 \tag{32.27}
\]
It retracts the patch to \(R\), fixes \(R\), and is gauge equivariant
for the equivariant reference atlas of Proposition
\ref{prop:v11v32-gauge-covariance}.
\end{proposition}

\begin{proof}
The derivative identities follow directly from (32.24), and the norm
bound is (32.3).  Equations (32.13) and (32.24) commute with source
conjugation.  At \(A=0\), every \(h_t\) is zero.
\end{proof}

If \(R\) is a commuting flat background, V19 provides the canonical
determinant frame at the retract.  V22 therefore gives on this patch
\[
 J_R=K_{h^R}F,\qquad dJ_R=F,
 \tag{32.28}
\]
and, from (22.15),
\[
 \|J_R\|_0
 \leq C_F[-\log(1-r)].
 \tag{32.29}
\]

\subsection{Adjacent references and current summability}

For adjacent reference fields \(R_i,R_{i+1}\), let
\[
 \eta_i=d_\infty(R_i,\widehat R_{i+1}).
 \tag{32.30}
\]
Assume all chart domains fit in one fixed pair
\(r_0+\eta_i\leq r<1\).  Theorem
\ref{thm:v11v32-transition} supplies V31 chart errors
\[
 \varepsilon_{0,i}\leq L_{\log,1}\eta_i,
 \tag{32.31}
\]
\[
 \varepsilon_{1,i}
 \leq(L_{\log,2}+L_{\log,1})\eta_i,
 \tag{32.32}
\]
\[
 \varepsilon_{2,i}
 \leq(L_{\log,3}+3L_{\log,2}+L_{\log,1})\eta_i.
 \tag{32.33}
\]

\begin{corollary}[Summable log-atlas row]
\label{cor:v11v32-summable-atlas}
If
\[
 \sum_i\eta_i<\infty,
 \tag{32.34}
\]
then all three chart-transition rows (32.31)--(32.33) are summable.
Under the uniform curvature and radial-homotopy stocks above, the
chart-induced part of the adjacent current card in V23 is summable.
If the remaining physical curvature differences are summable, the
patch currents have a projective limit.
\end{corollary}

\begin{proof}
The constants in (32.31)--(32.33) depend only on the fixed \(r<1\), so
(32.34) makes every transition row summable.  Apply Corollary
\ref{cor:v11v31-summable-chart} and the induced homotopy estimates
(31.27)--(31.29).  Theorem \ref{thm:v11v23-current-estimate} then adds
the physical curvature row.  Its telescoping corollary gives the
projective current.
\end{proof}

\subsection{Compatibility with the V12 gap tube}

Around a commuting flat reference \(R\), V12 gives the closed half-gap
condition
\[
 d_\infty(U,R)\leq\frac{\gamma_m}{16}.
 \tag{32.35}
\]
Since \(0<\gamma_m\leq1\), this lies inside the logarithm radius
\(r_0\leq1/16<1\).  Thus the principal-log chart and radial homotopy are
available everywhere on the certified V12 half-gap tube.  With
\(m=1\),
\[
 -\log(1-r_0)\leq-\log(15/16).
 \tag{32.36}
\]
This supplies an explicit homotopy-velocity constant in (32.27).

\begin{firewall}[Atlas coverage and topology]
\label{fw:v11v32-coverage}
The construction is local to a chord ball around a chosen reference.
It does not choose a global smooth flat reference \(R(U)\), cover
topologically nontrivial sectors, or prove that the overlap cocycle
between all reference patches is a coboundary.  Tree gauge in V14
produces near-identity representatives only on good contractible
blocks.  Probability concentration in those blocks and summability of
the adjacent reference differences (32.34) remain open.
\end{firewall}

\begin{assessment}[V32 advance]
\label{ass:v11v32}
A canonical, gauge-covariant chart and homotopy now exist on every V12
good-link tube.  The logarithm, exponential, and change-of-reference
maps have explicit volume-independent derivatives through the order
needed by V31, and a summable adjacent reference chord implies a
summable chart and homotopy row.  The next finite task is to analyze
overlaps of these logarithm patches: compute the exact difference of
their relative currents and reduce phase gluing to a quantitative
Cech cocycle with summable adjacent transition phases.
\end{assessment}
\section{Two-parameter transgression on overlaps of logarithm patches}
\label{sec:v11v33}

V32 constructs a current on each good-link patch.  Two patches may use
different flat references and hence different radial homotopies.  Their
currents have the same curvature, so their difference is closed; a
homotopy of homotopies gives an explicit primitive.  This note derives
that primitive, its norm, and its adjacent-scale comparison.

\subsection{Relative two-parameter formula}

Let \(F\) be a closed two-form on a finite regulator tube.  Let
\[
 H:[0,1]_s\times[0,1]_t\times\mathcal O\longrightarrow\mathcal T
 \tag{33.1}
\]
be smooth and satisfy
\[
 H(0,t,x)=h^\alpha_t(x),\qquad
 H(1,t,x)=h^\beta_t(x),
 \tag{33.2}
\]
\[
 H(s,1,x)=x,
 \tag{33.3}
\]
while
\[
 H(s,0,x)\in\mathcal F
 \tag{33.4}
\]
for a flat locus on which \(F\) restricts to zero.  Define
\[
 c_{\alpha\beta}(x)
 :=
 -\int_0^1\int_0^1
 F_{H(s,t,x)}
 \left(\partial_sH,\partial_tH\right)\,dt\,ds.
 \tag{33.5}
\]
Let
\[
 J_\alpha=K_{h^\alpha}F,\qquad
 J_\beta=K_{h^\beta}F.
 \tag{33.6}
\]

\begin{theorem}[Overlap transgression identity]
\label{thm:v11v33-overlap}
Under (33.1)--(33.6),
\[
 J_\beta-J_\alpha=dc_{\alpha\beta}.
 \tag{33.7}
\]
\end{theorem}

\begin{proof}
Pull \(F\) back by \(H\) and write its components as
\(F_{st},F_{tv},F_{sv}\), where \(v\) is a tangent leg on
\(\mathcal O\).  Since \(dF=0\) and the coordinate vector fields commute,
\[
 \partial_sF_{tv}-\partial_tF_{sv}+\partial_vF_{st}=0.
 \tag{33.8}
\]
Integrating (33.8) over the square gives
\[
 \begin{split}
 d\!\left(\int\!\!\int F_{st}\,dt\,ds\right)(v)
 ={}&
 -\int_0^1
 \left.F_{tv}\right|_{s=0}^{s=1}\,dt\\
 &+\int_0^1
 \left.F_{sv}\right|_{t=0}^{t=1}\,ds.
 \end{split}
 \tag{33.9}
\]
The first boundary is \(-(J_\beta-J_\alpha)(v)\).  At \(t=1\),
(33.3) makes \(\partial_sH=0\).  At \(t=0\), both
\(\partial_sH\) and \(D_xH\,v\) are tangent to \(\mathcal F\), so
(33.4) and \(F|_{\mathcal F}=0\) make the remaining boundary zero.
The minus sign in (33.5) then yields (33.7).
\end{proof}

This is a relative formula: the moving \(t=0\) endpoint is allowed
because it stays inside the zero-curvature flat locus.

\subsection{Interpolation of two flat references}

Let \(R_\alpha,R_\beta\) be commuting flat references in the same
torus sector.  Suppose
\[
 \eta_{\alpha\beta}
 :=
 d_\infty(R_\alpha,R_\beta)<1.
 \tag{33.10}
\]
The principal commuting logarithm defines a flat reference path
\[
 R_s=
 \exp\!\left(
 s\Log(R_\beta R_\alpha^{-1})
 \right)R_\alpha.
 \tag{33.11}
\]
For \(U\) in a common logarithm domain of every \(R_s\), set
\[
 H(s,t,U)
 =
 \exp\!\left(
 t\Log(UR_s^{-1})
 \right)R_s.
 \tag{33.12}
\]
Then (33.2)--(33.4) hold for the two V32 radial homotopies.

Let
\[
 a_r:=-\log(1-r),
 \qquad
 b_R:=
 \sup_s\|(\partial_sR_s)R_s^{-1}\|_\infty.
 \tag{33.13}
\]
For (33.11),
\[
 b_R
 =
 \|\Log(R_\beta R_\alpha^{-1})\|_\infty
 \leq-\log(1-\eta_{\alpha\beta}).
 \tag{33.14}
\]

\begin{lemma}[Reference-interpolation velocity]
\label{lem:v11v33-reference-velocity}
If
\[
 \sup_{s,U}d_\infty(U,R_s)\leq r<1,
 \tag{33.15}
\]
then the two-parameter map (33.12) obeys
\[
 \|\partial_tH\|_\infty\leq a_r,
 \tag{33.16}
\]
\[
 \|\partial_sH\|_\infty
 \leq
 b_R\left(1+\frac1{1-r}\right).
 \tag{33.17}
\]
\end{lemma}

\begin{proof}
Put \(A_s=\Log(UR_s^{-1})\).  The \(t\)-derivative of
\(e^{tA_s}R_s\) has norm at most \(\|A_s\|\), which is bounded by
(32.3).  For the \(s\)-derivative, the relative link
\(Q_s=UR_s^{-1}\) satisfies
\[
 \|\partial_sQ_s\|
 \leq b_R.
\]
Equation (32.4) gives
\[
 \|\partial_sA_s\|
 \leq(1-r)^{-1}b_R.
\]
Differentiate (33.12), apply the exponential Duhamel bound, and add the
derivative of \(R_s\).  This gives (33.17).
\end{proof}

\begin{corollary}[Quantitative overlap phase]
\label{cor:v11v33-overlap-bound}
If
\[
 |F_x(u,v)|\leq C_F\|u\|\|v\|
 \tag{33.18}
\]
throughout the two-parameter image, then
\[
 |c_{\alpha\beta}(U)|
 \leq
 C_Fa_rb_R\left(1+\frac1{1-r}\right).
 \tag{33.19}
\]
\end{corollary}

\begin{proof}
Insert (33.16)--(33.18) into (33.5) and integrate over a square of area
one.
\end{proof}

For nearby references, (33.14) makes the right-hand side tend linearly
to zero.  The estimate retains the finite flat-holonomy sector rather
than setting it to zero.

\subsection{Determinant-frame transition}

Let \(s_\alpha,s_\beta\) be the V22 determinant frames rephased so that
their connections are \(J_\alpha,J_\beta\).  Normalize them using the
canonical parallel V19 frame along the flat path (33.11).

\begin{proposition}[Explicit overlap transition]
\label{prop:v11v33-frame-transition}
On a simply connected overlap,
\[
 s_\beta=g_{\alpha\beta}s_\alpha,\qquad
 g_{\alpha\beta}=\exp(c_{\alpha\beta}),
 \tag{33.20}
\]
up to one constant phase fixed by the flat-path normalization.  In
particular,
\[
 g_{\beta\alpha}=g_{\alpha\beta}^{-1}.
 \tag{33.21}
\]
\end{proposition}

\begin{proof}
For line frames related by \(s_\beta=g s_\alpha\), their connection
forms satisfy
\[
 J_\beta-J_\alpha=g^{-1}dg=d\log g.
\]
Equation (33.7) is solved by \(g=\exp(c_{\alpha\beta})\).
The normalization fixes the integration constant.  Reversing the
\(s\)-orientation in (33.5) changes the sign of \(c\), proving
(33.21).
\end{proof}

On a triple overlap,
\[
 d(c_{\alpha\beta}+c_{\beta\gamma}+c_{\gamma\alpha})=0.
 \tag{33.22}
\]
Thus the sum is locally constant, and
\[
 g_{\alpha\beta}g_{\beta\gamma}g_{\gamma\alpha}
 \tag{33.23}
\]
is a constant phase.  With compatible determinant-frame
normalizations it equals one.  If a nontrivial constant remains, it is
a genuine Cech cocycle datum and cannot be removed by a local
curvature estimate.

\subsection{Adjacent-scale stability of the transgression}

Let \((F,H)\) and \((\widetilde F,\widetilde H)\) be two adjacent
two-parameter data sets represented in one chart.  Assume
\[
 \|F-\widetilde F\|_0\leq\epsilon_F,\qquad
 \|\widetilde F\|_0\leq C_F,
 \tag{33.24}
\]
that \(\widetilde F\) has base-point Lipschitz constant
\(C_{\rm Lip}\), and
\[
 \|H-\widetilde H\|\leq\epsilon_0,
 \tag{33.25}
\]
\[
 \|\partial_sH-\partial_s\widetilde H\|\leq\epsilon_s,\qquad
 \|\partial_tH-\partial_t\widetilde H\|\leq\epsilon_t.
 \tag{33.26}
\]
Suppose both \(s\)-velocities are bounded by \(B_s\) and both
\(t\)-velocities by \(B_t\).

\begin{theorem}[Adjacent overlap-phase estimate]
\label{thm:v11v33-phase-estimate}
The scalar transgressions satisfy
\[
 |c_{\alpha\beta}-\widetilde c_{\alpha\beta}|
 \leq
 B_sB_t\epsilon_F
 +C_{\rm Lip}B_sB_t\epsilon_0
 +C_FB_t\epsilon_s
 +C_FB_s\epsilon_t.
 \tag{33.27}
\]
\end{theorem}

\begin{proof}
In the integrand of (33.5), add and subtract successively:
the curvature at the same base point, the curvature after moving the
base point, the changed \(s\)-velocity, and the changed \(t\)-velocity.
The four differences are bounded by the four terms in (33.27).
Integrate over the unit square.
\end{proof}

Define the right-hand side of (33.27) to be
\(\zeta_{\alpha\beta,i}\) at adjacent scale \(i\).

\begin{corollary}[Projective transition phases]
\label{cor:v11v33-projective-phases}
If, on every overlap in a fixed finite cylinder cover,
\[
 \sum_i\zeta_{\alpha\beta,i}<\infty
 \tag{33.28}
\]
and the flat-path normalization phases have a summable adjacent
difference, then the transported functions
\(g_{\alpha\beta,i}\) converge uniformly.  Their limiting transitions
retain (33.21) and every exact triple-overlap identity satisfied at
finite cutoff.
\end{corollary}

\begin{proof}
Equation (33.27) makes the logarithmic transitions Cauchy after their
base phases are included.  Exponentiation is Lipschitz on the imaginary
axis.  Uniform limits preserve products and inverses, so the finite
cocycle identities pass to the limit.
\end{proof}

\subsection{A finite gluing criterion}

Let \(\mathcal N\) be the nerve graph of a finite patch cover on one
cylinder.  Assume the restricted transition functions have been pulled to
one common connected comparison domain, so every graph-cycle product is
pointwise defined.  Choose a spanning tree and a root patch.  Define patch phases
recursively along the tree using \(g_{\alpha\beta}\).

\begin{proposition}[Cycle-holonomy criterion]
\label{prop:v11v33-cycle}
The overlap frames glue to one global frame on the covered cylinder if
and only if the product of transition phases around every independent
non-tree cycle is one.  On a tree nerve, gluing is automatic.
\end{proposition}

\begin{proof}
Along the spanning tree, recursion defines a phase on every vertex
without ambiguity.  For a non-tree edge, the two endpoint phases agree
with its prescribed transition exactly when the product around the
unique cycle formed by that edge and the tree path is one.  These
fundamental cycles generate the cycle space.  Necessity follows by
telescoping the ratios of any global frame around a cycle.
\end{proof}

Without the common-domain hypothesis one must use the full Cech nerve,
including its higher intersections.  This is a finite certificate.  For a
cofinal atlas one must also control
the number and geometry of relevant cycles, or work cylinder by
cylinder with compatible restrictions.

\begin{firewall}[Unproved phase topology]
\label{fw:v11v33-topology}
A two-parameter flat interpolation has only been constructed for
references in one common commuting torus sector and with the common
logarithm bound (33.15).  References in different topological sectors
may not be connected inside \(\mathcal F\).  No estimate here proves
that every cycle holonomy in Proposition
\ref{prop:v11v33-cycle} is one, or that the phase comparison series
(33.28) is summable under the coupled measure.
\end{firewall}

\begin{assessment}[V33 advance]
\label{ass:v11v33}
The overlap ambiguity left in V21--V22 is now an explicit
two-parameter curvature integral.  It has a finite norm bound, a
four-row adjacent comparison, and a cycle-holonomy gluing certificate.
For nearby commuting flat references its size is controlled by their
principal-log chord.  The next finite task is to separate the
topological sector label from local logarithm patches and prove that
cylinder observables require only a locally finite phase nerve with
stable cycle generators.
\end{assessment}
\section{Finite cylinder phase nerves and the sector ledger}
\label{sec:v11v34}

V33 reduces phase gluing to overlap cycles but leaves open whether a
cofinal atlas has controllably many cycles.  The correct separation is
local.  On a fixed transported cylinder, compactness gives a finite
logarithm-patch cover of the declared good set.  Global topological
sectors are then summed through a separate tightness ledger; they are
not edges of the local patch nerve.

\subsection{Finite cover of a fixed good cylinder}

Let \(\mathcal Q_K\) be the finite-dimensional parameter space for a
fixed cylinder \(K\) after every regulator has been transported to the
same coordinates.  Let
\[
 \mathcal G_K\subset\mathcal Q_K
 \tag{34.1}
\]
be a closed good set on which the V12 gap and the V14 completion
conditions hold.  Since the gauge group is compact and the metric
variables are restricted to the ellipticity set (29.2),
\(\mathcal G_K\) is compact.

Assume every point of \(\mathcal G_K\) lies in a V32 logarithm patch
\[
 \mathcal B_{r_\alpha}(R_\alpha),\qquad r_\alpha<1,
 \tag{34.2}
\]
whose reference \(R_\alpha\) is flat in the relevant local sector.

\begin{theorem}[Finite local phase atlas]
\label{thm:v11v34-finite-atlas}
There are finitely many references
\(R_1,\ldots,R_N\) such that
\[
 \mathcal G_K
 \subset
 \bigcup_{\alpha=1}^N\mathcal B_{r_\alpha}(R_\alpha).
 \tag{34.3}
\]
Consequently the overlap nerve and its cycle space are finite on the
fixed cylinder.
\end{theorem}

\begin{proof}
The patches (34.2) form an open cover of the compact set
\(\mathcal G_K\).  Extract a finite subcover.  A finite cover has a
finite nerve.  Its one-skeleton has finitely many edges, so its cycle
space is finitely generated.
\end{proof}

The theorem does not prove the covering hypothesis.  It says that once
the good-block completion supplies local flat references, no infinite
patch inventory is needed for one fixed cylinder.

\subsection{Stability of the nerve}

For simplicity take a common patch radius \(r\) in a compact metric
comparison domain \(Q\).  For every nonempty subset
\(S\subset\{1,\ldots,N\}\), define
\[
 m_S(\mathbf R)
 :=
 \inf_{x\in Q}\max_{\alpha\in S}d(x,R_\alpha).
 \tag{34.4}
\]
The patches indexed by \(S\) have a common intersection exactly when
\[
 m_S(\mathbf R)<r
 \tag{34.5}
\]
for open balls.

\begin{lemma}[Lipschitz intersection margin]
\label{lem:v11v34-margin}
If
\[
 \max_\alpha d(R_\alpha,\widetilde R_\alpha)\leq\eta,
 \tag{34.6}
\]
then
\[
 |m_S(\mathbf R)-m_S(\widetilde{\mathbf R})|
 \leq\eta
 \tag{34.7}
\]
for every \(S\).
\end{lemma}

\begin{proof}
For each \(x\),
\[
 \left|
 \max_{\alpha\in S}d(x,R_\alpha)
 -\max_{\alpha\in S}d(x,\widetilde R_\alpha)
 \right|
 \leq\eta
\]
by the triangle inequality.  Taking the infimum gives one direction of
(34.7); exchanging the two reference families gives the other.
\end{proof}

\begin{theorem}[Eventual nerve stabilization]
\label{thm:v11v34-nerve-stability}
Suppose a limiting reference family \(\mathbf R_\infty\) satisfies the
strict finite margin
\[
 \min_{\varnothing\ne S\subset\{1,\ldots,N\}}
 |m_S(\mathbf R_\infty)-r|
 \geq\sigma>0.
 \tag{34.8}
\]
If
\[
 \max_\alpha d(R_{\alpha,i},R_{\alpha,\infty})<\frac\sigma2
 \tag{34.9}
\]
for all sufficiently large \(i\), then the entire Cech nerve at level
\(i\) equals the limiting nerve.  In particular, one fixed spanning
tree and one fixed set of independent cycle generators work at every
large regulator.
\end{theorem}

\begin{proof}
Lemma \ref{lem:v11v34-margin} moves every \(m_S\) by less than
\(\sigma/2\).  The sign of \(m_S-r\) therefore cannot change.
Equation (34.5) shows that every simplex is present or absent exactly as
in the limiting nerve.  The statements about the one-skeleton and cycle
generators follow.
\end{proof}

If the adjacent reference motions satisfy
\[
 \sum_i\max_\alpha
 d(R_{\alpha,i},\widehat R_{\alpha,i+1})<\infty,
 \tag{34.10}
\]
then the references are Cauchy and (34.9) follows whenever the strict
margin (34.8) holds.  Threshold intersections with zero margin must be
refined or tracked explicitly; they cannot be declared stable.

\subsection{Locality removes remote atlas dependence}

Let \(J_U(X)\) be a current paired with a tangent field \(X\) supported
in \(K\).  Assume its first variation has an exponentially local
bilinear kernel: for some \(C,\mu>0\),
\[
 |D J_U[Y](X)|
 \leq
 C\|X\|_\infty
 \sum_{e\in\operatorname{supp}X}
 \sum_{f}
 e^{-\mu d(e,f)}\|Y_f\|.
 \tag{34.11}
\]
Assume the lattice has polynomial volume growth.

\begin{proposition}[Cylinder buffer estimate]
\label{prop:v11v34-buffer}
If \(U,V\) agree on the radius-\(R\) neighbourhood of \(K\) and can be
joined by a path on which (34.11) holds, then for every
\(0<\mu'<\mu\),
\[
 |J_U(X)-J_V(X)|
 \leq
 C_{\mu,\mu'}|\operatorname{supp}X|
 e^{-\mu'R}\|X\|_\infty
 \sup_f\|U_f-V_f\|.
 \tag{34.12}
\]
\end{proposition}

\begin{proof}
Integrate \(DJ\) along the path from \(U\) to \(V\).  Its tangent
vanishes inside the \(R\)-neighbourhood of \(K\).  For
\(e\in\operatorname{supp}X\), the remaining sum has \(d(e,f)\geq R\).
Write
\[
 e^{-\mu d(e,f)}
 =
 e^{-\mu'R}e^{-(\mu-\mu')d(e,f)}
 e^{-\mu'(d(e,f)-R)}
\]
and discard the last factor.  Polynomial volume growth makes
\(\sum_f e^{-(\mu-\mu')d(e,f)}\) uniformly finite.  Sum over the support
of \(X\) and integrate the path.
\end{proof}

The same argument applies to the overlap scalar (33.5) if its
differentiated curvature kernel obeys the corresponding weighted bound.
Thus remote changes of a patch representative affect a fixed cylinder
phase only by an exponentially small buffer tail.  This does not erase
a global sector weight in the partition function.

\subsection{Disconnected sectors}

Let \(\Sigma_i\) be the finite or countable sector set at regulator
\(i\).  After a consistent identification of stable sector labels, let
\[
 Z_i^\sigma(\mathcal O)
 \tag{34.13}
\]
be the complex saturated cylinder coefficient in sector \(\sigma\),
including its prescribed topological phase.  The full coefficient is
\[
 Z_i(\mathcal O)=\sum_{\sigma\in\Sigma_i}
 Z_i^\sigma(\mathcal O).
 \tag{34.14}
\]
Absolute sector tightness means that for every \(\varepsilon>0\) there
is a finite set \(F\) of stable labels such that
\[
 \sup_i\sum_{\sigma\notin F}
 |Z_i^\sigma(\mathcal O)|<\varepsilon.
 \tag{34.15}
\]

\begin{theorem}[Sectorwise convergence plus tightness]
\label{thm:v11v34-sector-tightness}
Assume (34.15) and, for every fixed stable sector \(\sigma\),
\[
 Z_i^\sigma(\mathcal O)\longrightarrow
 Z_\infty^\sigma(\mathcal O).
 \tag{34.16}
\]
Then
\[
 \sum_\sigma|Z_\infty^\sigma(\mathcal O)|<\infty
 \tag{34.17}
\]
and
\[
 Z_i(\mathcal O)\longrightarrow
 \sum_\sigma Z_\infty^\sigma(\mathcal O).
 \tag{34.18}
\]
\end{theorem}

\begin{proof}
For a finite set \(F\), (34.16) permits passage to the limit in its
finite sum.  Fatou's lemma for counting measure and (34.15) give
\[
 \sum_{\sigma\notin F}|Z_\infty^\sigma|
 \leq
 \liminf_i\sum_{\sigma\notin F}|Z_i^\sigma|
 \leq\varepsilon.
\]
This proves absolute summability.  The difference in (34.18) is bounded
by the finite-\(F\) difference plus the two tails.  The first tends to
zero and the tails are at most \(2\varepsilon\).  Let
\(\varepsilon\downarrow0\).
\end{proof}

Sector phases must be fixed by the physical topological-angle convention
or another declared normalization.  A local determinant-line frame
inside one sector does not determine their relative values.

\subsection{Combined finite-cylinder certificate}

For a fixed cylinder, consider the following rows:
\begin{enumerate}[label=\textnormal{(\roman*)}]
\item a compact good set covered by finitely many V32 patches;
\item summable reference transitions and the strict nerve margin
(34.8);
\item vanishing cycle holonomies for the stabilized finite generator
set in Proposition \ref{prop:v11v33-cycle};
\item summable adjacent overlap-phase differences from (33.27);
\item an exponential buffer estimate of the form (34.12); and
\item sectorwise coefficient convergence with absolute tightness
(34.15).
\end{enumerate}

\begin{corollary}[Phase-coherent cylinder limit]
\label{cor:v11v34-cylinder-limit}
If rows (i)--(vi) hold and the determinant moduli and zero-mode
saturated minors converge in the norms required by V18, then the
fermionic saturated cylinder coefficients have a phase-coherent
sector-summed limit.
\end{corollary}

\begin{proof}
Rows (i)--(iv) give a stabilized finite phase nerve, trivialize its
cycle holonomies, and make the transition phases Cauchy.  Row (v) makes
the result independent of remote atlas choices in the cylinder limit.
The V18 hypothesis supplies the coefficient modulus through determinant
zeros.  Apply Theorem \ref{thm:v11v34-sector-tightness} to sum the
sectorwise limits.
\end{proof}

\begin{firewall}[Still missing]
\label{fw:v11v34-missing}
Compactness does not prove that the coupled measure concentrates on
\(\mathcal G_K\), that flat-reference patches cover every retained
configuration, or that the strict nerve margin and cycle-holonomy rows
hold.  No sector tightness estimate (34.15) has been obtained for the
coupled Standard Model--Einstein density.  The bosonic normalization,
stress-energy response, and Osterwalder--Schrader reconstruction remain
open.
\end{firewall}

\begin{assessment}[V34 advance]
\label{ass:v11v34}
The phase-topology problem is now finite on every fixed good cylinder.
A strict intersection margin stabilizes one finite Cech nerve and its
cycle generators, exponential locality removes remote atlas dependence,
and a separate absolute-tightness theorem handles disconnected
topological sectors.  The next mandatory V35 audit will check the
companion programmes before choosing between the good-set probability
row and the determinant-modulus comparison.
\end{assessment}
\section{V35 companion audit: response after stationary elimination}
\label{sec:v11v35}

This is the compulsory five-version checkpoint.  The active companion
records inspected are
\[
\begin{array}{c|c|r|l}
 \text{programme}&\text{version}&\text{bytes}&\text{SHA--256}\\ \hline
 \text{Yang--Mills}&32&458760&
 \texttt{C2794F960F65463F5CFECE3D25EA6039612A8E2CEBD8AA3B2B6BE5870269B0F0}\\
 \text{Navier--Stokes}&26&278283&
 \texttt{82C887F8C2C69E4F28FAD7C3DC8DE5B9099CB5A4BF431EBA1FFF3E91935978AD}
\end{array}
\tag{35.1}
\]

\begin{audit}[New YM result and unchanged NS result]
\label{audit:v11v35-companions}
YM V32 completes the translated-coordinate correction at zero internal
and external momentum.  Its independently generated cubic and quartic
Wilson words replay suffix formulas, and the corrected finite endpoint
card is nonzero.  It explicitly does not interpret one cell as a mass.
For a nonzero external momentum derivative, it keeps three distinct
rows: the direct quartic variation, the stationary-center acceleration,
and the Haar-density response.

NS V26 remains unchanged and supplies no new programme transfer beyond
the restricted-input warning recorded at V25 and V30.
\end{audit}

\subsection{The reduction-order lesson}

Let \(x\) denote retained variables and \(y\) eliminated variables.  A
stationary reduction uses
\[
 D_yS(x,y_*(x))=0.
 \tag{35.2}
\]
Even if \(y_*(0)=0\) and a measure density is stationary at \(x=0\),
their second response need not vanish.

\begin{lemma}[Value zero does not delete a second jet]
\label{lem:v11v35-second-jet}
There are smooth functions \(f\) with
\[
 f(0)=f'(0)=0,\qquad f''(0)\ne0.
 \tag{35.3}
\]
Therefore a row that vanishes at a symmetric base point cannot be
omitted from a second-response calculation without an independent
second-jet identity.
\end{lemma}

\begin{proof}
Take \(f(q)=q^2\).
\end{proof}

The point is elementary, but it blocks a recurrent invalid inference:
exact flatness or translation invariance at one background determines a
value, while the external response differentiates the family around
that background.

\subsection{Three finite response rows}

Suppose a reduced weight has schematic form
\[
 \mathcal W_{\rm red}(x)
 =
 e^{-S(x,y_*(x))}
 a(x,y_*(x))
 J(x,y_*(x)),
 \tag{35.4}
\]
where \(a\) is a determinant or fluctuation factor and \(J\) is a Haar,
coarea, or changing-density factor.  Then
\[
 D^2\log\mathcal W_{\rm red}
 =
 -D^2(S\circ y_*)
 +D^2\log(a\circ y_*)
 +D^2\log(J\circ y_*).
 \tag{35.5}
\]
Each composite derivative contains the jets of \(y_*\).

\begin{proposition}[No termwise symmetry deletion]
\label{prop:v11v35-no-deletion}
A symmetry that makes the assembled first derivative of
\(\log\mathcal W_{\rm red}\) vanish does not imply that any of the three
second derivatives in (35.5) vanishes separately.  Deleting one is
valid only after an exact identity for that row or after the complete
second derivatives have been assembled and cancelled.
\end{proposition}

\begin{proof}
The three terms in (35.5) are independent smooth scalar functions.
Their first derivatives may cancel, or each may be odd at the base
point, while their even quadratic coefficients remain arbitrary.
Lemma \ref{lem:v11v35-second-jet} gives an example for each row.
Linearity of differentiation justifies cancellation only after the
complete sum has been formed.
\end{proof}

\subsection{Transfer to the SM--Einstein route}

The determinant-curvature work in V22--V34 controls the fermionic phase
connection.  It is not the full logarithmic derivative of the coupled
reduced density.  A complete cylinder response must also retain:
\begin{enumerate}[label=\textnormal{(\roman*)}]
\item motion of every gauge, Higgs, lapse, shift, connection, or metric
variable eliminated by a stationary equation;
\item the determinant of the eliminated Hessian or its non-Gaussian
replacement;
\item Haar, Faddeev--Popov, coarea, and changing-Hilbert-density
Jacobians;
\item boundary and shell factors created by the local completion; and
\item the zero-mode saturated fermion minors from V18.
\end{enumerate}

\begin{decision}[V36 post-audit acquisition]
\label{dec:v11v35-next}
V36 will derive exact first and second jets of the stationary map by the
implicit function theorem, the Schur-complement Hessian of the reduced
action, and the first two logarithmic jets of Hessian determinants and
coarea densities.  These formulas will be kept as separate adjacent
comparison rows.
\end{decision}

\begin{firewall}[Audit boundary]
\label{fw:v11v35-boundary}
YM V32 supplies an order-of-operations correction, not an SM
stationary map or probability estimate.  Its corrected nonzero cell is
not imported as a physical constant.  NS V26 remains unrelated to the
coupled reduced-density derivatives.
\end{firewall}

\begin{assessment}[V35 audit]
\label{ass:v11v35}
The required audit is complete.  It redirects the next step from phase
topology to the reduced-density response: the continuum coefficient
cannot be certified from determinant curvature while stationary-center
and measure-density jets remain uncompiled.
\end{assessment}
\section{Stationary elimination and reduced-density response}
\label{sec:v11v36}

V35 requires the response of eliminated variables and measure factors to
be compiled before the coupled density is differentiated.  This note
does so in finite dimensions.  The formulas apply to a local gauge,
Higgs, connection, lapse, shift, or metric block whenever its stationary
Hessian is invertible.

\subsection{Implicit stationary map}

Let \(X,Y\) be finite-dimensional normed spaces and
\[
 S:\mathcal U\subset X\times Y\longrightarrow\mathbb R
 \tag{36.1}
\]
be \(C^4\).  Put
\[
 \mathcal F(x,y):=D_yS(x,y).
 \tag{36.2}
\]
Assume
\[
 \mathcal F(x,y_*(x))=0
 \tag{36.3}
\]
and that
\[
 A(x):=D_y\mathcal F(x,y_*(x))
      =D_y^2S(x,y_*(x))
 \tag{36.4}
\]
is invertible, with
\[
 \|A^{-1}\|\leq\kappa.
 \tag{36.5}
\]
The implicit function theorem gives a local \(C^2\) stationary map
\(y_*\).

\begin{theorem}[First and second stationary jets]
\label{thm:v11v36-stationary-jets}
For constant retained legs \(u,v\in X\),
\[
 y_u:=
 Dy_*[u]
 =
 -A^{-1}\mathcal F_x[u],
 \tag{36.6}
\]
and
\[
 \begin{split}
 y_{uv}:=D^2y_*[u,v]
 =-A^{-1}\bigl(
 &\mathcal F_{xx}[u,v]
 +\mathcal F_{xy}[u,y_v]\\
 &+\mathcal F_{yx}[y_u,v]
 +\mathcal F_{yy}[y_u,y_v]
 \bigr).
 \end{split}
 \tag{36.7}
\]
All derivatives of \(\mathcal F\) are evaluated at
\((x,y_*(x))\).
\end{theorem}

\begin{proof}
Differentiate (36.3) in the \(u\)-leg:
\[
 0=\mathcal F_x[u]+A y_u,
\]
which gives (36.6).  Differentiate this identity in the \(v\)-leg.
The derivative can hit the base \(x\), the stationary argument \(y_*\),
or \(y_u\).  Sorting the four second derivatives of \(\mathcal F\) and
the final term \(Ay_{uv}\) gives (36.7).
\end{proof}

If, for unit retained legs,
\[
 \|\mathcal F_x\|\leq B_1,\quad
 \|\mathcal F_{xx}\|\leq B_2,\quad
 \|\mathcal F_{xy}\|,\|\mathcal F_{yx}\|\leq B_{xy},\quad
 \|\mathcal F_{yy}\|\leq B_{yy},
 \tag{36.8}
\]
then
\[
 \|y_u\|\leq\kappa B_1,
 \tag{36.9}
\]
\[
 \|y_{uv}\|
 \leq
 \kappa\left(
 B_2+2B_{xy}\kappa B_1+B_{yy}\kappa^2B_1^2
 \right).
 \tag{36.10}
\]
Thus the stationary acceleration has a literal inverse-Hessian stock.

\subsection{Reduced action and the Schur complement}

Define
\[
 S_{\rm red}(x):=S(x,y_*(x)).
 \tag{36.11}
\]

\begin{theorem}[Reduced Hessian]
\label{thm:v11v36-schur}
The first derivative is the envelope derivative
\[
 D S_{\rm red}[u]=S_x[u],
 \tag{36.12}
\]
and the second derivative is
\[
 D^2S_{\rm red}[u,v]
 =
 S_{xx}[u,v]
 -
 S_{xy}[u]\,
 A^{-1}\,
 S_{yx}[v].
 \tag{36.13}
\]
Equation (36.13) is the Schur complement of the eliminated Hessian.
\end{theorem}

\begin{proof}
The chain rule gives
\[
 D S_{\rm red}[u]=S_x[u]+S_y[y_u].
\]
The second term vanishes by (36.3), proving (36.12).  Differentiate
(36.12) in the \(v\)-leg:
\[
 D^2S_{\rm red}[u,v]
 =
 S_{xx}[u,v]+S_{xy}[u,y_v].
\]
Insert (36.6), noting that
\(\mathcal F_x[v]=S_{yx}[v]\), to obtain (36.13).
\end{proof}

The acceleration \(y_{uv}\) cancels from the action Hessian because the
action is stationary in \(y\).  It does not cancel from a general
amplitude or density.

\subsection{A general amplitude after elimination}

Let \(\ell(x,y)\) be any \(C^2\) scalar, such as a logarithmic Haar
factor, a local counterterm, or the logarithm of a nonvanishing
saturated coefficient.  Put
\[
 \ell_{\rm red}(x)=\ell(x,y_*(x)).
 \tag{36.14}
\]

\begin{proposition}[Amplitude jets]
\label{prop:v11v36-amplitude}
One has
\[
 D\ell_{\rm red}[u]
 =
 \ell_x[u]+\ell_y[y_u],
 \tag{36.15}
\]
and
\[
 \begin{split}
 D^2\ell_{\rm red}[u,v]
 ={}&
 \ell_{xx}[u,v]
 +\ell_{xy}[u,y_v]
 +\ell_{yx}[y_u,v]\\
 &+\ell_{yy}[y_u,y_v]
 +\ell_y[y_{uv}].
 \end{split}
 \tag{36.16}
\]
\end{proposition}

\begin{proof}
Apply the first- and second-order chain rules to
\(\ell\circ(\operatorname{id},y_*)\).
\end{proof}

The last term in (36.16) is the stationary-acceleration row highlighted
by YM V32.  It vanishes only when \(\ell_y=0\) at the base point or
when an exact cancellation has been proved.

\subsection{Hessian determinant}

Suppose an exact Gaussian elimination or a declared saddle coefficient
contains
\[
 (\det A(x))^{-1/2},
 \tag{36.17}
\]
where \(A(x)\) is positive definite and already includes the total
dependence through \(y_*(x)\).  Write \(A_u=DA[u]\) and
\(A_{uv}=D^2A[u,v]\).

\begin{theorem}[Log-Hessian jets]
\label{thm:v11v36-logdet}
One has
\[
 D\log\det A[u]
 =
 \Tr(A^{-1}A_u),
 \tag{36.18}
\]
and
\[
 D^2\log\det A[u,v]
 =
 \Tr\left(
 A^{-1}A_{uv}
 -A^{-1}A_vA^{-1}A_u
 \right).
 \tag{36.19}
\]
The contribution of (36.17) is minus one half of these expressions.
\end{theorem}

\begin{proof}
Jacobi's formula gives (36.18).  Differentiate it in the \(v\)-leg and
use
\[
 D A^{-1}[v]=-A^{-1}A_vA^{-1}.
\]
This yields (36.19).
\end{proof}

The total derivative \(A_u\) includes
\[
 A_u=
 D_x(D_y^2S)[u]+D_y(D_y^2S)[y_u],
 \tag{36.20}
\]
and \(A_{uv}\) includes \(y_u,y_v,y_{uv}\).  Computing only the explicit
\(x\)-derivative of \(D_y^2S\) is incomplete.

If \(A\) acts on an \(n\)-dimensional eliminated block and
\[
 \|A^{-1}\|\leq\kappa,\qquad
 \|A_u\|\leq a_1,\qquad \|A_{uv}\|\leq a_2,
 \tag{36.21}
\]
then
\[
 |D\log\det A[u]|\leq n\kappa a_1,
 \tag{36.22}
\]
\[
 |D^2\log\det A[u,v]|
 \leq n(\kappa a_2+\kappa^2a_1^2).
 \tag{36.23}
\]

\subsection{Coarea and constraint densities}

Let a regular constraint map have normal Gram operator
\[
 G(x)=D C(x)D C(x)^*>0.
 \tag{36.24}
\]
The coarea factor is
\[
 J_{\rm co}(x)=(\det G(x))^{-1/2}.
 \tag{36.25}
\]

\begin{corollary}[Coarea jets]
\label{cor:v11v36-coarea}
With total derivatives through every stationary variable,
\[
 D\log J_{\rm co}[u]
 =
 -\frac12\Tr(G^{-1}G_u),
 \tag{36.26}
\]
\[
 D^2\log J_{\rm co}[u,v]
 =
 -\frac12\Tr\left(
 G^{-1}G_{uv}-G^{-1}G_vG^{-1}G_u
 \right).
 \tag{36.27}
\]
\end{corollary}

\begin{proof}
Apply Theorem \ref{thm:v11v36-logdet} to \(G\) and multiply by
\(-1/2\).
\end{proof}

The same formula applies to a Faddeev--Popov determinant on a patch
where it is nonzero, with its sign and power taken from the actual
measure convention.

\subsection{The complete local logarithmic ledger}

On a patch where every scalar factor is nonzero, a reduced local weight
may be written
\[
 \log\mathcal W_{\rm red}
 =
 -S_{\rm red}
 +\ell_{\rm ferm}
 +\ell_{\rm Haar}
 +\ell_{\rm dens}
 -\frac12\log\det A
 -\frac12\log\det G
 +\ell_{\rm shell}.
 \tag{36.28}
\]
Here \(\ell_{\rm ferm}\) includes the chosen determinant-line phase and
the modulus of the saturated fermion coefficient.

\begin{proposition}[Response completeness]
\label{prop:v11v36-ledger}
The first and second response of (36.28) is the sum of:
\begin{enumerate}[label=\textnormal{(\roman*)}]
\item the envelope and Schur-complement action terms
(36.12)--(36.13);
\item the composite amplitude terms (36.15)--(36.16) for every
\(\ell\)-row;
\item the Hessian-determinant terms (36.18)--(36.19); and
\item the coarea terms (36.26)--(36.27).
\end{enumerate}
No stationary acceleration appears in the action Hessian, but it appears
in every nonstationary amplitude through (36.16) and in the total
derivatives of \(A\) and \(G\).
\end{proposition}

\begin{proof}
Differentiate the exact sum (36.28) twice and apply the cited formulas
to each term.  There are no other product terms because the logarithm
has converted the product of nonzero scalar factors into a sum.
\end{proof}

At a fermion zero, \(\ell_{\rm ferm}\) is not defined.  One must instead
use the complementary minor or exterior-power saturated coefficient of
V18 and differentiate that coefficient directly.  The logarithmic
ledger is valid only on its nonvanishing locus.

\subsection{Adjacent inverse and determinant rows}

Let \(A,\widetilde A\) be adjacent positive Hessians with
\[
 \|A^{-1}\|,\|\widetilde A^{-1}\|\leq\kappa.
 \tag{36.29}
\]
Then
\[
 \|A^{-1}-\widetilde A^{-1}\|
 \leq\kappa^2\|A-\widetilde A\|.
 \tag{36.30}
\]

\begin{proof}
Use
\[
 A^{-1}-\widetilde A^{-1}
 =
 A^{-1}(\widetilde A-A)\widetilde A^{-1}.
\]
\end{proof}

Suppose
\[
 \|A-\widetilde A\|\leq\delta_0,\quad
 \|A_u-\widetilde A_u\|\leq\delta_1,\quad
 \|A_{uv}-\widetilde A_{uv}\|\leq\delta_2,
 \tag{36.31}
\]
and both first and second jets are bounded by \(a_1,a_2\).

\begin{theorem}[Adjacent log-determinant comparison]
\label{thm:v11v36-adjacent-logdet}
The first jets differ by at most
\[
 n(\kappa^2a_1\delta_0+\kappa\delta_1),
 \tag{36.32}
\]
and the second jets differ by at most
\[
 n\left[
 (\kappa^2a_2+2\kappa^3a_1^2)\delta_0
 +2\kappa^2a_1\delta_1
 +\kappa\delta_2
 \right].
 \tag{36.33}
\]
The same bounds hold for \(\log\det G\) with its own dimension and
inverse stocks.
\end{theorem}

\begin{proof}
Subtract (36.18) for the two Hessians.  Change the inverse and the first
jet one at a time; (36.30) gives (36.32).

For (36.19), the term \(A^{-1}A_{uv}\) contributes
\[
 n(\kappa^2a_2\delta_0+\kappa\delta_2).
\]
In the four-factor term \(A^{-1}A_vA^{-1}A_u\), changing either inverse
contributes \(\kappa^3a_1^2\delta_0\), and changing either first jet
contributes \(\kappa^2a_1\delta_1\).  Sum the four changes to obtain
(36.33).
\end{proof}

Thus a summable three-row Hessian card
\[
 \sum_i(\delta_{0,i}+\delta_{1,i}+\delta_{2,i})<\infty
 \tag{36.34}
\]
makes both determinant-response jets projectively Cauchy under uniform
inverse and derivative stocks.

\begin{firewall}[Reduction is not yet constructed]
\label{fw:v11v36-reduction}
No coupled Standard Model--Einstein stationary block has yet been
specified with a uniformly invertible Hessian (36.5), and an interacting
elimination need not be exactly Gaussian.  Gribov degeneracies may make
the coarea Gram operator singular.  The nonvanishing logarithmic ledger
does not cross fermion determinant zeros without the V18 saturated
minor construction.  Probability estimates for the inverse-Hessian,
ellipticity, and gap good sets remain open.
\end{firewall}

\begin{assessment}[V36 advance]
\label{ass:v11v36}
The post-audit reduced-response calculation is complete at the
deterministic finite-dimensional level.  Stationary velocity and
acceleration, the Schur-complement action Hessian, amplitude response,
Hessian determinant, and coarea jets are exact, and (36.32)--(36.34)
give a summable adjacent determinant row.  The next finite task is to
derive a probability transfer theorem showing how exponential moments
of local badness variables imply vanishing local failure intensity for
the simultaneous gap, ellipticity, and inverse-Hessian good set.
\end{assessment}
\section{Local badness moments and good-set probability transfer}
\label{sec:v11v37}

The deterministic estimates in V23--V36 hold on a simultaneous good
set: gauge completion is gapped, the metric is elliptic, stationary
Hessians and coarea Gram operators are invertible, and local derivative
stocks are bounded.  This note proves a probability transfer theorem
from exponential moments of local badness variables to the complement
bounds required for cylinder convergence.

\subsection{Local badness vector}

For a block \(B\), introduce nonnegative measurable variables
\[
 Z^{\rm pl}_{i,B},\quad
 Z^{\rm met}_{i,B},\quad
 Z^{\rm Hess}_{i,B},\quad
 Z^{\rm co}_{i,B},\quad
 Z^{\rm der}_{i,B}.
 \tag{37.1}
\]
They may be chosen, for example, to measure the normalized maximum
plaquette defect, the larger of
\(\|g\|\) and \(\|g^{-1}\|\), the inverse stationary-Hessian norm, the
inverse coarea-Gram norm, and the required coefficient-derivative
stock.  A singular inverse is assigned value \(+\infty\).

Given thresholds \(t_{j,i}\), define
\[
 \mathsf G_{i,B}
 =
 \bigcap_j\{Z^j_{i,B}\leq t_{j,i}\}.
 \tag{37.2}
\]
The thresholds must be chosen so that \(\mathsf G_{i,B}\) implies the
deterministic hypotheses already recorded; for the plaquette row this
includes the V14 inequality
\(A_*\epsilon_i\leq\gamma_m/16\).

\subsection{Exponential moment to local intensity}

\begin{lemma}[Exponential Markov bound]
\label{lem:v11v37-markov}
If \(Z\geq0\) and
\[
 \mathbb E e^{\theta Z}\leq M
 \tag{37.3}
\]
for \(\theta>0\), then
\[
 \mathbb P\{Z>t\}\leq M e^{-\theta t}.
 \tag{37.4}
\]
\end{lemma}

\begin{proof}
On \(\{Z>t\}\), \(e^{\theta Z}\geq e^{\theta t}\).  Apply Markov's
inequality to \(e^{\theta Z}\).
\end{proof}

Assume, uniformly over block anchors,
\[
 \mathbb E_{\nu_i}
 \exp(\theta_{j,i}Z^j_{i,B})
 \leq M_{j,i}.
 \tag{37.5}
\]
Define the local failure intensity
\[
 p_i:=
 \sup_B\nu_i(\mathsf G_{i,B}^{\,c}).
 \tag{37.6}
\]

\begin{theorem}[Simultaneous local good-set bound]
\label{thm:v11v37-local-intensity}
Equations (37.2) and (37.5) imply
\[
 p_i
 \leq
 \sum_j M_{j,i}e^{-\theta_{j,i}t_{j,i}}.
 \tag{37.7}
\]
No factor proportional to the total lattice volume occurs.
\end{theorem}

\begin{proof}
The complement of (37.2) is the union of the five threshold-failure
events.  Apply Lemma \ref{lem:v11v37-markov} to each and then the union
bound.  Taking the supremum over the anchor preserves the estimate.
\end{proof}

A fixed cylinder meeting \(N_K\) blocks therefore obeys
\[
 \nu_i\{\text{some block meeting }K\text{ is bad}\}
 \leq N_Kp_i.
 \tag{37.8}
\]
If a buffer grows and meets \(N_i\) blocks, the correct requirement is
\(N_ip_i\to0\), or the corresponding summable condition.  A
fixed-block estimate cannot be multiplied by the entire cofinal volume.

\subsection{Bad-complement expectation}

Let \(q>1\) and let \(q'=q/(q-1)\).

\begin{lemma}[Local complement by Hölder]
\label{lem:v11v37-holder}
For a random coefficient \(C_i\) and an event \(E_i\),
\[
 \mathbb E_{\nu_i}
 \left[|C_i|\mathbf1_{E_i}\right]
 \leq
 \|C_i\|_{L^q(\nu_i)}
 \nu_i(E_i)^{1/q'}.
 \tag{37.9}
\]
\end{lemma}

\begin{proof}
Apply Hölder's inequality to
\(|C_i|\) and \(\mathbf1_{E_i}\), noting that
\(\|\mathbf1_{E_i}\|_{q'}=\nu_i(E_i)^{1/q'}\).
\end{proof}

Combining (37.7)--(37.9), a uniform local \(L^q\) stock
\[
 \|C_i\|_{L^q(\nu_i)}\leq C_q
 \tag{37.10}
\]
gives
\[
 \mathbb E\left[
 |C_i|\mathbf1_{\mathsf G_{i,K}^{\,c}}
 \right]
 \leq
 C_q(N_Kp_i)^{1-1/q}.
 \tag{37.11}
\]
This is the probability row needed in the V7--V9 complement argument.

\subsection{Combination with adjacent deterministic estimates}

To compare two regulator measures, assume a compatible adjacent coupling
\(\pi_i\) of their transported cylinder fields.  Let \(C_i,C_{i+1}\)
be the two coefficients on that coupling and let
\(\mathsf G_i^{\rm adj}\) be the event on which all deterministic
adjacent hypotheses hold at both levels.  Suppose
\[
 |C_{i+1}-C_i|
 \leq\beta_i
 \quad\hbox{on }\mathsf G_i^{\rm adj},
 \tag{37.12}
\]
\[
 \|C_{i+1}-C_i\|_{L^q(\pi_i)}\leq D_q,
 \tag{37.13}
\]
and
\[
 \pi_i((\mathsf G_i^{\rm adj})^c)\leq\widehat p_i.
 \tag{37.14}
\]

\begin{theorem}[Good-set plus complement comparison]
\label{thm:v11v37-adjacent}
Under (37.12)--(37.14),
\[
 \mathbb E_{\pi_i}|C_{i+1}-C_i|
 \leq
 \beta_i+D_q\widehat p_i^{\,1-1/q}.
 \tag{37.15}
\]
If
\[
 \sum_i\beta_i<\infty,\qquad
 \sum_i\widehat p_i^{\,1-1/q}<\infty,
 \tag{37.16}
\]
then the transported coefficients are Cauchy in \(L^1\) on any common
projective coupling, and their expectations are Cauchy.
\end{theorem}

\begin{proof}
Split the expectation into the good event and its complement.
Equation (37.12) bounds the first part by \(\beta_i\).
Lemma \ref{lem:v11v37-holder}, applied under \(\pi_i\), bounds the second
by the second term of (37.15).  Summability makes the telescoped
\(L^1\) differences tend to zero.
\end{proof}

For the determinant current, \(\beta_i\) may be the V23 card.  For
metric coefficients it may be (29.23); for patch phases, (33.27); and
for stationary determinant factors, (36.33).  The theorem adds them
only after all required good events have been intersected.

\subsection{Almost-sure use requires a common coupling}

If all regulator fields are realized on one probability space and
\[
 \sum_i\mathbb P((\mathsf G_i^{\rm adj})^c)<\infty,
 \tag{37.17}
\]
then the first Borel--Cantelli lemma gives
\[
 \mathbb P\{
 (\mathsf G_i^{\rm adj})^c\text{ infinitely often}
 \}=0.
 \tag{37.18}
\]

\begin{proof}
The probability of the limsup is bounded by the tail of the convergent
series in (37.17), hence is zero.
\end{proof}

Separate measures \(\nu_i\) with no common coupling do not support the
almost-sure statement (37.18).  Their marginal tail bounds still give
the expectation estimate (37.11).

\subsection{A scale-choice corollary}

Suppose for constants \(M_j,\theta_j>0\),
\[
 M_{j,i}\leq M_j,\qquad \theta_{j,i}\geq\theta_j,
 \tag{37.19}
\]
and choose
\[
 t_{j,i}=a_j\log(i+2).
 \tag{37.20}
\]
Then
\[
 p_i\leq
 \sum_jM_j(i+2)^{-\theta_ja_j}.
 \tag{37.21}
\]

\begin{corollary}[Summable complement design]
\label{cor:v11v37-scale-choice}
If
\[
 \min_j\theta_ja_j
 >
 \frac{q}{q-1},
 \tag{37.22}
\]
then
\[
 \sum_i p_i^{\,1-1/q}<\infty.
 \tag{37.23}
\]
For a buffer with \(N_i\leq C(i+2)^d\), it suffices to replace
(37.22) by
\[
 \min_j\theta_ja_j
 >
 d+\frac{q}{q-1}
 \tag{37.24}
\]
to make
\[
 \sum_i(N_ip_i)^{\,1-1/q}<\infty.
 \tag{37.25}
\]
\end{corollary}

\begin{proof}
Equation (37.21) is a finite sum of powers.  Raising it to
\(1-1/q\) preserves a power with exponent
\((1-1/q)\min_j\theta_ja_j\), up to a fixed constant.
Condition (37.22) makes that exponent greater than one.
With \(N_i\leq C(i+2)^d\), the exponent becomes
\((1-1/q)(\min_j\theta_ja_j-d)\); (37.24) makes it greater than one.
\end{proof}

The thresholds in (37.20) grow.  If a deterministic theorem requires a
fixed small upper threshold, as the plaquette defect does, then the
badness variable must be normalized by the scale-dependent physical
threshold before applying this corollary.  One may not enlarge the
allowed plaquette defect past the V14 gap condition.

\subsection{Acquisition ledger}

For a fixed cylinder, a sufficient probabilistic input is now:
\begin{enumerate}[label=\textnormal{(\roman*)}]
\item exponential moments (37.5) for every normalized badness row;
\item threshold choices compatible with the deterministic gap,
ellipticity, inverse-Hessian, and coarea conditions;
\item a local \(L^q\) stock for the complete saturated coefficient;
\item a compatible adjacent coupling when \(L^1\) projective
convergence, rather than only marginal expectation control, is sought;
and
\item summability as in (37.16).
\end{enumerate}

\begin{firewall}[Moments are not derived]
\label{fw:v11v37-moments}
No exponential moment (37.5) has been proved for the coupled renewal
Standard Model--Einstein measure.  In particular, Euclidean Einstein
weights need not be positive or coercive without a specified contour,
gauge fixing, and conformal-mode treatment.  The theorem converts
moments into the required local failure bounds; it does not manufacture
the moments, the compatible coupling, or reflection positivity.
\end{firewall}

\begin{assessment}[V37 advance]
\label{ass:v11v37}
The probability interface is now explicit.  Five local badness moments
give a simultaneous failure intensity with no total-volume factor;
Hölder controls the discarded complement; and a compatible coupling
combines that complement with every deterministic adjacent card already
proved.  The next finite task is to identify a positive finite-cutoff
bosonic reference measure and state a coercive local action inequality
strong enough to produce the moment hypotheses without assuming the
Einstein continuum.
\end{assessment}
\section{Conditional coercivity and positive-reference moment transfer}
\label{sec:v11v38}

V37 assumes exponential moments.  At finite cutoff they can be derived
from a local conditional action inequality, but only under a positive
measure.  The chiral Standard Model weight may be complex, so this note
first constructs the estimate under a positive bosonic or
phase-quenched reference measure and then gives the exact reweighting
condition needed to transfer it.

\subsection{A local conditional moment theorem}

Let \(z_B\) denote all variables in a finite block and \(z_{B^c}\) the
exterior variables.  Let \(\mu_B\) be a probability reference measure
on the block variables.  Suppose a positive finite-cutoff measure has
conditional density
\[
 d\nu(z_B\mid z_{B^c})
 =
 \frac{e^{-S_B(z_B\mid z_{B^c})}}
 {Z_B(z_{B^c})}\,d\mu_B(z_B).
 \tag{38.1}
\]
Let \(Z^1,\ldots,Z^J\geq0\) be local badness variables.

Assume, uniformly in the exterior configuration,
\[
 S_B(z_B\mid z_{B^c})
 \geq
 \sum_{j=1}^Ja_jZ^j(z_B)-b,
 \tag{38.2}
\]
with \(a_j>0\).  Also assume there is a comparison set
\(\mathcal A_B(z_{B^c})\) such that
\[
 \mu_B(\mathcal A_B)\geq v>0,
 \qquad
 S_B\leq c
 \quad\hbox{on }\mathcal A_B.
 \tag{38.3}
\]

\begin{theorem}[Uniform conditional exponential moments]
\label{thm:v11v38-conditional-moment}
For every \(0\leq\theta_j\leq a_j\),
\[
 \mathbb E_\nu
 \left[
 \exp\left(\sum_j\theta_jZ^j\right)
 \,\middle|\,z_{B^c}
 \right]
 \leq
 v^{-1}e^{b+c}.
 \tag{38.4}
\]
The same bound holds after averaging over \(z_{B^c}\).
\end{theorem}

\begin{proof}
Equation (38.3) gives the uniform denominator bound
\[
 Z_B(z_{B^c})
 \geq
 \int_{\mathcal A_B}e^{-S_B}\,d\mu_B
 \geq ve^{-c}.
 \tag{38.5}
\]
By (38.2),
\[
 \exp\left(\sum_j\theta_jZ^j-S_B\right)
 \leq
 e^b\exp\left(-\sum_j(a_j-\theta_j)Z^j\right)
 \leq e^b.
 \tag{38.6}
\]
Integrating (38.6) against the probability measure \(\mu_B\) and
dividing by (38.5) proves (38.4).  Averaging preserves the bound.
\end{proof}

This theorem permits arbitrary interactions across the block boundary;
their effect must be absorbed in uniform constants \(b,c,v\).  A formal
lower bound on the action without the conditional normalization bound
(38.3) is insufficient.

\subsection{A positive regulated reference}

At a finite regulator, a candidate positive reference can be assembled
from:
\begin{enumerate}[label=\textnormal{(\roman*)}]
\item product Haar measure for compact gauge links;
\item a positive finite-dimensional scalar or Higgs reference measure;
\item metric or tetrad variables restricted to the oriented ellipticity
chamber (29.2), with a positive local reference density; and
\item a real stabilizing bosonic action for which (38.2)--(38.3) can be
checked.
\end{enumerate}
Compact cutoffs make all continuous badness variables bounded, hence
give finite exponential moments.  The cofinal problem is to make their
constants uniform while the cutoffs are removed or refined.

For an \(SU(N)\) plaquette \(U_p\), put
\[
 z_p=1-\frac1N\operatorname{ReTr}U_p.
 \tag{38.7}
\]

\begin{lemma}[Wilson energy controls the chord defect]
\label{lem:v11v38-wilson-chord}
For every unitary \(U_p\),
\[
 \|I-U_p\|^2\leq2N z_p.
 \tag{38.8}
\]
Consequently,
\[
 z_p\leq\frac{\epsilon^2}{2N}
 \quad\Longrightarrow\quad
 \|I-U_p\|\leq\epsilon.
 \tag{38.9}
\]
\end{lemma}

\begin{proof}
If \(e^{i\theta_k}\) are the eigenvalues, then
\[
 \|I-U_p\|^2
 =
 2\max_k(1-\cos\theta_k)
 \leq
 2\sum_{k=1}^N(1-\cos\theta_k)
 =2Nz_p.
\]
\end{proof}

Thus a Wilson action contribution \(\beta_i z_p\), combined with a
uniform conditional comparison set, can feed the V14 plaquette
threshold.  The estimate must be applied locally; correlations do not
permit replacing the conditional theorem by an independent-plaquette
calculation.

\subsection{Metric and inverse-Hessian coercivity rows}

A sufficient metric badness variable is
\[
 Z^{\rm met}
 =
 \|g\|+\|g^{-1}\|,
 \tag{38.10}
\]
with value \(+\infty\) at degeneracy.  A bound
\(Z^{\rm met}\leq t\) implies
\[
 t^{-1}I\leq g\leq tI.
 \tag{38.11}
\]
Similarly, for the stationary and coarea blocks,
\[
 Z^{\rm Hess}=\|A^{-1}\|,
 \qquad
 Z^{\rm co}=\|G^{-1}\|.
 \tag{38.12}
\]
Thresholds on (38.12) give the inverse stocks in V36.

\begin{proposition}[What local coercivity would suffice]
\label{prop:v11v38-sufficient-coercivity}
If a positive conditional action obeys
\[
 S_B
 \geq
 a_{\rm pl}\sum_{p\subset B}z_p
 +a_{\rm met}Z^{\rm met}
 +a_{\rm Hess}Z^{\rm Hess}
 +a_{\rm co}Z^{\rm co}
 +a_{\rm der}Z^{\rm der}
 -b
 \tag{38.13}
\]
and has a comparison set satisfying (38.3), then every badness moment
in V37 is bounded uniformly for any exponent not exceeding its
coefficient in (38.13).
\end{proposition}

\begin{proof}
Apply Theorem \ref{thm:v11v38-conditional-moment} with the five
displayed variables, setting unused exponents to zero.
\end{proof}

Equation (38.13) is a sufficient criterion, not a claim about the
Einstein action.  Inverse-Hessian nondegeneracy usually requires a
separate small-ball or Morse estimate rather than direct action
coercivity; the proposition records exactly how either estimate must
enter the moment ledger.

\subsection{Positive reweighting transfer}

Let \(\mu\) be a positive probability measure for which
\[
 \mathbb E_\mu e^{q\theta Z}\leq M
 \tag{38.14}
\]
and let \(r\geq0\) be a reweighting density.  Define
\[
 d\nu=\frac{r\,d\mu}{\mathbb E_\mu r}.
 \tag{38.15}
\]
Let \(p,q>1\) be conjugate exponents.

\begin{theorem}[Moment transfer under matter reweighting]
\label{thm:v11v38-reweight}
If
\[
 \|r\|_{L^p(\mu)}\leq R,\qquad
 \mathbb E_\mu r\geq z_0>0,
 \tag{38.16}
\]
then
\[
 \mathbb E_\nu e^{\theta Z}
 \leq
 \frac R{z_0}M^{1/q}.
 \tag{38.17}
\]
\end{theorem}

\begin{proof}
By Hölder,
\[
 \mathbb E_\mu(re^{\theta Z})
 \leq
 \|r\|_{L^p(\mu)}
 \left(\mathbb E_\mu e^{q\theta Z}\right)^{1/q}
 \leq RM^{1/q}.
\]
Divide by the denominator in (38.16).
\end{proof}

This permits a positive Higgs or vectorlike determinant modulus to be
added when its \(L^p\) norm and normalization are controlled.  The
constants must remain local and uniform in the regulator.

\subsection{Complex chiral weights}

Let a finite-cutoff chiral weight be
\[
 w=|w|e^{i\vartheta}.
 \tag{38.18}
\]
Its normalized absolute measure is
\[
 d\overline\nu
 =
 \frac{|w|\,d\mu}{\int|w|\,d\mu}.
 \tag{38.19}
\]
All probability statements in V37 may be made under
\(\overline\nu\), provided Theorem \ref{thm:v11v38-reweight} controls
the modulus.  The phase \(e^{i\vartheta}\), including the V22--V34
determinant-line construction, remains inside the cylinder
coefficient.

The physical normalized expectation, when defined, is
\[
 \langle\mathcal O\rangle_w
 =
 \frac{
 \mathbb E_{\overline\nu}
 [e^{i\vartheta}\mathcal O]
 }{
 \mathbb E_{\overline\nu}
 [e^{i\vartheta}]
 }.
 \tag{38.20}
\]

\begin{proposition}[Phase-denominator requirement]
\label{prop:v11v38-phase-denominator}
If
\[
 \left|
 \mathbb E_{\overline\nu}e^{i\vartheta}
 \right|\geq d_0>0
 \tag{38.21}
\]
uniformly and the numerator coefficients converge, then the normalized
expectations (38.20) converge whenever the phase denominators converge.
Without (38.21), absolute-measure probability estimates do not control
the normalized chiral expectation.
\end{proposition}

\begin{proof}
The quotient map is continuous on
\(\{(a,b):|b|\geq d_0\}\).  If the denominator approaches zero, bounded
numerators can have unbounded or undefined quotients, so no conclusion
follows.
\end{proof}

Anomaly cancellation from V24 removes the local gauge obstruction to
constructing \(\vartheta\); it does not prove the lower bound (38.21).

\subsection{From conditional moments to the V37 card}

Assume the following finite-cutoff acquisitions:
\begin{enumerate}[label=\textnormal{(\roman*)}]
\item a positive local reference and conditional densities (38.1);
\item uniform coercivity and comparison-set constants in
(38.2)--(38.3);
\item an \(L^p\) modulus reweighting bound (38.16);
\item thresholds compatible with V12, V14, V29, and V36; and
\item, for normalized chiral expectations, the phase-denominator row
(38.21).
\end{enumerate}

\begin{corollary}[Conditional production of local failure intensities]
\label{cor:v11v38-production}
Rows (i)--(iv) imply the exponential moments (37.5) under the positive
dominating measure, with explicit constants from (38.4) and (38.17).
Theorem \ref{thm:v11v37-local-intensity} then gives the simultaneous
local failure bound (37.7).  Row (v) permits passage from unnormalized
or absolute coefficients to normalized chiral expectations.
\end{corollary}

\begin{proof}
Apply Theorem \ref{thm:v11v38-conditional-moment} at the larger exponent
required by (38.14), then Theorem \ref{thm:v11v38-reweight}.  Insert the
resulting moment constants in V37.  Proposition
\ref{prop:v11v38-phase-denominator} handles the final quotient.
\end{proof}

\begin{firewall}[Einstein conformal and sign rows]
\label{fw:v11v38-einstein}
No positive, cutoff-uniform renewal Einstein action satisfying
(38.13) has been constructed.  The Euclidean conformal mode, orientation
changes, gauge fixing, and possible Gribov singularities obstruct a
naive positive measure.  Compact ellipticity cutoffs give finite
moments but do not justify their removal.  No modulus reweighting bound
or uniform chiral phase denominator has been proved.
\end{firewall}

\begin{assessment}[V38 advance]
\label{ass:v11v38}
The moment assumption in V37 now has a precise upstream source.  A
uniform conditional coercivity inequality plus a positive comparison
set produces all local exponential moments; matter moduli transfer by
one Hölder estimate; and the complex chiral phase is isolated with an
explicit denominator condition.  The next finite task is to construct
a conformal-mode-stabilized finite Einstein reference on a compact
ellipticity chamber and determine which constants can remain uniform
as the chamber expands.
\end{assessment}
\section{A positive conformal-mode-stabilized metric reference}
\label{sec:v11v39}

V38 asks for a positive finite-cutoff metric reference with uniform
ellipticity moments.  A simple ultralocal construction is available on
positive definite metric matrices.  It is a regulator reference, not a
claim that the Euclidean Einstein action is positive.  Its value is that
the compact ellipticity chamber can be expanded at fixed stabilization
strength with explicit uniform moments.

\subsection{One-site reference measure}

Let
\[
 \mathcal P_n=\{g\in{\rm Sym}_n:g>0\}
 \tag{39.1}
\]
with Lebesgue measure \(dg\) on the
\(n(n+1)/2\) independent entries.  Define
\[
 Z(g):=\operatorname{tr}g+\operatorname{tr}g^{-1},
 \tag{39.2}
\]
and, for \(\kappa>0\),
\[
 I_n(\kappa):=
 \int_{\mathcal P_n}e^{-\kappa Z(g)}\,dg.
 \tag{39.3}
\]

\begin{theorem}[Finite positive metric reference]
\label{thm:v11v39-reference}
For every \(\kappa>0\),
\[
 0<I_n(\kappa)<\infty.
 \tag{39.4}
\]
Hence
\[
 d\mu_\kappa(g)
 =
 I_n(\kappa)^{-1}e^{-\kappa Z(g)}\,dg
 \tag{39.5}
\]
is a positive probability measure on all of \(\mathcal P_n\).
\end{theorem}

\begin{proof}
Positivity is immediate because every nonempty open subset of
\(\mathcal P_n\) has positive Lebesgue measure.  For finiteness, split
the cone into
\(\{\|g\|_F\leq1\}\) and its complement.  The first set is bounded in
the finite-dimensional vector space \({\rm Sym}_n\), and the integrand
is at most one, so its integral is finite.  On the complement, positivity
of the eigenvalues gives
\[
 \operatorname{tr}g\geq\|g\|_F.
\]
Therefore
\[
 e^{-\kappa Z(g)}
 \leq e^{-\kappa\|g\|_F},
\]
whose integral over the ambient finite-dimensional vector space is
finite.  Restricting it to \(\mathcal P_n\) preserves finiteness.
\end{proof}

The inverse term is not needed for integrability at the boundary of the
cone, which has finite volume in bounded sets.  It is needed to give
uniform exponential control of metric degeneracy.

\subsection{Exact exponential moments}

\begin{theorem}[Ellipticity moment]
\label{thm:v11v39-moment}
For \(0\leq\theta<\kappa\),
\[
 \mathbb E_{\mu_\kappa}e^{\theta Z(g)}
 =
 \frac{I_n(\kappa-\theta)}{I_n(\kappa)}
 <\infty.
 \tag{39.6}
\]
Moreover,
\[
 \|g\|+\|g^{-1}\|\leq Z(g),
 \tag{39.7}
\]
so (39.6) controls the metric badness variable (38.10).
\end{theorem}

\begin{proof}
Multiplying the density in (39.5) by \(e^{\theta Z}\) changes
\(\kappa\) to \(\kappa-\theta\), giving (39.6) and its finiteness by
Theorem \ref{thm:v11v39-reference}.  For a positive matrix, the largest
eigenvalue is at most the sum of the eigenvalues, and the same is true
for the inverse.  Adding the two inequalities gives (39.7).
\end{proof}

Every polynomial moment of \(g\) and \(g^{-1}\) follows.  Indeed, for
\(p\geq0\) and \(\delta>0\),
\[
 x^p\leq
 \left(\frac{p}{e\delta}\right)^p e^{\delta x},
 \qquad x\geq0,
 \tag{39.8}
\]
with the usual value one at \(p=0\).  Thus coefficient maps with
polynomial growth in \(g,g^{-1}\) have finite moments under
\(\mu_\kappa\), provided a small part of the exponential budget is
reserved.

\subsection{Conformal direction}

In four dimensions write
\[
 g=e^{2\sigma}\bar g,\qquad\det\bar g=1.
 \tag{39.9}
\]
The arithmetic--geometric mean inequality gives
\[
 \operatorname{tr}\bar g\geq4,\qquad
 \operatorname{tr}\bar g^{-1}\geq4.
 \tag{39.10}
\]
Consequently,
\[
 Z(g)
 =
 e^{2\sigma}\operatorname{tr}\bar g
 +e^{-2\sigma}\operatorname{tr}\bar g^{-1}
 \geq
 8\cosh(2\sigma).
 \tag{39.11}
\]

\begin{proposition}[Two-sided conformal stabilization]
\label{prop:v11v39-conformal}
The reference action \(\kappa Z(g)\) tends to \(+\infty\) as
\(\sigma\to+\infty\) or \(\sigma\to-\infty\), uniformly in the
unit-determinant shape \(\bar g\).
\end{proposition}

\begin{proof}
Equation (39.11) and
\(\cosh(2\sigma)\to\infty\) in both directions prove the assertion.
\end{proof}

This repairs the sign of the reference measure only by adding a
non-Einstein stabilizer.  It does not solve the conformal-factor problem
for the target Einstein weight.

\subsection{Expansion of the ellipticity chamber}

Let
\[
 \mathcal K_R
 =
 \{g:R^{-1}I\leq g\leq RI\},
 \qquad R>1,
 \tag{39.12}
\]
and let \(\mu_{\kappa,R}\) be (39.5) conditioned on
\(\mathcal K_R\).  Put
\[
 I_{n,R}(\kappa)
 =
 \int_{\mathcal K_R}e^{-\kappa Z(g)}\,dg.
 \tag{39.13}
\]

\begin{theorem}[Uniform chamber-removal moments]
\label{thm:v11v39-chamber}
Fix \(\kappa>0\), \(0\leq\theta<\kappa\), and \(R_0>1\).  Then, for
all \(R\geq R_0\),
\[
 \mathbb E_{\mu_{\kappa,R}}e^{\theta Z}
 =
 \frac{I_{n,R}(\kappa-\theta)}{I_{n,R}(\kappa)}
 \leq
 \frac{I_n(\kappa-\theta)}
      {I_{n,R_0}(\kappa)}.
 \tag{39.14}
\]
Moreover,
\[
 \mu_{\kappa,R}\Longrightarrow\mu_\kappa
 \quad\hbox{as }R\to\infty,
 \tag{39.15}
\]
and every moment with exponent below \(\kappa\) converges.
\end{theorem}

\begin{proof}
The sets \(\mathcal K_R\) increase to \(\mathcal P_n\).
The numerator in (39.14) is at most the full integral, while the
denominator is at least its value at \(R_0\).  Monotone convergence gives
\[
 I_{n,R}(a)\uparrow I_n(a)
\]
for every \(a>0\), proving convergence of normalizations and of the
displayed exponential moments.  Bounded continuous test functions then
give (39.15).
\end{proof}

Thus the artificial hard ellipticity cutoff can be removed while
\(\kappa\) remains fixed.

\subsection{Why the stabilizer cannot be silently removed}

\begin{proposition}[Loss of a fixed exponential budget]
\label{prop:v11v39-kappa}
Let \(\kappa_i\downarrow0\).  No fixed \(\theta>0\) can satisfy the
moment identity (39.6) for all sufficiently large \(i\), because it
requires \(\theta<\kappa_i\).
\end{proposition}

\begin{proof}
For any fixed \(\theta>0\), eventually
\(\kappa_i-\theta\leq0\).  The numerator corresponding to
\(I_n(\kappa_i-\theta)\) is then not a decaying integral at infinity and
is infinite.  Hence the moment is unavailable.
\end{proof}

A continuum construction may keep a physical stabilizing term, replace
it with a different coercive mechanism, or prove cancellations by
methods not based on positive moments.  It may not use the fixed
\(\kappa\) estimates after sending \(\kappa\) to zero.

\subsection{Finite lattice and local conditioning}

On a finite set of metric sites \(\Lambda\), define the product measure
\[
 \mu_{\kappa,\Lambda}
 =
 \bigotimes_{x\in\Lambda}\mu_\kappa.
 \tag{39.16}
\]
It is positive and ultralocal.  If a real interaction \(V_\Lambda\)
is added and its block conditional action satisfies
\[
 V_B(g_B\mid g_{B^c})\geq-b
 \tag{39.17}
\]
together with a uniform comparison set of positive
\(\mu_{\kappa,B}\)-measure on which \(V_B\leq c\), then Theorem
\ref{thm:v11v38-conditional-moment} applies with the stabilizer
coefficients supplied by \(\kappa Z\).

\begin{corollary}[Metric row for the V37 good set]
\label{cor:v11v39-good-metric}
Under (39.17) and the comparison-set hypothesis, every
\(0<\theta<\kappa\) has a uniform local bound
\[
 \sup_{i,B}
 \mathbb E e^{\theta Z^{\rm met}_{i,B}}<\infty
 \tag{39.18}
\]
for blocks of fixed size, provided their interaction constants are
uniform.  The V37 Markov bound then yields a local ellipticity-failure
intensity.
\end{corollary}

\begin{proof}
Use (39.7) and reserve the difference \(\kappa-\theta>0\) in the
conditional numerator.  Equations (39.17) and the comparison set give
the \(b,c,v\) constants of V38.  Apply V37.
\end{proof}

For a site reflection, the product reference (39.16) factorizes across
the reflection plane and is reflection positive as an ultralocal
measure.  A coupled action preserves that property only if its
cross-plane terms admit a separate reflection-positive decomposition.

\subsection{Relation to the target programme}

The reference (39.16) supplies:
\begin{enumerate}[label=\textnormal{(\roman*)}]
\item a positive dominating metric measure;
\item two-sided conformal and ellipticity moments;
\item removal of the hard chamber at fixed \(\kappa\); and
\item polynomial coefficient moments for the V28 stencil.
\end{enumerate}
It does not yet supply stationary-Hessian, coarea, gauge-gap, or chiral
phase-denominator bounds.

\begin{firewall}[Reference versus Einstein dynamics]
\label{fw:v11v39-reference}
The stabilizer \(\kappa\operatorname{tr}(g+g^{-1})\) depends on the
chosen lattice frame and is not the Einstein--Hilbert action.  No
renormalization argument has shown that it is irrelevant, diffeomorphism
invariance has not been recovered, and Proposition
\ref{prop:v11v39-kappa} forbids using its moment budget after removing
it.  A physical continuum limit therefore requires a declared
stabilizer trajectory and counterterm analysis or a different positive
construction.
\end{firewall}

\begin{assessment}[V39 advance]
\label{ass:v11v39}
A fully positive finite metric reference now exists on the entire
positive cone, with exact exponential ellipticity moments and a
uniform hard-chamber removal.  Its conformal control is explicit, as is
the obstruction to sending the stabilizer strength to zero.  V40 is the
next mandatory companion audit; the following substantive step must
integrate the gauge-link Wilson energy and this metric reference into
one local dominating measure without claiming that the target Einstein
dynamics has already been recovered.
\end{assessment}
\section{V40 companion audit: exact Haar invariance and finite-cell scope}
\label{sec:v11v40}

This compulsory checkpoint inspected
\[
\begin{array}{c|c|r|l}
 \text{programme}&\text{version}&\text{bytes}&\text{SHA--256}\\ \hline
 \text{Yang--Mills}&33&471453&
 \texttt{ACC5EA1D076C855363B1317C668144914A65C5E245EC5B9D84B279D9BBBC344C}\\
 \text{Navier--Stokes}&26&278283&
 \texttt{82C887F8C2C69E4F28FAD7C3DC8DE5B9099CB5A4BF431EBA1FFF3E91935978AD}
\end{array}
\tag{40.1}
\]

\begin{audit}[Companion results]
\label{audit:v11v40-companions}
YM V33 completes one translated toron response through second external
order.  It checks all stationary-score jets through that order and
finds them zero because the two background colours coincide.  It also
proves that parameter-dependent left multiplication preserves product
Haar measure exactly, so every Haar Jacobian derivative vanishes in
that group chart.  The resulting number remains one polarization and
one internal momentum cell with its finite-grid weight.

NS V26 is unchanged and provides no new transfer.
\end{audit}

\subsection{Accepted Haar transfer}

Let \(G\) be compact with normalized Haar measure \(dU\).  For an
arbitrary parameter-dependent \(g(q)\in G\), define
\[
 L_q(U)=g(q)U.
 \tag{40.2}
\]

\begin{proposition}[Parameter-dependent Haar translation]
\label{prop:v11v40-haar}
For every fixed \(q\),
\[
 (L_q)_*dU=dU.
 \tag{40.3}
\]
Consequently the Radon--Nikodym derivative is identically one as a
function of \(q\), and all of its parameter derivatives vanish.  The
same holds for a product of independent link Haar measures with a
separate left multiplier on each link.
\end{proposition}

\begin{proof}
Equation (40.3) is left invariance of Haar measure.  Since the equality
holds for every parameter value, the Jacobian is the constant function
one; differentiating it gives zero.  Tensoring the identities proves the
product statement.
\end{proof}

This closes the Haar row only for a genuine group translation in product
Haar coordinates.  Gauge fixing, coarea reduction, a nonlinear
exponential-coordinate Lebesgue density, or mixing different links can
produce additional Jacobians and remains governed by V36.

\subsection{Accepted symmetry rule}

If an exact colour tensor factors through
\[
 [T_a,T_a]=0,
 \tag{40.4}
\]
then that complete tensor row vanishes.  The factorization must be
proved before norms or partial contractions are taken.

\begin{lemma}[Same-colour commutator closure]
\label{lem:v11v40-colour}
Suppose a multilinear response has the exact form
\[
 \mathcal R(X;B,B)
 =
 \mathcal L(X)\,\langle C,[B,B]\rangle
 \tag{40.5}
\]
for all retained \(X\).  Then \(\mathcal R(X;B,B)=0\) identically, and
all momentum derivatives that preserve the same colour factor also
vanish.
\end{lemma}

\begin{proof}
The Lie bracket is alternating, so \([B,B]=0\).  If momentum
differentiation acts only on \(\mathcal L\) and the kinematic
coefficients, the zero colour factor remains.
\end{proof}

For the full Standard Model species sum, different hypercharges and
noncommuting weak or colour legs do not generally reduce to (40.5).
The anomaly identities in V24 must still be applied to the complete
independent-leg tensor.

\subsection{Finite-cell scope}

A finite-grid response has the form
\[
 \mathcal R_M(q)
 =
 \frac1{M^4}\sum_{k\in\Gamma_M}\mathcal I_M(k,q).
 \tag{40.6}
\]
An exact toron value is the single summand
\(M^{-4}\mathcal I_M(0,q)\).

\begin{proposition}[One cell does not determine the mean]
\label{prop:v11v40-one-cell}
No bound or value for \(\mathcal I_M(0,q)\) alone determines the sign,
size, or limit of (40.6) without a bound on the remaining summands.
The toron contribution vanishes in magnitude as \(M^{-4}\) when its
integrand is uniformly bounded, but it must be retained at every finite
\(M\).
\end{proposition}

\begin{proof}
The remaining \(M^4-1\) summands can be chosen with arbitrary common
sign and magnitude while the zero summand is held fixed.  This proves
the first statement.  The second follows from its explicit grid weight.
Deleting the cell would change the exact finite sum.
\end{proof}

This is the same order-of-operations rule already used for flat
holonomies in V16 and for sector sums in V34.

\subsection{Direction after the audit}

V39 supplies a positive metric product reference.  Proposition
\ref{prop:v11v40-haar} supplies a positive gauge-link product reference
with no recentering Jacobian.  Their product is therefore the natural
dominating measure for the next finite calculation.

\begin{decision}[V41 combined reference]
\label{dec:v11v40-next}
V41 will combine link Haar measure, Wilson plaquette energy, and the V39
metric stabilizer.  It will derive the conditional block moment bound
with its exact crossing-plaquette penalty.  If that penalty prevents a
uniform small-plaquette tail, it will be recorded as a bottleneck rather
than hidden in a volume-independent constant.
\end{decision}

\begin{firewall}[Audit boundary]
\label{fw:v11v40-boundary}
YM V33 proves a vanishing stationary and Haar row only for its declared
same-colour group chart.  It supplies no Standard Model determinant
modulus, Einstein metric moment, nonzero-momentum mean, cylinder
comparison, or continuum measure.  NS V26 remains unchanged.
\end{firewall}

\begin{assessment}[V40 audit]
\label{ass:v11v40}
The required audit is complete.  Exact product-Haar translation is now
an accepted zero row, while gauge-fixed and coarea Jacobians remain
explicit.  The finite-cell warning prevents a toron regression from
being promoted to a continuum response.
\end{assessment}
\section{Combined positive gauge--metric reference and plaquette tails}
\label{sec:v11v41}

V40 combines product Haar measure with the stabilized metric reference
of V39.  A naive conditional-block comparison incurs a crossing-
plaquette penalty.  This note computes that penalty and then goes
upstream to a global entropy comparison, which yields a genuine
weak-coupling local plaquette tail for a translation-invariant positive
reference with uniformly positive gauge coefficients.

\subsection{Finite positive coupled reference}

Let \(\Lambda\) be a finite periodic four-dimensional lattice, with
oriented positive links \(E_\Lambda\) and plaquettes
\(\mathcal P_\Lambda\).  Let \(G\subset U(N)\) be compact.  Use
\[
 d\mu_{0,\Lambda}
 =
 \left(\bigotimes_{e\in E_\Lambda}dU_e\right)
 \otimes
 \left(\bigotimes_{x\in\Lambda}d\mu_{\kappa,R}(g_x)\right),
 \tag{41.1}
\]
where \(dU_e\) is normalized Haar measure and
\(\mu_{\kappa,R}\) is the V39 metric reference restricted to
\(\mathcal K_R\).

For
\[
 z_p(U)=1-\frac1N\operatorname{ReTr}U_p,
 \qquad 0\leq z_p\leq2,
 \tag{41.2}
\]
consider
\[
 S_{\rm g}(U,g)
 =
 \sum_{p\in\mathcal P_\Lambda}\beta_p(g)z_p(U),
 \tag{41.3}
\]
with
\[
 0<\beta_-\leq\beta_p(g)\leq\beta_+<\infty
 \tag{41.4}
\]
on the chamber.  Define
\[
 d\nu_\Lambda
 =
 Z_\Lambda^{-1}e^{-S_{\rm g}}d\mu_{0,\Lambda}.
 \tag{41.5}
\]
This is a positive probability measure.  Parameter-dependent link
recentering by left multiplication has unit Haar Jacobian by
Proposition \ref{prop:v11v40-haar}.

The bounds (41.4) are a finite-chamber hypothesis on the curved gauge
kinetic coefficients.  They must be verified for the chosen
discretization.

\subsection{The conditional boundary penalty}

Let \(E_B\) be the links varied in a finite block.  Split incident
plaquettes into interior plaquettes \(\mathcal P_B^\circ\), whose links
are all varied, and crossing plaquettes
\(\partial\mathcal P_B\).  Put
\[
 Z_B^\circ=\sum_{p\in\mathcal P_B^\circ}z_p,
 \qquad
 N_\partial=|\partial\mathcal P_B|.
 \tag{41.6}
\]
Suppose there is a block-link comparison set \(\mathcal A_B\) of Haar
measure \(v_B>0\) on which
\[
 Z_B^\circ\leq\delta_B.
 \tag{41.7}
\]

\begin{proposition}[Conditional moment with crossing cost]
\label{prop:v11v41-crossing}
For \(0\leq\theta\leq\beta_-\), the conditional gauge measure obeys
\[
 \mathbb E\left[
 e^{\theta Z_B^\circ}\mid U_{E_B^c},g
 \right]
 \leq
 v_B^{-1}
 \exp\left[
 \beta_+(\delta_B+2N_\partial)
 \right].
 \tag{41.8}
\]
Consequently,
\[
 \mathbb P\left\{
 Z_B^\circ>t\mid U_{E_B^c},g
 \right\}
 \leq
 v_B^{-1}
 \exp\left[
 -\theta t+\beta_+(\delta_B+2N_\partial)
 \right].
 \tag{41.9}
\]
\end{proposition}

\begin{proof}
The conditional action is nonnegative and satisfies
\[
 S_B\geq\beta_-Z_B^\circ.
\]
On \(\mathcal A_B\), the interior action is at most
\(\beta_+\delta_B\).  Every crossing plaquette has \(z_p\leq2\), so its
total action is at most \(2\beta_+N_\partial\).  The conditional
normalization is therefore at least
\[
 v_Be^{-\beta_+(\delta_B+2N_\partial)}.
\]
The numerator with \(e^{\theta Z_B^\circ}\) is at most one when
\(\theta\leq\beta_-\).  This proves (41.8), and exponential Markov gives
(41.9).
\end{proof}

For a one-plaquette threshold, the exponent in (41.9) need not decay:
the boundary cost is independent of the target defect and can be
larger.  Thus this crude conditional comparison does not prove the V14
maximum-plaquette tail.  The penalty may not be hidden in a
volume-independent constant.

\subsection{Global entropy comparison}

For \(0<\delta<1/4\), let
\[
 \mathcal A_\delta
 =
 \{U:\|U_e-I\|\leq\delta
       \text{ for every }e\in E_\Lambda\}.
 \tag{41.10}
\]
Let
\[
 v_G(\delta)
 =
 {\rm Haar}\{U\in G:\|U-I\|\leq\delta\}>0.
 \tag{41.11}
\]
By the product defect bound of V14, on \(\mathcal A_\delta\),
\[
 \|I-U_p\|\leq4\delta,
 \qquad
 z_p\leq4\delta.
 \tag{41.12}
\]
Therefore
\[
 Z_\Lambda
 \geq
 v_G(\delta)^{|E_\Lambda|}
 \exp\left(
 -4\beta_+\delta|\mathcal P_\Lambda|
 \right).
 \tag{41.13}
\]

Put
\[
 \overline z_\Lambda
 =
 \frac1{|\mathcal P_\Lambda|}
 \sum_{p\in\mathcal P_\Lambda}z_p,
 \qquad
 r_\Lambda=\frac{|E_\Lambda|}{|\mathcal P_\Lambda|}.
 \tag{41.14}
\]

\begin{theorem}[Finite-volume average-energy tail]
\label{thm:v11v41-average-tail}
For every \(\tau>0\) and \(0<\delta<1/4\),
\[
 \nu_\Lambda\{\overline z_\Lambda\geq\tau\}
 \leq
 \exp\left\{
 -|\mathcal P_\Lambda|
 \left[
 \beta_-\tau
 -4\beta_+\delta
 -r_\Lambda\log v_G(\delta)^{-1}
 \right]
 \right\}.
 \tag{41.15}
\]
\end{theorem}

\begin{proof}
On the event in question, (41.3)--(41.4) give
\[
 S_{\rm g}\geq
 \beta_-\tau|\mathcal P_\Lambda|.
\]
The numerator of its probability is therefore at most the exponential
of the negative right-hand side.  Divide by the partition-function
lower bound (41.13) and use (41.14).
\end{proof}

Since \(0\leq\overline z_\Lambda\leq2\),
\[
 \mathbb E\overline z_\Lambda
 \leq
 \tau+
 2\nu_\Lambda\{\overline z_\Lambda\geq\tau\}.
 \tag{41.16}
\]
For a periodic hypercubic lattice,
\[
 r_\Lambda=\frac{4}{6}=\frac23.
 \tag{41.17}
\]

\begin{corollary}[Thermodynamic local plaquette mean]
\label{cor:v11v41-local-mean}
Assume the joint measure is translation invariant.  Along periodic
volumes tending to infinity,
\[
 \limsup_{\Lambda\uparrow\mathbb Z^4}
 \mathbb E_{\nu_\Lambda}z_p
 \leq
 4\frac{\beta_+}{\beta_-}\delta
 +\frac{2}{3\beta_-}\log v_G(\delta)^{-1}
 \tag{41.18}
\]
for every \(0<\delta<1/4\).
\end{corollary}

\begin{proof}
Translation invariance makes
\(\mathbb E z_p=\mathbb E\overline z_\Lambda\).
In (41.15)--(41.16), take any
\[
 \tau>
 4\frac{\beta_+}{\beta_-}\delta
 +\frac{r_\Lambda}{\beta_-}\log v_G(\delta)^{-1}.
\]
The exponential probability tends to zero with the volume.  Then let
\(\tau\) decrease to the threshold and use (41.17).
\end{proof}

This is a mean-energy statement derived without independent
plaquettes.  It uses positivity and translation invariance of the full
gauge--metric reference.

\subsection{Weak-coupling rate}

A compact \(D_G\)-dimensional Lie group has constants
\(c_G,\delta_0>0\) such that
\[
 v_G(\delta)\geq c_G\delta^{D_G},
 \qquad0<\delta\leq\delta_0.
 \tag{41.19}
\]
This follows from a nonsingular exponential coordinate chart at the
identity and a positive smooth Haar density.

Assume
\[
 \frac{\beta_+}{\beta_-}\leq C_\beta,
 \qquad
 \beta_-\geq\delta_0^{-1}.
 \tag{41.20}
\]
Choose \(\delta=\beta_-^{-1}\) in (41.18).

\begin{corollary}[Explicit weak-coupling mean bound]
\label{cor:v11v41-weak}
There are constants \(C_0,C_1\), depending only on
\(G,C_\beta\), such that
\[
 \limsup_{\Lambda\uparrow\mathbb Z^4}
 \mathbb E z_p
 \leq
 \frac{C_0+C_1\log\beta_-}{\beta_-}.
 \tag{41.21}
\]
\end{corollary}

\begin{proof}
The first term in (41.18) is at most \(4C_\beta/\beta_-\).
Equation (41.19) gives
\[
 \log v_G(\beta_-^{-1})^{-1}
 \leq
 \log c_G^{-1}+D_G\log\beta_-.
\]
Insert this in the second term of (41.18).
\end{proof}

The logarithm is an entropy cost of placing every comparison link in a
small identity ball.  The estimate is not expected to be sharp.

\subsection{From mean energy to the V14 chord event}

Lemma \ref{lem:v11v38-wilson-chord} gives
\[
 \{\|I-U_p\|>\epsilon\}
 \subset
 \left\{z_p>\frac{\epsilon^2}{2N}\right\}.
 \tag{41.22}
\]
Markov's inequality and (41.21) yield
\[
 \limsup_{\Lambda\uparrow\mathbb Z^4}
 \nu_\Lambda\{\|I-U_p\|>\epsilon\}
 \leq
 \frac{2N}{\epsilon^2}
 \frac{C_0+C_1\log\beta_-}{\beta_-}.
 \tag{41.23}
\]

For a fixed block with \(N_{\mathcal P}(L)\) plaquettes, the union bound
gives
\[
 p_{\rm block}
 \leq
 \frac{2N\,N_{\mathcal P}(L)}{\epsilon^2}
 \frac{C_0+C_1\log\beta_-}{\beta_-}.
 \tag{41.24}
\]

\begin{theorem}[Sufficient summable gauge-gap schedule]
\label{thm:v11v41-schedule}
Let the V14 block scale, plaquette threshold, and lower gauge coefficient
at regulator \(i\) be \(L_i,\epsilon_i,\beta_{-,i}\), with
\[
 A_*(L_i)\epsilon_i\leq\frac{\gamma_m}{16}.
 \tag{41.25}
\]
If
\[
 \sum_i
 \frac{
 N_{\mathcal P}(L_i)
 [1+\log\beta_{-,i}]
 }{
 \beta_{-,i}\epsilon_i^2
 }
 <\infty,
 \tag{41.26}
\]
and the coefficient ratios and group constants are uniform, then the
thermodynamic local intensity of blocks failing the V14 plaquette
condition is summable.
\end{theorem}

\begin{proof}
Equation (41.25) makes every successful block a gapped-completion block.
Equation (41.24), with constants absorbed, bounds its failure
probability.  The series (41.26) sums those bounds.
\end{proof}

For the elementary filling estimate \(A_*(L)\leq C_dL^2\), choosing
\(\epsilon_i\) at the allowed order \(L_i^{-2}\) makes the denominator
in (41.26) contain an additional \(L_i^{-4}\); hence the required
growth of \(\beta_{-,i}\) must beat the combined plaquette entropy and
\(L_i^4\) gap cost.

\subsection{Scope of the positive result}

The estimate applies to (41.5), or to an additional positive
reweighting controlled by Theorem \ref{thm:v11v38-reweight}.  For a
complex chiral weight it applies under the absolute measure (38.19);
the phase-denominator row (38.21) remains necessary.

If the metric chamber expands, V39 keeps its reference moments uniform
at fixed \(\kappa\).  The coefficient bounds (41.4) must remain uniform
along that expansion; this is a separate property of the chosen curved
gauge action.

\begin{firewall}[Target dynamics and order of limits]
\label{fw:v11v41-target}
The measure (41.5) is a positive dominating reference, not yet the
target Standard Model--Einstein measure.  The thermodynamic limit is
taken before the weak-coupling conclusion in (41.21), and no
Osterwalder--Schrader limit has been constructed.  Uniform curved
coefficient bounds, matter-modulus reweighting, the chiral phase
denominator, the stationary-Hessian and coarea moments, and removal or
renormalization of the metric stabilizer are unproved.  Condition
(41.26) is a sufficient schedule, not a derivation of the physical
running coupling.
\end{firewall}

\begin{assessment}[V41 advance]
\label{ass:v11v41}
The first nontrivial gauge probability row is now proved for a positive
gauge--metric reference.  The naive conditional route exposes an
unacceptable crossing penalty; the global entropy comparison instead
gives a translation-invariant local plaquette mean of order
\((1+\log\beta_-)/\beta_-\), a chord-defect tail, and the explicit
summability schedule (41.26) that activates the V14 gapped completion.
The next finite task is to combine this gauge schedule with the V39
metric tail and V36 inverse-Hessian rows under one adjacent coupling,
then state the remaining measure-construction theorem as a minimal list
of acquisitions rather than a continuum claim.
\end{assessment}
\section{Master adjacent coupling and the cylinder-measure ledger}
\label{sec:v11v42}

V41 proves a summable gauge-gap schedule for a positive reference.
The remaining good-set rows cannot simply be assigned one fixed
ellipticity or inverse-Hessian threshold: under a fixed nondegenerate
reference, the probability of leaving a fixed compact set need not tend
to zero with the regulator.  The correct balance permits slowly growing
thresholds and charges the resulting deterministic stocks against the
adjacent approximation rate.  This note proves that balance and then
states a sufficient finite list for constructing the limiting cylinder
functional.

\subsection{The simultaneous adjacent good event}

Let \(\pi_i\) be a coupling of the level-\(i\) and transported
level-\(i+1\) fields on a fixed cylinder \(K\).  Define
\[
 \mathsf G_i
 =
 \mathsf G_i^{\rm gauge}
 \cap\mathsf G_i^{\rm met}
 \cap\mathsf G_i^{\rm Hess}
 \cap\mathsf G_i^{\rm co}
 \cap\mathsf G_i^{\rm der}
 \cap\mathsf G_i^{\rm atlas}.
 \tag{42.1}
\]
The rows mean, respectively: the V14 plaquette and gap condition; a
metric threshold for (38.10); inverse stationary-Hessian and coarea
thresholds for (38.12); the coefficient-derivative stocks of V28--V29;
and the logarithm-patch and nerve conditions of V32--V34.

Let
\[
 p_{j,i}:=\pi_i((\mathsf G_i^j)^c).
 \tag{42.2}
\]

\begin{lemma}[Combined failure intensity]
\label{lem:v11v42-union}
The complement probability satisfies
\[
 p_i:=\pi_i(\mathsf G_i^c)
 \leq
 \sum_jp_{j,i}.
 \tag{42.3}
\]
If \(p_{j,i}\leq M_{j,i}e^{-\theta_{j,i}t_{j,i}}\) for the moment rows,
and the gauge row obeys V41, then
\[
 \begin{split}
 p_i\leq{}&
 C_{\rm g}
 \frac{
 N_{\mathcal P}(L_i)[1+\log\beta_{-,i}]
 }{
 \beta_{-,i}\epsilon_i^2
 }\\
 &+
 \sum_{j\in\{{\rm met,Hess,co,der}\}}
 M_{j,i}e^{-\theta_{j,i}t_{j,i}}
 +p_{{\rm atlas},i}.
 \end{split}
 \tag{42.4}
\]
\end{lemma}

\begin{proof}
Equation (42.3) is the union bound.  Insert the gauge estimate (41.24)
and the exponential Markov estimates (37.4) for the other rows.
\end{proof}

The atlas failure includes configurations not covered by the declared
good logarithm patches and any unstable cycle certificate.  It is not
automatically included in the plaquette event unless the V14
completion-to-atlas implication has been proved for the chosen cover.

\subsection{Scale-dependent deterministic stocks}

Let \(\delta_i\) collect the native adjacent differences of links,
metrics, chart maps, homotopies, Hessians, coarea operators, shell
terms, and saturated minors after transport.  On \(\mathsf G_i\), all
the theorems of V23--V36 give a bound of the form
\[
 |C_{i+1}(\mathcal O)-C_i(\mathcal O)|
 \leq
 r_i,\qquad
 r_i\leq K_i\delta_i,
 \tag{42.5}
\]
where \(C_i(\mathcal O)\) is the complete saturated cylinder
coefficient and \(K_i\) is the compiled stock.  The product
\(K_i\delta_i\), rather than \(\delta_i\) alone, is the quantity that
must be summable.

\begin{proposition}[Growing-threshold balance]
\label{prop:v11v42-growing}
Suppose, for one badness row,
\[
 \mathbb E e^{\theta Z_i}\leq M,
 \tag{42.6}
\]
and choose
\[
 t_i=a\log(i+2).
 \tag{42.7}
\]
Assume its contribution to the deterministic stock satisfies
\[
 K_i\leq C(1+t_i)^r
 \tag{42.8}
\]
and its native adjacent error satisfies
\[
 \delta_i\leq C_\delta(i+2)^{-s}.
 \tag{42.9}
\]
For \(q>1\), both
\[
 \sum_iK_i\delta_i<\infty
 \tag{42.10}
\]
and
\[
 \sum_i
 \mathbb P\{Z_i>t_i\}^{\,1-1/q}<\infty
 \tag{42.11}
\]
hold whenever
\[
 s>1,\qquad
 \theta a\left(1-\frac1q\right)>1.
 \tag{42.12}
\]
\end{proposition}

\begin{proof}
Equations (42.7)--(42.9) bound the first summand by
\[
 C'(1+\log(i+2))^r(i+2)^{-s},
\]
whose series converges for \(s>1\).  Exponential Markov gives
\[
 \mathbb P\{Z_i>t_i\}^{1-1/q}
 \leq
 M^{1-1/q}
 (i+2)^{-\theta a(1-1/q)},
\]
and the second condition in (42.12) makes this series converge.
\end{proof}

This resolves the fixed-chamber obstruction at the level of a
sufficient schedule.  For example, the V36 inverse formulas grow
polynomially in a bound for \(\|A^{-1}\|\), and the V28 coefficients
often grow polynomially in \(\|g\|+\|g^{-1}\|\).  Logarithmic
thresholds preserve summability when the native adjacent errors have
any power strictly better than \(i^{-1}\).  The required power and
polynomial degree must be computed for the actual discretization.

\subsection{Master coefficient comparison}

Let \(q>1\).  Assume
\[
 \|C_{i+1}(\mathcal O)-C_i(\mathcal O)\|_{L^q(\pi_i)}
 \leq D_i.
 \tag{42.13}
\]
No uniformity of \(D_i\) is assumed.

\begin{theorem}[Complete adjacent cylinder criterion]
\label{thm:v11v42-master}
If
\[
 \sum_i r_i<\infty,
 \qquad
 \sum_iD_ip_i^{\,1-1/q}<\infty,
 \tag{42.14}
\]
then
\[
 \sum_i
 \mathbb E_{\pi_i}
 |C_{i+1}(\mathcal O)-C_i(\mathcal O)|
 <\infty.
 \tag{42.15}
\]
On a compatible projective coupling, the transported coefficients
converge in \(L^1\) and almost surely outside a null set after passage to
a subsequence; their expectations converge without taking a
subsequence.
\end{theorem}

\begin{proof}
Split the adjacent difference over \(\mathsf G_i\) and its complement.
Equation (42.5) bounds the first contribution by \(r_i\).
Hölder and (42.13) bound the second by
\(D_ip_i^{1-1/q}\).  Summation proves (42.15), and telescoping gives
\(L^1\) convergence on a compatible coupling.  Choose a subsequence
whose \(L^1\) distances to the limit are summable; Markov and
Borel--Cantelli give almost-sure convergence along it.  Finally,
\[
 |\mathbb EC_{i+1}-\mathbb EC_i|
 \leq\mathbb E|C_{i+1}-C_i|,
\]
so the expectations are Cauchy.
\end{proof}

Almost-sure convergence of the full sequence follows if the right-hand
side of (42.15) bounds adjacent differences on one common probability
space and their sum is finite almost surely.  The theorem does not infer
such a common realization from pairwise couplings alone.

\subsection{Compilation of the good-event increment}

A sufficient decomposition of \(r_i\) is
\[
 \begin{split}
 r_i={}&
 r_i^{\rm action}
 +r_i^{\rm modulus}
 +r_i^{\rm phase}
 +r_i^{\rm metric}\\
 &+r_i^{\rm stationary}
 +r_i^{\rm Hessdet}
 +r_i^{\rm coarea}
 +r_i^{\rm shell}
 +r_i^{\rm sector}.
 \end{split}
 \tag{42.16}
\]
The rows are populated as follows:
\begin{enumerate}[label=\textnormal{(\roman*)}]
\item V23, V31--V33 control the determinant-line current, chart, and
overlap phase;
\item V18 controls saturated minors and exterior powers through
fermion zeros;
\item V26--V29 control projector curvature and mixed gauge--metric
coefficients;
\item V36 controls stationary motion, Hessian determinants, and coarea
jets;
\item V13--V14 isolate completion-shell terms; and
\item V34 controls the sector sum under absolute tightness.
\end{enumerate}
Each entry in (42.16) is a transported difference of the complete row.
Cancellations such as V24 anomaly arithmetic or V40 Haar invariance are
applied before absolute values only when their exact tensor hypotheses
hold.

\begin{lemma}[Summation ledger]
\label{lem:v11v42-ledger}
If every row in (42.16) is nonnegative and summable, then \(r_i\) is
summable.  Omitting a row is valid only if it is absent from the
declared finite-cutoff density or has an exact zero identity.
\end{lemma}

\begin{proof}
A finite sum of convergent nonnegative series converges.  The second
statement follows because the triangle inequality used in (42.5)
requires a bound for every actual summand.
\end{proof}

\subsection{Normalization with a complex chiral phase}

Let
\[
 A_i(\mathcal O)
 =
 \mathbb E_{\overline\nu_i}
 [e^{i\vartheta_i}C_i(\mathcal O)],
 \qquad
 d_i=
 \mathbb E_{\overline\nu_i}e^{i\vartheta_i},
 \tag{42.17}
\]
where \(\overline\nu_i\) is the positive absolute measure of V38.
The normalized cylinder value is
\[
 W_i(\mathcal O)=\frac{A_i(\mathcal O)}{d_i}.
 \tag{42.18}
\]

\begin{proposition}[Stable quotient]
\label{prop:v11v42-quotient}
Suppose
\[
 |d_i|\geq d_0>0,\qquad |A_i(\mathcal O)|\leq M_{\mathcal O},
 \tag{42.19}
\]
and both \(A_i(\mathcal O)\) and \(d_i\) are Cauchy.  Then
\(W_i(\mathcal O)\) is Cauchy, with
\[
 |W_{i+1}-W_i|
 \leq
 \frac{|A_{i+1}-A_i|}{d_0}
 +\frac{M_{\mathcal O}|d_{i+1}-d_i|}{d_0^2}.
 \tag{42.20}
\]
\end{proposition}

\begin{proof}
Write the difference of the two quotients over
\(d_{i+1}d_i\), add and subtract \(A_i/d_{i+1}\), and use (42.19).
\end{proof}

The denominator row is therefore an adjacent comparison row as well as
a nonvanishing row.  A finite-cutoff lower bound with no convergence
estimate is insufficient.

\subsection{Cylinder functional and measure extension}

Let \(\mathfrak A_{\rm cyl}\) be a countable cylinder algebra closed
under conjugation and containing the constant one.  Assume its restriction is
uniformly dense in the continuous functions on each compact cylinder cutoff,
or in \(C_0\) after the stated tightness limit.  Suppose
\(W_i(\mathcal O)\) converges for every
\(\mathcal O\in\mathfrak A_{\rm cyl}\); define
\[
 W(\mathcal O)=\lim_iW_i(\mathcal O).
 \tag{42.21}
\]

\begin{theorem}[Passage of finite-cutoff axioms]
\label{thm:v11v42-axioms}
Assume the following estimates are uniform on each fixed cylinder:
\[
 W_i(1)=1,
 \tag{42.22}
\]
\[
 W_i(F^*F)\geq0,
 \tag{42.23}
\]
and
\[
 |W_i(F)|\leq C_K\|F\|_\infty
 \tag{42.24}
\]
for cylinder functions supported in \(K\).  Then \(W\) is normalized,
positive, and bounded on each cylinder.  If the finite-dimensional
marginals are projectively consistent and tight, \(W\) extends to a
probability measure on the projective field space.
\end{theorem}

\begin{proof}
Take limits in (42.22)--(42.24); the positive cone is closed.
For every finite cylinder, the bounded positive functional is
represented by a finite Borel probability measure.  Projective
consistency makes these marginals compatible.  Tightness and the
Kolmogorov extension theorem on the declared standard Borel projective
system produce the measure.
\end{proof}

For a complex chiral functional, (42.23) need not hold after the
fermions are integrated out.  One must either reconstruct from the full
boson--Grassmann Schwinger system with its correct positivity statement,
or prove that the normalized bosonic functional is positive.  Absolute
measure positivity alone is not the target assertion.

\subsection{Reflection positivity and reconstruction}

Let \(\Theta\) be a lattice reflection and
\(\mathfrak A_+\) the positive-time cylinder algebra.  If
\[
 W_i((\Theta F)F)\geq0
 \tag{42.25}
\]
for every \(F\in\mathfrak A_+\) once the cylinder is represented at
level \(i\), then
\[
 W((\Theta F)F)\geq0.
 \tag{42.26}
\]
The proof is again passage to the limit.  Euclidean covariance passes
when the finite-cutoff covariance defect tends to zero on every fixed
cylinder.  Clustering requires a separate uniform separation estimate;
it does not follow from projective convergence.

\begin{programme}[Minimal remaining measure acquisitions]
\label{prog:v11v42-acquisitions}
A sufficient nonduplicated ledger after V41 is:
\begin{enumerate}[label=\textnormal{(M\arabic*)}]
\item compatible adjacent couplings or an equivalent direct comparison
of all finite-cylinder marginals;
\item the gauge schedule (41.25)--(41.26) with curved coefficient
bounds;
\item exponential moments for metric, inverse-Hessian, coarea, and
derivative badness, with a growing-threshold balance such as
Proposition \ref{prop:v11v42-growing};
\item summable native adjacent errors after multiplication by every
scale-dependent deterministic stock;
\item saturated-minor modulus comparison and the V33 finite phase-nerve
cycle certificates;
\item absolute sector tightness and compatible topological phases;
\item a Cauchy, uniformly nonzero chiral phase denominator;
\item projective consistency and tightness of the normalized cylinder
marginals; and
\item reflection positivity, Euclidean covariance restoration,
regularity, and clustering estimates for reconstruction.
\end{enumerate}
\end{programme}

\begin{firewall}[The master theorem is conditional]
\label{fw:v11v42-conditional}
Theorem \ref{thm:v11v42-master} proves convergence once the displayed
rows are acquired; it does not prove their hypotheses for the target
coupled density.  In particular, no inverse-Hessian or coarea moment,
native adjacent rate, phase-denominator bound, reflection-positive
chiral regulator, or clustering estimate has been constructed.
The positive reference of V39--V41 remains stabilized and is not yet
identified with Einstein dynamics.
\end{firewall}

\begin{assessment}[V42 advance]
\label{ass:v11v42}
All deterministic and probabilistic rows now enter one adjacent
coupling theorem.  Logarithmically growing good-set thresholds can make
fixed-reference tails summable without demanding uniform compact
support, provided the native regulator differences beat \(i^{-1}\)
after their polynomial stocks are included.  The limiting cylinder
functional and the exact axioms that pass by closure have been
separated from reflection positivity and clustering.  The next finite
task is to acquire one previously missing analytic row rather than add
more bookkeeping: a quantitative small-singular-value estimate for the
stationary Hessian and coarea Gram operator under a positive
finite-dimensional reference.
\end{assessment}
\section{Small-singular-value control for Hessian and coarea rows}
\label{sec:v11v43}

V42 identifies inverse stationary-Hessian and coarea-Gram tails as
missing analytic inputs.  A finite-dimensional scalar shift gives an
exact anti-concentration estimate without independence of matrix
entries.  Its tail is polynomial rather than exponential, so this note
also derives the corresponding growing-threshold schedule and exposes
the cost of removing the shift.

\subsection{Hermitian scalar-shift anti-concentration}

Let \(A_0\) be an \(n\times n\) Hermitian matrix, possibly depending on
other fields.  Let \(t\) be a real auxiliary variable whose conditional
density, given all other fields, is bounded by \(\rho\).  Set
\[
 A(t)=A_0+tI.
 \tag{43.1}
\]

\begin{theorem}[Exact scalar-shift small ball]
\label{thm:v11v43-shift}
For every \(\varepsilon>0\),
\[
 \mathbb P\{
 s_{\min}(A(t))\leq\varepsilon
 \mid A_0
 \}
 \leq2n\rho\varepsilon.
 \tag{43.2}
\]
Consequently, for \(K>0\),
\[
 \mathbb P\{
 \|A(t)^{-1}\|>K
 \mid A_0
 \}
 \leq\frac{2n\rho}{K}.
 \tag{43.3}
\]
The same bounds hold after averaging over \(A_0\).
\end{theorem}

\begin{proof}
Let \(\lambda_1,\ldots,\lambda_n\) be the eigenvalues of \(A_0\).
Those of \(A(t)\) are \(\lambda_j+t\), hence
\[
 s_{\min}(A(t))
 =
 \min_j|\lambda_j+t|.
\]
The event in (43.2) is contained in the union of the \(n\) intervals
\[
 [-\lambda_j-\varepsilon,-\lambda_j+\varepsilon].
\]
Each has conditional probability at most \(2\rho\varepsilon\).
The union bound proves (43.2).  On the invertible set,
\(\|A^{-1}\|=s_{\min}(A)^{-1}\); set
\(\varepsilon=K^{-1}\) to obtain (43.3).
\end{proof}

No distributional hypothesis on \(A_0\) is used.  Exact gauge zero modes
must first be removed or gauge fixed: shifting a protected zero mode
would change the symmetry rather than estimate the physical reduced
Hessian.

\subsection{A removable finite shift}

If
\[
 t\sim{\rm Uniform}[-T,T],
 \tag{43.4}
\]
then \(\rho=(2T)^{-1}\) and
\[
 \mathbb P\{\|A(t)^{-1}\|>K\}
 \leq\frac n{TK}.
 \tag{43.5}
\]
The perturbation (43.1) is generated by adding
\[
 \frac t2\|y\|^2
 \tag{43.6}
\]
to the eliminated-variable action.  Thus it is an explicit finite
regulator, not a property of the unshifted target theory.

There is a deterministic alternative when \(A_0\geq0\).  For a ridge
\(\tau>0\),
\[
 A_\tau=A_0+\tau I\geq\tau I,
 \qquad
 \|A_\tau^{-1}\|\leq\tau^{-1}.
 \tag{43.7}
\]
This preserves positivity but again changes the action.

\begin{proposition}[Resolvent cost of ridge removal]
\label{prop:v11v43-ridge}
For \(\tau,\widetilde\tau>0\) and \(A_0\geq0\),
\[
 \|A_\tau^{-1}-A_{\widetilde\tau}^{-1}\|
 \leq
 \frac{|\tau-\widetilde\tau|}
 {\tau\widetilde\tau}.
 \tag{43.8}
\]
\end{proposition}

\begin{proof}
The resolvent identity gives
\[
 A_\tau^{-1}-A_{\widetilde\tau}^{-1}
 =
 (\widetilde\tau-\tau)
 A_\tau^{-1}A_{\widetilde\tau}^{-1}.
\]
Use (43.7).
\end{proof}

Hence a ridge sequence \(\tau_i\downarrow0\) is compatible with an
adjacent inverse limit only if the weighted differences
\[
 \sum_i
 \frac{|\tau_{i+1}-\tau_i|}
 {\tau_i\tau_{i+1}}
 \tag{43.9}
\]
are controlled together with the native Hessian differences.  The
series in (43.9) cannot converge for a monotone sequence tending to
zero, because its summands telescope exactly:
\[
 \frac{\tau_i-\tau_{i+1}}{\tau_i\tau_{i+1}}
 =
 \tau_{i+1}^{-1}-\tau_i^{-1}.
 \tag{43.10}
\]
Thus operator-norm convergence of the full inverse cannot be obtained
by this crude ridge bound as the ridge is removed.  Observables may
still converge after projection, cancellation, or integration, but
that requires a sharper target-specific argument.

\subsection{General polynomial small-ball schedule}

The scalar shift has exponent one.  To keep the scale balance visible,
assume more generally
\[
 \mathbb P\{s_{\min}(A_i)\leq\varepsilon\}
 \leq C_i\varepsilon^\alpha,
 \qquad 0<\varepsilon\leq1,
 \tag{43.11}
\]
for some \(\alpha>0\).  Then
\[
 \mathbb P\{\|A_i^{-1}\|>K\}
 \leq C_iK^{-\alpha},
 \qquad K\geq1.
 \tag{43.12}
\]

Suppose the deterministic good-set stock grows like
\[
 K_i^r
 \tag{43.13}
\]
when \(\|A_i^{-1}\|\leq K_i\), and the native adjacent error obeys
\[
 \delta_i\leq C_\delta(i+2)^{-s}.
 \tag{43.14}
\]
Choose
\[
 K_i=(i+2)^a.
 \tag{43.15}
\]

\begin{theorem}[Polynomial-tail balance]
\label{thm:v11v43-polynomial-balance}
Let \(q>1\), assume \(\sup_iC_i<\infty\), and suppose
\[
 a\alpha\left(1-\frac1q\right)>1,
 \tag{43.16}
\]
\[
 s-ar>1.
 \tag{43.17}
\]
Then
\[
 \sum_i
 \mathbb P\{\|A_i^{-1}\|>K_i\}^{\,1-1/q}<\infty
 \tag{43.18}
\]
and
\[
 \sum_iK_i^r\delta_i<\infty.
 \tag{43.19}
\]
A permissible exponent \(a\) exists whenever
\[
 s>
 1+\frac{r}{\alpha(1-1/q)}.
 \tag{43.20}
\]
\end{theorem}

\begin{proof}
Equation (43.12) and (43.15) bound the summand in (43.18) by a constant
times
\[
 (i+2)^{-a\alpha(1-1/q)},
\]
which is summable by (43.16).  Equations (43.13)--(43.15) bound the
deterministic summand by
\[
 C(i+2)^{ar-s},
\]
which is summable by (43.17).  The interval
\[
 \frac1{\alpha(1-1/q)}
 <a<
 \frac{s-1}{r}
\]
is nonempty precisely under (43.20), with the evident interpretation
when \(r=0\).
\end{proof}

This theorem replaces the logarithmic-threshold proposition in V42 for
inverse variables with polynomial small-ball tails.

\subsection{A vanishing scalar-smoothing window}

Let the smoothing width in (43.4) depend on the regulator:
\[
 T_i=(i+2)^{-b},\qquad b\geq0.
 \tag{43.21}
\]
Equation (43.5) and the threshold (43.15) give
\[
 \mathbb P\{\|A_i^{-1}\|>K_i\}
 \leq n(i+2)^{b-a}.
 \tag{43.22}
\]

\begin{corollary}[Cost of removing scalar smoothing]
\label{cor:v11v43-window}
The complement and deterministic series converge if
\[
 (a-b)\left(1-\frac1q\right)>1,
 \qquad
 s-ar>1.
 \tag{43.23}
\]
Such an \(a\) exists only if
\[
 b+\frac1{1-1/q}
 <
 \frac{s-1}{r}.
 \tag{43.24}
\]
\end{corollary}

\begin{proof}
Apply the same two power-series tests as in Theorem
\ref{thm:v11v43-polynomial-balance}, using exponent \(a-b\) in the
probability row.
\end{proof}

The smoothing width may therefore shrink only if the native adjacent
approximation is sufficiently faster than the inverse-stock growth.
This is a quantitative regulator-removal condition.

\subsection{Coarea Gram operators}

Let \(J\) be the derivative of a regular constraint map and
\[
 G=JJ^*.
 \tag{43.25}
\]
On the normal space,
\[
 \lambda_{\min}(G)=s_{\min}(J)^2,
 \qquad
 \|G^{-1}\|=s_{\min}(J)^{-2}.
 \tag{43.26}
\]

\begin{proposition}[Jacobian small ball to coarea tail]
\label{prop:v11v43-coarea}
If
\[
 \mathbb P\{s_{\min}(J_i)\leq\varepsilon\}
 \leq C\varepsilon^\alpha,
 \tag{43.27}
\]
then
\[
 \mathbb P\{\|G_i^{-1}\|>K\}
 \leq CK^{-\alpha/2}.
 \tag{43.28}
\]
Thus the polynomial balance theorem applies to the coarea row with
small-ball exponent \(\alpha/2\).
\end{proposition}

\begin{proof}
By (43.26),
\[
 \{\|G_i^{-1}\|>K\}
 =
 \{s_{\min}(J_i)<K^{-1/2}\}.
\]
Insert \(\varepsilon=K^{-1/2}\) in (43.27).
\end{proof}

Alternatively, a regularized Gram operator
\[
 G_\tau=JJ^*+\tau I
 \tag{43.29}
\]
has \(\|G_\tau^{-1}\|\leq\tau^{-1}\), but its determinant is a ridge
coarea factor, not the exact coarea density.  Removal again requires an
argument beyond the crude operator-norm estimate (43.8).

\subsection{Logarithmic versus inverse moments}

A tail \(P\{\|A^{-1}\|>K\}\leq CK^{-1}\) implies finite moments
\[
 \mathbb E\|A^{-1}\|^p<\infty
 \quad\hbox{for }0<p<1,
 \tag{43.30}
\]
but the bound alone supplies no moment of order \(p\geq1\).

\begin{proof}
For a nonnegative random variable \(X\),
\[
 \mathbb EX^p
 =
 p\int_0^\infty t^{p-1}\mathbb P\{X>t\}\,dt.
\]
Use the trivial probability bound below one and \(Ct^{-1}\) above one.
The tail integral converges exactly for \(p<1\).
\end{proof}

The V36 response contains several powers of \(A^{-1}\), so a scalar
shift does not by itself give the uniform \(L^q\) stock (42.13).
It does, however, give a valid truncated-good-set probability row to be
combined with separate bounds on the complete coefficient.

\subsection{Insertion into the master theorem}

For a stationary Hessian smoothed as in (43.1), Theorem
\ref{thm:v11v43-shift} supplies
\[
 p_{{\rm Hess},i}
 \leq\frac{2n_i\rho_i}{K_i}.
 \tag{43.31}
\]
For a constraint Jacobian satisfying (43.27), Proposition
\ref{prop:v11v43-coarea} supplies
\[
 p_{{\rm co},i}
 \leq C_iK_i^{-\alpha/2}.
 \tag{43.32}
\]
Choose thresholds by Theorem
\ref{thm:v11v43-polynomial-balance} and insert these two rows in
(42.4).  The V36 constants determine the power \(r\) of each inverse
threshold in the good-event increment.

\begin{firewall}[Regularization and transversality]
\label{fw:v11v43-regularization}
The scalar shift and ridge are auxiliary regulators.  No symmetry-
preserving removal has been proved for the Standard Model--Einstein
stationary Hessian, and exact gauge modes must be separated first.
No Jacobian small-ball estimate (43.27) has been proved for the physical
gauge-fixing or constraint map.  The exponent-one shift tail is too weak
to supply the high inverse moments required by the complete response
without truncation and a separate \(L^q\) estimate.
\end{firewall}

\begin{assessment}[V43 advance]
\label{ass:v11v43}
A genuine finite-dimensional inverse-Hessian probability estimate is
now available: a scalar shift with bounded conditional density gives
the exact \(K^{-1}\) tail, and a general polynomial small-ball exponent
has been integrated into the V42 scale balance.  The coarea exponent is
shown to halve when passing from a constraint Jacobian to its Gram
operator.  The calculation also proves that crude ridge operator-norm
control cannot survive monotone ridge removal.  The next finite task is
therefore upstream: isolate exact gauge zero modes and formulate the
stationary Hessian on a transverse slice before seeking a stronger
small-ball exponent.
\end{assessment}
\section{Gauge zero modes, rooted orbit Gram bounds, and transverse Hessians}
\label{sec:v11v44}

V43 shows that an inverse-Hessian estimate is meaningless until exact
symmetry zero modes are removed.  This note performs that reduction for
a compact lattice gauge action.  On a rooted finite block, the
gauge-orbit derivative has a deterministic lower singular-value bound
for every link field.  The stationary Hessian is then transported on
the orthogonal transverse fibre, where the V36 and V43 estimates apply.

\subsection{Gauge invariance forces Hessian zero modes}

Let a compact Lie group \(\mathcal G\) act smoothly and isometrically on
a finite-dimensional field manifold \(\mathcal M\).  For
\(\xi\in\mathfrak g\), let
\[
 R_y\xi=\xi^\#_y\in T_y\mathcal M
 \tag{44.1}
\]
be the infinitesimal orbit map.  Let \(S:\mathcal M\to\mathbb R\) be
\(\mathcal G\)-invariant and let \(y_*\) be critical:
\[
 dS_{y_*}=0.
 \tag{44.2}
\]
Write \(A=\operatorname{Hess}_{y_*}S\) as a self-adjoint operator.

\begin{theorem}[Exact orbit kernel]
\label{thm:v11v44-orbit-kernel}
For every \(\xi\in\mathfrak g\),
\[
 A R_{y_*}\xi=0.
 \tag{44.3}
\]
Thus the full Hessian is never invertible when the orbit through
\(y_*\) has positive dimension.
\end{theorem}

\begin{proof}
Gauge invariance gives the identity
\[
 dS_y(R_y\xi)=0
 \tag{44.4}
\]
for every \(y\).  Differentiate it at \(y_*\) in an arbitrary direction
\(v\):
\[
 \operatorname{Hess}_{y_*}S(v,R_{y_*}\xi)
 +dS_{y_*}(D_vR\,\xi)=0.
\]
The second term vanishes by (44.2).  Since \(v\) is arbitrary and the
Hessian is self-adjoint, (44.3) follows.
\end{proof}

A scalar shift \(A+tI\) would move these exact gauge modes and break the
identity.  Any symmetry-preserving shift must act only after the orbit
directions have been removed.

\subsection{The rooted lattice orbit derivative}

Let \(B=(V_B,E_B)\) be a finite connected graph with root \(o\), graph
diameter \(D_B\), and a compact gauge group \(G\) whose Lie algebra has
an \(\operatorname{Ad}\)-invariant inner product.  For a link field
\(U\), define the covariant graph gradient
\[
 (R_U\xi)_{xy}
 =
 \xi_x-\operatorname{Ad}_{U_{xy}}\xi_y,
 \qquad xy\in E_B,
 \tag{44.5}
\]
on the rooted gauge algebra
\[
 \mathfrak g_{B,o}
 =
 \{\xi:V_B\to\mathfrak g:\xi_o=0\}.
 \tag{44.6}
\]

\begin{theorem}[Uniform rooted orbit bound]
\label{thm:v11v44-rooted-bound}
For every link field \(U\) and every
\(\xi\in\mathfrak g_{B,o}\),
\[
 \|\xi\|_{\ell^2(V_B)}^2
 \leq
 |V_B|D_B\,
 \|R_U\xi\|_{\ell^2(E_B)}^2.
 \tag{44.7}
\]
Consequently,
\[
 s_{\min}(R_U)
 \geq(|V_B|D_B)^{-1/2},
 \tag{44.8}
\]
and the rooted orbit Gram operator
\[
 G_U:=R_U^*R_U
 \tag{44.9}
\]
satisfies
\[
 G_U\geq(|V_B|D_B)^{-1}I,
 \qquad
 \|G_U^{-1}\|\leq|V_B|D_B.
 \tag{44.10}
\]
\end{theorem}

\begin{proof}
For each vertex \(x\), choose a path of length at most \(D_B\) from
\(o\) to \(x\).  Starting with \(\xi_o=0\), equation (44.5) expresses
\(\xi_x\) as a sum of the edge defects along that path, each transported
by a product of adjoint actions.  Those transports are isometries.
Cauchy--Schwarz therefore gives
\[
 \|\xi_x\|^2
 \leq
 D_B\sum_{e\in{\rm path}(o,x)}
 \|(R_U\xi)_e\|^2
 \leq
 D_B\|R_U\xi\|_{\ell^2(E_B)}^2.
\]
Sum over the \(|V_B|\) vertices to obtain (44.7).
Equations (44.8)--(44.10) are its singular-value and Gram-operator
forms.
\end{proof}

This result is independent of plaquette curvature and has no Gribov
small-ball event.  Its price is the deterministic block factor
\(|V_B|D_B\).  For a four-dimensional block of side \(L\), the crude
stock grows at order \(L^5\).

The root removes the constant global gauge freedom.  If a physical
stabilizer remains after rooting, its Lie algebra must also be
factored out before (44.8) is interpreted as the smallest physical
orbit singular value.

\subsection{Orthogonal orbit and transverse projectors}

On the rooted algebra define
\[
 P_U^{\rm orb}
 =
 R_U(R_U^*R_U)^{-1}R_U^*,
 \qquad
 Q_U=I-P_U^{\rm orb}.
 \tag{44.11}
\]

\begin{proposition}[Exact orthogonal splitting]
\label{prop:v11v44-splitting}
The operator \(P_U^{\rm orb}\) is the orthogonal projector onto
\(\operatorname{ran}R_U\), and \(Q_U\) projects onto
\[
 \mathcal H_U^\perp=\ker R_U^*.
 \tag{44.12}
\]
At a critical point of a gauge-invariant action,
\[
 A=Q_UAQ_U.
 \tag{44.13}
\]
\end{proposition}

\begin{proof}
Invertibility of \(R_U^*R_U\) follows from Theorem
\ref{thm:v11v44-rooted-bound}.  Direct multiplication shows that
\(P_U^{\rm orb}\) is self-adjoint and idempotent, with range equal to
\(\operatorname{ran}R_U\).  Its orthogonal complement is
\(\ker R_U^*\).  Theorem \ref{thm:v11v44-orbit-kernel} gives
\(AR_U=0\); self-adjointness gives \(R_U^*A=0\).  These two identities
are equivalent to (44.13).
\end{proof}

The physical stationary Hessian is therefore
\[
 A_U^\perp
 :=
 A|_{\mathcal H_U^\perp}.
 \tag{44.14}
\]
Only \(A_U^\perp\) may enter the inverse and determinant formulas of
V36.

\subsection{Variation of the transverse fibre}

Let \(R,\widetilde R\) be two rooted orbit maps with
\[
 \|R\|,\|\widetilde R\|\leq M_R,\qquad
 s_{\min}(R),s_{\min}(\widetilde R)\geq\sigma>0,
 \tag{44.15}
\]
and
\[
 \|R-\widetilde R\|\leq\delta_R.
 \tag{44.16}
\]
Let \(P,\widetilde P\) be their orbit projectors.

\begin{theorem}[Orbit-projector comparison]
\label{thm:v11v44-projector-comparison}
One has
\[
 \|P-\widetilde P\|
 \leq
 \left(
 2M_R\sigma^{-2}
 +2M_R^3\sigma^{-4}
 \right)\delta_R.
 \tag{44.17}
\]
The same bound holds for \(Q-\widetilde Q\).
\end{theorem}

\begin{proof}
Put \(G=R^*R\) and
\(\widetilde G=\widetilde R^*\widetilde R\).  Then
\[
 \|G-\widetilde G\|
 \leq2M_R\delta_R,
 \tag{44.18}
\]
and
\[
 \|G^{-1}-\widetilde G^{-1}\|
 \leq
 \|G^{-1}\|\|G-\widetilde G\|
 \|\widetilde G^{-1}\|
 \leq2M_R\sigma^{-4}\delta_R.
 \tag{44.19}
\]
Subtract
\(RG^{-1}R^*\) and
\(\widetilde R\widetilde G^{-1}\widetilde R^*\), changing the left
factor, inverse, and right factor one at a time.  The two outer changes
contribute \(2M_R\sigma^{-2}\delta_R\); the inverse change contributes
\(2M_R^3\sigma^{-4}\delta_R\).  This proves (44.17).
Since \(Q=I-P\), the complementary estimate is identical.
\end{proof}

On the rooted lattice, Theorem \ref{thm:v11v44-rooted-bound} permits
\[
 \sigma^{-2}=|V_B|D_B.
 \tag{44.20}
\]
Thus the moving transverse-fibre cost is explicit and must be included
in an adjacent block schedule.

\subsection{Kato transport of transverse Hessians}

Set
\[
 d=\|Q-\widetilde Q\|<1.
 \tag{44.21}
\]
The V21 direct rotation gives a unitary \(\mathcal U\) with
\[
 \mathcal U\,\operatorname{ran}Q
 =
 \operatorname{ran}\widetilde Q.
 \tag{44.22}
\]
A literal bound from its polar formula is
\[
 \|\mathcal U-I\|
 \leq
 \omega(d):=
 \frac d{\sqrt{1-d^2}}
 +\frac1{\sqrt{1-d^2}}-1.
 \tag{44.23}
\]

\begin{proof}
The direct-rotation numerator is
\[
 C=\widetilde QQ+(I-\widetilde Q)(I-Q)
\]
and satisfies
\[
 C-I=(\widetilde Q-Q)(2Q-I),
 \qquad \|C-I\|\leq d.
\]
Its polar factor is
\[
 \mathcal U=C[I-(\widetilde Q-Q)^2]^{-1/2}.
\]
The inverse square-root factor has norm at most
\((1-d^2)^{-1/2}\) and differs from the identity by at most
\((1-d^2)^{-1/2}-1\).  Splitting
\(\mathcal U-I\) gives (44.23).
\end{proof}

Transport the second Hessian back to the first transverse fibre:
\[
 \widehat A^\perp
 =
 \mathcal U^*\widetilde A^\perp\mathcal U
 \quad\hbox{on }\operatorname{ran}Q.
 \tag{44.24}
\]

\begin{theorem}[Transverse Hessian comparison]
\label{thm:v11v44-hessian-comparison}
If
\[
 \|A\|,\|\widetilde A\|\leq M_A,\qquad
 \|A-\widetilde A\|\leq\delta_A,
 \tag{44.25}
\]
then
\[
 \|A^\perp-\widehat A^\perp\|
 \leq
 \delta_A+2M_A\omega(d).
 \tag{44.26}
\]
If both transverse inverses are bounded by \(\kappa\), then
\[
 \|(A^\perp)^{-1}-(\widehat A^\perp)^{-1}\|
 \leq
 \kappa^2[
 \delta_A+2M_A\omega(d)].
 \tag{44.27}
\]
\end{theorem}

\begin{proof}
Restriction cannot increase the norm of \(A-\widetilde A\).
Moreover,
\[
 \|\mathcal U^*\widetilde A\mathcal U-\widetilde A\|
 \leq
 \|(\mathcal U^*-I)\widetilde A\mathcal U\|
 +\|\widetilde A(\mathcal U-I)\|
 \leq2M_A\|\mathcal U-I\|.
\]
This proves (44.26).  Apply the inverse resolvent identity on the fixed
space \(\operatorname{ran}Q\) to obtain (44.27).
\end{proof}

The determinant factor is the pseudodeterminant
\[
 \det{}'A:=\det A^\perp.
 \tag{44.28}
\]
After Kato transport, its first two logarithmic derivatives are exactly
the V36 formulas on one fixed fibre.  If \(A^\perp\) is indefinite, its
absolute determinant and sign or phase must be treated separately.

\subsection{Symmetry-preserving small-ball regularization}

A transverse scalar smoothing has the form
\[
 A(t)=A+tQ.
 \tag{44.29}
\]
It leaves the orbit kernel exact and restricts to
\[
 A^\perp(t)=A^\perp+tI_{\mathcal H^\perp}.
 \tag{44.30}
\]

\begin{corollary}[Transverse scalar-shift tail]
\label{cor:v11v44-transverse-tail}
If the conditional density of \(t\) is bounded by \(\rho\) and
\(n_\perp=\dim\mathcal H^\perp\), then
\[
 \mathbb P\{
 \|(A^\perp(t))^{-1}\|>K
 \}
 \leq
 \frac{2n_\perp\rho}{K}.
 \tag{44.31}
\]
\end{corollary}

\begin{proof}
Apply Theorem \ref{thm:v11v43-shift} on the transverse fibre
(44.30).
\end{proof}

Equation (44.29) is an operator statement on a local slice.  A
gauge-invariant action realizing it globally is not automatic because
\(Q\) itself depends on the background.

\subsection{Coarea consequence for gauge orbits}

The rooted orbit Gram determinant satisfies the deterministic bound
\[
 \det G_U
 \geq
 (|V_B|D_B)^{-\dim\mathfrak g_{B,o}}.
 \tag{44.32}
\]
Hence its inverse and logarithm have no random small-ball event at fixed
block size.  Their scale growth is entirely explicit.

\begin{corollary}[Gauge-orbit coarea row]
\label{cor:v11v44-gauge-coarea}
For a rooted block, the gauge-orbit Gram inverse row in (42.4) may be
removed from the probability complement and placed in the deterministic
stock with
\[
 K_{{\rm orb},B}=|V_B|D_B.
 \tag{44.33}
\]
Other constraint Grams, including diffeomorphism, lapse, shift, or
nonlinear gauge-fixing maps, are not covered by this conclusion.
\end{corollary}

\begin{proof}
Equation (44.10) holds for every link field, so the failure probability
is zero.  Its deterministic inverse bound is (44.33).
\end{proof}

\begin{firewall}[Local slice and remaining zero modes]
\label{fw:v11v44-slice}
Rooting gives a local algebraic orbit splitting, not a global
gauge-fixing section.  Patch overlaps can still have Gribov topology,
and the full Standard Model--Einstein system also has diffeomorphism
and gravitational constraint directions.  No stronger small-ball
exponent for the transverse stationary Hessian has been proved, and no
global action realizing the smoothing (44.29) has been constructed.
The block factor in (44.33) must be charged in the cofinal schedule.
\end{firewall}

\begin{assessment}[V44 advance]
\label{ass:v11v44}
Exact gauge zero modes have been separated before inversion.  On every
rooted connected block, the gauge-orbit derivative has the uniform
singular-value bound (44.8), so its coarea Gram row is deterministic
rather than probabilistic.  Kato transport now compares stationary
Hessians on one transverse fibre, and scalar anti-concentration applies
without shifting protected gauge modes.  V45 is the next mandatory
companion audit; the following substantive note should examine whether
the YM finite response supplies a sharper transverse-Hessian or
polarization input before the programme turns to gravitational
constraint zero modes.
\end{assessment}
% =============================================================================
% BEGIN APPENDED ELEVENTH-SERIES V45 MATERIAL
% =============================================================================

\section{Companion audit at V45: quotient symmetry, toron data, and
rank-one source envelopes}
\label{sec:v11v45}

The fifth-version audit is now compulsory.  It is carried out against
the newest numeric Yang--Mills and Navier--Stokes notes visible in the
shared directory.  This audit has two purposes.  First, it prevents the
SM--Einstein programme from silently reproving a companion result.
Second, it separates structural information that survives passage
between programmes from endpoint-specific numerical information that
does not.

\subsection{Audited records}

\begin{definition}[V45 companion record]
\label{def:v11v45-record}
The files inspected for the present audit are
\[
\begin{array}{c|r|l}
\text{file}&\text{bytes}&\text{SHA--256}\\ \hline
\texttt{Renewal\_Yang\_Mills\_Mass\_Gap\_Research\_Notes\_V35.tex}
 &483056&
\texttt{E4AB11D1DCC9EC43955A92008AC9EFA838CA5FABA7BD3F92656762A6AFEE3C30}\\
\texttt{Renewal\_NS\_Stability\_Research\_Notes\_V26.tex}
 &278283&
\texttt{82C887F8C2C69E4F28FAD7C3DC8DE5B9099CB5A4BF431EBA1FFF3E91935978AD}.
\end{array}                                                     \tag{45.1}
\]
The Yang--Mills directory also contains a file with numeric suffix
V34 having the same byte count and digest as the displayed V35 file.
The payload of both ends with the section labelled compact-series V34;
there is no distinct V35 mathematical block in that payload.  We
therefore cite the visible V35 filename but attribute the new
mathematics to its compact V34 block.
\end{definition}

\subsection{The transferable quotient statement}

The Yang--Mills audit supplies a clean reason why a rooted tree used to
write coordinates cannot alter a gauge-invariant scalar response.  The
following finite-dimensional formulation is the part needed here.

\begin{theorem}[Slice-independent finite quotient integral]
\label{thm:v11v45-slice-independent}
Let \(X\) be a compact Hausdorff space, let a compact group \(G\) act
continuously on \(X\), and let \(\mu\) be a finite \(G\)-invariant
Borel measure.  Write \(\pi:X\to X/G\) for the orbit map.  If
\(F=f\circ\pi\) is bounded and gauge invariant, then
\[
              \int_X F\,d\mu=\int_{X/G} f\,d(\pi_*\mu).        \tag{45.2}
\]
Consequently any two measurable gauge-fixing parametrizations that
disintegrate the same measure \(\mu\) over \(\pi_*\mu\) give the same
integral of \(F\).

Suppose in addition that a homeomorphism \(S:X\to X\) preserves
\(\mu\) and maps each \(G\)-orbit onto a \(G\)-orbit.  Then it induces
a quotient homeomorphism \(\bar S:X/G\to X/G\), and
\[
 \int_X F\circ S\,d\mu
 =\int_{X/G} f\circ\bar S\,d(\pi_*\mu).                        \tag{45.3}
\]
In particular, if \(F_\alpha\circ S=F_{\sigma(\alpha)}\) for a finite
index permutation \(\sigma\), the corresponding expectation values
belong to the same \(\sigma\)-orbit, independently of the chosen
gauge coordinates.
\end{theorem}

\begin{proof}
Equation (45.2) is the defining identity for pushforward measure.
A gauge-fixing parametrization is only a representation of the same
quotient integral, so its choice cannot change that identity.

If \(S\) maps orbits to orbits, the rule
\(\bar S(\pi(x))=\pi(Sx)\) is well defined.  It is continuous because
\(\pi\circ S=\bar S\circ\pi\) and \(\pi\) is a quotient map; applying
the same argument to \(S^{-1}\) proves that \(\bar S\) is a
homeomorphism.  Measure preservation gives
\[
 \int_X F\circ S\,d\mu=\int_X F\,d\mu,
\]
while \(F\circ S=(f\circ\bar S)\circ\pi\), which proves (45.3).
The last assertion follows by substituting the covariance relation.
\end{proof}

\begin{corollary}[Axial polarization orbit]
\label{cor:v11v45-polarization}
Let a finite hypercubic gauge system have external momentum \(q e_0\).
If its measure, action, and quotient response construction are
invariant under permutations and sign changes of the three transverse
axes, then its two-index response has the form
\[
       {\cal W}(q e_0)
       =\operatorname{diag}({\cal W}_L,{\cal W}_T,
                            {\cal W}_T,{\cal W}_T).             \tag{45.4}
\]
\end{corollary}

\begin{proof}
Permutations of axes \(1,2,3\) make the three transverse diagonal
entries equal by Theorem~\ref{thm:v11v45-slice-independent}.  Reflecting
one axis changes the sign of every response component carrying that
axis exactly once, while preserving the measure and \(q e_0\).
Every off-diagonal component is therefore equal to its negative.
\end{proof}

This corollary is structural and can be reused in the SM gauge sector.
The companion's exact toron fractions, however, belong to its chosen
endpoint pair and finite parity cell.  They are regression data for
that construction, not universal Standard Model coefficients.

\subsection{What the toron calculation does and does not buy}

The new Yang--Mills block computes the longitudinal and transverse
zero-internal-momentum response through second external order and uses
hypercubic symmetry to close all four polarizations.  It also records
a nonzero longitudinal entry before the internal momentum sum.

\begin{proposition}[Cellwise Ward cancellation is not required]
\label{prop:v11v45-cellwise-ward}
Let a regulated polarization tensor be decomposed into finitely many
internal cells,
\[
                   \Pi_{\mu\nu}(q)=\sum_{k\in{\cal K}}
                                  \Pi_{\mu\nu}(q;k).            \tag{45.5}
\]
The Ward identity
\[
                         q^\mu\Pi_{\mu\nu}(q)=0                \tag{45.6}
\]
implies only
\[
          \sum_{k\in{\cal K}}q^\mu\Pi_{\mu\nu}(q;k)=0.         \tag{45.7}
\]
It does not imply
\(q^\mu\Pi_{\mu\nu}(q;k)=0\) for each \(k\).
\end{proposition}

\begin{proof}
Equation (45.7) follows from (45.5) by linearity.  The converse
cellwise conclusion would require linear independence of the summands
or a separate local identity, neither of which follows from the
vanishing of their sum.  Two nonzero opposite summands already give a
counterexample.
\end{proof}

Thus the SM--Einstein ledger must keep the Ward defect as a
sum-level row until all internal momentum orbits and their boundary
terms have been assembled.  The audited toron response can test the
zero-cell compiler, but it cannot close the gauge contribution to the
continuum stress tensor or the Einstein source.

\subsection{Restricted transfer of the Navier--Stokes rank-one card}

The current NS note computes pointwise rank-one norms for first and
second Oseen-kernel derivatives.  Its first-derivative refinement gives
no improvement; its second-derivative refinement is strictly smaller
only on an intermediate radial range.  The transferable principle is
the following elementary source decomposition.

\begin{lemma}[Positive rank-one source envelope]
\label{lem:v11v45-rank-one-envelope}
Let \(V\) be a finite-dimensional inner-product space, let
\(L:\operatorname{Sym}(V)\to W\) be linear, and set
\[
 c_{r1}(L):=\sup_{\lVert a\rVert=1}
                     \lVert L(a\otimes a)\rVert_W.             \tag{45.8}
\]
For every positive semidefinite \(T\in\operatorname{Sym}(V)\),
\[
                         \lVert L(T)\rVert_W
                    \le c_{r1}(L)\operatorname{tr}T.           \tag{45.9}
\]
For a general symmetric \(T\),
\[
                         \lVert L(T)\rVert_W
                    \le c_{r1}(L)\lVert T\rVert_{S_1},         \tag{45.10}
\]
where \(\lVert T\rVert_{S_1}\) is the sum of the absolute eigenvalues.
\end{lemma}

\begin{proof}
Choose an orthonormal eigenbasis and write
\(T=\sum_j\lambda_j a_j\otimes a_j\).  Linearity and the triangle
inequality give
\[
 \lVert L(T)\rVert_W
 \le\sum_j|\lambda_j|
       \lVert L(a_j\otimes a_j)\rVert_W
 \le c_{r1}(L)\sum_j|\lambda_j|.
\]
This is (45.10).  If \(T\ge0\), the last sum is
\(\operatorname{tr}T\), proving (45.9).
\end{proof}

\begin{transfer}[Use in the Einstein source ledger]
\label{trans:v11v45-rank-one}
When a regulated matter stress is a positive sum of rank-one currents,
Lemma~\ref{lem:v11v45-rank-one-envelope} can replace a full symmetric
operator norm by a rank-one norm, provided the trace or nuclear-norm
moment is carried explicitly.  This transfer is invalid for an
indefinite renormalized stress unless the positive and negative
spectral parts are separately controlled.  The NS numerical Oseen
constants themselves are dimension- and kernel-specific and are not
inserted into the four-dimensional SM--Einstein estimates.
\end{transfer}

\subsection{Audit decision}

\begin{decision}[V46 acquisition]
\label{dec:v11v45-next}
Theorem~\ref{thm:v11v45-slice-independent} identifies the correct
upstream repair for V44.  On a finite connected lattice, spanning-tree
gauge fixing should be proved as an exact global Haar
disintegration, with the root group retained or divided out
explicitly.  V46 will prove that statement, identify the independent
off-tree links, and show that the tree-gauge Faddeev--Popov factor is
constant.  This removes a finite-graph coordinate ambiguity before
returning to probabilistic transverse-Hessian bounds.  Diffeomorphism
and gravitational constraint orbits remain a separate acquisition.
\end{decision}

\begin{firewall}[Audit boundary]
\label{fire:v11v45-boundary}
The audit transfers quotient symmetry, the sum-level interpretation of
Ward cancellation, and a restricted rank-one source envelope.  It does
not import the Yang--Mills toron fractions as Standard Model constants,
does not infer a Ward identity from one internal cell, and does not
import Navier--Stokes regularity.  Neither companion note proves the
SM--Einstein continuum measure, reflection positivity, clustering, or
the semiclassical Einstein equation.
\end{firewall}

\begin{assessment}[Progress after V45]
\label{ass:v11v45-status}
The required companion inspection is complete.  Its main effect is a
change of order: the finite lattice gauge quotient will be closed
globally by exact tree coordinates before any further small-singular-
value estimate is attached to the transverse Hessian.
\end{assessment}

\subsection{Parent record}

This V45 note was formed by appending the preceding material to the
validated V44 source, whose record was
\[
\begin{array}{c|c}
\text{lines}&14060\\
\text{bytes}&412774\\
\text{SHA--256}&
\texttt{5243BC524F0EEAB2C4385983AB6028EE3788E06D89FC11185A892DA377328858}.
\end{array}                                                     \tag{45.11}
\]

% =============================================================================
% END APPENDED ELEVENTH-SERIES V45 MATERIAL
% =============================================================================
% =============================================================================
% BEGIN APPENDED ELEVENTH-SERIES V46 MATERIAL
% =============================================================================

\section{Exact spanning-tree gauge disintegration on a finite lattice}
\label{sec:v11v46}

V44 removed infinitesimal gauge zero modes by an orthogonal transverse
projector.  V45 showed that the finite quotient should first be closed
globally.  This section proves that closure for the internal gauge
group on any finite connected graph.  It also identifies precisely
what remains singular: the based quotient has no gauge copies, whereas
the final residual conjugation quotient can have stabilizer strata.

\subsection{Tree coordinates}

Let \(\Gamma=(V,E)\) be a finite connected graph with one orientation
chosen for each geometric edge.  Let \(G\) be a compact Lie group with
normalized Haar probability \(dg\).  Fix a root \(o\in V\) and a
spanning tree \(T\subset E\).  Reversing a chosen edge orientation
replaces \(U_e\) by \(U_e^{-1}\), which preserves Haar measure, so in
the proof we orient every tree edge away from \(o\).

The link configuration space and based gauge group are
\[
 {\cal A}_\Gamma=G^E,\qquad
 {\cal G}_o=\{k\in G^V:k_o=1\}\simeq G^{V\setminus\{o\}},       \tag{46.1}
\]
with action
\[
                    (k\cdot U)_e=k_{s(e)}U_e k_{t(e)}^{-1}.   \tag{46.2}
\]
For each vertex \(v\), define the tree transport \(h_v(U)\) recursively
by
\[
 h_o(U)=1,\qquad h_v(U)=h_{p(v)}(U)U_{p(v)v},                 \tag{46.3}
\]
where \(p(v)\) is the predecessor of \(v\) in the rooted tree.
For every chord \(e\in C:=E\setminus T\), put
\[
                W_e(U)=h_{s(e)}(U)U_eh_{t(e)}(U)^{-1}.        \tag{46.4}
\]

\begin{theorem}[Global based tree gauge]
\label{thm:v11v46-global-tree}
The map
\[
 \Theta_T:{\cal A}_\Gamma\longrightarrow
           G^{V\setminus\{o\}}\times G^C,\qquad
 U\longmapsto\bigl((h_v(U))_{v\ne o},(W_e(U))_{e\in C}\bigr)  \tag{46.5}
\]
is a diffeomorphism.  Its inverse is
\[
 U_{p(v)v}=h_{p(v)}^{-1}h_v,\qquad
 U_e=h_{s(e)}^{-1}W_eh_{t(e)}\quad(e\in C).                   \tag{46.6}
\]
For \(k\in{\cal G}_o\),
\[
 h_v(k\cdot U)=h_v(U)k_v^{-1},\qquad
 W_e(k\cdot U)=W_e(U).                                       \tag{46.7}
\]
Consequently the action of \({\cal G}_o\) is free, every based orbit
meets the tree slice
\[
                 {\cal S}_T:=\{U:U_e=1\text{ for }e\in T\}    \tag{46.8}
\]
exactly once, and
\[
                    {\cal A}_\Gamma/{\cal G}_o\simeq G^C.     \tag{46.9}
\]
\end{theorem}

\begin{proof}
Recursion (46.3) uniquely determines all \(h_v\).  Equations (46.6)
recover every tree link and every chord link, so they give a two-sided
smooth inverse to \(\Theta_T\).

For (46.7), induction from the root gives
\[
 h_v(k\cdot U)
 =h_{p(v)}(U)k_{p(v)}^{-1}
   k_{p(v)}U_{p(v)v}k_v^{-1}
 =h_v(U)k_v^{-1}.
\]
Substitution in (46.4) cancels \(k_{s(e)}\) and \(k_{t(e)}\), proving
that every \(W_e\) is based-gauge invariant.

Choose \(k_v=h_v(U)\).  Then the transformed tree link is
\[
 h_{p(v)}U_{p(v)v}h_v^{-1}=1,
\]
so every orbit meets \({\cal S}_T\).  If \(U\in{\cal S}_T\) and
\(k\cdot U\in{\cal S}_T\), recursion along the tree gives
\(k_v=k_{p(v)}\), hence \(k_v=1\) for all \(v\).  The intersection is
unique and the action is free.  On the slice the chord links are
exactly the coordinates \(W_e\), proving (46.9).
\end{proof}

\subsection{Exact Haar Jacobian}

\begin{theorem}[Product-Haar disintegration]
\label{thm:v11v46-haar}
Under \(\Theta_T\), normalized product Haar measure satisfies
\[
             \prod_{e\in E}dU_e
       =\left(\prod_{v\ne o}dh_v\right)
        \left(\prod_{e\in C}dW_e\right)                       \tag{46.10}
\]
as an identity of measures.  Hence for every integrable
based-gauge-invariant \(F(U)=f(W(U))\),
\[
                 \int_{G^E}F(U)\prod_{e\in E}dU_e
                  =\int_{G^C}f(W)\prod_{e\in C}dW_e.          \tag{46.11}
\]
Thus the spanning-tree Faddeev--Popov factor is the constant \(1\)
for normalized Haar measure.
\end{theorem}

\begin{proof}
Order the nonroot vertices so that predecessors occur before
descendants.  The triangular change
\[
                 U_{p(v)v}=h_{p(v)}^{-1}h_v
\]
preserves product Haar measure inductively: with all earlier variables
fixed, left multiplication by \(h_{p(v)}^{-1}\) sends \(dh_v\) to
\(dU_{p(v)v}\).  Conditional on all \(h_v\), each chord change
\[
                 U_e=h_{s(e)}^{-1}W_eh_{t(e)}
\]
also preserves Haar measure by left and right invariance.  Fubini's
theorem proves (46.10).  If \(F=f(W)\), integration over the
\(|V|-1\) normalized \(h\)-variables gives one and yields (46.11).
The last statement is precisely the measure Jacobian assertion.
\end{proof}

\begin{corollary}[Unnormalized convention]
\label{cor:v11v46-unnormalized}
If Haar measure has total mass \(\operatorname{vol}G\), the eliminated
based gauge volume is
\[
                         (\operatorname{vol}G)^{|V|-1}.        \tag{46.12}
\]
It is configuration independent and cancels from normalized
expectations.
\end{corollary}

\begin{proof}
Repeat Theorem~\ref{thm:v11v46-haar} without normalizing the
\(|V|-1\) gauge-coordinate integrals.
\end{proof}

This exact measure statement is stronger than a local determinant
calculation.  It remains valid at configurations with nontrivial
holonomy because it uses the tree recursion globally.

\subsection{Residual root symmetry and the genuine quotient}

The based quotient is not yet the quotient by the full gauge group.
Let \(g\in G^V\), set \(a=g_o\), and apply the same tree construction
to \(g\cdot U\).  Direct induction gives
\[
        h_v(g\cdot U)=a\,h_v(U)g_v^{-1},\qquad
        W_e(g\cdot U)=aW_e(U)a^{-1}.                           \tag{46.13}
\]

\begin{corollary}[Full finite gauge quotient]
\label{cor:v11v46-full-quotient}
There is a natural homeomorphism
\[
                   G^E/G^V\simeq G^C/\operatorname{Ad}G.      \tag{46.14}
\]
A fully gauge-invariant integrand becomes a simultaneous-conjugation
invariant function of the chord variables.
\end{corollary}

\begin{proof}
Theorem~\ref{thm:v11v46-global-tree} first divides by
\({\cal G}_o\).  Equation (46.13) shows that the remaining root value
acts on every chord by simultaneous conjugation.  Taking the second
quotient proves (46.14).  Compactness makes the continuous bijection
between the compact quotient and the Hausdorff orbit space a
homeomorphism.
\end{proof}

\begin{remark}[Where quotient singularities remain]
The action of \({\cal G}_o\) is free, so the based tree quotient has no
Gribov copies and is the smooth manifold \(G^C\).  Simultaneous
conjugation need not be free: central or commuting chord tuples have
larger stabilizers.  The stratification of
\(G^C/\operatorname{Ad}G\) is a genuine residual quotient feature, not
a failure of tree gauge.
\end{remark}

\subsection{Wilson words and independence of the tree}

\begin{proposition}[Chord generators]
\label{prop:v11v46-chord-generators}
For a chord \(e\), \(W_e\) is the holonomy of the fundamental loop
formed by the root-to-\(s(e)\) tree path, the chord \(e\), and the
reverse root-to-\(t(e)\) tree path.  Every closed-loop holonomy is,
up to conjugation at its base point, a word in the \(W_e^{\pm1}\).
\end{proposition}

\begin{proof}
The first assertion is (46.4) read as a path product.  Gauge the tree
links to the identity.  A closed path then contributes the identity on
each tree segment and a factor \(W_e\) or \(W_e^{-1}\) whenever it
crosses a chord according to or against its orientation.  Restoring
the original gauge conjugates the path product at its base point.
\end{proof}

\begin{corollary}[Tree independence of observables]
\label{cor:v11v46-tree-independent}
Let \(T\) and \(T'\) be two spanning trees.  The induced coordinate
map \(G^{E\setminus T}\to G^{E\setminus T'}\) preserves the quotient
probability measure.  Every integrable gauge-invariant Wilson
observable has the same expectation in the two coordinate systems.
\end{corollary}

\begin{proof}
Both coordinate measures are the pushforward of the same product Haar
measure by Theorem~\ref{thm:v11v46-haar}.  Apply
Theorem~\ref{thm:v11v45-slice-independent}; the Wilson functions agree
because both are evaluations of the same link-space observable.
\end{proof}

\subsection{Matter fields and the anomaly boundary}

\begin{corollary}[Invariant site matter]
\label{cor:v11v46-matter}
Let a finite-dimensional unitary representation of \(G\) act on a
site field \(\phi_v\), and suppose the site measure \(d\nu(\phi_v)\)
is invariant.  With
\[
                         \Phi_v=h_v(U)\phi_v,                  \tag{46.15}
\]
the combined link--matter measure disintegrates into normalized based
gauge volume times the tree-gauge variables \((W,\Phi)\).
Every integrable based-gauge-invariant bosonic observable can therefore
be integrated on the global tree slice with constant gauge factor.
\end{corollary}

\begin{proof}
For fixed \(h_v\), unitarity and invariance give
\(d\nu(\phi_v)=d\nu(\Phi_v)\).  Combine this sitewise identity with
(46.10).  Gauge invariance removes the \(h\)-variables exactly as in
(46.11).
\end{proof}

\begin{firewall}[Chiral measure]
\label{fire:v11v46-chiral}
Corollary~\ref{cor:v11v46-matter} applies directly to invariant
finite-dimensional bosonic measures and to vectorlike regulated
fermion measures with an invariant determinant.  A chiral fermion
phase is not declared invariant by this coordinate calculation.  Its
Jacobian and determinant line are governed by the anomaly and phase
ledgers already isolated earlier in this cycle.  The exact Standard
Model anomaly arithmetic removes the local perturbative obstruction;
the nonperturbative phase denominator and global compatibility remain
open.
\end{firewall}

\subsection{Consequence for the transverse Hessian programme}

Let \(S(U)=s(W(U))\) be a smooth based-gauge-invariant action.  In the
global coordinates \((h,W)\), it is independent of \(h\).  Therefore
at every critical point its coordinate Hessian has the block form
\[
             D^2(S\circ\Theta_T^{-1})
              =\begin{pmatrix}0&0\\0&D_W^2s\end{pmatrix}.      \tag{46.16}
\]

\begin{proposition}[Gauge kernel removed before probability]
\label{prop:v11v46-kernel-removed}
For the based finite lattice, all \((|V|-1)\dim G\) infinitesimal gauge
directions are removed exactly by tree coordinates.  Any small
singular value of \(D_W^2s\) is therefore a quotient-Hessian
phenomenon rather than an unfixed based-gauge zero mode.  No random
inverse moment of the based orbit Gram matrix is required.
\end{proposition}

\begin{proof}
Theorem~\ref{thm:v11v46-global-tree} identifies the based quotient
globally with \(G^C\).  Equation (46.16) separates the gauge-coordinate
zero block from the Hessian of the quotient action.  The orbit Gram
operator still controls comparison with the orthogonal slice of V44,
but its inverse has the deterministic graph bound proved there.
Thus it does not enter the probabilistic small-ball row.
\end{proof}

\begin{remark}[Coordinate determinants]
The determinant of \(D_W^2s\) in tree coordinates and the
pseudodeterminant of the ambient orthogonal Hessian need not be
numerically identical: their bases carry different induced metrics.
They are related by the square of the nonsingular tangent-coordinate
Jacobian at a critical point.  The Haar disintegration proves that the
measure Faddeev--Popov factor is constant; it does not authorize
dropping this metric Jacobian from an independently defined Gaussian
stationary-phase determinant.
\end{remark}

\begin{decision}[Next acquisition]
\label{dec:v11v46-next}
The internal based gauge quotient is now globally closed at every
finite lattice level.  The next upstream obstruction is the Einstein
constraint sector.  V47 will formulate a finite-dimensional regular
constraint theorem with an explicit Gram determinant, prove its
coarea disintegration and orthogonal tangent Hessian, and mark exactly
where gravitational constraint rank or noncompact orbit volume can
fail.  This will keep the gravitational row logically separate from
the compact tree-gauge result.
\end{decision}

\begin{firewall}[Finite-lattice boundary]
\label{fire:v11v46-boundary}
Tree gauge proves a global finite-graph statement for compact internal
gauge groups.  It does not prove uniform continuum bounds as the graph
grows, does not remove residual conjugation strata, and does not apply
unchanged to the noncompact diffeomorphism group or the Hamiltonian
constraint.  Reflection positivity, clustering, the chiral phase
limit, and the Einstein dynamics remain open.
\end{firewall}

\begin{assessment}[Progress after V46]
\label{ass:v11v46-status}
The based finite gauge quotient is no longer a local-slice
assumption.  It is the exact product-Haar quotient \(G^C\), with one
independent chord variable per fundamental cycle and constant
tree-gauge factor.  Probabilistic Hessian estimates may now be directed
only at the genuine quotient Hessian.
\end{assessment}

\subsection{Parent record}

This V46 note was formed by appending the preceding material to the
validated V45 source, whose record was
\[
\begin{array}{c|c}
\text{lines}&14321\\
\text{bytes}&423525\\
\text{SHA--256}&
\texttt{07547B7CF37E4E0C87F35F51DFD1E4937A73BD781DB7F3F67EB51A11B8BBD3F9}.
\end{array}                                                     \tag{46.17}
\]

% =============================================================================
% END APPENDED ELEVENTH-SERIES V46 MATERIAL
% =============================================================================
% =============================================================================
% BEGIN APPENDED ELEVENTH-SERIES V47 MATERIAL
% =============================================================================

\section{Regular constraint surfaces, coarea measure, and the constrained
Hessian}
\label{sec:v11v47}

The compact internal gauge sector has an exact global finite-graph
slice.  The Einstein constraint sector is different.  Before choosing
lapse, shift, or diffeomorphism gauges, one must first distinguish
three operations: imposing a regular constraint map, restricting the
action to its level set, and quotienting any first-class flows tangent
to that level set.  This section proves the finite-dimensional
statements needed for the first two operations.

\subsection{The regular level-set card}

Let \((M,g)\) be an \(n\)-dimensional Riemannian manifold and let
\[
                         F=(F^1,\ldots,F^m):M\to\mathbb R^m   \tag{47.1}
\]
be \(C^2\), with \(m\le n\).  Write
\[
 A_x:=D F_x:T_xM\to\mathbb R^m,\qquad
 G_x:=A_xA_x^*:\mathbb R^m\to\mathbb R^m,                     \tag{47.2}
\]
and
\[
                           J_F(x):=\sqrt{\det G_x}.            \tag{47.3}
\]
The adjoint uses \(g_x\) and the Euclidean metric.

\begin{theorem}[Regular constraint geometry]
\label{thm:v11v47-regular}
Assume \(0\) is a regular value of \(F\), and put
\(C=F^{-1}(0)\).  Then \(C\) is an embedded submanifold of codimension
\(m\),
\[
          T_xC=\ker A_x,\qquad N_xC=\operatorname{ran}A_x^*.  \tag{47.4}
\]
Moreover \(G_x\) is positive definite on \(C\), and the orthogonal
normal and tangent projectors are
\[
 P_x=A_x^*G_x^{-1}A_x,\qquad Q_x=I-P_x.                       \tag{47.5}
\]
\end{theorem}

\begin{proof}
Regularity says that \(A_x\) is surjective on \(C\), so the regular
value theorem gives the embedded submanifold and
\(T_xC=\ker A_x\).  In finite dimension,
\((\ker A_x)^\perp=\operatorname{ran}A_x^*\), proving the normal
identity.

For \(z\in\mathbb R^m\),
\[
                  \langle G_xz,z\rangle=\lVert A_x^*z\rVert^2.
\]
Surjectivity of \(A_x\) is equivalent to injectivity of \(A_x^*\);
hence \(G_x>0\).  Direct multiplication gives
\(P_x^*=P_x\), \(P_x^2=P_x\), and
\(\operatorname{ran}P_x=\operatorname{ran}A_x^*\).
Thus \(P_x\) is the orthogonal normal projector and \(Q_x\) the tangent
projector.
\end{proof}

\begin{corollary}[Quantitative normal right inverse]
\label{cor:v11v47-right-inverse}
If \(\lambda_{\min}(G_x)\ge\gamma^2>0\), then
\[
             R_x:=A_x^*G_x^{-1},\qquad A_xR_x=I,\qquad
             \lVert R_x\rVert\le\gamma^{-1}.                 \tag{47.6}
\]
\end{corollary}

\begin{proof}
The right-inverse identity is immediate.  Its squared operator norm is
\[
 \lVert R_x\rVert^2
 =\lVert G_x^{-1}A_xA_x^*G_x^{-1}\rVert
 =\lVert G_x^{-1}\rVert\le\gamma^{-2}.
\]
\end{proof}

\subsection{The exact coarea normalization}

Let \(dV_M\) be Riemannian volume and \(d\sigma_y\) the induced volume
on \(F^{-1}(y)\).

\begin{theorem}[Coarea constraint identities]
\label{thm:v11v47-coarea}
For every nonnegative measurable \(a\),
\[
 \int_M a(x)J_F(x)\,dV_M(x)
 =\int_{\mathbb R^m}
      \left(\int_{F^{-1}(y)}a(x)\,d\sigma_y(x)\right)dy.       \tag{47.7}
\]
Consequently the two formal delta conventions mean
\[
 \int_M a(x)\delta_0(F(x))J_F(x)\,dV_M(x)
       =\int_Ca\,d\sigma_0,                                   \tag{47.8}
\]
and
\[
 \int_M a(x)\delta_0(F(x))\,dV_M(x)
       =\int_C\frac{a}{J_F}\,d\sigma_0,                        \tag{47.9}
\]
whenever the displayed right-hand sides are integrable.
\end{theorem}

\begin{proof}
Equation (47.7) is the coarea theorem applied to \(F\).  To make the
delta notation precise, let
\(\rho\in C_c(\mathbb R^m)\) be nonnegative with integral one and set
\(\rho_\epsilon(y)=\epsilon^{-m}\rho(y/\epsilon)\).  Applying (47.7)
to \(a(x)\rho_\epsilon(F(x))\) gives
\[
 \int_M a\rho_\epsilon(F)J_F\,dV_M
 =\int_{\mathbb R^m}\rho_\epsilon(y)
       \int_{F^{-1}(y)}a\,d\sigma_y\,dy.                       \tag{47.10}
\]
On a compact set on which \(F\) is a submersion, local submersion
charts and a partition of unity make the inner fibre integral
continuous at \(y=0\).  Approximate-identity convergence gives (47.8).
Replacing \(a\) in (47.8) by \(a/J_F\), which is defined near the
regular level, gives (47.9).  General integrable cases follow by
localization and monotone or dominated convergence.
\end{proof}

\begin{warning}[The determinant cannot be chosen silently]
\label{warn:v11v47-convention}
The induced geometric measure on \(C\) is represented in ambient
variables by \(\delta_0(F)J_F\,dV_M\).  The bare constraint density
\(\delta_0(F)dV_M\) instead induces \(J_F^{-1}d\sigma_0\).
These define different measures unless \(J_F\) is constant.  A
finite-cutoff Einstein construction must state which measure its
canonical or covariant derivation produces.
\end{warning}

\subsection{Stability of the normal projector}

The next estimate is the constraint analogue of the orbit-projector
comparison in V44.

\begin{proposition}[Constraint-projector comparison]
\label{prop:v11v47-projector}
Let \(A,\widetilde A:H\to\mathbb R^m\) be surjective linear maps on a
finite-dimensional Hilbert space.  Suppose
\[
 \lVert A\rVert,\lVert\widetilde A\rVert\le M,\qquad
 \lambda_{\min}(AA^*),\lambda_{\min}(\widetilde A\widetilde A^*)
       \ge\gamma^2.                                           \tag{47.11}
\]
For
\(P=A^*(AA^*)^{-1}A\) and
\(\widetilde P=\widetilde A^*(\widetilde A\widetilde A^*)^{-1}
\widetilde A\),
\[
 \lVert P-\widetilde P\rVert
 \le\left(2M\gamma^{-2}+2M^3\gamma^{-4}\right)
                     \lVert A-\widetilde A\rVert.             \tag{47.12}
\]
The same bound holds for the tangent projectors \(Q=I-P\) and
\(\widetilde Q=I-\widetilde P\).
\end{proposition}

\begin{proof}
Put \(G=AA^*\) and \(\widetilde G=\widetilde A\widetilde A^*\).
Then
\[
 \lVert G-\widetilde G\rVert
 \le(\lVert A\rVert+\lVert\widetilde A\rVert)
                    \lVert A-\widetilde A\rVert
 \le2M\lVert A-\widetilde A\rVert,                            \tag{47.13}
\]
and the resolvent identity gives
\[
 \lVert G^{-1}-\widetilde G^{-1}\rVert
 \le2M\gamma^{-4}\lVert A-\widetilde A\rVert.                 \tag{47.14}
\]
Insert and subtract
\(\widetilde A^*G^{-1}A\) and
\(\widetilde A^*\widetilde G^{-1}A\) between \(P\) and
\(\widetilde P\).  The two exterior differences contribute
\(2M\gamma^{-2}\lVert A-\widetilde A\rVert\), and (47.14) contributes
\(2M^3\gamma^{-4}\lVert A-\widetilde A\rVert\).
Since \(Q-\widetilde Q=-(P-\widetilde P)\), the last assertion follows.
\end{proof}

\subsection{The correct constrained Hessian}

Let \(S\in C^2(M)\), and suppose \(x\in C\) is a critical point of
\(S|_C\).  By (47.4), there is a unique multiplier
\(\lambda=(\lambda_1,\ldots,\lambda_m)\) such that
\[
                         \nabla S(x)=A_x^*\lambda.             \tag{47.15}
\]
Define the Lagrangian
\[
                        {\cal L}=S-\sum_{a=1}^m\lambda_aF^a.   \tag{47.16}
\]

\begin{theorem}[Intrinsic Hessian at a constrained critical point]
\label{thm:v11v47-constrained-hessian}
For \(u,v\in T_xC\),
\[
             \operatorname{Hess}_C(S|_C)_x(u,v)
             =\operatorname{Hess}_M{\cal L}_x(u,v).            \tag{47.17}
\]
Equivalently, as a self-adjoint operator on \(T_xC\),
\[
 H_C
 =\left.Q_x\left(\operatorname{Hess}S_x
       -\sum_{a=1}^m\lambda_a\operatorname{Hess}F^a_x\right)
                         Q_x\right|_{T_xC}.                    \tag{47.18}
\]
\end{theorem}

\begin{proof}
Choose a smooth curve \(c(t)\subset C\) with \(c(0)=x\) and
\(\dot c(0)=u\).  Since \(F^a(c(t))=0\),
\[
 0=\frac{d^2}{dt^2}F^a(c(t))\big|_{0}
   =\operatorname{Hess}F^a_x(u,u)
      +\langle\nabla F^a(x),\nabla_{\dot c}\dot c(0)\rangle.  \tag{47.19}
\]
Likewise,
\[
 \frac{d^2}{dt^2}S(c(t))\big|_0
 =\operatorname{Hess}S_x(u,u)
   +\langle\nabla S(x),\nabla_{\dot c}\dot c(0)\rangle.        \tag{47.20}
\]
Insert (47.15) into the last term and use (47.19).  This gives
\[
 \frac{d^2}{dt^2}S(c(t))\big|_0
 =\operatorname{Hess}S_x(u,u)
       -\sum_a\lambda_a\operatorname{Hess}F^a_x(u,u).         \tag{47.21}
\]
Polarization proves the bilinear identity (47.17).  Since \(Q_x\) is
the identity on \(T_xC\), (47.18) follows.
\end{proof}

\begin{corollary}[Nondegeneracy on the constrained tangent]
\label{cor:v11v47-nondegenerate}
The stationary point \(x\) is nondegenerate as a critical point of
\(S|_C\) exactly when the operator \(H_C\) in (47.18) is invertible.
Ambient invertibility of \(\operatorname{Hess}S_x\) is neither
necessary nor sufficient for this condition.
\end{corollary}

\begin{proof}
Nondegeneracy on a manifold is, by definition, nonsingularity of its
intrinsic Hessian.  Theorem~\ref{thm:v11v47-constrained-hessian}
identifies that Hessian with \(H_C\).  The final statement follows
because normal blocks and the second fundamental form terms can change
ambient invertibility without changing, or while changing, the
tangent restriction.
\end{proof}

\subsection{What this says for an Einstein cutoff}

At a finite regulator, let \(q\) denote metric and matter variables and
let
\[
             {\cal C}(q)=({\cal H}(q),{\cal H}_1(q),
                    \ldots,{\cal H}_d(q))                     \tag{47.22}
\]
denote a finite list of Hamiltonian and momentum constraints.  If
\(D{\cal C}\) is surjective on the relevant zero set, the preceding
theorems apply with \(F={\cal C}\).

\begin{ledger}[Regular Einstein-constraint row]
\label{led:v11v47-einstein}
A finite-cutoff proof must provide, uniformly on its good sets:
\begin{enumerate}
\item a rank bound
\(\lambda_{\min}(D{\cal C}D{\cal C}^*)\ge\gamma_i^2\);
\item a declared choice between induced constraint measure and bare
delta measure, with the corresponding factor \(J_{\cal C}\) or
\(J_{\cal C}^{-1}\);
\item the multiplier vector solving
\(\nabla S=D{\cal C}^*\lambda\);
\item inverse and determinant bounds for the corrected tangent Hessian
\[
 Q\left(\operatorname{Hess}S
       -\sum_a\lambda_a\operatorname{Hess}{\cal C}^a\right)Q;
                                                                  \tag{47.23}
\]
\item a separate quotient or gauge-fixing argument for first-class
flows tangent to \({\cal C}^{-1}(0)\).
\end{enumerate}
\end{ledger}

\begin{firewall}[Regularity is not gravitational reduction]
\label{fire:v11v47-first-class}
Theorem~\ref{thm:v11v47-regular} imposes a regular level set.  It does
not divide the Hamiltonian or momentum gauge flows, prove closure of
their constraint algebra, or control a noncompact diffeomorphism
volume.  If the constraint rank drops, \(J_{\cal C}=0\), the projector
(47.5) ceases to exist, and the regular card cannot be used across
that stratum.  These are physical acquisition requirements, not
coordinate details.
\end{firewall}

\begin{decision}[Next acquisition]
\label{dec:v11v47-next}
V48 will add the missing finite-dimensional first-class reduction
card.  Given constraints and gauge conditions, it will identify the
Dirac/Faddeev--Popov matrix, prove local uniqueness by the inverse
function theorem, derive the determinant in the gauge-fixed coarea
formula, and distinguish compact finite orbit volume from noncompact
gravitational gauge volume.
\end{decision}

\begin{assessment}[Progress after V47]
\label{ass:v11v47-status}
The constraint surface, its exact coarea normalization, its tangent
projector, and its intrinsic stationary Hessian are now fixed without
gauge-fixing ambiguity.  The result exposes the next independent
obstruction: regular first-class orbit reduction and a uniform lower
bound for its gauge-fixing matrix.
\end{assessment}

\subsection{Parent record}

This V47 note was formed by appending the preceding material to the
validated V46 source, whose record was
\[
\begin{array}{c|c}
\text{lines}&14670\\
\text{bytes}&437132\\
\text{SHA--256}&
\texttt{84FB6581694C418E8EC43A09A8E2296C0928E6799C2CF757CB621283F140EAAF}.
\end{array}                                                     \tag{47.24}
\]

% =============================================================================
% END APPENDED ELEVENTH-SERIES V47 MATERIAL
% =============================================================================
% =============================================================================
% BEGIN APPENDED ELEVENTH-SERIES V48 MATERIAL
% =============================================================================

\section{First-class orbit reduction, Faddeev--Popov density, and the
Dirac matrix}
\label{sec:v11v48}

V47 imposed a regular constraint surface but deliberately did not
divide flows tangent to it.  The present section supplies the finite-
dimensional local quotient theorem.  The result exposes three separate
global requirements for gravity: a uniformly transverse gauge
condition, control of multiple intersections, and a quotient
normalization when the gauge group is noncompact.

\subsection{The transverse-slice criterion}

Let \(C\) be a smooth \(d\)-manifold carrying a smooth action of an
\(r\)-dimensional Lie group \(H\).  For a fixed basis
\((e_b)_{b=1}^r\) of its Lie algebra, write
\[
 R_x:\mathfrak h\to T_xC,\qquad
 R_xe_b=\left.\frac{d}{dt}\right|_{0}\exp(te_b)\cdot x.        \tag{48.1}
\]
Let \(\chi=(\chi^1,\ldots,\chi^r):C\to\mathbb R^r\) be a gauge
condition, \(\Sigma=\chi^{-1}(0)\), and define
\[
                         M_x:=D\chi_xR_x.                      \tag{48.2}
\]

\begin{theorem}[Local gauge slice]
\label{thm:v11v48-local-slice}
Suppose \(x\in\Sigma\), \(D\chi_x\) is surjective, and \(M_x\) is
invertible.  Then
\[
                   T_xC=T_x\Sigma\oplus\operatorname{ran}R_x. \tag{48.3}
\]
There are neighborhoods \(U\subset\mathfrak h\) of \(0\),
\(V\subset\Sigma\) of \(x\), and \(O\subset C\) of \(x\) for which
\[
                    \Phi:U\times V\to O,\qquad
                    (\xi,s)\mapsto\exp(\xi)\cdot s            \tag{48.4}
\]
is a diffeomorphism.  In particular, every orbit segment in \(O\)
meets \(V\) once.
\end{theorem}

\begin{proof}
The tangent space of the slice is
\(T_x\Sigma=\ker D\chi_x\).  If \(R_x\xi\in T_x\Sigma\), then
\(M_x\xi=0\), so \(\xi=0\).  Hence the two spaces in (48.3) intersect
trivially.  Their dimensions are \(d-r\) and \(r\), proving the direct
sum.  The differential of (48.4) at \((0,x)\) is
\[
                         (\xi,v)\longmapsto R_x\xi+v,          \tag{48.5}
\]
which is an isomorphism by (48.3).  The inverse function theorem gives
the claimed neighborhoods and diffeomorphism.
\end{proof}

\begin{corollary}[Quantitative infinitesimal transversality]
\label{cor:v11v48-transversality}
If \(s_{\min}(M_x)\ge\kappa>0\), then the gauge parameter required to
correct a linearized gauge-condition error \(b\in\mathbb R^r\) obeys
\[
                           \lVert M_x^{-1}b\rVert
                           \le\kappa^{-1}\lVert b\rVert.       \tag{48.6}
\]
\end{corollary}

\begin{proof}
This is the definition of the least singular value.
\end{proof}

\subsection{The finite-copy identity}

Assume now that \(C\) is Riemannian, \(H\) acts by measure-preserving
diffeomorphisms, and \(d\mu_C\) is an invariant measure.  Fix a Haar
volume form \(dh\) on \(H\), with total volume
\(\operatorname{vol}H\) when finite.  The determinant below is computed
relative to this Haar volume and Lebesgue volume on \(\mathbb R^r\).

For \(x\in C\), let
\[
 Z_x=\{h\in H:\chi(h\cdot x)=0\}.                              \tag{48.7}
\]

\begin{theorem}[Faddeev--Popov identity with copy count]
\label{thm:v11v48-fp}
Let \(\Omega\subset C\) be invariant.  Suppose that for every
\(x\in\Omega\), \(Z_x\) consists of exactly \(N\) points and every
zero is simple:
\[
                      \det M_{h\cdot x}\ne0\quad(h\in Z_x).    \tag{48.8}
\]
If \(f\) is an integrable \(H\)-invariant function supported in
\(\Omega\), then, when \(\operatorname{vol}H<\infty\),
\[
 \frac1{\operatorname{vol}H}\int_C f\,d\mu_C
 =\frac1N\int_C f(x)\delta_0(\chi(x))
                  |\det M_x|\,d\mu_C(x).                      \tag{48.9}
\]
If the gauge condition has one intersection per orbit, then \(N=1\).
\end{theorem}

\begin{proof}
For fixed \(x\), apply the change-of-variables formula for a delta
distribution to
\[
                  \theta_x:H\to\mathbb R^r,\qquad
                  \theta_x(h)=\chi(h\cdot x).
\]
At a zero \(h\), its derivative in the right-translated Lie algebra
frame is \(M_{h\cdot x}\).  Therefore
\[
 \int_H\delta_0(\chi(h\cdot x))|\det M_{h\cdot x}|\,dh=N.      \tag{48.10}
\]
Multiply by \(f(x)\), integrate in \(x\), and use Fubini.  Invariance
of \(d\mu_C\) and of \(f\) permits the substitution \(y=h\cdot x\);
the inner integral becomes independent of \(h\).  Thus
\[
 N\int_C f\,d\mu_C
 =(\operatorname{vol}H)\int_C f(y)\delta_0(\chi(y))
                         |\det M_y|\,d\mu_C(y),
\]
which is (48.9).
\end{proof}

\begin{corollary}[Measure induced on the gauge slice]
\label{cor:v11v48-slice-measure}
Assume \(d\mu_C=a\,d\sigma_C\), where \(d\sigma_C\) is Riemannian
volume and \(a\) is invariant.  Put
\[
 J_\chi^C(x)=
 \sqrt{\det\!\left(D\chi_x|_{T_xC}
             (D\chi_x|_{T_xC})^*\right)}.                     \tag{48.11}
\]
Under the hypotheses of Theorem~\ref{thm:v11v48-fp},
\[
 \frac1{\operatorname{vol}H}\int_C f\,a\,d\sigma_C
 =\frac1N\int_\Sigma
      f\,a\,\frac{|\det M|}{J_\chi^C}\,d\sigma_\Sigma.         \tag{48.12}
\]
\end{corollary}

\begin{proof}
Apply the coarea delta identity
Theorem~\ref{thm:v11v47-coarea} on \(C\) to the right-hand side of
(48.9).
\end{proof}

The quotient density is therefore the ratio
\(|\det M|/J_\chi^C\) in induced slice volume.  Writing only
\(|\det M|\) is correct in the ambient delta representation (48.9),
but omits the coarea denominator if one has already moved to
\(d\sigma_\Sigma\).

\subsection{First-class constraints and the Dirac determinant}

Let \((P,\omega)\) be a \(2n\)-dimensional symplectic manifold with
Poisson bracket \(\{\cdot,\cdot\}\).  Let
\({\cal C}_a\), \(1\le a\le r\), be regular constraints satisfying
\[
                  \{{\cal C}_a,{\cal C}_b\}
                    =f_{ab}^{\ \ c}{\cal C}_c.                \tag{48.13}
\]
Their Hamiltonian vector fields are tangent to
\(C={\cal C}^{-1}(0)\).  For gauge conditions \(\chi^a\), define
\[
                  M^a_{\ b}=\{\chi^a,{\cal C}_b\}.             \tag{48.14}
\]

\begin{theorem}[Dirac matrix criterion]
\label{thm:v11v48-dirac}
On the joint surface \({\cal C}=0=\chi\), the \(2r\times2r\)
Poisson matrix of
\(\Psi=({\cal C}_1,\ldots,{\cal C}_r,\chi^1,\ldots,\chi^r)\)
has block form
\[
 \Delta=\{\Psi_A,\Psi_B\}
 =\begin{pmatrix}
       0&-M^{\mathsf T}\\
       M&K
   \end{pmatrix},\qquad K^{ab}=\{\chi^a,\chi^b\}.              \tag{48.15}
\]
It is invertible exactly when \(M\) is invertible, and
\[
                              \det\Delta=(\det M)^2.           \tag{48.16}
\]
Thus the Faddeev--Popov factor is the positive square root of the
Dirac determinant:
\[
                              |\det M|=\sqrt{\det\Delta}.      \tag{48.17}
\]
\end{theorem}

\begin{proof}
Closure (48.13) makes the upper-left block vanish on \(C\), and
antisymmetry gives the other displayed blocks.  If \(M\) is
invertible, the equations
\[
 -M^{\mathsf T}v=p,\qquad Mu+Kv=q
\]
have a unique solution \((u,v)\) for every \((p,q)\), so \(\Delta\)
is invertible.  For invertible \(M\), block row operations of determinant one remove
\(K\):
\[
 \begin{pmatrix}I&0\\ KM^{-\mathsf T}&I\end{pmatrix}
 \begin{pmatrix}0&-M^{\mathsf T}\\M&K\end{pmatrix}
 =\begin{pmatrix}0&-M^{\mathsf T}\\M&0\end{pmatrix}.           \tag{48.18}
\]
The determinant of the last matrix is \((\det M)^2\).
Both sides are polynomials in the entries of \(M,K\), so (48.16)
holds also for singular \(M\); then it proves noninvertibility.
Equation (48.17) follows.
\end{proof}



\begin{corollary}[The gauge condition fixes first-class directions]
\label{cor:v11v48-first-class}
At a joint zero, \(M\) is precisely the matrix \(D\chi\,R\) of
Theorem~\ref{thm:v11v48-local-slice}, with \(R e_b=X_{{\cal C}_b}\).
Hence \(\det\Delta\ne0\) is equivalent to local transverse gauge
fixing.
\end{corollary}

\begin{proof}
By the definition of a Hamiltonian vector field,
\(D\chi^a(X_{{\cal C}_b})=\{\chi^a,{\cal C}_b\}\).
Apply Theorems~\ref{thm:v11v48-local-slice} and
\ref{thm:v11v48-dirac}.
\end{proof}

\subsection{Stability and the next probabilistic row}

\begin{proposition}[Stable inverse and determinant]
\label{prop:v11v48-stable}
Let \(M\) be invertible with \(s_{\min}(M)\ge\kappa\), and let
\(\widetilde M=M+E\) with \(\lVert E\rVert<\kappa\).  Then
\[
 s_{\min}(\widetilde M)\ge\kappa-\lVert E\rVert,\qquad
 \lVert\widetilde M^{-1}-M^{-1}\rVert
 \le\frac{\lVert E\rVert}
          {\kappa(\kappa-\lVert E\rVert)}.                    \tag{48.19}
\]
If \(r=\dim H\), then
\[
 \left|\log|\det\widetilde M|-\log|\det M|\right|
 \le\frac{r\,\lVert E\rVert}
          {\kappa-\lVert E\rVert}.                            \tag{48.20}
\]
\end{proposition}

\begin{proof}
The singular-value inequality gives the first bound.  The resolvent
identity
\(\widetilde M^{-1}-M^{-1}=-\widetilde M^{-1}EM^{-1}\)
gives the second.  Along \(M_t=M+tE\), the first bound gives
\(\lVert M_t^{-1}\rVert\le(\kappa-\lVert E\rVert)^{-1}\).
Jacobi's formula and
\(|\operatorname{tr}(M_t^{-1}E)|
 \le r\lVert M_t^{-1}\rVert\lVert E\rVert\)
integrated from \(0\) to \(1\) prove (48.20).  The determinant cannot
cross zero along this path, so the logarithm of its absolute value is
well defined.
\end{proof}

\begin{ledger}[Finite-cutoff gravitational gauge row]
\label{led:v11v48-gravity}
For a regulated ADM-like system, a usable first-class reduction must
supply:
\begin{enumerate}
\item closure of the discrete constraint algebra, or a quantified
closure defect that vanishes in the adjacent comparison;
\item a gauge map \(\chi_i\) and a good-set lower bound
\(s_{\min}(M_i)\ge\kappa_i\);
\item a copy count \(N_i\), or a fundamental region with controlled
boundary, rather than an implicit assumption \(N_i=1\);
\item the coarea factor \(J_{\chi_i}^{C_i}\) together with
\(|\det M_i|\);
\item an orbit-volume normalization compatible with adjacent cutoffs.
\end{enumerate}
\end{ledger}

\begin{firewall}[Noncompact and global obstructions]
\label{fire:v11v48-global}
Theorem~\ref{thm:v11v48-fp} is a finite-volume, finite-copy identity.
The diffeomorphism group relevant to gravity is noncompact and gauge
conditions can have multiple or continuous copies.  Local
invertibility of \(M\) neither supplies a finite global orbit volume
nor rules out Gribov horizons where \(\det M=0\).  BRST language does
not remove these analytic requirements; it repackages them.
\end{firewall}

\begin{decision}[Next acquisition]
\label{dec:v11v48-next}
V49 will combine the constraint Gram matrix of V47 and the
Faddeev--Popov matrix of V48 into one small-singular-value ledger.  It
will derive the exact exponent losses for inverse coarea factors,
Faddeev--Popov determinants, and their adjacent log increments, then
state a feasible threshold schedule for insertion into the V42 master
criterion.
\end{decision}

\begin{assessment}[Progress after V48]
\label{ass:v11v48-status}
Regular first-class reduction is now explicit at finite dimension.
The determinant in the ambient delta formula, its coarea correction on
the slice, its equality with the square root of the Dirac determinant,
and its dependence on copy count are fixed.  The open work is to obtain
uniform small-ball and global-volume control for the actual regulated
Einstein constraints.
\end{assessment}

\subsection{Parent record}

This V48 note was formed by appending the preceding material to the
validated V47 source, whose record was
\[
\begin{array}{c|c}
\text{lines}&15010\\
\text{bytes}&449333\\
\text{SHA--256}&
\texttt{E0251E4FC53377D50E06BEF55D4016E908B6463CA15EABFED8040EDF7D205CDF}.
\end{array}                                                     \tag{48.21}
\]

% =============================================================================
% END APPENDED ELEVENTH-SERIES V48 MATERIAL
% =============================================================================

% =============================================================================
% BEGIN APPENDED ELEVENTH-SERIES V49 MATERIAL
% =============================================================================

\section{Small-ball exponents for constraint and gauge determinants}
\label{sec:v11v49}

The regular constraint and first-class reduction cards contain three
matrix families:
\[
 A=D{\cal C},\qquad B=D\chi|_{T C},\qquad
 M=D\chi\,R=\{\chi,{\cal C}\}.                                \tag{49.1}
\]
Their associated factors are
\[
 J_{\cal C}=\sqrt{\det(AA^*)},\qquad
 J_\chi^C=\sqrt{\det(BB^*)},\qquad
 \Delta_{\rm FP}=|\det M|.                                   \tag{49.2}
\]
This section determines what a least-singular-value estimate actually
buys for these factors.  Dimension losses are kept visible because
they decide whether the V42 series can converge.

\subsection{Products of singular values}

For a surjective map \(L:H\to\mathbb R^k\), set
\[
                       J(L):=\sqrt{\det(LL^*)}
                             =\prod_{a=1}^k s_a(L),            \tag{49.3}
\]
where \(s_1(L)\ge\cdots\ge s_k(L)>0\).

\begin{lemma}[Determinant envelopes]
\label{lem:v11v49-envelope}
If \(\lVert L\rVert\le B\) and \(s_{\min}(L)\ge\eta\), then
\[
                  \eta^k\le J(L)\le B^k,\qquad
                  J(L)^{-p}\le\eta^{-kp}\quad(p>0).           \tag{49.4}
\]
For square \(M\) of size \(r\),
\[
                  |\det M|=J(M),                              \tag{49.5}
\]
so the same bounds hold with \(k=r\).
\end{lemma}

\begin{proof}
Every singular value lies between \(\eta\) and \(B\); multiply the
\(k\) factors in (49.3).  For a square matrix, the product of singular
values is the absolute determinant.
\end{proof}

\begin{theorem}[Inverse determinant moments from one small-ball row]
\label{thm:v11v49-inverse-moment}
Let \(L\) be a random surjective map into \(\mathbb R^k\), almost
surely, and suppose for \(0<\epsilon\le1\)
\[
              \mathbb P\{s_{\min}(L)\le\epsilon\}
                         \le C\epsilon^\alpha.                 \tag{49.6}
\]
Then
\[
                      \mathbb E[J(L)^{-p}]<\infty             \tag{49.7}
\]
is guaranteed by this hypothesis whenever
\[
                               kp<\alpha.                      \tag{49.8}
\]
More precisely,
\[
 \mathbb E[J(L)^{-p}]
 \le 1+\frac{Ckp}{\alpha-kp}                                  \tag{49.9}
\]
if (49.6) holds with the displayed \(C\) and
\(J(L)^{-p}\le s_{\min}(L)^{-kp}\) is used.
\end{theorem}

\begin{proof}
Lemma~\ref{lem:v11v49-envelope} gives
\(J(L)^{-p}\le s_{\min}(L)^{-kp}\).  For
\(X=s_{\min}(L)^{-1}\), the layer-cake formula gives
\[
 \mathbb E[X^{kp}]
 =1+kp\int_1^\infty t^{kp-1}\mathbb P\{X>t\}\,dt
 \le1+Ckp\int_1^\infty t^{kp-1-\alpha}\,dt.                  \tag{49.10}
\]
The integral equals \((\alpha-kp)^{-1}\) exactly under (49.8).
\end{proof}

\begin{warning}[Dimension loss]
\label{warn:v11v49-dimension}
A tail exponent \(\alpha\) for only the least singular value controls
the full inverse \(k\)-dimensional Jacobian by the sufficient condition
\(\alpha>kp\).  The weaker condition \(\alpha>p\) would require
separate lower control of the other \(k-1\) singular values, or a
joint singular-value density estimate.  It cannot be inferred from
(49.6).
\end{warning}

\subsection{Logarithmic and absolute adjacent variation}

\begin{lemma}[Log-Jacobian differential]
\label{lem:v11v49-log}
Let \(L_t\) be a \(C^1\) path of surjective maps into
\(\mathbb R^k\).  Then
\[
 \frac d{dt}\log J(L_t)
 =\operatorname{tr}\!\left[
    L_t^*(L_tL_t^*)^{-1}\dot L_t\right],                      \tag{49.11}
\]
and
\[
 \left|\frac d{dt}\log J(L_t)\right|
 \le k\,s_{\min}(L_t)^{-1}\lVert\dot L_t\rVert.               \tag{49.12}
\]
\end{lemma}

\begin{proof}
Differentiate
\(\log J(L_t)=\frac12\log\det(L_tL_t^*)\).
Jacobi's formula and cyclicity of trace give (49.11).  The operator
\(L_t^*(L_tL_t^*)^{-1}\) is the Moore--Penrose right inverse and has
norm \(s_{\min}(L_t)^{-1}\).  The trace bound
\(|\operatorname{tr}X|\le k\lVert X\rVert\) proves (49.12).
\end{proof}

\begin{proposition}[Good-set adjacent bounds]
\label{prop:v11v49-adjacent}
Let \(L,\widetilde L:H\to\mathbb R^k\) obey
\[
 s_{\min}(L),s_{\min}(\widetilde L)\ge2K^{-1},\qquad
 \delta:=\lVert\widetilde L-L\rVert\le K^{-1}.                \tag{49.13}
\]
Then the segment \(L_t=L+t(\widetilde L-L)\) satisfies
\(s_{\min}(L_t)\ge K^{-1}\), and
\[
       |\log J(\widetilde L)-\log J(L)|\le kK\delta.           \tag{49.14}
\]
If also \(kK\delta\le1\), then
\[
       |J(\widetilde L)^{-1}-J(L)^{-1}|
                     \le 2kK^{k+1}\delta.                     \tag{49.15}
\]
If instead only \(\lVert L\rVert,\lVert\widetilde L\rVert\le B\)
is used, then
\[
                  |J(\widetilde L)-J(L)|
                     \le kB^{k-1}\delta.                      \tag{49.16}
\]
\end{proposition}

\begin{proof}
The singular-value perturbation inequality yields the segment lower
bound.  Integrating Lemma~\ref{lem:v11v49-log} proves (49.14).
By Lemma~\ref{lem:v11v49-envelope}, both inverse Jacobians are at most
\(K^k\) along the segment.  Write the quotient of the endpoint
Jacobians as \(e^z\), where \(|z|\le kK\delta\le1\), and use
\(|e^z-1|\le2|z|\).  This gives (49.15).

For (49.16), use
\(J(L)=\prod_{a=1}^k s_a(L)\), the telescoping product identity, and
\(|s_a(\widetilde L)-s_a(L)|\le\delta\).  Each remaining factor is at
most \(B\).
\end{proof}

For \(M\), (49.16) shows that the Faddeev--Popov determinant itself
does not require \(M^{-1}\) for an absolute comparison.  Invertibility
of \(M\) is nevertheless required to maintain a valid slice and a
constant copy count, while a logarithmic determinant comparison costs
one power \(K_M\).

\subsection{The combined reduction density}

In the induced constraint-and-gauge slice convention of V47--V48, a
typical geometric density has the schematic form
\[
                  {\cal R}
     =J(A)^{\varepsilon_C}
       \frac{|\det M|}{J(B)},\qquad
                  \varepsilon_C\in\{0,-1,1\}.                 \tag{49.17}
\]
Here \(\varepsilon_C=1\) represents the induced constraint measure in
an ambient delta formula, \(\varepsilon_C=-1\) represents a bare
constraint delta after coarea, and \(\varepsilon_C=0\) means that the
constraint measure has already been fixed intrinsically.

\begin{proposition}[Deterministic stock of the reduction density]
\label{prop:v11v49-stock}
Suppose on a common good set
\[
\begin{split}
 &s_{\min}(A),s_{\min}(\widetilde A)\ge2K_A^{-1},\\
 &s_{\min}(B),s_{\min}(\widetilde B)\ge2K_B^{-1},\\
 &s_{\min}(M),s_{\min}(\widetilde M)\ge2K_M^{-1},              \tag{49.18}\\
 &\lVert A\rVert,\lVert\widetilde A\rVert\le B_A,\quad
  \lVert B\rVert,\lVert\widetilde B\rVert\le B_B,\quad
  \lVert M\rVert,\lVert\widetilde M\rVert\le B_M.
\end{split}
\]
Let the row dimensions of \(A,B,M\) be \(m,r,r\), respectively.
When the three perturbations are small enough for
Proposition~\ref{prop:v11v49-adjacent}, the logarithmic increment obeys
\[
 |\Delta\log{\cal R}|
 \le mK_A\delta_A\,\mathbf1_{\{\varepsilon_C\ne0\}}
       +rK_M\delta_M+rK_B\delta_B.                             \tag{49.19}
\]
For absolute increments, the dangerous inverse-Jacobian stocks are
\[
                   K_A^{m+1}\delta_A
       \quad\text{if }\varepsilon_C=-1,\qquad
                   K_B^{r+1}\delta_B,                         \tag{49.20}
\]
up to constants and upper determinant factors.  Variation of
\(|\det M|\) costs only
\(rB_M^{r-1}\delta_M\), while maintaining the slice costs the separate
threshold \(K_M\).
\end{proposition}

\begin{proof}
Take logarithms in (49.17) and apply (49.14) to each nonzero factor.
For inverse Jacobians use (49.15).  For the numerator determinant use
(49.16).  Product differences are expanded one factor at a time; all
unchanged positive factors are bounded by
Lemma~\ref{lem:v11v49-envelope}.  This introduces only the stated
upper determinant factors and constants.
\end{proof}

\subsection{Feasible polynomial thresholds}

Let \(j\) index any one of the singular-value rows \(A,B,M\).  Suppose
\[
 \mathbb P\{s_{\min}(L_{j,i})\le\epsilon\}
       \le C_j(i+2)^{d_j}\epsilon^{\alpha_j},                 \tag{49.21}
\]
the bad-set observable has an \(L^{q_j}\) envelope growing at most as
\((i+2)^{b_j}\), and the native adjacent error is
\[
                         \delta_{j,i}\le C'_j(i+2)^{-s_j}.     \tag{49.22}
\]
Choose \(K_{j,i}=(i+2)^{a_j}\).  Let \(\ell_j\) be the power of \(K_j\)
used by the good-set deterministic comparison: \(\ell_j=1\) for a
log-Jacobian row, \(\ell_j=k_j+1\) for an inverse \(k_j\)-Jacobian
absolute row.

\begin{theorem}[Threshold interval]
\label{thm:v11v49-threshold}
The good-set series and the Hölder-controlled complement series both
converge if
\[
 \frac{b_j+d_j(1-1/q_j)+1}
      {\alpha_j(1-1/q_j)}<a_j<
                         \frac{s_j-1}{\ell_j}.                 \tag{49.23}
\]
Such an \(a_j\) exists exactly when
\[
 \alpha_j(1-1/q_j)(s_j-1)>
        \ell_j\bigl[b_j+d_j(1-1/q_j)+1\bigr].                 \tag{49.24}
\]
For finitely many rows, choosing one \(a_j\) in each nonempty interval
and taking the union of their bad events gives a summable combined
constraint-and-gauge contribution to the V42 master criterion.
\end{theorem}

\begin{proof}
Equation (49.21) at \(\epsilon=K_{j,i}^{-1}\) gives a bad probability
bounded by
\(C_j(i+2)^{d_j-a_j\alpha_j}\).  Hölder therefore bounds the bad
summand by
\[
 C(i+2)^{b_j+(d_j-a_j\alpha_j)(1-1/q_j)},                     \tag{49.25}
\]
which is summable under the left inequality in (49.23).  The good-set
summand is bounded by
\[
 C(i+2)^{-s_j+a_j\ell_j},                                     \tag{49.26}
\]
summable under the right inequality.  The interval is nonempty exactly
when (49.24) holds.  A finite union preserves summability.
\end{proof}

\begin{remark}[Bookkeeping convention]
The factor \(d_j(1-1/q_j)\) is the price of polynomial growth in the
small-ball prefactor.  It must remain visible unless it has already
been included in a redefined bad-envelope exponent.
\end{remark}

\subsection{Insertion into the SM--Einstein master criterion}

\begin{ledger}[New V42 rows]
\label{led:v11v49-master}
The gravitational reduction adds the following rows to the master
adjacent increment:
\[
\begin{array}{c|c|c}
\text{row}&\text{good-set stock}&\text{required tail}\\ \hline
\log J_{\cal C}&K_A\delta_A&
 \mathbb P(s_{\min}A<K_A^{-1})\\
J_{\cal C}^{-1}&K_A^{m+1}\delta_A&
 \mathbb P(s_{\min}A<K_A^{-1})\\
\log|\det M|&K_M\delta_M&
 \mathbb P(s_{\min}M<K_M^{-1})\\
|\det M|&B_M^{r-1}\delta_M&
 \text{invertibility event still required}\\
(J_\chi^C)^{-1}&K_B^{r+1}\delta_B&
 \mathbb P(s_{\min}B<K_B^{-1})\\
N_i,\ \operatorname{vol}H_i&\text{discrete/global variation}&
 \text{no singular-value tail suffices}.
\end{array}                                                     \tag{49.27}
\]
The final row is topological and global.  It cannot be paid for by
improving a matrix small-ball exponent.
\end{ledger}

\begin{firewall}[No acquired exponent]
\label{fire:v11v49-no-exponent}
V49 computes the exponent requirements but does not prove (49.21) for
a regulated Einstein constraint map or Faddeev--Popov matrix.  In
particular, the scalar-shift argument of V43 applies only when an
independent scalar direction shifts the relevant matrix with bounded
conditional density.  No such direction has yet been constructed for
the gravitational constraint algebra.
\end{firewall}

\begin{decision}[V50 audit]
\label{dec:v11v49-next}
The next version is the mandatory companion audit.  It will first
check whether the YM or NS programmes have advanced beyond the records
in V45.  The following substantive version will then pursue the
weakest missing input exposed here: a conditional transversality
mechanism for \(M\), or, if the companions provide no such mechanism,
a deterministic reference-gauge perturbation theorem that reduces the
small-ball question to a quantified random matrix perturbation.
\end{decision}

\begin{assessment}[Progress after V49]
\label{ass:v11v49-status}
The coarea and Faddeev--Popov determinants now have explicit moment and
adjacent-comparison costs.  The dimension penalty is unavoidable under
a least-singular-value tail alone, and the feasibility inequality
(49.24) states exactly how strong the native rate and small-ball
exponent must be.  The remaining obstacle is acquisition of those
rates for the actual Einstein regulator, plus global copy and orbit
volume control.
\end{assessment}

\subsection{Parent record}

This V49 note was formed by appending the preceding material to the
validated V48 source, whose record was
\[
\begin{array}{c|c}
\text{lines}&15342\\
\text{bytes}&461258\\
\text{SHA--256}&
\texttt{98D4B59FF5E9C5AA8EE5BFBBBBF42574DBAFC380595028934909C3484EFE2905}.
\end{array}                                                     \tag{49.28}
\]

% =============================================================================
% END APPENDED ELEVENTH-SERIES V49 MATERIAL
% =============================================================================

% =============================================================================
% BEGIN APPENDED ELEVENTH-SERIES V50 MATERIAL
% =============================================================================

\section{Companion audit at V50: positive-homotopy concentration and
tree-chart transfer}
\label{sec:v11v50}

The second compulsory audit of this cycle finds a substantive change
in the Yang--Mills companion.  Its V35 now contains a genuine new
block, rather than the byte-identical V34 carry-forward observed at
V45.  The Navier--Stokes companion remains unchanged.

\subsection{Audited records}

\begin{definition}[V50 companion record]
\label{def:v11v50-record}
The current companion files are
\[
\begin{array}{c|r|l}
\text{file}&\text{bytes}&\text{SHA--256}\\ \hline
\texttt{Renewal\_Yang\_Mills\_Mass\_Gap\_Research\_Notes\_V35.tex}
 &498660&
\texttt{23F674C142CA9B298546508F7A3E9EAF2CE2E154078C32B0BA2FC8660AD25098}\\
\texttt{Renewal\_NS\_Stability\_Research\_Notes\_V26.tex}
 &278283&
\texttt{82C887F8C2C69E4F28FAD7C3DC8DE5B9099CB5A4BF431EBA1FFF3E91935978AD}.
\end{array}                                                     \tag{50.1}
\]
The Yang--Mills digest and byte count differ from the V45 audit.  Its
new mathematical payload is labelled compact-series V35.  The NS
record is identical to (45.1), so its rank-one transfer is not
repeated.
\end{definition}

\subsection{The new Yang--Mills input}

On the periodic four-dimensional lattice, the companion considers
the positive homotopy
\[
 V_t=V_{\rm W}+t{\cal O}_1,\qquad 0\le t\le1,                 \tag{50.2}
\]
where both terms are nonnegative.  Its global partition-function lower
bound on the event that all links lie in a Haar ball yields, for every
translation-invariant infinite-volume subsequential limit,
\[
 {\mathbb E}_{\beta,\infty,t}z_p
 \le m_{\beta,t}
 :=\frac{C_t+\frac23\log(c_G^{-1})+\frac83\log\beta}{\beta},
 \qquad C_t=8+384t\le392.                                    \tag{50.3}
\]
Here \(z_p=\frac16\lVert I-U_p\rVert_{\rm HS}^2\), and \(c_G>0\)
is the local \(SU(3)\) Haar-ball constant.  Consequently, for a finite
plaquette block \(B\),
\[
 \mathbb P_{\beta,\infty,t}
 \left\{\max_{p\in B}\lVert I-U_p\rVert_{\rm op}>\epsilon\right\}
 \le\frac{6|B|m_{\beta,t}}{\epsilon^2}.                       \tag{50.4}
\]
No independence of plaquettes is used.

\begin{proposition}[Compatibility with exact tree coordinates]
\label{prop:v11v50-compatible}
Let \(K\) be a finite contractible cubical block with rooted spanning
tree \(T\).  Suppose each fundamental chord loop has a fixed cellular
contraction using at most \(A_K\) plaquettes.  On the event
\[
                   \max_{p\subset K}
                   \lVert I-U_p\rVert_{\rm op}\le\epsilon,     \tag{50.5}
\]
the global based tree coordinate of V46 satisfies
\[
                    \max_{e\in E(K)\setminus T}
                    \lVert I-W_e\rVert_{\rm op}
                    \le A_K\epsilon.                          \tag{50.6}
\]
All tree links equal the identity on the tree slice.
\end{proposition}

\begin{proof}
By Proposition~\ref{prop:v11v46-chord-generators}, \(W_e\) is the
holonomy of the fundamental loop of \(e\).  The chosen contraction
writes that holonomy as an ordered product of at most \(A_K\)
conjugates of plaquette holonomies or their inverses.  Operator distance
from the identity is invariant under conjugation and inversion.
For unitary \(Q_1,\ldots,Q_N\),
\[
 \lVert Q_1\cdots Q_N-I\rVert_{\rm op}
 \le\sum_{j=1}^N\lVert Q_j-I\rVert_{\rm op}.                  \tag{50.7}
\]
Using (50.5) proves (50.6).  The assertion about tree links is
Theorem~\ref{thm:v11v46-global-tree}.
\end{proof}

\begin{corollary}[Uniform principal-chart failure bound]
\label{cor:v11v50-chart}
If \(A_K\epsilon\le\delta_{\rm ch}\), where
\(\delta_{\rm ch}\) is a fixed principal-logarithm radius, then the
failure probability of the tree-gauged link chart on \(K\) is at most
\[
                         \frac{6|B(K)|m_{\beta,t}}{\epsilon^2},\tag{50.8}
\]
uniformly for \(0\le t\le1\).  For a sequence
\((K_i,\beta_i,\epsilon_i)\), the chart failures are summable whenever
\[
 \sum_i\frac{|B(K_i)|(1+\log\beta_i)}
                 {\beta_i\epsilon_i^2}<\infty,\qquad
                         A_{K_i}\epsilon_i\le\delta_{\rm ch}. \tag{50.9}
\]
\end{corollary}

\begin{proof}
On the complement of the plaquette-bad event, (50.6) puts every chord
in the principal chart.  Apply (50.4) and the bound
\(C_t\le392\).  Summing gives (50.9).
\end{proof}

\subsection{Local curvature versus global holonomy}

\begin{proposition}[The torus obstruction is a separate row]
\label{prop:v11v50-torus}
Small plaquette curvature on a periodic lattice does not imply that
all chord coordinates of V46 are close to the identity.
\end{proposition}

\begin{proof}
Choose commuting constants \(H_\mu\in G\) and put
\(U_{x,\mu}=H_\mu\) on every link in direction \(\mu\).
Every plaquette is
\(H_\mu H_\nu H_\mu^{-1}H_\nu^{-1}=I\), while a Wilson line winding
around direction \(\mu\) is \(H_\mu^L\), which can be far from the
identity.  Based gauge transformations leave its conjugacy class
unchanged.  Thus zero local curvature can coexist with nontrivial
global chord holonomy.
\end{proof}

This proves why the block event (50.5) is useful only on a
contractible acquisition domain.  Sector tightness and toron control
remain independent entries in the V42 ledger.

\subsection{Transfer to a matrix transversality problem}

The concentration estimate becomes relevant to V49 only after a
deterministic Lipschitz bridge from link charts to the matrix under
study.

\begin{lemma}[Chart-to-matrix bridge]
\label{lem:v11v50-matrix-bridge}
Let \(L(U)\) be a matrix-valued function on the tree-gauged block and
assume
\[
       \lVert L(U)-L(I)\rVert
       \le {\cal L}_K
           \max_{e\in E(K)\setminus T}\lVert W_e-I\rVert_{\rm op}.
                                                                  \tag{50.10}
\]
If \(s_{\min}(L(I))\ge\kappa_0>0\), then on (50.5)
\[
                  s_{\min}(L(U))
                  \ge\kappa_0-{\cal L}_KA_K\epsilon.           \tag{50.11}
\]
In particular, if
\({\cal L}_KA_K\epsilon\le\kappa_0/2\), then
\(s_{\min}(L(U))\ge\kappa_0/2\).
\end{lemma}

\begin{proof}
Combine (50.6) with (50.10), then apply the least-singular-value
perturbation inequality
\[
 |s_{\min}(X)-s_{\min}(Y)|\le\lVert X-Y\rVert.
\]
\end{proof}

\begin{corollary}[Reference transversality with explicit failure]
\label{cor:v11v50-reference}
Under the hypotheses of Lemma~\ref{lem:v11v50-matrix-bridge} and
\({\cal L}_KA_K\epsilon\le\kappa_0/2\),
\[
 \mathbb P_{\beta,\infty,t}
       \{s_{\min}(L(U))<\kappa_0/2\}
 \le\frac{6|B(K)|m_{\beta,t}}{\epsilon^2}.                    \tag{50.12}
\]
\end{corollary}

\begin{proof}
The event on the left is contained in the plaquette-bad event by
Lemma~\ref{lem:v11v50-matrix-bridge}.  Apply (50.4).
\end{proof}

\begin{transfer}[What can enter the SM--Einstein ledger]
\label{trans:v11v50-use}
For the compact internal gauge sector, (50.12) can control any
finite-block Hessian, constraint, or response matrix whose reference
value has a proved gap and whose link dependence has a proved
Lipschitz constant.  For a mixed SM--Einstein matrix, the same argument
also needs a metric-and-matter chart event; gauge plaquette
concentration alone does not control those variables.  For the
gravitational Faddeev--Popov matrix \(M=\{\chi,{\cal C}\}\), the
reference gap and metric Lipschitz bridge remain to be acquired.
\end{transfer}

\subsection{Audit decision}

\begin{decision}[V51 acquisition]
\label{dec:v11v50-next}
V51 will prove the multi-field reference-matrix theorem needed by
Transfer~\ref{trans:v11v50-use}.  It will combine independent gauge,
metric, matter, and derivative chart radii into one deterministic
singular-value gap, then combine their failure bounds without assuming
independence.  The output will be an explicit sufficient schedule for
the \(A,B,M\) rows of V49.  The theorem will not assert that the
Einstein reference gap exists; it will isolate that finite computation
as a checkable acquisition.
\end{decision}

\begin{firewall}[Audit boundary]
\label{fire:v11v50-boundary}
The V35 Yang--Mills bound is a weak-coupling local curvature estimate
uniform along a positive endpoint homotopy.  It does not control
noncontractible holonomy, the chiral determinant phase, gravitational
constraint rank, a noncompact orbit volume, adjacent cylinder
comparison, reflection positivity, or clustering.  The unchanged NS
note contributes no new transfer at V50.
\end{firewall}

\begin{assessment}[Progress after V50]
\label{ass:v11v50-status}
The mandatory audit is complete and changes the next step
constructively.  The internal gauge chart now has an explicit
probability bound, while Lemma~\ref{lem:v11v50-matrix-bridge} states
the exact deterministic information still required to turn that chart
into a singular-value estimate.
\end{assessment}

\subsection{Parent record}

This V50 note was formed by appending the preceding material to the
validated V49 source, whose record was
\[
\begin{array}{c|c}
\text{lines}&15705\\
\text{bytes}&474366\\
\text{SHA--256}&
\texttt{139115605A7A18DD5FB5369C72C9E84A581E7D52B65F87AB971D6FE32383D858}.
\end{array}                                                     \tag{50.13}
\]

% =============================================================================
% END APPENDED ELEVENTH-SERIES V50 MATERIAL
% =============================================================================
% =============================================================================
% BEGIN APPENDED ELEVENTH-SERIES V51 MATERIAL
% =============================================================================

\section{Certified multi-field reference atlases for transversality}
\label{sec:v11v51}

A single reference matrix controls a neighborhood of one background.
The metric reference of V39 supplies ellipticity and moment bounds, but
it does not force every metric to be close to the identity at fixed
stabilizer strength.  The correct finite-cutoff replacement is a
certified finite atlas over a compact good chamber.  This section proves
that reduction and combines its probability with the gauge chart of
V50.

\subsection{A product-chamber gap theorem}

Let \(X_j\) be metric spaces, \(1\le j\le s\), and let
\[
                         {\cal K}=\prod_{j=1}^s{\cal K}_j      \tag{51.1}
\]
with each \({\cal K}_j\) compact.  Let
\(L:{\cal K}\to{\cal L}(E,F)\) be a continuous family of linear maps
between fixed finite-dimensional Hilbert spaces.  Assume the
coordinate Lipschitz estimate
\[
 \lVert L(x)-L(y)\rVert
       \le\sum_{j=1}^s\Lambda_jd_j(x_j,y_j).                  \tag{51.2}
\]

\begin{theorem}[Finite atlas certificate]
\label{thm:v11v51-atlas}
For each \(j\), let \({\cal N}_j\subset{\cal K}_j\) be a finite
\(\rho_j\)-net.  Suppose that for every
\(a\in{\cal N}:=\prod_j{\cal N}_j\),
\[
                         s_{\min}(L(a))\ge\kappa_{\rm net}.    \tag{51.3}
\]
Then for every \(x\in{\cal K}\),
\[
 s_{\min}(L(x))
       \ge\kappa_{\rm net}-\sum_{j=1}^s\Lambda_j\rho_j.        \tag{51.4}
\]
In particular, if
\[
                    \sum_j\Lambda_j\rho_j
                         \le\frac12\kappa_{\rm net},           \tag{51.5}
\]
then \(s_{\min}(L(x))\ge\kappa_{\rm net}/2\) throughout
\({\cal K}\).
\end{theorem}

\begin{proof}
Compactness gives finite nets.  For \(x\in{\cal K}\), choose
\(a_j\in{\cal N}_j\) with \(d_j(x_j,a_j)\le\rho_j\).
The least-singular-value perturbation inequality and (51.2) give
\[
 s_{\min}(L(x))
 \ge s_{\min}(L(a))-\lVert L(x)-L(a)\rVert
 \ge\kappa_{\rm net}-\sum_j\Lambda_j\rho_j.
\]
This is (51.4), and (51.5) gives the last assertion.
\end{proof}

\begin{corollary}[Simultaneous certificate]
\label{cor:v11v51-simultaneous}
Let \(L^\ell\), \(1\le\ell\le q\), be finitely many matrix families
with coordinate constants \(\Lambda_j^\ell\).  If
\[
 s_{\min}(L^\ell(a))\ge\kappa_\ell
 \quad\text{on every net point},\qquad
 \sum_j\Lambda_j^\ell\rho_j\le\kappa_\ell/2,                 \tag{51.6}
\]
then all families have the simultaneous chamber bounds
\[
                       s_{\min}(L^\ell(x))\ge\kappa_\ell/2.   \tag{51.7}
\]
\end{corollary}

\begin{proof}
Apply Theorem~\ref{thm:v11v51-atlas} to each of the finitely many
families.
\end{proof}

The intended families are the constraint derivative \(A=D{\cal C}\),
the slice coarea derivative \(B=D\chi|_{TC}\), the Faddeev--Popov
matrix \(M=D\chi R\), and the corrected tangent Hessian of V47.
Equation (51.6) is a finite interval-arithmetic acquisition, not a
continuum assertion.

\subsection{Changing dimensions at adjacent cutoffs}

The number of constraints and variables can grow with the regulator.
A comparison of matrices of different sizes is meaningless until an
identification and a new-mode block are specified.

\begin{lemma}[Reference block extension]
\label{lem:v11v51-block}
Let
\[
 L_i^0:E_i\to F_i,\qquad N_i^0:E_i^{\rm new}\to F_i^{\rm new}
\]
be surjective and satisfy
\[
 s_{\min}(L_i^0)\ge\kappa_i,\qquad
 s_{\min}(N_i^0)\ge\nu_i.                                    \tag{51.8}
\]
Then
\[
 \widehat L_{i+1}^0
 :=L_i^0\oplus N_i^0:
 E_i\oplus E_i^{\rm new}\to F_i\oplus F_i^{\rm new}
\]
has
\[
                s_{\min}(\widehat L_{i+1}^0)
                    =\min\{\kappa_i,\nu_i\}.                  \tag{51.9}
\]
If
\(L_{i+1}=\widehat L_{i+1}^0+E_i^{\rm mix}\), then
\[
 s_{\min}(L_{i+1})
 \ge\min\{\kappa_i,\nu_i\}-\lVert E_i^{\rm mix}\rVert.        \tag{51.10}
\]
\end{lemma}

\begin{proof}
For \(u\oplus v\),
\[
 \lVert(L_i^0\oplus N_i^0)(u\oplus v)\rVert^2
 \ge\kappa_i^2\lVert u\rVert^2+\nu_i^2\lVert v\rVert^2,
\]
which gives the lower bound in (51.9).  Unit vectors approaching a
least singular vector in either summand give equality.  The
perturbation inequality proves (51.10).
\end{proof}

Thus uniform transversality requires a gap for the newly introduced
constraint modes and a small old--new mixing norm.  Stability of the
old block alone cannot prove it.

\subsection{The metric chamber already available}

For the one-site positive metric reference of V39, recall
\[
 Z(g)=\operatorname{tr}g+\operatorname{tr}g^{-1},\qquad
 {\mathbb E}_{\mu_\kappa}e^{\theta Z}
 =\frac{I_n(\kappa-\theta)}{I_n(\kappa)}
 =:C_{\kappa,\theta},\quad0<\theta<\kappa.                   \tag{51.11}
\]

\begin{lemma}[Finite-block ellipticity chamber]
\label{lem:v11v51-metric-chamber}
For a block \(\Lambda\) under the product reference
\(\mu_{\kappa,\Lambda}\), define
\[
 {\cal E}_{\Lambda,T}
     =\{g:T^{-1}I\le g_x\le TI\text{ for every }x\in\Lambda\},
 \qquad T\ge2n.                                               \tag{51.12}
\]
Then
\[
 \mu_{\kappa,\Lambda}({\cal E}_{\Lambda,T}^c)
 \le|\Lambda|C_{\kappa,\theta}e^{-\theta T}.                 \tag{51.13}
\]
The chamber in (51.12) is compact.
\end{lemma}

\begin{proof}
If \(Z(g_x)\le T\), then every eigenvalue \(\lambda\) of \(g_x\)
obeys
\[
              \lambda\le\operatorname{tr}g_x\le T,\qquad
              \lambda^{-1}\le\operatorname{tr}g_x^{-1}\le T.
\]
Thus failure of (51.12) implies \(Z(g_x)>T\) for some \(x\).
Exponential Markov and (51.11) give
\[
 \mathbb P\{Z(g_x)>T\}\le C_{\kappa,\theta}e^{-\theta T}.
\]
Take the union bound.  Closedness and boundedness in the
finite-dimensional product of symmetric matrices prove compactness.
\end{proof}

For interacting positive references satisfying the conditional moment
hypotheses of V38--V39, the same proof uses the corresponding uniform
constant in place of \(C_{\kappa,\theta}\).

\subsection{A simultaneous probability theorem}

At cutoff \(i\), let the bosonic background be split into gauge,
metric, matter, and discrete derivative data,
\[
                           X_i=(U_i,g_i,\phi_i,D_i).            \tag{51.14}
\]
Choose good chambers
\({\cal K}_{U,i},{\cal K}_{g,i},{\cal K}_{\phi,i},{\cal K}_{D,i}\)
and write
\[
\begin{split}
 p_{U,i}&=\mathbb P\{U_i\notin{\cal K}_{U,i}\},\\
 p_{g,i}&=\mathbb P\{g_i\notin{\cal K}_{g,i}\},\\
 p_{\phi,i}&=\mathbb P\{\phi_i\notin{\cal K}_{\phi,i}\},\\
 p_{D,i}&=\mathbb P\{D_i\notin{\cal K}_{D,i}\}.               \tag{51.15}
\end{split}
\]

\begin{theorem}[Certified hard gap with dependent fields]
\label{thm:v11v51-hard-gap}
Suppose Corollary~\ref{cor:v11v51-simultaneous} certifies on the product
good chamber
\[
 s_{\min}(A_i)\ge\gamma_{A,i},\quad
 s_{\min}(B_i)\ge\gamma_{B,i},\quad
 s_{\min}(M_i)\ge\gamma_{M,i},\quad
 s_{\min}(H_i^C)\ge\gamma_{H,i}.                              \tag{51.16}
\]
Then
\[
 \mathbb P\!\left\{
  \min\!\left(
   \frac{s_{\min}(A_i)}{\gamma_{A,i}},
   \frac{s_{\min}(B_i)}{\gamma_{B,i}},
   \frac{s_{\min}(M_i)}{\gamma_{M,i}},
   \frac{s_{\min}(H_i^C)}{\gamma_{H,i}}
       \right)<1\right\}
 \le p_{U,i}+p_{g,i}+p_{\phi,i}+p_{D,i}.                      \tag{51.17}
\]
No independence among the four field families is required.
\end{theorem}

\begin{proof}
Every certified bound holds on the intersection of the four good
events.  Hence the event on the left is contained in the union of
their complements.  The union bound proves (51.17).
\end{proof}

For the internal gauge chamber on a contractible block, V50 gives
\[
 p_{U,i}\le
 \frac{6|B(K_i)|m_{\beta_i,t_i}}{\epsilon_i^2},
 \qquad A_{K_i}\epsilon_i\le\delta_{\rm ch}.                  \tag{51.18}
\]
For the metric product reference, Lemma~\ref{lem:v11v51-metric-chamber}
gives
\[
 p_{g,i}\le|\Lambda_i|C_{\kappa_i,\theta_i}
                                  e^{-\theta_iT_i}.           \tag{51.19}
\]
Moment or conditional-moment estimates can supply the last two rows.

\begin{corollary}[Summable certified-gap schedule]
\label{cor:v11v51-schedule}
If
\[
 \sum_i\left[
 \frac{|B(K_i)|(1+\log\beta_i)}
      {\beta_i\epsilon_i^2}
 +|\Lambda_i|C_{\kappa_i,\theta_i}e^{-\theta_iT_i}
 +p_{\phi,i}+p_{D,i}\right]<\infty,                           \tag{51.20}
\]
and the finite atlas certificates (51.6) hold at every \(i\), then all
four failures in (51.17) occur only finitely often under any compatible
common coupling.
\end{corollary}

\begin{proof}
Equations (51.18)--(51.20) and Theorem
\ref{thm:v11v51-hard-gap} make the joint failure probabilities
summable.  The first Borel--Cantelli lemma requires no independence.
\end{proof}

\subsection{Insertion into the adjacent-cutoff ledger}

Set
\[
 K_{A,i}=\gamma_{A,i}^{-1},\quad
 K_{B,i}=\gamma_{B,i}^{-1},\quad
 K_{M,i}=\gamma_{M,i}^{-1},\quad
 K_{H,i}=\gamma_{H,i}^{-1}.                                  \tag{51.21}
\]
On the certified event, all deterministic estimates of V43, V47--V49
apply with these thresholds.  If, for example,
\(\gamma_{A,i}\asymp(i+2)^{-a_A}\), then the inverse constraint
Jacobian row costs
\[
                         (i+2)^{a_A(m_i+1)}\delta_{A,i},       \tag{51.22}
\]
where the constraint dimension \(m_i\) must itself be tracked.  A
growing \(m_i\) generally makes a polynomial ansatz insufficient unless
the determinant is normalized per degree of freedom or a stronger
joint singular-value estimate replaces the crude least-value bound.

\begin{warning}[Hard gaps do not control the bad integrand]
\label{warn:v11v51-bad-integrand}
Equation (51.17) controls the probability of leaving the certified
chamber.  It does not bound an inverse determinant or inverse Hessian
on that complement.  The V42 bad-set term still requires an
\(L^q\) envelope, an alternate nonsingular chart, or a regulated
definition whose removal is controlled.  Multiplying an undefined
inverse by a small failure probability is not an estimate.
\end{warning}

\begin{ledger}[Finite computations now isolated]
\label{led:v11v51-computations}
At each finite cutoff, the remaining checkable transversality
acquisition is:
\begin{enumerate}
\item specify compact chambers for all bosonic backgrounds;
\item bound the coordinate Lipschitz constants of \(A,B,M,H^C\);
\item construct finite nets or interval boxes;
\item certify their least singular values at every box;
\item certify new-mode gaps and old--new mixing as dimensions grow;
\item attach an alternate chart or bad-set \(L^q\) envelope.
\end{enumerate}
The first five items are finite, although possibly large, computations.
Uniform asymptotics and the sixth item remain analytic.
\end{ledger}

\begin{firewall}[Reference versus target]
\label{fire:v11v51-target}
The metric probability (51.19) uses the positive stabilizing reference
of V39.  Its fixed exponential budget cannot be used after the
stabilizer is removed, and no counterterm argument has shown that this
reference flows to Einstein dynamics.  The atlas theorem reduces
transversality to certified finite data; it does not certify those data
for a chosen SM--Einstein regulator.
\end{firewall}

\begin{decision}[Next acquisition]
\label{dec:v11v51-next}
The next step will address the warning rather than refine the good-set
gap.  V52 will prove a multi-chart quotient formula with a gauge-
invariant partition of unity.  Its purpose is to replace a single
Faddeev--Popov matrix that can cross a Gribov horizon by finitely many
local gauges on a compact regular quotient, while exposing the overlap
and stabilizer conditions needed for a well-defined bad-set envelope.
\end{decision}

\begin{assessment}[Progress after V51]
\label{ass:v11v51-status}
Gauge concentration and metric ellipticity moments now feed one
certified transversality theorem.  A finite atlas can prove uniform
gaps for the constraint, coarea, Faddeev--Popov, and constrained-
Hessian matrices on the complete good chamber, including new modes at
successive cutoffs.  The remaining local obstruction is coverage of
configurations where any one chosen gauge matrix becomes singular.
\end{assessment}

\subsection{Parent record}

This V51 note was formed by appending the preceding material to the
validated V50 source, whose record was
\[
\begin{array}{c|c}
\text{lines}&15956\\
\text{bytes}&483895\\
\text{SHA--256}&
\texttt{D047F74371F17A9122863771DC4E33B85329B2A5E9F20D556124032839AD63DD}.
\end{array}                                                     \tag{51.23}
\]

% =============================================================================
% END APPENDED ELEVENTH-SERIES V51 MATERIAL
% =============================================================================
% =============================================================================
% BEGIN APPENDED ELEVENTH-SERIES V52 MATERIAL
% =============================================================================

\section{Multi-chart gauge reduction on a compact regular quotient}
\label{sec:v11v52}

A single gauge condition may cross a Faddeev--Popov horizon even when
the underlying quotient is regular.  On a compact free quotient this
is a coordinate failure that can be removed by finitely many local
sections and a partition of unity.  This section proves the resulting
measure formula and identifies the conditions under which it supplies
the bad-chart envelope missing in V51.

\subsection{Canonical quotient measure}

Let a compact Lie group \(H\) act freely and smoothly on a manifold
\(C\).  Since a compact action is proper, the orbit space
\[
                             Q=C/H                             \tag{52.1}
\]
is a manifold and \(\pi:C\to Q\) is a principal \(H\)-bundle.
Let \(d\mu_C\) be a finite \(H\)-invariant Borel measure and let \(dh\)
be Haar measure of finite mass \(\operatorname{vol}H\).

\begin{definition}[Quotient measure]
\label{def:v11v52-quotient}
Define
\[
                        d\nu_Q
                  :=\frac1{\operatorname{vol}H}\,\pi_*d\mu_C. \tag{52.2}
\]
For every integrable gauge-invariant function
\(F=f\circ\pi\),
\[
 \frac1{\operatorname{vol}H}\int_CF\,d\mu_C
                         =\int_Qf\,d\nu_Q.                    \tag{52.3}
\]
\end{definition}

Equation (52.3) is immediate from pushforward and fixes the answer
before any gauge coordinate is chosen.

\subsection{Finite section atlas}

Assume henceforth that \(Q\) is compact.  Choose a finite principal-
bundle trivializing cover \((U_a)_{a=1}^N\) and smooth local sections
\[
                            s_a:U_a\to C.                      \tag{52.4}
\]
Every \(x\in\pi^{-1}(U_a)\) has a unique representation
\[
                    x=h_a(x)\cdot s_a(\pi(x)),\qquad h_a(x)\in H.
                                                                  \tag{52.5}
\]
Choose a smooth partition of unity
\[
            0\le\rho_a\le1,\qquad
            \operatorname{supp}\rho_a\Subset U_a,\qquad
            \sum_{a=1}^N\rho_a=1.                             \tag{52.6}
\]

For a test function \(\varphi\) on \(H\), let the Haar delta at the
identity be defined by
\[
                         \int_H\delta_e(h)\varphi(h)\,dh
                                  =\varphi(e).                 \tag{52.7}
\]

\begin{theorem}[Multi-section quotient formula]
\label{thm:v11v52-multisection}
For every integrable gauge-invariant \(F=f\circ\pi\),
\[
 \boxed{
 \int_Qf\,d\nu_Q
 =\sum_{a=1}^N\frac1{\operatorname{vol}H}
   \int_{\pi^{-1}(U_a)}
     F(x)\rho_a(\pi(x))\,
     \operatorname{vol}H\,\delta_e(h_a(x))\,d\mu_C(x).}       \tag{52.8}
\]
Equivalently, in the local disintegration
\[
             d\mu_C(x)=dh\,d\nu_a(q),\qquad
             x=h\cdot s_a(q),                                \tag{52.9}
\]
one has
\[
 \int_Qf\,d\nu_Q
 =\sum_a\int_{U_a}f(q)\rho_a(q)\,d\nu_Q(q).                   \tag{52.10}
\]
The result is independent of the sections, cover, and partition.
\end{theorem}

\begin{proof}
In a principal trivialization, invariance of \(d\mu_C\) implies that
its conditional measure on each fibre is a scalar multiple of Haar
measure.  The scalar defines the base measure; comparing pushforwards
shows that it is \(d\nu_Q\).  Thus (52.9) holds with the Haar factor
and quotient normalization consistent with (52.2).  The factor
\(\operatorname{vol}H\,\delta_e\) in (52.8) removes the normalized
fibre average and evaluates the invariant integrand at the section.
It gives
\[
 \int_{U_a}f(q)\rho_a(q)\,d\nu_Q(q).
\]
Summing and using (52.6) proves (52.10), hence (52.8).
Independence follows because every choice equals the canonical
left-hand side of (52.3).
\end{proof}

If normalized Haar probability is used, the two displayed
\(\operatorname{vol}H\) factors are both one.

\subsection{Recovery of local Faddeev--Popov formulas}

Suppose on \(\pi^{-1}(U_a)\) a smooth gauge condition
\(\chi_a:C\to\mathbb R^r\) has \(N_a\) simple intersections with each
orbit meeting \(\operatorname{supp}\rho_a\).  Put
\[
                        M_a=D\chi_aR.                          \tag{52.11}
\]

\begin{corollary}[Partitioned Faddeev--Popov formula]
\label{cor:v11v52-fp}
Under the preceding hypotheses,
\[
 \boxed{
 \int_Qf\,d\nu_Q
 =\sum_{a=1}^N\frac1{N_a}
  \int_C F(x)\rho_a(\pi(x))
       \delta_0(\chi_a(x))|\det M_a(x)|\,d\mu_C(x).}          \tag{52.12}
\]
If \(d\mu_C=w\,d\sigma_C\), the corresponding induced-slice formula is
\[
 \int_Qf\,d\nu_Q
 =\sum_{a=1}^N\frac1{N_a}
  \int_{\Sigma_a}
  f\,\rho_a\,w\,
  \frac{|\det M_a|}{J_{\chi_a}^C}\,d\sigma_{\Sigma_a},        \tag{52.13}
\]
where \(\Sigma_a=\{\chi_a=0\}\cap\pi^{-1}(U_a)\).
\end{corollary}

\begin{proof}
Apply Theorem~\ref{thm:v11v48-fp} to the invariant integrand
\(F(\rho_a\circ\pi)\) on each chart and divide by the fixed copy count
\(N_a\).  Sum over \(a\).  The coarea identity on \(C\) gives
(52.13).
\end{proof}

\begin{warning}[Copy counts belong to the density]
\label{warn:v11v52-copies}
If the number of simple intersections changes inside a chart, then a
zero has collided, escaped through the chart boundary, or crossed a
nontransverse point.  Formula (52.12) cannot retain a constant
\(N_a\) across that event.  Shrinking the chart, changing the section,
or treating the transition region separately is necessary.
\end{warning}

\subsection{Uniform avoidance of local horizons}

The compact support condition in (52.6) turns pointwise
transversality into a uniform bound.

\begin{proposition}[Finite-chart determinant floor]
\label{prop:v11v52-floor}
Suppose each \(M_a\) and \(B_a=D\chi_a|_{TC}\) is continuous and
surjective on the compact set
\[
 {\cal Z}_a
 :=\{x:\chi_a(x)=0,\ \pi(x)\in\operatorname{supp}\rho_a\}.    \tag{52.14}
\]
Then there are constants \(\kappa_{M,a},\kappa_{B,a}>0\) such that
\[
 s_{\min}(M_a)\ge\kappa_{M,a},\qquad
 s_{\min}(B_a)\ge\kappa_{B,a}\quad\text{on }{\cal Z}_a.       \tag{52.15}
\]
For finitely many charts, their minima
\[
 \kappa_M=\min_a\kappa_{M,a},\qquad
 \kappa_B=\min_a\kappa_{B,a}                                 \tag{52.16}
\]
are positive.  Hence the slice densities in (52.13) have no
Faddeev--Popov or slice-coarea singularity on the partition supports.
\end{proposition}

\begin{proof}
The least singular value is continuous in the matrix entries.
It is strictly positive at every point by the hypotheses.  A positive
continuous function on a compact set has a positive minimum.
Taking the minimum over finitely many charts preserves positivity.
The determinant bounds follow from
Lemma~\ref{lem:v11v49-envelope}.
\end{proof}

\begin{corollary}[Repair of a single-chart horizon]
\label{cor:v11v52-repair}
On a compact free quotient, a zero of one \(M_a\) outside
\({\cal Z}_a\) does not enter the quotient integral.  If the partition
is subordinate to charts satisfying
Proposition~\ref{prop:v11v52-floor}, the complete formula
(52.12) remains nonsingular even though no single gauge condition is
transverse everywhere.
\end{corollary}

\begin{proof}
The factor \(\rho_a\circ\pi\) vanishes outside the relevant chart
support.  On its support, (52.15) gives transversality.  The finite sum
therefore contains no singular local density.
\end{proof}

This is the legitimate replacement for multiplying a divergent inverse
by a small horizon probability.  It covers the complete compact regular
chamber with nonsingular local formulas.

\subsection{Overlap compatibility}

On \(U_a\cap U_b\), two sections differ by a smooth transition map
\(t_{ab}:U_a\cap U_b\to H\):
\[
                         s_b(q)=t_{ab}(q)\cdot s_a(q).         \tag{52.17}
\]

\begin{proposition}[Gauge-invariant gluing]
\label{prop:v11v52-gluing}
If \(F=f\circ\pi\), then
\[
                         F(s_a(q))=F(s_b(q))=f(q).             \tag{52.18}
\]
Consequently all overlap dependence is confined to the local density
representatives, and the partitioned integral is unchanged under a
refinement of the cover or a change of sections.
\end{proposition}

\begin{proof}
Both sections project to \(q\), proving (52.18).  A refinement replaces
\(\rho_a\) by functions whose sum is \(\rho_a\), so (52.10) is
unchanged.  A section change acts within each fibre and cannot alter
the quotient function or the invariant conditional Haar measure.
\end{proof}

For gauge-covariant fields rather than scalar observables, (52.17)
becomes an actual transition function.  Their local representatives
must satisfy the associated-bundle cocycle; scalar gluing alone does
not construct them.

\subsection{Constraint and stationary-phase compatibility}

The quotient construction occurs after imposing the regular constraint
surface of V47.  On each slice, the physical tangent is
\[
 T_xC=
 T_x\Sigma_a\oplus\operatorname{ran}R_x.                      \tag{52.19}
\]
A stationary-phase Hessian must first use the multiplier-corrected
constraint Hessian \(H_C\) from (47.18), then restrict or transport it
to a quotient complement.  Changing charts conjugates the physical
quadratic form by the derivative of the transition map, so its
signature and nondegeneracy are invariant; its coordinate determinant
also carries the corresponding metric Jacobian.

\begin{ledger}[Multi-chart adjacent data]
\label{led:v11v52-adjacent}
For comparison of cutoffs \(i\) and \(i+1\), the quotient atlas adds:
\begin{enumerate}
\item a common refinement of the two quotient covers on the retained
cylinder variables;
\item partition increments
\(\sum_a\lVert\rho_{a,i+1}-\rho_{a,i}\rVert\);
\item transition-map derivative bounds on overlaps;
\item stable local copy counts \(N_{a,i}\);
\item the floors \(\kappa_{M,a,i},\kappa_{B,a,i}\);
\item boundary mass for the part not covered by the regular compact
chamber.
\end{enumerate}
The first five are deterministic atlas stocks.  Only the last is a
bad-set probability row.
\end{ledger}

\begin{firewall}[Where the theorem stops]
\label{fire:v11v52-boundary}
Compactness and freeness are essential here.  If stabilizers occur,
\(Q\) is stratified and the principal-bundle atlas fails at the
singular strata.  If \(H\) is noncompact, there is no finite
\(\operatorname{vol}H\) normalization in (52.2).  The full
diffeomorphism group has both difficulties in relevant sectors.  V52
therefore closes local gauge horizons only on a compact regular
finite-cutoff chamber; it does not construct the global gravitational
quotient.
\end{firewall}

\begin{decision}[Next acquisition]
\label{dec:v11v52-next}
V53 will treat noncompact gauge volume by an exhaustion criterion.
Rather than assign a formal infinite constant, it will state sufficient
conditions under which normalized orbit-window integrals converge and
are independent of the exhaustion.  It will also show why this
criterion fails when an invariant observable is integrated against an
unfixed infinite orbit without a quotient disintegration.
\end{decision}

\begin{assessment}[Progress after V52]
\label{ass:v11v52-status}
On a compact free regular chamber, Gribov horizons of individual gauge
coordinates no longer force an inverse-determinant bad set.  A finite
partition of unity glues nonsingular local Faddeev--Popov formulas into
the canonical quotient measure.  Stabilizer strata, noncompact orbit
volume, and the chamber boundary remain genuine unresolved rows.
\end{assessment}

\subsection{Parent record}

This V52 note was formed by appending the preceding material to the
validated V51 source, whose record was
\[
\begin{array}{c|c}
\text{lines}&16321\\
\text{bytes}&496618\\
\text{SHA--256}&
\texttt{EEDE9939B634DF67D080CADF0107B0A218FDF1466C025D611F957AE0A18B621A}.
\end{array}                                                     \tag{52.20}
\]

% =============================================================================
% END APPENDED ELEVENTH-SERIES V52 MATERIAL
% =============================================================================
% =============================================================================
% BEGIN APPENDED ELEVENTH-SERIES V53 MATERIAL
% =============================================================================

\section{Exact removal of finite-dimensional noncompact orbit volume}
\label{sec:v11v53}

Formal division by an infinite gauge-group volume is not a definition.
For a free proper action of a finite-dimensional unimodular group, an
exact finite-window identity replaces it.

\subsection{Invariant disintegration}

Let \(H\) be a locally compact second-countable unimodular group with
left Haar measure \(dh\).  Let it act freely and properly on \(C\), and
assume \(Q=C/H\) is compact.  Let \(d\mu_C\) be an invariant
sigma-finite Radon measure with local disintegrations
\[
 d\mu_C(h\cdot s_a(q))=dh\,d\nu_a(q).                         \tag{53.1}
\]

\begin{lemma}[Compatibility of base measures]
\label{lem:v11v53-base}
The measures \(d\nu_a\) agree on overlaps and define a global measure
\(d\nu_Q\), up to the inverse of the overall normalization of \(dh\).
\end{lemma}

\begin{proof}
On an overlap, \(s_b(q)=t_{ab}(q)\cdot s_a(q)\), so the two fibre
coordinates differ by right translation by \(t_{ab}(q)\).
Unimodularity makes left Haar measure right invariant.  Comparing the
two representations of \(d\mu_C\) gives \(d\nu_a=d\nu_b\).
Rescaling Haar measure produces the stated common inverse rescaling.
\end{proof}

\subsection{Partitioned orbit windows}

Choose the finite bundle atlas and partition \(\rho_a\) of V52.  Let
\(W\subset H\) be measurable with
\[
                         0<|W|_H:=\int_H\mathbf1_W\,dh<\infty. \tag{53.2}
\]
For \(x\in\pi^{-1}(U_a)\), write
\[
                         x=h_a(x)\cdot s_a(\pi(x))             \tag{53.3}
\]
and define
\[
 w_W(x)=\sum_a\rho_a(\pi(x))\mathbf1_W(h_a(x)),\qquad
 \widehat w_W=|W|_H^{-1}w_W.                                 \tag{53.4}
\]

\begin{theorem}[Constant orbit mass]
\label{thm:v11v53-orbit-mass}
For every \(x\in C\),
\[
 \int_Hw_W(g\cdot x)\,dg=|W|_H,\qquad
 \int_H\widehat w_W(g\cdot x)\,dg=1.                         \tag{53.5}
\]
\end{theorem}

\begin{proof}
For every active chart,
\(h_a(g\cdot x)=g h_a(x)\).  Unimodularity and right translation give
\[
 \int_H\mathbf1_W(g h_a(x))\,dg=|W|_H.
\]
Multiply by \(\rho_a(\pi(x))\), sum, and use
\(\sum_a\rho_a=1\).
\end{proof}

\begin{theorem}[Window-independent quotient integral]
\label{thm:v11v53-window}
If \(F=f\circ\pi\) and \(f\in L^1(Q,\nu_Q)\), then
\[
 \boxed{\int_CF(x)\widehat w_W(x)\,d\mu_C(x)
                       =\int_Qf(q)\,d\nu_Q(q).}               \tag{53.6}
\]
The result is independent of \(W\), the sections, and the partition.
\end{theorem}

\begin{proof}
Using (53.1) and (53.4),
\[
\begin{split}
 \int_CF\widehat w_W\,d\mu_C
 &=\sum_a\int_{U_a}f(q)\rho_a(q)
   \left[|W|_H^{-1}\int_H\mathbf1_W(h)\,dh\right]d\nu_Q(q)\\
 &=\sum_a\int_{U_a}f(q)\rho_a(q)\,d\nu_Q(q)
 =\int_Qf\,d\nu_Q.
\end{split}                                                    \tag{53.7}
\]
Every admissible choice equals the same final integral.
\end{proof}

\begin{corollary}[Normalized expectation]
\label{cor:v11v53-normalized}
If \(S=s\circ\pi\) and
\(0<\int_Qe^{-s}d\nu_Q<\infty\), then
\[
 \frac{\int_CF e^{-S}w_W\,d\mu_C}
      {\int_Ce^{-S}w_W\,d\mu_C}
 =\frac{\int_Qf e^{-s}\,d\nu_Q}
       {\int_Qe^{-s}\,d\nu_Q}.                               \tag{53.8}
\]
Thus finite orbit windows cancel exactly.
\end{corollary}

\begin{proof}
Apply Theorem~\ref{thm:v11v53-window} to numerator and denominator.
The common factor \(|W|_H\) cancels.
\end{proof}

\subsection{Why the unfixed integral diverges}

\begin{proposition}[Infinite invariant fibre]
\label{prop:v11v53-divergence}
If \(H\) has infinite Haar volume, \(f\ge0\), and
\(\int_Qf\,d\nu_Q>0\), then
\[
                         \int_C(f\circ\pi)\,d\mu_C=\infty.     \tag{53.9}
\]
\end{proposition}

\begin{proof}
The integrand is constant on every fibre.  Tonelli's theorem inserts
the factor \(\int_Hdh=\infty\) over a base set of positive integral.
\end{proof}

An unfixed ratio of two infinite integrals therefore does not define
(53.8).  Quotient disintegration or a finite orbit window must be
inserted before normalization.

\subsection{Agreement with local Faddeev--Popov gauges}

Combining V52 with Theorem~\ref{thm:v11v53-window} gives
\[
 \int_Qf\,d\nu_Q
 =\int_CF\widehat w_W\,d\mu_C
 =\sum_a\frac1{N_a}\int_C
  F(\rho_a\circ\pi)\delta_0(\chi_a)|\det M_a|\,d\mu_C.        \tag{53.10}
\]
The window formula and the local transverse-delta formula are two
representations of the same quotient measure.  Neither contains an
infinite group-volume constant.

\begin{proposition}[No adjacent orbit-volume row]
\label{prop:v11v53-adjacent}
Changing \(W\) at a fixed cutoff does not change a normalized
gauge-invariant expectation.  Across adjacent cutoffs, orbit-window
volumes contribute zero once both measures are reduced to quotient
measures and the retained cylinder observable is gauge invariant.
\end{proposition}

\begin{proof}
The first assertion is
Corollary~\ref{cor:v11v53-normalized}.  The quotient expressions on the
right of (53.8) contain no window or window volume, so neither can
enter their adjacent difference.
\end{proof}

\subsection{Application boundary}

A finite-dimensional gravity regulator may use a noncompact
lapse-shift or frame gauge group.  The theorem applies only after
checking freeness, properness, unimodularity, compactness of the
selected quotient chamber, and invariant disintegration of its
measure.

\begin{firewall}[Continuum diffeomorphisms]
\label{fire:v11v53-continuum}
The full smooth diffeomorphism group is infinite dimensional and is not
a locally compact group carrying the Haar measure used here.
Gravitational configurations can have isometries, so the action need
not be free, and properness can fail.  V53 is a finite-dimensional
regulator theorem; it does not construct the continuum
diffeomorphism quotient.
\end{firewall}

\begin{firewall}[Modular correction]
\label{fire:v11v53-modular}
For a nonunimodular group, right translation in the proof introduces
the modular function.  A constant orbit window then requires a
modularly corrected density.  Formula (53.6) is invalid without that
correction.
\end{firewall}

\begin{ledger}[Remaining quotient acquisitions]
\label{led:v11v53-remaining}
After the finite orbit volume is removed, the regulated Einstein
quotient still requires:
\begin{enumerate}
\item regular constraints and the coarea convention of V47;
\item first-class closure and local transversality from V48;
\item determinant thresholds and certified gaps from V49--V51;
\item a finite nonsingular atlas as in V52;
\item verification of the V53 group and measure hypotheses;
\item control of stabilizer and rank-deficient strata;
\item adjacent convergence of the resulting quotient measures.
\end{enumerate}
Under the first five items, noncompact finite-dimensional orbit volume
is no longer a probabilistic error.
\end{ledger}

\begin{decision}[Next acquisition]
\label{dec:v11v53-next}
V54 will address singular strata.  It will prove a tubular
codimension estimate for their reduced measure and show separately why
codimension alone does not control inverse determinants near a
rank-deficient set.
\end{decision}

\begin{assessment}[Progress after V53]
\label{ass:v11v53-status}
For a free proper unimodular finite-dimensional gauge action with
compact quotient, infinite orbit volume can be removed exactly and
independently of the chosen finite window.  The unresolved geometric
row is the behavior near nonfree or rank-deficient strata, followed by
the actual adjacent quotient limit.
\end{assessment}

\subsection{Parent record}

This V53 note was formed by appending the preceding material to the
validated V52 source, whose record was
\[
\begin{array}{c|c}
\text{lines}&16647\\
\text{bytes}&508669\\
\text{SHA--256}&
\texttt{C962B9F3349201B746DBB4C07ECFE87DA54C2C66DF67A8D6FB53FD5B99ED796C}.
\end{array}                                                     \tag{53.11}
\]

% =============================================================================
% END APPENDED ELEVENTH-SERIES V53 MATERIAL
% =============================================================================
% =============================================================================
% BEGIN APPENDED ELEVENTH-SERIES V54 MATERIAL
% =============================================================================

\section{Singular strata: tubular mass, vanishing order, and inverse
determinants}
\label{sec:v11v54}

V52--V53 close regular finite-dimensional quotient charts, including
noncompact orbit volume under explicit hypotheses.  The next question
is whether configurations with stabilizers or deficient constraint
rank can be discarded in measure and, more strongly, whether inverse
Jacobians remain integrable near them.  These are different questions.
Positive codimension answers the first under absolute continuity.
The second also requires a quantitative vanishing order.

\subsection{Tubular mass of a smooth stratum}

Let \((Q,g)\) be an \(n\)-dimensional Riemannian manifold and let
\(S\subset Q\) be a compact embedded submanifold of codimension
\(c\ge1\).  Write
\[
                         T_r(S)=\{x\in Q:d(x,S)<r\}.            \tag{54.1}
\]

\begin{theorem}[Tubular volume bound]
\label{thm:v11v54-tube}
There are \(r_0>0\) and \(C_S<\infty\) such that
\[
                         \operatorname{Vol}_Q(T_r(S))
                              \le C_S r^c,\qquad 0<r<r_0.      \tag{54.2}
\]
If \(d\mu=w\,dV_Q\) with \(w\in L^p(Q)\), \(1<p\le\infty\), then
\[
 \mu(T_r(S))
 \le \lVert w\rVert_{L^p}C_S^{\,1-1/p}
                         r^{\,c(1-1/p)}.                      \tag{54.3}
\]
For \(p=\infty\), the exponent is \(c\).
\end{theorem}

\begin{proof}
The tubular-neighborhood theorem identifies \(T_{r_0}(S)\) with a
normal disk bundle by the normal exponential map.  Its Jacobian is
continuous and bounded on the compact disk bundle of radius \(r_0\).
In each fibre, the normal ball has \(c\)-dimensional volume proportional
to \(r^c\).  Integrating over compact \(S\) proves (54.2).
Hölder's inequality gives
\[
 \mu(T_r(S))\le\lVert w\rVert_p
              \operatorname{Vol}(T_r(S))^{1-1/p},
\]
which is (54.3).
\end{proof}

\begin{corollary}[Finite stratification]
\label{cor:v11v54-stratified}
If \({\cal S}=\bigcup_{a=1}^N S_a\) is a finite union of compact
embedded strata and \(c_*=\min_a\operatorname{codim}S_a\), then
\[
                     \mu(T_r({\cal S}))\le C r^\beta,\qquad
              \beta=c_*(1-1/p),                              \tag{54.4}
\]
for sufficiently small \(r\), with the usual value
\(\beta=c_*\) when \(p=\infty\).
\end{corollary}

\begin{proof}
Use \(T_r({\cal S})\subset\bigcup_aT_r(S_a)\), apply
Theorem~\ref{thm:v11v54-tube}, and sum the finite constants.
\end{proof}

For \(p=1\), absolute continuity still implies
\(\mu(T_r(S))\to0\), but no positive power rate follows from an
\(L^1\) density alone.

\subsection{From distance to a singular-value tail}

Let \(L(x)\) be a continuous matrix family on a neighborhood of
\({\cal S}\), with \(s(x)=s_{\min}(L(x))\).

\begin{theorem}[Vanishing-order transfer]
\label{thm:v11v54-order}
Assume that for some \(a,\nu,r_0,s_0>0\),
\[
 s(x)\ge a\,d(x,{\cal S})^\nu
       \quad\text{when }d(x,{\cal S})<r_0,\qquad
 s(x)\ge s_0
       \quad\text{outside }T_{r_0}({\cal S}).                 \tag{54.5}
\]
If (54.4) holds, then for sufficiently small \(\epsilon\),
\[
             \mu\{s_{\min}(L)\le\epsilon\}
                  \le C a^{-\beta/\nu}
                         \epsilon^{\beta/\nu}.                \tag{54.6}
\]
Thus the small-ball exponent acquired from the stratum is
\[
                              \alpha=\frac{\beta}{\nu}.        \tag{54.7}
\]
\end{theorem}

\begin{proof}
For
\(\epsilon<\min\{s_0,ar_0^\nu\}\), (54.5) implies
\[
 \{s\le\epsilon\}\subset
 T_{(\epsilon/a)^{1/\nu}}({\cal S}).
\]
Apply (54.4) at this radius.
\end{proof}

\begin{corollary}[Inverse Jacobian criterion]
\label{cor:v11v54-inverse}
Let \(L:H\to\mathbb R^k\) be surjective off \({\cal S}\) and obey
Theorem~\ref{thm:v11v54-order}.  The preceding information guarantees
\[
                         \mathbb E[J(L)^{-q}]<\infty          \tag{54.8}
\]
whenever
\[
                              kq<\frac{\beta}{\nu}.            \tag{54.9}
\]
\end{corollary}

\begin{proof}
Use the small-ball exponent (54.7) in
Theorem~\ref{thm:v11v49-inverse-moment}.
\end{proof}

This criterion applies separately to the constraint coarea matrix
\(A\), the slice matrix \(B\), the Faddeev--Popov matrix \(M\), and the
constrained Hessian.  Their singular sets and vanishing orders need
not agree.

\subsection{Sharp local obstruction}

\begin{proposition}[Divergence at excessive vanishing order]
\label{prop:v11v54-divergence}
Suppose \(S\) is a codimension-\(c\) compact stratum, the density
satisfies \(w\ge w_0>0\) on a tube, and
\[
                         s_{\min}(L(x))
                              \le b\,d(x,S)^\nu.               \tag{54.10}
\]
Then
\[
                         \int s_{\min}(L)^{-u}\,d\mu=\infty   \tag{54.11}
\]
whenever
\[
                               \nu u\ge c.                     \tag{54.12}
\]
\end{proposition}

\begin{proof}
Tubular coordinates have volume density bounded below after reducing
\(r_0\).  Therefore the integral in (54.11) is bounded below by a
positive constant times
\[
                         \int_0^{r_0}r^{c-1-\nu u}\,dr.
\]
This integral diverges exactly when \(c-1-\nu u\le-1\), proving
(54.12).
\end{proof}

\begin{proposition}[Codimension alone gives no inverse moment]
\label{prop:v11v54-flat}
For every \(c\ge1\), there is a smooth nonnegative scalar matrix family
with zero set a codimension-\(c\) submanifold and with no positive
inverse moment locally.
\end{proposition}

\begin{proof}
In tubular coordinates \((y,z)\), \(z\in\mathbb R^c\), set
\[
 s(y,z)=
 \begin{cases}
 \exp(-|z|^{-2}),&z\ne0,\\
 0,&z=0.
 \end{cases}                                                   \tag{54.13}
\]
This is smooth and vanishes precisely at \(z=0\).  For every \(u>0\),
the inverse integral contains
\[
                    \int_0^{r_0}e^{u/r^2}r^{c-1}\,dr=\infty. \tag{54.14}
\]
Thus even arbitrarily high codimension does not control inverse
moments without a finite-order lower bound such as (54.5).
\end{proof}

The corresponding small-ball mass decays only as a negative power of
\(\log(1/\epsilon)\), rather than as a power of \(\epsilon\).

\subsection{A feasible tube schedule}

Suppose the bad observable has \(L^q\) envelope
\(D_i\lesssim(i+2)^b\), \(q>1\), the native good-set error is
\(\delta_i\lesssim(i+2)^{-s}\), and a matrix stock costs
\(d(x,{\cal S}_i)^{-\nu\ell}\).  Assume the tube mass is uniformly
bounded by \(Cr^\beta\).  Choose
\[
                             r_i=(i+2)^{-a}.                   \tag{54.15}
\]

\begin{theorem}[Tube-threshold interval]
\label{thm:v11v54-schedule}
The good-set and bad-tube series both converge if
\[
 \frac{b+1}{\beta(1-1/q)}<a<
                              \frac{s-1}{\nu\ell}.             \tag{54.16}
\]
Such an \(a\) exists exactly when
\[
              \beta(1-1/q)(s-1)>
                              \nu\ell(b+1).                    \tag{54.17}
\]
\end{theorem}

\begin{proof}
Outside the tube the good summand is at most
\[
                 C(i+2)^{-s+a\nu\ell},
\]
which is summable under the right inequality.  Hölder and the tube
bound give a bad summand at most
\[
                 C(i+2)^{b-a\beta(1-1/q)},
\]
summable under the left inequality.  The interval is nonempty exactly
under (54.17).
\end{proof}

For an inverse \(k\)-Jacobian absolute increment, V49 has
\(\ell=k+1\).  For a logarithmic determinant increment, \(\ell=1\).
The feasibility gap can therefore be much narrower for the unnormalized
density than for its logarithm.

\subsection{Meaning for stabilizers and constraint rank}

A positive-codimension stabilizer stratum has zero reduced measure
under the hypotheses of Theorem~\ref{thm:v11v54-tube}.  This is enough
for bounded gauge-invariant observables.  It does not automatically
produce a principal-bundle quotient across the stratum, nor does it
control determinants used in unbounded observables or stationary
phase.  Those require the vanishing-order card.

\begin{ledger}[Singular-stratum acquisition]
\label{led:v11v54-strata}
For each finite Einstein cutoff one must now identify:
\begin{enumerate}
\item a finite or locally finite Whitney-type stratification of the
rank-deficient and stabilizer locus;
\item the codimension of each relevant stratum;
\item an \(L^p\) bound for the reduced density in tubular coordinates;
\item a lower vanishing-order estimate (54.5) for every required
matrix;
\item uniform tube constants and orders as the cutoff grows;
\item a tube schedule satisfying (54.17).
\end{enumerate}
Measure zero proves only the bounded-observable part of this list.
\end{ledger}

\begin{firewall}[No automatic algebraic order]
\label{fire:v11v54-algebraic}
Finite regulator matrices may be analytic or polynomial in chosen
coordinates, but a uniform lower bound of the form (54.5) still needs
a quantitative Lojasiewicz exponent and constant.  Neither analyticity
nor codimension supplies cutoff-uniform values automatically.
No such bound has yet been proved for the regulated gravitational
constraint or Faddeev--Popov matrices.
\end{firewall}

\begin{decision}[Mandatory V55 audit]
\label{dec:v11v54-next}
The next version is the compulsory companion audit.  It will inspect
the newest YM and NS notes and determine whether their current
concentration, rank-one, or finite-cell calculations improve the
tube exponent, the good-chamber probability, or the adjacent
determinant schedule.  V56 will then make a substantive acquisition
using that audit.
\end{decision}

\begin{assessment}[Progress after V54]
\label{ass:v11v54-status}
Singular strata are no longer represented by a bare measure-zero
assertion.  Their codimension and density regularity yield an explicit
tube exponent, while a separate vanishing order converts that exponent
to a singular-value tail.  The exact condition for inverse
integrability and the adjacent tube schedule are now visible.
\end{assessment}

\subsection{Parent record}

This V54 note was formed by appending the preceding material to the
validated V53 source, whose record was
\[
\begin{array}{c|c}
\text{lines}&16884\\
\text{bytes}&516923\\
\text{SHA--256}&
\texttt{C8EC5E48265C387FF43C1A2E3AD781C9C0DA6341964CEB0D4D75878760B88CD4}.
\end{array}                                                     \tag{54.18}
\]

% =============================================================================
% END APPENDED ELEVENTH-SERIES V54 MATERIAL
% =============================================================================
% =============================================================================
% BEGIN APPENDED ELEVENTH-SERIES V55 MATERIAL
% =============================================================================

\section{Companion audit at V55: unchanged records and the analytic
singular-locus route}
\label{sec:v11v55}

The compulsory audit finds no new companion version after V50.  Its
value is therefore diagnostic: the existing Yang--Mills chart theorem
is compared with the new singular-stratum requirements of V54.

\subsection{Audited records}

\begin{definition}[V55 companion record]
\label{def:v11v55-record}
The files inspected are
\[
\begin{array}{c|r|l}
\text{file}&\text{bytes}&\text{SHA--256}\\ \hline
\texttt{Renewal\_Yang\_Mills\_Mass\_Gap\_Research\_Notes\_V35.tex}
 &498660&
\texttt{23F674C142CA9B298546508F7A3E9EAF2CE2E154078C32B0BA2FC8660AD25098}\\
\texttt{Renewal\_NS\_Stability\_Research\_Notes\_V26.tex}
 &278283&
\texttt{82C887F8C2C69E4F28FAD7C3DC8DE5B9099CB5A4BF431EBA1FFF3E91935978AD}.
\end{array}                                                     \tag{55.1}
\]
Both byte counts and digests equal those in the V50 audit.  Hence no
new YM or NS theorem is imported at V55.
\end{definition}

\subsection{Effect on the V54 acquisition}

The YM V35 theorem still supplies the summable principal-link-chart
event on contractible blocks.  Inside that event, every compact-group
link can be written
\[
                             U_e=\exp X_e                    \tag{55.2}
\]
with \(X_e\) in a fixed bounded Lie-algebra ball.  Finite products,
adjoint actions, Wilson words, and derivatives of finite regulated
actions are real analytic functions of these coordinates.  The
positive metric chamber of V39 is likewise a compact subset of the
open positive cone, where inversion and determinant are analytic.

\begin{proposition}[Fixed-cutoff analytic reduction]
\label{prop:v11v55-analytic}
At a fixed finite cutoff, suppose a constraint, coarea,
Faddeev--Popov, or Hessian matrix is assembled by finitely many sums,
products, inverses with denominators bounded away from zero, matrix
exponentials, and derivatives of analytic regulator actions.  On a
compact product of the V35 gauge chart and V39 metric chamber, its
entries are real analytic in finitely many real coordinates.
\end{proposition}

\begin{proof}
Finite sums, products, and derivatives preserve real analyticity.
The matrix exponential is entire.  Matrix inversion is analytic on
the open set where the determinant is nonzero, and the assumed
denominator floor keeps the compact chamber inside that open set.
Composition proves the claim.
\end{proof}

This observation does not itself give a singular-value exponent, but
it excludes the arbitrary smooth flat example (54.13) once the matrix
family and chamber satisfy the stated analytic hypotheses.  The
appropriate next tool is the finite-dimensional Lojasiewicz distance
inequality for the analytic function
\[
                    \det(L(x)L(x)^*)
       =\sum_I|\det L_I(x)|^2,                                \tag{55.3}
\]
where the sum is over maximal minors.

\begin{transfer}[What the companions now provide]
\label{trans:v11v55}
The YM chart supplies probability of entering analytic compact link
coordinates.  The NS rank-one calculation supplies no information
about the zero set of the Einstein constraint or Faddeev--Popov
matrices.  Neither companion supplies cutoff-uniform Lojasiewicz
constants.  Thus the useful transfer is the analytic coordinate
domain, not a singular-value exponent copied from another programme.
\end{transfer}

\begin{decision}[V56 acquisition]
\label{dec:v11v55-next}
V56 will apply the Lojasiewicz distance inequality to (55.3).  It will
prove that every fixed-cutoff analytic surjectivity matrix has a
finite vanishing-order lower bound for its least singular value on a
compact chamber, derive the resulting tube and inverse-Jacobian
exponents, and isolate uniformity of the constants and exponent as the
remaining continuum problem.
\end{decision}

\begin{firewall}[Audit boundary]
\label{fire:v11v55-boundary}
The unchanged YM concentration theorem controls entry into a
contractible local link chart; it does not control global holonomy or
the gravitational matrix rank.  Fixed-cutoff analyticity does not
produce a cutoff-uniform exponent, and it does not show that the
rank-deficient set has positive codimension.  The unchanged NS result
adds no gravitational reduction estimate.
\end{firewall}

\begin{assessment}[Progress after V55]
\label{ass:v11v55-status}
The mandatory audit is complete.  It does not repeat either companion
calculation.  It identifies a new upstream bridge from their existing
chart control to V54: analytic maximal minors can supply a finite
vanishing order at each fixed cutoff.
\end{assessment}

\subsection{Parent record}

This V55 note was formed by appending the preceding material to the
validated V54 source, whose record was
\[
\begin{array}{c|c}
\text{lines}&17192\\
\text{bytes}&527526\\
\text{SHA--256}&
\texttt{607656AC8994C06205FE86D99950C9E49CB7C059FB4D0336D50E7859FDC0A50D}.
\end{array}                                                     \tag{55.4}
\]

% =============================================================================
% END APPENDED ELEVENTH-SERIES V55 MATERIAL
% =============================================================================
% =============================================================================
% BEGIN APPENDED ELEVENTH-SERIES V56 MATERIAL
% =============================================================================

\section{Analytic rank loci and fixed-cutoff Lojasiewicz exponents}
\label{sec:v11v56}

V54 requires a finite vanishing order near rank-deficient strata.
V55 identifies real-analytic finite-cutoff coordinates.  The present
section joins those statements.  It also replaces projector-defined
matrices by analytic carrier matrices, avoiding a circular use of
\((AA^*)^{-1}\) near the set where \(A\) loses rank.

\subsection{The analytic distance inequality}

\begin{lemma}[Compact Lojasiewicz distance inequality]
\label{lem:v11v56-loja}
Let \(\Omega\subset\mathbb R^N\) be open, let
\(f:\Omega\to[0,\infty)\) be real analytic, and let
\(K\Subset\Omega\) be compact.  Put \(Z=f^{-1}(0)\).  There are
\(C>0\) and \(\vartheta>0\) such that
\[
 f(x)\ge C\,\min\{1,d(x,Z)\}^{\vartheta},
                         \qquad x\in K.                       \tag{56.1}
\]
If \(Z\cap K=\varnothing\), then \(f\) has a positive minimum on
\(K\).
\end{lemma}

\begin{proof}
At each \(z\in Z\cap K\), the local Lojasiewicz distance theorem for a
real-analytic function gives a neighborhood \(U_z\), a constant
\(C_z>0\), and a finite exponent \(\vartheta_z>0\) with
\[
 f(x)\ge C_zd(x,Z)^{\vartheta_z}\quad(x\in U_z).
\]
Compactness of \(Z\cap K\) gives a finite subcover.  Increase all
exponents to their maximum and reduce the constants, after shrinking
the neighborhoods so that the distance is at most one.  On the compact
remainder of \(K\), separated from \(Z\), continuity gives a positive
minimum.  Combining the two bounds proves (56.1).  The last assertion
is the same compact-minimum argument.
\end{proof}

The exponent in this theorem is finite at each fixed dimension.  The
proof gives no cutoff-uniform upper bound for it.

\subsection{Maximal minors and the least singular value}

Let \(L(x):E\to\mathbb R^k\), \(k\le\dim E\), have real-analytic matrix
entries.  Define
\[
 {\cal D}_L(x):=\det(L(x)L(x)^*)
              =\sum_{I}|\det L_I(x)|^2,                       \tag{56.2}
\]
where \(I\) runs over the \(k\)-column minors.  The second equality is
Cauchy--Binet, and
\[
                 {\cal D}_L(x)=0
             \quad\Longleftrightarrow\quad\operatorname{rank}L(x)<k.
                                                                  \tag{56.3}
\]

\begin{theorem}[Fixed-cutoff singular-value order]
\label{thm:v11v56-singular-order}
Let \(K\Subset\Omega\) be compact and suppose
\(\lVert L(x)\rVert\le B\) on \(K\), with \(B>0\).  Let
\[
                         Z_L=\{x:\operatorname{rank}L(x)<k\}. \tag{56.4}
\]
There are \(a>0\) and \(\nu<\infty\) such that
\[
                  s_{\min}(L(x))
       \ge a\,\min\{1,d(x,Z_L)\}^{\nu},\qquad x\in K.         \tag{56.5}
\]
If \(Z_L\cap K=\varnothing\), \(s_{\min}(L)\) has a positive uniform
floor on \(K\).
\end{theorem}

\begin{proof}
Apply Lemma~\ref{lem:v11v56-loja} to the analytic nonnegative function
\({\cal D}_L\).  Since
\[
 \sqrt{{\cal D}_L}
       =\prod_{j=1}^ks_j(L)
       \le B^{k-1}s_{\min}(L),                                \tag{56.6}
\]
inequality (56.1) gives
\[
 s_{\min}(L)
 \ge B^{1-k}\sqrt C\,
      \min\{1,d(x,Z_L)\}^{\vartheta/2}.
\]
Take \(a=B^{1-k}\sqrt C\) and \(\nu=\vartheta/2\).
If \(Z_L\cap K\) is empty, continuity of the least singular value gives
the positive floor.
\end{proof}

\begin{corollary}[Measure-zero rank locus]
\label{cor:v11v56-null}
Suppose \(\Omega_0\) is a connected component meeting \(K\) and
\(L(x_0)\) has full row rank at one point \(x_0\in\Omega_0\).
Then \({\cal D}_L\) is not identically zero on \(\Omega_0\), and its
zero set has Lebesgue measure zero and empty interior there.
\end{corollary}

\begin{proof}
Full rank at \(x_0\) gives \({\cal D}_L(x_0)>0\), so analyticity
prevents \({\cal D}_L\) from vanishing identically on the connected
component.  The zero set of a nonzero real-analytic function has
Lebesgue measure zero and empty interior.
\end{proof}

This corollary proves fixed-cutoff almost-everywhere regularity only
after one full-rank reference point is exhibited on every relevant
connected component.

\subsection{Analytic carriers for constrained reduction}

The tangent projector
\(Q=I-A^*(AA^*)^{-1}A\) is undefined on the singular set of \(A\).
It should not be used to prove a distance inequality across that set.

First consider the slice derivative.  Let
\[
 A:E\to\mathbb R^m,\qquad C:E\to\mathbb R^r,\qquad
 {\cal R}:=\begin{pmatrix}A\\C\end{pmatrix}:
 E\to\mathbb R^{m+r}.                                        \tag{56.7}
\]

\begin{lemma}[Stacked carrier for the slice]
\label{lem:v11v56-stacked}
If \(A\) is surjective, then the restriction
\(C|_{\ker A}:\ker A\to\mathbb R^r\) is surjective exactly when
\({\cal R}\) is surjective.
\end{lemma}

\begin{proof}
If \({\cal R}\) is surjective, solve
\({\cal R}u=(0,b)\) for arbitrary \(b\); then
\(u\in\ker A\) and \(Cu=b\).  Conversely, choose \(u_0\) with
\(Au_0=a\).  Surjectivity of \(C|_{\ker A}\) supplies
\(v\in\ker A\) with \(Cv=b-Cu_0\).  Then
\({\cal R}(u_0+v)=(a,b)\).
\end{proof}

Thus the analytic maximal minors of \({\cal R}\), together with those
of \(A\), control both constraint regularity and slice-coarea
regularity without forming \(Q\).

For the stationary Hessian, let
\[
 H=\operatorname{Hess}S-\sum_{a=1}^m\lambda_a
                         \operatorname{Hess}{\cal C}^a        \tag{56.8}
\]
be the multiplier-corrected ambient Hessian and define the bordered
KKT operator
\[
 {\cal K}:=
 \begin{pmatrix}
        H&A^*\\
        A&0
 \end{pmatrix}:E\oplus\mathbb R^m\to E\oplus\mathbb R^m.     \tag{56.9}
\]

For analytic continuation across the rank-deficient set, the
multiplier \(\lambda\) is treated as an independent auxiliary
coordinate and the stationarity equation
\(\nabla S-A^*\lambda=0\) is imposed jointly.  One must not solve for
\(\lambda\) using \((AA^*)^{-1}\) before applying the analytic
argument.

\begin{theorem}[Bordered carrier for constrained nondegeneracy]
\label{thm:v11v56-kkt}
Assume \(A\) is surjective and \(H=H^*\).  The constrained Hessian
\[
                         H_C=QH Q|_{\ker A}                   \tag{56.10}
\]
is invertible on \(\ker A\) exactly when \({\cal K}\) is invertible.
\end{theorem}

\begin{proof}
If \({\cal K}(u,\eta)=0\), then
\[
                         Au=0,\qquad Hu+A^*\eta=0.            \tag{56.11}
\]
Applying \(Q\) to the second equation gives \(H_Cu=0\).  If \(H_C\)
is invertible, \(u=0\); then \(A^*\eta=0\), and surjectivity of \(A\)
makes \(A^*\) injective, so \(\eta=0\).  Thus \({\cal K}\) is
injective and hence invertible.

Conversely, if \(H_C\) is singular, choose nonzero
\(u\in\ker A\) with \(QHu=0\).  Then
\(Hu\in(\ker A)^\perp=\operatorname{ran}A^*\), so there is
\(\eta\) with \(Hu+A^*\eta=0\).  Hence
\({\cal K}(u,\eta)=0\), and \({\cal K}\) is singular.
\end{proof}

The entries of \({\cal K}\) are analytic whenever \(A,H,\lambda\) are
analytic.  No inverse of \(AA^*\) occurs in (56.9).

\begin{corollary}[Analytic carrier family]
\label{cor:v11v56-carriers}
On an analytic regulator chart, fixed-cutoff failures of the following
conditions are zero sets of analytic maximal-minor functions:
\[
\begin{array}{c|c}
\text{condition}&\text{analytic carrier}\\ \hline
\text{constraint regularity}&A=D{\cal C}\\
\text{slice regularity}&{\cal R}=(A,D\chi)\\
\text{gauge transversality}&M=D\chi R\\
\text{constrained nondegeneracy}&{\cal K}\text{ of (56.9)}\\
\text{Dirac nondegeneracy}&\Delta\text{ of (48.15)}.
\end{array}                                                     \tag{56.12}
\]
Theorem~\ref{thm:v11v56-singular-order} applies to each carrier on a
compact chamber.
\end{corollary}

\begin{proof}
Analyticity follows from Proposition~\ref{prop:v11v55-analytic}.
The equivalences are Lemma~\ref{lem:v11v56-stacked},
Theorem~\ref{thm:v11v56-kkt}, and
Theorem~\ref{thm:v11v48-dirac}.  Apply
Theorem~\ref{thm:v11v56-singular-order}.
\end{proof}

\subsection{Tube and moment exponents}

Assume a carrier rank locus on a compact semianalytic chamber admits a
finite Whitney stratification whose smallest relevant codimension is
\(c_*\).  Such a stratification exists for a relatively compact
semianalytic set.  Let the reduced density belong to \(L^p\),
\(1<p\le\infty\), and put
\[
                         \beta=c_*(1-1/p).                    \tag{56.13}
\]
Let \(\vartheta_L\) be an exponent in (56.1) for
\({\cal D}_L\), so \(\nu_L=\vartheta_L/2\).

\begin{theorem}[Analytic small-ball exponent]
\label{thm:v11v56-small-ball}
Under these hypotheses,
\[
 \mu\{s_{\min}(L)\le\epsilon\}
              \le C_L\epsilon^{\alpha_L},\qquad
 \alpha_L=\frac{2\beta}{\vartheta_L}
          =\frac{\beta}{\nu_L},                              \tag{56.14}
\]
for sufficiently small \(\epsilon\).  If \(L\) has \(k\) rows, then
\[
                          \mathbb E[J(L)^{-q}]<\infty         \tag{56.15}
\]
is guaranteed whenever
\[
                         kq<\alpha_L.                         \tag{56.16}
\]
\end{theorem}

\begin{proof}
Theorem~\ref{thm:v11v56-singular-order} supplies the distance lower
bound with order \(\nu_L\).  The tubular estimate
Corollary~\ref{cor:v11v54-stratified} supplies exponent \(\beta\).
Apply Theorem~\ref{thm:v11v54-order}, then
Corollary~\ref{cor:v11v54-inverse}.
\end{proof}

\begin{warning}[The carrier dimension matters]
\label{warn:v11v56-carrier-dimension}
For \({\cal K}\), the determinant has size \(\dim E+m\), but the
physical stationary determinant is taken only on \(\ker A\).
The bordered determinant is used here to locate and quantify
nondegeneracy, not as a replacement for the physical determinant.
Conversion between their inverse moments also carries factors from
\(AA^*\), obtained by Schur-complement identities on the regular set.
Those factors must remain in the coarea ledger.
\end{warning}

\subsection{Uniformity is the continuum problem}

At every fixed cutoff \(i\), the preceding theorem produces constants
\[
                    C_i,\quad\vartheta_i,\quad c_i,\quad p_i. \tag{56.17}
\]
The master adjacent criterion requires quantitative control of these
sequences, not merely their existence.

\begin{ledger}[Uniform analytic acquisition]
\label{led:v11v56-uniform}
For each carrier \(A,{\cal R},M,{\cal K},\Delta\), a continuum proof
must establish:
\begin{enumerate}
\item one full-rank reference point in every relevant connected
component;
\item a compact semianalytic chamber with a controlled density;
\item a lower bound on stratum codimension;
\item an upper bound on the Lojasiewicz exponent \(\vartheta_i\);
\item a lower bound on its constant \(C_i\);
\item compatibility of these bounds with the V54 threshold interval;
\item a treatment of the chamber complement using V42.
\end{enumerate}
\end{ledger}

\begin{proposition}[Why fixed-cutoff existence is insufficient]
\label{prop:v11v56-nonuniform}
Even if every \(\vartheta_i\) is finite, the small-ball exponents
\(\alpha_i=2\beta_i/\vartheta_i\) can tend to zero.  In that case no
fixed inverse moment and no uniform polynomial threshold interval
follows from (56.14).
\end{proposition}

\begin{proof}
If \(\vartheta_i\to\infty\) while \(\beta_i\) stays bounded, then
\(\alpha_i\to0\).  The inverse criterion \(k_iq<\alpha_i\) and the
left side of the threshold condition (49.23) eventually fail for
every fixed positive \(q\) or threshold exponent.
\end{proof}

Effective real-algebraic bounds can make \(\vartheta_i\) explicit when
the carrier entries have uniformly controlled polynomial degree and
dimension.  In the present programme both dimension and stencil
complexity grow, so such a bound must be derived for the actual
regulator rather than quoted as uniform.

\begin{firewall}[What V56 proves]
\label{fire:v11v56-boundary}
V56 rules out infinitely flat rank collapse at each fixed analytic
cutoff and supplies analytic carriers that avoid inverse-projector
circularity.  It does not prove favorable codimension, a usable
cutoff-uniform Lojasiewicz exponent, a lower constant, or the
full-rank reference computations for the Einstein constraint system.
It therefore does not close the inverse-determinant or continuum
measure rows.
\end{firewall}

\begin{decision}[Next acquisition]
\label{dec:v11v56-next}
V57 will attack the uniformity problem through locality.  It will
derive a block-banded lower-bound theorem: a uniform local inf-sup
constant plus a controlled inter-block coupling gives a global
singular-value floor independent of the number of blocks.  This avoids
using a global Lojasiewicz exponent when the regulated constraint and
gauge matrices admit a stable local decomposition.
\end{decision}

\begin{assessment}[Progress after V56]
\label{ass:v11v56-status}
Every fixed-cutoff analytic rank problem now has a finite
distance-to-singular-set order, and the constrained slice and Hessian
have analytic carriers that remain defined on the singular locus.
The bottleneck has narrowed from existence of an exponent to uniform
local structure, reference ranks, and constants as the regulator
grows.
\end{assessment}

\subsection{Parent record}

This V56 note was formed by appending the preceding material to the
validated V55 source, whose record was
\[
\begin{array}{c|c}
\text{lines}&17321\\
\text{bytes}&532908\\
\text{SHA--256}&
\texttt{55C81E0B06A350C12BCA9355EC10F426ED479121584EA41BAE5A133D7C56C693}.
\end{array}                                                     \tag{56.18}
\]

% =============================================================================
% END APPENDED ELEVENTH-SERIES V56 MATERIAL
% =============================================================================

% =============================================================================
% BEGIN APPENDED ELEVENTH-SERIES V57 MATERIAL
% =============================================================================

\section{Local inf-sup assembly and volume-independent singular-value
floors}
\label{sec:v11v57}

The global analytic exponent in V56 can deteriorate as the number of
degrees of freedom grows.  Locality offers a different route.  If a
regulated carrier matrix is a uniformly invertible local part plus a
small finite-range coupling, its least singular value has a
volume-independent floor.  More generally, local solvers can be
assembled into an approximate right inverse whose defect is corrected
by a Neumann series.

\subsection{Block diagonal core}

Let
\[
                 E=\bigoplus_{b\in{\cal B}}E_b,\qquad
                 F=\bigoplus_{b\in{\cal B}}F_b                 \tag{57.1}
\]
be finite Hilbert direct sums.  Let
\(D=\bigoplus_bD_b:E\to F\), with each
\(D_b:E_b\to F_b\) surjective.

\begin{lemma}[Uniform local core]
\label{lem:v11v57-core}
If
\[
                           s_{\min}(D_b)\ge\kappa>0
                           \quad\text{for every }b,            \tag{57.2}
\]
then \(D\) is surjective and
\[
                              s_{\min}(D)\ge\kappa,            \tag{57.3}
\]
independently of \(|{\cal B}|\).
\end{lemma}

\begin{proof}
For \(u=(u_b)\),
\[
                         \lVert Du\rVert^2
                    =\sum_b\lVert D_bu_b\rVert^2.
\]
On the orthogonal complement of \(\ker D=\bigoplus_b\ker D_b\),
each \(u_b\) lies in \((\ker D_b)^\perp\), so (57.2) gives
\[
                         \lVert Du\rVert^2
                    \ge\kappa^2\sum_b\lVert u_b\rVert^2.
\]
This is (57.3), and blockwise surjectivity gives surjectivity of \(D\).
\end{proof}

\subsection{Finite-range coupling and the block Schur test}

Write a coupling \(C:E\to F\) in blocks \(C_{ab}:E_b\to F_a\), and set
\[
 r_C=\sup_a\sum_b\lVert C_{ab}\rVert,\qquad
 c_C=\sup_b\sum_a\lVert C_{ab}\rVert.                         \tag{57.4}
\]

\begin{lemma}[Block Schur bound]
\label{lem:v11v57-schur}
One has
\[
                              \lVert C\rVert\le\sqrt{r_Cc_C}. \tag{57.5}
\]
\end{lemma}

\begin{proof}
Put \(x_b=\lVert u_b\rVert\).  For every row,
\[
 \lVert(Cu)_a\rVert
 \le\sum_b\lVert C_{ab}\rVert x_b.
\]
Cauchy--Schwarz with weights \(\lVert C_{ab}\rVert\) gives
\[
 \lVert(Cu)_a\rVert^2
 \le\left(\sum_b\lVert C_{ab}\rVert\right)
     \left(\sum_b\lVert C_{ab}\rVert x_b^2\right).
\]
Sum in \(a\), use the row bound \(r_C\), then the column bound \(c_C\):
\[
 \lVert Cu\rVert^2
 \le r_C\sum_bx_b^2\sum_a\lVert C_{ab}\rVert
 \le r_Cc_C\lVert u\rVert^2.
\]
\end{proof}

\begin{theorem}[Local-core global gap]
\label{thm:v11v57-global-gap}
Under (57.2), let \(L=D+C\).  If
\[
                              \sqrt{r_Cc_C}<\kappa,            \tag{57.6}
\]
then \(L\) is surjective and
\[
                     s_{\min}(L)
                       \ge\kappa-\sqrt{r_Cc_C}>0.             \tag{57.7}
\]
The bound is independent of the number of blocks.
\end{theorem}

\begin{proof}
The least-singular-value perturbation inequality gives
\[
                       s_{\min}(D+C)
                       \ge s_{\min}(D)-\lVert C\rVert.
\]
Use Lemmas~\ref{lem:v11v57-core} and
\ref{lem:v11v57-schur}.  A positive least row singular value implies
surjectivity.
\end{proof}

\begin{corollary}[Bounded interaction degree]
\label{cor:v11v57-degree}
If every block row and column contains at most \(d\) nonzero coupling
blocks and \(\lVert C_{ab}\rVert\le\epsilon\), then
\[
                       s_{\min}(D+C)\ge\kappa-d\epsilon.       \tag{57.8}
\]
\end{corollary}

\begin{proof}
Here \(r_C,c_C\le d\epsilon\).  Apply
Theorem~\ref{thm:v11v57-global-gap}.
\end{proof}

\subsection{Assembly by an approximate right inverse}

A differential constraint operator need not have a block-diagonal
dominant part.  Patch solvers are more naturally expressed as an
approximate right inverse.

\begin{theorem}[Defect-corrected right inverse]
\label{thm:v11v57-right-inverse}
Let \(L:E\to F\) be linear between finite-dimensional Hilbert spaces.
Suppose \(R_0:F\to E\) obeys
\[
                       \lVert I_F-LR_0\rVert\le\eta<1.         \tag{57.9}
\]
Then \(L\) is surjective and has the exact right inverse
\[
                R=R_0(LR_0)^{-1}
                 =R_0\sum_{n=0}^{\infty}(I_F-LR_0)^n,         \tag{57.10}
\]
with
\[
                    \lVert R\rVert
                    \le\frac{\lVert R_0\rVert}{1-\eta},
 \qquad
                    s_{\min}(L)
                    \ge\frac{1-\eta}{\lVert R_0\rVert}.       \tag{57.11}
\]
\end{theorem}

\begin{proof}
The Neumann series in (57.10) converges because of (57.9) and is the
inverse of \(LR_0\).  Hence \(LR=I_F\), proving surjectivity and the
norm bound.  The Moore--Penrose right inverse \(L^\dagger\) has minimal
operator norm among right inverses, so
\[
 s_{\min}(L)^{-1}=\lVert L^\dagger\rVert\le\lVert R\rVert.
\]
This proves the last inequality.
\end{proof}

Let \(P_b:F\to F_b\) be local analysis maps, \(J_b:E_b\to E\) local
synthesis maps, and \(R_b:F_b\to E_b\) local solvers.  The natural
parametrix is
\[
                              R_0=\sum_bJ_bR_bP_b.             \tag{57.12}
\]

\begin{corollary}[Patchwise inf-sup certificate]
\label{cor:v11v57-patch}
If
\[
               \left\|\sum_bJ_bR_bP_b\right\|\le C_R\kappa^{-1},
 \qquad
               \left\|I-L\sum_bJ_bR_bP_b\right\|\le\eta<1,    \tag{57.13}
\]
then
\[
                              s_{\min}(L)
                       \ge\frac{(1-\eta)\kappa}{C_R}.          \tag{57.14}
\]
If \(C_R,\eta,\kappa\) are uniform, so is the global gap.
\end{corollary}

\begin{proof}
Insert (57.12)--(57.13) into
Theorem~\ref{thm:v11v57-right-inverse}.
\end{proof}

The constants \(C_R\) and \(\eta\) encode overlap multiplicity,
commutators with the partition of unity, boundary corrections, and
old--new mode mixing.  They cannot be suppressed as generic locality
constants.

\subsection{Application to the analytic carriers}

For each carrier in (56.12), either theorem can replace the
distance-to-singular-set estimate on a good chamber:

\begin{enumerate}
\item For \(A=D{\cal C}\) and the stacked carrier
\({\cal R}=(D{\cal C},D\chi)\), construct local right inverses for
constraint and gauge-condition data.
\item For \(M=D\chi R\), a local block core requires a gauge whose
principal symbol is uniformly invertible after removal of exact
residual modes.
\item For the KKT carrier \({\cal K}\), patch solvers must control both
the multiplier and field variables.
\item For the Dirac carrier \(\Delta\), the gap follows from \(M\), but
the determinant identity still carries the growing gauge dimension.
\end{enumerate}

\begin{proposition}[Good-chamber deterministic replacement]
\label{prop:v11v57-replacement}
If a carrier \(L_i(x)\) on the complete cutoff-\(i\) good chamber has
a parametrix satisfying
\[
             \sup_x\lVert R_{0,i}(x)\rVert\le R_*,
 \qquad
             \sup_x\lVert I-L_i(x)R_{0,i}(x)\rVert\le\eta_*<1,\tag{57.15}
\]
then
\[
                 \inf_xs_{\min}(L_i(x))
                    \ge(1-\eta_*)R_*^{-1}                    \tag{57.16}
\]
uniformly in \(i\).  No Lojasiewicz exponent or singular tube is needed
inside that chamber.
\end{proposition}

\begin{proof}
Apply Theorem~\ref{thm:v11v57-right-inverse} pointwise and take the
uniform bounds.
\end{proof}

This is stronger than a small-ball estimate.  Randomness remains only
in the probability of leaving the chamber.

\subsection{Low modes and derivative operators}

Locality alone does not imply (57.15).  For example, the periodic
forward difference on a one-dimensional lattice is local and has
Fourier singular values
\[
                              2|\sin(\pi k/L)|.                \tag{57.17}
\]
After removing its exact constant mode, the least singular value is
asymptotic to \(2\pi/L\), so it vanishes with the diameter.

\begin{warning}[Poincare scale]
\label{warn:v11v57-poincare}
Constraint and gauge carriers with derivative principal parts can have
physical long-wavelength modes.  A patchwise solver produces a uniform
global gap only if a mass, boundary anchoring, elliptic complement, or
scale-weighted norm removes the \(L^{-1}\) Poincare loss.  Declaring
finite interaction range is insufficient.
\end{warning}

A scale-weighted formulation may still be uniform.  If the target norm
on a mode of momentum \(p\) includes the natural factor \(|p|\), then
the derivative can have a dimensionless inf-sup constant.  The
cylinder ledger must state those norms before comparing determinants
or Hessians.

\subsection{Adjacent stability}

\begin{corollary}[Adjacent carrier stability]
\label{cor:v11v57-adjacent}
If \(s_{\min}(L_i)\ge\gamma_i\) on a common chamber and an identified
adjacent carrier satisfies
\[
                              \lVert L_{i+1}-L_i\rVert
                                   \le\delta_i<\gamma_i,       \tag{57.18}
\]
then
\[
                              s_{\min}(L_{i+1})
                                   \ge\gamma_i-\delta_i.       \tag{57.19}
\]
If \(\inf_i\gamma_i>0\) and \(\delta_i\to0\), the adjacent family has a
uniform eventual gap.
\end{corollary}

\begin{proof}
Use the least-singular-value perturbation inequality.
\end{proof}

When dimensions grow, the block-extension identification of V51 must
precede (57.18).

\begin{warning}[A gap does not normalize a growing determinant]
\label{warn:v11v57-determinant}
Even a uniform floor \(\gamma>0\) gives only
\(J(L_i)^{-1}\le\gamma^{-k_i}\).  If the row dimension \(k_i\) grows,
this stock is exponential in \(k_i\).  Local inf-sup control therefore
closes inverse operators but does not by itself close extensive
determinants.  A per-volume logarithmic normalization, local
determinant ratio, or cancellation with ghosts and coarea factors is
still required.
\end{warning}

\begin{ledger}[Locality acquisition]
\label{led:v11v57-locality}
For the chosen Einstein regulator, a volume-independent carrier gap
requires:
\begin{enumerate}
\item explicit patch spaces and scale-weighted norms;
\item local right inverses with a common bound;
\item bounded overlap and synthesis constants;
\item a defect estimate below one, including boundary commutators;
\item treatment of exact gauge and physical low modes;
\item adjacent compatibility of the patch parametrices.
\end{enumerate}
No current note has supplied these data for the full gravitational
constraint complex.
\end{ledger}

\begin{decision}[Next acquisition]
\label{dec:v11v57-next}
V58 will address Warning~\ref{warn:v11v57-determinant}.  It will replace
raw growing determinants by relative log determinants, prove a
trace-class adjacent bound, and state the summability condition needed
for a normalized determinant density.  The construction will retain
the anomaly and phase ledgers rather than treating a modulus
determinant as the full chiral weight.
\end{decision}

\begin{assessment}[Progress after V57]
\label{ass:v11v57-status}
Local gaps now assemble into a global volume-independent inf-sup bound
under explicit coupling or parametrix-defect conditions.  This offers
a route around global analytic exponents, while the periodic derivative
example records the long-wavelength obstruction.  The next independent
problem is the extensive determinant rather than inverse-operator
stability.
\end{assessment}

\subsection{Parent record}

This V57 note was formed by appending the preceding material to the
validated V56 source, whose record was
\[
\begin{array}{c|c}
\text{lines}&17693\\
\text{bytes}&546663\\
\text{SHA--256}&
\texttt{982791BA1BCD921D115C3C847B78E34354BA7637D867CDBE91BC7DDB3C3623E0}.
\end{array}                                                     \tag{57.20}
\]

% =============================================================================
% END APPENDED ELEVENTH-SERIES V57 MATERIAL
% =============================================================================
% =============================================================================
% BEGIN APPENDED ELEVENTH-SERIES V58 MATERIAL
% =============================================================================

\section{Relative determinant convergence and trace-class adjacent
increments}
\label{sec:v11v58}

A uniform inverse-operator gap does not control a raw determinant when
the matrix dimension grows.  The correct adjacent object is a relative
determinant after the reference contribution of newly introduced modes
has been divided out.  Trace-class summability then produces a finite
nonzero limit.  If only Hilbert--Schmidt summability is available, the
linear trace must be extracted as a counterterm.

\subsection{One relative determinant increment}

\begin{lemma}[Trace-class logarithm bound]
\label{lem:v11v58-log}
Let \(T\) be a finite-dimensional complex operator with
\(\lVert T\rVert\le\rho<1\).  The principal logarithm is
\[
                       \log(I+T)
                 =\sum_{n=1}^{\infty}\frac{(-1)^{n+1}}nT^n,   \tag{58.1}
\]
and
\[
 \left|\log\det(I+T)\right|
 =\left|\operatorname{Tr}\log(I+T)\right|
 \le\frac{\lVert T\rVert_1}{1-\rho}.                          \tag{58.2}
\]
Here \(\lVert\cdot\rVert_1\) is the trace norm.
\end{lemma}

\begin{proof}
The operator series converges absolutely because \(\lVert T\rVert<1\).
Moreover,
\[
 \lVert T^n\rVert_1
 \le\lVert T\rVert^{n-1}\lVert T\rVert_1.
\]
Thus
\[
 |\operatorname{Tr}\log(I+T)|
 \le\sum_{n\ge1}\frac{\rho^{n-1}}n\lVert T\rVert_1
 \le\frac{\lVert T\rVert_1}{1-\rho}.
\]
The disk \(\{z:|z-1|\le\rho\}\) avoids the nonpositive real axis and
zero, so the principal logarithm and determinant identity are
well-defined.
\end{proof}

\begin{corollary}[Positive Gram increment]
\label{cor:v11v58-gram}
Let \(G,\widetilde G>0\) act on the same finite-dimensional Hilbert
space and put
\[
                 T=G^{-1/2}(\widetilde G-G)G^{-1/2}.          \tag{58.3}
\]
If \(\lVert T\rVert\le\rho<1\), then
\[
 \log\frac{\det\widetilde G}{\det G}
                   =\operatorname{Tr}\log(I+T),               \tag{58.4}
\]
and its absolute value is bounded by (58.2).  For a coarea Jacobian
\(J=\sqrt{\det G}\), the relative log increment is one half of
(58.4).
\end{corollary}

\begin{proof}
The factorization
\[
                 \widetilde G=G^{1/2}(I+T)G^{1/2}
\]
gives (58.4) by multiplicativity of the determinant.  Take one half
for \(J\).
\end{proof}

\subsection{Dimension growth and reference new modes}

Let \(G_i>0\) act on \(F_i\).  Suppose
\[
                         F_{i+1}\simeq F_i\oplus F_i^{\rm new} \tag{58.5}
\]
and choose a positive reference block \(R_i>0\) on
\(F_i^{\rm new}\).  Put
\[
                         \widehat G_i=G_i\oplus R_i.           \tag{58.6}
\]
Assume the actual next Gram matrix has the relative representation
\[
            G_{i+1}=\widehat G_i^{1/2}(I+T_i)
                                  \widehat G_i^{1/2}.          \tag{58.7}
\]

\begin{theorem}[Relative determinant product]
\label{thm:v11v58-product}
Suppose
\[
             \sup_i\lVert T_i\rVert\le\rho<1,\qquad
             \sum_i\lVert T_i\rVert_1<\infty.                 \tag{58.8}
\]
Define
\[
 {\cal Z}_n
 =\frac{\det G_n}
        {\det G_0\prod_{i=0}^{n-1}\det R_i}.                  \tag{58.9}
\]
Then
\[
             \log{\cal Z}_n
               =\sum_{i=0}^{n-1}\operatorname{Tr}\log(I+T_i) \tag{58.10}
\]
converges absolutely.  Consequently
\[
                         {\cal Z}_n\longrightarrow{\cal Z}_\infty
                         \quad\text{with}\quad
                         0<{\cal Z}_\infty<\infty.            \tag{58.11}
\]
For Jacobians \(J_i=\sqrt{\det G_i}\), the same theorem holds with
all logarithms divided by two.
\end{theorem}

\begin{proof}
Taking determinants in (58.7) gives
\[
 \det G_{i+1}
 =\det G_i\,\det R_i\,\det(I+T_i).
\]
Multiplication telescopes to (58.10).  Lemma~\ref{lem:v11v58-log} and
(58.8) make the series absolutely convergent.  Its exponential is
finite and nonzero, proving (58.11).
\end{proof}

The reference blocks \(R_i\) are part of the renormalization
prescription.  Changing them changes \({\cal Z}_\infty\) unless the
change has its own convergent relative product or is absorbed into a
declared local counterterm.

\subsection{Only Hilbert--Schmidt control}

For a finite matrix \(T\), define
\[
                  \det{}_2(I+T):=\det(I+T)e^{-\operatorname{Tr}T}.
                                                                  \tag{58.12}
\]

\begin{theorem}[Quadratically regularized increment]
\label{thm:v11v58-det2}
If \(\lVert T\rVert\le\rho<1\), then
\[
 \log\det{}_2(I+T)
 =\sum_{n=2}^{\infty}\frac{(-1)^{n+1}}n\operatorname{Tr}(T^n)
                                                                  \tag{58.13}
\]
and
\[
 \left|\log\det{}_2(I+T)\right|
 \le\frac{\lVert T\rVert_2^2}{2(1-\rho)},                     \tag{58.14}
\]
where \(\lVert\cdot\rVert_2\) is the Hilbert--Schmidt norm.
Consequently, if
\[
              \sup_i\lVert T_i\rVert\le\rho<1,\qquad
              \sum_i\lVert T_i\rVert_2^2<\infty,              \tag{58.15}
\]
then \(\prod_i\det{}_2(I+T_i)\) converges to a finite nonzero limit.
\end{theorem}

\begin{proof}
Subtracting \(\operatorname{Tr}T\) removes the \(n=1\) term from
(58.1).  For \(n\ge2\),
\[
 |\operatorname{Tr}T^n|
 \le\lVert T^2\rVert_1\lVert T\rVert^{n-2}
 \le\lVert T\rVert_2^2\rho^{n-2}.
\]
Since \(1/n\le1/2\),
\[
 \sum_{n\ge2}\frac{\rho^{n-2}}n
 \le\frac1{2(1-\rho)}.
\]
This proves (58.14); summing it under (58.15) proves convergence.
\end{proof}

\begin{warning}[The linear trace is a counterterm]
\label{warn:v11v58-linear}
The ordinary determinant decomposes as
\[
                 \log\det(I+T)
                  =\operatorname{Tr}T+\log\det{}_2(I+T).      \tag{58.16}
\]
Hilbert--Schmidt summability controls only the second term.  The series
\(\sum_i\operatorname{Tr}T_i\) must converge, cancel between fields,
or be subtracted by an explicitly justified local counterterm.
Discarding it silently changes the measure.
\end{warning}

\subsection{Combined boson, ghost, and fermion factors}

A regulated gauge--matter density contains several determinant factors.
On a common branch, write its relative logarithm schematically as
\[
                 \Delta_i^{\det}
       =\sum_{\alpha}\sigma_\alpha
          \operatorname{Tr}\log(I+T_{\alpha,i}),              \tag{58.17}
\]
where \(\sigma_\alpha\) includes integer or half-integer powers and
fermionic signs.

\begin{proposition}[Combined trace cancellation]
\label{prop:v11v58-cancellation}
Suppose \(\sup_{\alpha,i}\lVert T_{\alpha,i}\rVert\le\rho<1\) and
\[
 \sum_i\left|
      \sum_\alpha\sigma_\alpha\operatorname{Tr}T_{\alpha,i}
      \right|<\infty,\qquad
 \sum_{i,\alpha}|\sigma_\alpha|
                    \lVert T_{\alpha,i}\rVert_2^2<\infty.     \tag{58.18}
\]
Then
\[
                         \sum_i\Delta_i^{\det}                \tag{58.19}
\]
converges absolutely after grouping all fields at each common cutoff.
\end{proposition}

\begin{proof}
Use (58.16) for each \(\alpha\).  The grouped linear terms are summable
by the first condition in (58.18).  The regularized remainders are
summable by (58.14) and the second condition.
\end{proof}

The grouping at a common cutoff is essential.  Separately divergent
bosonic and fermionic trace series cannot be reordered before their
local cancellation has been proved.

\subsection{Complex chiral determinants}

For a complex chiral operator, (58.1) selects a local logarithm branch
whenever \(\lVert T_i\rVert<1\).  The resulting adjacent phase
increment is
\[
                    \operatorname{ImTr}\log(I+T_i).           \tag{58.20}
\]
Trace-class summability controls this local phase along one chain of
cutoffs.

\begin{firewall}[Determinant line and anomaly]
\label{fire:v11v58-chiral}
A convergent product of local principal-branch increments does not
trivialize the determinant line around closed loops in field space.
The exact Standard Model perturbative anomaly cancellation established
earlier removes the local anomaly polynomial, but global anomaly,
overlap transition, and phase-denominator rows remain.  Replacing the
chiral determinant by its modulus would erase precisely this
information and is not used here.
\end{firewall}

\subsection{Relation to extensive free energy}

A per-volume limit
\[
                   |F_i|^{-1}\log\det G_i\longrightarrow p   \tag{58.21}
\]
controls the extensive free-energy density.  It does not imply
Theorem~\ref{thm:v11v58-product}: an error \(o(|F_i|)\) can still
diverge and alter every finite cylinder weight.

\begin{proposition}[Relative summability is stronger]
\label{prop:v11v58-stronger}
There are positive scalar determinants with a convergent per-volume
logarithm but a divergent relative product after subtraction of that
volume term.
\end{proposition}

\begin{proof}
Let \(\det G_n=\exp(np+\sqrt n)\) and take volume \(n\).
Then \(n^{-1}\log\det G_n\to p\), but after removing \(\exp(np)\) the
relative determinant is \(\exp(\sqrt n)\), which diverges.
\end{proof}

Thus pressure control and cylinder-weight control are separate ledger
entries.

\subsection{The V42 determinant row}

\begin{ledger}[Relative determinant acquisition]
\label{led:v11v58-acquisition}
For every coarea, Faddeev--Popov, stationary, ghost, and matter
determinant, the adjacent construction must specify:
\begin{enumerate}
\item the old-mode identification and positive or invertible new-mode
reference block;
\item the relative perturbation \(T_i\) in a common fibre;
\item a uniform operator-norm radius below one;
\item trace-norm summability, or Hilbert--Schmidt square summability
plus the linear trace counterterm;
\item common-cutoff cancellation identities between field species;
\item branch and overlap compatibility for complex determinants;
\item the effect of changing the reference blocks.
\end{enumerate}
Under these conditions, the determinant contribution becomes a
summable adjacent scalar rather than an exponential dimension stock.
\end{ledger}

\begin{warning}[Locality is not trace class]
\label{warn:v11v58-locality}
If \(T_i\) has \(N_i\) local blocks of size \(O(\epsilon_i)\), its
operator norm may be \(O(\epsilon_i)\) while its trace norm is
\(O(N_i\epsilon_i)\).  Finite interaction range and the V57 gap
therefore do not imply (58.8).  The refinement rate must beat the
number of newly affected local modes, or local counterterms must remove
the extensive part.
\end{warning}

\begin{decision}[Next acquisition]
\label{dec:v11v58-next}
V59 will extract local trace counterterms from finite-range adjacent
operators.  It will prove a bulk--boundary decomposition under a
finite propagation hypothesis and state the rate at which the
renormalized remainder becomes trace summable.  The result will be a
deterministic criterion, not a claim that the required Einstein
counterterms have already been identified.
\end{decision}

\begin{assessment}[Progress after V58]
\label{ass:v11v58-status}
Growing determinants are now expressed as relative adjacent products.
Trace-class increments give an ordinary nonzero limit; Hilbert--Schmidt
increments give a regularized limit only after the linear trace is
retained as a counterterm.  This removes the crude exponential
inverse-Jacobian stock when the stronger relative hypotheses can be
verified.
\end{assessment}

\subsection{Parent record}

This V58 note was formed by appending the preceding material to the
validated V57 source, whose record was
\[
\begin{array}{c|c}
\text{lines}&18055\\
\text{bytes}&558635\\
\text{SHA--256}&
\texttt{FC5E98CCA619889976C664834EE39C1164DF55F66CF30F7B0A79EB3CFF62D663}.
\end{array}                                                     \tag{58.22}
\]

% =============================================================================
% END APPENDED ELEVENTH-SERIES V58 MATERIAL
% =============================================================================
% =============================================================================
% BEGIN APPENDED ELEVENTH-SERIES V59 MATERIAL
% =============================================================================

\section{Finite-propagation counterterms and bulk--boundary trace
decomposition}
\label{sec:v11v59}

V58 isolates trace powers that must be subtracted when an adjacent
relative operator is not trace class.  Locality makes those powers
local, but only after boundary and defect layers are counted.  This
section gives the exact decomposition and its summability criterion.

\subsection{Higher regularized determinants}

For an integer \(p\ge2\), define in finite dimension
\[
 \det{}_p(I+T)
 :=\det(I+T)
   \exp\left\{\sum_{n=1}^{p-1}
                    \frac{(-1)^n}{n}\operatorname{Tr}(T^n)\right\}.
                                                                  \tag{59.1}
\]

\begin{theorem}[Schatten-\(p\) remainder]
\label{thm:v11v59-detp}
If \(\lVert T\rVert\le\rho<1\), then
\[
 \log\det{}_p(I+T)
 =\sum_{n=p}^{\infty}
       \frac{(-1)^{n+1}}n\operatorname{Tr}(T^n),              \tag{59.2}
\]
and
\[
 \left|\log\det{}_p(I+T)\right|
 \le\frac{\lVert T\rVert_p^p}{p(1-\rho)}.                    \tag{59.3}
\]
Consequently a sequence with a common \(\rho<1\) and
\[
                              \sum_i\lVert T_i\rVert_p^p<\infty\tag{59.4}
\]
has a convergent product of \(\det{}_p(I+T_i)\).
\end{theorem}

\begin{proof}
The exponential in (59.1) cancels the first \(p-1\) terms of the
logarithm series.  For \(n\ge p\),
\[
 |\operatorname{Tr}T^n|
 \le\lVert T^n\rVert_1
 \le\lVert T\rVert_p^p\lVert T\rVert^{n-p}
 \le\lVert T\rVert_p^p\rho^{n-p}.
\]
Since \(1/n\le1/p\), summation proves (59.3).  Equation (59.4) then
gives absolute convergence of the logarithmic product.
\end{proof}

The removed trace powers \(1\le n<p\) are not discarded.  They are
the counterterm ledger.

\subsection{Finite propagation}

Let \(\Lambda\) be a finite subset of a graph with graph distance
\(d\), and let
\[
                         {\cal H}_\Lambda
                    =\bigoplus_{x\in\Lambda}\mathbb C^d_0     \tag{59.5}
\]
with fixed internal dimension \(d_0\).  Write \(T_{xy}\) for the block
kernel.

\begin{definition}[Propagation radius]
\label{def:v11v59-propagation}
The operator \(T\) has propagation at most \(R\) if
\[
                              T_{xy}=0\quad\text{when }d(x,y)>R.
                                                                  \tag{59.6}
\]
\end{definition}

\begin{lemma}[Locality of trace powers]
\label{lem:v11v59-power-local}
If \(T\) has propagation \(R\), then
\[
 \operatorname{tr}_{\mathbb C^{d_0}}(T^n)_{xx}               \tag{59.7}
\]
depends only on kernel blocks \(T_{yz}\) with
\(y,z\in B_{nR}(x)\).  Moreover,
\[
       \left|\operatorname{tr}_{\mathbb C^{d_0}}(T^n)_{xx}\right|
                         \le d_0\lVert T\rVert^n.              \tag{59.8}
\]
\end{lemma}

\begin{proof}
The diagonal block is
\[
 (T^n)_{xx}
 =\sum_{x_1,\ldots,x_{n-1}}
       T_{xx_1}T_{x_1x_2}\cdots T_{x_{n-1}x}.                 \tag{59.9}
\]
Every nonzero path has steps of length at most \(R\), so all vertices
lie in \(B_{nR}(x)\).  Compression to the \(x\)-block has norm at most
\(\lVert T^n\rVert\le\lVert T\rVert^n\); its trace has absolute value
at most \(d_0\) times that norm.
\end{proof}

\subsection{Exact bulk--boundary decomposition}

Let \({\cal E}\subset\Lambda\) contain the outer boundary, interface,
and local observable defect sites.  Define its \(r\)-neighborhood
inside \(\Lambda\) by
\[
                         {\cal E}^{(r)}
       =\{x\in\Lambda:d(x,{\cal E})\le r\}.                   \tag{59.10}
\]
Assume that outside \({\cal E}^{(nR)}\), the radius-\(nR\) kernel
environment is translation-equivalent to a fixed bulk environment.
Then the diagonal trace (59.7) has a common value \(\tau_n\) there.

\begin{theorem}[Bulk trace and boundary remainder]
\label{thm:v11v59-bulk}
If \(\lVert T\rVert\le\epsilon\), then for every \(n\ge1\),
\[
                 \operatorname{Tr}(T^n)
                       =|\Lambda|\tau_n+{\cal B}_n,           \tag{59.11}
\]
where
\[
                         |{\cal B}_n|
                   \le2d_0|{\cal E}^{(nR)}|\epsilon^n.        \tag{59.12}
\]
\end{theorem}

\begin{proof}
Sum the diagonal block traces over \(x\).  For
\(x\notin{\cal E}^{(nR)}\), the summand equals \(\tau_n\).
Thus
\[
 {\cal B}_n
 =\sum_{x\in{\cal E}^{(nR)}}
   \left[\operatorname{tr}(T^n)_{xx}-\tau_n\right].           \tag{59.13}
\]
By (59.8), the first term has magnitude at most
\(d_0\epsilon^n\).  The same bound holds for \(\tau_n\), since it is a
diagonal trace in the bulk environment.  Summation proves (59.12).
\end{proof}

The counterterm \(|\Lambda|\tau_n\) is an extensive local bulk term.
The remainder is supported in a layer whose thickness grows only as
\(nR\).

\begin{corollary}[Summable counterterm remainder]
\label{cor:v11v59-summable}
For adjacent operators \(T_i\) with
\(\lVert T_i\rVert\le\epsilon_i\le\rho<1\), propagation \(R_i\),
internal dimension \(d_{0,i}\), and exceptional set \({\cal E}_i\),
the boundary parts of all counterterms below order \(p\) are summable
if
\[
 \sum_i d_{0,i}\sum_{n=1}^{p-1}
       \frac{|{\cal E}_i^{(nR_i)}|\epsilon_i^n}{n}<\infty.    \tag{59.14}
\]
If also
\[
                              \sum_i\lVert T_i\rVert_p^p<\infty,\tag{59.15}
\]
then the determinant product converges after subtracting the bulk
counterterms
\[
                 \sum_{n=1}^{p-1}
                   \frac{(-1)^{n+1}}n|\Lambda_i|\tau_{n,i}.  \tag{59.16}
\]
\end{corollary}

\begin{proof}
Theorem~\ref{thm:v11v59-bulk} and (59.14) control the difference
between the exact trace counterterms and (59.16).  Theorem
\ref{thm:v11v59-detp} and (59.15) control the regularized remainder.
Adding the two absolutely convergent series proves the claim.
\end{proof}

\subsection{A usable volume-rate condition}

If \(\operatorname{rank}T_i\le d_{0,i}|\Lambda_i|\), then
\[
                 \lVert T_i\rVert_p^p
             \le d_{0,i}|\Lambda_i|\epsilon_i^p.              \tag{59.17}
\]

\begin{corollary}[Power-counting schedule]
\label{cor:v11v59-rate}
Suppose
\[
 |\Lambda_i|\lesssim(i+2)^v,\quad
 |{\cal E}_i^{((p-1)R_i)}|\lesssim(i+2)^e,\quad
 d_{0,i}\lesssim(i+2)^d,\quad
 \epsilon_i\lesssim(i+2)^{-s}.                               \tag{59.18}
\]
Then (59.15) follows if
\[
                              ps>v+d+1,                       \tag{59.19}
\]
and the boundary-counterterm condition (59.14) follows from
\[
                              s>e+d+1.                        \tag{59.20}
\]
\end{corollary}

\begin{proof}
Equation (59.17) bounds the Schatten summand by
\(C(i+2)^{v+d-ps}\), which is summable under (59.19).
Since \(\epsilon_i\le1\), the \(n=1\) boundary term dominates the
finite sum in (59.14) up to a constant; it is bounded by
\(C(i+2)^{e+d-s}\), summable under (59.20).
\end{proof}

The boundary rate is often the stricter condition.  Increasing the
regularization order \(p\) helps the bulk Schatten remainder but does
not improve the first trace boundary term.

\subsection{Background-dependent local densities}

Translation invariance is not required for locality.  Suppose a local
counterterm density \(c_{n,i}(x)\), depending only on the background in
\(B_{nR_i}(x)\), satisfies
\[
 \operatorname{tr}(T_i^n)_{xx}=c_{n,i}(x)
       \quad\text{outside }{\cal E}_i^{(nR_i)}.               \tag{59.21}
\]
Then Theorem~\ref{thm:v11v59-bulk} holds with
\[
                         |\Lambda_i|\tau_{n,i}
                    \quad\text{replaced by}\quad
                         \sum_{x\in\Lambda_i}c_{n,i}(x),       \tag{59.22}
\]
and the same boundary bound if
\(|c_{n,i}(x)|\le d_{0,i}\epsilon_i^n\).

These local densities must be expanded in the allowed regulator
counterterm basis.  For the SM--Einstein programme this includes,
subject to symmetries and dimension, vacuum, curvature, gauge,
Higgs, Yukawa, and higher-order irrelevant terms.  The proof above
establishes locality of the trace subtraction; it does not prove that
a finite renormalizable basis suffices for the proposed renewal
regulator.

\subsection{Combined-field cancellation}

Let \(\alpha\) label bosons, ghosts, and fermions.  At trace order
\(n<p\), define the grouped local density
\[
                         c_{n,i}^{\rm tot}(x)
                   =\sum_\alpha\sigma_\alpha
                         c_{\alpha,n,i}(x).                   \tag{59.23}
\]

\begin{proposition}[Local cancellation before summation]
\label{prop:v11v59-local-cancel}
If \(c_{n,i}^{\rm tot}(x)=0\) pointwise in every common bulk chart,
then the extensive counterterm at order \(n\) cancels exactly between
species.  Only the summed boundary and defect remainders remain.
\end{proposition}

\begin{proof}
Sum (59.23) over \(x\).  Pointwise cancellation makes the bulk sum
zero.  The exceptional-layer remainders are not translation-equivalent
to the bulk and remain bounded by the species-summed version of
(59.12).
\end{proof}

An integrated cancellation without a common regulator and common local
chart is weaker and does not justify rearranging separately divergent
species traces.

\subsection{Cylinder compatibility}

For a fixed cylinder observable supported in \(K\), include \(K\) in
\({\cal E}_i\).  Finite propagation then confines its determinant
response through trace order \(p-1\) to
\(K^{((p-1)R_i)}\).  Outer-boundary effects are separate and vanish
only under (59.14).

\begin{ledger}[Determinant counterterm row]
\label{led:v11v59-row}
A complete adjacent determinant proof must now provide:
\begin{enumerate}
\item a relative operator \(T_i\) with finite propagation or a
quantified off-diagonal decay substitute;
\item its operator and Schatten-\(p\) rates;
\item the local densities \(c_{n,i}(x)\), \(1\le n<p\);
\item their expansion in symmetry-allowed counterterms;
\item common-cutoff cancellations between species;
\item boundary and observable-defect layer counts;
\item summability conditions (59.14)--(59.15);
\item complex phase compatibility from V58.
\end{enumerate}
\end{ledger}

\begin{firewall}[No counterterms identified]
\label{fire:v11v59-boundary}
V59 proves that trace-power subtractions are local under finite
propagation and gives exact bulk--boundary bounds.  It does not compute
the \(c_{n,i}\) for the renewal SM--Einstein regulator, show that all
divergences fit a finite physical counterterm basis, or prove their
renormalization conditions.  Nonlocal chiral phases and massless
long-range resolvents may require decay estimates beyond finite
propagation.
\end{firewall}

\begin{decision}[Mandatory V60 audit]
\label{dec:v11v59-next}
The next version is the compulsory YM/NS audit.  It will check whether
either companion has produced a nonzero-momentum, decay, or boundary
estimate relevant to (59.14)--(59.23).  V61 will then continue with
the strongest available route, using off-diagonal decay if finite
propagation is unavailable.
\end{decision}

\begin{assessment}[Progress after V59]
\label{ass:v11v59-status}
Schatten regularization is now tied to explicit local counterterms.
Finite propagation turns every removed trace power into a bulk local
density plus a quantified boundary layer, and the refinement rates
needed for both pieces are stated.  The remaining work is to derive
these data from the actual SM--Einstein adjacent operators.
\end{assessment}

\subsection{Parent record}

This V59 note was formed by appending the preceding material to the
validated V58 source, whose record was
\[
\begin{array}{c|c}
\text{lines}&18407\\
\text{bytes}&570558\\
\text{SHA--256}&
\texttt{76CDEB9013975F218923928BC865BA4751EC43B61069F8FEA8EC4F08007EA708}.
\end{array}                                                     \tag{59.24}
\]

% =============================================================================
% END APPENDED ELEVENTH-SERIES V59 MATERIAL
% =============================================================================
% =============================================================================
% BEGIN APPENDED ELEVENTH-SERIES V60 MATERIAL
% =============================================================================

\section{Companion audit at V60: locality without new decay data}
\label{sec:v11v60}

The mandatory audit again finds unchanged companion files:
\[
\begin{array}{c|r|l}
\text{file}&\text{bytes}&\text{SHA--256}\\ \hline
\texttt{Renewal\_Yang\_Mills\_Mass\_Gap\_Research\_Notes\_V35.tex}
 &498660&
\texttt{23F674C142CA9B298546508F7A3E9EAF2CE2E154078C32B0BA2FC8660AD25098}\\
\texttt{Renewal\_NS\_Stability\_Research\_Notes\_V26.tex}
 &278283&
\texttt{82C887F8C2C69E4F28FAD7C3DC8DE5B9099CB5A4BF431EBA1FFF3E91935978AD}.
\end{array}                                                     \tag{60.1}
\]
These are identical to the V55 records.  No companion theorem is
repeated.

\subsection{Audit consequence}

The YM V35 action and its local translated vertices have finite
lattice range before an inverse covariance or a determinant logarithm
is formed.  Thus V59 applies directly to trace polynomials of bare
finite-range variations.  The nonzero-momentum response requested by
that programme, however, contains shifted covariance inverses and is
not itself finite propagation.

The NS V26 Oseen-kernel calculations likewise concern kernels with
tails rather than compact support.  Their rank-one constants do not
supply a determinant counterterm or a gravitational resolvent decay
rate.

\begin{proposition}[Where finite propagation is lost]
\label{prop:v11v60-inverse}
Let \(A\) be an invertible finite-range matrix and let \(\dot A\) be a
finite-range perturbation.  Then
\[
                    \frac d{dt}(A+t\dot A)^{-1}\big|_{t=0}
                              =-A^{-1}\dot A A^{-1}.           \tag{60.2}
\]
Even when \(A\) and \(\dot A\) have finite propagation, the right-hand
side generally does not.
\end{proposition}

\begin{proof}
Differentiate
\((A+t\dot A)(A+t\dot A)^{-1}=I\).  A finite-range inverse would
require additional algebraic structure; a generic inverse of a
banded connected matrix has nonzero entries at arbitrary separation.
\end{proof}

Therefore the next step must quantify decay rather than assume exact
support.

\begin{decision}[V61 acquisition]
\label{dec:v11v60-next}
V61 will prove truncation theorems for exponentially and polynomially
decaying block kernels on bounded-growth graphs.  It will convert
off-diagonal tails into finite-propagation approximants, quantify the
operator, trace, and trace-power errors, and state schedules that feed
V59.  No decay rate will be attributed to the Einstein or chiral
resolvents until separately proved.
\end{decision}

\begin{firewall}[Audit boundary]
\label{fire:v11v60-boundary}
The unchanged YM note supplies a positive local chart and finite-range
bare action, but no nonzero-momentum inverse-kernel decay theorem.  The
unchanged NS note supplies pointwise Oseen derivative norms, not a
four-dimensional SM--Einstein determinant estimate.  Neither changes
the unresolved chiral phase, gravitational low-mode, or adjacent
quotient rows.
\end{firewall}

\begin{assessment}[Progress after V60]
\label{ass:v11v60-status}
The required companion inspection is complete.  It confirms that V59
covers bare local trace powers and that resolvent-generated terms
require a new quasi-local truncation estimate, which is taken up next.
\end{assessment}

\subsection{Parent record}

This V60 note was formed by appending the preceding material to the
validated V59 source, whose record was
\[
\begin{array}{c|c}
\text{lines}&18751\\
\text{bytes}&582526\\
\text{SHA--256}&
\texttt{1A55630CEB48BB1592A3B8F7D5E53B03F5BBC44DB306AC4090E42BA53B95C808}.
\end{array}                                                     \tag{60.3}
\]

% =============================================================================
% END APPENDED ELEVENTH-SERIES V60 MATERIAL
% =============================================================================
% =============================================================================
% BEGIN APPENDED ELEVENTH-SERIES V61 MATERIAL
% =============================================================================

\section{Quasi-local kernels, finite-range truncation, and resolvent
decay}
\label{sec:v11v61}

V59 applies directly to finite-propagation relative operators.  The
inverse identity in V60 shows why resolvents need not have finite
range.  The present section quantifies their replacement by truncated
kernels and distinguishes exponential decay, which preserves a
logarithmic boundary layer, from polynomial decay, which can require a
much larger layer.

\subsection{An abstract tail-Schur theorem}

Let \(\Lambda\) be a finite graph and
\({\cal H}_\Lambda=\bigoplus_{x\in\Lambda}\mathbb C^{d_0}\).
For a block kernel \(T=(T_{xy})\), let
\[
                 T^{[R]}_{xy}=\mathbf1_{\{d(x,y)\le R\}}T_{xy}.
                                                                  \tag{61.1}
\]
Define the row and column tails
\[
\begin{split}
 a_r(R)&=\sup_x\sum_{y:d(x,y)>R}\lVert T_{xy}\rVert,\\
 a_c(R)&=\sup_y\sum_{x:d(x,y)>R}\lVert T_{xy}\rVert.           \tag{61.2}
\end{split}
\]

\begin{theorem}[Finite-range truncation]
\label{thm:v11v61-truncation}
One has
\[
 \lVert T-T^{[R]}\rVert
                     \le\sqrt{a_r(R)a_c(R)}.                  \tag{61.3}
\]
Moreover,
\[
 \lVert T-T^{[R]}\rVert_1
 \le d_0|\Lambda|\sqrt{a_r(R)a_c(R)}.                         \tag{61.4}
\]
If \(\lVert T\rVert,\lVert T^{[R]}\rVert\le M\), then
\[
 \left|\operatorname{Tr}(T^n)
       -\operatorname{Tr}\bigl((T^{[R]})^n\bigr)\right|
 \le nM^{n-1}d_0|\Lambda|
                    \sqrt{a_r(R)a_c(R)}.                      \tag{61.5}
\]
\end{theorem}

\begin{proof}
Apply the block Schur estimate
Lemma~\ref{lem:v11v57-schur} to the tail kernel to obtain (61.3).
Its rank is at most \(d_0|\Lambda|\), so
\(\lVert A\rVert_1\le\operatorname{rank}(A)\lVert A\rVert\)
gives (61.4).  Finally,
\[
 T^n-S^n=\sum_{j=0}^{n-1}T^j(T-S)S^{n-1-j},\qquad S=T^{[R]}. \tag{61.6}
\]
Use \(|\operatorname{Tr}A|\le\lVert A\rVert_1\), the ideal property of
the trace norm, and (61.4).
\end{proof}

The volume factor in (61.4) is real.  Operator-norm quasi-locality does
not by itself give trace-class adjacent convergence.

\subsection{Decay on polynomial-growth graphs}

Assume uniformly in \(x\) and \(\Lambda\) that
\[
                 |\{y:d(x,y)=r\}|\le C_g(1+r)^{D-1}.          \tag{61.7}
\]

\begin{theorem}[Exponential tail]
\label{thm:v11v61-exponential}
If
\[
                         \lVert T_{xy}\rVert
                    \le\epsilon e^{-\mu d(x,y)},              \tag{61.8}
\]
then there is \(C_{\mu,D,g}\) independent of \(\Lambda\) such that
\[
 a_r(R),a_c(R)
 \le C_{\mu,D,g}\epsilon e^{-\mu R/2}.                        \tag{61.9}
\]
Consequently
\[
 \lVert T-T^{[R]}\rVert
 \le C_{\mu,D,g}\epsilon e^{-\mu R/2}.                        \tag{61.10}
\]
\end{theorem}

\begin{proof}
Using (61.7), each row tail is at most
\[
 C_g\epsilon\sum_{r>R}(1+r)^{D-1}e^{-\mu r}.
\]
The polynomial factor is bounded by a constant times \(e^{\mu r/2}\),
and the remaining geometric series proves (61.9).  The same argument
applies to columns.  Insert the result in (61.3).
\end{proof}

\begin{theorem}[Polynomial tail]
\label{thm:v11v61-polynomial}
If for some \(\sigma>0\),
\[
                         \lVert T_{xy}\rVert
           \le\epsilon(1+d(x,y))^{-(D+\sigma)},               \tag{61.11}
\]
then
\[
 a_r(R),a_c(R)\le C_{\sigma,D,g}\epsilon(1+R)^{-\sigma},      \tag{61.12}
\]
and
\[
 \lVert T-T^{[R]}\rVert
          \le C_{\sigma,D,g}\epsilon(1+R)^{-\sigma}.          \tag{61.13}
\]
\end{theorem}

\begin{proof}
The shell sum is bounded by a constant times
\[
 \epsilon\sum_{r>R}(1+r)^{-1-\sigma}
 \le C\epsilon(1+R)^{-\sigma}.
\]
Apply the same estimate to rows and columns, then (61.3).
\end{proof}

\subsection{Truncation schedules for trace powers}

Let \(N_i=d_{0,i}|\Lambda_i|\), and let \(\eta_i>0\) be a desired
trace-norm error.

\begin{corollary}[Exponential truncation radius]
\label{cor:v11v61-exp-radius}
Under (61.8), it suffices to choose
\[
 R_i\ge\frac2{\mu_i}
   \log\left(\frac{C_iN_i\epsilon_i}{\eta_i}\right)_+         \tag{61.14}
\]
to obtain
\[
                    \lVert T_i-T_i^{[R_i]}\rVert_1\le\eta_i. \tag{61.15}
\]
Here \(\log(x)_+=\max\{0,\log x\}\), with an additive integer rounding.
\end{corollary}

\begin{proof}
Combine (61.4) and (61.9), then solve the resulting exponential
inequality for \(R_i\).
\end{proof}

\begin{corollary}[Polynomial truncation radius]
\label{cor:v11v61-poly-radius}
Under (61.11), it suffices to choose
\[
 R_i\ge
 \left(\frac{C_iN_i\epsilon_i}{\eta_i}\right)^{1/\sigma_i}-1 \tag{61.16}
\]
for (61.15).
\end{corollary}

\begin{proof}
Combine (61.4) and (61.12) and solve for \(R_i\).
\end{proof}

If \(\sum_i\eta_i<\infty\), equation (61.5) makes the difference
between every fixed finite list of exact and truncated trace powers
summable, provided their operator norms are uniformly bounded.
The truncated operators then enter V59 with propagation \(R_i\).

Exponential decay costs a logarithmic growth of \(R_i\) in the volume.
Polynomial decay can make \(R_i\) a positive power of the volume,
possibly as large as the system itself.  In that regime the
bulk--boundary separation of V59 gives no gain.

\subsection{Weighted Schur algebra}

For a block kernel \(K\), define
\[
 \lVert K\rVert_{{\cal S}_\mu}
 :=\max\left\{
 \sup_x\sum_y e^{\mu d(x,y)}\lVert K_{xy}\rVert,\,
 \sup_y\sum_x e^{\mu d(x,y)}\lVert K_{xy}\rVert
 \right\}.                                                    \tag{61.17}
\]

\begin{lemma}[Weighted algebra property]
\label{lem:v11v61-algebra}
For compatible kernels,
\[
                   \lVert KL\rVert_{{\cal S}_\mu}
             \le\lVert K\rVert_{{\cal S}_\mu}
                \lVert L\rVert_{{\cal S}_\mu}.                \tag{61.18}
\]
Moreover,
\[
                         \lVert K_{xy}\rVert
            \le\lVert K\rVert_{{\cal S}_\mu}e^{-\mu d(x,y)}. \tag{61.19}
\]
\end{lemma}

\begin{proof}
The triangle inequality for graph distance gives
\[
 e^{\mu d(x,z)}
 \le e^{\mu d(x,y)}e^{\mu d(y,z)}.
\]
Insert the kernel product, sum first in \(z\) and then in \(y\), and
use the two row bounds.  The column proof is identical.  A single term
in the row sum gives (61.19).
\end{proof}

\begin{theorem}[Exponential inverse from a gapped local core]
\label{thm:v11v61-inverse}
Let \(A=D+K\), where \(D\) is block diagonal,
\[
                         \sup_x\lVert D_x^{-1}\rVert\le\kappa^{-1},
                                                                  \tag{61.20}
\]
and
\[
                         q:=\kappa^{-1}
                             \lVert K\rVert_{{\cal S}_\mu}<1. \tag{61.21}
\]
Then \(A\) is invertible and
\[
 \lVert A^{-1}\rVert_{{\cal S}_\mu}
                    \le\frac{\kappa^{-1}}{1-q},               \tag{61.22}
\]
so
\[
                         \lVert(A^{-1})_{xy}\rVert
                    \le\frac{\kappa^{-1}}{1-q}
                                e^{-\mu d(x,y)}.               \tag{61.23}
\]
\end{theorem}

\begin{proof}
Factor
\[
                         A=D(I+D^{-1}K).
\]
The diagonal kernel \(D^{-1}\) has weighted Schur norm at most
\(\kappa^{-1}\), so
\(\lVert D^{-1}K\rVert_{{\cal S}_\mu}\le q<1\).
The Neumann series gives
\[
 A^{-1}=\sum_{n=0}^{\infty}(-D^{-1}K)^nD^{-1}.
\]
Use Lemma~\ref{lem:v11v61-algebra} to sum its weighted norm and obtain
(61.22); (61.19) gives (61.23).
\end{proof}

This theorem is deterministic and volume independent, but its
hypothesis is stronger than an unweighted spectral gap.  It is a
weighted diagonal-dominance condition.

\begin{corollary}[Derivative of a quasi-local inverse]
\label{cor:v11v61-derivative}
If \(A^{-1}\) obeys (61.22) and a perturbation \(\dot A\) has finite
weighted Schur norm, then
\[
                  \frac d{dt}(A+t\dot A)^{-1}\big|_{t=0}
                             =-A^{-1}\dot A A^{-1}             \tag{61.24}
\]
belongs to the same weighted algebra, with norm at most
\[
              \lVert A^{-1}\rVert_{{\cal S}_\mu}^2
                    \lVert\dot A\rVert_{{\cal S}_\mu}.        \tag{61.25}
\]
\end{corollary}

\begin{proof}
Use (60.2) and the algebra property (61.18).
\end{proof}

\subsection{Massless and gauge modes}

The condition \(q<1\) fails when \(D+K\) has a massless collective mode
or when an exact gauge mode has not been removed.  Based tree gauge and
the constraint quotient remove specified exact modes, but physical
massless photon and graviton channels remain.

\begin{warning}[No exponential decay from locality alone]
\label{warn:v11v61-massless}
A finite-range operator with a vanishing spectral gap generally has a
polynomial Green kernel.  V61 does not assign exponential decay to the
hypercharge, photon, graviton, or unbroken constraint channels.  Those
sectors require infrared weights, finite-volume zero-mode conditions,
or polynomial versions of the determinant schedule.
\end{warning}

\subsection{Insertion into V59}

\begin{ledger}[Quasi-local counterterm data]
\label{led:v11v61-data}
To use the V59 bulk--boundary theorem for an inverse-generated
relative operator, one must prove:
\begin{enumerate}
\item exponential or polynomial row and column tail bounds;
\item a truncation radius satisfying (61.14) or (61.16);
\item summability of the chosen trace errors \(\eta_i\);
\item V59 boundary summability with the enlarged radius \(R_i\);
\item a Schatten-\(p\) bound for the untruncated or truncated remainder;
\item separate infrared treatment of massless and residual modes.
\end{enumerate}
\end{ledger}

\begin{firewall}[No acquired physical decay rate]
\label{fire:v11v61-boundary}
The weighted Neumann theorem proves exponential decay only under
(61.20)--(61.21).  Those inequalities have not been verified for the
full SM--Einstein Hessian, chiral operator, or constraint carrier.
The polynomial alternative is a truncation theorem conditional on a
kernel bound; it does not derive that bound from the target dynamics.
\end{firewall}

\begin{decision}[Next acquisition]
\label{dec:v11v61-next}
V62 will isolate infrared modes by a low/high spectral split.  It will
show that a high-sector gap gives quasi-local determinant control while
the finite or controlled low sector is retained as an explicit
cylinder variable.  This avoids assigning a false mass to photon or
graviton modes.
\end{decision}

\begin{assessment}[Progress after V61]
\label{ass:v11v61-status}
Finite-range counterterm analysis now extends to quasi-local kernels
with quantified truncation errors.  Exponential decay gives logarithmic
range growth, polynomial decay gives a potentially volume-sized range,
and a weighted local-core criterion supplies one rigorous source of
exponential resolvent decay.  Infrared modes remain the next separate
obstruction.
\end{assessment}

\subsection{Parent record}

This V61 note was formed by appending the preceding material to the
validated V60 source, whose record was
\[
\begin{array}{c|c}
\text{lines}&18852\\
\text{bytes}&586526\\
\text{SHA--256}&
\texttt{0D602EBA5497FB9AEE9DB067D0D6D0B0BACA402B110154B0615E5366D175DB67}.
\end{array}                                                     \tag{61.26}
\]

% =============================================================================
% END APPENDED ELEVENTH-SERIES V61 MATERIAL
% =============================================================================
% =============================================================================
% BEGIN APPENDED ELEVENTH-SERIES V62 MATERIAL
% =============================================================================

\section{Infrared retention and exact high-mode Schur reduction}
\label{sec:v11v62}

Massless photon and graviton channels prevent a uniform inverse bound
on the complete quadratic operator.  The correct finite-cutoff
operation is to retain a low-mode sector and integrate only a
gapped high-mode block.  The determinant and inverse then factor
through an exact Schur complement.  No mass is inserted into the
retained sector.

\subsection{The Feshbach--Schur card}

Let \(E=E_L\oplus E_H\) and write an operator as
\[
 A=\begin{pmatrix}
       A_{LL}&B\\
       C&D
   \end{pmatrix}.                                             \tag{62.1}
\]
Assume \(D:E_H\to E_H\) is invertible and define
\[
                              S=A_{LL}-BD^{-1}C.               \tag{62.2}
\]

\begin{theorem}[Exact high-mode elimination]
\label{thm:v11v62-schur}
One has the block factorization
\[
 A=
 \begin{pmatrix}I&BD^{-1}\\0&I\end{pmatrix}
 \begin{pmatrix}S&0\\C&D\end{pmatrix}.                        \tag{62.3}
\]
Consequently,
\[
                              \det A=\det D\,\det S,           \tag{62.4}
\]
and \(A\) is invertible exactly when \(S\) is invertible.  In that
case,
\[
 A^{-1}=
 \begin{pmatrix}
 S^{-1}&-S^{-1}BD^{-1}\\
 -D^{-1}CS^{-1}&
 D^{-1}+D^{-1}CS^{-1}BD^{-1}
 \end{pmatrix}.                                               \tag{62.5}
\]
\end{theorem}

\begin{proof}
Direct block multiplication proves (62.3).  The first factor has
determinant one, while the second is block triangular, proving
(62.4).  If \(A(u,v)=0\), the high equation gives
\(v=-D^{-1}Cu\), and the low equation becomes \(Su=0\).
This proves the invertibility equivalence.  Solving the two block
equations for arbitrary right-hand side gives (62.5).
\end{proof}

For self-adjoint \(A\), one has \(C=B^*\), \(D=D^*\), and
\(S=S^*\).  Small eigenvalues are allowed in \(S\); only \(D\) is
required to have a gap.

\begin{corollary}[High-gap bounds]
\label{cor:v11v62-gap}
If \(\lVert D^{-1}\rVert\le\gamma^{-1}\), then
\[
                  \lVert S-A_{LL}\rVert
                       \le\gamma^{-1}\lVert B\rVert\lVert C\rVert.
                                                                  \tag{62.6}
\]
No inverse of \(S\) is needed to define the effective low operator.
\end{corollary}

\begin{proof}
Use (62.2) and submultiplicativity.
\end{proof}

\subsection{Exact Gaussian reduction}

\begin{theorem}[Bosonic high-mode integral]
\label{thm:v11v62-gaussian}
Let \(A=A^*>0\) on a finite-dimensional real space
\(E_L\oplus E_H\), with \(\dim E_H=N_H\).  For fixed \(u\in E_L\),
\[
\begin{split}
 &\int_{E_H}
 \exp\left\{-\frac12
 \left\langle\binom uv,
 A\binom uv\right\rangle\right\}\,dv\\
 &\hspace{2cm}
 =(2\pi)^{N_H/2}(\det D)^{-1/2}
                    \exp\{-\tfrac12\langle u,Su\rangle\}.     \tag{62.7}
\end{split}
\]
\end{theorem}

\begin{proof}
Complete the square:
\[
 \left\langle\binom uv,A\binom uv\right\rangle
 =\langle u,Su\rangle+
   \langle v+D^{-1}Cu,D(v+D^{-1}Cu)\rangle.                  \tag{62.8}
\]
Translate \(v\) and apply the standard positive Gaussian integral.
\end{proof}

For a complex fermionic quadratic form, finite-dimensional Berezin
integration gives the corresponding algebraic factor
\[
                              \det D                           \tag{62.9}
\]
and leaves the fermionic Schur operator \(S\) on the retained modes.
Equation (62.4) preserves the complex phase; taking absolute values
before the factorization is not permitted in the chiral sector.

\subsection{Stable adjacent low operator}

Consider two block matrices with primed and unprimed blocks.  Suppose
\[
 \lVert D^{-1}\rVert,\lVert(D')^{-1}\rVert\le\gamma^{-1},     \tag{62.10}
\]
and write
\[
 \delta_A=\lVert A'_{LL}-A_{LL}\rVert,\quad
 \delta_B=\lVert B'-B\rVert,\quad
 \delta_C=\lVert C'-C\rVert,\quad
 \delta_D=\lVert D'-D\rVert.                                 \tag{62.11}
\]

\begin{proposition}[Schur-complement increment]
\label{prop:v11v62-increment}
If
\[
 \lVert B\rVert,\lVert B'\rVert\le M_B,\qquad
 \lVert C\rVert,\lVert C'\rVert\le M_C,                       \tag{62.12}
\]
then
\[
 \lVert S'-S\rVert
 \le\delta_A+\gamma^{-1}(M_C\delta_B+M_B\delta_C)
                   +M_BM_C\gamma^{-2}\delta_D.                \tag{62.13}
\]
\end{proposition}

\begin{proof}
Insert and subtract \(B'D^{-1}C'\):
\[
\begin{split}
 B'(D')^{-1}C'-BD^{-1}C
 &=(B'-B)(D')^{-1}C'\\
 &\quad+B\bigl((D')^{-1}-D^{-1}\bigr)C'
       +BD^{-1}(C'-C).
\end{split}
\]
Use
\[
              (D')^{-1}-D^{-1}=-(D')^{-1}(D'-D)D^{-1}
\]
and (62.10)--(62.12).  Relabeling the endpoint bounds symmetrically
gives (62.13).
\end{proof}

If the low rank \(r=\dim E_L\) is fixed and both low Schur operators
have singular value at least \(\sigma\), then their log determinants
obey
\[
                   |\Delta\log\det S|
          \le r\sigma^{-1}\lVert S'-S\rVert                  \tag{62.14}
\]
along any segment retaining the same floor.  If \(S\) is intentionally
massless, its determinant is not used; the low field remains explicit
in the cylinder measure.

\subsection{High relative determinant}

At adjacent cutoffs, apply V58 to the high blocks:
\[
 D_{i+1}
 =\widehat D_i^{1/2}(I+T_i^H)\widehat D_i^{1/2},              \tag{62.15}
\]
in the positive bosonic case, or the corresponding two-sided
invertible factorization for a complex fermionic block.  Trace-class
summability of \(T_i^H\), or the V59 regularized alternative, controls
the high determinant.  The complete finite determinant is
\[
                         \det A_i=\det D_i\,\det S_i.          \tag{62.16}
\]
Thus high ultraviolet renormalization and low infrared dynamics are
separate factors.

\begin{proposition}[No false full gap]
\label{prop:v11v62-no-false-gap}
A uniform bound on \(D_i^{-1}\) and a convergent relative determinant
for \(D_i\) make no assertion that \(A_i\) has a uniform gap.
The complete operator is singular exactly when \(S_i\) is singular.
\end{proposition}

\begin{proof}
This is the invertibility equivalence in
Theorem~\ref{thm:v11v62-schur}.
\end{proof}

\subsection{Choosing the low sector}

Let \(A_i^0\) be a self-adjoint reference operator, and let \(P_i\) be
its spectral projection onto \(|\lambda|\le\Lambda_i\), with
\(Q_i=I-P_i\).  Suppose the actual operator is
\(A_i=A_i^0+V_i\) and
\[
             \inf\{|\lambda|:\lambda\in
               \operatorname{spec}(Q_iA_i^0Q_i)\}
                           \ge\Gamma_i,\qquad
             \lVert Q_iV_iQ_i\rVert\le\eta_i<\Gamma_i.       \tag{62.17}
\]

\begin{lemma}[Reference spectral split]
\label{lem:v11v62-reference}
The high block \(D_i=Q_iA_iQ_i\) is invertible and
\[
                         \lVert D_i^{-1}\rVert
                        \le(\Gamma_i-\eta_i)^{-1}.            \tag{62.18}
\]
\end{lemma}

\begin{proof}
The least singular value changes by at most the operator norm of the
perturbation.  Hence
\(s_{\min}(D_i)\ge\Gamma_i-\eta_i>0\).
\end{proof}

The low projection may be chosen by momentum, by a reference spectral
window, or by explicit harmonic and gauge modes.  Its compatibility
with locality is not automatic: a sharp spectral projection is
generally nonlocal.

\subsection{The infrared counting tradeoff}

On a \(d\)-dimensional periodic box of volume \(V_i\), the number of
Fourier modes with \(|p|\le\Lambda_i\) is of order
\[
                               r_i\asymp V_i\Lambda_i^d        \tag{62.19}
\]
away from exceptional degeneracies.  Keeping \(r_i\) bounded forces
\(\Lambda_i\lesssim V_i^{-1/d}\), but then the high gap of a
first-order massless operator is also \(O(V_i^{-1/d})\).

\begin{warning}[Low-rank versus high-gap balance]
\label{warn:v11v62-balance}
For a massless field, one cannot generally keep both the retained rank
and the inverse high-sector gap uniformly bounded as the volume grows.
The factors \(r_i\) and \(\gamma_i^{-1}\) must both enter the adjacent
schedule.  A fixed positive mass would remove this tradeoff but would
change the target photon or graviton theory.
\end{warning}

A multiscale cylinder construction can instead retain successively
more low modes while integrating high shells one at a time.  Then the
relevant requirement is summability of shell Schur increments, not a
single volume-independent high gap.

\subsection{Gauge, constraint, and topology modes}

Exact internal gauge modes are removed by V46 before the split.
First-class gravitational modes are removed locally by V48--V53 on
regular quotient charts.  Remaining low modes can include:
\begin{enumerate}
\item physical photon and graviton momenta;
\item torons and noncontractible Wilson lines;
\item harmonic representatives of constraint complexes;
\item Higgs or fermion zero modes in special backgrounds;
\item modes approaching a singular quotient stratum.
\end{enumerate}
Each class has a different measure and transition law and must be
identified rather than counted under one generic infrared rank.

\begin{ledger}[Infrared-retention data]
\label{led:v11v62-data}
At every cutoff the programme must now provide:
\begin{enumerate}
\item projections \(P_i,Q_i\) compatible with the retained cylinder
algebra;
\item the high-block gap \(\Gamma_i-\eta_i\);
\item decay and relative determinant bounds for \(D_i^{-1}\);
\item low--high coupling bounds \(B_i,C_i\);
\item adjacent Schur estimates from (62.13);
\item a projectively consistent measure for the growing low sector;
\item special treatment of exact zero modes and determinant phases.
\end{enumerate}
\end{ledger}

\begin{firewall}[No infrared solution claimed]
\label{fire:v11v62-boundary}
V62 gives an exact finite-dimensional separation.  It does not prove
the required high gap for the SM--Einstein operator, locality of a
chosen spectral projection, convergence of the low-mode measure, or
clustering in massless sectors.  It does not turn photon or graviton
modes into massive modes.
\end{firewall}

\begin{decision}[Next acquisition]
\label{dec:v11v62-next}
V63 will construct the projective low-mode Gaussian card.  Given
compatible covariance compressions and summable Schur corrections, it
will prove consistency of finite-dimensional characteristic
functionals and state the tightness condition needed for a distribution-
valued massless limit.
\end{decision}

\begin{assessment}[Progress after V62]
\label{ass:v11v62-status}
Ultraviolet determinant control and infrared dynamics are now
separated exactly.  The high block can use the quasi-local and
relative-determinant machinery of V58--V61, while every near-zero mode
is retained in the Schur complement.  The unresolved task is the
projectively consistent and tight limit of that growing low sector.
\end{assessment}

\subsection{Parent record}

This V62 note was formed by appending the preceding material to the
validated V61 source, whose record was
\[
\begin{array}{c|c}
\text{lines}&19209\\
\text{bytes}&597978\\
\text{SHA--256}&
\texttt{CB0A41907E641EC779A2E8E0A784EA0C6D22D2497A84CB99E93AA07049A90DD4}.
\end{array}                                                     \tag{62.20}
\]

% =============================================================================
% END APPENDED ELEVENTH-SERIES V62 MATERIAL
% =============================================================================
% =============================================================================
% BEGIN APPENDED ELEVENTH-SERIES V63 MATERIAL
% =============================================================================

\section{Projective Gaussian low modes and distributional tightness}
\label{sec:v11v63}

The Schur complement of V62 leaves a growing low-mode sector.  This
section gives the exact consistency criterion for centered Gaussian
measures, a summable-defect replacement, and a separate tightness
criterion in a Hilbert distribution space.  These statements concern
a positive Gaussian reference sector; interacting and indefinite
gravitational weights require later transfer arguments.

\subsection{Exact finite-dimensional consistency}

Let
\[
                         E_1\subset E_2\subset\cdots           \tag{63.1}
\]
be nested finite-dimensional real Hilbert spaces, with orthogonal
projections \(P_i:E_{i+1}\to E_i\).  Let \(\mu_i\) be a centered
Gaussian probability on \(E_i\) with covariance \(C_i\ge0\):
\[
 \widehat\mu_i(f)
 :=\int_{E_i}e^{i\langle f,x\rangle}\,d\mu_i(x)
 =\exp\{-\tfrac12\langle f,C_if\rangle\}.                     \tag{63.2}
\]

\begin{theorem}[Gaussian projective criterion]
\label{thm:v11v63-projective}
The measures satisfy
\[
                              (P_i)_*\mu_{i+1}=\mu_i           \tag{63.3}
\]
if and only if
\[
                              P_iC_{i+1}P_i^*=C_i.             \tag{63.4}
\]
Under (63.4), there is a unique probability measure on the coordinate
inverse limit whose \(E_i\)-marginal is \(\mu_i\).
\end{theorem}

\begin{proof}
For \(f\in E_i\), the characteristic function of
\((P_i)_*\mu_{i+1}\) is
\[
 \exp\{-\tfrac12
       \langle P_i^*f,C_{i+1}P_i^*f\rangle\}
 =\exp\{-\tfrac12
       \langle f,P_iC_{i+1}P_i^*f\rangle\}.
\]
Equality with (63.2) for every \(f\) is equivalent to (63.4), because
quadratic forms determine self-adjoint operators.  The final assertion
is the Kolmogorov extension theorem applied to the consistent finite
marginals.
\end{proof}

The inverse-limit measure initially lives on a coordinate product.
Additional tightness is required to show that it is supported in a
specified distribution space.

\subsection{Summable consistency defects}

Exact equality (63.4) is stronger than the adjacent programme normally
produces.  Let \(J_{i,j}:E_j\to E_i\) be the inclusion for \(i\ge j\),
and define the covariance compressed from level \(i\) to \(E_j\) by
\[
                         C_i^{(j)}=J_{i,j}^*C_iJ_{i,j}.        \tag{63.5}
\]

\begin{theorem}[Cylinder covariance limit]
\label{thm:v11v63-defect}
Suppose for every fixed \(j\) there are nonnegative
\(\epsilon_i^{(j)}\), \(i\ge j\), such that
\[
 \lVert C_{i+1}^{(j)}-C_i^{(j)}\rVert
                       \le\epsilon_i^{(j)},\qquad
 \sum_{i=j}^{\infty}\epsilon_i^{(j)}<\infty.                 \tag{63.6}
\]
Then \(C_i^{(j)}\) converges in operator norm to a positive operator
\(C_\infty^{(j)}\).  The limits are mutually compatible:
\[
                         P_{k,j}C_\infty^{(k)}P_{k,j}^*
                              =C_\infty^{(j)},\qquad j\le k,   \tag{63.7}
\]
where \(P_{k,j}:E_k\to E_j\).  Hence
\[
 \Phi_\infty(f)
 =\exp\{-\tfrac12\langle f,C_\infty^{(j)}f\rangle\},
 \qquad f\in E_j,                                             \tag{63.8}
\]
is a well-defined positive-definite characteristic functional on
\(\bigcup_jE_j\), and determines consistent Gaussian cylinder
marginals.
\end{theorem}

\begin{proof}
Equation (63.6) makes \(C_i^{(j)}\) Cauchy in the finite-dimensional
operator norm, hence convergent.  Positivity is closed under norm
limits.  For \(j\le k\), finite-level compressions satisfy
\[
                 P_{k,j}C_i^{(k)}P_{k,j}^*=C_i^{(j)}
\]
by their definition from the same \(C_i\).  Passing to the limit gives
(63.7).  Each finite-dimensional function in (63.8) is a centered
Gaussian characteristic function.  Compatibility makes its value
independent of the containing \(E_j\) and gives consistent cylinder
marginals by Theorem~\ref{thm:v11v63-projective}.
\end{proof}

\begin{corollary}[Quantitative characteristic increment]
\label{cor:v11v63-characteristic}
For \(f\in E_j\),
\[
 |\widehat\mu_{i+1}(J_{i+1,j}f)
       -\widehat\mu_i(J_{i,j}f)|
 \le\frac12\lVert f\rVert^2\epsilon_i^{(j)}.                 \tag{63.9}
\]
\end{corollary}

\begin{proof}
For \(a,b\ge0\),
\(|e^{-a/2}-e^{-b/2}|\le\frac12|a-b|\).  Apply this to the two
covariance quadratic forms and use (63.6).
\end{proof}

This is a direct Gaussian instance of the V42 adjacent cylinder
criterion.

\subsection{Covariance from the Schur operator}

If a positive low Schur operator \(S_i>0\) is integrated as a Gaussian,
then
\[
                              C_i=S_i^{-1}.                    \tag{63.10}
\]
For identified retained spaces, the resolvent identity gives
\[
 C_{i+1}^{(j)}-C_i^{(j)}
 =-\,C_{i+1}^{(j)}
       \bigl(S_{i+1}^{(j)}-S_i^{(j)}\bigr)C_i^{(j)}            \tag{63.11}
\]
when the compressed inverse relation is exact.

\begin{proposition}[Schur-to-covariance summability]
\label{prop:v11v63-resolvent}
Suppose on a fixed cylinder \(E_j\),
\[
 \lVert (S_i^{(j)})^{-1}\rVert\le K_j
 \quad(i\ge j),\qquad
 \sum_{i=j}^{\infty}
     \lVert S_{i+1}^{(j)}-S_i^{(j)}\rVert<\infty.             \tag{63.12}
\]
Then (63.6) holds with
\[
 \epsilon_i^{(j)}
 =K_j^2\lVert S_{i+1}^{(j)}-S_i^{(j)}\rVert.                 \tag{63.13}
\]
\end{proposition}

\begin{proof}
Apply the resolvent identity and submultiplicativity.
\end{proof}

For genuinely massless fields, the bound \(K_j\) can depend on the
fixed cylinder scale \(j\).  Cylinder convergence permits this;
uniform tightness is a separate question.

\subsection{Tightness in a distribution Hilbert space}

Let \({\cal H}\) be a separable Hilbert space containing every \(E_i\)
through isometries, and regard \(\mu_i\) as measures on \({\cal H}\)
supported on \(E_i\).  Let \(Q_i\) be their covariance operators
extended by zero.

\begin{theorem}[Spectral-tail tightness]
\label{thm:v11v63-tight}
Let \((e_n)\) be an orthonormal basis of \({\cal H}\), and let
\(\Pi_N\) project onto its first \(N\) vectors.  If
\[
 \sup_i\operatorname{Tr}Q_i<\infty,\qquad
 \lim_{N\to\infty}\sup_i
       \operatorname{Tr}\bigl((I-\Pi_N)Q_i\bigr)=0,           \tag{63.14}
\]
then \((\mu_i)\) is tight on \({\cal H}\).
\end{theorem}

\begin{proof}
Choose increasing \(N_k\) so that
\[
 \sup_i\operatorname{Tr}((I-\Pi_{N_k})Q_i)
                         \le 2^{-3k}\varepsilon.
\]
For \(X_i\sim\mu_i\), Markov's inequality gives
\[
 \mathbb P\{\lVert(I-\Pi_{N_k})X_i\rVert>2^{-k}\}
 \le2^{2k}\operatorname{Tr}((I-\Pi_{N_k})Q_i)
 \le2^{-k}\varepsilon.                                      \tag{63.15}
\]
The finite-dimensional heads \(\Pi_{N_1}X_i\) have uniformly bounded
second moments by the first condition, so choose a common radius
containing them with probability at least \(1-\varepsilon\).
The set of vectors with that head bound and all tail bounds in
(63.15) is totally bounded: truncate at \(N_k\) and use the
\(2^{-k}\) tail.  Its closure is compact.  A union bound shows that it
has \(\mu_i\)-mass at least \(1-2\varepsilon\), proving tightness.
\end{proof}

\begin{corollary}[Compact Sobolev embedding criterion]
\label{cor:v11v63-sobolev}
Let \({\cal H}_0\) embed compactly into \({\cal H}\).  If
\[
                         \sup_i\mathbb E\lVert X_i\rVert_{{\cal H}_0}^2
                              <\infty,                         \tag{63.16}
\]
then \((\mu_i)\) is tight on \({\cal H}\).
\end{corollary}

\begin{proof}
Closed \({\cal H}_0\)-balls are relatively compact in \({\cal H}\).
Markov's inequality gives a common ball with arbitrarily large
probability.
\end{proof}

On a compact spatial domain one may take
\({\cal H}_0=H^{-s_0}\) and \({\cal H}=H^{-s}\) with \(s>s_0\);
the embedding is compact.  The required moment is
\[
              \mathbb E\lVert X_i\rVert_{H^{-s_0}}^2
       =\operatorname{Tr}\bigl((1+\Delta)^{-s_0}C_i\bigr),    \tag{63.17}
\]
under the chosen mode identification.

\subsection{Limit identification}

\begin{theorem}[Tight Gaussian realization]
\label{thm:v11v63-realization}
Assume the cylinder covariance limits of
Theorem~\ref{thm:v11v63-defect} and tightness on \({\cal H}\).
Then \(\mu_i\) converges weakly on \({\cal H}\) to the unique centered
Gaussian probability whose covariance quadratic form on
\(\bigcup_jE_j\) is \(C_\infty\).
\end{theorem}

\begin{proof}
Tightness gives weakly convergent subsequences by Prokhorov's theorem.
For every cylinder vector \(f\), (63.8)--(63.9) give convergence of
the characteristic functions to \(\Phi_\infty(f)\).  Every
subsequential limit therefore has the same characteristic functional
on a dense set.  Continuity extends equality to \({\cal H}\), and
characteristic functions determine Borel probabilities on a separable
Hilbert space.  All subsequences have the same limit, so the full
sequence converges.
\end{proof}

\subsection{Exact zero modes}

If \(S_i\) has an exact zero mode on a noncompact linear field, the
Lebesgue Gaussian normalization is infinite and (63.10) is undefined.

\begin{proposition}[Zero-mode alternatives]
\label{prop:v11v63-zero}
A normalized Gaussian treatment of an exact zero mode requires one of
the following additional structures:
\begin{enumerate}
\item quotienting a genuine gauge direction;
\item fixing or conditioning the zero coordinate;
\item compactifying it with a finite invariant measure;
\item adding a source-neutrality condition that removes it from the
observable algebra;
\item proving an interacting potential that is integrable in that
direction.
\end{enumerate}
A positive mass regulator is a sixth finite-cutoff option, but its
removal requires a separate uniform limit.
\end{proposition}

\begin{proof}
Along an unmodified zero direction the quadratic density is constant,
so its integral against Lebesgue measure is infinite.  Each listed
operation changes the space, measure, observable algebra, or action so
that this infinite factor is absent or replaced by an integrable one.
Without one of them, no probability normalization exists.
\end{proof}

\subsection{SM--Einstein low-mode ledger}

\begin{ledger}[Projective infrared data]
\label{led:v11v63-data}
For each physical low sector the programme must specify:
\begin{enumerate}
\item nested retained spaces and projection maps;
\item positive reference covariances or a different constructive
measure for indefinite sectors;
\item summable compressed covariance defects (63.6);
\item a distribution space and one of the tightness criteria
(63.14) or (63.16);
\item exact zero-mode treatment from
Proposition~\ref{prop:v11v63-zero};
\item compatibility with global holonomy and topological sectors;
\item transfer from the positive reference to the interacting chiral
SM--Einstein weight.
\end{enumerate}
\end{ledger}

\begin{firewall}[Gaussian reference boundary]
\label{fire:v11v63-boundary}
V63 constructs a positive Gaussian low-mode limit under covariance and
tightness hypotheses.  The Einstein conformal direction is not a
positive Gaussian under the target Euclidean action, fermionic
variables require Berezin rather than probability measures, and the
interacting chiral phase still requires the stable denominator of
V38.  None of these target transfers is proved here.
\end{firewall}

\begin{decision}[Next acquisition]
\label{dec:v11v63-next}
V64 will treat reflection positivity and clustering for the assembled
reference measure.  It will prove that reflection positivity passes
under cylinder convergence, while clustering requires a separate
uniform covariance or connected-correlation decay bound and fails for
uncontrolled retained zero modes.
\end{decision}

\begin{assessment}[Progress after V63]
\label{ass:v11v63-status}
The retained infrared sector now has a rigorous projective Gaussian
route: summable covariance defects determine all cylinder marginals,
and a spectral-tail or compact-Sobolev estimate realizes them as a
tight distribution-valued probability.  Exact zero modes and transfer
to the interacting SM--Einstein measure remain explicit acquisitions.
\end{assessment}

\subsection{Parent record}

This V63 note was formed by appending the preceding material to the
validated V62 source, whose record was
\[
\begin{array}{c|c}
\text{lines}&19545\\
\text{bytes}&609400\\
\text{SHA--256}&
\texttt{FAC208E151550CD867C24DB2C8050FA0DE69AA7D6199D0343B98C91E685F6572}.
\end{array}                                                     \tag{63.18}
\]

% =============================================================================
% END APPENDED ELEVENTH-SERIES V63 MATERIAL
% =============================================================================
% =============================================================================
% BEGIN APPENDED ELEVENTH-SERIES V64 MATERIAL
% =============================================================================

\section{Reflection positivity, marginal stability, and clustering
criteria}
\label{sec:v11v64}

The cylinder measure obtained in V63 is constructive only if the
Euclidean positivity and long-distance axioms can be passed to the
limit.  Reflection positivity is closed under convergence of the
relevant cylinder expectations.  Clustering is not: it requires a
separate connected-correlation or transfer-gap estimate.

\subsection{Reflection-positive cylinder forms}

Let \(\theta\) be a time reflection acting on the field space, and let
\({\cal A}_+\) be the algebra of bounded continuous cylinder functions
supported in the positive-time half.  Define the antilinear reflection
\[
                         (\Theta F)(\phi)
                    =\overline{F(\theta\phi)}.                 \tag{64.1}
\]
For a probability \(\mu\), put
\[
                         {\cal Q}_\mu(F)
                    =\int(\Theta F)(\phi)F(\phi)\,d\mu(\phi). \tag{64.2}
\]

\begin{definition}[Reflection positivity]
\label{def:v11v64-rp}
The measure \(\mu\) is reflection positive on \({\cal A}_+\) if
\[
                              {\cal Q}_\mu(F)\ge0
                              \quad(F\in{\cal A}_+).           \tag{64.3}
\]
\end{definition}

\begin{theorem}[Closure under cylinder convergence]
\label{thm:v11v64-limit}
Let \(\mu_i\) be reflection-positive probabilities on compatible
finite-cutoff field spaces.  Suppose that for every bounded continuous
cylinder function \(G\),
\[
                              \int G\,d\mu_i\longrightarrow
                              \int G\,d\mu.                    \tag{64.4}
\]
Assume also that \(\Theta\) preserves the common cylinder algebra.
Then \(\mu\) is reflection positive.
\end{theorem}

\begin{proof}
For \(F\in{\cal A}_+\), the product \((\Theta F)F\) is another bounded
continuous cylinder function.  Hence
\[
 {\cal Q}_\mu(F)
 =\lim_i{\cal Q}_{\mu_i}(F)\ge0.
\]
\end{proof}

\begin{warning}[Unbounded fields]
\label{warn:v11v64-unbounded}
For polynomial or other unbounded cylinder functions, weak or bounded-
cylinder convergence is insufficient.  One also needs uniform
integrability of \((\Theta F)F\), supplied for example by a common
moment strictly above its degree.
\end{warning}

\subsection{Marginalization and the high-mode split}

\begin{proposition}[Reflection-positive marginals]
\label{prop:v11v64-marginal}
Let \(\mu\) be reflection positive and let
\(\pi\) be a measurable retained-field map commuting with reflection:
\[
                              \pi\theta=\theta\pi.             \tag{64.5}
\]
Then the marginal \(\nu=\pi_*\mu\) is reflection positive on every
positive-time function of the retained variables.
\end{proposition}

\begin{proof}
For \(F\) on the retained space, \(F\circ\pi\in{\cal A}_+\) and
\[
 {\cal Q}_\nu(F)
 =\int(\Theta F)(\pi\phi)F(\pi\phi)\,d\mu(\phi)
 ={\cal Q}_\mu(F\circ\pi)\ge0.
\]
\end{proof}

Thus exact high-mode integration in V62 preserves reflection
positivity when the low/high projection commutes with reflection.
A sharp momentum split that mixes positive and negative time need not
satisfy (64.5).

\begin{proposition}[Safe factorized reweighting]
\label{prop:v11v64-reweight}
Let \(\mu\) be reflection positive and let \(V_+\) be a real
positive-time cylinder function such that
\[
                         0<Z=\int e^{-V_+-\Theta V_+}\,d\mu<\infty.
                                                                  \tag{64.6}
\]
Then
\[
                         d\nu=Z^{-1}e^{-V_+-\Theta V_+}\,d\mu \tag{64.7}
\]
is reflection positive.
\end{proposition}

\begin{proof}
For \(F\in{\cal A}_+\),
\[
 {\cal Q}_\nu(F)
 =Z^{-1}{\cal Q}_\mu(e^{-V_+}F)\ge0.                          \tag{64.8}
\]
\end{proof}

Cross-plane interactions require their own positive decomposition;
they are not covered merely because the action is real.  A complex
chiral determinant is not a positive density, so its Euclidean
fermionic positivity must be formulated through the regulated
fermionic algebra and its reflection operation rather than by applying
Proposition~\ref{prop:v11v64-reweight} to a complex scalar weight.

\subsection{Gaussian clustering}

Let \(\mu_C\) be a centered real Gaussian field with covariance
bilinear form \(C\).  For real test functions define Weyl cylinders
\[
                              W(f)=e^{i\phi(f)}.               \tag{64.9}
\]
Their expectations are
\[
                              {\mathbb E}W(f)
                              =e^{-C(f,f)/2}.                 \tag{64.10}
\]

\begin{theorem}[Gaussian Weyl clustering]
\label{thm:v11v64-gaussian}
Let \(\tau_xg\) denote a spatial translate.  Then
\[
\begin{split}
 &{\mathbb E}[W(f)W(\tau_xg)]
       -{\mathbb E}W(f)\,{\mathbb E}W(g)\\
 &\quad=e^{-\frac12[C(f,f)+C(g,g)]}
             \left(e^{-C(f,\tau_xg)}-1\right).                \tag{64.11}
\end{split}
\]
Consequently, if \(C(f,\tau_xg)\to0\) as \(|x|\to\infty\), the Weyl
observables cluster.  Quantitatively,
\[
 |\operatorname{Cov}(W(f),W(\tau_xg))|
 \le e^{|C(f,\tau_xg)|}|C(f,\tau_xg)|.                       \tag{64.12}
\]
\end{theorem}

\begin{proof}
The product is \(W(f+\tau_xg)\).  Insert (64.10), expand the quadratic
form, and factor the two one-point expectations to obtain (64.11).
The inequality
\(|e^{-z}-1|\le e^{|z|}|z|\) for real \(z\) gives (64.12).
\end{proof}

Wick's theorem extends the conclusion to polynomial cylinders when
all translated cross-covariances tend to zero and the required moments
are finite.

\begin{corollary}[Inherited decay rate]
\label{cor:v11v64-rate}
If for fixed compactly supported tests
\[
 |C(f,\tau_xg)|\le A_{f,g}e^{-m|x|},                         \tag{64.13}
\]
then the connected Weyl correlation decays exponentially with rate
\(m\), up to a bounded prefactor.  A polynomial covariance bound gives
the same polynomial clustering order.
\end{corollary}

\begin{proof}
Insert the covariance bound in (64.12).  The exponential prefactor is
uniformly bounded because the cross-covariance is bounded by
Cauchy--Schwarz.
\end{proof}

\subsection{Reflection positivity does not imply clustering}

\begin{proposition}[Persistent-phase counterexample]
\label{prop:v11v64-counterexample}
Let a lattice field be the random constant
\[
                              \phi_x=X,\qquad
                 \mathbb P\{X=1\}=\mathbb P\{X=-1\}=\tfrac12. \tag{64.14}
\]
Its law is translation invariant and reflection positive, but it does
not cluster:
\[
                    {\mathbb E}\phi_0=0,\qquad
                    {\mathbb E}\phi_0\phi_x=1                 \tag{64.15}
\]
for every \(x\).
\end{proposition}

\begin{proof}
Reflection leaves a constant configuration fixed.  Hence for every
positive-time function \(F\),
\[
                         {\cal Q}(F)={\mathbb E}|F(X)|^2\ge0.
\]
Equation (64.15) is immediate.  Its connected two-point function never
tends to zero.
\end{proof}

This example is also a zero-mode warning: tightness and reflection
positivity can coexist with a nondecaying global mode.

\subsection{Transfer-operator spectral criterion}

Reflection positivity can be used to construct a Hilbert space by
quotienting \({\cal A}_+\) by the null space of (64.2).  Assume a unit
time translation induces a self-adjoint contraction \(T\) on this
space and fixes the vacuum vector \(\Omega\).

\begin{theorem}[Transfer gap implies exponential clustering]
\label{thm:v11v64-transfer}
Suppose
\[
              \lVert T|_{\Omega^\perp}\rVert\le e^{-m}<1.     \tag{64.16}
\]
For vectors \(a,b\) generated by bounded cylinder observables,
\[
 |\langle a,T^nb\rangle
       -\langle a,\Omega\rangle\langle\Omega,b\rangle|
 \le e^{-mn}\lVert a-\langle\Omega,a\rangle\Omega\rVert
             \lVert b-\langle\Omega,b\rangle\Omega\rVert.     \tag{64.17}
\]
\end{theorem}

\begin{proof}
Decompose \(a\) and \(b\) by the orthogonal projection onto
\(\operatorname{span}\{\Omega\}\).  Since \(T\Omega=\Omega\) and \(T\)
is self-adjoint, the connected expression is the matrix element of
\(T^n\) between the two orthogonal components.  Apply Cauchy--Schwarz
and (64.16).
\end{proof}

The converse identification of a physical mass gap requires the full
Osterwalder--Schrader reconstruction hypotheses and control of spatial
momentum; it is not supplied by (64.16) alone.

\subsection{Adjacent passage of decay}

Suppose finite-cutoff connected correlations obey
\[
                      |G_i^c(A,\tau_xB)|\le h(x)              \tag{64.18}
\]
with one \(h(x)\to0\), uniformly in \(i\), and suppose both the joint
and one-point cylinder expectations converge.  Then
\[
                      |G_\infty^c(A,\tau_xB)|\le h(x).        \tag{64.19}
\]

\begin{proposition}[Uniform clustering passes to the limit]
\label{prop:v11v64-passage}
Under the preceding hypotheses, the limiting measure clusters for
\(A,B\).  If \(h(x)\le Ce^{-m|x|}\), it inherits that exponential
bound.
\end{proposition}

\begin{proof}
For each fixed \(x\), pass to the limit in the finite-cutoff connected
correlation.  Inequality (64.19) follows.  Then let \(|x|\to\infty\).
\end{proof}

Pointwise clustering at each cutoff without a common \(h\) is
insufficient because the separation scale can diverge with the cutoff.

\subsection{SM--Einstein positivity and clustering ledger}

\begin{ledger}[Axiom passage]
\label{led:v11v64-data}
A complete continuum construction must establish:
\begin{enumerate}
\item a reflection map compatible with the lattice refinement and
gauge quotient;
\item finite-cutoff reflection positivity for the full regulated
boson--fermion observable algebra;
\item cylinder convergence plus uniform integrability for unbounded
fields;
\item reflection-compatible low/high marginalization;
\item a uniform connected-correlation envelope or a reconstructed
transfer gap;
\item removal or controlled treatment of global zero and phase modes;
\item compatibility with topological-sector mixing.
\end{enumerate}
\end{ledger}

\begin{firewall}[No positivity or mass-gap claim]
\label{fire:v11v64-boundary}
The ultralocal positive metric reference of V39 is reflection positive,
but the complete coupled renewal action has not been given a
cross-plane positive decomposition.  The chiral determinant requires
a fermionic reflection proof, and no uniform connected-correlation
decay or transfer gap has been acquired for the SM--Einstein measure.
V64 proves passage criteria, not those physical inputs.
\end{firewall}

\begin{decision}[Mandatory V65 audit]
\label{dec:v11v64-next}
The next version is the compulsory companion audit.  It will check
whether YM has advanced its nonzero-momentum or mass-gap row and
whether NS has acquired a uniform decay estimate.  V66 will then
attack the remaining reference-to-target transfer or clustering input
indicated by that audit.
\end{decision}

\begin{assessment}[Progress after V64]
\label{ass:v11v64-status}
Reflection positivity now passes rigorously through cylinder limits and
reflection-compatible marginals.  Gaussian clustering is reduced to
covariance decay, and general exponential clustering to a uniform
transfer gap.  The constant-field counterexample proves why neither
positivity nor tightness can replace the missing decay estimate.
\end{assessment}

\subsection{Parent record}

This V64 note was formed by appending the preceding material to the
validated V63 source, whose record was
\[
\begin{array}{c|c}
\text{lines}&19895\\
\text{bytes}&622097\\
\text{SHA--256}&
\texttt{EE80728A858D3964C9694E6194D50044322D7F08FC98AD180AAEACB49FAB6E3A}.
\end{array}                                                     \tag{64.20}
\]

% =============================================================================
% END APPENDED ELEVENTH-SERIES V64 MATERIAL
% =============================================================================

% =============================================================================
% BEGIN APPENDED ELEVENTH-SERIES V65 MATERIAL
% =============================================================================

\section{Companion audit at V65: the first nonzero momentum card and
sign-changing response}
\label{sec:v11v65}

The mandatory audit finds a new Yang--Mills V36 and an unchanged
Navier--Stokes V26.

\subsection{Audited records}

\begin{definition}[V65 companion record]
\label{def:v11v65-record}
The inspected files are
\[
\begin{array}{c|r|l}
\text{file}&\text{bytes}&\text{SHA--256}\\ \hline
\texttt{Renewal\_Yang\_Mills\_Mass\_Gap\_Research\_Notes\_V36.tex}
 &511923&
\texttt{2B737CB7D11C2FED1B9719824ED13E6591F4940611AEEB2E9B1ACE1636F33853}\\
\texttt{Renewal\_NS\_Stability\_Research\_Notes\_V26.tex}
 &278283&
\texttt{82C887F8C2C69E4F28FAD7C3DC8DE5B9099CB5A4BF431EBA1FFF3E91935978AD}.
\end{array}                                                     \tag{65.1}
\]
The YM record has advanced from V35; the NS record is unchanged.
\end{definition}

\subsection{New mathematical content}

The YM note evaluates one self-inverse nonzero internal orbit with
independent internal momentum \(\kappa\) and external differentiation
parameter \(q\).  Its exact jet ring keeps
\[
                    e^{i\alpha\cdot\kappa}
                    \quad\text{separate from}\quad
                    \frac{d^j}{dq^j}e^{iq\alpha\cdot e_\mu}.  \tag{65.2}
\]
It passes exact Hermitian, plus/minus, stationary-score, Haar, and
first-jet parity gates and reproduces the earlier toron card through
the generalized code path.

\begin{proposition}[Independent momentum labels are compulsory]
\label{prop:v11v65-labels}
Let a response word contain phases
\[
                 \prod_{\ell=1}^N
                 e^{i\eta_\ell\kappa\cdot a_\ell}
                 e^{i\sigma_\ell q\cdot a_\ell}.              \tag{65.3}
\]
External differentiation at fixed internal momentum gives
\[
 \partial_q^j\big|_{q=0}(65.3)
 =\left(\prod_\ell e^{i\eta_\ell\kappa\cdot a_\ell}\right)
  \left[i\sum_\ell\sigma_\ell a_\ell\right]^j.               \tag{65.4}
\]
Replacing the independent labels \(\eta_\ell,\sigma_\ell\) by one
phase label changes (65.4) unless their declared leg assignments
coincide.
\end{proposition}

\begin{proof}
The internal factor in (65.3) is constant with respect to \(q\).
Differentiate only the external exponential product.  A merged label
differentiates phase contributions assigned solely to \(\kappa\) or
omits contributions assigned solely to \(q\), so it agrees only in the
exceptional case of identical assignments.
\end{proof}

This rule transfers directly to the SM gauge and matter vertices:
loop momentum, external probe momentum, and background translation
jets must remain independent through every noncommutative word.

\subsection{The sign-change consequence}

For the audited transverse channel, the YM endpoint difference obeys
\[
                        \Delta^{[0]}(0)<0,\qquad
                        \Delta^{[0]}(\pi e_2)>0.               \tag{65.5}
\]
Continuity yields at least one zero between them.

\begin{corollary}[No one-cell sign certificate]
\label{cor:v11v65-no-sign}
Neither endpoint value in (65.5), and neither finite set of point
values without an enclosure between them, determines the sign or value
of the Brillouin-zone integral.
\end{corollary}

\begin{proof}
A continuous periodic function can be modified by a smooth bump
supported away from any prescribed finite sample set.  The
modification leaves all sampled values fixed while changing its
integral by an arbitrary signed amount.  Thus sample values alone do
not certify the mean.
\end{proof}

The remaining gauge response requires exact orbit pairing followed by
a certified global quadrature or an analytic cancellation theorem.

\subsection{Reflection terminology}

The YM V36 “reflection first-jet” gate states that certain finite
momentum derivatives vanish or transform with the expected sign under
coordinate reflection.  It is an algebraic covariance check on one
cell response.

\begin{warning}[Parity is not reflection positivity]
\label{warn:v11v65-reflection}
A coordinate-reflection parity identity does not imply
\[
                         \int(\Theta F)F\,d\mu\ge0             \tag{65.6}
\]
for all positive-time observables.  The latter is the
Osterwalder--Schrader condition of V64 and depends on the complete
measure and its cross-plane interactions.  The YM gate is retained as
a compiler regression, not imported as continuum reflection
positivity.
\end{warning}

\subsection{Transfer to the determinant and infrared ledgers}

The audited cell uses inverses of positive finite matrices at
\(\kappa\) and \(\kappa+q e_1\).  Its smooth momentum dependence follows
from a positive finite-cutoff Hessian floor.  This supports the
analytic carrier and high-sector framework of V56 and V62 at that
endpoint.  It does not establish a physical high-sector gap for the
full SM--Einstein target.

The exact nonzero cell is valuable as a regression point for a global
quadrature, while its sign change shows why the quadrature error must
be rigorous and smaller than the claimed final coefficient margin.

\begin{decision}[V66 acquisition]
\label{dec:v11v65-next}
V66 will prove an analytic-strip trapezoidal quadrature theorem for
periodic matrix-response functions.  It will derive the strip width
from a complex-momentum inverse floor, give an exponential grid error,
and add interval evaluation error.  This supplies a route from exact
toron and nonzero regression points to a certified Brillouin mean,
conditional on a uniform complex strip for the chosen regulator.
\end{decision}

\begin{firewall}[Audit boundary]
\label{fire:v11v65-boundary}
The YM V36 note computes one nonzero orbit, not the complete momentum
mean or Ward telescope.  Its positive endpoint floor and parity gates
do not prove the target chiral SM determinant, Einstein constraint
gap, reflection positivity, clustering, or a mass gap.  The unchanged
NS note contributes no new decay or quadrature input.
\end{firewall}

\begin{assessment}[Progress after V65]
\label{ass:v11v65-status}
The required audit is complete.  It contributes an exact nonzero
regression point, enforces independent internal and external phases,
and proves that the response integrand changes sign.  The next
well-posed acquisition is therefore global certified quadrature rather
than extrapolation from local samples.
\end{assessment}

\subsection{Parent record}

This V65 note was formed by appending the preceding material to the
validated V64 source, whose record was
\[
\begin{array}{c|c}
\text{lines}&20237\\
\text{bytes}&634060\\
\text{SHA--256}&
\texttt{8FAB919F828BD037AB4575D1AEF5A63702877260A32C12443A42BD3CF7C7C693}.
\end{array}                                                     \tag{65.7}
\]

% =============================================================================
% END APPENDED ELEVENTH-SERIES V65 MATERIAL
% =============================================================================
% =============================================================================
% BEGIN APPENDED ELEVENTH-SERIES V66 MATERIAL
% =============================================================================

\section{Certified analytic quadrature for Brillouin-zone response}
\label{sec:v11v66}

Exact toron and nonzero-orbit values are regression data, not a
momentum integral.  A periodic analytic strip supplies a rigorous
global certificate.  The tensor trapezoidal rule then converges
exponentially, with an error determined by the strip width and a
uniform complex bound.

\subsection{Fourier decay in a periodic polystrip}

Let \(f\) be \(2\pi\)-periodic in each of \(d\) variables and analytic
on
\[
 {\cal S}_{\boldsymbol a}
 =\{z\in\mathbb C^d:|\operatorname{Im}z_j|<a_j,\ 1\le j\le d\}.
                                                                  \tag{66.1}
\]
Assume \(|f(z)|\le M\) there.  Its normalized Fourier coefficients are
\[
 \widehat f(n)=\frac1{(2\pi)^d}
   \int_{[0,2\pi]^d}f(k)e^{-in\cdot k}\,dk.                   \tag{66.2}
\]

\begin{lemma}[Polystrip Fourier bound]
\label{lem:v11v66-fourier}
For every \(n\in\mathbb Z^d\),
\[
                         |\widehat f(n)|
                    \le M\exp\left\{-\sum_{j=1}^da_j|n_j|\right\}.
                                                                  \tag{66.3}
\]
\end{lemma}

\begin{proof}
For each coordinate with \(n_j\ne0\), shift its contour to
\(\operatorname{Im}z_j=-a_j\operatorname{sgn}(n_j)\), first with
\(a_j\) replaced by \(a_j-\varepsilon\).  Periodicity cancels the
vertical sides, and analyticity permits the shift.  The Fourier
exponential gains
\(e^{-(a_j-\varepsilon)|n_j|}\).  Multiply the coordinate factors,
bound \(f\) by \(M\), and let \(\varepsilon\downarrow0\).
\end{proof}

\subsection{Tensor trapezoidal error}

For integers \(N_j\ge1\), define
\[
 Q_{\boldsymbol N}(f)
 =\frac1{\prod_jN_j}
   \sum_{r_1=0}^{N_1-1}\cdots\sum_{r_d=0}^{N_d-1}
   f\left(\frac{2\pi r_1}{N_1},\ldots,
          \frac{2\pi r_d}{N_d}\right).                        \tag{66.4}
\]
Let
\[
                         q_j=e^{-a_jN_j}.                     \tag{66.5}
\]

\begin{theorem}[Certified periodic trapezoidal rule]
\label{thm:v11v66-trapezoid}
The normalized Brillouin mean
\[
                         I(f)=\frac1{(2\pi)^d}
                                  \int_{[0,2\pi]^d}f(k)\,dk   \tag{66.6}
\]
obeys
\[
 \boxed{
 |Q_{\boldsymbol N}(f)-I(f)|
 \le M\left[
     \prod_{j=1}^d\frac{1+q_j}{1-q_j}-1
        \right].}                                             \tag{66.7}
\]
For an isotropic strip and grid, \(a_j=a\), \(N_j=N\),
\[
 |Q_N(f)-I(f)|
 \le M\left[\left(\frac{1+e^{-aN}}{1-e^{-aN}}\right)^d-1\right].
                                                                  \tag{66.8}
\]
\end{theorem}

\begin{proof}
Discrete character orthogonality gives the exact aliasing identity
\[
                         Q_{\boldsymbol N}(f)
                    =\sum_{m\in\mathbb Z^d}
                         \widehat f(N_1m_1,\ldots,N_dm_d).     \tag{66.9}
\]
The integral is \(\widehat f(0)\).  Apply
Lemma~\ref{lem:v11v66-fourier} and sum:
\[
 \sum_{m\in\mathbb Z^d}
      \prod_jq_j^{|m_j|}
 =\prod_j\left(1+2\sum_{\ell\ge1}q_j^\ell\right)
 =\prod_j\frac{1+q_j}{1-q_j}.                                \tag{66.10}
\]
Removing \(m=0\) proves (66.7).
\end{proof}

\subsection{A complex inverse strip from a real gap}

Let \(H(z)\) be an analytic \(n\times n\) matrix on a polystrip and
suppose that on the real torus
\[
                         s_{\min}(H(k))\ge\gamma>0.            \tag{66.11}
\]
Assume the complex displacement bound
\[
 \lVert H(k+iy)-H(k)\rVert
                  \le\sum_{j=1}^dL_j|y_j|.                   \tag{66.12}
\]

\begin{theorem}[Complex-momentum inverse floor]
\label{thm:v11v66-strip}
If the strip widths satisfy
\[
                              \sum_jL_ja_j\le\theta\gamma
                              \quad\text{for some }0<\theta<1,\tag{66.13}
\]
then throughout \({\cal S}_{\boldsymbol a}\),
\[
                 H(z)\text{ is invertible},\qquad
                 \lVert H(z)^{-1}\rVert
                      \le\frac1{(1-\theta)\gamma}.            \tag{66.14}
\]
\end{theorem}

\begin{proof}
For \(z=k+iy\), write
\[
 H(z)=H(k)\left[I+H(k)^{-1}(H(z)-H(k))\right].
\]
The norm of the bracketed perturbation is at most \(\theta<1\) by
(66.11)--(66.13).  Its Neumann inverse has norm at most
\((1-\theta)^{-1}\), which proves (66.14).
\end{proof}

For Laurent-polynomial \(H\), a bound of the form (66.12) follows by
integrating its complex derivatives along \(k+ity\).  The derivative
bound must be evaluated on the same strip; a real-axis derivative
bound alone is insufficient.

\subsection{Bounding a matrix response}

Suppose a response is a sum of \(R\) trace words
\[
                 f(z)=\sum_{\ell=1}^R c_\ell
                       \operatorname{Tr}
                       \prod_{r=1}^{m_\ell}A_{\ell r}(z),      \tag{66.15}
\]
where factors are analytic vertices or inverse matrices controlled by
Theorem~\ref{thm:v11v66-strip}.

\begin{proposition}[Complex response envelope]
\label{prop:v11v66-envelope}
If
\[
                         \lVert A_{\ell r}(z)\rVert
                                  \le M_{\ell r}               \tag{66.16}
\]
on the strip and the trace matrices have size \(n_\ell\), then
\[
                         |f(z)|
 \le\sum_{\ell=1}^R|c_\ell|n_\ell
                          \prod_{r=1}^{m_\ell}M_{\ell r}
 =:M_f.                                                       \tag{66.17}
\]
\end{proposition}

\begin{proof}
Use
\[
 |\operatorname{Tr}(A_1\cdots A_m)|
 \le n\prod_r\lVert A_r\rVert
\]
term by term.
\end{proof}

Equations (66.13), (66.17), and (66.7) are a complete deterministic
quadrature certificate.

\subsection{Interval evaluation}

Suppose each grid value is enclosed by a complex interval
\({\cal I}_r\) with center \(v_r\) and radius at most
\(\eta_r\).  Let
\[
                         \eta_{\rm eval}
             =\frac1{\prod_jN_j}\sum_r\eta_r.                 \tag{66.18}
\]

\begin{corollary}[Total certified error]
\label{cor:v11v66-interval}
Let \(Q_{\boldsymbol N}^{\rm comp}\) be the average of the interval
centers.  Then
\[
 |I(f)-Q_{\boldsymbol N}^{\rm comp}|
 \le\eta_{\rm eval}
 +M_f\left[\prod_j\frac{1+e^{-a_jN_j}}
                            {1-e^{-a_jN_j}}-1\right].          \tag{66.19}
\]
If the resulting real interval excludes zero, it certifies the sign of
the real Brillouin mean.
\end{corollary}

\begin{proof}
The interval radii bound the difference between the exact and computed
grid averages.  Add the analytic aliasing error (66.7).
\end{proof}

Exact rational or algebraic evaluations at special orbits, such as YM
V34 and V36, should be included as regression checks on the interval
compiler.  They do not reduce (66.19) unless used in an interpolation
or symmetry theorem with its own proved error.

\subsection{Exact Ward telescopes before quadrature}

Let \(G\) be periodic and let a grid shift
\[
                              h_j=\frac{2\pi m_j}{N_j}         \tag{66.20}
\]
map the tensor grid to itself.

\begin{lemma}[Discrete periodic telescope]
\label{lem:v11v66-telescope}
For
\[
                              (\Delta_hG)(k)=G(k+h)-G(k),
\]
one has
\[
                              I(\Delta_hG)=0,\qquad
                              Q_{\boldsymbol N}(\Delta_hG)=0. \tag{66.21}
\]
\end{lemma}

\begin{proof}
Translation invariance of periodic Lebesgue measure gives the integral
identity.  The shift permutes the grid points, so the two discrete sums
are identical.
\end{proof}

Thus exact Ward differences and orbit-pair cancellations should be
removed algebraically before applying (66.19).  Numerical quadrature
should certify only the irreducible fundamental-domain integrand.

\subsection{A summable quadrature schedule}

Suppose at cutoff \(i\) the response dimension \(d\) is fixed, the
strip is isotropic with width \(a_i\), and
\(M_i\lesssim(i+2)^m\).  For sufficiently large \(a_iN_i\), the
bracket in (66.8) is bounded by \(C_de^{-a_iN_i}\).

\begin{corollary}[Adjacent summability]
\label{cor:v11v66-schedule}
If for some \(\delta>0\),
\[
                         a_iN_i\ge(m+1+\delta)\log(i+2)        \tag{66.22}
\]
and the interval errors are summable, then the quadrature errors are
summable.
\end{corollary}

\begin{proof}
The analytic error is bounded by
\[
 C(i+2)^m e^{-a_iN_i}
 \le C(i+2)^{-1-\delta}.
\]
Sum this and the interval errors.
\end{proof}

If the complex gap shrinks, Theorem~\ref{thm:v11v66-strip} makes
\(a_i\) shrink as well, and the grid must grow accordingly.  A real
positive floor without the complex derivative bound does not establish
(66.22).

\subsection{SM--Einstein momentum ledger}

\begin{ledger}[Global response certificate]
\label{led:v11v66-data}
For each gauge, ghost, Higgs, fermion, and metric momentum response,
the continuum programme must supply:
\begin{enumerate}
\item independent internal, external, and background phase labels;
\item algebraic orbit pairing and Ward telescopes;
\item a real high-sector inverse floor;
\item a complex displacement bound and certified strip widths;
\item a uniform trace-word envelope \(M_i\);
\item interval enclosures of the irreducible grid values;
\item a quadrature schedule satisfying (66.22);
\item adjacent compatibility of counterterms and complex phases.
\end{enumerate}
\end{ledger}

\begin{firewall}[No Brillouin mean computed]
\label{fire:v11v66-boundary}
V66 proves the quadrature theorem but does not establish a complex
strip or evaluate the global grid for the full SM--Einstein response.
Zeros of a chiral determinant, massless infrared modes, or a shrinking
constraint gap can collapse the strip.  The YM V36 sign change remains
a regression and obstruction, not a completed mean.
\end{firewall}

\begin{decision}[Next acquisition]
\label{dec:v11v66-next}
V67 will formalize symmetry-orbit reduction before quadrature.  It will
prove the fundamental-domain weighting formula for a finite momentum
symmetry group, include stabilizer weights, and state the exact
compatibility conditions for complex conjugation, \(k\mapsto-k\), and
Ward-paired cells.
\end{decision}

\begin{assessment}[Progress after V66]
\label{ass:v11v66-status}
A periodic matrix response can now be promoted from exact sample values
to a certified Brillouin mean once a complex inverse strip is proved.
The error is explicit, interval arithmetic is included, and exact Ward
telescopes are removed before numerical work.  Acquisition of the
physical strip and execution of the full SM--Einstein grid remain open.
\end{assessment}

\subsection{Parent record}

This V66 note was formed by appending the preceding material to the
validated V65 source, whose record was
\[
\begin{array}{c|c}
\text{lines}&20416\\
\text{bytes}&641130\\
\text{SHA--256}&
\texttt{BF2C1FA024E3F15F4F691FE3923E052B34FE312D86DAF7051C13DDB79091691E}.
\end{array}                                                     \tag{66.23}
\]

% =============================================================================
% END APPENDED ELEVENTH-SERIES V66 MATERIAL
% =============================================================================
% =============================================================================
% BEGIN APPENDED ELEVENTH-SERIES V67 MATERIAL
% =============================================================================

\section{Momentum-orbit reduction with stabilizer weights}
\label{sec:v11v67}

Symmetry should remove exact repetitions before interval quadrature.
On a finite grid, orbit sizes are not constant: axes, corners, torons,
and self-inverse momenta have nontrivial stabilizers.  This section
gives the exact weights and the covariant tensor version.

\subsection{Invariant scalar sums}

Let a finite group \({\mathfrak G}\) act on a finite set \(X\).
For \(x\in X\), write
\[
 {\mathfrak G}_x=\{g:g x=x\},\qquad
 |{\mathfrak G}x|=\frac{|{\mathfrak G}|}{|{\mathfrak G}_x|}. \tag{67.1}
\]
Choose one representative from each orbit, forming \(R\subset X\).

\begin{theorem}[Exact orbit-weighted average]
\label{thm:v11v67-orbit}
If \(f(gx)=f(x)\), then
\[
 \boxed{
 \frac1{|X|}\sum_{x\in X}f(x)
 =\sum_{r\in R}
    \frac{|{\mathfrak G}|}
         {|X|\,|{\mathfrak G}_r|}\,f(r).}                    \tag{67.2}
\]
The weights are positive and sum to one.
\end{theorem}

\begin{proof}
Partition \(X\) into disjoint orbits.  Invariance makes \(f\) constant
on each orbit, and orbit--stabilizer gives its cardinality (67.1).
Applying the formula to \(f=1\) proves that the weights sum to one.
\end{proof}

\begin{corollary}[Weighted interval error]
\label{cor:v11v67-interval}
If the value at representative \(r\) is enclosed with radius
\(\eta_r\), the orbit-reduced grid average has evaluation radius
\[
                  \eta_{\rm orb}
     =\sum_{r\in R}
       \frac{|{\mathfrak G}|}
            {|X|\,|{\mathfrak G}_r|}\eta_r.                  \tag{67.3}
\]
Together with the analytic aliasing error of V66 this gives a certified
mean.
\end{corollary}

\begin{proof}
Use (67.2) and the triangle inequality for a positive weighted sum.
\end{proof}

Using a generic orbit size at a fixed point overcounts that point.
The stabilizer denominator in (67.2) is therefore compulsory.

\subsection{Covariant tensor responses}

Let \(\rho:{\mathfrak G}\to{\rm GL}(V)\) be a finite-dimensional
representation and suppose
\[
                              F(gx)=\rho(g)F(x).               \tag{67.4}
\]

\begin{lemma}[Stabilizer compatibility]
\label{lem:v11v67-stabilizer}
For every \(h\in{\mathfrak G}_x\),
\[
                              \rho(h)F(x)=F(x).                \tag{67.5}
\]
Thus \(F(x)\) lies in the stabilizer-fixed subspace
\[
                  V^{{\mathfrak G}_x}
      =\{v:\rho(h)v=v\text{ for every }h\in{\mathfrak G}_x\}. \tag{67.6}
\]
\end{lemma}

\begin{proof}
Set \(g=h\) in (67.4) and use \(hx=x\).
\end{proof}

\begin{theorem}[Covariant orbit reconstruction]
\label{thm:v11v67-covariant}
For each representative \(r\), choose coset representatives
\({\mathfrak G}/{\mathfrak G}_r\).  Then
\[
                    \sum_{x\in{\mathfrak G}r}F(x)
       =\sum_{[g]\in{\mathfrak G}/{\mathfrak G}_r}\rho(g)F(r).
                                                                  \tag{67.7}
\]
The right-hand side is independent of the coset representatives.
\end{theorem}

\begin{proof}
The orbit map sends the cosets bijectively to \({\mathfrak G}r\).
If \(g\) is replaced by \(gh\) with \(h\in{\mathfrak G}_r\), then
\[
                         \rho(gh)F(r)=\rho(g)F(r)
\]
by Lemma~\ref{lem:v11v67-stabilizer}.  Thus the sum is well defined
and equals the orbit sum by covariance.
\end{proof}

For a two-index polarization tensor,
\(\rho(g)\) acts on both indices.  Scalar trace or projected channel
formulas may use (67.2) only after the tensor covariance and the
external-data stabilizer have been identified.

\subsection{The external stabilizer group}

Let \(q\) and a polarization label \(e\) be fixed external data.
Only
\[
 {\mathfrak G}_{q,e}
 =\{g\in{\mathfrak G}:gq=q,\ ge=e
       \text{ with the declared tensor action}\}             \tag{67.8}
\]
acts within that scalar response channel.

\begin{warning}[Do not average different external channels]
\label{warn:v11v67-external}
A symmetry carrying \((q,e)\) to \((gq,ge)\ne(q,e)\) relates two
components of the full tensor response.  It does not make either
component invariant.  Such a symmetry can be used through
Theorem~\ref{thm:v11v67-covariant}, but not by applying the scalar
formula (67.2) to one fixed channel.
\end{warning}

This is the group-theoretic form of the longitudinal/transverse
distinction in the audited YM toron cards.

\subsection{Momentum reversal and self-inverse points}

On the periodic grid, let \(\iota k=-k\) modulo \(2\pi\).

\begin{proposition}[Conjugate pairing]
\label{prop:v11v67-conjugate}
Suppose
\[
                              f(-k)=\overline{f(k)}.           \tag{67.9}
\]
Every two-point orbit \(\{k,-k\}\) contributes
\(2\operatorname{Re}f(k)\).  At a self-inverse momentum
\(k=-k\), equation (67.9) forces \(f(k)\in\mathbb R\), and the point is
counted once.
\end{proposition}

\begin{proof}
For a two-point orbit, add (67.9) to its conjugate.  At a fixed point,
\(f(k)=\overline{f(k)}\).  Orbit--stabilizer assigns orbit size one.
\end{proof}

On an even tensor grid, self-inverse coordinates are \(0\) or \(\pi\).
These include the YM toron and the V36 orbit.  Doubling their values
would be an exact factor-of-two error.

\begin{corollary}[Odd cancellation]
\label{cor:v11v67-odd}
If \(f(-k)=-f(k)\), the complete grid sum vanishes.  Every
self-inverse value is zero.
\end{corollary}

\begin{proof}
Two-point orbits cancel, while at a fixed point
\(f(k)=-f(k)\).
\end{proof}

\subsection{Compatibility with the tensor grid}

Signed coordinate permutations preserve an isotropic \(N^d\) grid.
For an anisotropic grid, a permutation exchanging axes \(j,\ell\)
preserves it only when \(N_j=N_\ell\).  More generally, a proposed
symmetry reduction is exact only when it permutes the actual grid set.

\begin{proposition}[Orbit reduction preserves quadrature]
\label{prop:v11v67-quadrature}
If \({\mathfrak G}\) preserves the tensor grid and \(f\) is invariant,
the orbit-reduced value (67.2) is exactly
\(Q_{\boldsymbol N}(f)\).  Therefore its analytic quadrature error is
still exactly the V66 bound; orbit reduction adds no approximation.
\end{proposition}

\begin{proof}
Equation (67.2) is merely a regrouping of the same finite sum.
\end{proof}

If the grid is not preserved, interpolating a transformed point back
to the grid adds a new numerical error and is not an exact symmetry
reduction.

\subsection{Ward shifts and group orbits}

Let \(\Delta_hG(k)=G(k+h)-G(k)\), with \(h\) a grid shift as in V66.
The full-grid telescope vanishes exactly.  An orbit representative set
need not be stable under \(k\mapsto k+h\).

\begin{lemma}[Safe order of operations]
\label{lem:v11v67-order}
A Ward difference may be removed before orbit reduction whenever its
full-grid sum is zero.  If it is retained after orbit reduction, the
representative weights must be applied to both \(k\) and \(k+h\)
through their actual orbits; pairing representatives without this
check need not telescope.
\end{lemma}

\begin{proof}
The first statement is the exact identity (66.21).  The second follows
because the shift can carry a point to an orbit with a different
stabilizer and hence a different weight in (67.2).  Only the fully
weighted sums reproduce the permuted full grid.
\end{proof}

\subsection{Fundamental domains in the continuum}

For the continuous torus with invariant Haar measure, let
\({\cal F}\) be a measurable fundamental region whose distinct
translates overlap only on a null boundary.  If \(f\) is invariant,
\[
                         \int_{\mathbb T^d}f(k)\,dk
                    =|{\mathfrak G}|\int_{\cal F}f(k)\,dk.    \tag{67.10}
\]
Positive-dimensional fixed sets require separate strata if they have
nonzero measure; for finite linear symmetry groups on a torus, proper
fixed subspaces have Lebesgue measure zero.

\begin{proof}
The translates of \({\cal F}\) partition the torus modulo a null set.
Use invariance and add their equal integrals.
\end{proof}

The finite-grid formula (67.2) remains preferable for computation
because it handles all boundary and stabilizer points exactly.

\subsection{Compiler and continuum ledgers}

\begin{ledger}[Symmetry certificate]
\label{led:v11v67-data}
Before global response quadrature, each channel must record:
\begin{enumerate}
\item the finite group acting on internal momentum;
\item its simultaneous action on external momentum, polarization,
colour, spin, chirality, and background indices;
\item the subgroup fixing the declared external channel;
\item orbit representatives and exact stabilizer orders;
\item conjugation or sign rules under \(k\mapsto-k\);
\item self-inverse momenta and their real or vanishing gates;
\item Ward shifts removed on the full grid;
\item regression against unreduced small grids.
\end{enumerate}
\end{ledger}

\begin{firewall}[Symmetry does not determine the mean]
\label{fire:v11v67-boundary}
Orbit reduction shortens an exact grid sum and enforces cancellations.
It does not supply the complex analytic strip, interval enclosures, or
the values on the irreducible domain.  The sign-changing YM response
still requires a global certified evaluation.  Symmetry of a
finite-cell compiler also does not prove continuum reflection
positivity or clustering.
\end{firewall}

\begin{decision}[Next acquisition]
\label{dec:v11v67-next}
V68 will formulate the full summed Ward identity as a discrete
cochain telescope.  It will separate exact incidence cancellation from
anomaly terms, boundary terms, and regulator defects, and give the
summability condition under which a vanishing finite-cutoff Ward defect
passes to the continuum cylinder functional.
\end{decision}

\begin{assessment}[Progress after V67]
\label{ass:v11v67-status}
Momentum symmetry reduction is now exact and stabilizer aware.
Self-inverse points are counted once, conjugate pairs become real
contributions, tensor channels are transformed rather than conflated,
and Ward telescopes are removed in the safe full-grid order.  The
irreducible response values and complex strip remain to be acquired.
\end{assessment}

\subsection{Parent record}

This V67 note was formed by appending the preceding material to the
validated V66 source, whose record was
\[
\begin{array}{c|c}
\text{lines}&20759\\
\text{bytes}&652337\\
\text{SHA--256}&
\texttt{AB71F1D7039A4AA6C90D8C4040C0199CFA07ABAEE7D78D6B3302F7444274C4F5}.
\end{array}                                                     \tag{67.11}
\]

% =============================================================================
% END APPENDED ELEVENTH-SERIES V67 MATERIAL
% =============================================================================
% =============================================================================
% BEGIN APPENDED ELEVENTH-SERIES V68 MATERIAL
% =============================================================================

\section{Discrete Ward telescopes, anomaly rows, and continuum passage}
\label{sec:v11v68}

A cell response need not satisfy a Ward identity separately.  The
identity belongs to the complete finite-cutoff measure and, in momentum
space, often appears only after all shifted cells are summed.  This
section states the exact finite-dimensional change-of-variables
identity and separates anomaly, regulator, and boundary defects.

\subsection{The cochain telescope}

Let \(K\) be a finite oriented cell complex and
\[
                         d_p:C^p(K)\to C^{p+1}(K)             \tag{68.1}
\]
its coboundary.  With the chosen cell inner products, let
\(d_p^*\) be the adjoint.

\begin{lemma}[Closed-complex divergence sum]
\label{lem:v11v68-divergence}
If \(K\) has no boundary and \({\bf1}\in C^0(K)\) is constant, then for
every one-cochain \(J\),
\[
                         \sum_{v\in K^0}(d_0^*J)(v)
          =\langle{\bf1},d_0^*J\rangle
          =\langle d_0{\bf1},J\rangle=0.                     \tag{68.2}
\]
On a complex with boundary, the same incidence cancellation leaves
only the oriented boundary flux.
\end{lemma}

\begin{proof}
The coboundary of a constant zero-cochain vanishes.  Adjointness gives
(68.2).  If boundary cells are omitted from the adjoint pairing, the
uncancelled incidences are precisely the boundary faces, which is the
discrete Stokes formula.
\end{proof}

For a periodic momentum grid, Fourier transformation turns (68.2)
into the vanishing sum of a periodic finite difference.  This is the
cochain origin of the grid telescope (66.21).

\subsection{Finite-dimensional Ward identity}

Let \(X\) be an oriented finite-dimensional field manifold with volume
form \(dV\), and let \(R_\epsilon\) be the infinitesimal field vector
generated by a gauge parameter \(\epsilon\).  Let the regulated weight
be a nowhere-zero complex density \(w\), normalized by
\[
                              Z=\int_Xw\,dV\ne0.               \tag{68.3}
\]
Assume integration by parts is valid, either because \(X\) is compact
or because the boundary term vanishes.

\begin{theorem}[Ward identity with density variation]
\label{thm:v11v68-ward}
For every smooth integrable observable \(O\),
\[
 \boxed{
 {\mathbb E}_w[R_\epsilon O]
 +{\mathbb E}_w\!\left[
 O\left(\operatorname{div}R_\epsilon+
              R_\epsilon\log w\right)\right]=0.}             \tag{68.4}
\]
Here expectations are normalized by \(Z\), and
\(R_\epsilon\log w\) means \(w^{-1}R_\epsilon w\), so no global
logarithm branch is required.
\end{theorem}

\begin{proof}
The divergence theorem gives
\[
 0=\int_X\operatorname{div}(OwR_\epsilon)\,dV
  =\int_Xw\left[
 R_\epsilon O+
 O\operatorname{div}R_\epsilon+
 O\,w^{-1}R_\epsilon w\right]dV.
\]
Divide by \(Z\).
\end{proof}

If the bosonic action, volume form, and regulated fermion density are
all invariant, the second expectation vanishes and
\[
                              {\mathbb E}[R_\epsilon O]=0.     \tag{68.5}
\]

\subsection{Anomaly and regulator decomposition}

Write the logarithmic density variation as
\[
 \operatorname{div}R_\epsilon+R_\epsilon\log w
       ={\cal A}_\epsilon+{\cal E}_\epsilon
                             +d_0^*J_\epsilon.                \tag{68.6}
\]
Here \({\cal A}\) is the local anomaly density, \({\cal E}\) is a
regulator or gauge-fixing defect, and the last term is a cochain
divergence.

\begin{corollary}[Summed finite-cutoff Ward equation]
\label{cor:v11v68-summed}
On a periodic complex, the complete summed identity is
\[
 {\mathbb E}_w[R_\epsilon O]
 +{\mathbb E}_w[O{\cal A}_\epsilon]
 +{\mathbb E}_w[O{\cal E}_\epsilon]
 +{\mathbb E}_w[O\,d_0^*J_\epsilon]=0.                       \tag{68.7}
\]
For \(O=1\), the unweighted divergence sum vanishes by
Lemma~\ref{lem:v11v68-divergence}.  For local \(O\), discrete
integration by parts moves \(d_0\) onto \(O\) and leaves contact terms
at its support.
\end{corollary}

\begin{proof}
Insert (68.6) into (68.4).  Apply adjointness to the divergence term;
when \(O=1\), use (68.2).
\end{proof}

The exact Standard Model anomaly arithmetic proved earlier in this
cycle cancels the perturbative local gauge-anomaly coefficients after
all chiral species are grouped.  It does not remove
\({\cal E}_\epsilon\), a boundary flux, or global determinant-line
holonomy.

\begin{warning}[Cellwise cancellation is stronger]
\label{warn:v11v68-cell}
Equation (68.7) is a statement about the complete field and momentum
sum.  If
\[
                              {\cal W}(q)=\sum_k{\cal W}(q;k),
\]
then \(q^\mu{\cal W}_{\mu\nu}(q)=0\) requires only cancellation of the
sum over \(k\).  It does not require each cell contraction to vanish.
The sign-changing YM V36 card is consistent with such a summed
identity.
\end{warning}

\subsection{Momentum-shift form}

Suppose a regulator defect is a periodic momentum difference
\[
                   {\cal E}(k)=G(k+h)-G(k)+R(k),              \tag{68.8}
\]
where \(h\) permutes the finite grid.  Then
\[
                         \sum_k{\cal E}(k)=\sum_kR(k).         \tag{68.9}
\]

\begin{proposition}[Exact Ward telescope with remainder]
\label{prop:v11v68-momentum}
The shifted part of (68.8) cancels exactly before quadrature or orbit
reduction.  If interval arithmetic gives
\[
                              \sum_k|R(k)|\le\varepsilon,      \tag{68.10}
\]
then the complete discrete Ward defect is bounded by
\(\varepsilon\), apart from separately displayed anomaly and boundary
terms.
\end{proposition}

\begin{proof}
The shift \(k\mapsto k+h\) is a permutation, so the two sums of \(G\)
are equal.  Apply the triangle inequality to the remainder.
\end{proof}

A numerical near-cancellation of the two large shifted sums is weaker
than this algebraic removal and pays unnecessary interval width.

\subsection{Boundary removal}

On a finite region \(\Lambda_i\), suppose the Ward identity leaves a
boundary current \(J_i\).  Let a cylinder observable \(O\) be supported
in a fixed set \(K\), and suppose
\[
       |{\mathbb E}[O\,J_i(e)]|
        \le C_O h(d(e,K))                                    \tag{68.11}
\]
for boundary edges \(e\), with \(h(r)\to0\).

\begin{lemma}[Boundary-flux criterion]
\label{lem:v11v68-boundary}
If
\[
                         |\partial\Lambda_i|
               \sup_{e\in\partial\Lambda_i}h(d(e,K))
                              \longrightarrow0,               \tag{68.12}
\]
then the boundary term in the Ward identity for \(O\) vanishes.
\end{lemma}

\begin{proof}
Sum (68.11) over the oriented boundary edges and use the triangle
inequality.
\end{proof}

Exponential clustering makes (68.12) immediate for polynomial boundary
growth and linearly receding boundaries.  Polynomial clustering
requires its decay exponent to beat the boundary growth.

\subsection{Passage to the continuum functional}

Let \({\cal L}_i(O)\) denote the complete left side of the desired Ward
identity at cutoff \(i\), after anomaly grouping and exact telescopes.
Suppose
\[
                              |{\cal L}_i(O)|\le r_i(O).        \tag{68.13}
\]

\begin{theorem}[Ward passage]
\label{thm:v11v68-passage}
Assume for every cylinder \(O\):
\begin{enumerate}
\item each expectation entering \({\cal L}_i(O)\) converges to the
corresponding continuum expectation;
\item the cutoff Ward operator \(R_{\epsilon,i}O\) converges in that
expectation;
\item \(r_i(O)\to0\).
\end{enumerate}
Then the limiting cylinder functional satisfies the continuum Ward
identity
\[
                              {\cal L}_\infty(O)=0.            \tag{68.14}
\]
\end{theorem}

\begin{proof}
Pass to the limit in every term of \({\cal L}_i(O)\).  The triangle
inequality and (68.13) give
\[
                         |{\cal L}_\infty(O)|
                    \le\limsup_i r_i(O)=0.
\]
\end{proof}

\begin{proposition}[Summable increments do not force zero]
\label{prop:v11v68-summable}
A bound
\[
                    |{\cal L}_{i+1}(O)-{\cal L}_i(O)|
                              \le\epsilon_i(O),\qquad
                    \sum_i\epsilon_i(O)<\infty               \tag{68.15}
\]
proves existence of a limiting Ward defect but does not prove that it
is zero.
\end{proposition}

\begin{proof}
A summable-increment sequence converges, but its limit can be any
constant.  For example, \({\cal L}_i=1\) has zero increments and
nonzero limit.
\end{proof}

To prove zero from adjacent estimates, one needs an exact zero at every
finite cutoff, the direct remainder bound (68.13), or a renormalization
condition fixing the limiting constant.

\subsection{Diffeomorphism and Einstein consistency}

For the gravitational sector, the analogous identity couples the
diffeomorphism constraint, gauge fixing, metric measure, and matter
stress tensor.  In the continuum target it must imply the compatibility
condition
\[
                              \nabla^\mu T_{\mu\nu}=0          \tag{68.16}
\]
that matches the contracted Bianchi identity
\(\nabla^\mu G_{\mu\nu}=0\).  A discrete regulator can violate either
identity by \({\cal E}_{\epsilon,i}\); recovery requires
Theorem~\ref{thm:v11v68-passage}, not only pointwise convergence of an
effective action.

\begin{ledger}[Ward acquisition]
\label{led:v11v68-data}
For each internal gauge and gravitational symmetry, the programme must
provide:
\begin{enumerate}
\item the finite-cutoff change-of-variables vector field and its
measure divergence;
\item grouped local anomaly coefficients;
\item exact cochain or momentum telescopes;
\item interval bounds for irreducible regulator remainders;
\item boundary-current decay;
\item convergence of transformed cylinder observables;
\item a remainder tending to zero, rather than only summable
increments;
\item global determinant-line and topological-sector compatibility.
\end{enumerate}
\end{ledger}

\begin{firewall}[Ward identity not acquired]
\label{fire:v11v68-boundary}
V68 proves the finite-dimensional identity and its passage criterion.
It does not show that the chosen SM--Einstein regulator has an invariant
measure, a vanishing gravitational regulator defect, a disappearing
boundary flux, or a globally trivial chiral phase.  The earlier anomaly
arithmetic closes only the local perturbative species sum.
\end{firewall}

\begin{decision}[Next acquisition]
\label{dec:v11v68-next}
V69 will return to the complex chiral reweighting bottleneck.  It will
derive a stable reference-to-target expectation bound from numerator
and denominator increments, state a nonvanishing phase-denominator
criterion, and show how determinant-line overlap errors enter the
adjacent series without being hidden by modulus estimates.
\end{decision}

\begin{assessment}[Progress after V68]
\label{ass:v11v68-status}
The Ward row now has an exact finite-cutoff origin and a rigorous
continuum passage theorem.  Momentum shifts and cochain divergences are
removed algebraically, while anomaly, regulator, and boundary terms
remain explicit.  Summability alone is correctly distinguished from
vanishing of the limiting defect.
\end{assessment}

\subsection{Parent record}

This V68 note was formed by appending the preceding material to the
validated V67 source, whose record was
\[
\begin{array}{c|c}
\text{lines}&21060\\
\text{bytes}&663112\\
\text{SHA--256}&
\texttt{0BAF1082E462F996E9FDD8FA956E2E32D6680C1775D4E51BFF01B88B6BF839DA}.
\end{array}                                                     \tag{68.17}
\]

% =============================================================================
% END APPENDED ELEVENTH-SERIES V68 MATERIAL
% =============================================================================
% =============================================================================
% BEGIN APPENDED ELEVENTH-SERIES V69 MATERIAL
% =============================================================================

\section{Complex chiral reweighting, phase denominators, and determinant
line gluing}
\label{sec:v11v69}

V42 already gives the algebraic stability estimate for a quotient of
a complex numerator by a nonzero denominator.  The unresolved input is
a uniform lower bound for that denominator and a globally compatible
chiral phase.  This section supplies usable sufficient criteria and
places overlap errors in the adjacent measure ledger.

\subsection{Complex measures on a fixed cylinder}

Let \(\sigma_i\) be finite complex measures on a common measurable
space, with
\[
                         Z_i=\sigma_i(1)\ne0,\qquad
                         \omega_i=Z_i^{-1}\sigma_i.            \tag{69.1}
\]

\begin{theorem}[Total-variation reference-to-target convergence]
\label{thm:v11v69-tv}
Suppose
\[
                         \sum_i\lVert\sigma_{i+1}-\sigma_i
                                      \rVert_{\rm TV}<\infty. \tag{69.2}
\]
Then \(\sigma_i\) converges in total variation to a finite complex
measure \(\sigma_\infty\), and \(Z_i\to Z_\infty=\sigma_\infty(1)\).
If
\[
                              |Z_\infty|\ge\zeta>0,            \tag{69.3}
\]
then, eventually,
\[
                              |Z_i|\ge\zeta/2,                 \tag{69.4}
\]
and \(\omega_i\to\omega_\infty=Z_\infty^{-1}\sigma_\infty\) in total
variation.
\end{theorem}

\begin{proof}
Finite complex measures form a Banach space in total variation, so
(69.2) makes \(\sigma_i\) Cauchy.  Evaluation at the constant function
one has norm one, hence \(Z_i\to Z_\infty\), which gives (69.4).
Finally,
\[
 \left\|\frac{\sigma_i}{Z_i}
          -\frac{\sigma_\infty}{Z_\infty}\right\|_{\rm TV}
 \le\frac{\lVert\sigma_i-\sigma_\infty\rVert_{\rm TV}}{|Z_i|}
 +\lVert\sigma_\infty\rVert_{\rm TV}
       \frac{|Z_i-Z_\infty|}{|Z_iZ_\infty|},                 \tag{69.5}
\]
and both terms tend to zero.
\end{proof}

For changing cutoff spaces, this theorem is applied to every fixed
cylinder pushforward.  Compatibility of the limiting pushforwards is
still required.

\begin{corollary}[Bounded cylinder expectations]
\label{cor:v11v69-cylinder}
Under Theorem~\ref{thm:v11v69-tv}, every bounded cylinder observable
\(F\) satisfies
\[
                         \omega_i(F)\longrightarrow
                         \omega_\infty(F),                    \tag{69.6}
\]
with error at most
\(\lVert F\rVert_\infty
 \lVert\omega_i-\omega_\infty\rVert_{\rm TV}\).
\end{corollary}

\begin{proof}
This is the dual characterization of total variation.
\end{proof}

Total-variation convergence is stronger than the minimum cylinderwise
convergence needed by V42, but it gives a clean sufficient target.

\subsection{Modulus measure and the phase mean}

Let a positive reference probability be \(\mu\), and write an
integrable chiral weight as
\[
                         w=|w|e^{i\vartheta}.                 \tag{69.7}
\]
Assume \(0<M:=\int|w|\,d\mu<\infty\), and define the modulus-reweighted
probability
\[
                              d\nu=M^{-1}|w|\,d\mu.            \tag{69.8}
\]
Then
\[
                         Z=\int w\,d\mu
                           =M\,z_\vartheta,\qquad
                         z_\vartheta={\mathbb E}_\nu
                                      e^{i\vartheta}.         \tag{69.9}
\]

\begin{theorem}[Circular concentration criterion]
\label{thm:v11v69-circular}
If there is a deterministic angle \(\vartheta_0\) such that
\[
                  {\mathbb E}_\nu
                  |e^{i\vartheta}-e^{i\vartheta_0}|
                              \le\delta<1,                    \tag{69.10}
\]
then
\[
                              |z_\vartheta|\ge1-\delta,\qquad
                              |Z|\ge M(1-\delta).              \tag{69.11}
\]
\end{theorem}

\begin{proof}
The triangle inequality gives
\[
 |z_\vartheta-e^{i\vartheta_0}|
 \le{\mathbb E}_\nu
     |e^{i\vartheta}-e^{i\vartheta_0}|
 \le\delta.
\]
The reverse triangle inequality proves (69.11).
\end{proof}

\begin{corollary}[Arc-probability criterion]
\label{cor:v11v69-arc}
Suppose for some \(0\le\alpha<\pi/2\) and \(0\le\varepsilon<1\),
\[
                \nu\{|\vartheta-\vartheta_0|_{\mathbb T}
                                  \le\alpha\}\ge1-\varepsilon. \tag{69.12}
\]
Then
\[
              |z_\vartheta|
       \ge (1-\varepsilon)\cos\alpha-\varepsilon.             \tag{69.13}
\]
The bound is positive when
\((1-\varepsilon)\cos\alpha>\varepsilon\).
\end{corollary}

\begin{proof}
Rotate by \(e^{-i\vartheta_0}\) and take the real part.  On the good arc
the real part is at least \(\cos\alpha\); on its complement it is at
least \(-1\).  Thus
\[
 \operatorname{Re}(e^{-i\vartheta_0}z_\vartheta)
 \ge(1-\varepsilon)\cos\alpha-\varepsilon.
\]
The modulus is at least a positive real part.
\end{proof}

These criteria are sufficient but strong.  A bounded phase variance
modulo no chosen branch does not by itself prevent cancellation around
the circle.

\subsection{Stability along adjacent cutoffs}

Let
\[
                         z_i={\mathbb E}_{\nu_i}e^{i\vartheta_i}.
                                                                  \tag{69.14}
\]

\begin{proposition}[Anchored phase denominator]
\label{prop:v11v69-anchor}
If at some \(i_0\),
\[
                         |z_{i_0}|\ge\zeta_0>0,\qquad
                         \sum_{i=i_0}^{\infty}|z_{i+1}-z_i|
                                  \le\delta<\zeta_0,           \tag{69.15}
\]
then
\[
                         \inf_{i\ge i_0}|z_i|
                              \ge\zeta_0-\delta>0              \tag{69.16}
\]
and the limit \(z_\infty\) obeys the same lower bound.
\end{proposition}

\begin{proof}
For every \(n\ge i_0\),
\[
 |z_n-z_{i_0}|
 \le\sum_{i=i_0}^{n-1}|z_{i+1}-z_i|
 \le\delta.
\]
Apply the reverse triangle inequality and then pass to the limit.
\end{proof}

This theorem needs a quantitative anchor and a total variation smaller
than it.  Mere summability of phase increments does not prevent the
path from crossing zero.

\begin{lemma}[One coupling bound for the phase mean]
\label{lem:v11v69-coupling}
Let \((\Theta_i,\Theta_{i+1})\) be any coupling of the two phase
variables, and let \(\nu_i,\nu_{i+1}\) be their laws after a common
comparison has been chosen.  Then
\[
 |{\mathbb E}e^{i\Theta_{i+1}}
       -{\mathbb E}e^{i\Theta_i}|
 \le{\mathbb E}|e^{i\Theta_{i+1}}-e^{i\Theta_i}|
 \le{\mathbb E}\min\{2,|\Theta_{i+1}-\Theta_i|\}.             \tag{69.17}
\]
\end{lemma}

\begin{proof}
Use the triangle inequality and
\(|e^{ia}-e^{ib}|\le\min\{2,|a-b|\}\).
\end{proof}

When the modulus measures themselves differ, the coupling must include
their transport or Radon--Nikodym comparison; comparing phase functions
under only one of the two laws omits a measure increment.

\subsection{Determinant-line gluing}

Let \((U_a)\) cover the bosonic background space.  Suppose local
chiral determinants \(w_a\) satisfy on overlaps
\[
                              w_b=g_{ab}w_a,\qquad
                              g_{ab}:U_a\cap U_b\to U(1),      \tag{69.18}
\]
with
\[
                              g_{ab}g_{bc}g_{ca}=1            \tag{69.19}
\]
on triple overlaps.  These are transition functions of a determinant
line bundle.

\begin{theorem}[Scalar phase exists exactly under trivialization]
\label{thm:v11v69-trivial}
There are nowhere-zero local functions \(t_a:U_a\to U(1)\) such that
\[
                              \widetilde w_a=t_aw_a            \tag{69.20}
\]
agree on every overlap if and only if
\[
                              g_{ab}=t_b^{-1}t_a.              \tag{69.21}
\]
In that case the \(\widetilde w_a\) define one global complex weight.
\end{theorem}

\begin{proof}
Agreement means
\(t_bw_b=t_aw_a\).  Using (69.18) and cancelling \(w_a\) where
nonzero gives \(t_bg_{ab}=t_a\), which is (69.21).  Conversely,
(69.21) inserted into (69.18) gives
\(t_bw_b=t_aw_a\), so the local functions glue.
\end{proof}

Equation (69.21) is the Cech coboundary condition.  Local perturbative
anomaly cancellation constrains the curvature of this line; it does
not by itself prove trivial holonomy and (69.21).

\begin{proposition}[Overlap mismatch enters the numerator]
\label{prop:v11v69-overlap}
Suppose chosen local scalar representatives obey on active overlaps
\[
                              |t_bw_b-t_aw_a|\le\delta_i.      \tag{69.22}
\]
If a partition of unity is changed while gluing a bounded observable
\(F\), the induced numerator ambiguity is bounded by
\[
                         C_{\rm ov}\lVert F\rVert_\infty
                         \delta_i\,\mu_i({\cal O}_i),          \tag{69.23}
\]
where \({\cal O}_i\) is the union of active overlaps and
\(C_{\rm ov}\) is the finite overlap multiplicity.
\end{proposition}

\begin{proof}
Subtract the two partitioned local sums.  Terms off the overlaps agree.
On an overlap, use (69.22), the fact that partition weights lie in
\([0,1]\), and sum at most \(C_{\rm ov}\) active terms.  Integration
over \({\cal O}_i\) proves (69.23).
\end{proof}

Thus approximate phase gluing requires the series of overlap
ambiguities to be summable and, for a unique target, to vanish in the
limit.  An exact nontrivial holonomy cannot be repaired by refining the
cover.

\subsection{Reference-to-target ledger}

For a fixed cylinder, write the complex target measure as
\[
                              d\sigma_i=w_i\,d\mu_i^+          \tag{69.24}
\]
with positive reference \(\mu_i^+\).  A direct sufficient increment
bound is
\[
\begin{split}
 \lVert\sigma_{i+1}-\sigma_i\rVert_{\rm TV}
 &\le
 \int|w_{i+1}-w_i^{\uparrow}|\,d\mu_{i+1}^+\\
 &\quad+
 \lVert w_i^{\uparrow}\rVert_{L^q(\mu_{i+1}^+)}
 \left\|\frac{d(\mu_i^+)^\uparrow}
              {d\mu_{i+1}^+}-1\right\|_{L^{q'}(\mu_{i+1}^+)},
                                                                  \tag{69.25}
\end{split}
\]
whenever the displayed transport and densities exist.  The first row
contains modulus, local phase, and determinant-line overlap increments;
the second contains the positive reference measure comparison.

\begin{ledger}[Chiral target acquisition]
\label{led:v11v69-data}
The full target construction must provide:
\begin{enumerate}
\item a positive modulus reference with uniform moments;
\item summable cylinder pushforward increments such as (69.25);
\item a phase anchor and bound (69.15), or a uniform circular
concentration criterion;
\item exact or vanishing-error determinant-line gluing;
\item global anomaly and topological-sector compatibility;
\item a nonzero limiting denominator;
\item a fermionic reflection-positivity proof distinct from scalar
positivity of the complex weight.
\end{enumerate}
\end{ledger}

\begin{firewall}[The sign problem remains]
\label{fire:v11v69-boundary}
No current estimate proves (69.10), (69.12), or (69.15) for the full
chiral Standard Model determinant coupled to the fluctuating metric.
The exact local anomaly arithmetic does not bound the phase mean away
from zero.  Total-variation convergence without (69.3) also does not
define normalized target expectations.
\end{firewall}

\begin{decision}[Mandatory V70 audit]
\label{dec:v11v69-next}
The next version is the required companion audit.  It will inspect
whether the YM orbit-pairing programme or NS decay programme has
advanced and test any new result against the phase, quadrature, and
clustering ledgers.  V71 will then pursue the strongest remaining
reference-to-target estimate.
\end{decision}

\begin{assessment}[Progress after V69]
\label{ass:v11v69-status}
The complex target measure now has a rigorous convergence route and
three explicit denominator mechanisms.  Determinant-line transition
functions are treated as gluing data rather than an invisible phase,
and their overlap error enters the numerator variation.  The missing
physical input is a uniform noncancellation estimate for the actual
chiral phase.
\end{assessment}

\subsection{Parent record}

This V69 note was formed by appending the preceding material to the
validated V68 source, whose record was
\[
\begin{array}{c|c}
\text{lines}&21397\\
\text{bytes}&674874\\
\text{SHA--256}&
\texttt{DC3D764005743B78E1FF9F987FD161F97B4A45105046C98F792A8D5AACDEC1F5}.
\end{array}                                                     \tag{69.26}
\]

% =============================================================================
% END APPENDED ELEVENTH-SERIES V69 MATERIAL
% =============================================================================

% =============================================================================
% BEGIN APPENDED ELEVENTH-SERIES V70 MATERIAL
% =============================================================================

\section{Companion audit at V70: exact Ward cycles and repair of the
three-sample derivative lemma}
\label{sec:v11v70}

The mandatory audit finds Yang--Mills V38, containing new V37 and V38
blocks, while Navier--Stokes remains at V26.

\subsection{Audited records}

\begin{definition}[V70 companion record]
\label{def:v11v70-record}
The inspected files are
\[
\begin{array}{c|r|l}
\text{file}&\text{bytes}&\text{SHA--256}\\ \hline
\texttt{Renewal\_Yang\_Mills\_Mass\_Gap\_Research\_Notes\_V38.tex}
 &536904&
\texttt{ED5ACB149AF22C720DD516913B5A8AD3BFA15FE710C106F93FFE5E4DC3238390}\\
\texttt{Renewal\_NS\_Stability\_Research\_Notes\_V26.tex}
 &278283&
\texttt{82C887F8C2C69E4F28FAD7C3DC8DE5B9099CB5A4BF431EBA1FFF3E91935978AD}.
\end{array}                                                     \tag{70.1}
\]
The YM file has advanced by two substantive blocks.  The NS file is
unchanged.
\end{definition}

\subsection{Accepted V37 transfer}

The YM V37 block proves signed hypercubic covariance only for the
complete measured scalar, not for raw tree-coordinate matrices.  It
then derives exact orbit--stabilizer grid weights and a cyclewise Ward
telescope at nonzero periodic external momentum.  These results agree
with the independently derived general cards in V67--V68 and supply a
concrete positive-endpoint instance.  They are cited rather than
reproved.

Its main remaining step is the uniform passage from nonzero grid
momenta to formal derivatives at zero external momentum.

\subsection{The V38 notation defect}

The V38 source states a lemma for \(A\in C^3\) with zero values at
\(h,2h,3h\), but its displayed formulas literally use
\(A''\), \(A''''\), and \(A''''''\).  The proof invokes the fundamental
theorem of calculus, first and second forward differences, and Cauchy
bounds only through order three.  Its final constants also correspond
to derivatives of orders one, two, and three.  Thus the doubled
apostrophes are a source-escaping defect and the displayed lemma cannot
be imported literally.

The corrected statement is as follows.

\begin{lemma}[Correct three-zero derivative card]
\label{lem:v11v70-three-zero}
Let \(A\in C^3([0,3h])\) and suppose
\[
                              A(h)=A(2h)=A(3h)=0.              \tag{70.2}
\]
Then
\[
\begin{split}
 |A(0)|&\le h\sup_{[0,h]}|A'|,\\
 |A'(0)|&\le\frac32h\sup_{[0,2h]}|A''|,\\
 |A''(0)|&\le2h\sup_{[0,3h]}|A'''|.                          \tag{70.3}
\end{split}
\]
\end{lemma}

\begin{proof}
The first row follows from
\[
                              A(0)=-\int_0^hA'(s)\,ds.
\]
Since \(A(2h)-A(h)=0\),
\[
                              0=\frac1h\int_h^{2h}A'(s)\,ds.
\]
Subtract \(A'(0)\) inside the average and use
\[
                              |A'(s)-A'(0)|
                        \le s\sup_{[0,2h]}|A''|.
\]
The average of \(s\) on \([h,2h]\) is \(3h/2\), proving the second row.

The second forward difference also vanishes:
\[
 0=\frac{A(3h)-2A(2h)+A(h)}{h^2}
   =\int_0^1\int_0^1A''(h(1+u+v))\,du\,dv.                  \tag{70.4}
\]
Subtract \(A''(0)\) and use
\[
 |A''(h(1+u+v))-A''(0)|
 \le h(1+u+v)\sup_{[0,3h]}|A'''|.
\]
The average of \(1+u+v\) on the unit square is \(2\), proving the last
row.
\end{proof}

\begin{corollary}[Correct analytic Ward-limit constants]
\label{cor:v11v70-cauchy}
If \(A\) is analytic and bounded by \(B\) on every complex circle of
radius \(\rho/2\) centered in \([0,3h]\), then
\[
\begin{split}
 |A(0)|&\le\frac{2B}{\rho}h,\\
 |A'(0)|&\le\frac{12B}{\rho^2}h,\\
 |A''(0)|&\le\frac{96B}{\rho^3}h.                            \tag{70.5}
\end{split}
\]
\end{corollary}

\begin{proof}
Cauchy's estimate gives
\[
                         \sup|A^{(r)}|
                    \le r!(2/\rho)^rB,\qquad r=1,2,3.
\]
Insert these three bounds in (70.3).
\end{proof}

These are the numerical constants displayed by YM V38, now attached
to the correct derivative orders.

\subsection{Scope of the accepted Ward result}

The audited V38 constructs a common external-momentum strip and a
volume-independent scalar bound for the positive one-loop Gaussian
endpoint homotopy.  Three exact nonzero grid Ward values then give
(70.5) with \(h=2\pi/M\).  This closes the normalized longitudinal
formal Ward rows for that declared homotopy, subject to the corrected
notation.

\begin{warning}[External strip versus internal quadrature strip]
\label{warn:v11v70-strips}
The V38 strip is in the complex external scalar \(z\), uniformly over
real internal momentum.  The V66 trapezoidal theorem needs a complex
strip in every internal Brillouin coordinate.  One strip cannot be
substituted for the other.
\end{warning}

\begin{warning}[Positive endpoint versus chiral target]
\label{warn:v11v70-target}
The audited scalar bound uses positive finite matrices along a
Wilson-to-positive endpoint homotopy.  A normalized complex chiral
target response also contains a phase denominator.  A real-axis lower
bound for that denominator is insufficient for Cauchy's theorem; it
must remain nonzero on the complex external strip.
\end{warning}

\subsection{Audit decision}

\begin{decision}[V71 acquisition]
\label{dec:v11v70-next}
V71 will combine Lemma~\ref{lem:v11v70-three-zero} with V69.  It will
derive a zero-free complex strip for a normalized numerator
\(N(z)/Z(z)\) from an anchored denominator and a derivative bound,
then prove an \(O(M^{-1})\) formal Ward limit for the complex target
whenever three exact nonzero grid Ward identities and a uniform
numerator bound are available.
\end{decision}

\begin{firewall}[Audit boundary]
\label{fire:v11v70-boundary}
The V37--V38 result concerns the positive Gaussian YM homotopy and does
not provide the SM chiral phase denominator, an internal-momentum
quadrature, nonlinear cylinder comparison, reflection positivity,
clustering, or the Einstein constraint limit.  The notation repair
above validates the intended calculus but does not enlarge the
companion's physical scope.  NS supplies no new input.
\end{firewall}

\begin{assessment}[Progress after V70]
\label{ass:v11v70-status}
The required companion audit is complete.  It imports the concrete
cyclewise Ward structure, repairs the malformed derivative notation,
and isolates the exact additional condition needed for the chiral
SM target: a zero-free normalized denominator throughout the external
complex strip.
\end{assessment}

\subsection{Parent record}

This V70 note was formed by appending the preceding material to the
validated V69 source, whose record was
\[
\begin{array}{c|c}
\text{lines}&21770\\
\text{bytes}&687468\\
\text{SHA--256}&
\texttt{F09A6DA4FD3116C7D6F8B6971E7F7EF3822E4DEAD04E8973F769ABB38ADD1078}.
\end{array}                                                     \tag{70.6}
\]

% =============================================================================
% END APPENDED ELEVENTH-SERIES V70 MATERIAL
% =============================================================================
% =============================================================================
% BEGIN APPENDED ELEVENTH-SERIES V71 MATERIAL
% =============================================================================

\section{A zero-free chiral strip and the normalized Ward limit}
\label{sec:v11v71}

The positive YM Ward-limit card uses a uniformly bounded analytic
response.  For the chiral target, the response is a quotient of an
analytic numerator by a phase-sensitive partition denominator.  This
section gives a local zero-free criterion for that quotient and the
resulting Ward-limit schedule.

\subsection{Protecting the denominator}

Let \(N(z)\) and \(Z(z)\) be analytic on an open neighborhood of the
closed \(r\)-neighborhood of the segment \([0,3h]\):
\[
                         {\cal R}_{h,r}
          =\{z\in\mathbb C:d(z,[0,3h])\le r\}.                \tag{71.1}
\]
Assume
\[
                 |Z(0)|\ge\zeta>0,\qquad
                 \sup_{{\cal R}_{h,r}}|Z'(z)|\le L.          \tag{71.2}
\]

\begin{theorem}[Anchored zero-free rectangle]
\label{thm:v11v71-zero-free}
If
\[
                              L(3h+r)\le\theta\zeta
                              \quad\text{for some }0<\theta<1,\tag{71.3}
\]
then
\[
                         |Z(z)|\ge(1-\theta)\zeta
                              \quad(z\in{\cal R}_{h,r}).       \tag{71.4}
\]
Consequently \(A(z)=N(z)/Z(z)\) is analytic there, and if
\[
                         \sup_{{\cal R}_{h,r}}|N(z)|\le M_N,  \tag{71.5}
\]
then
\[
                         \sup_{{\cal R}_{h,r}}|A(z)|
                         \le B:=\frac{M_N}{(1-\theta)\zeta}.  \tag{71.6}
\]
\end{theorem}

\begin{proof}
For \(z\in{\cal R}_{h,r}\), join zero to \(z\) by a path of length at
most \(3h+r\) inside the declared analytic neighborhood.  Then
\[
 |Z(z)-Z(0)|\le L(3h+r)\le\theta\zeta.
\]
The reverse triangle inequality gives (71.4).  Division by a zero-free
analytic function proves analyticity of \(A\), and (71.6) is immediate.
\end{proof}

A stronger real-axis phase criterion can protect an entire strip.

\begin{corollary}[Uniform vertical protection]
\label{cor:v11v71-vertical}
Suppose for every real \(x\),
\[
                         |Z(x)|\ge\zeta,\qquad
                         \sup_{|\operatorname{Im}z|\le r}|Z'(z)|
                                  \le L,                      \tag{71.7}
\]
and \(Lr\le\theta\zeta\).  Then
\[
                         |Z(x+iy)|\ge(1-\theta)\zeta
                              \quad(|y|\le r).                 \tag{71.8}
\]
\end{corollary}

\begin{proof}
Integrate \(Z'\) along the vertical segment from \(x\) to \(x+iy\).
\end{proof}

The real-axis lower bound in (71.7) can in principle come from the
circular or arc criteria of V69 uniformly in \(x\).  It has not yet
been established for the target determinant.

\subsection{Normalized three-sample Ward card}

Assume the exact nonzero grid identities give
\[
                              A(h)=A(2h)=A(3h)=0.              \tag{71.9}
\]

\begin{theorem}[Complex normalized Ward estimate]
\label{thm:v11v71-ward}
Under (71.1)--(71.6) and (71.9),
\[
\begin{split}
 |A(0)|&\le \frac{B}{r}h,\\
 |A'(0)|&\le \frac{3B}{r^2}h,\\
 |A''(0)|&\le \frac{12B}{r^3}h.                              \tag{71.10}
\end{split}
\]
\end{theorem}

\begin{proof}
Every circle of radius \(r\) centered on \([0,3h]\) lies in
\({\cal R}_{h,r}\).  Cauchy's estimate and (71.6) give
\[
                         \sup_{[0,3h]}|A^{(m)}|
                              \le m!B r^{-m},\qquad m=1,2,3.  \tag{71.11}
\]
Insert these bounds into the corrected three-zero lemma
\ref{lem:v11v70-three-zero}.
\end{proof}

For \(r=\rho/2\), (71.10) becomes the constants
\(2B/\rho\), \(12B/\rho^2\), and \(96B/\rho^3\) recorded in V70.

\subsection{Cutoff-dependent schedule}

At momentum-grid size \(M\), put \(h_M=2\pi/M\).  Allow
\[
                         \zeta_M,\quad M_{N,M},\quad r_M       \tag{71.12}
\]
to depend on the cutoff, and fix a common \(\theta<1\).  Define
\[
                         B_M=\frac{M_{N,M}}
                                  {(1-\theta)\zeta_M}.         \tag{71.13}
\]

\begin{corollary}[Ward-limit schedule]
\label{cor:v11v71-schedule}
If the denominator condition (71.3) and the three exact identities
hold for all sufficiently large \(M\), and
\[
 \frac{h_MB_M}{r_M}\to0,\qquad
 \frac{h_MB_M}{r_M^2}\to0,\qquad
 \frac{h_MB_M}{r_M^3}\to0,                                  \tag{71.14}
\]
then the normalized formal rows \(A_M(0)\), \(A_M'(0)\), and
\(A_M''(0)\) all vanish.  The last condition implies the first two
when \(r_M\) is uniformly bounded above.
\end{corollary}

\begin{proof}
Apply (71.10) at each \(M\) and take the limit.
\end{proof}

For power laws
\[
 \zeta_M\gtrsim M^{-z},\qquad
 M_{N,M}\lesssim M^n,\qquad
 r_M\gtrsim M^{-a},                                          \tag{71.15}
\]
the second-derivative row closes under the sufficient inequality
\[
                              n+z+3a<1.                       \tag{71.16}
\]
A phase denominator that decays exponentially in the volume cannot
satisfy this polynomial schedule without compensating numerator or
strip estimates.

\subsection{Obtaining the derivative bound for \(Z\)}

Suppose
\[
                              Z(z)=\int w(z,\phi)\,d\mu(\phi) \tag{71.17}
\]
on a fixed positive reference measure, and \(w\) is analytic in \(z\).
If there is an integrable \(G(\phi)\) with
\[
                         |\partial_zw(z,\phi)|\le G(\phi)
                              \quad(z\in{\cal R}_{h,r}),       \tag{71.18}
\]
then differentiation under the integral gives
\[
                              |Z'(z)|\le\int G\,d\mu.          \tag{71.19}
\]

\begin{proof}
Dominated differentiation for analytic parameter integrals gives
\(Z'=\int\partial_zw\,d\mu\).  Take absolute values.
\end{proof}

If the positive reference measure itself depends on \(z\), its
logarithmic derivative contributes a score term and must be included in
\(G\).  This is the same stationary and density row that the YM
complete response keeps before taking the Ward sum.

\subsection{Determinant-line compatibility}

Cauchy's theorem applies to one analytic scalar \(Z(z)\).  Local
determinant representatives related by nontrivial transition functions
do not define such a scalar until the gluing condition (69.21) is
satisfied.

\begin{proposition}[Analytic gluing requirement]
\label{prop:v11v71-gluing}
Suppose the zero-free background chamber has analytic transition
functions \(g_{ab}(z,\phi)\) and analytic trivializers \(t_a\)
satisfying \(g_{ab}=t_b^{-1}t_a\).  If the locally rephased integrands
\(t_aw_a\) agree and obey one common domination bound, then (71.17)
defines a global analytic denominator.  Without the coboundary
condition, local nonvanishing does not define a global \(Z(z)\).
\end{proposition}

\begin{proof}
The coboundary identity makes the local integrands equal on overlaps
by Theorem~\ref{thm:v11v69-trivial}; hence they glue to one analytic
function of field and \(z\).  Dominated analytic integration gives a
global analytic \(Z\).  If the line is nontrivial, no choice of local
phases makes the representatives a global scalar, so the integral in
(71.17) lacks the declared meaning.
\end{proof}

\subsection{Approximate Ward samples}

If the three nonzero grid Ward values are not exact but satisfy
\[
                              |A(nh)|\le\varepsilon_h,
                              \qquad n=1,2,3,                 \tag{71.20}
\]
then interpolation errors are amplified by inverse powers of \(h\)
when derivatives are reconstructed.  In particular, if \(P\) is the
quadratic interpolant through these three values, direct Lagrange
calculation gives
\[
 |P(0)|\le7\varepsilon_h,\quad
 |P'(0)|\le8\varepsilon_hh^{-1},\quad
 |P''(0)|\le4\varepsilon_hh^{-2}.                             \tag{71.21}
\]

\begin{proposition}[Defect rate required for derivative Ward rows]
\label{prop:v11v71-defect}
In addition to the analytic remainder bounds, approximate samples can
give vanishing rows through order two only if
\[
                         \varepsilon_h=o(h^2)                 \tag{71.22}
\]
for the second derivative, with the sharper analytic interpolation
remainder included separately.
\end{proposition}

\begin{proof}
The quadratic interpolant contributes the three quantities in
(71.21).  The second derivative term vanishes only when
\(\varepsilon_hh^{-2}\to0\).  The zero-interpolant remainder is
controlled by the third derivative as in V70.
\end{proof}

Thus a regulator defect merely tending to zero may be too slow for the
formal second Ward derivative.

\subsection{Target Ward acquisition}

\begin{ledger}[Complex Ward-limit data]
\label{led:v11v71-data}
For the chiral SM target, one must provide:
\begin{enumerate}
\item exact anomaly-grouped nonzero grid Ward identities, or defects
satisfying (71.22);
\item a global analytic determinant-line trivialization on the
relevant chamber;
\item a real-axis phase denominator anchor;
\item a derivative or vertical-variation bound protecting a complex
strip;
\item a numerator envelope \(M_{N,M}\);
\item the scale condition (71.14);
\item control of the chamber complement and topological sectors.
\end{enumerate}
\end{ledger}

\begin{firewall}[Conditional extension only]
\label{fire:v11v71-boundary}
V71 extends the corrected YM calculus to a normalized complex quotient,
but none of the denominator, analytic gluing, or numerator hypotheses
has been proved for the complete SM determinant on fluctuating
Einstein backgrounds.  The theorem does not solve the sign problem or
establish the target Ward identity.
\end{firewall}

\begin{decision}[Next acquisition]
\label{dec:v11v71-next}
V72 will investigate a usable phase anchor through cumulants.  It will
derive a noncancellation criterion from an absolutely summable
connected-cumulant series for the phase, identify the required volume
scaling, and show why an extensive phase variance produces an
exponentially small denominator unless local cancellations remove it.
\end{decision}

\begin{assessment}[Progress after V71]
\label{ass:v11v71-status}
The normalized chiral Ward limit now has an explicit zero-free-strip
criterion and a sharp scale ledger.  The phase denominator, numerator
growth, and strip shrinkage appear in one condition; approximate Ward
zeros must converge faster than \(h^2\) for the second formal row.
The remaining core input is quantitative phase noncancellation.
\end{assessment}

\subsection{Parent record}

This V71 note was formed by appending the preceding material to the
validated V70 source, whose record was
\[
\begin{array}{c|c}
\text{lines}&21971\\
\text{bytes}&694575\\
\text{SHA--256}&
\texttt{CC68C305213F6FFCBFA72CFBEF5561CA2E93DB7DD62977DFEDA75760E329E2AB}.
\end{array}                                                     \tag{71.23}
\]

% =============================================================================
% END APPENDED ELEVENTH-SERIES V71 MATERIAL
% =============================================================================
% =============================================================================
% BEGIN APPENDED ELEVENTH-SERIES V72 MATERIAL
% =============================================================================

\section{Phase cumulants, shell increments, and the extensive sign
obstruction}
\label{sec:v11v72}

The phase denominator is
\[
                         z_\vartheta={\mathbb E}_\nu e^{i\Theta},
                                                                  \tag{72.1}
\]
after a global determinant-line phase and modulus measure have been
chosen.  Cumulants give a nonzero denominator when their series is
valid through unit argument.  They also show sharply why an extensive
phase variance normally produces exponential cancellation.

\subsection{An absolutely convergent cumulant criterion}

Let \(\Theta\) be a real random variable.  Write
\(\kappa_n(\Theta)\) for its real cumulants.

\begin{theorem}[Cumulant noncancellation]
\label{thm:v11v72-cumulant}
Assume:
\begin{enumerate}
\item \({\mathbb E}e^{r|\Theta|}<\infty\) for some \(r>1\);
\item the Taylor series
\[
                         K(t)=\sum_{n=1}^{\infty}
                       \kappa_n(\Theta)\frac{(it)^n}{n!}      \tag{72.2}
\]
has radius of convergence strictly larger than one.
\end{enumerate}
Then
\[
                         {\mathbb E}e^{it\Theta}=e^{K(t)}
                              \quad(|t|\le1),                  \tag{72.3}
\]
so in particular
\[
 \boxed{
 |{\mathbb E}e^{i\Theta}|
 =\exp\left\{\sum_{m=1}^{\infty}
          \frac{(-1)^m\kappa_{2m}(\Theta)}{(2m)!}\right\}
 \ge\exp\left\{-\sum_{m=1}^{\infty}
          \frac{|\kappa_{2m}(\Theta)|}{(2m)!}\right\}>0.}     \tag{72.4}
\]
\end{theorem}

\begin{proof}
The exponential moment makes the characteristic function analytic in
a complex neighborhood containing \([-1,1]\).  Near zero, its
logarithm has Taylor series (72.2), so the characteristic function
equals \(e^{K(t)}\).  The assumed radius makes \(e^{K}\) analytic on a
connected neighborhood of the closed unit disk.  The identity theorem
extends the equality to \(|t|\le1\), proving (72.3).  Since the
cumulants are real, odd terms in \(K(1)\) are imaginary and even terms
are real with alternating signs.  Taking the real part gives the
equality in (72.4), and the triangle inequality gives its lower bound.
\end{proof}

The radius assumption is substantive.  Formal cumulants at zero do not
prove that the characteristic function has no zero before \(t=1\).

\begin{corollary}[Uniform cumulant budget]
\label{cor:v11v72-uniform}
For a cutoff family \(\Theta_i\), a uniform bound
\[
                         \sup_i\sum_{m\ge1}
                         \frac{|\kappa_{2m}(\Theta_i)|}{(2m)!}
                              \le C_\Theta                    \tag{72.5}
\]
together with the analytic hypotheses of
Theorem~\ref{thm:v11v72-cumulant} gives
\[
                         \inf_i|{\mathbb E}e^{i\Theta_i}|
                              \ge e^{-C_\Theta}>0.             \tag{72.6}
\]
\end{corollary}

\begin{proof}
Apply (72.4) uniformly.
\end{proof}

\subsection{Connected local cumulants}

Suppose on a finite lattice
\[
                              \Theta_\Lambda=\sum_{x\in\Lambda}\theta_x.
                                                                  \tag{72.7}
\]
Multilinearity of joint cumulants gives
\[
 \kappa_n(\Theta_\Lambda)
 =\sum_{x_1,\ldots,x_n\in\Lambda}
                  \kappa(\theta_{x_1},\ldots,\theta_{x_n}).    \tag{72.8}
\]

\begin{proposition}[Extensive cumulant bound]
\label{prop:v11v72-extensive}
Assume translation invariance and
\[
 \sum_{x_2,\ldots,x_n}
 |\kappa(\theta_0,\theta_{x_2},\ldots,\theta_{x_n})|
                              \le c_n.                         \tag{72.9}
\]
Then
\[
                         |\kappa_n(\Theta_\Lambda)|
                              \le|\Lambda|c_n.                 \tag{72.10}
\]
Consequently (72.4) gives only
\[
 |{\mathbb E}e^{i\Theta_\Lambda}|
 \ge\exp\left\{-|\Lambda|
        \sum_{m\ge1}\frac{c_{2m}}{(2m)!}\right\}.             \tag{72.11}
\]
\end{proposition}

\begin{proof}
Fix \(x_1\) in (72.8), translate it to zero, and apply (72.9) to the
remaining sum.  There are \(|\Lambda|\) choices of \(x_1\).
Insert the resulting even-cumulant bounds in (72.4).
\end{proof}

Absolute summability of connected correlations therefore proves an
extensive phase free energy, not a volume-uniform denominator.  A
uniform bound needs further cancellation making the total even
cumulant budget \(O(1)\).

\subsection{The Gaussian obstruction}

\begin{proposition}[Exact Gaussian phase mean]
\label{prop:v11v72-gaussian}
If \(\Theta\) is Gaussian with mean \(m\) and variance \(v\), then
\[
                         {\mathbb E}e^{i\Theta}
                              =e^{im-v/2},\qquad
                         |{\mathbb E}e^{i\Theta}|=e^{-v/2}.    \tag{72.12}
\]
If \(v_\Lambda\ge c|\Lambda|\), the denominator is at most
\(e^{-c|\Lambda|/2}\).
\end{proposition}

\begin{proof}
Only the first two cumulants are nonzero.  Insert
\(\kappa_1=m\), \(\kappa_2=v\) into (72.2).
\end{proof}

Thus even perfect Gaussian control is harmful when the phase variance
is extensive.  The required input is cancellation or localization of
the total phase, not merely a central limit theorem.

\subsection{Shellwise phase increments}

Let
\[
                              \Theta_n=\sum_{j=1}^nX_j.        \tag{72.13}
\]

\begin{theorem}[Independent shell product]
\label{thm:v11v72-product}
If the \(X_j\) are independent, then
\[
                         {\mathbb E}e^{i\Theta_n}
                           =\prod_{j=1}^n\phi_j,\qquad
                         \phi_j={\mathbb E}e^{iX_j}.           \tag{72.14}
\]
If no \(\phi_j\) vanishes and
\[
                              \sum_j|1-\phi_j|<\infty,         \tag{72.15}
\]
the product converges to a nonzero limit.
\end{theorem}

\begin{proof}
Independence gives (72.14).  Condition (72.15) is the standard
absolute convergence criterion for an infinite product:
for all sufficiently large \(j\), \(|1-\phi_j|<1/2\), so the principal
series \(\sum_j\log\phi_j\) converges absolutely by
\(|\log\phi_j|\le2|1-\phi_j|\).  Its exponential is nonzero; the finite
initial product is nonzero by assumption.
\end{proof}

Exact independence is not expected between renewal shells.  A robust
nonfactorized alternative follows directly from V69.

\begin{proposition}[Coupled shell anchor]
\label{prop:v11v72-shell-anchor}
For arbitrary dependent increments,
\[
 \left|{\mathbb E}e^{i\Theta_n}-1\right|
 \le\sum_{j=1}^n{\mathbb E}
             |e^{iX_j}-1|
 \le\sum_{j=1}^n{\mathbb E}\min\{2,|X_j|\}.                  \tag{72.16}
\]
Hence if
\[
                              \sum_j{\mathbb E}
                              \min\{2,|X_j|\}\le\delta<1,      \tag{72.17}
\]
then
\[
                              \inf_n|{\mathbb E}e^{i\Theta_n}|
                                  \ge1-\delta.                 \tag{72.18}
\]
\end{proposition}

\begin{proof}
The telescoping identity
\[
 e^{i\Theta_n}-1
 =\sum_{j=1}^ne^{i\Theta_{j-1}}(e^{iX_j}-1)
\]
and \(|e^{i\Theta_{j-1}}|=1\) give (72.16).  The reverse triangle
inequality gives (72.18).
\end{proof}

This criterion is strong but needs neither independence nor a cumulant
branch.  It is suited to an adjacent construction if each shell phase
increment is absolutely summable in circular distance.

\subsection{Conditional shell factorization}

Let \({\cal F}_j\) be the shell filtration.  If
\[
                         {\mathbb E}[e^{iX_j}\mid{\cal F}_{j-1}]
                              =\phi_j                         \tag{72.19}
\]
is deterministic, then iteration gives (72.14) without independence.
If the conditional factor is random, its correlation with
\(e^{i\Theta_{j-1}}\) remains and cannot be replaced by its mean.

\begin{proof}
Under (72.19),
\[
 {\mathbb E}e^{i\Theta_j}
 ={\mathbb E}\left[
 e^{i\Theta_{j-1}}
 {\mathbb E}(e^{iX_j}\mid{\cal F}_{j-1})\right]
 =\phi_j{\mathbb E}e^{i\Theta_{j-1}}.
\]
Iterate.
\end{proof}

\subsection{Phase free-energy density}

When the thermodynamic limit exists, define
\[
                         f_{\rm ph}
 =-\lim_{|\Lambda|\to\infty}
       |\Lambda|^{-1}\log|{\mathbb E}e^{i\Theta_\Lambda}|.    \tag{72.20}
\]
A strictly positive \(f_{\rm ph}\) means exponential denominator
decay.  The V71 polynomial Ward schedule then fails unless the
numerator or quotient is reorganized so the same exponential factor
cancels before separate bounds are taken.

\begin{warning}[A lower bound can still be too small]
\label{warn:v11v72-small}
The positive lower bound in (72.11) is finite-volume rigorous but
exponentially small.  It does not provide the uniform or polynomial
\(\zeta_M\) required in (71.15)--(71.16).  Proving existence of a
nonzero denominator separately at every cutoff is insufficient.
\end{warning}

\subsection{Determinant-line and branch requirements}

Cumulants of a real lift \(\Theta\) require a global choice of phase.
If only \(e^{i\Theta}\) is defined, different lifts change ordinary
cumulants even though the circular mean (72.1) is unchanged.  The
determinant-line trivialization of V69 must therefore precede an
ordinary cumulant expansion.  Alternatively, use the circular
increment criterion (72.17), which depends only on phase ratios.

\begin{ledger}[Phase-anchor acquisition]
\label{led:v11v72-data}
A cumulant route to the chiral denominator must provide:
\begin{enumerate}
\item a global determinant-line phase or circular shell ratios;
\item exponential moments and a cumulant radius beyond \(t=1\);
\item connected-cumulant bounds with all factorial constants;
\item cancellation reducing the total even cumulant budget from
\(O(|\Lambda|)\) to \(O(1)\), or an alternative direct quotient;
\item uniformity in the metric, gauge, Higgs, and topological sectors;
\item compatibility with the external complex strip of V71.
\end{enumerate}
The shell route may replace items two through four by (72.17).
\end{ledger}

\begin{firewall}[Anomaly cancellation is not phase concentration]
\label{fire:v11v72-boundary}
Cancellation of the local perturbative gauge anomaly controls the
gauge variation of the determinant phase.  It does not imply that the
phase itself has \(O(1)\) variance or cumulants.  No estimate in the
current notes proves (72.5) or (72.17) for the full Standard Model
phase on fluctuating metrics.
\end{firewall}

\begin{decision}[Next acquisition]
\label{dec:v11v72-next}
V73 will derive a phase-localization theorem from a coboundary
decomposition.  If the shell phase is an exact discrete divergence plus
a summable remainder, the bulk contribution cancels and only boundary,
topological, and remainder terms enter (72.17).  This is the natural
upstream mechanism to seek after local anomaly cancellation.
\end{decision}

\begin{assessment}[Progress after V72]
\label{ass:v11v72-status}
The phase denominator now has two rigorous construction routes:
an absolutely convergent cumulant series with an \(O(1)\) even budget,
or summable circular shell increments.  The Gaussian example and
extensive cumulant bound show precisely why ordinary clustering alone
does not solve the sign problem.
\end{assessment}

\subsection{Parent record}

This V72 note was formed by appending the preceding material to the
validated V71 source, whose record was
\[
\begin{array}{c|c}
\text{lines}&22286\\
\text{bytes}&705593\\
\text{SHA--256}&
\texttt{83291171F6FF52F9179487AC26039428D4C6065562968ED8755653A96EDECC29}.
\end{array}                                                     \tag{72.21}
\]

% =============================================================================
% END APPENDED ELEVENTH-SERIES V72 MATERIAL
% =============================================================================
% =============================================================================
% BEGIN APPENDED ELEVENTH-SERIES V73 MATERIAL
% =============================================================================

\section{Coboundary localization of phase increments}
\label{sec:v11v73}

An extensive sum of local phase densities normally has extensive
cumulants.  A possible escape is an exact cochain divergence or cutoff
coboundary.  Such terms telescope algebraically.  This section proves
the resulting phase bound and keeps harmonic/topological terms
separate.

\subsection{Spatial cochain decomposition}

Let \(\Lambda\) be a finite oriented cell complex, let \(j\) be a real
one-cochain, and let \(r\) be a real zero-cochain.  Consider the phase
\[
                         \Theta_\Lambda
                    =\langle{\bf1},d_0^*j+r\rangle.           \tag{73.1}
\]

\begin{theorem}[Bulk divergence localization]
\label{thm:v11v73-divergence}
If \(\Lambda\) is closed or periodic, then
\[
                         \Theta_\Lambda=\sum_{x\in\Lambda}r(x).
                                                                  \tag{73.2}
\]
If \(\Lambda\) has boundary, then
\[
                         \Theta_\Lambda
                    =\operatorname{Flux}_{\partial\Lambda}(j)
                       +\sum_{x\in\Lambda}r(x).               \tag{73.3}
\]
\end{theorem}

\begin{proof}
Adjointness gives
\[
                         \langle{\bf1},d_0^*j\rangle
                              =\langle d_0{\bf1},j\rangle.
\]
This is zero on a closed complex.  On a complex with boundary, the
uncancelled incidence terms are exactly the oriented boundary flux,
by the discrete Stokes formula.
\end{proof}

Thus a genuine divergence has no bulk phase free energy.  A local
density that is merely small or has zero mean does not satisfy this
identity.

\subsection{Circular remainder control}

At renewal step \(i\), suppose the phase ratio is
\[
 e^{iX_i}
 =e^{i(B_i+R_i)},                                            \tag{73.4}
\]
where \(B_i\) is a boundary flux, \(R_i\) is a regulator remainder, and
integer multiples of \(2\pi\) have been omitted because their
exponential is one.

\begin{lemma}[Boundary and remainder increment]
\label{lem:v11v73-increment}
One has
\[
 |e^{iX_i}-1|
 \le\min\{2,|B_i|\}+\min\{2,|R_i|\}.                         \tag{73.5}
\]
\end{lemma}

\begin{proof}
Factor
\[
 e^{i(B_i+R_i)}-1
 =e^{iB_i}(e^{iR_i}-1)+(e^{iB_i}-1)
\]
and use \(|e^{it}-1|\le\min\{2,|t|\}\).
\end{proof}

\begin{theorem}[Coboundary phase anchor]
\label{thm:v11v73-anchor}
Let
\[
                         \Theta_n=\Theta_{i_0}
                                      +\sum_{i=i_0}^{n-1}X_i. \tag{73.6}
\]
If
\[
 |{\mathbb E}e^{i\Theta_{i_0}}|\ge\zeta_0>0                 \tag{73.7}
\]
and
\[
 \sum_{i=i_0}^{\infty}
 \left[
 {\mathbb E}\min\{2,|B_i|\}
 +{\mathbb E}\min\{2,|R_i|\}
 \right]\le\delta<\zeta_0,                                  \tag{73.8}
\]
then
\[
                         \inf_{n\ge i_0}
                         |{\mathbb E}e^{i\Theta_n}|
                              \ge\zeta_0-\delta>0.             \tag{73.9}
\]
\end{theorem}

\begin{proof}
The telescoping identity for phase exponentials gives
\[
 |{\mathbb E}e^{i\Theta_n}
       -{\mathbb E}e^{i\Theta_{i_0}}|
 \le\sum_{i=i_0}^{n-1}{\mathbb E}|e^{iX_i}-1|.
\]
Apply Lemma~\ref{lem:v11v73-increment}, then the reverse triangle
inequality.
\end{proof}

No independence is used.  The smallness in (73.8), rather than only
finiteness, protects the anchored denominator.

\subsection{A boundary-decay schedule}

Suppose \(B_i\) is a sum over a boundary set \(\partial\Lambda_i\) and
\[
                         {\mathbb E}|j_i(e)|
                      \le C_i h_i(d(e,K))                    \tag{73.10}
\]
for a fixed core \(K\).

\begin{corollary}[Summable boundary phase]
\label{cor:v11v73-boundary}
If
\[
 \sum_i|\partial\Lambda_i|C_i
       \sup_{e\in\partial\Lambda_i}h_i(d(e,K))<\infty,        \tag{73.11}
\]
then
\[
                         \sum_i{\mathbb E}\min\{2,|B_i|\}
                              <\infty.                         \tag{73.12}
\]
\end{corollary}

\begin{proof}
The triangle inequality bounds
\[
 {\mathbb E}|B_i|
 \le|\partial\Lambda_i|C_i
       \sup_{e\in\partial\Lambda_i}h_i(d(e,K)).
\]
Use \(\min\{2,|B_i|\}\le|B_i|\) and sum.
\end{proof}

Exponential current decay beats polynomial boundary growth.
Polynomial decay must beat both the boundary dimension and the renewal
frequency.

\subsection{Cutoff coboundaries}

A second localization can occur along the refinement direction.
Suppose, modulo \(2\pi\),
\[
                         X_i=F_{i+1}-F_i+R_i.                 \tag{73.13}
\]
Then
\[
                         \sum_{i=i_0}^{n-1}X_i
                 =F_n-F_{i_0}+\sum_{i=i_0}^{n-1}R_i
                              \pmod{2\pi}.                    \tag{73.14}
\]

\begin{proposition}[Renewal telescope]
\label{prop:v11v73-renewal}
If
\[
 \sum_i{\mathbb E}|e^{i(F_{i+1}-F_i)}-1|<\infty,\qquad
 \sum_i{\mathbb E}|e^{iR_i}-1|<\infty,                       \tag{73.15}
\]
then \(e^{i\Theta_n}\) is Cauchy in \(L^1\), hence converges in
\(L^1\) to a unit-modulus random variable.
\end{proposition}

\begin{proof}
Equation (73.13) and the factorization used in
Lemma~\ref{lem:v11v73-increment} bound the \(L^1\) distance between
successive exponentials by the two summands in (73.15).  Their series
is finite, so the phase sequence is \(L^1\)-Cauchy.  Its almost-surely
convergent subsequence has unit modulus, and \(L^1\) convergence
preserves that property.
\end{proof}

The expectation of the limiting unit phase can still be zero.
The anchored small-tail condition of
Theorem~\ref{thm:v11v73-anchor} is required for noncancellation.

\subsection{Harmonic and topological terms}

On a periodic complex, a closed cochain need not be exact.  A Hodge
decomposition contains a harmonic component representing nontrivial
cohomology.  Accordingly a phase decomposition can have the form
\[
                         \Theta_\Lambda
 =\langle{\bf1},d_0^*j\rangle
       +\Theta_{\rm harm}+\Theta_{\rm top}+R.                 \tag{73.16}
\]
The divergence term vanishes, but the harmonic and topological terms
remain.

\begin{proposition}[Sector phase mixture]
\label{prop:v11v73-sector}
Suppose the modulus measure decomposes into sectors \(s\) with
probabilities \(p_s\), and the residual phase in sector \(s\) has mean
\(z_s\).  Then
\[
                         z_\vartheta=\sum_sp_sz_s.            \tag{73.17}
\]
Even if \(|z_s|=1\) for every \(s\), the total denominator can vanish.
\end{proposition}

\begin{proof}
Equation (73.17) is the law of total expectation.  Two sectors with
equal probabilities and phases \(1\) and \(-1\) give zero.
\end{proof}

A sectorwise construction therefore also needs tightness of the sector
probabilities and a common phase arc or another noncancellation
mechanism.

\subsection{What anomaly cancellation supplies}

Let \(\alpha\) be the determinant-phase one-form on field space.  Local
perturbative anomaly cancellation states that its gauge-direction
curvature or local variation cancels after the Standard Model species
sum.  If one could strengthen this to
\[
                         \alpha=dF+\alpha_{\rm harm}+r         \tag{73.18}
\]
with a summable remainder \(r\), then pathwise phase increments would
have the renewal form (73.13).  The harmonic holonomy would remain a
global anomaly row.

\begin{warning}[Closed is not exact]
\label{warn:v11v73-closed}
Vanishing local curvature makes a connection locally flat.  It does
not make its holonomy trivial on a multiply connected field or gauge
quotient.  Therefore perturbative anomaly cancellation does not by
itself prove the coboundary decomposition needed in
Theorem~\ref{thm:v11v73-anchor}.
\end{warning}

\subsection{Boundary size is still extensive enough to hurt}

If boundary phase cumulants scale as \(|\partial\Lambda|\), the
denominator can decay like
\(\exp\{-c|\partial\Lambda|\}\).  This is subextensive compared with
volume but still tends to zero and can violate V71.  Uniform
noncancellation requires the total circular boundary increment to be
summable or confined to a finite retained sector.

\begin{ledger}[Coboundary phase acquisition]
\label{led:v11v73-data}
A successful phase-localization proof must provide:
\begin{enumerate}
\item an exact spatial divergence or cutoff-coboundary formula;
\item bounds for its boundary flux;
\item a summable circular regulator remainder;
\item explicit harmonic and topological phase variables;
\item sector probabilities and a noncancellation condition;
\item determinant-line holonomy control;
\item a tail smaller than the anchored phase mean.
\end{enumerate}
\end{ledger}

\begin{firewall}[Conditional mechanism]
\label{fire:v11v73-boundary}
No current calculation writes the full chiral SM determinant phase on
a fluctuating metric as (73.16) or (73.18).  Local anomaly cancellation
suggests where to search but does not prove phase localization,
boundary summability, or trivial global holonomy.
\end{firewall}

\begin{decision}[Next acquisition]
\label{dec:v11v73-next}
V74 will build a sector-conditioned projective measure theorem.  It
will combine tight sector probabilities, sectorwise cylinder limits,
and a uniform phase-arc condition into a global nonzero complex
functional, while showing how probability can escape through an
uncontrolled topological-sector tail.
\end{decision}

\begin{assessment}[Progress after V73]
\label{ass:v11v73-status}
A coboundary phase now has a precise payoff: its bulk cancels, and a
summable boundary-plus-remainder tail protects an existing phase
anchor without independence.  Harmonic holonomy and sector mixing are
proved to be separate possible sources of exact cancellation.
\end{assessment}

\subsection{Parent record}

This V73 note was formed by appending the preceding material to the
validated V72 source, whose record was
\[
\begin{array}{c|c}
\text{lines}&22621\\
\text{bytes}&717343\\
\text{SHA--256}&
\texttt{9FFDEFC9F95FDE3A34968713D9B1C2BDE394DF95CAA5A7BD79E4FAD87548BD07}.
\end{array}                                                     \tag{73.19}
\]

% =============================================================================
% END APPENDED ELEVENTH-SERIES V73 MATERIAL
% =============================================================================
% =============================================================================
% BEGIN APPENDED ELEVENTH-SERIES V74 MATERIAL
% =============================================================================

\section{Sector-conditioned projective limits and global phase
noncancellation}
\label{sec:v11v74}

A phase may be controlled within each topological sector and still
cancel after sectors are mixed.  Conversely, sectorwise convergence
does not give a global measure if probability escapes to sectors whose
labels diverge.  This section gives the needed tightness and common-
arc conditions.

\subsection{Discrete sector tightness}

Let \({\cal S}\) be a countable sector set and let
\(p_i\) be probability masses on \({\cal S}\).

\begin{definition}[Sector tightness]
\label{def:v11v74-tight}
The family is tight if for every \(\varepsilon>0\) there is a finite
\({\cal S}_\varepsilon\subset{\cal S}\) such that
\[
                         \sup_i\sum_{s\notin{\cal S}_\varepsilon}
                                      p_i(s)<\varepsilon.      \tag{74.1}
\]
\end{definition}

\begin{theorem}[Pointwise plus tight implies total variation]
\label{thm:v11v74-probability}
Suppose \(p_i(s)\to p(s)\) for every \(s\) and (74.1) holds.  Then
\(p\) is a probability mass and
\[
                              \sum_s|p_i(s)-p(s)|\longrightarrow0.
                                                                  \tag{74.2}
\]
\end{theorem}

\begin{proof}
On a finite set \({\cal S}_\varepsilon\), pointwise convergence is
uniform after summation.  Fatou's lemma and (74.1) show
\[
 \sum_{s\notin{\cal S}_\varepsilon}p(s)\le\varepsilon.
\]
Thus \(\sum_sp(s)\ge1-\varepsilon\), while Fatou gives
\(\sum_sp(s)\le1\); hence the total mass is one.  For large \(i\), the
finite-core \(l^1\) difference is below \(\varepsilon\), and both
tails have mass at most \(2\varepsilon\).  This proves (74.2) after
letting \(\varepsilon\downarrow0\).
\end{proof}

\begin{warning}[Escape example]
\label{warn:v11v74-escape}
On \({\cal S}=\mathbb N\), the probabilities \(p_i=\delta_i\) converge
pointwise to zero in every fixed sector, but their limiting pointwise
mass is zero.  No global probability limit exists on \(\mathbb N\).
This is the sector analogue of mass escaping to infinity.
\end{warning}

\subsection{Sectorwise complex cylinder functionals}

For each \(i,s\), let \(\nu_{i,s}\) be a positive conditional
probability and let \(u_{i,s}\) be a unit-modulus phase.  Define
\[
 \tau_{i,s}(F)
       ={\mathbb E}_{\nu_{i,s}}[u_{i,s}F],\qquad
 z_{i,s}=\tau_{i,s}(1),                                      \tag{74.3}
\]
and
\[
 N_i(F)=\sum_sp_i(s)\tau_{i,s}(F),\qquad
 Z_i=N_i(1)=\sum_sp_i(s)z_{i,s}.                             \tag{74.4}
\]
For bounded \(F\),
\[
                              |\tau_{i,s}(F)|\le\lVert F\rVert_\infty.
                                                                  \tag{74.5}
\]

\begin{theorem}[Sectorwise-to-global cylinder convergence]
\label{thm:v11v74-global}
Assume:
\begin{enumerate}
\item \(p_i\to p\) in \(l^1({\cal S})\);
\item for every fixed \(s\) and bounded cylinder \(F\),
\(\tau_{i,s}(F)\to\tau_s(F)\);
\item the cylinder identifications are compatible with the sector
labels.
\end{enumerate}
Then
\[
                         N_i(F)\longrightarrow
                         N_\infty(F):=\sum_sp(s)\tau_s(F),    \tag{74.6}
\]
and \(Z_i\to Z_\infty=\sum_sp(s)z_s\).
\end{theorem}

\begin{proof}
Fix \(\varepsilon>0\) and choose a finite sector set carrying all but
\(\varepsilon\) of \(p\) and, by \(l^1\) convergence, all but
\(2\varepsilon\) of \(p_i\) for large \(i\).  On that finite set,
sectorwise convergence permits passage to the limit term by term.
The two tails are bounded by (74.5), giving an error at most
\(3\varepsilon\lVert F\rVert_\infty\).  Let
\(\varepsilon\downarrow0\).  Take \(F=1\) for the denominator.
\end{proof}

\subsection{A common phase half-plane}

Sectorwise nonzero means can cancel.  A common oriented half-plane
prevents this.

\begin{theorem}[Sector phase alignment]
\label{thm:v11v74-alignment}
Suppose there are \(\vartheta_0\in\mathbb R\), \(\zeta>0\), and a
sector core \({\cal S}_0\) such that
\[
             \operatorname{Re}(e^{-i\vartheta_0}z_{i,s})
                              \ge\zeta
                   \quad(s\in{\cal S}_0),                    \tag{74.7}
\]
while
\[
                              \sum_{s\notin{\cal S}_0}p_i(s)
                                   \le\varepsilon.             \tag{74.8}
\]
Then
\[
                              |Z_i|
             \ge\zeta(1-\varepsilon)-\varepsilon.             \tag{74.9}
\]
In particular the global denominator is uniformly nonzero when the
right side is positive.
\end{theorem}

\begin{proof}
Rotate (74.4) by \(e^{-i\vartheta_0}\) and take the real part.
The core contributes at least \(\zeta(1-\varepsilon)\).  Since
\(|z_{i,s}|\le1\), the tail contributes at least
\(-\varepsilon\).  The modulus is at least a positive real part.
\end{proof}

\begin{corollary}[Global aligned sectors]
\label{cor:v11v74-all}
If (74.7) holds for every sector, then
\[
                              |Z_i|\ge\zeta                   \tag{74.10}
\]
without a tail estimate.
\end{corollary}

\begin{proof}
Set \(\varepsilon=0\).
\end{proof}

\begin{proposition}[Exact sector cancellation remains possible]
\label{prop:v11v74-cancel}
Two sectors with \(p_+=p_-=1/2\), \(z_+=1\), and \(z_-=-1\) have
\(Z=0\), although both sector denominators have unit modulus.
\end{proposition}

\begin{proof}
Insert the data into (74.4).
\end{proof}

Thus sectorwise phase control must include relative phase alignment.

\subsection{Adjacent total-variation estimate by sectors}

Let \(\tau_{i,s}\) now denote finite complex measures with
\(\lVert\tau_{i,s}\rVert_{\rm TV}\le1\), and define
\[
                              \sigma_i=\sum_sp_i(s)\tau_{i,s}. \tag{74.11}
\]

\begin{proposition}[Sectorwise adjacent bound]
\label{prop:v11v74-adjacent}
After identifying common sector and cylinder spaces,
\[
 \lVert\sigma_{i+1}-\sigma_i\rVert_{\rm TV}
 \le\lVert p_{i+1}-p_i\rVert_{\ell^1}
       +\sum_sp_i(s)
          \lVert\tau_{i+1,s}-\tau_{i,s}\rVert_{\rm TV}.       \tag{74.12}
\]
Therefore both series on the right being summable imply the
total-variation hypothesis (69.2).
\end{proposition}

\begin{proof}
Add and subtract
\(\sum_sp_i(s)\tau_{i+1,s}\).  The first difference has total variation
at most
\[
                         \sum_s|p_{i+1}(s)-p_i(s)|
                              \lVert\tau_{i+1,s}\rVert_{\rm TV},
\]
and the second is bounded by the weighted sum displayed in (74.12).
\end{proof}

If sector labels refine rather than remain fixed, a pushforward map
between sector sets must be applied before (74.12).  Splitting one
coarse sector among many fine sectors is part of the probability
increment, not a within-sector field fluctuation.

\subsection{Sectorwise projective consistency}

Let \(P_{i,j}\) project fine fields to a fixed cylinder \(j\).  A
sector label map
\[
                              q_{i,j}:{\cal S}_i\to{\cal S}_j \tag{74.13}
\]
must satisfy composition under refinement.  Exact projective
consistency requires both
\[
                         (q_{i+1,i})_*p_{i+1}=p_i             \tag{74.14}
\]
and conditional consistency after all fine sectors mapping to a coarse
sector are mixed with their conditional probabilities.

\begin{warning}[Conditionals do not project one by one]
\label{warn:v11v74-conditionals}
A fine topological sector can split a coarse sector.  Its individual
field marginal need not equal the coarse conditional measure.
Only the probability-weighted mixture over the fibre of
\(q_{i+1,i}\) is required to project correctly.
\end{warning}

This prevents an overly strong and generally false sectorwise
consistency demand.

\subsection{Topological phases}

For an integer charge \(Q=s\) and theta angle \(\theta\), a sector
factor \(e^{i\theta s}\) rotates successive sectors.  Even if all
within-sector determinant phases are aligned, the theta factors can
violate (74.7).  A symmetric broad charge distribution at
\(\theta=\pi\) can cause severe cancellation.

The required estimate is therefore on the characteristic function of
the sector charge,
\[
                              \sum_sp_i(s)e^{i\theta s},       \tag{74.15}
\]
possibly multiplied by the residual sector phase mean.  Tightness of
\(p_i\) alone does not bound (74.15) away from zero.

\subsection{Global normalized target}

\begin{theorem}[Sector-assembled complex functional]
\label{thm:v11v74-functional}
Under Theorem~\ref{thm:v11v74-global}, suppose
\[
                              |Z_\infty|\ge\zeta_\infty>0.    \tag{74.16}
\]
Then
\[
                              \Omega_i(F)=\frac{N_i(F)}{Z_i}
                  \longrightarrow
                              \Omega_\infty(F)
                              =\frac{N_\infty(F)}{Z_\infty}   \tag{74.17}
\]
for every bounded cylinder \(F\).  A uniform version of
Theorem~\ref{thm:v11v74-alignment} is a sufficient condition for
(74.16).
\end{theorem}

\begin{proof}
The numerators and denominators converge by
Theorem~\ref{thm:v11v74-global}.  The denominator is eventually bounded
away from zero by (74.16).  Apply the stable quotient estimate of V42.
\end{proof}

\subsection{SM--Einstein sector ledger}

Relevant sector data can include gauge bundles, instanton charge,
torons, spin structures, residual conjugacy strata, and gravitational
topology.  They need not all share one discrete label set, but their
joint finite-cutoff label must admit a tight projective description.

\begin{ledger}[Sector acquisition]
\label{led:v11v74-data}
The global complex continuum construction requires:
\begin{enumerate}
\item refinement-compatible sector maps;
\item tight sector probabilities or summable \(l^1\) increments;
\item sector-conditioned cylinder convergence;
\item within-sector phase anchors;
\item a common phase half-plane on a high-probability sector core;
\item control of theta and harmonic holonomy factors;
\item a nonzero assembled denominator;
\item compatibility with reflection and cluster decompositions.
\end{enumerate}
\end{ledger}

\begin{firewall}[No sector estimates acquired]
\label{fire:v11v74-boundary}
No current SM--Einstein calculation supplies tight probabilities for
all gauge and gravitational sectors, a common phase arc, or the
characteristic-function bound (74.15).  V74 proves the assembly theorem
and the failure modes; it does not establish the physical sector
inputs.
\end{firewall}

\begin{decision}[Mandatory V75 audit]
\label{dec:v11v74-next}
The next version is the compulsory YM/NS companion audit.  It will
inspect any new YM quadrature or sector result and any NS tail result.
V76 will then continue with the strongest missing global-measure
input identified by that audit.
\end{decision}

\begin{assessment}[Progress after V74]
\label{ass:v11v74-status}
Sectorwise convergence now assembles into a global complex cylinder
functional under tightness and a nonzero mixed denominator.  A common
phase half-plane prevents cancellation, while explicit examples show
why sectorwise nonzero phases and pointwise sector convergence are not
enough.
\end{assessment}

\subsection{Parent record}

This V74 note was formed by appending the preceding material to the
validated V73 source, whose record was
\[
\begin{array}{c|c}
\text{lines}&22937\\
\text{bytes}&727714\\
\text{SHA--256}&
\texttt{7DE9AE42CDBB96FCF0BADC7459834AF51D68DE43F972413AFB3980CF63681D9C}.
\end{array}                                                     \tag{74.18}
\]

% =============================================================================
% END APPENDED ELEVENTH-SERIES V74 MATERIAL
% =============================================================================
% =============================================================================
% BEGIN APPENDED ELEVENTH-SERIES V75 MATERIAL
% =============================================================================

\section{Companion audit at V75: no new sector or phase estimate}
\label{sec:v11v75}

The compulsory audit finds the same companion records as V70:
\[
\begin{array}{c|r|l}
\text{file}&\text{bytes}&\text{SHA--256}\\ \hline
\texttt{Renewal\_Yang\_Mills\_Mass\_Gap\_Research\_Notes\_V38.tex}
 &536904&
\texttt{ED5ACB149AF22C720DD516913B5A8AD3BFA15FE710C106F93FFE5E4DC3238390}\\
\texttt{Renewal\_NS\_Stability\_Research\_Notes\_V26.tex}
 &278283&
\texttt{82C887F8C2C69E4F28FAD7C3DC8DE5B9099CB5A4BF431EBA1FFF3E91935978AD}.
\end{array}                                                     \tag{75.1}
\]
No new companion theorem is imported.

The corrected YM V38 longitudinal Ward-limit result remains useful on
its declared positive endpoint, but it supplies neither topological
sector probabilities nor a chiral phase denominator.  The NS
rank-one kernel card likewise supplies no sector entropy estimate.

\begin{decision}[V76 acquisition]
\label{dec:v11v75-next}
V76 will derive a sector-tightness theorem from a lower action cost per
topological charge and an upper entropy growth per sector.  It will
give an explicit exponential tail and a small-angle characteristic-
function bound.  Application to the SM--Einstein regulator will remain
conditional on a genuine lattice charge, an admissibility event, and
uniform energy--entropy constants.
\end{decision}

\begin{firewall}[Audit boundary]
\label{fire:v11v75-boundary}
The unchanged companions do not close the topological, phase,
reflection-positivity, clustering, or Einstein-sector rows.  The YM
complex external strip is not an internal sector estimate, and the NS
decay constants do not count gauge or gravitational topologies.
\end{firewall}

\begin{assessment}[Progress after V75]
\label{ass:v11v75-status}
The required audit is complete and confirms that the next unresolved
input is sector probability rather than another local response
identity.
\end{assessment}

\subsection{Parent record}

This V75 note was formed by appending the preceding material to the
validated V74 source, whose record was
\[
\begin{array}{c|c}
\text{lines}&23274\\
\text{bytes}&739495\\
\text{SHA--256}&
\texttt{457979837E72CC6E132944C1B62A0DF393E3AB37A329307BECBDA78544C4F6A3}.
\end{array}                                                     \tag{75.2}
\]

% =============================================================================
% END APPENDED ELEVENTH-SERIES V75 MATERIAL
% =============================================================================
% =============================================================================
% BEGIN APPENDED ELEVENTH-SERIES V76 MATERIAL
% =============================================================================

\section{Energy--entropy confinement of topological charge}
\label{sec:v11v76}

V74 reduced global assembly to sector tightness and phase control.
This version supplies a uniform sufficient condition for the first
of those inputs.  The argument applies to the positive amplitude
measure.  Complex determinant and theta phases are inserted only
after the probability estimate has been proved.

Let \(I_i\) be a finite or countable charge set contained in
\(\mathbb Z\), let \(W_i(q)\ge0\) be the unnormalised amplitude weight
of sector \(q\), and put
\[
                 {\cal Z}_i^{\rm abs}=\sum_{q\in I_i}W_i(q),
        \qquad
                 p_i(q)=\frac{W_i(q)}{{\cal Z}_i^{\rm abs}}.  \tag{76.1}
\]

\begin{theorem}[Uniform sector tail from an energy--entropy gap]
\label{thm:v11v76-tail}
Suppose that constants \(A,z_0,\delta>0\), independent of \(i\), obey
\[
 W_i(q)\le A e^{-\delta |q|}
 \quad(q\in I_i),\qquad
 {\cal Z}_i^{\rm abs}\ge z_0.                                \tag{76.2}
\]
Extend \(p_i\) by zero outside \(I_i\).  Then, for every integer
\(Q\ge1\),
\[
 p_i\{|q|\ge Q\}
 \le
 \frac{2A}{z_0(1-e^{-\delta})}\,e^{-\delta Q}.                \tag{76.3}
\]
In particular the family \((p_i)\) is uniformly tight on
\(\mathbb Z\).  Moreover, for \(0\le\lambda<\delta\),
\[
 \mathbb E_i e^{\lambda |q|}
 \le {A\over z_0}
       {1+e^{-(\delta-\lambda)}
        \over 1-e^{-(\delta-\lambda)}}.                       \tag{76.4}
\]
\end{theorem}

\begin{proof}
By (76.1)--(76.2),
\[
 p_i\{|q|\ge Q\}
 \le {A\over z_0}
       \sum_{\substack{n\in\mathbb Z\\ |n|\ge Q}}
             e^{-\delta |n|}
 = {2A\over z_0}\,{e^{-\delta Q}\over1-e^{-\delta}},
\]
which proves (76.3).  The same comparison gives
\[
 \mathbb E_i e^{\lambda |q|}
 \le {A\over z_0}\sum_{n\in\mathbb Z}
             e^{-(\delta-\lambda)|n|}
 = {A\over z_0}
       {1+e^{-(\delta-\lambda)}
        \over1-e^{-(\delta-\lambda)}}.
\]
No limit interchange is used.
\end{proof}

\begin{corollary}[Action cost minus sector entropy]
\label{cor:v11v76-energy-entropy}
Let
\[
 W_i(q)=\int_{X_i(q)}e^{-S_i(x)}\,dm_i(x).                    \tag{76.5}
\]
Assume uniformly in \(i\) and \(q\) that
\[
 S_i(x)\ge \tau |q|-C\quad(x\in X_i(q)),\qquad
 m_i(X_i(q))\le A_0e^{b|q|},                                 \tag{76.6}
\]
where \(\tau>b\).  Suppose also that a measurable
\(B_i\subset X_i(0)\) satisfies
\[
 m_i(B_i)\ge m_0>0,\qquad S_i(x)\le C_0\quad(x\in B_i).        \tag{76.7}
\]
Then Theorem~\ref{thm:v11v76-tail} holds with
\[
 A=A_0e^C,\qquad \delta=\tau-b,\qquad z_0=m_0e^{-C_0}.        \tag{76.8}
\]
\end{corollary}

\begin{proof}
Equations (76.5)--(76.6) give
\[
 W_i(q)\le e^{C-\tau|q|}m_i(X_i(q))
          \le A_0e^C e^{-(\tau-b)|q|}.
\]
The integral over \(B_i\) gives
\({\cal Z}_i^{\rm abs}\ge W_i(0)\ge m_0e^{-C_0}\).
\end{proof}

\begin{remark}[Weighted reference measures]
\label{rem:v11v76-reference}
For noncompact field coordinates, the raw volume
\(m_i(X_i(q))\) in (76.6) may be infinite.  The same proof applies
after writing \(S_i=S_i^{(0)}+V_i\), replacing the volume by
\[
 H_i(q)=\int_{X_i(q)}e^{-S_i^{(0)}}\,dm_i,
\]
and assuming \(H_i(q)\le A_0e^{b|q|}\) together with
\(V_i\ge\tau|q|-C\).  Thus the substantive input is a sector
partition-function entropy bound, not finiteness of coordinate
Lebesgue volume.
\end{remark}

\subsection{Quantitative theta window}

Let
\[
                         \phi_i(\theta)
           =\sum_qp_i(q)e^{i\theta q}.                        \tag{76.9}
\]
Set \(r=e^{-\delta}\) and \(B=A/z_0\).  Direct summation gives
\[
 \mathbb E_i|q|
 \le {2Br\over(1-r)^2}=:M_1,\qquad
 \mathbb E_iq^2
 \le {2Br(1+r)\over(1-r)^3}=:M_2.                            \tag{76.10}
\]

\begin{proposition}[Noncancellation near zero theta]
\label{prop:v11v76-theta}
Under Theorem~\ref{thm:v11v76-tail},
\[
 |\phi_i(\theta)|
 \ge 1-|\theta|M_1,\qquad
 |\phi_i(\theta)|
 \ge 1-\frac{\theta^2M_2}{2}.                                \tag{76.11}
\]
Consequently the topological theta factor has a uniform nonzero
expectation whenever either right-hand side is positive.
\end{proposition}

\begin{proof}
The inequalities
\[
 |e^{it}-1|\le |t|,\qquad
 \cos t\ge1-\frac{t^2}{2}
\]
and (76.10) imply
\[
 |\phi_i(\theta)-1|
 \le|\theta|\mathbb E_i|q|\le|\theta|M_1
\]
and
\[
 \operatorname{Re}\phi_i(\theta)
 =\mathbb E_i\cos(\theta q)
 \ge1-\frac{\theta^2}{2}\mathbb E_iq^2
 \ge1-\frac{\theta^2M_2}{2}.
\]
Since \(|z|\ge\operatorname{Re}z\), (76.11) follows.
\end{proof}

\begin{warning}[The window is necessarily local]
\label{warn:v11v76-large-theta}
Exponential tails do not prevent a characteristic-function zero at a
larger angle.  For example, the probability assigning mass \(1/2\)
to each of \(q=0,1\) has finite exponential moments of every order but
\(\phi(\pi)=0\).  A prescribed physical theta angle outside the
window (76.11) requires additional sector structure.
\end{warning}

\subsection{Insertion of residual determinant phases}

Suppose the sector-conditioned residual phase mean is \(a_i(q)\) with
\(|a_i(q)|\le1\), and define
\[
 D_i(\theta)=\sum_qp_i(q)a_i(q)e^{i\theta q}.                 \tag{76.12}
\]

\begin{proposition}[Perturbation of the theta characteristic function]
\label{prop:v11v76-residual}
If
\[
                    \sum_qp_i(q)|a_i(q)-1|\le\varepsilon_i,  \tag{76.13}
\]
then
\[
 |D_i(\theta)|\ge|\phi_i(\theta)|-\varepsilon_i.              \tag{76.14}
\]
Thus (76.11) gives a nonzero full sector denominator whenever its
chosen lower bound exceeds \(\varepsilon_i\).
\end{proposition}

\begin{proof}
The difference between (76.12) and (76.9) is at most the left-hand
side of (76.13), by the triangle inequality.
\end{proof}

Condition (76.13) is stronger than sectorwise nonvanishing.  It asks
for average alignment with a single trivialisation and therefore
interfaces directly with the determinant-line problem of V69 and the
boundary localisation problem of V73.

\begin{ledger}[Concrete acquisition required]
\label{led:v11v76-acquisition}
To apply Corollary~\ref{cor:v11v76-energy-entropy} uniformly to an
SM--Einstein refinement, one must construct:
\begin{enumerate}
\item an integer or finitely generated charge on every admissible
      cutoff configuration;
\item refinement-compatible sector maps;
\item a uniform lower action cost \(\tau|q|-C\);
\item a sector entropy exponent \(b<\tau\);
\item a cutoff-independent comparison set giving (76.7);
\item a residual phase estimate of the form (76.13).
\end{enumerate}
Gauge admissibility, gravitational topology, and chiral
determinant-line holonomy enter different rows and cannot be inferred
from the abstract tail theorem.
\end{ledger}

\begin{firewall}[What V76 does not establish]
\label{fire:v11v76-boundary}
No attached manuscript currently proves the uniform constants in
(76.6)--(76.7) for the coupled theory.  The Yang--Mills local chart
controls a neighbourhood of a reference field but does not by itself
define global topological charge or bound the number of sectors.
The Einstein action also has conformal and topology-sum difficulties
that are not covered by a positive action-cost hypothesis.
\end{firewall}

\begin{decision}[V77 acquisition]
\label{dec:v11v76-next}
The next version will formulate a noncircular admissibility and
refinement criterion under which a lattice topological charge is
locally constant, stable under interpolation, and compatible with
sector maps.  This is the first acquisition row in
Ledger~\ref{led:v11v76-acquisition}.
\end{decision}

\begin{assessment}[Progress after V76]
\label{ass:v11v76-status}
A uniform energy advantage over sector entropy now yields explicit
charge tails, exponential moments, and a certified theta interval.
The remaining physical work is cleanly separated into construction of
the charge, proof of the uniform energy--entropy constants, and
alignment of the residual chiral phase.
\end{assessment}

\subsection{Parent record}

This V76 note was formed by appending the preceding material to the
validated V75 source, whose record was
\[
\begin{array}{c|c}
\text{lines}&23341\\
\text{bytes}&742204\\
\text{SHA--256}&
\texttt{7847960B90DF71A6B8FC104B888020866570CDE33A0CAFD5B00F58E63BD91899}.
\end{array}                                                     \tag{76.15}
\]

% =============================================================================
% END APPENDED ELEVENTH-SERIES V76 MATERIAL
% =============================================================================
% =============================================================================
% BEGIN APPENDED ELEVENTH-SERIES V77 MATERIAL
% =============================================================================

\section{Admissible index sectors and refinement stability}
\label{sec:v11v77}

The first acquisition row of V76 is a charge that survives refinement.
A useful construction must not declare sectors compatible by
definition.  This version gives a finite-dimensional certificate:
the charge is the inertia index of a Hermitian carrier, admissibility
is a spectral gap, and refinement compatibility follows from an
explicit gapped homotopy.

Let \(X_i\) be a cutoff configuration space and let
\[
                    H_i:X_i\longrightarrow {\rm Herm}(N_i)    \tag{77.1}
\]
be norm-continuous.  Define the admissible set
\[
 {\cal A}_i=\{x\in X_i:0\notin\sigma(H_i(x))\}.                \tag{77.2}
\]
Choose an integer \(r_i\) and set
\[
                    Q_i(x)=n_-(H_i(x))-r_i,\qquad x\in{\cal A}_i,
                                                                    \tag{77.3}
\]
where \(n_-\) counts negative eigenvalues with multiplicity.

\begin{theorem}[Spectral definition of discrete sectors]
\label{thm:v11v77-index}
The map \(Q_i:{\cal A}_i\to\mathbb Z\) is locally constant.  It is
constant along every continuous admissible path.  Consequently each
set
\[
                    {\cal A}_i(q)=\{x\in{\cal A}_i:Q_i(x)=q\} \tag{77.4}
\]
is both open and closed relative to \({\cal A}_i\).
\end{theorem}

\begin{proof}
Fix \(x_0\in{\cal A}_i\) and put
\(g_0=\min_{\lambda\in\sigma(H_i(x_0))}|\lambda|>0\).
Norm continuity gives a neighbourhood \(U\) of \(x_0\) on which
\(\|H_i(x)-H_i(x_0)\|<g_0/2\).  Weyl's eigenvalue perturbation
inequality shows that no eigenvalue can cross zero on \(U\), so the
negative inertia and \(Q_i\) are constant there.  A locally constant
integer-valued function is constant on the connected image of a path.
The final assertion follows because each fibre of a locally constant
map and its complement are open.
\end{proof}

\begin{proposition}[Quantitative perturbation certificate]
\label{prop:v11v77-perturb}
Let \(x,y\in X_i\).  If
\[
 g_i(x):=\min_{\lambda\in\sigma(H_i(x))}|\lambda|>0,\qquad
 \|H_i(y)-H_i(x)\|<g_i(x),                                   \tag{77.5}
\]
then \(y\in{\cal A}_i\) and \(Q_i(y)=Q_i(x)\).  More generally, if a
path \(x(t)\) admits points
\(0=t_0<\cdots<t_m=1\) such that
\[
 \|H_i(x(t))-H_i(x(t_j))\|<g_i(x(t_j))
 \quad(t_j\le t\le t_{j+1}),                                 \tag{77.6}
\]
then the path is admissible and its charge is constant.
\end{proposition}

\begin{proof}
For (77.5), Weyl's inequality preserves the sign of every ordered
eigenvalue.  Apply the same argument on each interval in (77.6) and
induct over \(j\).
\end{proof}

\subsection{Exact refinement certificate}

Let \(E_i:{\cal A}_i\to X_{i+1}\) be a chosen refinement map.

\begin{theorem}[Gapped refinement homotopy]
\label{thm:v11v77-refine}
Suppose that, for every \(x\in{\cal A}_i\), there is a continuous
Hermitian path \(K_{i,x}(t)\), \(0\le t\le1\), such that
\[
 \begin{split}
 K_{i,x}(0)&=H_{i+1}(E_ix),\\
 K_{i,x}(1)&=H_i(x)\oplus J_i,\\
 0&\notin\sigma(K_{i,x}(t))\quad(0\le t\le1),                 \tag{77.7}\\
 n_-(J_i)&=r_{i+1}-r_i.
 \end{split}
\]
Then \(E_ix\in{\cal A}_{i+1}\) and
\[
                              Q_{i+1}(E_ix)=Q_i(x).            \tag{77.8}
\]
Thus \(E_i\) induces the identity map on the integer charge labels.
\end{theorem}

\begin{proof}
Negative inertia is constant along the gapped path, so
\[
 \begin{split}
 Q_{i+1}(E_ix)
 &=n_-(H_{i+1}(E_ix))-r_{i+1}\\
 &=n_-(H_i(x))+n_-(J_i)-r_{i+1}\\
 &=n_-(H_i(x))-r_i=Q_i(x).
 \end{split}
\]
\end{proof}

The condition in (77.7) is checkable without knowing the sector label.
For example, choose a reference block diagonal path \(K^0_{i,x}(t)\)
whose gap is at least \(g>0\).  If
\[
                 \sup_{x,t}\|K_{i,x}(t)-K^0_{i,x}(t)\|<g,     \tag{77.9}
\]
then Weyl's inequality proves the required gap.  Interval matrix
bounds can certify (77.9) on a compact local chart.

\begin{corollary}[Composition over many refinements]
\label{cor:v11v77-composition}
If Theorem~\ref{thm:v11v77-refine} holds at every adjacent pair, then
for \(j>i\),
\[
 Q_j(E_{j-1}\cdots E_i x)=Q_i(x).                             \tag{77.10}
\]
\end{corollary}

\begin{proof}
Iterate (77.8).
\end{proof}

\subsection{Approximate coarse projection}

Exact embeddings alone do not control a general fine field.  Let
\(P_i:X_{i+1}\to X_i\) be a coarse projection and let
\({\cal G}_{i+1}\subset{\cal A}_{i+1}\) be a good set on which
\(P_ix\in{\cal A}_i\).

\begin{proposition}[Projection stability from a gapped comparison]
\label{prop:v11v77-project}
Suppose that for each \(x\in{\cal G}_{i+1}\) there is a gapped
Hermitian homotopy from \(H_{i+1}(x)\) to
\(H_i(P_ix)\oplus J_i\), with the same reference block as in
Theorem~\ref{thm:v11v77-refine}.  Then
\[
                             Q_{i+1}(x)=Q_i(P_ix)
                 \quad(x\in{\cal G}_{i+1}).                  \tag{77.11}
\]
If \(\mu_{i+1}({\cal G}_{i+1}^c)\le\varepsilon_i\), then the
pushforward charge laws satisfy
\[
 \big\|(Q_{i+1})_\#\mu_{i+1}
       -(Q_i\circ P_i)_\#\mu_{i+1}\big\|_{\rm TV}
 \le\varepsilon_i.                                           \tag{77.12}
\]
\end{proposition}

\begin{proof}
The inertia calculation proving (77.8) gives (77.11).
Couple the two charge random variables using the same \(x\).  They
differ only on \({\cal G}_{i+1}^c\), so the coupling inequality gives
(77.12).
\end{proof}

\begin{warning}[Admissibility loss is a real defect]
\label{warn:v11v77-admissibility}
A continuous interpolation in the full configuration space proves
nothing if it crosses the zero-spectrum wall.  Likewise,
\(\varepsilon_i\to0\) alone does not give a projective charge law;
for a telescoping continuum construction one normally needs
\(\sum_i\varepsilon_i<\infty\), together with the ordinary adjacent
measure defect.  Neither requirement follows from local constancy.
\end{warning}

\subsection{Relation to gauge and chiral topology}

For a lattice gauge field, \(H_i\) may be a Hermitian Wilson--Dirac
kernel and (77.3) an overlap-type index after the conventional
normalisation.  Plaquette admissibility is useful only after it has
been shown to imply the spectral gap in (77.2).  For a chiral
determinant, the same gap can define an index sector, but it does not
trivialise the determinant line inside that sector.

\begin{ledger}[Certification chain]
\label{led:v11v77-chain}
A noncircular regulator proof should provide, in this order:
\begin{enumerate}
\item a concrete finite matrix \(H_i(x)\);
\item a computable lower bound for its spectral gap on an admissible
      set;
\item the normalisation \(r_i\) and agreement with the intended
      continuum charge on smooth sampled fields;
\item a gapped adjacent-refinement homotopy;
\item a summable probability for leaving the good projection set;
\item only then, the energy--entropy estimate sector by sector.
\end{enumerate}
\end{ledger}

\begin{firewall}[Gravitational sectors]
\label{fire:v11v77-gravity}
The inertia construction does not classify changes of manifold
topology, spin structure, or all gravitational instanton data.
Those labels require their own refinement category.  A Dirac index on
a fixed spin manifold supplies one discrete invariant, not a sum over
geometries.
\end{firewall}

\begin{decision}[V78 acquisition]
\label{dec:v11v77-next}
V78 will transfer sector tails and charge compatibility from a
positive reference law to an interacting amplitude law.  The target
is an explicit Radon--Nikodym comparison and an adjacent refinement
bound that keeps the bad-admissibility probability and the measure
defect separate.
\end{decision}

\begin{assessment}[Progress after V77]
\label{ass:v11v77-status}
Topological sector compatibility has been reduced to a spectral gap
and a gapped refinement homotopy.  The resulting charge is exactly
integer-valued and stable, while failures of admissibility enter as a
quantified probability defect rather than an assumed equality.
\end{assessment}

\subsection{Parent record}

This V77 note was formed by appending the preceding material to the
validated V76 source, whose record was
\[
\begin{array}{c|c}
\text{lines}&23607\\
\text{bytes}&750857\\
\text{SHA--256}&
\texttt{7C4EC49F08171B0039BC67B5E063C4E99624FE1C52BA1A3D0EC66E9354ED2833}.
\end{array}                                                     \tag{77.13}
\]

% =============================================================================
% END APPENDED ELEVENTH-SERIES V77 MATERIAL
% =============================================================================
% =============================================================================
% BEGIN APPENDED ELEVENTH-SERIES V78 MATERIAL
% =============================================================================

\section{Transfer of sector control to an interacting amplitude law}
\label{sec:v11v78}

V76 controls a positive sector law once its weights have an
energy--entropy gap, while V77 identifies the admissibility defect on
which charges can change.  In practice these estimates are often
proved first under a Gaussian or product reference law.  This version
records exactly what survives positive reweighting.

Let \(\nu_i\) be a reference probability and let
\[
 \mu_i(dx)=f_i(x)\nu_i(dx),\qquad
 f_i={g_i\over Z_i},\qquad
 Z_i=\int g_i\,d\nu_i>0,                                     \tag{78.1}
\]
where \(g_i\ge0\).  Fix \(r\in(1,\infty]\), and write
\(r'=r/(r-1)\), with \(r'=1\) when \(r=\infty\).

\begin{theorem}[\(L^r\) change-of-measure transfer]
\label{thm:v11v78-holder}
Suppose
\[
                         \|g_i\|_{L^r(\nu_i)}\le G,\qquad
                         Z_i\ge z>0                           \tag{78.2}
\]
uniformly in \(i\).  Then, for every measurable event \(E\),
\[
             \mu_i(E)\le {G\over z}\,
                    \nu_i(E)^{1/r'}.                         \tag{78.3}
\]
If a reference charge \(Q_i\) satisfies
\[
             \nu_i\{|Q_i|\ge n\}\le C e^{-\delta n}
                    \quad(n\ge1),                            \tag{78.4}
\]
then
\[
             \mu_i\{|Q_i|\ge n\}
             \le {G\over z}C^{1/r'}
                    e^{-\delta n/r'}.                        \tag{78.5}
\]
\end{theorem}

\begin{proof}
Holder's inequality and (78.1)--(78.2) give
\[
 \mu_i(E)
 ={1\over Z_i}\int_Eg_i\,d\nu_i
 \le {G\over z}\nu_i(E)^{1/r'}.
\]
Insert (78.4).
\end{proof}

\begin{corollary}[Exponential moments after interaction]
\label{cor:v11v78-moment}
Suppose, more strongly, that
\[
                  \sup_i\mathbb E_{\nu_i}e^{a|Q_i|}<\infty
                                                                  \tag{78.6}
\]
for every \(a<\delta\).  Then
\[
                  \sup_i\mathbb E_{\mu_i}e^{\lambda|Q_i|}
                  <\infty
            \quad\hbox{for every }\lambda<\delta/r'.          \tag{78.7}
\]
\end{corollary}

\begin{proof}
For finite \(r\), Holder's inequality gives
\[
 \mathbb E_{\mu_i}e^{\lambda|Q_i|}
 \le {G\over z}
       \left(\mathbb E_{\nu_i}e^{\lambda r'|Q_i|}\right)^{1/r'}.
\]
Choose \(\lambda r'<\delta\).  The \(r=\infty\) case follows by direct
domination.
\end{proof}

The exponent loss \(1/r'\) cannot generally be removed.  Reweighting
can concentrate the target on a rare reference event up to the exact
Holder power.

\begin{corollary}[Admissibility defect transfer]
\label{cor:v11v78-admissible}
If \({\cal G}_i\) is a good admissibility event and
\(\nu_i({\cal G}_i^c)\le\varepsilon_i\), then
\[
               \mu_i({\cal G}_i^c)
               \le {G\over z}\varepsilon_i^{1/r'}.            \tag{78.8}
\]
Thus summability under the interacting law follows from
\[
                            \sum_i\varepsilon_i^{1/r'}<\infty, \tag{78.9}
\]
which is strictly stronger than \(\sum_i\varepsilon_i<\infty\) when
\(r<\infty\).
\end{corollary}

\subsection{Adjacent refinement of the reweighted laws}

Let \(P_i:X_{i+1}\to X_i\).  First suppose the references are exactly
projective:
\[
                              (P_i)_\#\nu_{i+1}=\nu_i.          \tag{78.10}
\]
Let
\[
 \bar f_i(y)=
 \mathbb E_{\nu_{i+1}}\!\left[f_{i+1}\mid P_i=y\right].       \tag{78.11}
\]

\begin{theorem}[Conditional-density refinement identity]
\label{thm:v11v78-martingale}
Under (78.10),
\[
                       (P_i)_\#\mu_{i+1}(dy)
                       =\bar f_i(y)\nu_i(dy),                 \tag{78.12}
\]
and, with
\[
 d_{\rm b}(\alpha,\beta)
 :=\sup_{\|F\|_\infty\le1}
       \left|\int F\,d\alpha-\int F\,d\beta\right|,           \tag{78.13}
\]
one has the exact identity
\[
 d_{\rm b}\big((P_i)_\#\mu_{i+1},\mu_i\big)
                    =\|\bar f_i-f_i\|_{L^1(\nu_i)}.           \tag{78.14}
\]
Consequently
\[
          \sum_i\|\bar f_i-f_i\|_{L^1(\nu_i)}<\infty          \tag{78.15}
\]
is a sufficient adjacent projective-Cauchy condition for all bounded
coarse observables.
\end{theorem}

\begin{proof}
For bounded \(F\),
\[
 \int F\,d(P_i)_\#\mu_{i+1}
 =\int F(P_ix)f_{i+1}(x)\,d\nu_{i+1}(x)
 =\int F(y)\bar f_i(y)\,d\nu_i(y),
\]
which proves (78.12).  The dual characterisation of the \(L^1\) norm,
with \(F\) equal to the measurable sign of
\(\bar f_i-f_i\), proves (78.14).  Summing (78.14) proves the Cauchy
claim.
\end{proof}

For approximate reference consistency, put
\(\rho_i=(P_i)_\#\nu_{i+1}\), choose a common dominating measure
\(\lambda_i\), and write
\[
                   \rho_i=a_i\lambda_i,\qquad
                   \nu_i=b_i\lambda_i.                       \tag{78.16}
\]
Let \(\bar f_i\) be the conditional density in (78.11), now defined
\(\rho_i\)-almost everywhere and extended by zero where necessary.

\begin{proposition}[Separate reference and interaction defects]
\label{prop:v11v78-two-defects}
If \(0\le\bar f_i\le K\), then
\[
 \begin{split}
 d_{\rm b}\big((P_i)_\#\mu_{i+1},\mu_i\big)
 &=\int|\bar f_i a_i-f_i b_i|\,d\lambda_i\\
 &\le K\int|a_i-b_i|\,d\lambda_i
       +\int|\bar f_i-f_i|b_i\,d\lambda_i.                   \tag{78.17}
 \end{split}
\]
\end{proposition}

\begin{proof}
The first equality is the dual \(L^1\) formula.  Pointwise,
\[
 |\bar f_i a_i-f_i b_i|
 \le\bar f_i|a_i-b_i|+b_i|\bar f_i-f_i|.
\]
Integrate and use the bound on \(\bar f_i\).
\end{proof}

The first term in (78.17) is a reference-measure defect.  The second
is an interaction or counterterm defect.  Treating their sum as a
single numerical residual would obscure which construction has
failed.

\subsection{Charge-law and theta convergence}

Assume the good set \({\cal G}_{i+1}\) in
Proposition~\ref{prop:v11v77-project} satisfies
\[
              Q_{i+1}(x)=Q_i(P_ix)
                        \quad(x\in{\cal G}_{i+1}).            \tag{78.18}
\]
Put
\[
 \eta_i=d_{\rm b}\big((P_i)_\#\mu_{i+1},\mu_i\big),\qquad
 \beta_i=\mu_{i+1}({\cal G}_{i+1}^c).                         \tag{78.19}
\]

\begin{theorem}[Adjacent characteristic-function bound]
\label{thm:v11v78-characteristic}
For every real \(\theta\),
\[
 \left|
 \mathbb E_{\mu_{i+1}}e^{i\theta Q_{i+1}}
 -\mathbb E_{\mu_i}e^{i\theta Q_i}
 \right|
 \le 2\beta_i+\eta_i.                                        \tag{78.20}
\]
If \(\sum_i(\beta_i+\eta_i)<\infty\), the characteristic functions
converge uniformly in \(\theta\in\mathbb R\).
\end{theorem}

\begin{proof}
On \({\cal G}_{i+1}\), the two phases in (78.18) agree.  On its
complement their difference has modulus at most two.  Therefore
\[
 \left|
 \mathbb E_{\mu_{i+1}}
   \left(e^{i\theta Q_{i+1}}-e^{i\theta Q_i\circ P_i}\right)
 \right|\le2\beta_i.
\]
The remaining difference is at most \(\eta_i\) by (78.13), applied to
the bounded coarse function \(e^{i\theta Q_i}\).  This proves
(78.20); summation gives uniform Cauchy convergence.
\end{proof}

Combining (78.8), (78.17), and (78.20) gives a concrete route from
reference estimates to a limiting theta characteristic function.
Nonvanishing still requires an anchor such as V76 at one scale plus a
total adjacent defect smaller than the anchor.

\begin{firewall}[Positive amplitude only]
\label{fire:v11v78-boundary}
The density \(g_i\) in (78.1) is nonnegative.  A complex chiral
determinant cannot be used as this Radon--Nikodym density before its
absolute amplitude and phase have been separated.  Bounds on
\(\|g_i\|_{L^r}\) and \(Z_i\) are additional uniform analytic inputs;
neither follows from pointwise convergence of regulated actions.
\end{firewall}

\begin{decision}[V79 acquisition]
\label{dec:v11v78-next}
V79 will combine a finite-scale phase anchor with the summable
charge, measure, and residual determinant-phase increments.  The aim
is a single noncancellation theorem whose hypotheses are individually
checkable and whose conclusion is stable under the continuum limit.
\end{decision}

\begin{assessment}[Progress after V78]
\label{ass:v11v78-status}
Reference charge tails now transfer to interacting positive amplitude
laws with the sharp Holder exponent loss.  Exact and approximate
reference refinement defects have been separated from interaction
defects, and all three enter a uniform adjacent theta estimate.
\end{assessment}

\subsection{Parent record}

This V78 note was formed by appending the preceding material to the
validated V77 source, whose record was
\[
\begin{array}{c|c}
\text{lines}&23848\\
\text{bytes}&759598\\
\text{SHA--256}&
\texttt{A0E84312A28DD04A878DF8A18C63FBA5B1997B9A385BBD745E7924EB60B21827}.
\end{array}                                                     \tag{78.21}
\]

% =============================================================================
% END APPENDED ELEVENTH-SERIES V78 MATERIAL
% =============================================================================
% =============================================================================
% BEGIN APPENDED ELEVENTH-SERIES V79 MATERIAL
% =============================================================================

\section{Anchored noncancellation under summable refinement defects}
\label{sec:v11v79}

The preceding versions supply separate estimates for sector tails,
charge stability, and positive amplitude refinement.  They become a
continuum complex measure only if their combined phase denominator
cannot drift through zero.  The correct global statement is an
anchored telescoping theorem.

Let \(\mu_i\) be positive amplitude probabilities on \(X_i\), let
\(P_i:X_{i+1}\to X_i\), and let \(u_i:X_i\to\mathbb C\) satisfy
\[
                                  |u_i|\le1.                  \tag{79.1}
\]
The complex denominator is
\[
                                  D_i=\int u_i\,d\mu_i.        \tag{79.2}
\]
Define the amplitude refinement defect
\[
 \eta_i=d_{\rm b}\big((P_i)_\#\mu_{i+1},\mu_i\big)            \tag{79.3}
\]
using (78.13), and the transported phase defect
\[
 \rho_i=\int
        |u_{i+1}(x)-u_i(P_ix)|\,d\mu_{i+1}(x).                \tag{79.4}
\]

\begin{theorem}[Anchored denominator theorem]
\label{thm:v11v79-anchor}
For every \(i\),
\[
                                  |D_{i+1}-D_i|
                                  \le\rho_i+\eta_i.            \tag{79.5}
\]
Fix \(i_0\), put
\[
                        R_{i_0}=\sum_{i\ge i_0}(\rho_i+\eta_i),
                                                                    \tag{79.6}
\]
and suppose
\[
                             R_{i_0}<|D_{i_0}|.               \tag{79.7}
\]
Then \(D_i\) converges to a limit \(D_\infty\), and
\[
 |D_j|\ge |D_{i_0}|-\sum_{i=i_0}^{j-1}(\rho_i+\eta_i)>0,
 \qquad
 |D_\infty|\ge|D_{i_0}|-R_{i_0}>0.                           \tag{79.8}
\]
\end{theorem}

\begin{proof}
Insert the intermediate expectation under \(\mu_{i+1}\):
\[
 \begin{split}
 |D_{i+1}-D_i|
 &\le\left|\int(u_{i+1}-u_i\circ P_i)\,d\mu_{i+1}\right|\\
 &\quad+
 \left|\int u_i\,d(P_i)_\#\mu_{i+1}-\int u_i\,d\mu_i\right|\\
 &\le\rho_i+\eta_i.
 \end{split}
\]
The series in (79.6) makes \((D_i)\) Cauchy.  The reverse triangle
inequality and telescoping give (79.8).
\end{proof}

\subsection{Charge and residual determinant phase}

Suppose, in a chosen determinant-line trivialisation,
\[
                         u_i(x)=e^{i\theta Q_i(x)}a_i(x),
                  \qquad |a_i(x)|\le1.                       \tag{79.9}
\]
Let \({\cal G}_{i+1}\) be the charge-compatible good set of V77 and
put
\[
 \beta_i=\mu_{i+1}({\cal G}_{i+1}^c),\qquad
 \alpha_i=\int|a_{i+1}(x)-a_i(P_ix)|\,d\mu_{i+1}(x).          \tag{79.10}
\]

\begin{proposition}[Resolved phase-defect budget]
\label{prop:v11v79-budget}
Under (79.9)--(79.10),
\[
                                  \rho_i\le\alpha_i+2\beta_i. \tag{79.11}
\]
Consequently a sufficient form of (79.7) is
\[
       \sum_{i\ge i_0}(\alpha_i+2\beta_i+\eta_i)<|D_{i_0}|.  \tag{79.12}
\]
\end{proposition}

\begin{proof}
For \(x\in{\cal G}_{i+1}\), charge compatibility gives
\(Q_{i+1}(x)=Q_i(P_ix)\), so the integrand in (79.4) is at most
\(|a_{i+1}-a_i\circ P_i|\).  On the complement, insert and subtract
\(e^{i\theta Q_{i+1}}a_i\circ P_i\); the residual term is unchanged
and the difference of two unit phases is at most two.  Integration
gives (79.11).
\end{proof}

The three series in (79.12) have different origins:
\(\alpha_i\) is determinant-phase transport,
\(\beta_i\) is failure of the gapped charge comparison, and
\(\eta_i\) is positive amplitude inconsistency.  A numerical bound on
their sum is useful only after each term has been estimated in its
own norm.

\subsection{Ways to certify the finite anchor}

\begin{lemma}[High-probability phase arc]
\label{lem:v11v79-arc}
Suppose an event \(G\subset X_{i_0}\), an angle \(\varphi\), and
\(0\le\gamma<\pi/2\) satisfy
\[
 \mu_{i_0}(G^c)\le\varepsilon,\qquad
 \operatorname{Re}\!\left(e^{-i\varphi}u_{i_0}(x)\right)
                    \ge\cos\gamma\quad(x\in G).              \tag{79.13}
\]
Then
\[
 |D_{i_0}|
 \ge (1-\varepsilon)\cos\gamma-\varepsilon.                  \tag{79.14}
\]
\end{lemma}

\begin{proof}
On \(G^c\), (79.1) gives
\(\operatorname{Re}(e^{-i\varphi}u_{i_0})\ge-1\).  Hence
\[
 \operatorname{Re}(e^{-i\varphi}D_{i_0})
 \ge(1-\varepsilon)\cos\gamma-\varepsilon.
\]
The modulus is at least this real part.
\end{proof}

A second anchor follows directly from the sector theorem.

\begin{lemma}[Small-theta anchor with residual phase]
\label{lem:v11v79-theta}
At scale \(i_0\), suppose
\[
 \mathbb E_{\mu_{i_0}}Q_{i_0}^2\le M_2,\qquad
 \mathbb E_{\mu_{i_0}}|a_{i_0}-1|\le\varepsilon_0.            \tag{79.15}
\]
Then
\[
                 |D_{i_0}|
                 \ge1-\frac{\theta^2M_2}{2}-\varepsilon_0.   \tag{79.16}
\]
\end{lemma}

\begin{proof}
Proposition~\ref{prop:v11v76-theta} bounds the pure charge
characteristic function below by \(1-\theta^2M_2/2\).
The replacement of \(1\) by \(a_{i_0}\) changes its expectation by at
most \(\varepsilon_0\).
\end{proof}

Thus an entirely explicit sufficient condition for continuum
noncancellation is
\[
 {\,\theta^2M_2\over2}+\varepsilon_0
 +\sum_{i\ge i_0}(\alpha_i+2\beta_i+\eta_i)<1.                \tag{79.17}
\]
It remains a conditional inequality until all constants are obtained
from the regulator.

\subsection{Normalized cylinder functionals}

Let \(F\) be a bounded cylinder observable represented at every
\(i\ge k\) by \(F_i\), with
\[
                                  F_{i+1}=F_i\circ P_i.        \tag{79.18}
\]
Put
\[
                                  N_i(F)=\int F_i u_i\,d\mu_i.
                                                                    \tag{79.19}
\]

\begin{theorem}[Complex cylinder limit]
\label{thm:v11v79-cylinder}
For \(i\ge k\),
\[
 |N_{i+1}(F)-N_i(F)|
             \le\|F\|_\infty(\rho_i+\eta_i).                 \tag{79.20}
\]
If (79.6)--(79.7) hold, then
\[
                         \Omega_\infty(F)
            =\lim_{i\to\infty}{N_i(F)\over D_i}              \tag{79.21}
\]
exists for every bounded cylinder \(F\), and
\[
 |\Omega_\infty(F)|
 \le {\|F\|_\infty\over |D_{i_0}|-R_{i_0}}.                  \tag{79.22}
\]
\end{theorem}

\begin{proof}
Repeat the proof of (79.5) with the common factor \(F_{i+1}\) and use
(79.18), obtaining (79.20).  Hence the numerators are Cauchy.
The denominators converge and have the positive lower bound (79.8),
so quotient convergence proves (79.21).  Since
\(|N_i(F)|\le\|F\|_\infty\), passage to the limit gives (79.22).
\end{proof}

\begin{warning}[Line-bundle comparison]
\label{warn:v11v79-line}
The subtraction in (79.4) and (79.10) is meaningful only after phases
at adjacent cutoffs have been placed in one trivialisation, or after
a specified unitary parallel transport has been inserted.  Local
anomaly cancellation can make the determinant connection flat, but
V73 already showed that flatness does not eliminate global holonomy.
The numbers \(\alpha_i\) must include this transport defect.
\end{warning}

\begin{firewall}[Positivity and clustering remain separate]
\label{fire:v11v79-boundary}
Theorem~\ref{thm:v11v79-cylinder} constructs a bounded normalized
complex cylinder functional.  It does not make that functional
reflection positive, prove Osterwalder--Schrader reconstruction, or
supply a mass gap.  Those conclusions require the separate symmetry,
positivity, and clustering hypotheses isolated in V64.
\end{firewall}

\begin{decision}[Mandatory V80 audit]
\label{dec:v11v79-next}
The next version is the compulsory companion audit.  It will inspect
the newest Yang--Mills and Navier--Stokes notes for a new charge-gap,
phase-transport, decay, or refinement estimate.  The following
substantive version will attack the strongest remaining input after
that audit.
\end{decision}

\begin{assessment}[Progress after V79]
\label{ass:v11v79-status}
A finite-scale phase certificate now survives to the continuum when
the determinant transport, charge-admissibility, and positive measure
defects have a summable budget smaller than the anchor.  Under the
same budget, every bounded cylinder numerator and its normalized
complex expectation converge.
\end{assessment}

\subsection{Parent record}

This V79 note was formed by appending the preceding material to the
validated V78 source, whose record was
\[
\begin{array}{c|c}
\text{lines}&24128\\
\text{bytes}&768592\\
\text{SHA--256}&
\texttt{47CEF75B0A43F2D2AD43353253B05D766F3393C837C1A6FE4DE21F17FB2111E4}.
\end{array}                                                     \tag{79.23}
\]

% =============================================================================
% END APPENDED ELEVENTH-SERIES V79 MATERIAL
% =============================================================================
% =============================================================================
% BEGIN APPENDED ELEVENTH-SERIES V80 MATERIAL
% =============================================================================

\section{Companion audit at V80: measured symmetry before global transport}
\label{sec:v11v80}

The compulsory audit records
\[
\begin{array}{c|r|r|l}
\text{file}&\text{lines}&\text{bytes}&\text{SHA--256}\\ \hline
\texttt{Renewal\_Yang\_Mills\_Mass\_Gap\_Research\_Notes\_V40.tex}
&16447&560619&
\texttt{F0B16C4DCCF204359A62342A87322FAE793567A22E3EF63AC813FEB4DF51951B}\\
\texttt{Renewal\_NS\_Stability\_Research\_Notes\_V26.tex}
&5319&278283&
\texttt{82C887F8C2C69E4F28FAD7C3DC8DE5B9099CB5A4BF431EBA1FFF3E91935978AD}.
\end{array}                                                     \tag{80.1}
\]

\subsection{New Yang--Mills information}

YM V39 evaluates the complete \(M=2\) cubic bubble grid at all sixteen
corners and proves that the raw bubble fails a residual hypercubic
orbit identity.  The failure is not a physical symmetry violation:
the quartic tadpole, stationary-coordinate, and Haar/quotient-density
rows must be added before the measured response is compressed by
orbits.

YM V40 then constructs the exact
\((\mathbb Z/2\mathbb Z)^4\) Walsh jet algebra for batching those
corner rows.  It also formally repairs the apostrophe-doubling
derivative defect already identified independently in SM V70.
Its acceptance criterion is exact cancellation of every orbit
residual after full measured recombination.  The measured batch has
not yet been executed.

Two conclusions transfer to this programme.
\begin{enumerate}
\item Symmetry and phase transport estimates must be attached to the
      complete regulated density, including coordinate and quotient
      Jacobians.  A diagram-level estimate need not transform as a
      physical scalar.
\item Exact finite-group character transforms can certify a
      finite-cutoff anchor, but an \(M=2\) corner calculation does not
      control the internal-momentum alias tail or the refinement
      series in (79.12).
\end{enumerate}

The new YM files explicitly state that they do not acquire the
uniform sector action cost, entropy, continuum measure, or clustering
inputs required here.

\subsection{Navier--Stokes information}

The NS file is unchanged.  Its last result gives closed-form
pointwise rank-one norms for first and second derivatives of the Oseen
kernel.  These constants refine a nonlinear stability estimate but
do not supply a gauge-sector tail, determinant phase comparison, or
Einstein refinement bound.  No NS theorem is imported.

\begin{warning}[No premature symmetry compression]
\label{warn:v11v80-compression}
A partial collection of diagrams or coordinate terms cannot be used
to certify the phase anchor in Lemma~\ref{lem:v11v79-arc} merely
because the final physical observable is symmetric.  The interval
enclosure must be applied after the full regulated density is
recombined and the exact Ward and orbit residuals vanish.
\end{warning}

\begin{decision}[V81 goes upstream to phase transport]
\label{dec:v11v80-next}
The dominant term still missing from (79.12) is
\(\alpha_i\), the adjacent determinant-phase transport defect.
V81 will derive a trace-norm Lipschitz estimate for the phase of a
relative Fredholm determinant on a zero-free path, including an
explicit singular-value denominator and a shell-summability
criterion.  This estimate must be applied to the fully recombined
regulated operator, in accordance with the V80 audit.
\end{decision}

\begin{firewall}[Audit boundary]
\label{fire:v11v80-boundary}
Neither companion closes sector tightness, determinant-line
triviality, reflection positivity, clustering, the coupled Einstein
measure, or the full continuum.  The Walsh corner algebra is a
finite-cutoff evaluation device and does not replace the global
analytic estimates of V76--V79.
\end{firewall}

\begin{assessment}[Progress after V80]
\label{ass:v11v80-status}
The audit adds a concrete ordering constraint: certify symmetries and
phase bounds only for the complete measured density.  With no new
sector input available, the programme turns upstream to the analytic
estimate needed to make determinant-phase increments summable.
\end{assessment}

\subsection{Parent record}

This V80 note was formed by appending the preceding material to the
validated V79 source, whose record was
\[
\begin{array}{c|c}
\text{lines}&24394\\
\text{bytes}&777357\\
\text{SHA--256}&
\texttt{E6FC7517EC872D1E1E8FFCAE810832BDFD1308A533BC2383CA456262DCCE0DE9}.
\end{array}                                                     \tag{80.2}
\]

% =============================================================================
% END APPENDED ELEVENTH-SERIES V80 MATERIAL
% =============================================================================
% =============================================================================
% BEGIN APPENDED ELEVENTH-SERIES V81 MATERIAL
% =============================================================================

\section{Determinant-phase transport from relative operator bounds}
\label{sec:v11v81}

The unresolved quantity \(\alpha_i\) in (79.12) compares determinant
phases at adjacent cutoffs.  This version derives it from a zero-free
operator path and a trace-ideal increment.  All operators below are
already stabilised onto a common Hilbert space; new reference blocks
and their known phases must be included in that stabilisation.

\subsection{Trace-class determinant phase}

Let
\[
                         A(t)=I+K(t),\qquad 0\le t\le1,       \tag{81.1}
\]
where \(K(t)\) is continuously differentiable in trace norm and
\(A(t)\) is invertible.  Since its Fredholm determinant is nonzero, it
has a continuous lifted argument \(\Theta(t)\) along this interval.

\begin{theorem}[Phase-length bound]
\label{thm:v11v81-phase-length}
Under the preceding assumptions,
\[
 \Theta'(t)
 =\operatorname{Im}\operatorname{Tr}
       \big(A(t)^{-1}K'(t)\big),                              \tag{81.2}
\]
and therefore
\[
 |\Theta(1)-\Theta(0)|
 \le\int_0^1\|A(t)^{-1}\|\,\|K'(t)\|_1\,dt.                  \tag{81.3}
\]
For the unit determinant phases
\[
                  v(t)={\det A(t)\over|\det A(t)|},
\]
one has
\[
 |v(1)-v(0)|
 \le\min\left\{2,\,
       \int_0^1\|A(t)^{-1}\|\,\|K'(t)\|_1\,dt\right\}.        \tag{81.4}
\]
\end{theorem}

\begin{proof}
The trace-norm differentiability formula for the Fredholm determinant
is
\[
 {d\over dt}\log\det A(t)
 =\operatorname{Tr}\big(A(t)^{-1}K'(t)\big).
\]
It follows first for finite-rank \(K(t)\) from the ordinary determinant
identity.  Trace-norm finite-rank approximation, continuity of the
Fredholm determinant, and
\[
 |\operatorname{Tr}(BC)|\le\|B\|\,\|C\|_1
\]
pass the identity to (81.1).  Taking imaginary parts and integrating
gives (81.2)--(81.3).  Finally,
\[
 |e^{ia}-e^{ib}|\le\min\{2,|a-b|\},
\]
which proves (81.4).
\end{proof}

\begin{corollary}[One-step singular-value certificate]
\label{cor:v11v81-one-step}
Let \(A_0=I+K_0\) satisfy
\[
                        s_{\min}(A_0)\ge\sigma>0,            \tag{81.5}
\]
and let \(\Delta=K_1-K_0\) obey
\[
                  \|\Delta\|\le\kappa<\sigma,\qquad
                  \|\Delta\|_1=d.                           \tag{81.6}
\]
Then the segment \(A(t)=A_0+t\Delta\) is invertible and
\[
 \begin{split}
 |\Theta_1-\Theta_0|
 &\le
 \begin{cases}
 {d\over\kappa}\log{\sigma\over\sigma-\kappa},&\kappa>0,\\
 0,&\kappa=0,
 \end{cases}                                                 \tag{81.7}\\
 &\le {d\over\sigma-\kappa}.
 \end{split}
\]
The same right-hand sides, truncated at two, bound the chord between
the two unit phases.
\end{corollary}

\begin{proof}
Weyl's singular-value inequality gives
\[
 s_{\min}(A_0+t\Delta)\ge\sigma-t\kappa>0.
\]
Use the segment in (81.3), for which \(K'(t)=\Delta\), and integrate
\(d/(\sigma-t\kappa)\).  The last inequality follows from
\(\log(\sigma/(\sigma-\kappa))\le
\kappa/(\sigma-\kappa)\).
\end{proof}

The trace norm in (81.6), rather than a bare operator norm, controls
the total determinant phase.  An operator-norm-small perturbation
repeated in an increasing number of modes can have an order-one or
divergent determinant phase.

\subsection{The regularised determinant}

For Hilbert--Schmidt \(K(t)\), define the second regularised
determinant \(\det_2(I+K)\).  On a zero-free path its logarithmic
derivative is
\[
 {d\over dt}\log\det_2(I+K(t))
 =
 \operatorname{Tr}\left[
 \big((I+K(t))^{-1}-I\big)K'(t)\right].                      \tag{81.8}
\]

\begin{theorem}[Hilbert--Schmidt phase bound]
\label{thm:v11v81-det2}
If \(K(t)\) is continuously differentiable in Hilbert--Schmidt norm
and \(I+K(t)\) is invertible, then the lifted phase
\(\Theta_2=\arg\det_2(I+K)\) obeys
\[
 |\Theta_2(1)-\Theta_2(0)|
 \le\int_0^1
 \|(I+K(t))^{-1}\|\,\|K(t)\|_2\,\|K'(t)\|_2\,dt.             \tag{81.9}
\]
For \(K(t)=K_0+t\Delta\), under (81.5)--(81.6) with
\(\Delta\) Hilbert--Schmidt,
\[
 |\Theta_2(1)-\Theta_2(0)|
 \le {(\|K_0\|_2+\|\Delta\|_2)\|\Delta\|_2
             \over\sigma-\kappa}.                           \tag{81.10}
\]
\end{theorem}

\begin{proof}
Since
\[
 (I+K)^{-1}-I=-(I+K)^{-1}K,
\]
the first factor in (81.8) is Hilbert--Schmidt and
\[
 |\operatorname{Tr}(BC)|\le\|B\|_2\|C\|_2
\]
gives (81.9).  On the segment,
\(\|K(t)\|_2\le\|K_0\|_2+\|\Delta\|_2\) and
\(\|(I+K(t))^{-1}\|\le(\sigma-\kappa)^{-1}\), proving (81.10).
\end{proof}

For \(\det_p\) with \(p>2\), the corresponding derivative subtracts
the first \(p-1\) trace terms.  Its estimate must be made in the
matching Schatten ideals.  No trace-class conclusion may be inferred
from a Hilbert--Schmidt estimate without this subtraction.

\subsection{Random shell transport}

At adjacent cutoffs, let \(A_i(P_ix)\) and \(A_{i+1}(x)\) be stabilised
onto the same space and set
\[
                  \Delta_i(x)=A_{i+1}(x)-A_i(P_ix).          \tag{81.11}
\]
Let \(v_i\) denote their unit Fredholm determinant phases.  Suppose a
good event \({\cal H}_{i+1}\) satisfies
\[
 \begin{split}
 s_{\min}(A_i(P_ix))&\ge\sigma_i,\\
 \|\Delta_i(x)\|&\le\vartheta\sigma_i,
                  \qquad 0<\vartheta<1
                  \quad(x\in{\cal H}_{i+1}),                 \tag{81.12}\\
 \mu_{i+1}({\cal H}_{i+1}^c)&\le\chi_i.
 \end{split}
\]

\begin{theorem}[Expected adjacent phase estimate]
\label{thm:v11v81-random}
Under (81.11)--(81.12),
\[
 \int|v_{i+1}(x)-v_i(P_ix)|\,d\mu_{i+1}(x)
 \le {1\over(1-\vartheta)\sigma_i}
      \int_{{\cal H}_{i+1}}\|\Delta_i(x)\|_1\,d\mu_{i+1}(x)
      +2\chi_i.                                              \tag{81.13}
\]
In particular determinant-phase transport is summable if
\[
 \sum_i{\mathbb E_{\mu_{i+1}}
       [\,{\bf1}_{{\cal H}_{i+1}}\|\Delta_i\|_1]\over\sigma_i}
 <\infty,\qquad
 \sum_i\chi_i<\infty.                                       \tag{81.14}
\]
\end{theorem}

\begin{proof}
On the good event, Corollary~\ref{cor:v11v81-one-step} gives the chord
bound
\[
 |v_{i+1}-v_i\circ P_i|
 \le{\|\Delta_i\|_1\over(1-\vartheta)\sigma_i}.
\]
On the complement it is at most two.  Integrate.
\end{proof}

\begin{corollary}[Insertion into the V79 budget]
\label{cor:v11v81-budget}
Suppose the residual phase in (79.9) is
\[
                            a_i=e^{ic_i}v_i,                 \tag{81.15}
\]
where \(c_i\) is a real lifted local counterterm phase.  Put
\[
 \zeta_i=\int\min\{2,|c_{i+1}(x)-c_i(P_ix)|\}\,
                         d\mu_{i+1}(x).                      \tag{81.16}
\]
Then the phase term in (79.10) satisfies
\[
 \alpha_i
 \le { {\mathbb E}_{\mu_{i+1}}
       [\,{\bf1}_{{\cal H}_{i+1}}\|\Delta_i\|_1]
       \over(1-\vartheta)\sigma_i}
       +2\chi_i+\zeta_i.                                    \tag{81.17}
\]
Thus (79.12) follows from the corresponding summability of
(81.17), \(2\beta_i\), and \(\eta_i\), together with the finite anchor.
\end{corollary}

\begin{proof}
Use
\[
 |e^{ic_{i+1}}v_{i+1}-e^{ic_i\circ P_i}v_i\circ P_i|
 \le |v_{i+1}-v_i\circ P_i|
     +|e^{ic_{i+1}}-e^{ic_i\circ P_i}|
\]
and \(|e^{is}-e^{it}|\le\min\{2,|s-t|\}\), then apply (81.13).
\end{proof}

\subsection{Where the gap can come from}

The uniform factor \(\sigma_i^{-1}\) in (81.17) is the main
bottleneck.  V62 supplies the appropriate architecture: split off
exact zero and low modes, treat the high block through a Feshbach
operator, and apply V81 only where a relative high-mode gap is
available.  A physical massless photon or graviton cannot be assigned
a fictitious gap to improve (81.17).

In a nonzero-index chiral sector, the unprimed determinant vanishes
because of exact zero modes.  One must first form the correct
zero-mode-saturated correlation or primed determinant and prove that
its remaining carrier is zero-free.  Its phase can then be treated by
the same argument.

\begin{warning}[Global determinant line]
\label{warn:v11v81-holonomy}
Equations (81.2)--(81.17) control phase variation inside a chosen
line-bundle trivialisation.  If adjacent charts have transition phase
\(\tau_i(x)\), its circular increment must be added to \(\zeta_i\).
A flat connection with nontrivial holonomy can make every local
integral small while retaining a global phase mismatch.
\end{warning}

\begin{firewall}[No acquired shell estimate]
\label{fire:v11v81-boundary}
V81 proves the analytic conversion from relative operator increments
to phase increments.  The attached sources do not yet prove the
weighted trace-norm series (81.14), the good-event probabilities
\(\chi_i\), or a zero-free primed chiral carrier uniformly across the
coupled SM--Einstein sectors.  The full determinant phase denominator
therefore remains open.
\end{firewall}

\begin{decision}[V82 acquisition]
\label{dec:v11v81-next}
V82 will go one step further upstream and derive trace-norm shell
bounds from quasi-local kernel decay and boundary-layer counting.
It will state the scale exponents needed for (81.14), separately for
massive high modes and retained massless low modes.
\end{decision}

\begin{assessment}[Progress after V81]
\label{ass:v11v81-status}
The abstract phase term \(\alpha_i\) now has an explicit operator
certificate.  A zero-free singular-value bound and a summable
trace-class shell increment control the Fredholm phase; a
Hilbert--Schmidt analogue controls \(\det_2\); bad gaps and local
counterterm phases enter as separate additive errors.
\end{assessment}

\subsection{Parent record}

This V81 note was formed by appending the preceding material to the
validated V80 source, whose record was
\[
\begin{array}{c|c}
\text{lines}&24508\\
\text{bytes}&782204\\
\text{SHA--256}&
\texttt{281970D9B2AEF95D39A21B9E3D50540F74E2C6D01CAB90D59F5182B5B66C7DF9}.
\end{array}                                                     \tag{81.18}
\]

% =============================================================================
% END APPENDED ELEVENTH-SERIES V81 MATERIAL
% =============================================================================
% =============================================================================
% BEGIN APPENDED ELEVENTH-SERIES V82 MATERIAL
% =============================================================================

\section{Quasi-local shell bounds for the phase-transport norm}
\label{sec:v11v82}

V81 reduces determinant-phase transport to a weighted trace norm.
This version derives that norm from a kernel estimate and shows the
different scale cost of a bulk perturbation and a boundary or defect
perturbation.

Let \(\Lambda\) be a finite metric lattice, with fibre
\(\mathbb C^d\), and let \(T\) act on
\(\ell^2(\Lambda;\mathbb C^d)\) with block kernel \(T(x,y)\).
Write \(\|\cdot\|_{1,d}\) and \(\|\cdot\|_{2,d}\) for the
finite-dimensional trace and Hilbert--Schmidt norms of a block.

\begin{lemma}[Kernel entry bounds for trace ideals]
\label{lem:v11v82-entry}
For every such \(T\),
\[
 \|T\|_1\le\sum_{x,y\in\Lambda}\|T(x,y)\|_{1,d},\qquad
 \|T\|_2^2=\sum_{x,y\in\Lambda}\|T(x,y)\|_{2,d}^2.            \tag{82.1}
\]
\end{lemma}

\begin{proof}
Let \(E_{xy}\) be the rank-one matrix from site \(y\) to site \(x\).
Then
\[
                         T=\sum_{x,y}E_{xy}\otimes T(x,y).
\]
Since \(\|E_{xy}\|_1=1\), the trace-norm triangle inequality gives the
first statement.  The block matrix units are orthonormal for the
Hilbert--Schmidt inner product, which gives the second.
\end{proof}

Assume the lattices have uniform polynomial growth: for some
\(D\ge1\) and \(c_D\),
\[
 \#\{y:n\le d(x,y)<n+1\}\le c_D(1+n)^{D-1}.                  \tag{82.2}
\]
Define the finite constants
\[
 \begin{split}
 S_{\exp}(m)&=c_D\sum_{n\ge0}(1+n)^{D-1}e^{-mn},\\
 S_{\rm pol}(p)&=c_D\sum_{n\ge0}(1+n)^{D-1-p},
                         \qquad p>D.                         \tag{82.3}
 \end{split}
\]

\begin{theorem}[Boundary-layer trace-norm estimate]
\label{thm:v11v82-boundary}
Let \(B\subset\Lambda\).  Suppose \(T(x,y)=0\) unless
\(x\in B\) or \(y\in B\), and
\[
                 \|T(x,y)\|_{1,d}
                 \le Ca e^{-m d(x,y)}.                      \tag{82.4}
\]
Then
\[
                              \|T\|_1
                    \le2Ca\,|B|S_{\exp}(m).                 \tag{82.5}
\]
If (82.4) is replaced by
\(Ca(1+d(x,y))^{-p}\) with \(p>D\), then
\[
                              \|T\|_1
                    \le2Ca\,|B|S_{\rm pol}(p).              \tag{82.6}
\]
\end{theorem}

\begin{proof}
The sum in (82.1) over pairs with \(x\in B\) is at most
\(Ca|B|S_{\exp}(m)\), by (82.2)--(82.3).  The same bound holds for
pairs with \(y\in B\); counting their intersection twice is harmless.
This proves (82.5).  The polynomial calculation is identical.
\end{proof}

\begin{corollary}[Bulk cost]
\label{cor:v11v82-bulk}
If the same kernel bound holds without the support restriction, then
\[
 \|T\|_1\le Ca\,|\Lambda|S_{\exp}(m)                         \tag{82.7}
\]
or, in the polynomial case with \(p>D\),
\[
 \|T\|_1\le Ca\,|\Lambda|S_{\rm pol}(p).                     \tag{82.8}
\]
Thus quasi-locality alone does not remove the volume factor.
\end{corollary}

\begin{proof}
Repeat the first half of the proof of
Theorem~\ref{thm:v11v82-boundary} with \(B=\Lambda\), without the
second row/column copy.
\end{proof}

This distinction is decisive for (81.14).  A relative cutoff
increment must be localised to a shrinking defect layer, or its bulk
part must be cancelled by a local counterterm.  Exponential off-
diagonal decay does not make a translation-invariant bulk increment
trace class uniformly in volume.

\subsection{Hilbert--Schmidt comparison}

\begin{proposition}[Boundary-layer Hilbert--Schmidt estimate]
\label{prop:v11v82-HS}
Under the support assumption of
Theorem~\ref{thm:v11v82-boundary}, suppose instead that
\[
                 \|T(x,y)\|_{2,d}
                 \le Ca e^{-m d(x,y)}.                      \tag{82.9}
\]
Then
\[
                  \|T\|_2
                  \le Ca\,[\,2|B|S_{\exp}(2m)\,]^{1/2}.      \tag{82.10}
\]
Without boundary support, the right side is
\(Ca[\,|\Lambda|S_{\exp}(2m)\,]^{1/2}\).
\end{proposition}

\begin{proof}
Square (82.9), sum by (82.1), and use the same row and column
counting as in (82.5).
\end{proof}

The square-root gain in (82.10) is why a regularised determinant can
be preferable.  It must, however, be combined with the factor
\(\|K_i\|_2\) in (81.10).  If \(\|K_i\|_2\) itself grows like the
square root of the volume, the product can again be extensive.

\subsection{Dyadic scale exponents}

Let the lattice scale be \(i\), and suppose
\[
 |\Lambda_i|\le C_\Lambda 2^{Di},\qquad
 |B_i|\le C_B2^{(D-1)i},                                    \tag{82.11}
\]
where \(B_i\) has a uniformly bounded thickness in lattice units.
Assume on the good event in V81 that
\[
 \|\Delta_i(x;y,z)\|_{1,d}
       \le C_\Delta 2^{-si}e^{-m d(y,z)},\qquad
 \sigma_i\ge C_\sigma2^{-gi}.                               \tag{82.12}
\]

\begin{theorem}[Trace-class shell summability]
\label{thm:v11v82-exponents}
If \(\Delta_i\) is supported in the rows or columns of \(B_i\), then
\[
 {\|\Delta_i\|_1\over\sigma_i}
 \le {2C_\Delta C_BS_{\exp}(m)\over C_\sigma}
        2^{-(s-D+1-g)i}.                                    \tag{82.13}
\]
Consequently
\[
                                  s>D-1+g                    \tag{82.14}
\]
is sufficient for the deterministic part of (81.14).  For a bulk
increment, the corresponding sufficient condition is
\[
                                  s>D+g.                     \tag{82.15}
\]
\end{theorem}

\begin{proof}
Insert (82.11)--(82.12) into (82.5), divide by the gap, and sum the
resulting geometric series.  Equation (82.7) gives the bulk exponent.
\end{proof}

For a physical four-dimensional bulk, (82.15) demands decay faster
than \(2^{-(4+g)i}\).  This makes explicit why a naive comparison of
full cutoff matrices is usually too costly.  Exact local
renormalisation must first remove the extensive contribution.

If a defect has codimension \(c\), so that
\[
                              |B_i|\le C_B2^{(D-c)i},         \tag{82.16}
\]
the same proof replaces (82.14) by
\[
                              s>D-c+g.                       \tag{82.17}
\]
This links the singular-stratum geometry of V54 to the phase budget.

\subsection{Random amplitudes and bad decay events}

Suppose (82.12) holds with the deterministic amplitude replaced by a
nonnegative random \(A_i(x)\):
\[
 \|\Delta_i(x;y,z)\|_{1,d}
       \le A_i(x)2^{-si}e^{-m d(y,z)}.                       \tag{82.18}
\]
Then
\[
 {\mathbb E[
 {\bf1}_{{\cal H}_{i+1}}\|\Delta_i\|_1]\over\sigma_i}
 \le {2C_BS_{\exp}(m)\over C_\sigma}
       2^{-(s-D+1-g)i}
       \mathbb E[{\bf1}_{{\cal H}_{i+1}}A_i].                \tag{82.19}
\]
Uniform \(L^1\) control of \(A_i\) therefore preserves (82.14).
An \(L^r\) estimate can be transferred from a reference measure by
V78, but incurs its Holder exponent and lower-partition-function
requirement.

\subsection{Stabilisation and low modes}

Let \(J_i\) be the reference operator on the new modes.  The
stabilised comparison has the form
\[
       A_i^{\rm stab}=A_i\oplus J_i,\qquad
       \Delta_i=A_{i+1}-A_i^{\rm stab}.                      \tag{82.20}
\]
If \(J_i\) is chosen poorly, \(\Delta_i\) has an order-one diagonal
entry on every new mode and (82.7) is unavoidable.  The reference
block must match the principal high-frequency symbol, with its
known determinant phase absorbed consistently into the counterterm
of (81.15).

For massless modes, exponential kernel decay is generally false.
Polynomial decay with exponent \(p\le D\) makes the row sum in (82.3)
diverge with volume.  Those modes must remain in the low Feshbach
block of V62.  V82 applies to the massive or elliptic high block after
a genuine gap has produced the decay constant \(m>0\).

\begin{warning}[Decay constants must be uniform]
\label{warn:v11v82-uniform}
A Combes--Thomas or resolvent estimate whose decay rate
\(m_i\downarrow0\) can make \(S_{\exp}(m_i)\) grow and destroy
(82.14).  The scale exponent must include that growth.  Likewise,
a boundary layer of increasing lattice thickness changes the count
in (82.11).
\end{warning}

\begin{firewall}[No renormalised kernel acquired]
\label{fire:v11v82-boundary}
Theorems~\ref{thm:v11v82-boundary} and
\ref{thm:v11v82-exponents} convert a stated quasi-local remainder into
the V81 phase norm.  The current programme has not yet constructed
the fully measured, counterterm-subtracted SM--Einstein relative
kernel with the decay (82.12), nor proved the high-block gap or the
random amplitude bound in (82.19).
\end{firewall}

\begin{decision}[V83 acquisition]
\label{dec:v11v82-next}
V83 will formulate the local counterterm cancellation needed to turn
an extensive bulk phase derivative into the boundary or defect
remainder of V82.  It will separate exact local trace matching from
summable truncation error and from global determinant-line holonomy.
\end{decision}

\begin{assessment}[Progress after V82]
\label{ass:v11v82-status}
Quasi-local kernel control now implies explicit trace-class and
Hilbert--Schmidt shell bounds.  Boundary support saves one power of
the refinement scale, codimension \(c\) saves \(c\) powers, and a
bulk comparison remains extensive.  The exact summability thresholds
exhibit the required balance among decay, geometry, and the
singular-value loss.
\end{assessment}

\subsection{Parent record}

This V82 note was formed by appending the preceding material to the
validated V81 source, whose record was
\[
\begin{array}{c|c}
\text{lines}&24817\\
\text{bytes}&792303\\
\text{SHA--256}&
\texttt{E3CC3D4CB5D7169A822E6C226FED8166749BD0C19AC860A3D307ED32C8500D1B}.
\end{array}                                                     \tag{82.21}
\]

% =============================================================================
% END APPENDED ELEVENTH-SERIES V82 MATERIAL
% =============================================================================
% =============================================================================
% BEGIN APPENDED ELEVENTH-SERIES V83 MATERIAL
% =============================================================================

\section{Exact local counterterm cancellation of the extensive phase}
\label{sec:v11v83}

The bulk estimate (82.7) is generally too large for the V81 phase
budget.  Renormalisation can remove it only when the extensive phase
variation is the differential of an allowed local functional.  This
version separates that exact cancellation from the quasi-local
remainder and from determinant-line topology.

Let \(U_i\) be a smooth finite-dimensional admissible chart in cutoff
configuration space.  On a zero-free determinant chart define the
real one-form
\[
 \omega_i(x)[h]
 =\operatorname{Im}\operatorname{Tr}
       \big(A_i(x)^{-1}D A_i(x)[h]\big).                     \tag{83.1}
\]
It is the differential of a locally lifted determinant argument.
Suppose it has a decomposition
\[
                              \omega_i=b_i+r_i+e_i,           \tag{83.2}
\]
where \(b_i\) is the intended local bulk term, \(r_i\) is a boundary
or defect remainder, and \(e_i\) is a matching error.

\begin{theorem}[Exact counterterm removal]
\label{thm:v11v83-cancel}
Assume there is a real local functional \(L_i:U_i\to\mathbb R\) such
that
\[
                                      dL_i=b_i.              \tag{83.3}
\]
Define the renormalised lifted phase
\[
                       \Phi_i=\arg\det A_i-L_i.              \tag{83.4}
\]
Then, for every piecewise \(C^1\) path \(x:[0,1]\to U_i\),
\[
 \Phi_i(x(1))-\Phi_i(x(0))
       =\int_0^1(r_i+e_i)_{x(t)}[\dot x(t)]\,dt,              \tag{83.5}
\]
and the unit phases obey
\[
 \left|e^{i\Phi_i(x(1))}-e^{i\Phi_i(x(0))}\right|
 \le\min\left\{2,\int_0^1
       \big(|r_i[\dot x]|+|e_i[\dot x]|\big)\,dt\right\}.     \tag{83.6}
\]
\end{theorem}

\begin{proof}
Differentiate (83.4) along the path.  Equations (83.1)--(83.3) give
\[
 {d\over dt}\Phi_i(x(t))
 =(\omega_i-b_i)_{x(t)}[\dot x(t)]
 =(r_i+e_i)_{x(t)}[\dot x(t)].
\]
Integration proves (83.5), and the unit-circle chord inequality proves
(83.6).
\end{proof}

The cancellation in (83.5) is exact.  Estimating \(b_i\) and
\(dL_i\) separately by absolute values would restore the extensive
volume term and lose the purpose of the counterterm.

\subsection{Local exactness and the curvature test}

On a star-shaped chart \(U\) in a vector space, closedness is enough
to construct the counterterm.

\begin{theorem}[Poincare construction on a regulator chart]
\label{thm:v11v83-poincare}
Let \(b\) be a continuously differentiable real one-form on a
star-shaped open set \(U\), with centre \(x_0\), and suppose
\[
                                      db=0.                  \tag{83.7}
\]
Then
\[
 L(x)=\int_0^1
 b_{x_0+t(x-x_0)}[x-x_0]\,dt                                \tag{83.8}
\]
satisfies \(dL=b\) and \(L(x_0)=0\).
\end{theorem}

\begin{proof}
Put \(v=x-x_0\).  Differentiating (83.8) in direction \(h\) gives
\[
 DL(x)[h]
 =\int_0^1\left\{
 t(Db)_{x_0+tv}[h,v]+b_{x_0+tv}[h]\right\}\,dt.
\]
Closedness says
\((Db)[h,v]=(Db)[v,h]\).  The integrand is therefore
\[
 {d\over dt}\{t\,b_{x_0+tv}[h]\}.
\]
Its integral is \(b_x[h]\).
\end{proof}

\begin{corollary}[Curvature obstruction]
\label{cor:v11v83-curvature}
If \(db_i\ne0\) on a chart, no scalar counterterm can satisfy
(83.3) there.  More generally, if a loop \(\Gamma\) bounds a smooth
surface \(\Sigma\subset U_i\), then
\[
                              \oint_\Gamma b_i
                              =\int_\Sigma db_i.             \tag{83.9}
\]
Thus a nonzero local anomaly curvature produces a path-dependent
phase mismatch.
\end{corollary}

\begin{proof}
Every exact one-form is closed.  The second statement is Stokes'
theorem.
\end{proof}

For an anomaly-free multiplet, cancellation of the sum of local
curvatures is the relevant test.  Cancellation species by species is
unnecessary, while cancellation of only the integrated anomaly on
one path is insufficient.

\subsection{Local density matching and scale cost}

Suppose the bulk form and the implemented counterterm form admit site
densities on an interior set \(I_i\):
\[
 b_i[h]=\sum_{x\in I_i}b_i(x)[h],\qquad
 dL_i^{\rm app}[h]=\sum_{x\in I_i}\ell_i(x)[h].              \tag{83.10}
\]
If
\[
                         |b_i(x)[h]-\ell_i(x)[h]|
                         \le\epsilon_i\,q_i(x,h),            \tag{83.11}
\]
then
\[
 |(b_i-dL_i^{\rm app})[h]|
       \le\epsilon_i\sum_{x\in I_i}q_i(x,h).                 \tag{83.12}
\]
In particular, a uniform bound \(q_i\le Q\) costs
\(Q\epsilon_i|I_i|\).  With \(|I_i|=O(2^{Di})\), summability requires
the per-site matching error to decay faster than \(2^{-Di}\), after
including the singular-value loss.  Approximate equality of local
densities is therefore not interchangeable with the exact identity
(83.3).

If the mismatch is supported on a codimension-\(c\) layer, the count
improves to \(O(2^{(D-c)i})\), in agreement with (82.17).

\begin{proposition}[Counterterm plus quasi-local remainder]
\label{prop:v11v83-shell}
Assume (83.3).  Suppose that along the adjacent interpolation
\(x_i(t)\),
\[
 \int_0^1|r_i(x_i(t))[\dot x_i(t)]|\,dt\le R_i,\qquad
 \int_0^1|e_i(x_i(t))[\dot x_i(t)]|\,dt\le E_i.              \tag{83.13}
\]
Then its renormalised determinant phase increment is at most
\[
                                     \min\{2,R_i+E_i\}.       \tag{83.14}
\]
If \(R_i\) is obtained from the V82 boundary-kernel estimate and
\(\sum_i(R_i+E_i)<\infty\), this part of \(\alpha_i\) is summable.
\end{proposition}

\begin{proof}
This is (83.6) applied to the adjacent interpolation.
\end{proof}

\subsection{Counterterm covariance and refinement}

An allowed \(L_i\) must be a functional of the complete regulated
field and measure data.  Write the adjacent counterterm mismatch as
\[
 \delta_i^L(x)
 =L_{i+1}(x)-L_i(P_ix)-L_i^{\rm new}(x),                    \tag{83.15}
\]
where \(L_i^{\rm new}\) is the declared reference contribution of the
new high modes.  Even if \(dL_i=b_i\) separately at each scale, the
circular quantity
\[
                   \min\{2,|\delta_i^L(x)-c_i^{\rm ref}|\}   \tag{83.16}
\]
must be included in the \(\zeta_i\) term of (81.16), with
\(c_i^{\rm ref}\) the stabilising determinant-block phase.
Running coefficients, discretisation changes, and quotient
Jacobians all enter (83.15).

\begin{lemma}[Gauge-invariant descent]
\label{lem:v11v83-descent}
Let a compact gauge group act smoothly on \(U_i\).  If \(b_i\) is
gauge invariant and horizontal,
\[
 g^*b_i=b_i,\qquad b_i(\xi^\#)=0                             \tag{83.17}
\]
for every infinitesimal gauge vector \(\xi^\#\), and if the
construction (83.8) uses a gauge-equivariant contraction to a fixed
orbit, then \(L_i\) descends to the quotient, up to an additive
constant on each connected component.
\end{lemma}

\begin{proof}
Horizontality and invariance make \(b_i\) basic, hence it is the
pullback of a one-form on the quotient.  Apply
Theorem~\ref{thm:v11v83-poincare} on the contracted quotient chart.
Two primitives of the same one-form differ by a constant on each
connected component.
\end{proof}

A gauge-fixed coordinate expression that omits the Faddeev--Popov,
Haar, or quotient-density contribution need not satisfy (83.17).
The V80 audit therefore applies directly to the counterterm
calculation.

\subsection{Global periods}

Even when \(db_i=0\), a global primitive exists only if every period
vanishes:
\[
                               \oint_\Gamma b_i=0
                    \quad\text{for all closed }\Gamma.       \tag{83.18}
\]
A local counterterm changes a connection by an exact form and hence
cannot change its loop periods.  These periods are precisely the
global anomaly or determinant-line holonomy obstruction that was
left outside V81.

\begin{warning}[Renormalisation cannot be chosen to force an answer]
\label{warn:v11v83-renormalisation}
The imaginary local counterterms allowed in (83.4) are constrained by
gauge symmetry, diffeomorphism symmetry, reflection structure, and
the chosen physical theta parameters.  An arbitrary nonlocal
primitive along one refinement path could make (83.5) small but would
define a different, path-dependent theory.  It is not an admissible
renormalisation.
\end{warning}

\begin{firewall}[Missing local cohomology calculation]
\label{fire:v11v83-boundary}
The Standard Model representation is expected to cancel the familiar
perturbative gauge-anomaly polynomial, but V83 does not substitute
that expectation for a regulator-level proof of (83.7), the quotient
descent (83.17), the refinement mismatch (83.16), or the global
period condition (83.18).  Mixed gauge--gravitational and purely
gravitational phases must be included in the combined calculation.
\end{firewall}

\begin{decision}[V84 acquisition]
\label{dec:v11v83-next}
V84 will treat the global part: a flat determinant line will be
trivialised exactly when its holonomy representation is trivial.
A finite good-cover and spanning-tree criterion will turn transition
phases and cycle residuals into an adjacent transport budget.
\end{decision}

\begin{assessment}[Progress after V83]
\label{ass:v11v83-status}
The extensive phase can now be removed without a volume estimate
precisely when its local bulk one-form is the differential of an
allowed counterterm.  Curvature, per-site matching error, refinement
mismatch, quotient descent, and global periods have been exposed as
separate verifiable conditions.
\end{assessment}

\subsection{Parent record}

This V83 note was formed by appending the preceding material to the
validated V82 source, whose record was
\[
\begin{array}{c|c}
\text{lines}&25099\\
\text{bytes}&802204\\
\text{SHA--256}&
\texttt{4472B2A4C6B556DA35878EE999F79E3FDC4E8AF27D374A050EAAEBCA684BDFFF}.
\end{array}                                                     \tag{83.19}
\]

% =============================================================================
% END APPENDED ELEVENTH-SERIES V83 MATERIAL
% =============================================================================
% =============================================================================
% BEGIN APPENDED ELEVENTH-SERIES V84 MATERIAL
% =============================================================================

\section{Global determinant-line holonomy and finite cycle tests}
\label{sec:v11v84}

V83 removes an exact local bulk phase, but a flat determinant
connection can still have nontrivial periods.  This version gives the
global criterion and a finite good-cover implementation suitable for
the phase budget of V79.

Let \(L\to X\) be a Hermitian line bundle over a path-connected
space, with a flat unitary connection \(\nabla\).  Fix \(x_0\in X\).
Parallel transport defines
\[
                  {\rm Hol}_{x_0}:\pi_1(X,x_0)\longrightarrow U(1).
                                                                    \tag{84.1}
\]

\begin{theorem}[Flat-line trivialisation criterion]
\label{thm:v11v84-holonomy}
There exists a nowhere-zero global parallel unit section of \(L\) if
and only if
\[
                         {\rm Hol}_{x_0}([\Gamma])=1
                    \quad\text{for every closed loop }\Gamma.       \tag{84.2}
\]
When it exists, such a section is unique up to one constant phase.
\end{theorem}

\begin{proof}
A global parallel section returns to itself after transport around
every loop, so (84.2) is necessary.  Conversely, choose a unit vector
\(v_0\in L_{x_0}\).  For \(x\in X\), transport \(v_0\) along any path
from \(x_0\) to \(x\).  Two paths differ by a loop, and (84.2) makes
the result path-independent.  The resulting unit section is parallel.
Two such sections have a locally constant ratio in \(U(1)\), hence a
constant ratio because \(X\) is path-connected.
\end{proof}

\begin{corollary}[Local counterterms do not change holonomy]
\label{cor:v11v84-counterterm}
A phase redefinition by a single-valued function
\(e^{if}:X\to U(1)\) changes a local connection one-form by \(df\)
and leaves every holonomy in (84.1) unchanged.
\end{corollary}

\begin{proof}
The integral of \(df\) around a closed loop is zero.  Equivalently,
the endpoint factors in a gauge transformation cancel on a loop.
\end{proof}

Thus the local cancellation of V83 can remove exact bulk terms but
cannot repair a global anomaly.

\subsection{Good-cover cycle certificate}

Let \(\{U_a\}_{a\in V}\) be a finite good cover of \(X\): every
nonempty finite intersection is contractible.  Flatness permits a
parallel unit section \(s_a\) on each \(U_a\).  On a nonempty overlap,
\[
                              s_b=g_{ab}s_a,\qquad g_{ab}\in U(1),
                                                                    \tag{84.3}
\]
and \(g_{ab}\) is constant.  The constants satisfy
\[
                    g_{ba}=g_{ab}^{-1},\qquad
                    g_{ab}g_{bc}g_{ca}=1                    \tag{84.4}
\]
on triple overlaps.

Let \(G\) be the connected one-skeleton of the cover nerve and choose
a spanning tree \(T\subset G\).

\begin{theorem}[Spanning-tree holonomy test]
\label{thm:v11v84-tree}
There are constant rephasings of the sections \(s_a\) for which
\[
                              g_{ab}=1\quad\text{on every edge of }T.
                                                                    \tag{84.5}
\]
After this gauge choice, the transition on each off-tree edge
\(e\in G\setminus T\) equals the holonomy around its fundamental
cycle \(C_e\), up to the orientation convention.  Therefore \(L\)
has a global parallel unit section if and only if
\[
                              \prod_{f\in C_e}g_f=1
                  \quad(e\in G\setminus T).                 \tag{84.6}
\]
\end{theorem}

\begin{proof}
Choose a root.  Rephase its section arbitrarily, then proceed outward
along the unique tree paths, choosing the phase at each new vertex so
that its incoming tree transition is one.  This proves (84.5).
For an off-tree edge, its edge together with the unique tree path
between its endpoints forms \(C_e\).  All tree factors are one, so its
remaining transition is the cycle product.  The fundamental cycles
generate the graph cycle space.  The cocycle relations (84.4) fill
the two-simplices of the nerve, so triviality on the fundamental
cycles is equivalent to path-independent gluing.  This is precisely
the construction in Theorem~\ref{thm:v11v84-holonomy}.
\end{proof}

\begin{corollary}[Simply connected admissible component]
\label{cor:v11v84-simply}
If an admissible component \(X\) is simply connected and the
determinant connection is flat, then its determinant line has a
global parallel unit section.
\end{corollary}

\begin{proof}
Its fundamental group and hence every holonomy are trivial.
\end{proof}

Simple connectivity must be proved for the actual quotient
configuration component.  A local star-shaped gauge slice does not
imply that the global quotient is simply connected.

\subsection{Quantitative cycle residual}

The same tree construction is useful before exact cancellation has
been certified.  Let a connected finite graph carry oriented edge
phases \(z_{ab}\in U(1)\), with \(z_{ba}=z_{ab}^{-1}\).  Rephase the
vertices so all tree edges equal one, and let
\[
                              h_e=\prod_{f\in C_e}z_f
                       \quad(e\notin T).                    \tag{84.7}
\]
Define
\[
                              {\cal H}(G,T)
                       =\sum_{e\notin T}|h_e-1|.             \tag{84.8}
\]

\begin{proposition}[Finite gluing-defect bound]
\label{prop:v11v84-defect}
After the tree gauge, every off-tree transition equals \(h_e\).
For any simple nerve path \(\gamma\),
\[
 \left|\prod_{e\in\gamma}z_e-1\right|
 \le\sum_{\substack{e\in\gamma\\e\notin T}}|h_e-1|
 \le{\cal H}(G,T),                                          \tag{84.9}
\]
provided the endpoints are compared in the tree-induced
trivialisation.
\end{proposition}

\begin{proof}
The first statement is the algebraic part of
Theorem~\ref{thm:v11v84-tree}.  For unit complex numbers,
\[
 |w_1\cdots w_m-1|\le\sum_{j=1}^m|w_j-1|,
\]
by repeated insertion and the triangle inequality.  Tree factors are
one, and a simple path uses each off-tree edge at most once.
\end{proof}

The number (84.8) is invariant under the initial vertex rephasings
because each \(h_e\) is a cycle product.  It depends on the chosen
cycle basis only through an equivalent finite collection; exact
vanishing is basis-independent.

\subsection{Insertion into adjacent cutoff transport}

Let \(X_i(q)\) be an admissible charge component, let \(L_i(q)\) be
its determinant line, and suppose a unitary refinement comparison
\[
                 {\cal R}_i(x):
                 L_{i+1}(Q_{i+1}(x))
                 \longrightarrow L_i(Q_i(P_ix))             \tag{84.10}
\]
is defined on the good charge set.  When both lines have global
parallel unit sections \(s_{i+1}\) and \(s_i\), write
\[
                 {\cal R}_i(x)s_{i+1}(x)
                    =\tau_i(x)s_i(P_ix),\qquad \tau_i(x)\in U(1).
                                                                    \tag{84.11}
\]

\begin{theorem}[Holonomy contribution to the phase budget]
\label{thm:v11v84-budget}
Suppose:
\begin{enumerate}
\item the local operator and counterterm estimate of V81--V83 gives a
      phase chord \(d_i^{\rm loc}(x)\);
\item a good-cover comparison protocol crosses a simple nerve path
      with cycle defect bounded by \({\cal H}_i(q)\);
\item the refinement scalar in (84.11) has residual
      \(d_i^{\rm ref}(x)=|\tau_i(x)-1|\).
\end{enumerate}
Then the full adjacent determinant-phase chord is at most
\[
             d_i^{\rm loc}(x)+{\cal H}_i(q)
                         +d_i^{\rm ref}(x),                  \tag{84.12}
\]
truncated at two.  Hence its expectation is summable if the
expectations of these three terms are summable.
\end{theorem}

\begin{proof}
Transport first within the fine chart system, then through
\({\cal R}_i\), and finally within the coarse chart system.  Insert
the intermediate unit phases.  The local segment is bounded by
\(d_i^{\rm loc}\), Proposition~\ref{prop:v11v84-defect} bounds the
chart-gluing segment, and (84.11) bounds the refinement segment.
The triangle inequality gives (84.12).
\end{proof}

Exact global triviality makes \({\cal H}_i(q)=0\).  It does not force
\(d_i^{\rm ref}=0\): coherent trivialisations must also be selected
across cutoffs.

\subsection{Disconnected sector constants}

Theorem~\ref{thm:v11v84-holonomy} determines a parallel section only
up to a constant phase on each connected charge component.  If the
configuration space is
\[
                             X_i=\coprod_qX_i(q),             \tag{84.13}
\]
the constants for distinct \(q\) are independent.  Their relative
choice is physical when theta angles or zero-mode orientations are
present.  The common-half-plane hypothesis in V74 and the finite
anchor in V79 are exactly the additional cross-sector data.

\begin{warning}[Flat is not trivial]
\label{warn:v11v84-flat}
Vanishing perturbative anomaly curvature proves only flatness.
A loop can still have holonomy \(-1\), as in a possible
\(\mathbb Z_2\) global anomaly.  No local polynomial counterterm
changes that sign.  The cycle products (84.6), or an equivalent
bordism or index calculation, must be evaluated for the complete
fermion representation.
\end{warning}

\begin{firewall}[Uncomputed SM--Einstein periods]
\label{fire:v11v84-boundary}
V84 proves the criterion but does not compute the holonomy
representation of the regulated Standard Model coupled to dynamical
gravity.  Gauge transformations outside the identity component,
spin-structure changes, diffeomorphism loops, and mixed
gauge--gravitational cycles remain to be classified.  The determinant
line and its refinement maps in (84.10) have not yet been constructed
globally.
\end{firewall}

\begin{decision}[Mandatory V85 audit]
\label{dec:v11v84-next}
The next version is the compulsory companion audit.  It will check
whether the YM measured Walsh batch or the NS rank-one programme has
advanced and whether either supplies a complete-density symmetry or
decay input.  V86 will then consolidate the acquired local and global
phase criteria into the next continuum bottleneck.
\end{decision}

\begin{assessment}[Progress after V84]
\label{ass:v11v84-status}
Local anomaly cancellation has been separated from global
determinant-line triviality.  A finite good-cover spanning tree now
reduces global gluing to explicit cycle products, and approximate
cycle and refinement residuals enter the V79 phase budget with a
rigorous chord bound.
\end{assessment}

\subsection{Parent record}

This V84 note was formed by appending the preceding material to the
validated V83 source, whose record was
\[
\begin{array}{c|c}
\text{lines}&25384\\
\text{bytes}&812304\\
\text{SHA--256}&
\texttt{683E10071AD6EEC32877C669BC21E947D098218E8EBB8C7EF91F0082769B297D}.
\end{array}                                                     \tag{84.14}
\]

% =============================================================================
% END APPENDED ELEVENTH-SERIES V84 MATERIAL
% =============================================================================
% =============================================================================
% BEGIN APPENDED ELEVENTH-SERIES V85 MATERIAL
% =============================================================================

\section{Companion audit at V85: unchanged endpoints and an Einstein pivot}
\label{sec:v11v85}

The compulsory audit finds no change since V80:
\[
\begin{array}{c|r|r|l}
\text{file}&\text{lines}&\text{bytes}&\text{SHA--256}\\ \hline
\texttt{Renewal\_Yang\_Mills\_Mass\_Gap\_Research\_Notes\_V40.tex}
&16447&560619&
\texttt{F0B16C4DCCF204359A62342A87322FAE793567A22E3EF63AC813FEB4DF51951B}\\
\texttt{Renewal\_NS\_Stability\_Research\_Notes\_V26.tex}
&5319&278283&
\texttt{82C887F8C2C69E4F28FAD7C3DC8DE5B9099CB5A4BF431EBA1FFF3E91935978AD}.
\end{array}                                                     \tag{85.1}
\]
The YM endpoint still stops before the fully measured Walsh batch,
internal alias control, sector estimates, continuum measure, and
clustering.  The NS endpoint still supplies only rank-one Oseen-kernel
constants for its own stability ledger.  Neither adds an operator
shell estimate, determinant-line period calculation, or coupled
gravity input.

Versions V81--V84 now give a complete conditional route from a
zero-free relative operator and local counterterm to a globally
trivialised, summably transported determinant phase.  Continuing to
restate that route without a regulated carrier would be circular.
The programme therefore moves to an earlier coupled-theory
obstruction.

\begin{decision}[V86 Einstein positive-measure test]
\label{dec:v11v85-next}
V86 will analyse the Euclidean Einstein--Hilbert conformal direction.
It will prove the high-frequency conformal unboundedness on a compact
background and show why the positive amplitude hypotheses used in
V76--V79 cannot be applied to the bare Einstein action.  It will then
state the coercivity and contour data that any renewal regulator must
supply before the gauge and determinant estimates can be coupled to
gravity.
\end{decision}

\begin{firewall}[Audit boundary]
\label{fire:v11v85-boundary}
The audit supplies no new physical estimate.  In particular, the
conditional phase theorem is not a construction of the Einstein
reference law, and the local Yang--Mills chart cannot control the
gravitational conformal mode.
\end{firewall}

\begin{assessment}[Progress after V85]
\label{ass:v11v85-status}
The required cross-program check is complete.  Because the companion
endpoints are unchanged, the next noncircular step is the positive
Einstein amplitude on which all sector, reweighting, and phase
expectations depend.
\end{assessment}

\subsection{Parent record}

This V85 note was formed by appending the preceding material to the
validated V84 source, whose record was
\[
\begin{array}{c|c}
\text{lines}&25666\\
\text{bytes}&823382\\
\text{SHA--256}&
\texttt{E07429B751E098294B843020786640280D40C7CCCC792BD67218C39300919144}.
\end{array}                                                     \tag{85.2}
\]

% =============================================================================
% END APPENDED ELEVENTH-SERIES V85 MATERIAL
% =============================================================================
% =============================================================================
% BEGIN APPENDED ELEVENTH-SERIES V86 MATERIAL
% =============================================================================

\section{The Euclidean Einstein conformal obstruction}
\label{sec:v11v86}

The probability and reweighting theorems of V76--V79 require a
normalised positive amplitude law.  For the bare Euclidean
Einstein--Hilbert action, that input fails already inside one fixed
topology and one conformal class.

Let \((M,\bar g)\) be a closed smooth Riemannian manifold of dimension
\(d\ge3\), and define
\[
 S_{\rm EH}^{E}(g)
 =-{1\over16\pi G}\int_M(R_g-2\Lambda)\,dV_g,
                    \qquad G>0.                             \tag{86.1}
\]
For \(\phi\in C^\infty(M,\mathbb R)\), put
\(g_\phi=e^{2\phi}\bar g\).

\begin{lemma}[Exact conformal formula]
\label{lem:v11v86-conformal}
One has
\[
 \begin{split}
 S_{\rm EH}^{E}(g_\phi)
 =-{1\over16\pi G}\bigg[
 &\int_Me^{(d-2)\phi}
   \left\{\bar R+(d-1)(d-2)|\bar\nabla\phi|^2\right\}
                         dV_{\bar g}\\
 &-2\Lambda\int_Me^{d\phi}\,dV_{\bar g}\bigg].               \tag{86.2}
 \end{split}
\]
\end{lemma}

\begin{proof}
The conformal scalar-curvature identity is
\[
 R_{e^{2\phi}\bar g}
 =e^{-2\phi}\left[
 \bar R-2(d-1)\bar\Delta\phi
 -(d-1)(d-2)|\bar\nabla\phi|^2\right],                       \tag{86.3}
\]
while \(dV_{g_\phi}=e^{d\phi}dV_{\bar g}\).  Since \(M\) has no
boundary,
\[
 \int_M-2(d-1)e^{(d-2)\phi}\bar\Delta\phi\,dV_{\bar g}
 =2(d-1)(d-2)\int_Me^{(d-2)\phi}
                    |\bar\nabla\phi|^2dV_{\bar g}.           \tag{86.4}
\]
Combining (86.3)--(86.4) leaves one positive copy of the gradient
term inside the square bracket in (86.2).  The overall Euclidean
minus sign makes that conformal kinetic term negative.
\end{proof}

\begin{theorem}[High-frequency conformal unboundedness]
\label{thm:v11v86-unbounded}
For every nonempty coordinate ball in \(M\), there is a sequence
\(\phi_n\), supported there up to an additive constant, such that
\[
 \sup_n\|\phi_n\|_{L^\infty}<\infty,\qquad
 \int_Me^{d\phi_n}\,dV_{\bar g}
             =\int_MdV_{\bar g},                            \tag{86.5}
\]
and
\[
                         S_{\rm EH}^{E}(g_{\phi_n})
                         \longrightarrow-\infty.             \tag{86.6}
\]
All metrics \(g_{\phi_n}\) lie on the same smooth manifold and in the
same conformal path component.
\end{theorem}

\begin{proof}
Choose coordinates on a relatively compact ball and a nonzero smooth
bump \(\eta\) supported there.  For a fixed \(a>0\), set
\[
                       \psi_n(x)=a\eta(x)\sin(nx^1).          \tag{86.7}
\]
Metric equivalence on the coordinate support, the inequality
\(|u+v|^2\ge\frac12|u|^2-|v|^2\), and
\[
 \int_M\eta^2\cos^2(nx^1)\,dV_{\bar g}
            \longrightarrow\frac12\int_M\eta^2\,dV_{\bar g} \tag{86.8}
\]
give constants \(c,C>0\) such that
\[
                         \int_M|\bar\nabla\psi_n|^2dV_{\bar g}
                         \ge ca^2n^2-C.                      \tag{86.9}
\]
Equation (86.8) follows from
\(\cos^2 t=(1+\cos2t)/2\) and the Riemann--Lebesgue lemma in the
coordinate chart.

Choose
\[
 c_n={1\over d}\log
 { \int_MdV_{\bar g}\over
   \int_Me^{d\psi_n}dV_{\bar g}},\qquad
 \phi_n=\psi_n+c_n.                                         \tag{86.10}
\]
Since \(\|\psi_n\|_\infty\le a\|\eta\|_\infty\), the constants \(c_n\)
are uniformly bounded.  This proves (86.5) and the first part of
(86.5).  The weights \(e^{(d-2)\phi_n}\) are uniformly bounded above
and below by positive constants.  Hence (86.9) implies
\[
 \int_Me^{(d-2)\phi_n}|\bar\nabla\phi_n|^2dV_{\bar g}
                         \ge c_1n^2-C_1.                    \tag{86.11}
\]
The scalar-curvature term containing \(\bar R\) is uniformly bounded
because \(\phi_n\) is uniformly bounded.  The cosmological term is
constant by (86.5).  Insert (86.11) into (86.2) to obtain
\(S_{\rm EH}^{E}(g_{\phi_n})\le-c_2n^2+C_2\), proving (86.6).
The path \(t\mapsto e^{2t\phi_n}\bar g\) consists of smooth positive
metrics, so topology and conformal path component do not change.
\end{proof}

\begin{corollary}[Sector energy cannot repair the bare action]
\label{cor:v11v86-sector}
No estimate of the form
\[
                         S_{\rm EH}^{E}(g)\ge\tau|Q(g)|-C    \tag{86.12}
\]
with finite \(C\) can hold on the full conformal orbit of a fixed
topological sector, for any charge \(Q\) constant on that orbit.
\end{corollary}

\begin{proof}
The sequence in Theorem~\ref{thm:v11v86-unbounded} has a fixed charge
and action tending to \(-\infty\), contradicting (86.12).
\end{proof}

This does not contradict the gauge-sector energy--entropy mechanism
of V76.  It says that the gravitational factor cannot be appended to
that positive measure without a separate treatment of the conformal
direction.

\subsection{A finite-cutoff divergence diagnostic}

Unboundedness alone does not logically prove divergence of every
possible regulated integral; the reference measure can shrink near
the bad sequence.  The following statement gives the additional
condition under which the cutoff instability is quantitative.

\begin{proposition}[Mode-slab lower bound]
\label{prop:v11v86-slab}
Suppose a cutoff includes the conformal coordinate
\(\phi=t\eta\sin(nx^1)+c(t,n)\), with fixed volume as in (86.10), and
its reference measure has density at least \(m_0>0\) on
\(t\in[a,2a]\).  If all remaining coordinates range over a set of
reference mass at least \(m_1>0\), and their contribution to the
action is bounded above there by \(C_1\), then the cutoff partition
function satisfies
\[
                         Z_n\ge m_0m_1a\,
                         e^{c n^2-C}                         \tag{86.13}
\]
for constants \(c,C>0\).  In particular these partition functions
cannot have a uniform upper bound as \(n\to\infty\).
\end{proposition}

\begin{proof}
The proof of Theorem~\ref{thm:v11v86-unbounded} is uniform for
\(t\in[a,2a]\), after reducing the positive constant in (86.11).
Thus the action on the declared slab is at most \(-cn^2+C\).
Integrating \(e^{-S}\) over a set of reference mass at least
\(m_0m_1a\) proves (86.13).
\end{proof}

The hypothesis on the remaining coordinates prevents a positive
interaction from cancelling the conformal instability on the chosen
slab.  It must be checked in a coupled regulator rather than assumed.

\subsection{Why common shortcuts do not close the problem}

\begin{enumerate}
\item Fixing total volume removes the cosmological variation but not
      the local high-frequency sequence, as (86.5)--(86.6) show.
\item Diffeomorphism gauge fixing does not remove a general conformal
      rescaling; the scalar curvature and local volume density change.
\item Compactifying every finite lattice variable makes each
      finite-cutoff integral finite, but the lower action bound can
      deteriorate like \(n^2\), as (86.13) records.
\item Rotating the conformal contour can replace the divergent
      direction by an oscillatory complex integral, but then the
      positive probability framework of V76--V79 and reflection
      positivity require new proofs.
\item Adding a positive higher-derivative term can dominate high
      frequencies, but it changes the regulated action; removal of
      that term must have uniform estimates and may reintroduce the
      instability.
\end{enumerate}

\subsection{Required replacement datum}

A constructive continuum route must declare one of the following
mathematical objects before coupling the earlier gauge and fermion
estimates:
\begin{enumerate}
\item a positive gravitational reference law with a coercive
      conformal quadratic form, plus an interacting
      Radon--Nikodym density satisfying the V78 \(L^r\) and lower
      normalisation bounds;
\item a rigorously specified complex conformal contour, with
      projective convergence, noncancellation, and an independent
      reconstruction theorem;
\item a modified ultraviolet gravitational action with uniform
      tightness and a controlled flow to a stated continuum fixed
      point.
\end{enumerate}
For the first route, a typical finite-cutoff coercivity input is
\[
 S_i^{\rm grav}(\phi,h^\perp)
 \ge a\|\phi\|_{H^s_i}^2+S_i^\perp(h^\perp)-C,
                        \qquad a>0,                          \tag{86.14}
\]
with covariance tails tight in a fixed negative Sobolev space.
Equation (86.14) is not a property of the bare action (86.1).

\begin{warning}[Analytic continuation is additional structure]
\label{warn:v11v86-contour}
A formal substitution \(\phi\mapsto i\phi\) does not by itself define
a common integration cycle for nonlinear metrics, avoid determinant
zeros, or preserve the reflection operation.  The contour and its
homology class must be stable under cutoff refinement, and its thimble
phases enter the denominator budget of V79.
\end{warning}

\begin{firewall}[Einstein measure remains open]
\label{fire:v11v86-boundary}
V86 proves an obstruction, not a construction.  No attached source
provides the coercive law (86.14), a controlled complex contour, or a
higher-derivative fixed-point measure with the required projective,
reflection, and clustering estimates.  Therefore the full
SM--Einstein continuum is not yet obtained.
\end{firewall}

\begin{decision}[V87 acquisition]
\label{dec:v11v86-next}
V87 will develop the first replacement route: a coercive Gaussian
conformal reference with trace-compatible covariance.  It will prove
tightness, transfer under \(L^r\) reweighting, and an explicit
adjacent covariance criterion, while keeping equivalence to the
Einstein interaction as a separate unresolved renormalisation step.
\end{decision}

\begin{assessment}[Progress after V86]
\label{ass:v11v86-status}
The bare Euclidean Einstein--Hilbert action has been ruled out as the
unexamined positive base measure: bounded, fixed-volume,
fixed-topology conformal oscillations drive it to minus infinity.
The continuum programme must now supply a coercive reference,
a controlled contour, or a modified ultraviolet action before its
gauge, sector, and determinant estimates can be assembled.
\end{assessment}

\subsection{Parent record}

This V86 note was formed by appending the preceding material to the
validated V85 source, whose record was
\[
\begin{array}{c|c}
\text{lines}&25742\\
\text{bytes}&826586\\
\text{SHA--256}&
\texttt{89D63BEC9BE1FE11F9AE2E6E85B5C9F8C1484842621FB301E856B35BF591B975}.
\end{array}                                                     \tag{86.15}
\]

% =============================================================================
% END APPENDED ELEVENTH-SERIES V86 MATERIAL
% =============================================================================
% =============================================================================
% BEGIN APPENDED ELEVENTH-SERIES V87 MATERIAL
% =============================================================================

\section{A coercive conformal reference and its continuum law}
\label{sec:v11v87}

V86 shows that a positive construction cannot start from the bare
Euclidean Einstein--Hilbert conformal action.  This version builds a
coercive Gaussian reference and states the exact condition under
which an interacting positive amplitude remains tight after
reweighting.

Let \({\cal H}\) be a real separable Hilbert space with orthonormal
basis \((e_n)_{n\ge1}\), and let \(\lambda_n\ge1\) increase to
infinity.  For \(s>0\), define
\[
 \|x\|_{-s}^2=\sum_{n\ge1}\lambda_n^{-s}|x_n|^2,\qquad
 x_n=(x,e_n),                                                \tag{87.1}
\]
and let \({\cal H}^{-s}\) be the completion.  Choose \(c_n\ge0\) and
independent standard real Gaussians \(\xi_n\).  At cutoff \(N_i\), put
\[
                       \phi_i=\sum_{n\le N_i}c_n\xi_ne_n.     \tag{87.2}
\]

\begin{theorem}[Gaussian projective continuum]
\label{thm:v11v87-gaussian}
Suppose
\[
                              \sum_{n\ge1}\lambda_n^{-s}c_n^2
                              <\infty.                        \tag{87.3}
\]
Then:
\begin{enumerate}
\item \(\phi_i\) converges in
      \(L^2(\Omega;{\cal H}^{-s})\) and almost surely to a random
      \(\phi\in{\cal H}^{-s}\);
\item the laws \(\nu_i\) are exactly projective under coordinate
      truncation;
\item the embedded laws are tight on \({\cal H}^{-s}\) and converge
      weakly to the centred Gaussian law \(\nu\) of \(\phi\);
\item their covariance tails satisfy
\[
 \mathbb E\|\phi-\phi_i\|_{-s}^2
       =\sum_{n>N_i}\lambda_n^{-s}c_n^2\longrightarrow0.     \tag{87.4}
\]
\end{enumerate}
\end{theorem}

\begin{proof}
For \(j>i\),
\[
 \mathbb E\|\phi_j-\phi_i\|_{-s}^2
 =\sum_{N_i<n\le N_j}\lambda_n^{-s}c_n^2.
\]
Condition (87.3) makes the Gaussian series Cauchy in
\(L^2(\Omega;{\cal H}^{-s})\).  The standard independent-series
criterion, applied after (87.3), gives almost sure convergence as
well.  Coordinate truncation of (87.2) exactly recovers every earlier
law.  Convergence in probability implies convergence in distribution
and tightness.  Orthogonality gives (87.4).
\end{proof}

If \(c_n^2=a_n^{-1}\), the reference action is
\[
                         S_i^{(0)}(\phi)
                         ={1\over2}\sum_{n\le N_i}a_n\phi_n^2.
                                                                    \tag{87.5}
\]
It is coercive in every retained conformal mode.  When
\(\lambda_n\) are eigenvalues of \(1+\bar\Delta\),
\(a_n\asymp\lambda_n^p\), and Weyl growth holds in dimension \(d\),
condition (87.3) follows from
\[
                                      p+s>{d\over2}.          \tag{87.6}
\]
This is a reference regularity condition, not an identification of
the limiting dynamics.

\subsection{Uniform Gaussian exponential integrability}

Let \(Q\) be the covariance of \(\phi\) as an
\({\cal H}^{-s}\)-valued Gaussian.  Its eigenvalues in the normalised
coordinates of (87.1) are
\[
                              q_n=\lambda_n^{-s}c_n^2.        \tag{87.7}
\]
Condition (87.3) says \(Q\) is trace class.

\begin{theorem}[Quadratic exponential moment]
\label{thm:v11v87-exponential}
For
\[
                              0\le\alpha<{1\over2\|Q\|},     \tag{87.8}
\]
one has uniformly in \(i\)
\[
 \mathbb E_{\nu_i}e^{\alpha\|\phi_i\|_{-s}^2}
 =\prod_{n\le N_i}(1-2\alpha q_n)^{-1/2}
 \le\prod_{n\ge1}(1-2\alpha q_n)^{-1/2}<\infty.              \tag{87.9}
\]
\end{theorem}

\begin{proof}
The one-dimensional Gaussian integral gives the finite product.
Choose \(\rho<1\) with \(2\alpha q_n\le\rho\) for all \(n\).
Since
\[
                         -\log(1-x)\le{x\over1-\rho}
                         \quad(0\le x\le\rho),
\]
the logarithm of the infinite product is bounded by a constant times
\(\sum_nq_n<\infty\).
\end{proof}

\subsection{Positive interacting amplitude}

Let \(V_i\) be a real measurable interaction and define
\[
 g_i=e^{-V_i},\qquad
 Z_i=\mathbb E_{\nu_i}g_i,\qquad
 \mu_i={g_i\over Z_i}\nu_i.                                 \tag{87.10}
\]

\begin{theorem}[Relative lower bound gives \(L^r\) control]
\label{thm:v11v87-relative}
Fix \(r>1\).  Suppose constants \(\alpha,C,C_0,m_0\), independent of
\(i\), satisfy
\[
 V_i(\phi)\ge-\alpha\|\phi\|_{-s}^2-C,                       \tag{87.11}
\]
\[
 \nu_i\{V_i\le C_0\}\ge m_0>0,                              \tag{87.12}
\]
and
\[
                                      r\alpha<{1\over2\|Q\|}.\tag{87.13}
\]
Then
\[
 \|g_i\|_{L^r(\nu_i)}
 \le e^C
 \left(\mathbb E_{\nu_i}
       e^{r\alpha\|\phi_i\|_{-s}^2}\right)^{1/r}
 \le G_r<\infty,                                            \tag{87.14}
\]
and
\[
                                      Z_i\ge m_0e^{-C_0}=:z_0.
                                                                    \tag{87.15}
\]
Consequently the laws \(\mu_i\) are uniformly tight on
\({\cal H}^{-s}\).
\end{theorem}

\begin{proof}
Equation (87.11) gives
\(g_i^r\le e^{rC}e^{r\alpha\|\phi_i\|_{-s}^2}\);
Theorem~\ref{thm:v11v87-exponential} and (87.13) give (87.14).
On the event in (87.12), \(g_i\ge e^{-C_0}\), proving (87.15).
For a compact-tightness set \(K\), Theorem~\ref{thm:v11v78-holder}
gives
\[
 \mu_i(K^c)\le {G_r\over z_0}\nu_i(K^c)^{1/r'}.
\]
The reference laws are tight by
Theorem~\ref{thm:v11v87-gaussian}; choose \(K\) to make the right side
arbitrarily small.
\end{proof}

The comparison norm in (87.11) is deliberately explicit.  A negative
interaction quadratic in a stronger norm cannot be controlled by
(87.9).  The high-frequency conformal term of V86 is exactly such a
failure for an inadequately coercive reference.

\subsection{Adjacent interacting convergence}

Under exact Gaussian projection, define the conditional density
\[
 \bar f_i
 =\mathbb E_{\nu_{i+1}}
       \left[{e^{-V_{i+1}}\over Z_{i+1}}\mid P_i\phi\right],
 \qquad
 f_i={e^{-V_i}\over Z_i}.                                   \tag{87.16}
\]

\begin{theorem}[Tight projective gravitational law]
\label{thm:v11v87-projective}
Assume the hypotheses of
Theorem~\ref{thm:v11v87-relative} and
\[
                         \sum_i\|\bar f_i-f_i\|_{L^1(\nu_i)}
                         <\infty.                            \tag{87.17}
\]
Then every bounded cylinder expectation under \(\mu_i\) converges.
The embedded laws converge weakly on \({\cal H}^{-s}\) to a unique
probability measure \(\mu\).
\end{theorem}

\begin{proof}
Theorem~\ref{thm:v11v78-martingale} and (87.17) make every bounded
cylinder expectation Cauchy.  Uniform tightness follows from
Theorem~\ref{thm:v11v87-relative}, so every subsequence has a weakly
convergent further subsequence.  All such limits have the same
finite-coordinate distributions, because cylinder expectations
converge.  Coordinate projections along a complete orthonormal basis
generate the Borel sigma algebra of the separable Hilbert space, so
the subsequential limit is unique.  Tightness plus uniqueness implies
convergence of the full sequence.
\end{proof}

\subsection{Approximate covariance refinement}

Exact coefficients in (87.2) are the cleanest choice.  For a varying
triangular array \(c_{i,n}\), couple adjacent Gaussians with the same
\(\xi_n\), extend the coarse field by zero, and put
\[
 \epsilon_i^2
 =\sum_{n\le N_i}\lambda_n^{-s}
          (c_{i+1,n}-c_{i,n})^2
  +\sum_{N_i<n\le N_{i+1}}\lambda_n^{-s}c_{i+1,n}^2.         \tag{87.18}
\]

\begin{proposition}[Coupled covariance defect]
\label{prop:v11v87-covariance}
For every bounded \(1\)-Lipschitz function \(F\) on
\({\cal H}^{-s}\),
\[
                 |\mathbb E F(\phi_{i+1})
                   -\mathbb E F(\phi_i)|
                 \le\epsilon_i.                             \tag{87.19}
\]
Hence \(\sum_i\epsilon_i<\infty\) is a sufficient bounded-Lipschitz
Cauchy criterion for the reference laws.
\end{proposition}

\begin{proof}
Under the declared coupling,
\[
 \mathbb E\|\phi_{i+1}-\phi_i\|_{-s}^2=\epsilon_i^2.
\]
The Lipschitz property and Cauchy--Schwarz give (87.19).
\end{proof}

The new-mode term in (87.18) is absent from the pushforward comparison
of old coordinates but present when both cutoffs are embedded in one
field space.  These two notions must not be conflated.

\subsection{Coupling to the remaining fields}

For gauge, fermion, Higgs, ghost, and transverse metric variables
\(y_i\), the reference can be conditional:
\[
                         \nu_i(d\phi,dy)
                         =\nu_i^\phi(d\phi)\nu_i^y(dy\mid\phi).
                                                                    \tag{87.20}
\]
Theorems V78 and V87 remain applicable if their \(L^r\), lower
normalisation, tightness, and conditional-density estimates hold for
the joint law.  Uniform conformal control alone does not imply
tightness of \(y_i\).

\begin{warning}[Reference theory versus Einstein theory]
\label{warn:v11v87-equivalence}
The Gaussian law (87.2) solves a measure-existence problem only for
the declared coercive reference.  To recover an Einstein continuum,
one must prove that the renormalised interactions \(V_i\) satisfy
(87.11)--(87.17) and that their limiting Ward identities, relevant
couplings, and observables define the intended gravitational
dynamics.  Choosing \(a_n\) large enough for tightness is not itself
a universality proof.
\end{warning}

\begin{firewall}[Unacquired relative bound]
\label{fire:v11v87-boundary}
The bare conformal Einstein interaction from V86 does not presently
satisfy the uniform relative lower bound (87.11) with a demonstrated
choice of covariance, nor is the adjacent density series (87.17)
known.  Gauge fixing, Faddeev--Popov determinants, metric positivity,
and quotient geometry have not yet been incorporated into the joint
law (87.20).
\end{firewall}

\begin{decision}[V88 acquisition]
\label{dec:v11v87-next}
V88 will address the next prerequisite for (87.20): a quantitative
local diffeomorphism slice.  It will use an elliptic gauge carrier,
an inf--sup or spectral gap, and a Faddeev--Popov determinant bound to
separate genuine metric variables from gauge directions without
assuming a global slice.
\end{decision}

\begin{assessment}[Progress after V87]
\label{ass:v11v87-status}
A coercive conformal Gaussian reference now has an exact projective
continuum, Sobolev tightness, and uniform quadratic exponential
moments.  A relative lower interaction bound and a positive comparison
set transfer tightness, while a summable conditional-density defect
gives a unique interacting limit.  Establishing those hypotheses for
renormalised Einstein dynamics remains the central open step.
\end{assessment}

\subsection{Parent record}

This V87 note was formed by appending the preceding material to the
validated V86 source, whose record was
\[
\begin{array}{c|c}
\text{lines}&26016\\
\text{bytes}&837301\\
\text{SHA--256}&
\texttt{94039775964E51276706A1571A7D40F6089BD7D1843A99CCEA546CA0820EEB9C}.
\end{array}                                                     \tag{87.21}
\]

% =============================================================================
% END APPENDED ELEVENTH-SERIES V87 MATERIAL
% =============================================================================
% =============================================================================
% BEGIN APPENDED ELEVENTH-SERIES V88 MATERIAL
% =============================================================================

\section{A quantitative local diffeomorphism slice}
\label{sec:v11v88}

The joint law in (87.20) requires a quotient of metric variables by
diffeomorphisms.  This version gives a finite-cutoff local slice
criterion and a dimension-sensitive Faddeev--Popov bound.  It does not
assume that the local slices form one global gauge.

Let \(E_i\) be a finite-dimensional metric-coordinate space,
\(V_i\) a finite-dimensional infinitesimal diffeomorphism space after
removing declared stabiliser directions, and
\[
                         {\cal T}_i(\xi,h)\in E_i             \tag{88.1}
\]
a \(C^1\) local action with \({\cal T}_i(0,h)=h\).  Let
\(C_i:E_i\to V_i\) be a gauge condition.  Near a reference
\(h_0\), define
\[
 F_h(\xi)=C_i({\cal T}_i(\xi,h)),\qquad
 M_0=D_\xi F_{h_0}(0).                                      \tag{88.2}
\]
The matrix \(M_0\) is the finite-cutoff Faddeev--Popov carrier.

\begin{theorem}[Quantitative slice by contraction]
\label{thm:v11v88-slice}
Suppose
\[
                         s_{\min}(M_0)\ge\gamma>0.            \tag{88.3}
\]
Let \(B_r\subset V_i\) be the closed radius-\(r\) ball.  For a given
\(h\), assume
\[
 \sup_{\xi\in B_r}\|D_\xi F_h(\xi)-M_0\|\le\eta<\gamma,       \tag{88.4}
\]
and
\[
                              \|F_h(0)\|\le(\gamma-\eta)r.    \tag{88.5}
\]
Then there is a unique \(\xi(h)\in B_r\) satisfying
\[
                              C_i({\cal T}_i(\xi(h),h))=0.    \tag{88.6}
\]
Moreover,
\[
                              \|\xi(h)\|
                              \le{\|F_h(0)\|\over\gamma-\eta}.\tag{88.7}
\]
\end{theorem}

\begin{proof}
Define
\[
                              \Psi_h(\xi)=\xi-M_0^{-1}F_h(\xi).
                                                                    \tag{88.8}
\]
By (88.3)--(88.4),
\[
                              \|D\Psi_h(\xi)\|
                              \le q:={\eta\over\gamma}<1.     \tag{88.9}
\]
Also
\[
 \|\Psi_h(\xi)\|
 \le\|M_0^{-1}F_h(0)\|+q\|\xi\|
 \le(1-q)r+qr=r,
\]
where (88.5) was used.  Hence \(\Psi_h\) is a contraction of \(B_r\)
into itself and has a unique fixed point.  The fixed-point equation
is (88.6).  Finally,
\[
 \|\xi(h)\|
 \le\|M_0^{-1}F_h(0)\|+q\|\xi(h)\|,
\]
which gives (88.7).
\end{proof}

The theorem is uniform only if \(\gamma\), \(\eta/\gamma\), and the
chart radius do not deteriorate with the cutoff.  Killing fields and
other stabilisers make \(\gamma=0\) unless their directions are
removed or treated as an explicit compact residual group.

\subsection{Faddeev--Popov determinant control}

At a slice point \(h\), write
\[
                         M(h)=D_\xi F_h(0),\qquad
                         R(h)=(M(h)-M_0)M_0^{-1}.            \tag{88.10}
\]

\begin{theorem}[Relative determinant bound and sign]
\label{thm:v11v88-FP}
Suppose
\[
                              \|R(h)\|\le q<1.                \tag{88.11}
\]
Then \(M(h)\) is invertible, the real determinant
\(\det M(h)\) has the same sign as \(\det M_0\), and
\[
 \left|\log\left|
       {\det M(h)\over\det M_0}\right|\right|
 \le{\|R(h)\|_1\over1-q}.                                   \tag{88.12}
\]
Consequently, if \(\|R(h)\|_1\le B\),
\[
 e^{-B/(1-q)}
 \le\left|{\det M(h)\over\det M_0}\right|
 \le e^{B/(1-q)}.                                           \tag{88.13}
\]
\end{theorem}

\begin{proof}
The segment \(I+tR\) is invertible because
\(\|tR\|<1\).  Hence the real determinant cannot change sign along
\(M(t)=(I+tR)M_0\).  The convergent trace expansion gives
\[
 \log\det(I+R)=\sum_{n\ge1}{(-1)^{n+1}\over n}\operatorname{Tr}(R^n)
\]
with the continuous real branch along the segment.  Since
\[
 |\operatorname{Tr}(R^n)|
 \le\|R\|_1\|R\|^{n-1},
\]
its absolute value is at most
\[
 \|R\|_1\sum_{n\ge1}{q^{n-1}\over n}
 \le{\|R\|_1\over1-q}.
\]
Exponentiation proves (88.13).
\end{proof}

An operator-norm bound alone gives
\(\|R\|_1\le(\dim V_i)\|R\|\), which is extensive.  A uniform
Radon--Nikodym bound for the quotient density therefore requires a
relative trace-ideal estimate or an exact local determinant
counterterm, just as in V81--V83.

\begin{corollary}[Local \(L^r\) quotient-density input]
\label{cor:v11v88-density}
On a chart where \(\|R(h)\|_1\le B\), the relative positive
Faddeev--Popov weight
\[
                              w_i^{\rm FP}(h)
                              =\left|{\det M(h)\over\det M_0}\right|
                                                                    \tag{88.14}
\]
and its reciprocal are uniformly bounded by \(e^{B/(1-q)}\).
Multiplying a reference density by \(w_i^{\rm FP}\) therefore
preserves every finite \(L^r\) bound up to this factor.
\end{corollary}

\begin{proof}
Use (88.13) and the elementary inequality
\(\|wf\|_{L^r}\le\|w\|_\infty\|f\|_{L^r}\).
\end{proof}

\subsection{Local quotient formula}

Assume the map
\[
                    {\cal P}_i:S_i\times B_r\longrightarrow U_i,
              \qquad {\cal P}_i(s,\xi)={\cal T}_i(\xi,s),    \tag{88.15}
\]
is a \(C^1\) diffeomorphism from a slice chart onto its image.  The
ordinary change-of-variables theorem gives
\[
 \int_{U_i}f(h)\,dh
 =\int_{S_i\times B_r}
       f({\cal T}_i(\xi,s))
       |\det D{\cal P}_i(s,\xi)|\,ds\,d\xi.                  \tag{88.16}
\]
For a gauge-invariant \(f\), the orbit integral can be separated only
after the full Jacobian in (88.16) has been computed.  The
Faddeev--Popov determinant is one factor of this Jacobian; the orbit
metric, stabiliser volume, and coordinate density are additional
factors.  This is the gravitational analogue of the complete-density
warning imported from YM V39--V40.

\begin{proposition}[Local Faddeev--Popov identity]
\label{prop:v11v88-delta}
Suppose \(F_h:B_r\to V_i\) has the unique zero \(\xi(h)\) of
Theorem~\ref{thm:v11v88-slice}, and its derivative is invertible
there.  For every continuous test function \(a\) supported in \(B_r\),
\[
 \int_{B_r}a(\xi)\,
       \delta(F_h(\xi))\,|\det D_\xi F_h(\xi)|\,d\xi
                         =a(\xi(h)).                         \tag{88.17}
\]
\end{proposition}

\begin{proof}
Apply the finite-dimensional change-of-variables formula to
\(y=F_h(\xi)\) near its unique zero.  Contributions away from the
zero vanish against \(\delta(y)\), and the Jacobian cancels.
\end{proof}

Equation (88.17) is local.  If an orbit meets the gauge slice more
than once, summing over all zeros counts the copies with their signs
or absolute Jacobians.  Declaring the right side to be one globally
would assume away the Gribov problem.

\subsection{Elliptic source of the gap}

For a de Donder-type metric gauge at a smooth background, the
continuum principal part of \(M_0\) is a vector Laplacian.  A
finite-cutoff estimate of the form
\[
 \|M_0\xi\|_{L^2}\ge\gamma\|\xi\|_{H^2}                     \tag{88.18}
\]
on the orthogonal complement of Killing fields supplies (88.3) after
norm calibration.  Lower curvature terms can be absorbed when their
relative bound is below one.  The constant \(\gamma\) can fail
uniformly on degenerating geometries, long necks, or backgrounds
approaching new isometries.

\begin{ledger}[Local gravitational quotient inputs]
\label{led:v11v88-inputs}
A usable cutoff chart must provide:
\begin{enumerate}
\item a declared stabiliser decomposition;
\item an elliptic lower bound such as (88.18);
\item a nonlinear derivative variation \(\eta<\gamma\);
\item a residual gauge bound (88.5);
\item injectivity of the product chart (88.15);
\item relative trace-class control in (88.12);
\item the remaining orbit and coordinate Jacobian factors;
\item refinement compatibility of all these data.
\end{enumerate}
\end{ledger}

\begin{warning}[Local uniqueness is not a global slice]
\label{warn:v11v88-gribov}
The contraction theorem excludes a second solution only inside
\(B_r\).  It neither excludes a distant diffeomorphism returning to
the slice nor proves that every orbit meets the selected atlas.
Global quotient construction requires an atlas with multiplicity and
overlap control.
\end{warning}

\begin{firewall}[No uniform gravitational slice acquired]
\label{fire:v11v88-boundary}
The current sources do not supply a cutoff-uniform diffeomorphism
Faddeev--Popov gap, trace-norm relative determinant bound, global
orbit atlas, or refinement-compatible stabiliser treatment on the
support of the coercive reference law from V87.  V88 is the local
criterion, not its verification for the coupled theory.
\end{firewall}

\begin{decision}[V89 acquisition]
\label{dec:v11v88-next}
V89 will assemble overlapping local slices by a partition of unity
and an orbit-multiplicity correction.  It will prove when the
resulting quotient integral is independent of the slice atlas and
derive the overlap defect that must be summable under refinement.
\end{decision}

\begin{assessment}[Progress after V88]
\label{ass:v11v88-status}
A quantitative Faddeev--Popov gap now yields a unique local gauge
representative, an explicit chart radius, constant determinant sign,
and a dimension-sensitive relative determinant bound.  The result
clarifies exactly which local estimates are needed before the
coercive gravitational reference can descend to a quotient measure.
\end{assessment}

\subsection{Parent record}

This V88 note was formed by appending the preceding material to the
validated V87 source, whose record was
\[
\begin{array}{c|c}
\text{lines}&26335\\
\text{bytes}&848605\\
\text{SHA--256}&
\texttt{3176B95AB8C61ED1655F991634783CFAB6213BB02D005D674FF67DDE2DB4C9A3}.
\end{array}                                                     \tag{88.19}
\]

% =============================================================================
% END APPENDED ELEVENTH-SERIES V88 MATERIAL
% =============================================================================
% =============================================================================
% BEGIN APPENDED ELEVENTH-SERIES V89 MATERIAL
% =============================================================================

\section{Gluing local gravitational slices into a quotient measure}
\label{sec:v11v89}

V88 supplies a locally unique gauge representative and its Jacobian.
A global measure requires overlap consistency and correction for
orbits that meet a gauge condition more than once.  This version gives
the finite-atlas gluing theorem and the resulting refinement defect.

Let \(Q_i=X_i/{\cal G}_i\) be a cutoff orbit space and let
\(\{U_a\}_{a\in A_i}\) be a finite measurable cover by quotient images
of slice charts.  Fix a positive dominating measure \(\lambda_i\) on
\(Q_i\).  Suppose chart \(a\) produces the local quotient density
\(j_a\ge0\) on \(U_a\).

\begin{theorem}[Exact quotient-density gluing]
\label{thm:v11v89-glue}
There is a measure \(\nu_i\) on \(Q_i\) whose restriction to every
\(U_a\) is \(j_a\lambda_i\) if and only if
\[
                              j_a=j_b
        \quad\lambda_i\text{-almost everywhere on }U_a\cap U_b.
                                                                    \tag{89.1}
\]
When (89.1) holds and \((\chi_a)\) is any measurable partition of
unity subordinate to the cover, then for bounded measurable \(F\),
\[
 \int_{Q_i}F\,d\nu_i
 =\sum_a\int_{U_a}\chi_a(q)F(q)j_a(q)\,d\lambda_i(q),        \tag{89.2}
\]
and the result is independent of the partition.
\end{theorem}

\begin{proof}
Necessity follows by restricting the same measure to an overlap.
Under (89.1), define \(j(q)=j_a(q)\) when \(q\in U_a\); this is
well-defined almost everywhere.  Then \(\nu_i=j\lambda_i\), and
\[
 \sum_a\chi_a(q)j_a(q)=j(q)\sum_a\chi_a(q)=j(q)
\]
almost everywhere, proving (89.2) and partition independence.
\end{proof}

The theorem applies only after the full chart Jacobian has been
included.  Agreement of the Faddeev--Popov determinants alone does
not imply (89.1).

\subsection{Orbit multiplicity}

Let a gauge condition on an orbit \(q\) have finitely many regular
zeros \(\xi_1,\ldots,\xi_{N(q)}\) in the declared gauge domain.
The multi-root version of (88.17) is
\[
 \int a(\xi)\delta(F_q(\xi))
             |\det DF_q(\xi)|\,d\xi
                      =\sum_{k=1}^{N(q)}a(\xi_k).            \tag{89.3}
\]

\begin{proposition}[Gribov multiplicity correction]
\label{prop:v11v89-multiplicity}
For a gauge-invariant integrand, (89.3) counts the orbit \(N(q)\)
times.  If \(1\le N(q)<\infty\) is measurable and constant along the
orbit, insertion of
\[
                                      {1\over N(q)}          \tag{89.4}
\]
gives one positive copy of that orbit.  If a finite stabiliser
\(H_q\) remains and the orbit-coordinate measure counts every
stabiliser copy, the additional factor is \(1/|H_q|\).
\end{proposition}

\begin{proof}
Gauge invariance makes the value of the physical integrand identical
at all roots in (89.3), so their sum is \(N(q)\) times that value.
Equation (89.4) cancels the multiplicity.  The same finite counting
argument applies to the stabiliser.
\end{proof}

Using \(\det DF\) without absolute value computes an oriented degree,
which can vanish through cancellations and is not a positive quotient
density.  The signed construction may be useful for topological
identities, but it cannot silently replace the positive amplitude law
of V87.

\begin{warning}[Infinite or jumping multiplicity]
\label{warn:v11v89-multiplicity}
If \(N(q)\) is infinite, nonmeasurable, or unstable under refinement,
(89.4) does not define the required correction.  When roots merge,
the Faddeev--Popov determinant vanishes and the local estimate of V88
also fails.  These configurations belong to a singular atlas stratum
whose probability must be controlled separately.
\end{warning}

\subsection{Approximate overlap consistency}

Suppose exact equality in (89.1) is not yet known.  Define on the
union of overlaps
\[
 D_i(q)=
 \max_{\substack{a,b\\q\in U_a\cap U_b}}|j_a(q)-j_b(q)|,     \tag{89.5}
\]
with \(D_i(q)=0\) where only one chart is active.  For two subordinate
partitions \(\chi\) and \(\widetilde\chi\), let
\[
 j_\chi=\sum_a\chi_aj_a,\qquad
 j_{\widetilde\chi}=\sum_a\widetilde\chi_aj_a.               \tag{89.6}
\]

\begin{theorem}[Atlas-dependence bound]
\label{thm:v11v89-atlas}
One has pointwise almost everywhere
\[
                              |j_\chi-j_{\widetilde\chi}|
                              \le2D_i,                        \tag{89.7}
\]
and hence
\[
 \sup_{\|F\|_\infty\le1}
 \left|\int F(j_\chi-j_{\widetilde\chi})\,d\lambda_i\right|
                       \le2\int D_i\,d\lambda_i.             \tag{89.8}
\]
\end{theorem}

\begin{proof}
At a fixed \(q\), choose any active chart \(b\).  Because both
partitions sum to one,
\[
 j_\chi-j_{\widetilde\chi}
 =\sum_a(\chi_a-\widetilde\chi_a)(j_a-j_b).
\]
Every nonzero summand uses an active chart overlapping \(b\), so
\[
 |j_\chi-j_{\widetilde\chi}|
 \le D_i\sum_a|\chi_a-\widetilde\chi_a|
 \le2D_i.
\]
Integration proves (89.8).
\end{proof}

Thus the overlap mismatch must be controlled in \(L^1\), not merely
pointwise on every fixed field.  If the number of charts grows, the
maximum form (89.5) avoids an artificial chart-count factor, but it
still requires a common dominating quotient measure.

\subsection{Refinement of quotient measures}

Assume the cutoff projection is equivariant:
\[
                        P_i(gx)=\rho_i(g)P_i(x),             \tag{89.9}
\]
so it descends to
\[
                     \bar P_i:Q_{i+1}\longrightarrow Q_i,
                      \qquad \bar P_i([x])=[P_i(x)].         \tag{89.10}
\]
Let \(\nu_i\) be ideal quotient probabilities and
\(\widetilde\nu_i\) the probabilities assembled from the implemented
slice atlas, multiplicity, and stabiliser factors.  Put
\[
 d_{\rm b}(\widetilde\nu_i,\nu_i)\le\epsilon_i,\qquad
 d_{\rm b}((\bar P_i)_\#\nu_{i+1},\nu_i)\le\delta_i.         \tag{89.11}
\]

\begin{theorem}[Quotient refinement error]
\label{thm:v11v89-refine}
Under (89.9)--(89.11),
\[
 d_{\rm b}((\bar P_i)_\#\widetilde\nu_{i+1},
                         \widetilde\nu_i)
                    \le\epsilon_{i+1}+\delta_i+\epsilon_i.  \tag{89.12}
\]
Consequently, if
\[
                              \sum_i(\epsilon_i+\delta_i)<\infty,
                                                                    \tag{89.13}
\]
all bounded quotient-cylinder expectations are Cauchy.
\end{theorem}

\begin{proof}
Pushforward is a contraction for \(d_{\rm b}\), because
\(\|F\circ\bar P_i\|_\infty\le1\) whenever \(\|F\|_\infty\le1\).
Insert the two ideal measures and use the triangle inequality:
\[
 \begin{split}
 d_{\rm b}((\bar P_i)_\#\widetilde\nu_{i+1},\widetilde\nu_i)
 &\le d_{\rm b}(\widetilde\nu_{i+1},\nu_{i+1})\\
 &\quad+d_{\rm b}((\bar P_i)_\#\nu_{i+1},\nu_i)
       +d_{\rm b}(\nu_i,\widetilde\nu_i).
 \end{split}
\]
This is (89.12); summation gives the Cauchy statement.
\end{proof}

A concrete \(\epsilon_i\) can combine:
\[
 \epsilon_i^{\rm ov}
       =2\int D_i\,d\lambda_i,\quad
 \epsilon_i^{\rm mult},\quad
 \epsilon_i^{\rm stab},\quad
 \epsilon_i^{\rm sing},                                    \tag{89.14}
\]
representing overlap, multiplicity, stabiliser-volume, and singular-
stratum errors.  Their sum controls \(\epsilon_i\) after normalisation
has also been bounded away from zero.

\subsection{Interaction with the coercive reference}

Let \(\nu_i^{(0)}\) be the quotient pushforward of the Gaussian
reference from V87 after adding transverse metric variables.  If the
atlas density relative to an ideal local reference is bounded above
and below by V88 and the multiplicity factors satisfy
\[
                         1\le N_i(q)\le N_*<\infty,          \tag{89.15}
\]
then these quotient factors preserve finite \(L^r\) bounds up to a
constant.  The interacting density must still satisfy the joint
relative lower bound and comparison-set condition of V87.

\begin{lemma}[Tightness passes to the quotient]
\label{lem:v11v89-tight}
If \(X_i\) embed in a common Polish space \(X\), the group action is
continuous, \(Q=X/{\cal G}\) is Hausdorff and metrizable, and
\(\pi:X\to Q\) is continuous, then tightness of probabilities
\(\mu_i\) on \(X\) implies tightness of \(\pi_\#\mu_i\) on \(Q\).
\end{lemma}

\begin{proof}
For any \(\varepsilon>0\), choose compact \(K\subset X\) with
\(\mu_i(K)\ge1-\varepsilon\) uniformly.  Its continuous image
\(\pi(K)\) is compact and has quotient probability at least
\(1-\varepsilon\).
\end{proof}

The Hausdorff condition can fail for nonproper group actions or
degenerating metrics.  It is therefore part of the construction, not
a harmless topological convention.

\begin{firewall}[No global diffeomorphism quotient acquired]
\label{fire:v11v89-boundary}
V89 gives the gluing and error formulas, but the SM--Einstein
regulator still lacks a proven finite slice atlas covering a
high-probability set, measurable finite multiplicities, compatible
stabiliser volumes, overlap equality or summable mismatch, and an
equivariant refinement map on quotient spaces.
\end{firewall}

\begin{decision}[Mandatory V90 audit]
\label{dec:v11v89-next}
The next version is the compulsory YM/NS audit.  It will test whether
new fully measured symmetry or kernel information changes the
quotient construction.  V91 will then assemble the current positive
gravity, quotient, sector, and determinant hypotheses into one
conditional continuum theorem and identify the first still
unacquired physical input.
\end{decision}

\begin{assessment}[Progress after V89]
\label{ass:v11v89-status}
Local diffeomorphism slices now glue to an atlas-independent positive
quotient measure exactly when their full densities agree on overlaps.
Gribov copies and stabilisers have explicit counting factors, and
every remaining atlas error propagates into the adjacent projective
budget through (89.12).
\end{assessment}

\subsection{Parent record}

This V89 note was formed by appending the preceding material to the
validated V88 source, whose record was
\[
\begin{array}{c|c}
\text{lines}&26619\\
\text{bytes}&858515\\
\text{SHA--256}&
\texttt{3B6ADB566278A4E44A1A02F9F81B7EF18D4D803C8AE17B3DB3DDA18E66E3734B}.
\end{array}                                                     \tag{89.16}
\]

% =============================================================================
% END APPENDED ELEVENTH-SERIES V89 MATERIAL
% =============================================================================
% =============================================================================
% BEGIN APPENDED ELEVENTH-SERIES V90 MATERIAL
% =============================================================================

\section{Companion audit at V90: no change to the quotient inputs}
\label{sec:v11v90}

The compulsory audit again finds
\[
\begin{array}{c|r|r|l}
\text{file}&\text{lines}&\text{bytes}&\text{SHA--256}\\ \hline
\texttt{Renewal\_Yang\_Mills\_Mass\_Gap\_Research\_Notes\_V40.tex}
&16447&560619&
\texttt{F0B16C4DCCF204359A62342A87322FAE793567A22E3EF63AC813FEB4DF51951B}\\
\texttt{Renewal\_NS\_Stability\_Research\_Notes\_V26.tex}
&5319&278283&
\texttt{82C887F8C2C69E4F28FAD7C3DC8DE5B9099CB5A4BF431EBA1FFF3E91935978AD}.
\end{array}                                                     \tag{90.1}
\]
Both files are byte-for-byte unchanged from V85.

YM V40 still supports the requirement that the full quotient and
coordinate density be recombined before symmetry testing, but its
measured Walsh batch remains an announced next step.  It supplies no
diffeomorphism slice, Gribov multiplicity, gravitational reference
law, sector tail, or global determinant-line holonomy.  NS V26
remains confined to Oseen-kernel rank-one constants and adds no
coupled-continuum input.

\begin{decision}[V91 conditional assembly theorem]
\label{dec:v11v90-next}
V91 will assemble the positive gravitational law of V87, the quotient
defects of V88--V89, the sector laws of V76--V78, and the determinant
phase budget of V79--V84.  The theorem will conclude existence of a
normalised complex cylinder functional under explicit summable
hypotheses.  Reflection positivity, clustering, Ward identities, and
identification with the full SM--Einstein dynamics will remain
separate conclusion gates.
\end{decision}

\begin{firewall}[Audit boundary]
\label{fire:v11v90-boundary}
No companion result verifies a missing hypothesis of the forthcoming
assembly theorem.  The audit changes no claim status.
\end{firewall}

\begin{assessment}[Progress after V90]
\label{ass:v11v90-status}
The required audit is complete.  The unchanged companions leave the
programme ready for a precise conditional assembly statement, which
will prevent measure existence, phase noncancellation, and physical
identification from being conflated.
\end{assessment}

\subsection{Parent record}

This V90 note was formed by appending the preceding material to the
validated V89 source, whose record was
\[
\begin{array}{c|c}
\text{lines}&26904\\
\text{bytes}&869011\\
\text{SHA--256}&
\texttt{402ECBAD2B11DFB37E6670DBF6AE51C7EB50CA57157566D4E36ED49455790742}.
\end{array}                                                     \tag{90.2}
\]

% =============================================================================
% END APPENDED ELEVENTH-SERIES V90 MATERIAL
% =============================================================================
% =============================================================================
% BEGIN APPENDED ELEVENTH-SERIES V91 MATERIAL
% =============================================================================

\section{Conditional assembly of the SM--Einstein cylinder continuum}
\label{sec:v11v91}

Versions V76--V90 developed the missing global-measure chain in
separate pieces.  This version assembles them into one theorem.
The result is deliberately conditional: it proves that the listed
cutoff estimates suffice, while preserving the distinction between a
constructed complex cylinder functional and its identification as the
physical SM--Einstein continuum.

Let
\[
       X_1\xleftarrow{P_1}X_2\xleftarrow{P_2}X_3
       \xleftarrow{}\cdots                                  \tag{91.1}
\]
be Polish cutoff quotient spaces with continuous projections, and let
\[
 X_\infty=\{(x_i):P_ix_{i+1}=x_i\}
              \subset\prod_iX_i                             \tag{91.2}
\]
be their inverse limit.  Let \(\mu_i\) be positive amplitude
probabilities on \(X_i\).  Define
\[
 \eta_i=d_{\rm b}((P_i)_\#\mu_{i+1},\mu_i),                  \tag{91.3}
\]
where \(d_{\rm b}\) is (78.13).

\begin{theorem}[Positive projective limit from summable defects]
\label{thm:v11v91-positive}
If
\[
                                      \sum_i\eta_i<\infty,   \tag{91.4}
\]
then for every fixed \(k\), the coarse marginals
\[
                         \mu_{k}^{(j)}
                  =(P_{k,j})_\#\mu_j,\qquad j\ge k,          \tag{91.5}
\]
converge in total variation to a probability
\(\mu_k^\infty\), and
\[
                         (P_k)_\#\mu_{k+1}^\infty
                         =\mu_k^\infty.                      \tag{91.6}
\]
The consistent family defines a Borel probability
\(\mu_\infty\) on \(X_\infty\).
\end{theorem}

\begin{proof}
Pushforward contracts \(d_{\rm b}\).  Therefore, for \(j>i\ge k\),
\[
 d_{\rm b}(\mu_k^{(j)},\mu_k^{(i)})
 \le\sum_{n=i}^{j-1}\eta_n.                                 \tag{91.7}
\]
Finite signed measures are complete in total variation, so the limit
\(\mu_k^\infty\) exists and has mass one.  Applying the continuous
pushforward to the limit proves (91.6).  For each \(n\), push
\(\mu_n^\infty\) to the first \(n\) coordinates by
\[
 x_n\longmapsto(P_{1,n}x_n,\ldots,P_{n-1,n}x_n,x_n).
\]
Equation (91.6) makes these finite-dimensional laws consistent.
The countable projective-limit theorem gives a probability on the
product.  Every constraint \(P_ix_{i+1}=x_i\) has probability one;
their countable intersection is (91.2), so the measure is supported
on \(X_\infty\).
\end{proof}

The hypotheses leading to (91.4) are now explicit: V87 supplies a
coercive tight reference and \(L^r\) transfer, while V88--V89 supply
local quotient charts and the atlas estimate
\(\epsilon_{i+1}+\delta_i+\epsilon_i\).  Those estimates must be
proved for the same interacting law \(\mu_i\).

\subsection{Sector limit}

Let \(Q_i:X_i\to\mathbb Z\), and let
\({\cal G}_{i+1}\) be a set on which
\[
                              Q_{i+1}=Q_i\circ P_i.           \tag{91.8}
\]
Put
\[
                              \beta_i=\mu_{i+1}({\cal G}_{i+1}^c).
                                                                    \tag{91.9}
\]

\begin{theorem}[Limiting sector law]
\label{thm:v11v91-sector}
If
\[
                              \sum_i(2\beta_i+\eta_i)<\infty, \tag{91.10}
\]
then the charge laws \(p_i=(Q_i)_\#\mu_i\) converge in total variation
to a probability \(p_\infty\) on \(\mathbb Z\).  If, uniformly in
\(i\),
\[
                              p_i\{|q|\ge n\}\le Ce^{-\delta n},
                                                                    \tag{91.11}
\]
then
\[
                              p_\infty\{|q|\ge n\}\le Ce^{-\delta n}.
                                                                    \tag{91.12}
\]
\end{theorem}

\begin{proof}
The coupling argument of V77 and the amplitude defect give
\[
                              d_{\rm b}(p_{i+1},p_i)
                              \le2\beta_i+\eta_i.             \tag{91.13}
\]
Thus the laws are Cauchy in total variation.  Their limit has mass
one.  Total-variation convergence permits passage to the probability
of the fixed tail event, proving (91.12).
\end{proof}

The exponential hypothesis (91.11) follows from the V76
energy--entropy gap, possibly transferred through V78.  The index and
gapped refinement homotopy of V77 provide (91.8).

\subsection{Complex phase and cylinder functional}

Let \(u_i:X_i\to\mathbb C\), \(|u_i|\le1\), be the complete regulated
phase after determinant-line transport, local counterterms, theta
factors, and reference-block phases.  Define
\[
 \rho_i=\int|u_{i+1}-u_i\circ P_i|\,d\mu_{i+1},\qquad
 D_i=\int u_i\,d\mu_i.                                      \tag{91.14}
\]

\begin{theorem}[Normalised complex cylinder continuum]
\label{thm:v11v91-complex}
Fix \(i_0\) and suppose
\[
 R_{i_0}:=\sum_{i\ge i_0}(\rho_i+\eta_i)<|D_{i_0}|.          \tag{91.15}
\]
For a bounded cylinder function
\(F=f\circ\pi_k\) on \(X_\infty\), where \(f\) is bounded on \(X_k\),
set for \(i\ge k\)
\[
 N_i(F)=\int_{X_i}f(P_{k,i}x)u_i(x)\,d\mu_i(x).              \tag{91.16}
\]
Then the limits
\[
 D_\infty=\lim_iD_i,\qquad
 N_\infty(F)=\lim_iN_i(F),\qquad
 \Omega_\infty(F)={N_\infty(F)\over D_\infty}               \tag{91.17}
\]
exist and satisfy
\[
 |D_\infty|\ge|D_{i_0}|-R_{i_0}>0,\qquad
 |\Omega_\infty(F)|
 \le{\|F\|_\infty\over|D_{i_0}|-R_{i_0}}.                   \tag{91.18}
\]
The value of \(\Omega_\infty(F)\) is independent of the cutoff level
used to represent the same cylinder function.
\end{theorem}

\begin{proof}
The adjacent denominator estimate is (79.5):
\[
                              |D_{i+1}-D_i|\le\rho_i+\eta_i.
\]
For (91.16), compatibility of the cylinder pullback and the same
insertion argument give
\[
 |N_{i+1}(F)-N_i(F)|
       \le\|F\|_\infty(\rho_i+\eta_i).                       \tag{91.19}
\]
Thus both sequences are Cauchy.  The reverse triangle inequality and
(91.15) give the denominator bound in (91.18); the numerator bound
follows from \(|u_i|\le1\).  Two representations of the same cylinder
function agree after pullback to any common later level, so their
limits agree.
\end{proof}

A concrete sufficient estimate for \(\rho_i\) is
\[
 \rho_i\le
 { {\mathbb E}[{\bf1}_{{\cal H}_{i+1}}\|\Delta_i\|_1]
       \over(1-\vartheta)\sigma_i}
 +2\chi_i+\zeta_i+{\cal H}_i+d_i^{\rm ref}
 +2\beta_i,                                                  \tag{91.20}
\]
where the terms are, respectively, the zero-free relative operator
increment of V81--V82, the bad-gap probability, the counterterm
mismatch of V83, the determinant-line cycle defect and refinement
scalar of V84, and the charge-admissibility defect.  Formula (91.20)
is schematic only in the choice of common trivialisation; every term
has a precise definition in the cited versions.

\subsection{Conditional passage of physical identities}

Theorem~\ref{thm:v11v91-complex} gives a bounded functional on the
cylinder algebra.  The remaining physical conclusions have separate
stability statements.

\begin{theorem}[Passage gates]
\label{thm:v11v91-gates}
Assume Theorem~\ref{thm:v11v91-complex}.
\begin{enumerate}
\item If a reflection \(\Theta\) preserves the cylinder algebra and
\[
                       \Omega_i((\Theta F)F)\ge0             \tag{91.21}
\]
for every positive-time cylinder \(F\) at every sufficiently fine
cutoff, then
\[
                       \Omega_\infty((\Theta F)F)\ge0.       \tag{91.22}
\]
\item If cutoff Ward operators satisfy, for a limiting cylinder
      insertion \({\cal W}F\),
\[
 |\Omega_i({\cal W}_iF)|\le\varepsilon_i,\qquad
 \varepsilon_i\to0,\qquad
 \Omega_i({\cal W}_iF)-\Omega_i({\cal W}F)\to0,              \tag{91.23}
\]
then
\[
                              \Omega_\infty({\cal W}F)=0.    \tag{91.24}
\]
\item If translated cylinder functions obey a cutoff-uniform
      clustering bound
\[
 |\Omega_i(F\,\tau_aG)-\Omega_i(F)\Omega_i(\tau_aG)|
             \le C_{F,G}e^{-m|a|}+\epsilon_i,\qquad
             \epsilon_i\to0,                               \tag{91.25}
\]
then the same estimate without \(\epsilon_i\) holds for
\(\Omega_\infty\).
\end{enumerate}
\end{theorem}

\begin{proof}
Each claim follows by taking the already established cylinder limit.
For the product term in (91.25), convergence of the two factors gives
convergence of their product.  The bound is uniform in the cutoff, so
the limit preserves its right-hand side.
\end{proof}

Reflection positivity in (91.21) is a property of the complete
normalised functional, not of the positive amplitude \(\mu_i\) alone.
Ward convergence requires the anomaly, regulator, and boundary
defects separated in V68.  A positive \(m\) in (91.25) cannot be
borrowed from a gapped high block when physical massless fields remain
in the low block.

\subsection{Acquisition ledger for the actual theory}

\begin{ledger}[Conditions that have theorems but not physical inputs]
\label{led:v11v91-open}
The conditional assembly requires the following still-unacquired
regulator estimates:
\begin{enumerate}
\item a renormalised gravitational interaction satisfying the
      relative lower bound, lower normalisation, and conditional
      density series of V87 despite the conformal obstruction of V86;
\item a high-probability diffeomorphism slice atlas with uniform
      elliptic gap, finite multiplicity, full Jacobian overlap
      agreement, and equivariant refinement as in V88--V89;
\item a global lattice topological charge with gapped refinement and
      a uniform energy advantage over sector entropy as in V76--V77;
\item a zero-free, stabilised chiral determinant carrier with the
      trace-ideal shell series, exact local counterterm matching, and
      trivial global holonomy of V81--V84;
\item a finite phase anchor whose margin exceeds the complete tail
      (91.15);
\item exact or vanishing-defect gauge and diffeomorphism Ward
      identities for the complete measured density;
\item reflection positivity and a uniform clustering mechanism
      compatible with the massless photon and graviton sectors;
\item universality and parameter matching identifying the limit with
      the Standard Model coupled to Einstein gravity.
\end{enumerate}
\end{ledger}

\begin{theorem}[Precise claim status]
\label{thm:v11v91-status}
If every item in Ledger~\ref{led:v11v91-open} is supplied with the
summability and anchor inequalities stated above, then
Theorems~\ref{thm:v11v91-positive}--\ref{thm:v11v91-gates} construct
the corresponding projective positive law, nonzero normalised complex
cylinder functional, sector law, and inherited physical identities.
The attached sources do not currently supply every item, so the full
SM--Einstein continuum has not been proved.
\end{theorem}

\begin{proof}
The constructive implication is the concatenation of the cited
theorems, with (91.20) inserted into (91.15).  The second sentence is
the collection of the explicit firewalls in V76--V90 and is not a
mathematical inference beyond those recorded inputs.
\end{proof}

\begin{decision}[V92 acquisition]
\label{dec:v11v91-next}
The first unresolved input in dependency order is item one: a
renormalised positive gravitational interaction relative to the
coercive law.  V92 should test a higher-derivative conformal
stabilisation and determine whether its regulator-removal limit can
satisfy a uniform V87 relative bound; if it cannot, the programme
must select and analyse a complex contour instead.
\end{decision}

\begin{assessment}[Progress after V91]
\label{ass:v11v91-status}
The positive quotient measure, sector law, determinant phase, and
normalised complex cylinder limits now sit in one exact implication
chain.  The theorem identifies the leading physical obstruction as
the renormalised gravitational amplitude, followed by the global
diffeomorphism quotient and chiral phase estimates; it does not
mislabel these conditional inputs as completed constructions.
\end{assessment}

\subsection{Parent record}

This V91 note was formed by appending the preceding material to the
validated V90 source, whose record was
\[
\begin{array}{c|c}
\text{lines}&26974\\
\text{bytes}&871902\\
\text{SHA--256}&
\texttt{922DEE3A514042C57E81A295C75230D2950A53546B4B6E142E2DD61DDB0629BF}.
\end{array}                                                     \tag{91.26}
\]

% =============================================================================
% END APPENDED ELEVENTH-SERIES V91 MATERIAL
% =============================================================================
% =============================================================================
% BEGIN APPENDED ELEVENTH-SERIES V92 MATERIAL
% =============================================================================

\section{Testing higher-derivative removal in the conformal sector}
\label{sec:v11v92}

V91 identifies a uniform positive gravitational amplitude as the
first missing construction.  A natural proposal is to stabilise the
negative two-derivative conformal term by a positive four-derivative
term and then remove the latter.  This version tests that proposal at
the Gaussian Hessian level.

Let \(\lambda_n>0\) be the nonconstant eigenvalues of a positive
Laplacian on the fixed-volume conformal subspace, with
\(\lambda_n\to\infty\).  Consider the quadratic action
\[
 S_{\varepsilon,m,N}(\phi)
 ={1\over2}\sum_{n\le N}
       a_{\varepsilon,m}(\lambda_n)\phi_n^2,\qquad
 a_{\varepsilon,m}(\lambda)
       =\varepsilon\lambda^2-\kappa\lambda+m^2,              \tag{92.1}
\]
where \(\varepsilon>0\), \(\kappa>0\), and \(m^2\ge0\).
The negative term models the conformal Einstein Hessian and the
positive leading term models a curvature-squared regulator.

\begin{theorem}[Coercivity threshold]
\label{thm:v11v92-threshold}
For every \(\lambda\ge0\),
\[
 a_{\varepsilon,m}(\lambda)
 =\varepsilon\left(\lambda-{\kappa\over2\varepsilon}\right)^2
       +m^2-{\kappa^2\over4\varepsilon}.                    \tag{92.2}
\]
Consequently the uniform continuum spectral bound
\[
                         a_{\varepsilon,m}(\lambda)\ge a_0>0
                    \quad(\lambda\ge0)                      \tag{92.3}
\]
holds if and only if
\[
                         m^2\ge{\kappa^2\over4\varepsilon}+a_0.
                                                                    \tag{92.4}
\]
With \(m\) fixed, (92.3) fails for all sufficiently small
\(\varepsilon\).
\end{theorem}

\begin{proof}
Equation (92.2) is completion of the square.  Its minimum over
\([0,\infty)\) occurs at
\(\lambda=\kappa/(2\varepsilon)\) and equals
\(m^2-\kappa^2/(4\varepsilon)\), proving the equivalence.
\end{proof}

For a discrete spectrum, (92.4) is sufficient.  Whenever the spectrum
has eigenvalues uniformly close in relative scale to
\(\kappa/(2\varepsilon)\), the same order \(m^2\asymp\varepsilon^{-1}\)
is also necessary.

\begin{corollary}[Growing unstable band]
\label{cor:v11v92-band}
If \(m\) is fixed and \(\varepsilon\downarrow0\), then
\[
 a_{\varepsilon,m}(\lambda)<0
 \quad\text{for }\lambda\in
 \left(
 {\kappa-\sqrt{\kappa^2-4\varepsilon m^2}\over2\varepsilon},
 {\kappa+\sqrt{\kappa^2-4\varepsilon m^2}\over2\varepsilon}
 \right),                                                    \tag{92.5}
\]
once \(\kappa^2>4\varepsilon m^2\).  The upper edge is asymptotic to
\(\kappa/\varepsilon\).  Thus the negative modes cannot be confined
to one cutoff-independent finite low block.
\end{corollary}

\begin{proof}
Solve the quadratic inequality
\(\varepsilon\lambda^2-\kappa\lambda+m^2<0\).  The root asymptotics
follow by expansion at \(\varepsilon=0\).
\end{proof}

This is stronger than saying that a fourth-order term controls the
highest frequency at each fixed \(\varepsilon\).  Its removal exposes
an increasing band of negative Hessian directions.

\subsection{The \(L^r\) reweighting test}

Use as a positive reference
\[
 d\nu_{\varepsilon,m,N}
 \propto\exp\left[
 -{1\over2}\sum_{n\le N}
       (\varepsilon\lambda_n^2+m^2)\phi_n^2\right]d\phi,     \tag{92.6}
\]
and place the Einstein quadratic term in the interaction:
\[
                         g_N(\phi)
 =\exp\left({\kappa\over2}
             \sum_{n\le N}\lambda_n\phi_n^2\right).          \tag{92.7}
\]

\begin{theorem}[Exact Gaussian reweighting criterion]
\label{thm:v11v92-Lr}
For \(r>1\),
\[
 \|g_N\|_{L^r(\nu_{\varepsilon,m,N})}<\infty
\]
if and only if
\[
                         \varepsilon\lambda_n^2+m^2
                         >r\kappa\lambda_n
                    \quad\text{for every }n\le N.           \tag{92.8}
\]
A cutoff-uniform sufficient continuum condition is
\[
                         m^2>{r^2\kappa^2\over4\varepsilon}. \tag{92.9}
\]
No fixed \(m\) satisfies (92.9) as
\(\varepsilon\downarrow0\).
\end{theorem}

\begin{proof}
The \(r\)-th power of (92.7) changes the coefficient of \(\phi_n^2\)
in the Gaussian exponent to
\[
                         \varepsilon\lambda_n^2+m^2
                               -r\kappa\lambda_n.
\]
The product integral is finite exactly when every coefficient is
positive, proving (92.8).  Completing the square with
\(\kappa\) replaced by \(r\kappa\) gives (92.9).
\end{proof}

Thus the simple higher-derivative reference does not provide the
uniform V87 \(L^r\) hypothesis while its coefficient is removed.
This is an exact Gaussian obstruction, not merely a failure of a
particular estimate.

\subsection{Why the compensating mass does not recover Einstein
fluctuations}

Choose the minimal uniform stabilisation
\[
                         m_\varepsilon^2
 ={r^2\kappa^2\over4\varepsilon}+m_0^2.                     \tag{92.10}
\]
For each fixed eigenmode,
\[
 \varepsilon\lambda_n^2+m_\varepsilon^2
                         \longrightarrow\infty
                    \quad(\varepsilon\downarrow0).          \tag{92.11}
\]
Hence its reference variance tends to zero:
\[
 \mathbb E_{\nu_{\varepsilon,m_\varepsilon}}
                              |\phi_n|^2
 ={1\over\varepsilon\lambda_n^2+m_\varepsilon^2}
                         \longrightarrow0.                  \tag{92.12}
\]

\begin{proposition}[Degenerate removal limit]
\label{prop:v11v92-degenerate}
Under (92.10), every fixed finite collection of conformal modes
converges in law under the reference measure to zero.  Therefore this
uniformly coercive removal scheme does not approach a nondegenerate
Einstein conformal Gaussian.
\end{proposition}

\begin{proof}
The modes are centred independent Gaussians and their variances tend
to zero by (92.12).  Their joint covariance matrix therefore tends to
zero, which implies convergence in probability and in law to the
origin.
\end{proof}

A field rescaling that keeps (92.12) nonzero also rescales its
couplings and must be analysed as a different continuum limit.

\subsection{Finite cutoff versus uniform continuum}

At any fixed \(N\), one may choose \(\varepsilon\), \(m\), or a compact
field domain so that the integral is finite.  The master theorem V91
requires estimates uniform along the refinement sequence.  If
\(\lambda_{N_i}\to\infty\) and \(\varepsilon_i\to0\), then the spectral
window near \(r\kappa/(2\varepsilon_i)\) eventually enters the cutoff
unless the removal is faster than the spectral range in a way that
leaves a large lower-frequency negative band.  Neither ordering
restores (92.8) for all retained modes with fixed \(m\).

Nonlinear positive curvature-squared terms can change large-field
behaviour, but their Hessian at a background still contains the
competition (92.1).  A uniform global lower bound would have to
supply an additional mechanism that overcomes this growing negative
band.

\begin{theorem}[No naive positive-regulator removal]
\label{thm:v11v92-no-go}
Within the quadratic conformal model (92.1), the following three
requirements cannot hold simultaneously:
\begin{enumerate}
\item \(\varepsilon_i\to0\);
\item a cutoff-uniform \(L^r\) bound for the Einstein reweighting
      (92.7), for some fixed \(r>1\);
\item a nondegenerate fixed-mode Gaussian limit with bounded
      \(m_i^2\).
\end{enumerate}
\end{theorem}

\begin{proof}
Requirements one and two force
\(m_i^2\ge r^2\kappa^2/(4\varepsilon_i)\) up to a positive margin by
Theorem~\ref{thm:v11v92-Lr}, contradicting boundedness in requirement
three.  If the mass is allowed to grow at that rate,
Proposition~\ref{prop:v11v92-degenerate} gives degeneracy.
\end{proof}

\subsection{Remaining routes}

The theorem rules out only the naive positive fourth-order removal.
It leaves three logically distinct possibilities:
\begin{enumerate}
\item retain nonzero higher-derivative couplings and prove that the
      resulting fixed point has the required low-energy Einstein
      limit;
\item find a nonlinear positive reference and interaction split whose
      relative bound is not the Gaussian comparison (92.7);
\item abandon positivity in the conformal coordinate and construct a
      refinement-compatible complex contour with a direct
      noncancellation and reconstruction theorem.
\end{enumerate}

The first option is a universality problem, the second a new
coercivity problem, and the third a complex-measure problem.  They
cannot be merged into one unspecified analytic continuation.

\begin{warning}[Matter does not automatically stabilise the Hessian]
\label{warn:v11v92-matter}
Integrating out matter can generate curvature-squared terms and a
cosmological potential, but the signs and running coefficients must be
computed for the complete regulated multiplet.  A finite number of
low-momentum corrections does not remove the growing band
(92.5), and a diagram-level sign is not a complete-density
coercivity proof.
\end{warning}

\begin{firewall}[Still no positive Einstein amplitude]
\label{fire:v11v92-boundary}
V92 eliminates one simple regulator-removal strategy.  It does not
construct an alternative positive gravitational interaction, prove
an ultraviolet fixed point, or define the complex conformal contour.
Item one of Ledger~\ref{led:v11v91-open} therefore remains open.
\end{firewall}

\begin{decision}[V93 acquisition]
\label{dec:v11v92-next}
V93 should formulate a finite-dimensional complex conformal contour
criterion: uniform thimble coercivity, avoidance of determinant and
metric-degeneracy singularities, stable relative homology under
refinement, and a denominator anchor.  That route must be developed
without importing the positive-measure conclusions of V87 where they
do not apply.
\end{decision}

\begin{assessment}[Progress after V92]
\label{ass:v11v92-status}
The positive higher-derivative proposal has passed a decisive
spectral test and fails in its naive removal form.  Uniform
\(L^r\) reweighting forces a compensating mass of order
\(\varepsilon^{-1}\), which collapses every fixed conformal mode.
The next branch must retain a different fixed-point dynamics, find a
new coercive split, or construct a complex contour directly.
\end{assessment}

\subsection{Parent record}

This V92 note was formed by appending the preceding material to the
validated V91 source, whose record was
\[
\begin{array}{c|c}
\text{lines}&27310\\
\text{bytes}&884491\\
\text{SHA--256}&
\texttt{311646693CB249E1DE6764527C1F634D0AE64D4DA4C70A690D25C3EF92D6EA0D}.
\end{array}                                                     \tag{92.13}
\]

% =============================================================================
% END APPENDED ELEVENTH-SERIES V92 MATERIAL
% =============================================================================
% =============================================================================
% BEGIN APPENDED ELEVENTH-SERIES V93 MATERIAL
% =============================================================================

\section{Finite-dimensional complex conformal contours}
\label{sec:v11v93}

V92 rules out the naive removal of a positive fourth-order conformal
regulator within a uniform positive-reweighting scheme.  The remaining
direct branch is a complex integration cycle.  This version states
the finite-dimensional analytic conditions under which such a cycle
defines an absolutely convergent integral and makes its phase budget
explicit.

Let \(D\subset\mathbb C^N\) be open, let
\(\Sigma\subset D\) be the singular divisor of the complete regulated
integrand, and let \(S:D\setminus\Sigma\to\mathbb C\) be holomorphic.
Let
\[
                       \gamma:\mathbb R^N\longrightarrow
                       D\setminus\Sigma                      \tag{93.1}
\]
be a proper \(C^1\) oriented immersion.  Its image
\(\Gamma=\gamma(\mathbb R^N)\) is the candidate contour.  Write
\[
 J_\gamma(t)=
 \det_{\mathbb C}\left({\partial\gamma^a\over\partial t^b}\right)
                                                                    \tag{93.2}
\]
for the pullback coefficient of
\(dz^1\wedge\cdots\wedge dz^N\).

\begin{theorem}[Coercive-contour integrability]
\label{thm:v11v93-integrability}
Suppose constants \(a,c>0\), \(b,C\ge0\), \(p>0\), and \(m\ge0\)
satisf
\[
 \operatorname{Re}S(\gamma(t))
                       \ge a\|\gamma(t)\|^p-b,               \tag{93.3}
\]
\[
 \|\gamma(t)\|\ge c\|t\|-C,\qquad
 |J_\gamma(t)|\le C(1+\|t\|)^m.                             \tag{93.4}
\]
If a holomorphic observable \(F\) obeys on \(\Gamma\)
\[
                         |F(z)|\le C_F(1+\|z\|)^q            \tag{93.5}
\]
for some \(q\ge0\), then
\[
                  \int_\Gamma e^{-S(z)}F(z)\,
                     dz^1\cdots dz^N                        \tag{93.6}
\]
is absolutely convergent.
\end{theorem}

\begin{proof}
Pull (93.6) back by \(\gamma\).  Its absolute integrand is at most
\[
 C C_F e^b(1+\|\gamma(t)\|)^{m+q}
                  e^{-a\|\gamma(t)\|^p}.
\]
Proper linear escape in (93.4), with a harmless adjustment on a
compact ball, bounds this by a polynomial in \(\|t\|\) times
\(e^{-a'( \|t\|)^p}\) for some \(a'>0\).  That function is integrable
on \(\mathbb R^N\).
\end{proof}

Condition (93.3) concerns the complete real part on the actual
contour.  Positivity of the quadratic Hessian at one critical point
does not imply global coercivity, and coercivity before including
Jacobian or determinant singularities does not imply (93.6).

\subsection{Amplitude and residual phase on a contour}

Define
\[
 Z_\Gamma^{\rm abs}
 =\int_{\mathbb R^N}
       e^{-\operatorname{Re}S(\gamma(t))}
       |J_\gamma(t)|\,dt,                                   \tag{93.7}
\]
and, when this is positive and finite,
\[
 d\mu_\Gamma(t)
 ={e^{-\operatorname{Re}S(\gamma(t))}
       |J_\gamma(t)|\over Z_\Gamma^{\rm abs}}\,dt.           \tag{93.8}
\]
Away from the zero set of \(J_\gamma\), set
\[
 u_\Gamma(t)
 =e^{-i\operatorname{Im}S(\gamma(t))}
       {J_\gamma(t)\over|J_\gamma(t)|}.                      \tag{93.9}
\]
Assign any value of modulus at most one on the null Jacobian set.

\begin{proposition}[Exact contour phase decomposition]
\label{prop:v11v93-decomposition}
For \(F=1\),
\[
 Z_\Gamma=Z_\Gamma^{\rm abs}D_\Gamma,\qquad
 D_\Gamma=\mathbb E_{\mu_\Gamma}u_\Gamma.                   \tag{93.10}
\]
For bounded \(F\),
\[
 {\int_\Gamma e^{-S}F\,dz^1\cdots dz^N\over Z_\Gamma}
 ={ \mathbb E_{\mu_\Gamma}
       [\,u_\Gamma F\circ\gamma\,]\over D_\Gamma},           \tag{93.11}
\]
provided \(D_\Gamma\ne0\).
\end{proposition}

\begin{proof}
Equations (93.7)--(93.9) factor the pulled-back complex density into
its absolute value and unit phase.  Integrating proves (93.10);
the same factorisation with \(F\) proves (93.11).
\end{proof}

Thus contour coercivity establishes a positive amplitude law but does
not establish the nonzero denominator.  Lemma~\ref{lem:v11v79-arc}
applies to (93.9): if the combined action and Jacobian phase lies in a
strict half-plane on a high-probability set, it supplies an anchor.

\begin{corollary}[Single-thimble phase criterion]
\label{cor:v11v93-thimble}
Suppose \(\operatorname{Im}S\) is constant on \(\Gamma\) and the
Jacobian phase satisfies, on an event of
\(\mu_\Gamma\)-probability at least \(1-\varepsilon\),
\[
 \operatorname{Re}\left(
 e^{-i\varphi}{J_\gamma\over|J_\gamma|}\right)\ge\cos\alpha,
                  \qquad 0\le\alpha<{\pi\over2}.            \tag{93.12}
\]
Then
\[
 |D_\Gamma|\ge(1-\varepsilon)\cos\alpha-\varepsilon.         \tag{93.13}
\]
\end{corollary}

\begin{proof}
Absorb the constant action phase into \(\varphi\) and apply
Lemma~\ref{lem:v11v79-arc}.
\end{proof}

A constant imaginary action is not enough if the residual Jacobian
phase fluctuates broadly.  Multiple thimbles can also cancel even
when each individual denominator is nonzero.

\subsection{Holomorphic deformation invariance}

Let \(\Gamma_0\) and \(\Gamma_1\) be two contours in
\(D\setminus\Sigma\), joined by an oriented \((N+1)\)-chain \(B_R\)
after truncation at radius \(R\).  Let
\[
                 \omega_F=e^{-S(z)}F(z)\,
                    dz^1\wedge\cdots\wedge dz^N.             \tag{93.14}
\]
Holomorphy gives \(d\omega_F=0\).

\begin{theorem}[Relative-homology deformation]
\label{thm:v11v93-deformation}
Suppose:
\begin{enumerate}
\item the deformation chain does not meet \(\Sigma\);
\item \(\Gamma_1-\Gamma_0\) together with an outer boundary
      \(A_R\) is \(\partial B_R\);
\item the end contribution obeys
\[
                              \int_{A_R}|\omega_F|\longrightarrow0.
                                                                    \tag{93.15}
\]
\end{enumerate}
Then
\[
                              \int_{\Gamma_1}\omega_F
                              =\int_{\Gamma_0}\omega_F.       \tag{93.16}
\]
\end{theorem}

\begin{proof}
Stokes' theorem gives
\[
 0=\int_{B_R}d\omega_F
  =\int_{\Gamma_1\cap B_R}\omega_F
   -\int_{\Gamma_0\cap B_R}\omega_F+\int_{A_R}\omega_F.
\]
Absolute convergence and (93.15) permit \(R\to\infty\), proving
(93.16).
\end{proof}

The theorem requires relative homology in the complement of the full
singular divisor.  A deformation crossing a pole, a branch divisor,
a metric-degeneracy locus, or a gauge-fixing singularity can change
the integral by residues or by a jump of homology class.

\subsection{Steepest-ascent ends}

For the standard Hermitian metric on \(\mathbb C^N\), consider
\[
                              {dz^a\over d\tau}
                              =\overline{\partial_aS(z)}.     \tag{93.17}
\]

\begin{lemma}[Action along a thimble end]
\label{lem:v11v93-flow}
Along every solution of (93.17),
\[
 {d\over d\tau}\operatorname{Re}S(z(\tau))
             =\sum_a|\partial_aS|^2\ge0,\qquad
 {d\over d\tau}\operatorname{Im}S(z(\tau))=0.               \tag{93.18}
\]
\end{lemma}

\begin{proof}
Holomorphy gives
\[
 {dS\over d\tau}
 =\sum_a\partial_aS\,\overline{\partial_aS}
 =\sum_a|\partial_aS|^2,
\]
which is real and nonnegative.
\end{proof}

Equation (93.18) explains the constant action phase on an ideal
Lefschetz thimble.  It does not prove that the flow is complete,
proper, avoids \(\Sigma\), or has the uniform escape rate (93.3).

\subsection{The rotated Gaussian diagnostic}

The one-mode negative conformal Gaussian has
\[
                              S(z)=-{\kappa\over2}z^2,
                       \qquad \kappa>0.                     \tag{93.19}
\]
The real integral diverges.  On the contour \(z=iy\),
\[
 \int_{i\mathbb R}e^{-S(z)}\,dz
 =i\int_{\mathbb R}e^{-\kappa y^2/2}\,dy
 =i\sqrt{2\pi\over\kappa}.                                  \tag{93.20}
\]
For \(N\) independently rotated negative modes, the orientation phase
is \(i^N\).  Hence a cutoff that adds negative modes changes the
contour phase even though every amplitude integral is Gaussian.

\begin{warning}[Dimension phase must be stabilised]
\label{warn:v11v93-dimension}
The factor \(i^N\) in (93.20) is an extensive cutoff phase.  It must
be included in the reference-block phase and counterterm transport of
V81--V83.  Ignoring it makes the adjacent denominator oscillate with
the number of retained conformal modes.
\end{warning}

\subsection{Finite-cutoff contour ledger}

A candidate SM--Einstein contour \(\Gamma_i\) must provide:
\begin{enumerate}
\item a proper parameterisation and a uniform coercive bound of the
      form (93.3)--(93.4);
\item a full singular divisor containing gauge, ghost, fermion,
      coordinate, and metric degeneracies;
\item a proof that \(\Gamma_i\) and its required deformations avoid
      that divisor;
\item a controlled Jacobian phase and a nonzero anchor for
      \(D_{\Gamma_i}\);
\item a declared relative-homology class and vanishing ends as in
      (93.15);
\item a stabilised orientation phase when the cutoff dimension
      changes;
\item compatibility with the antilinear reflection used for
      Osterwalder--Schrader positivity.
\end{enumerate}

\begin{firewall}[No continuum contour acquired]
\label{fire:v11v93-boundary}
V93 proves finite-dimensional sufficiency statements.  It does not
construct the coupled SM--Einstein contours, prove a uniform
coercivity constant as \(N_i\to\infty\), control Stokes jumps, bound
the residual Jacobian phase, or establish reflection positivity.
The complex-contour branch therefore remains open.
\end{firewall}

\begin{decision}[V94 acquisition]
\label{dec:v11v93-next}
V94 will formulate adjacent refinement of relative contour classes.
It will separate deformation error, singular-set clearance,
orientation phase, and tail-at-infinity error, then derive the
summable cylinder and denominator estimate needed by V91 without
assuming a positive Radon--Nikodym comparison.
\end{decision}

\begin{assessment}[Progress after V93]
\label{ass:v11v93-status}
The complex branch now has a precise finite-cutoff foundation:
coercivity gives absolute convergence, phase quenching exposes the
denominator, relative homology gives deformation invariance, and the
rotated Gaussian identifies the extensive orientation phase that must
be transported across cutoffs.
\end{assessment}

\subsection{Parent record}

This V93 note was formed by appending the preceding material to the
validated V92 source, whose record was
\[
\begin{array}{c|c}
\text{lines}&27601\\
\text{bytes}&895395\\
\text{SHA--256}&
\texttt{E8B61B5C6083C1A28D6370599D1E6F1E9989E86A4ED333F3A5D3D0C73A532CFE}.
\end{array}                                                     \tag{93.21}
\]

% =============================================================================
% END APPENDED ELEVENTH-SERIES V93 MATERIAL
% =============================================================================
% =============================================================================
% BEGIN APPENDED ELEVENTH-SERIES V94 MATERIAL
% =============================================================================

\section{Relative contour classes under cutoff refinement}
\label{sec:v11v94}

A convergent contour at each cutoff does not yet define a continuum.
The fine cycle must be compared with the coarse cycle after adding the
new modes.  This version gives the exact stabilisation and the
adjacent error estimate without using a positive
Radon--Nikodym derivative.

Let \(\Gamma_i\subset D_i\setminus\Sigma_i\) have real dimension
\(N_i\), and let
\[
                   \omega_i=e^{-S_i(z)}
                      dz^1\wedge\cdots\wedge dz^{N_i}        \tag{94.1}
\]
denote a renormalised holomorphic top form, including all determinant
and coordinate factors.  For the \(m_i=N_{i+1}-N_i\) new variables,
choose a reference cycle \(T_i\) and a holomorphic top form \(\rho_i\)
such that
\[
                              c_i:=\int_{T_i}\rho_i\ne0.      \tag{94.2}
\]
Define the normalised stabilising form and product cycle
\[
 \widehat\rho_i=c_i^{-1}\rho_i,\qquad
 \widetilde\Gamma_i=\Gamma_i\times T_i,\qquad
 \widetilde\omega_i=\omega_i\wedge\widehat\rho_i.            \tag{94.3}
\]

\begin{lemma}[Exact new-mode stabilisation]
\label{lem:v11v94-stabilise}
For every bounded integrable coarse observable \(F\),
\[
 \int_{\widetilde\Gamma_i}
       (F\circ\pi_i)\widetilde\omega_i
                         =\int_{\Gamma_i}F\omega_i.          \tag{94.4}
\]
\end{lemma}

\begin{proof}
Absolute integrability permits Fubini's theorem.  The new-mode factor
is \(c_i^{-1}\int_{T_i}\rho_i=1\).
\end{proof}

The full complex number \(c_i\), including its orientation phase, is
used in (94.3).  For rotated negative Gaussian modes it contains the
factor \(i^{m_i}\) found in V93.  Replacing \(c_i\) by \(|c_i|\)
would introduce a spurious adjacent phase.

\subsection{Exact relative-homology consistency}

Embed \(\widetilde\Gamma_i\) and \(\Gamma_{i+1}\) in a common complex
domain \(D_{i+1}\setminus\Sigma_{i+1}\).  Let
\(\Pi_i:D_{i+1}\to D_i\) be the holomorphic coarse-coordinate map.

\begin{theorem}[Exact contour refinement]
\label{thm:v11v94-exact}
Suppose:
\begin{enumerate}
\item \(\Gamma_{i+1}\) and \(\widetilde\Gamma_i\) are in the same
      relative homology class in
      \(D_{i+1}\setminus\Sigma_{i+1}\), with vanishing boundary at
      infinity;
\item
\[
                              \omega_{i+1}=\widetilde\omega_i
                                                                    \tag{94.5}
\]
      as holomorphic top forms in a neighbourhood of a connecting
      chain;
\item \(F\circ\Pi_i\) is holomorphic there.
\end{enumerate}
Then
\[
 \int_{\Gamma_{i+1}}(F\circ\Pi_i)\omega_{i+1}
                  =\int_{\Gamma_i}F\omega_i.                 \tag{94.6}
\]
\end{theorem}

\begin{proof}
Apply Theorem~\ref{thm:v11v93-deformation} to the closed top form
\((F\circ\Pi_i)\widetilde\omega_i\), then apply
Lemma~\ref{lem:v11v94-stabilise}.
\end{proof}

Condition (94.5) is exact Wilsonian consistency.  In a renormalised
construction it is normally replaced by a summable form mismatch.

\subsection{An adjacent error theorem}

Let \(B_{i,R}\) be a truncated chain from
\(\widetilde\Gamma_i\) to \(\Gamma_{i+1}\).  Suppose singular
neighbourhoods have been excised, leaving an additional oriented
boundary \(H_{i,R}\), and let \(A_{i,R}\) be the outer boundary.
Choose a comparison holomorphic form
\(\omega_i^{\rm cmp}\) near the chain.  For every bounded holomorphic
coarse cylinder \(F\), assume
\[
 \int_{\Gamma_{i+1}}
 |F\circ\Pi_i|\,
 |\omega_{i+1}-\omega_i^{\rm cmp}|
                  \le\|F\|_\infty a_i,                      \tag{94.7}
\]
\[
 \limsup_{R\to\infty}
 \int_{A_{i,R}}|F\circ\Pi_i|\,|\omega_i^{\rm cmp}|
                  \le\|F\|_\infty b_i,                      \tag{94.8}
\]
\[
 \limsup_{R\to\infty}
 \int_{H_{i,R}}|F\circ\Pi_i|\,|\omega_i^{\rm cmp}|
                  \le\|F\|_\infty s_i,                      \tag{94.9}
\]
and
\[
 \left|
 \int_{\widetilde\Gamma_i}(F\circ\pi_i)
       (\omega_i^{\rm cmp}-\widetilde\omega_i)
 \right|
                  \le\|F\|_\infty r_i.                      \tag{94.10}
\]

\begin{theorem}[Adjacent complex contour bound]
\label{thm:v11v94-adjacent}
Under (94.7)--(94.10),
\[
 \left|
 \int_{\Gamma_{i+1}}(F\circ\Pi_i)\omega_{i+1}
 -\int_{\Gamma_i}F\omega_i
 \right|
 \le\|F\|_\infty(a_i+b_i+s_i+r_i).                          \tag{94.11}
\]
\end{theorem}

\begin{proof}
First replace \(\omega_{i+1}\) on the fine contour by
\(\omega_i^{\rm cmp}\), at cost (94.7).  Stokes' theorem on
\(B_{i,R}\), applied to the closed form
\((F\circ\Pi_i)\omega_i^{\rm cmp}\), moves the integral to
\(\widetilde\Gamma_i\).  The outer and excision boundaries cost
(94.8)--(94.9).  Replace the comparison form by
\(\widetilde\omega_i\) at cost (94.10), and use (94.4).
The triangle inequality gives (94.11).
\end{proof}

The four terms have distinct meanings:
\[
 \begin{array}{c|l}
 a_i&\text{fine action, determinant, and counterterm mismatch},\\
 b_i&\text{failure of uniform decay at infinity},\\
 s_i&\text{residual boundary around an excised singular set},\\
 r_i&\text{reference-block and coarse-form mismatch}.
 \end{array}                                                  \tag{94.12}
\]
They should not be compressed into one unanalysed numerical residual.

\subsection{Singular clearance is a topological gate}

If the connecting chain remains in
\(D_{i+1}\setminus\Sigma_{i+1}\), no excision boundary is needed and
\(s_i=0\).  If it crosses \(\Sigma_{i+1}\), the two cycles may lie in
different relative homology classes.

\begin{proposition}[Uniform clearance criterion]
\label{prop:v11v94-clearance}
Let \(H_i:[0,1]\times\widetilde\Gamma_i\to D_{i+1}\) be a proper
contour homotopy.  If
\[
 \inf_{t,z}\operatorname{dist}
       (H_i(t,z),\Sigma_{i+1})\ge d_i>0                     \tag{94.13}
\]
and the end estimate is uniform in \(t\), then the homotopy produces
no singular boundary and no Stokes jump.
\end{proposition}

\begin{proof}
The image of the homotopy lies in the complement of the singular
divisor by (94.13).  Truncate it at infinity and apply
Theorem~\ref{thm:v11v93-deformation}; the only extra boundary is the
outer end, which vanishes uniformly.
\end{proof}

A small value of \(d_i\) is not itself an error in the integral, but
it magnifies derivative and interval bounds used to prove (94.7).
At \(d_i=0\), continuity alone gives no comparison.  If residues or
vanishing cycles arise, their contributions must be computed and
added as new saddle sectors.

\subsection{Summable complex-cylinder construction}

Let
\[
                              \varepsilon_i=a_i+b_i+s_i+r_i. \tag{94.14}
\]
Assume the forms have been divided by declared nonzero reference
normalisations so that their total variations are uniformly bounded.
Put
\[
                              Z_i=\int_{\Gamma_i}\omega_i.    \tag{94.15}
\]

\begin{theorem}[Contour-projective cylinder limit]
\label{thm:v11v94-limit}
If
\[
                       \sum_{i\ge i_0}\varepsilon_i<|Z_{i_0}|,\tag{94.16}
\]
then \(Z_i\) converges to \(Z_\infty\ne0\), with
\[
 |Z_\infty|\ge|Z_{i_0}|-\sum_{i\ge i_0}\varepsilon_i.        \tag{94.17}
\]
For every bounded holomorphic cylinder \(F\), its unnormalised
integrals converge and
\[
 \Omega_\infty(F)
 =\lim_i
 {\int_{\Gamma_i}(F\circ\Pi_{k,i})\omega_i\over Z_i}         \tag{94.18}
\]
exists.
\end{theorem}

\begin{proof}
Set \(F=1\) in (94.11) and telescope to obtain convergence and
(94.17).  For a cylinder \(F\), (94.11) gives an adjacent bound
\(\|F\|_\infty\varepsilon_i\), so its numerator is Cauchy.  Divide by
the nonzero limiting denominator.
\end{proof}

Unlike V91, this theorem does not first construct a positive
interacting probability on real metrics.  Its positivity surrogate
is the phase-quenched contour amplitude (93.8), which can vary with
the contour and is not deformation invariant.

\subsection{Observable restriction}

The holomorphy of \(F\circ\Pi_i\) is essential to the Stokes argument.
Products involving complex conjugation, absolute values, or an
antilinear reflection are not holomorphic functions of the same
variables.

\begin{warning}[No Osterwalder--Schrader conclusion from contour
homology]
\label{warn:v11v94-OS}
Equation (94.18) applies directly to holomorphic cylinder
observables.  Reflection positivity concerns
\((\Theta F)F\), where \(\Theta\) normally includes complex
conjugation.  Its contour comparison requires an additional
antilinear symmetry and cannot be inferred from relative homology of
holomorphic forms.
\end{warning}

\begin{firewall}[Missing refinement chain]
\label{fire:v11v94-boundary}
No explicit SM--Einstein cycles, normalised new-mode forms,
singular-clearance homotopies, or summable estimates
(94.7)--(94.10) are currently available.  The theorem supplies the
required architecture but does not establish the complex continuum.
\end{firewall}

\begin{decision}[Mandatory V95 audit]
\label{dec:v11v94-next}
The next version is the compulsory companion audit.  It will inspect
the newest YM and NS notes for a measured symmetry, decay, or contour
input.  V96 will then study the antilinear reflection condition needed
to keep the complex contour branch compatible with
Osterwalder--Schrader positivity.
\end{decision}

\begin{assessment}[Progress after V94]
\label{ass:v11v94-status}
Adjacent complex contours now have an exact stabilisation rule and a
summable comparison theorem.  Fine-integrand mismatch, infinity ends,
singular excision, and reference-block mismatch are separated, while
loss of singular clearance is correctly treated as a possible Stokes
jump rather than an ordinary small error.
\end{assessment}

\subsection{Parent record}

This V94 note was formed by appending the preceding material to the
validated V93 source, whose record was
\[
\begin{array}{c|c}
\text{lines}&27916\\
\text{bytes}&906247\\
\text{SHA--256}&
\texttt{1A41EC02DEEF2A827EDB00BC452EF80271AB6CB63E251196F08976EFFF87AC68}.
\end{array}                                                     \tag{94.19}
\]

% =============================================================================
% END APPENDED ELEVENTH-SERIES V94 MATERIAL
% =============================================================================
% =============================================================================
% BEGIN APPENDED ELEVENTH-SERIES V95 MATERIAL
% =============================================================================

\section{Companion audit at V95: the YM re-treeing obstruction}
\label{sec:v11v95}

The compulsory audit finds a new Yang--Mills endpoint:
\[
\begin{array}{c|r|r|l}
\text{file}&\text{lines}&\text{bytes}&\text{SHA--256}\\ \hline
\texttt{Renewal\_Yang\_Mills\_Mass\_Gap\_Research\_Notes\_V41.tex}
&16811&574386&
\texttt{191ACA59453A84990AABF333FA79E534A760B1B1A6F22AFA326B5D9E6AE1502A}\\
\texttt{Renewal\_NS\_Stability\_Research\_Notes\_V26.tex}
&5319&278283&
\texttt{82C887F8C2C69E4F28FAD7C3DC8DE5B9099CB5A4BF431EBA1FFF3E91935978AD}.
\end{array}                                                     \tag{95.1}
\]
The older YM V40 file is also still present in the shared folder; it
is not modified here.  V41 contains V40 as its prefix and is the
newest companion record.

\subsection{New Yang--Mills result}

YM V41 executes the quartic Walsh corner pass and obtains a nonzero
exact residual for the proposed bubble--tadpole orbit cancellation.
It therefore stops the premature twelve-representative compression.
The unresolved operation is a background-dependent re-treeing map:
embed reduced variables into all links, apply the lattice symmetry,
then gauge-fix back to the original tree.  Factorwise Haar invariance
does not determine this map or its transported boundary channel.

The companion proves the exact local coordinate identity
\[
 K_y=J^{-T}K_xJ^{-1},\qquad
 \rho_y=\rho_x|\det J|^{-1},                                \tag{95.2}
\]
so
\[
                    {1\over2}\log\det K_y-\log\rho_y
                    ={1\over2}\log\det K_x-\log\rho_x.       \tag{95.3}
\]
This directly reinforces V88--V89: the Hessian, quotient density,
channel transport, and re-treeing Jacobian form one measured object.
None may be omitted when testing overlap equality or reflection
symmetry.

The V41 calculation does not supply a continuum contour, a global
quotient atlas, sector tightness, or clustering.  Its exact failed
identity is negative information that changes the order of
certification, not a new continuum input.

\subsection{Navier--Stokes result}

NS V26 is unchanged and remains confined to rank-one Oseen-kernel
bounds.  It supplies no contour-reflection or quotient-re-treeing
estimate.

\begin{decision}[V96 reflection with full coordinate transport]
\label{dec:v11v95-next}
V96 will formulate an antilinear reflection criterion for the complex
contour.  In light of YM V41, the reflection map must act on the full
slice atlas, including re-treeing, boundary channels, Hessian
carriers, quotient densities, and orientation Jacobians.  The theorem
will distinguish reality of the partition function from the stronger
positive-kernel factorisation required by
Osterwalder--Schrader positivity.
\end{decision}

\begin{firewall}[Audit boundary]
\label{fire:v11v95-boundary}
The new companion residual does not close a hypothesis of V94.
It instead rules out a shortcut: bare coordinate reflection or orbit
permutation is insufficient before the complete measured re-treeing
intertwiner has been certified.
\end{firewall}

\begin{assessment}[Progress after V95]
\label{ass:v11v95-status}
The audit imports one decisive structural constraint from YM V41:
reflection and symmetry must intertwine the fully transported
coordinate-density system.  The complex-contour branch proceeds with
that stronger carrier.
\end{assessment}

\subsection{Parent record}

This V95 note was formed by appending the preceding material to the
validated V94 source, whose record was
\[
\begin{array}{c|c}
\text{lines}&28213\\
\text{bytes}&916612\\
\text{SHA--256}&
\texttt{641026D6DF466E448AAEFFAB344DBADC5D6D2CEDFD9CDD91F2066D88E70BFE3F}.
\end{array}                                                     \tag{95.4}
\]

% =============================================================================
% END APPENDED ELEVENTH-SERIES V95 MATERIAL
% =============================================================================
% =============================================================================
% BEGIN APPENDED ELEVENTH-SERIES V96 MATERIAL
% =============================================================================

\section{Antilinear contour symmetry and reflection positivity}
\label{sec:v11v96}

Relative contour homology in V94 applies to holomorphic observables.
Osterwalder--Schrader reflection contains complex conjugation and
requires a separate argument.  This version first proves a reality
criterion and then states the stronger positive-kernel condition.

Let \(\vartheta:D\to D\) be an antilinear involution preserving the
singular complement.  For a holomorphic top form \(\omega\), define
the holomorphic reflected form
\[
                         \vartheta^\sharp\omega
                         :=\overline{\vartheta^*\omega}.      \tag{96.1}
\]
The orientation sign of the real contour is included in the statement
\(\vartheta_\#\Gamma=\Gamma\).

\begin{theorem}[Reality from antilinear contour symmetry]
\label{thm:v11v96-reality}
Suppose
\[
 \vartheta_\#\Gamma=\Gamma,\qquad
 \vartheta^\sharp\omega=\omega,                             \tag{96.2}
\]
and \(\omega\) is absolutely integrable on \(\Gamma\).  Then
\[
                              Z_\Gamma=\int_\Gamma\omega
                              \in\mathbb R.                  \tag{96.3}
\]
More generally, if an observable satisfies
\[
                              \vartheta^\sharp(F\omega)
                              =F\omega,                      \tag{96.4}
\]
then its unnormalised contour integral is real.
\end{theorem}

\begin{proof}
Complex conjugation and change of variables give
\[
 \overline{\int_\Gamma\omega}
 =\int_{\vartheta_\#\Gamma}\vartheta^\sharp\omega
 =\int_\Gamma\omega.
\]
The same calculation applies to \(F\omega\).
\end{proof}

The theorem says nothing about the sign of \(Z_\Gamma\), and reality
of every reflected expectation does not imply positivity.

\subsection{The positive-kernel criterion}

Let \({\cal A}_+\) be the algebra of cylinder observables supported at
positive Euclidean time, and define
\[
                         (\Theta F)(x)
                         =\overline{F(\vartheta x)}.          \tag{96.5}
\]
After integrating the negative-time and bulk variables, suppose the
unnormalised reflected form can be represented on a positive-time
space \(X_+\) by a measurable kernel \(K\):
\[
 {\cal Q}(F)
 =\int_{X_+\times X_+}
       \overline{F(x)}K(x,y)F(y)\,d\lambda(x)d\lambda(y).     \tag{96.6}
\]

\begin{theorem}[Kernel characterisation of cutoff reflection
positivity]
\label{thm:v11v96-kernel}
Assume (96.6) is finite for bounded cylinder \(F\).  Then
\[
                              {\cal Q}(F)\ge0
                         \quad\text{for every }F\in{\cal A}_+ \tag{96.7}
\]
if and only if \(K\) defines a positive semidefinite integral form on
the cylinder subspace.  A sufficient condition is the Gram
factorisation
\[
 K(x,y)=\int_A\overline{k_a(x)}\,k_a(y)\,d\rho(a),           \tag{96.8}
\]
where \(\rho\) is a positive measure and the integrals are absolutely
convergent.  In that case
\[
 {\cal Q}(F)
 =\int_A\left|
        \int_{X_+}k_a(y)F(y)\,d\lambda(y)
          \right|^2d\rho(a)\ge0.                            \tag{96.9}
\]
\end{theorem}

\begin{proof}
The first statement is the definition of a positive semidefinite
integral kernel on the declared test subspace.  Insert (96.8) into
(96.6), apply Fubini, and obtain (96.9).
\end{proof}

If \({\cal Q}(1)=Z>0\), the normalised functional
\(\Omega={\cal Q}/Z\) satisfies Osterwalder--Schrader positivity.
If \(Z=0\), normalisation is unavailable; if only reality is known,
the sign remains undetermined.

\begin{counterexample}[Reality is weaker than positivity]
\label{cnt:v11v96-real-not-positive}
On a two-point positive-time space with counting measure, take
\[
                              K=\begin{pmatrix}1&2\\2&1\end{pmatrix}.
                                                                    \tag{96.10}
\]
The kernel is real symmetric and hence satisfies the analogue of
(96.2), but its eigenvalues are \(3\) and \(-1\).  For
\(F=(1,-1)\),
\[
                              {\cal Q}(F)=-2<0.               \tag{96.11}
\]
\end{counterexample}

Thus an antilinearly invariant thimble can make its contribution real
while still violating reflection positivity.

\subsection{Factorisation before contour deformation}

A common sufficient structure is a time-zero boundary variable
\(a\) with positive measure \(\rho\), together with conjugate
half-space amplitudes:
\[
 \omega_{\rm full}
 \leadsto
 d\rho(a)\,
 \overline{k_a(x_-)}\,k_a(x_+)\,
 d\lambda_-(x_-)\,d\lambda_+(x_+).                          \tag{96.12}
\]
The negative contour must be the reflected image of the positive
contour with the conjugate orientation and full Jacobian.  Integrating
(96.12) gives (96.8).

\begin{proposition}[Reflection-compatible contour deformation]
\label{prop:v11v96-deform}
Suppose the positive half-contour is deformed through
\(\Gamma_+(t)\) without crossing the singular divisor and with
vanishing ends.  Deform the negative half simultaneously by
\[
                              \Gamma_-(t)=\vartheta\Gamma_+(t).
                                                                    \tag{96.13}
\]
If the time-zero measure stays positive and the half-amplitudes remain
holomorphic and conjugate as in (96.12), then the Gram form
(96.9) and reflection positivity are preserved.
\end{proposition}

\begin{proof}
Holomorphic deformation preserves each half-amplitude for fixed
boundary datum by Theorem~\ref{thm:v11v93-deformation}.
Equation (96.13) preserves their conjugate relation.  The same
positive boundary measure therefore gives the Gram representation at
every \(t\).
\end{proof}

An arbitrary deformation of the full contour may preserve the total
partition function while destroying the half-space factorisation.
Relative homology of the full cycle is therefore insufficient.

\subsection{Full slice-atlas intertwining}

Let \(a\) label a positive-time quotient chart and
\(\bar a\) its reflected negative-time chart.  A valid reflection
carrier must include a chart map
\[
             R_a:S_a\longrightarrow S_{\bar a},             \tag{96.14}
\]
the transported boundary channel, the re-gauging or re-treeing map,
and the complete local density.  If
\[
             \omega_a=j_a(z)\,dz^1\cdots dz^n,              \tag{96.15}
\]
then the reflected density condition is
\[
                         R_a^\sharp\omega_{\bar a}=\omega_a. \tag{96.16}
\]
Written in coordinates, (96.16) contains the determinant of
\(DR_a\).  Testing only
\(S_{\bar a}(R_az)=\overline{S_a(z)}\) omits that determinant and is
not enough.

\begin{theorem}[Coordinate independence of the reflected form]
\label{thm:v11v96-coordinate}
Suppose local half-space amplitudes in two slice coordinates are
related by a \(C^1\) re-treeing map with the full change-of-variables
density, and their Hessian and quotient density obey (95.2).
Then the reflected sesquilinear form (96.6) is unchanged by switching
between those coordinates.
\end{theorem}

\begin{proof}
The change-of-variables formula transports the integration form,
including \(|\det J|\).  Equation (95.2) cancels the inverse Hessian
Jacobian against the quotient-density Jacobian in the one-loop
factor.  The transported boundary channel composes \(F\) with the
same map.  Therefore both arguments and the full density in (96.6)
are transformed together, leaving the integral unchanged.
\end{proof}

The hypotheses are exactly what YM V41 found missing from its
premature corner symmetry test.

\subsection{Approximate reflection positivity and the limit}

\begin{theorem}[Vanishing negative part]
\label{thm:v11v96-approx}
Let \(\Omega_i\) converge on the reflection-stable cylinder algebra
as in V91 or V94.  Suppose
\[
                         \Omega_i((\Theta F)F)
                         \ge-\epsilon_i\|F\|_\infty^2,
                  \qquad \epsilon_i\longrightarrow0,        \tag{96.17}
\]
for every fixed positive-time cylinder \(F\).  Then
\[
                         \Omega_\infty((\Theta F)F)\ge0.     \tag{96.18}
\]
\end{theorem}

\begin{proof}
Take the limit in (96.17).  Cylinder convergence applies to
\((\Theta F)F\), and the right-hand side tends to zero.
\end{proof}

A useful sufficient estimate is an operator decomposition
\(K_i=K_i^+-E_i\), where \(K_i^+\) is positive and
\[
 |\langle F,E_iF\rangle|
       \le\epsilon_i\|F\|_\infty^2.                          \tag{96.19}
\]
The norm in (96.19) must be uniform on the declared cylinder class.
Entrywise small kernel errors can accumulate with volume.

\subsection{The rotated conformal direction}

For a time-paired conformal mode, rotating both halves independently
can produce orientation phases whose product is real but negative.
A reflection-compatible choice must pair the two orientations through
(96.13) and leave a positive time-zero factor.  Modes supported across
the reflection plane cannot be assigned to separate thimbles without
checking the resulting boundary kernel.

\begin{warning}[Ghost determinants and gauge observables]
\label{warn:v11v96-ghost}
A positive numerical Faddeev--Popov determinant on one chart does not
prove (96.8).  Fermionic and ghost Berezin integrals carry their own
reflection operation, and positivity is normally required only on the
gauge-invariant observable algebra.  The complete transported
fermion--ghost kernel must be used.
\end{warning}

\begin{firewall}[No positive contour kernel acquired]
\label{fire:v11v96-boundary}
No current SM--Einstein regulator supplies the reflection-paired
contours, positive time-zero measure, full re-treeing density
identity, and Gram factorisation (96.8) uniformly in the cutoff.
Reality of a proposed contour would not close this gap.
\end{firewall}

\begin{decision}[V97 acquisition]
\label{dec:v11v96-next}
V97 will quantify stability of a positive reflection kernel under
cutoff perturbations.  It will distinguish perturbations controlled
relative to a strictly positive transfer kernel from massless or
gauge null directions, where no uniform spectral margin exists.
\end{decision}

\begin{assessment}[Progress after V96]
\label{ass:v11v96-status}
The complex contour now has a precise reflection gate.  Antilinear
invariance proves reality, while a positive Gram kernel proves
Osterwalder--Schrader positivity.  Reflection-compatible contour
deformations preserve the Gram form only when the full slice,
re-treeing, boundary-channel, and density transport is included.
\end{assessment}

\subsection{Parent record}

This V96 note was formed by appending the preceding material to the
validated V95 source, whose record was
\[
\begin{array}{c|c}
\text{lines}&28315\\
\text{bytes}&920732\\
\text{SHA--256}&
\texttt{B76ED0BA9675E71EADEFB1DB842D88CC3259207DA1FE5E96E16C91E378291872}.
\end{array}                                                     \tag{96.20}
\]

% =============================================================================
% END APPENDED ELEVENTH-SERIES V96 MATERIAL
% =============================================================================
% =============================================================================
% BEGIN APPENDED ELEVENTH-SERIES V97 MATERIAL
% =============================================================================

\section{Stability of the reflection kernel}
\label{sec:v11v97}

V96 reduces cutoff reflection positivity to positivity of a
half-space kernel.  This version determines which perturbation norms
preserve that positivity.  A strict spectral margin gives an absolute
operator-norm test; massless and gauge-null directions require a
relative quadratic-form estimate or an exact Schur complement.

Let \({\cal H}_+\) be the completion of the positive-time test
functions in a reference \(L^2\) norm.  Let \(K_0\) be a bounded
self-adjoint operator with
\[
                              K_0\ge0,                       \tag{97.1}
\]
and let \(E=E^*\) be the cutoff or contour perturbation.

\begin{theorem}[Strict-margin stability]
\label{thm:v11v97-gap}
If
\[
                              K_0\ge\gamma I,\qquad
                              \|E\|\le\theta\gamma
                         \quad(0\le\theta\le1),              \tag{97.2}
\]
then
\[
                              K_0+E\ge(1-\theta)\gamma I.     \tag{97.3}
\]
\end{theorem}

\begin{proof}
For every \(f\in{\cal H}_+\),
\[
 \langle f,(K_0+E)f\rangle
 \ge\gamma\|f\|^2-\|E\|\|f\|^2
 \ge(1-\theta)\gamma\|f\|^2.
\]
\end{proof}

The condition is useful for a massive finite-volume block with a
uniform transfer-kernel gap.  It is unavailable when the spectrum of
\(K_0\) reaches zero.

\begin{theorem}[Gapless relative-form stability]
\label{thm:v11v97-relative}
Suppose
\[
 |\langle f,Ef\rangle|
       \le\theta\langle f,K_0f\rangle
                  \quad\text{for every }f\in{\cal H}_+,
                  \qquad 0\le\theta\le1.                    \tag{97.4}
\]
Then
\[
                              K_0+E\ge(1-\theta)K_0\ge0.      \tag{97.5}
\]
Moreover, \(E\) has no matrix elements between
\(\ker K_0\) and its orthogonal complement.
\end{theorem}

\begin{proof}
The first claim follows immediately from
\[
 \langle f,(K_0+E)f\rangle
 \ge\langle f,K_0f\rangle-|\langle f,Ef\rangle|.
\]
Let \(n\in\ker K_0\).  Equation (97.4) gives
\(\langle n,En\rangle=0\).  For arbitrary \(g\), apply (97.4) to
\(n+tg\), first with real \(t\) and then with imaginary \(t\).
The right side is \(O(|t|^2)\), while a nonzero real or imaginary part
of \(\langle n,Eg\rangle\) would give a linear term.  Hence
\(\langle n,Eg\rangle=0\).
\end{proof}

The final assertion shows why (97.4) is strong: exact null directions
must decouple.  For gauge null vectors this decoupling should follow
from an exact Ward identity or from passing to the quotient, not from
a numerical smallness claim.

\begin{counterexample}[Absolute smallness fails without a margin]
\label{cnt:v11v97-gapless}
On \(\ell^2(\mathbb N)\), let
\[
                              K_0e_n={1\over n}e_n.           \tag{97.6}
\]
For every \(\varepsilon>0\), choose \(n>\varepsilon^{-1}\) and set
\[
                              E=-\varepsilon
                                   |e_n\rangle\langle e_n|.  \tag{97.7}
\]
Then \(\|E\|=\varepsilon\), but
\[
                              \langle e_n,(K_0+E)e_n\rangle
                              ={1\over n}-\varepsilon<0.     \tag{97.8}
\]
\end{counterexample}

Thus an error tending to zero in operator norm can still violate
positivity at every cutoff if the reference margin tends to zero
faster.

\subsection{Exact low/null block criterion}

Decompose
\[
                         {\cal H}_+={\cal N}\oplus{\cal R},
 \qquad
 K_0=\begin{pmatrix}0&0\\0&A\end{pmatrix},\qquad
 A\ge\gamma I_{\cal R}.                                     \tag{97.9}
\]
Write
\[
 K_0+E=
 \begin{pmatrix}
 B&C^*\\
 C&A+D
 \end{pmatrix}.                                              \tag{97.10}
\]

\begin{theorem}[Schur complement for reflected null modes]
\label{thm:v11v97-schur}
If \(A+D>0\), then
\[
 K_0+E\ge0
 \quad\Longleftrightarrow\quad
 B-C^*(A+D)^{-1}C\ge0.                                      \tag{97.11}
\]
In particular, if \(B=0\), positivity forces \(C=0\).
\end{theorem}

\begin{proof}
For \(n\in{\cal N}\), \(r\in{\cal R}\), complete the square:
\[
 \begin{split}
 &\left\langle
 \binom nr,
 \begin{pmatrix}B&C^*\\C&A+D\end{pmatrix}
 \binom nr\right\rangle\\
 &\quad=
 \langle r+(A+D)^{-1}Cn,\,
        (A+D)[r+(A+D)^{-1}Cn]\rangle\\
 &\qquad+
 \langle n,[B-C^*(A+D)^{-1}C]n\rangle.
 \end{split}
\]
This proves (97.11).  If \(B=0\), the second term is
\(-\|(A+D)^{-1/2}Cn\|^2\), so it can be nonnegative for every \(n\)
only when \(C=0\).
\end{proof}

Physical massless modes need not lie in the exact kernel at finite
volume, but their margins can vanish in the continuum.  They should
be retained in \({\cal N}\) or a low block and treated by (97.11),
while a strict estimate is used only on \({\cal R}\).

\subsection{Approximate positivity}

If no exact positive comparison is available, an additive form bound
still quantifies the failure.

\begin{proposition}[Negative-part bound]
\label{prop:v11v97-negative}
If
\[
                              \|E\|\le\varepsilon,            \tag{97.12}
\]
then
\[
                    \langle f,(K_0+E)f\rangle
                              \ge-\varepsilon\|f\|^2.        \tag{97.13}
\]
If the normalising denominator satisfies \(Z_i\ge\zeta>0\), the
normalised reflected form has negative part at most
\(\varepsilon\|f\|^2/\zeta\).
\end{proposition}

\begin{proof}
Use \(K_0\ge0\) and
\(\langle f,Ef\rangle\ge-\|E\|\|f\|^2\), then divide by \(Z_i\).
\end{proof}

For the limit theorem V96, \(\varepsilon_i/\zeta_i\to0\) is enough for
each fixed test vector.  Exact positivity at finite cutoff requires
one of Theorems~\ref{thm:v11v97-gap},
\ref{thm:v11v97-relative}, or
\ref{thm:v11v97-schur}.

\subsection{Kernel errors from contour comparison}

Let \(K_i^0\) be a reflection-positive kernel obtained from a paired
contour factorisation, and let \(E_i\) contain the fine-integrand,
end, singular-excision, reference-block, and slice-atlas mismatches.
The scalar bound (94.11) for each fixed \(F\) does not by itself give
a uniform operator estimate.

\begin{lemma}[Uniform bilinear error criterion]
\label{lem:v11v97-bilinear}
Suppose a Hermitian sesquilinear error form \(e_i\) satisfies
\[
                         |e_i(f,g)|
                         \le\varepsilon_i\|f\|\,\|g\|        \tag{97.14}
\]
on a dense test subspace.  Then it extends to a bounded self-adjoint
operator \(E_i\) with \(\|E_i\|\le\varepsilon_i\).
\end{lemma}

\begin{proof}
Boundedness gives a unique continuous extension.  The Riesz
representation theorem yields \(E_i\), Hermiticity makes it
self-adjoint, and the operator norm is at most \(\varepsilon_i\).
\end{proof}

To use the relative theorem, the stronger weighted estimate is
\[
 |e_i(f,f)|
       \le\theta_i\langle f,K_i^0f\rangle,\qquad
                              \theta_i\le1.                  \tag{97.15}
\]
This is a statement about the full quadratic form and cannot be
replaced by pointwise closeness of the kernel entries.

\subsection{A common overall contour phase}

Sometimes the unnormalised reflected form has
\[
                              {\cal Q}_i(F)
                              =c_i{\cal P}_i(F),             \tag{97.16}
\]
where \(c_i\ne0\) is independent of \(F\) and
\({\cal P}_i(F)\ge0\).  Then
\[
                              Z_i={\cal Q}_i(1)
                              =c_i{\cal P}_i(1),             \tag{97.17}
\]
and, if \({\cal P}_i(1)>0\),
\[
                 {{\cal Q}_i(F)\over Z_i}
                 ={{\cal P}_i(F)\over{\cal P}_i(1)}\ge0.     \tag{97.18}
\]
Thus an extensive orientation phase can cancel in normalised
expectations only when it is truly common to the entire reflected
form.  Observable-dependent thimble phases cannot be removed this
way.

\begin{warning}[Sector and saddle mixtures]
\label{warn:v11v97-mixture}
A sum of positive kernels with nonnegative coefficients is positive.
A sum with real or complex coefficients of mixed sign need not be.
Reality from pairing conjugate saddles does not control the signs of
the resulting coefficients.
\end{warning}

\subsection{Acquisition test}

A reflection-positive complex-contour regulator should prove:
\begin{enumerate}
\item an exact Gram kernel \(K_i^0\) after full re-treeing and quotient
      transport;
\item a strictly positive margin on the massive complement;
\item exact Ward decoupling or a positive Schur complement on gauge
      and massless low directions;
\item a uniform bilinear or relative-form bound for every contour and
      refinement error;
\item a positive normalisation, possibly after cancelling one common
      orientation phase.
\end{enumerate}

\begin{firewall}[No reflection-kernel estimate acquired]
\label{fire:v11v97-boundary}
The current programme has not constructed \(K_i^0\), proved a
relative form bound (97.15), or evaluated the null-block Schur
complement for the coupled gauge--gravity contour.  Hence reflection
positivity remains an open input.
\end{firewall}

\begin{decision}[V98 acquisition]
\label{dec:v11v97-next}
V98 will analyse topological and thimble-sector mixtures in the
reflection kernel.  It will prove when theta phases are removable by
a unitary half-space conjugation and exhibit why conjugate-sector
pairing alone gives reality but not positivity.
\end{decision}

\begin{assessment}[Progress after V97]
\label{ass:v11v97-status}
Reflection positivity now has the correct perturbative norms.
Absolute operator smallness works only with a strict margin;
gapless directions require relative form control, and exact null
directions obey a Schur-complement constraint that forces their
coupling to vanish unless a positive null-block term absorbs it.
\end{assessment}

\subsection{Parent record}

This V97 note was formed by appending the preceding material to the
validated V96 source, whose record was
\[
\begin{array}{c|c}
\text{lines}&28615\\
\text{bytes}&931897\\
\text{SHA--256}&
\texttt{EF3752BFAA44170FC4150AC251CA27115F9EC8E1C77798B53CF5531C97DF57B4}.
\end{array}                                                     \tag{97.19}
\]

% =============================================================================
% END APPENDED ELEVENTH-SERIES V97 MATERIAL
% =============================================================================
% =============================================================================
% BEGIN APPENDED ELEVENTH-SERIES V98 MATERIAL
% =============================================================================

\section{Theta phases and saddle mixtures in the reflection kernel}
\label{sec:v11v98}

A nonzero sector denominator from V76--V79 is not enough for
reflection positivity.  The theta and thimble phases also enter the
half-space kernel.  This version identifies the exact case in which
they are harmless: the phase must be a unitary coboundary between the
two reflected halves.

Let \(K_0\) be a positive semidefinite kernel on \(X_+\), and let
\(q:X_+\to\mathbb Z\) be measurable.  Define
\[
 K_\theta(x,y)
       =e^{-i\theta q(x)}K_0(x,y)e^{i\theta q(y)}.            \tag{98.1}
\]

\begin{theorem}[Half-space theta conjugation]
\label{thm:v11v98-unitary}
For every real \(\theta\), the kernel \(K_\theta\) is positive
semidefinite.  If \(U_\theta\) is multiplication by
\(e^{i\theta q}\), then
\[
                              K_\theta=U_\theta^*K_0U_\theta. \tag{98.2}
\]
\end{theorem}

\begin{proof}
For every test function \(F\),
\[
 \langle F,K_\theta F\rangle
 =\langle U_\theta F,K_0U_\theta F\rangle\ge0.
\]
This is (98.2).
\end{proof}

The theorem applies when the full topological charge on a reflected
configuration decomposes as
\[
                      Q(x_-,x_+)=q(x_+)-q(\vartheta x_-),    \tag{98.3}
\]
with no residual time-zero term.  Then
\(e^{i\theta Q}\) supplies exactly the two half-space multipliers in
(98.1).

\subsection{The phase-cocycle criterion}

Let \(\chi:X_+\times X_+\to U(1)\) satisfy
\[
 \chi(x,x)=1,\qquad
 \chi(y,x)=\overline{\chi(x,y)}.                             \tag{98.4}
\]

\begin{theorem}[Unit-modulus coboundary test]
\label{thm:v11v98-cocycle}
The following are equivalent:
\begin{enumerate}
\item
\[
                       \chi(x,y)\chi(y,z)=\chi(x,z)
                       \quad\text{for all }x,y,z;            \tag{98.5}
\]
\item there is \(v:X_+\to U(1)\) such that
\[
                              \chi(x,y)=\overline{v(x)}v(y); \tag{98.6}
\]
\item multiplication of any positive kernel \(K_0\) by \(\chi\)
      is unitary conjugation by a multiplication operator.
\end{enumerate}
\end{theorem}

\begin{proof}
Fix \(x_0\) and set \(v(x)=\chi(x_0,x)\).  Equations
(98.4)--(98.5) give
\[
 \chi(x,y)=\chi(x,x_0)\chi(x_0,y)
          =\overline{v(x)}v(y),
\]
so one implies two.  Formula (98.6) gives the unitary conjugation
directly, proving two implies three.  Apply three to the constant
rank-one positive kernel and read off its phase factor to recover
(98.6), which implies (98.5).
\end{proof}

Thus the safe theta phase is not merely reflection-real; it is a
one-coboundary on the half-space kernel.  Failure of (98.5) measures
an interface or global topological obstruction.

\subsection{A time-zero remainder}

Suppose instead
\[
 Q(x_-,a,x_+)
 =q(x_+)-q(\vartheta x_-)+r(a),                              \tag{98.7}
\]
where \(a\) is time-zero boundary data.  A Gram representation then
has the form
\[
 K_\theta(x,y)
 =\int_A e^{i\theta r(a)}
        \overline{k_a(x)e^{i\theta q(x)}}
        k_a(y)e^{i\theta q(y)}\,d\rho(a).                   \tag{98.8}
\]

\begin{proposition}[Boundary-charge positivity]
\label{prop:v11v98-boundary}
The kernel (98.8) is positive if
\[
                              e^{i\theta r(a)}\ge0
                    \quad\text{for }\rho\text{-almost every }a.
                                                                    \tag{98.9}
\]
Because the factor has unit modulus, (98.9) means it equals one almost
everywhere, apart from a single common phase that cancels from all
normalised reflected expectations.
\end{proposition}

\begin{proof}
Under (98.9), absorb the nonnegative factor into the positive measure
in the Gram representation and apply
Theorem~\ref{thm:v11v96-kernel}.  A unit-modulus nonnegative real
number is one.  If the same constant phase multiplies the entire
quadratic form and its denominator, use (97.16)--(97.18).
\end{proof}

The condition is sufficient rather than necessary when the functions
\(k_a\) have linear dependencies.  Without such special
cancellations, a varying complex boundary coefficient is a direct
positivity obstruction.

\subsection{Reality of conjugate sectors}

Let \(K_q\ge0\) be sector kernels and suppose
\[
                              K_{-q}=K_q.                    \tag{98.10}
\]
For symmetric positive weights \(p_{-q}=p_q\), theta pairing gives
\[
 K_\theta
 =p_0K_0+\sum_{q\ge1}2p_q\cos(\theta q)K_q.                 \tag{98.11}
\]
This kernel is Hermitian and real in the reflected form, but some
coefficients can be negative.

\begin{counterexample}[Conjugate pairing is not positive]
\label{cnt:v11v98-pairing}
On \(\mathbb C^2\), let
\[
 K_0=\begin{pmatrix}1&0\\0&0\end{pmatrix},\qquad
 K_1=K_{-1}=\begin{pmatrix}0&0\\0&1\end{pmatrix},            \tag{98.12}
\]
with \(p_0,p_1>0\).  Then
\[
 K_\theta=
 \begin{pmatrix}p_0&0\\0&2p_1\cos\theta\end{pmatrix}.        \tag{98.13}
\]
For \(\pi/2<\theta<3\pi/2\), this is real but has a negative
eigenvalue.
\end{counterexample}

A lower bound on the characteristic function
\(\sum_qp_qe^{i\theta q}\) controls the denominator but not every
eigenchannel of the kernel.  Therefore the small-theta anchor of V76
does not replace the coboundary or kernel-positivity test.

\subsection{Positive mixtures and common phases}

\begin{proposition}[Safe kernel mixtures]
\label{prop:v11v98-mixture}
Let \(K_\alpha\ge0\).
\begin{enumerate}
\item If \(w_\alpha\ge0\) and
      \(\sum_\alpha w_\alpha\|K_\alpha\|<\infty\), then
\[
                              K=\sum_\alpha w_\alpha K_\alpha\ge0.
                                                                    \tag{98.14}
\]
\item If all terms share a common nonzero phase \(c\),
      so the unnormalised kernel is
      \(c\sum_\alpha w_\alpha K_\alpha\), that phase cancels in the
      normalised reflected form.
\item Pairing coefficients into real numbers is insufficient if their
      signs differ on linearly independent positive kernels.
\end{enumerate}
\end{proposition}

\begin{proof}
The quadratic form of (98.14) is a convergent sum of nonnegative
numbers.  The second statement is (97.16)--(97.18).  The
counterexample (98.12)--(98.13) proves the third.
\end{proof}

\subsection{Thimble coefficients}

A finite-dimensional Lefschetz decomposition has the schematic form
\[
                              \Gamma=\sum_\sigma n_\sigma
                                      {\cal J}_\sigma,       \tag{98.15}
\]
where \(n_\sigma\in\mathbb Z\) are intersection numbers.  Each thimble
can have constant action phase and a positive amplitude, yet the
coefficient entering the reflected kernel includes
\[
 n_\sigma e^{-i\operatorname{Im}S(z_\sigma)}
          \times\hbox{orientation and residual Jacobian phase}.    \tag{98.16}
\]
Antilinear symmetry pairs conjugate coefficients and makes their sum
real.  Proposition~\ref{prop:v11v98-mixture} shows why this does not
prove positivity.

A safe thimble construction must instead show one of:
\begin{enumerate}
\item the paired thimbles combine before half-space integration into
      a single Gram factor;
\item all resulting independent Gram kernels have nonnegative
      coefficients after removal of one common phase;
\item the entire theta or saddle phase is the half-space coboundary
      (98.6).
\end{enumerate}

\subsection{Refinement stability}

Suppose the cutoff phase cocycle is
\[
 \chi_i(x,y)=\overline{v_i(x)}v_i(y)                         \tag{98.17}
\]
and its adjacent transport obeys
\[
 \int|v_{i+1}(x)-v_i(P_ix)|^2\,d\mu_{i+1}(x)
                              \le\delta_i^2.                 \tag{98.18}
\]
For a bounded positive reference kernel operator \(K_i^0\), the
induced conjugation error on a fixed test function satisfies
\[
 \begin{split}
 &|\langle F,U_{i+1}^*K_i^0U_{i+1}F\rangle
   -\langle F,U_i^*K_i^0U_iF\rangle|\\
 &\qquad\le
 2\|K_i^0\|\,\|F\|_\infty\|F\|_{L^2}\,\delta_i,             \tag{98.19}
 \end{split}
\]
after placing both cutoffs under the same transported amplitude law.

\begin{proof}
Insert one intermediate factor and use
\[
 |\langle a,Ka\rangle-\langle b,Kb\rangle|
 \le\|K\|(\|a\|+\|b\|)\|a-b\|.
\]
Multiplication by unit phases preserves the \(L^2\) norm, while
(98.18) and boundedness of \(F\) bound the difference.
\end{proof}

The measure-transport error must be added separately.  Summability of
\(\delta_i\) is enough for cylinder convergence, but exact positivity
at each cutoff already follows from the coboundary form (98.17).

\begin{warning}[Topological charge must split at the reflection
plane]
\label{warn:v11v98-splitting}
The global integer charge of V77 need not admit the local
decomposition (98.3).  Boundary Chern--Simons terms, eta invariants,
or gravitational signature defects can contribute \(r(a)\) in
(98.7).  Their phases must be included in the boundary kernel rather
than discarded modulo an integer.
\end{warning}

\begin{firewall}[No theta coboundary acquired]
\label{fire:v11v98-boundary}
The current SM--Einstein regulator has no constructed reflection-plane
charge splitting, boundary eta-invariant control, or thimble Gram
decomposition.  V98 supplies the criterion and counterexample, not
its verification for the physical theta and gravitational sectors.
\end{firewall}

\begin{decision}[V99 acquisition]
\label{dec:v11v98-next}
V99 will consolidate the complex-contour branch into a dependency
ledger and prove a combined conditional theorem using V93--V98.
It will identify the first constructive object required after the
cycle rollover: an explicit reflection-paired conformal contour for
one regulated block, including its singular divisor and new-mode
normalisation.
\end{decision}

\begin{assessment}[Progress after V98]
\label{ass:v11v98-status}
Theta phases preserve the reflection kernel when they are exact
half-space coboundaries, hence unitary conjugations.  Time-zero
remainders and independent sector or thimble coefficients require
positive boundary weights.  Conjugate pairing proves reality only,
as the explicit two-channel counterexample demonstrates.
\end{assessment}

\subsection{Parent record}

This V98 note was formed by appending the preceding material to the
validated V97 source, whose record was
\[
\begin{array}{c|c}
\text{lines}&28928\\
\text{bytes}&942249\\
\text{SHA--256}&
\texttt{E3A8C675BE75CB4B3A945CB0F19F95053F459D4932E9865E276F5D628E9B1B99}.
\end{array}                                                     \tag{98.20}
\]

% =============================================================================
% END APPENDED ELEVENTH-SERIES V98 MATERIAL
% =============================================================================
% =============================================================================
% BEGIN APPENDED ELEVENTH-SERIES V99 MATERIAL
% =============================================================================

\section{A projective imaginary-Gaussian conformal block}
\label{sec:v11v99}

The contour criteria V93--V98 are conditional for the coupled theory.
This version constructs their simplest nontrivial component exactly:
a growing family of negative quadratic conformal modes integrated on
imaginary axes, with the extensive orientation phase removed by the
declared new-mode normalisation.

Let \(\kappa_n>0\), let
\[
                              S_N(z)
                              =-{1\over2}\sum_{n=1}^N
                                      \kappa_n z_n^2,        \tag{99.1}
\]
and choose
\[
                              \Gamma_N=(i\mathbb R)^N.       \tag{99.2}
\]
Define the holomorphic top form
\[
 d\omega_N(z)
 =\prod_{n=1}^N
   \left[
    \left({\kappa_n\over2\pi}\right)^{1/2}
       {dz_n\over i}\,
       e^{\kappa_n z_n^2/2}
   \right].                                                  \tag{99.3}
\]

\begin{theorem}[Exact normalisation and projectivity]
\label{thm:v11v99-gaussian}
For every \(N\),
\[
                              \int_{\Gamma_N}d\omega_N=1.    \tag{99.4}
\]
If \(P_N:\Gamma_{N+1}\to\Gamma_N\) forgets the last coordinate, then
\[
                              (P_N)_\#\omega_{N+1}=\omega_N. \tag{99.5}
\]
Consequently every polynomial cylinder observable has an exactly
consistent expectation.
\end{theorem}

\begin{proof}
Write \(z_n=iy_n\).  Then \(dz_n/i=dy_n\) and
\[
 e^{\kappa_nz_n^2/2}=e^{-\kappa_ny_n^2/2}.
\]
Each factor in (99.3) is a normalised real Gaussian measure in
\(y_n\), proving (99.4).  Fubini's theorem and integration of the
last normalised factor prove (99.5).
\end{proof}

The division by \(i\) in every factor is the reference-block
normalisation (94.2)--(94.3).  Without it the raw contour integral has
the phase \(i^N\), so adjacent cutoffs are not projective as complex
measures.

\begin{corollary}[Cylinder moments]
\label{cor:v11v99-moments}
Under \(\omega_N\),
\[
 \mathbb E z_n=0,\qquad
 \mathbb E z_nz_m=-{\delta_{nm}\over\kappa_n},               \tag{99.6}
\]
and Wick's rule holds with covariance
\(-\operatorname{diag}(\kappa_n^{-1})\).
\end{corollary}

\begin{proof}
Since \(z_n=iy_n\) and
\(\mathbb E y_ny_m=\delta_{nm}/\kappa_n\), (99.6) follows.
The ordinary real Gaussian Wick formula supplies the higher moments.
\end{proof}

The negative holomorphic covariance is not a probability covariance
for real conformal fields.  It records that the continuum variable
lives on an imaginary contour.

\subsection{Sobolev tightness of the contour amplitude}

Let \(\lambda_n\ge1\) define the Hilbert scale (87.1), and regard
\(z=iy\) as a random variable under the positive pulled-back Gaussian
amplitude.

\begin{theorem}[Imaginary Gaussian continuum]
\label{thm:v11v99-tight}
If
\[
                              \sum_{n\ge1}
                              {\lambda_n^{-s}\over\kappa_n}
                              <\infty,                       \tag{99.7}
\]
then
\[
                              z_N=i\sum_{n\le N}y_ne_n       \tag{99.8}
\]
converges in \(L^2({\cal H}^{-s}_{\mathbb C})\) and almost surely.
Its amplitude laws are tight, and the exact projective cylinder
functional of Theorem~\ref{thm:v11v99-gaussian} has an infinite
Gaussian realisation supported on \(i{\cal H}^{-s}\).
\end{theorem}

\begin{proof}
Exactly as in Theorem~\ref{thm:v11v87-gaussian},
\[
 \mathbb E\|z_M-z_N\|_{-s}^2
 =\sum_{N<n\le M}{\lambda_n^{-s}\over\kappa_n}.
\]
Condition (99.7) gives \(L^2\), almost-sure, and distributional
convergence.  Multiplication by \(i\) is an isometry of the complex
Hilbert space.
\end{proof}

\subsection{A limited reflection check}

If all variables in a finite block lie on the reflection plane, define
\[
                              (\Theta F)(iy)=\overline{F(iy)}. \tag{99.9}
\]

\begin{proposition}[Reflection-plane positivity]
\label{prop:v11v99-plane}
For every polynomial \(F\) of the finite block,
\[
                         \int_{\Gamma_N}(\Theta F)F\,d\omega_N
                         =\int_{\mathbb R^N}|F(iy)|^2d\nu_N(y)
                         \ge0,                              \tag{99.10}
\]
where \(\nu_N\) is the positive product Gaussian in the proof of
Theorem~\ref{thm:v11v99-gaussian}.
\end{proposition}

\begin{proof}
Pull (99.10) back by \(z=iy\) and use (99.3).
\end{proof}

This proposition covers only a decoupled reflection-plane block.
A spacetime conformal kinetic operator couples positive and negative
halves; its reflected kernel must satisfy V96--V97.  Equation (99.10)
cannot be tensorially assigned to nonlocal interacting modes without
that calculation.

\subsection{Perturbing the exact block}

Let \(V_N\) be holomorphic near \(\Gamma_N\) and define
\[
                              d\omega_N^V
                 ={e^{-V_N}\over Z_N^V}\,d\omega_N,\qquad
 Z_N^V=\int_{\Gamma_N}e^{-V_N}d\omega_N.                    \tag{99.11}
\]
After pullback to \(y\), this is generally a complex reweighting.
A sufficient absolute-integrability condition is
\[
                         \operatorname{Re}V_N(iy)
                         \ge-\alpha\sum_{n\le N}\kappa_ny_n^2-C,
                    \qquad \alpha<{1\over2}.                \tag{99.12}
\]
Indeed the total real Gaussian exponent then retains coefficient
\((1/2-\alpha)\kappa_n\).

\begin{proposition}[Uniform perturbative denominator criterion]
\label{prop:v11v99-perturb}
Suppose, uniformly in \(N\),
\[
                    \mathbb E_{\nu_N}|e^{-V_N(iy)}-1|
                    \le\epsilon<1.                          \tag{99.13}
\]
Then
\[
                              |Z_N^V-1|\le\epsilon,\qquad
                              |Z_N^V|\ge1-\epsilon.           \tag{99.14}
\]
If the conditional-density increments satisfy
\[
 \sum_N\mathbb E_{\nu_{N+1}}
 \left|
 e^{-V_{N+1}(iy)}
 -e^{-V_N(iP_Ny)}
 \right|<\infty,                                            \tag{99.15}
\]
then the unnormalised cylinder numerators and denominators are
Cauchy.
\end{proposition}

\begin{proof}
Equation (99.14) is the triangle inequality applied to
\(Z_N^V-1=\mathbb E_{\nu_N}(e^{-V_N}-1)\).
For a bounded coarse cylinder \(F\), exact Gaussian projectivity gives
the adjacent difference
\[
 \left|
 \mathbb E_{\nu_{N+1}}
 F(iP_Ny)
 [e^{-V_{N+1}}-e^{-V_N\circ P_N}]
 \right|
 \le\|F\|_\infty
 \mathbb E_{\nu_{N+1}}
 |e^{-V_{N+1}}-e^{-V_N\circ P_N}|.
\]
Summation proves the claim.
\end{proof}

Condition (99.13) is strong but noncircular: it can be checked under
the explicit positive Gaussian amplitude.  It is not known for the
nonlinear Einstein, gauge, ghost, and matter interaction.

\subsection{Conditional complex-contour assembly}

\begin{theorem}[Contour branch implication]
\label{thm:v11v99-assembly}
Assume a sequence of coupled cutoff contours has:
\begin{enumerate}
\item the coercivity and singular avoidance of V93;
\item normalised new-mode blocks and summable adjacent errors of V94;
\item a denominator anchor stronger than the total error tail;
\item reflection-paired half contours with the Gram or vanishing
      negative-part condition of V96--V97;
\item theta and saddle phases satisfying the coboundary or positive
      mixture criterion of V98;
\item convergent Ward insertions and a uniform clustering estimate as
      in V91.
\end{enumerate}
Then the holomorphic cylinder expectations have a nonzero normalised
projective limit, and the reflected cylinder algebra is positive in
that limit.  Ward identities and clustering pass to the limit under
their stated uniform hypotheses.
\end{theorem}

\begin{proof}
Items one through three give Theorem~\ref{thm:v11v94-limit}.
Items four and five give cutoff reflection positivity or a negative
part tending to zero; Theorem~\ref{thm:v11v96-approx} passes it to the
limit.  The passage statements for Ward identities and clustering are
Theorem~\ref{thm:v11v91-gates}.
\end{proof}

\begin{firewall}[The constructed block is not the coupled theory]
\label{fire:v11v99-boundary}
Theorem~\ref{thm:v11v99-gaussian} constructs only independent
quadratic negative conformal modes.  It has no nonlinear metric
positivity, diffeomorphism quotient, Faddeev--Popov singular divisor,
fermion determinant, gauge interaction, or cross-time kinetic
kernel.  The conditions of
Theorem~\ref{thm:v11v99-assembly} remain unverified for the full
SM--Einstein regulator.
\end{firewall}

\begin{decision}[V100 audit, consolidation, and rollover]
\label{dec:v11v99-next}
V100 will perform the compulsory companion audit and record the final
dependency position of this cycle.  It will then be frozen as
\texttt{\_f11}.  A new compact V1 will begin with context and a summary
of the major acquired results, then continue by testing interacting
holomorphic perturbations of the exact imaginary-Gaussian block
against singular clearance, reflection factorisation, and adjacent
summability.
\end{decision}

\begin{assessment}[Progress after V99]
\label{ass:v11v99-status}
The complex branch now contains one exact projective conformal
building block.  Normalising each imaginary-axis Gaussian by its
orientation phase removes the \(i^N\) drift, condition (99.7) gives
an infinite Sobolev-valued amplitude, and (99.13)--(99.15) state
direct perturbative tests.  Coupling this block to the full regulated
theory remains open.
\end{assessment}

\subsection{Parent record}

This V99 note was formed by appending the preceding material to the
validated V98 source, whose record was
\[
\begin{array}{c|c}
\text{lines}&29241\\
\text{bytes}&953021\\
\text{SHA--256}&
\texttt{46D62EB66BC2627EDDB08AAF67D864A0FE80185EEC8554076327EEE19A1A0897}.
\end{array}                                                     \tag{99.16}
\]

% =============================================================================
% END APPENDED ELEVENTH-SERIES V99 MATERIAL
% =============================================================================
% =============================================================================
% BEGIN APPENDED ELEVENTH-SERIES V100 MATERIAL
% =============================================================================

\section{V100 companion audit and eleventh-cycle consolidation}
\label{sec:v11v100}

The compulsory audit records the current companion endpoints:
\[
\begin{array}{c|r|r|l}
\text{file}&\text{lines}&\text{bytes}&\text{SHA--256}\\ \hline
\texttt{Renewal\_Yang\_Mills\_Mass\_Gap\_Research\_Notes\_V41.tex}
&16811&574386&
\texttt{191ACA59453A84990AABF333FA79E534A760B1B1A6F22AFA326B5D9E6AE1502A}\\
\texttt{Renewal\_NS\_Stability\_Research\_Notes\_V26.tex}
&5319&278283&
\texttt{82C887F8C2C69E4F28FAD7C3DC8DE5B9099CB5A4BF431EBA1FFF3E91935978AD}.
\end{array}                                                    \tag{100.1}
\]
Both are unchanged from V95.  YM V41 still requires the explicit
spectator re-treeing and transported-channel certificate before its
corner response can be symmetry-compressed.  NS V26 remains at its
rank-one Oseen-kernel block.  Neither supplies a new continuum
contour, gravitational quotient, sector, phase, or reflection input.

\subsection{Results acquired in the closing block}

The V76--V99 block established the following implication chain.
\begin{enumerate}
\item An energy advantage over sector entropy gives exponential charge
      tails, exponential moments, and a small-theta denominator
      window.
\item A gapped Hermitian index carrier gives integer sectors and an
      exact refinement criterion; \(L^r\) change of measure transfers
      their tails and admissibility probabilities.
\item A finite denominator anchor survives summable amplitude, charge,
      and determinant-phase increments.
\item Relative trace-class or Hilbert--Schmidt operator increments,
      local exact counterterms, and trivial determinant-line holonomy
      yield a quantitative phase-transport budget.
\item The bare Euclidean Einstein--Hilbert action is unbounded below
      along bounded, fixed-volume, fixed-topology conformal
      oscillations.
\item A coercive Gaussian reference has a projective Sobolev-valued
      continuum, but the Einstein interaction still needs a uniform
      relative lower bound and conditional-density series.
\item A local diffeomorphism slice follows from a quantitative
      Faddeev--Popov gap; a global positive quotient additionally
      requires full-density overlap equality, multiplicity and
      stabiliser corrections, and equivariant refinement.
\item A naive positive fourth-order conformal regulator cannot be
      removed while retaining both uniform \(L^r\) Einstein
      reweighting and nondegenerate fixed conformal modes.
\item A complex contour is absolutely convergent under coercivity,
      invariant under singularity-free relative homology, and
      projective under summable fine-form, end, excision, and
      reference-block errors.
\item Antilinear invariance gives reality; reflection positivity
      requires a Gram kernel.  Gapless directions need relative form
      control, and theta phases are harmless when they are exact
      half-space coboundaries.
\item The negative quadratic conformal block has an exact normalised
      imaginary-axis Gaussian contour with a projective infinite
      Sobolev-valued amplitude.
\end{enumerate}

\begin{theorem}[Eleventh-cycle claim boundary]
\label{thm:v11v100-boundary}
The cited results give rigorous conditional constructions and one
exact quadratic conformal contour block.  They do not construct the
nonlinear, quotient-complete, reflection-positive, clustered
SM--Einstein continuum.  In particular, the coupled contour,
singular divisor, diffeomorphism atlas, determinant-line transport,
sector energy--entropy constants, Ward limits, and physical
universality identification remain unacquired.
\end{theorem}

\begin{proof}
Each positive statement is proved in V76--V99.  The missing inputs are
the explicit firewalls and acquisition ledgers of those versions.
No theorem in either companion file supplies them.
\end{proof}

\subsection{Rollover instruction}

This source has reached V100.  In accordance with the standing file
policy, it is the terminal active version of the eleventh compact
cycle and is to be preserved as
\[
 \texttt{Renewal\_SM\_Einstein\_Continuum\_Research\_Notes\_V100\_f11.tex}.
                                                                    \tag{100.2}
\]
The next active note restarts at V1.  Its opening section will give
the programme context and a compact summary of the major mathematical
results, then continue with interacting holomorphic perturbations of
the exact imaginary-Gaussian conformal block.

\begin{decision}[Twelfth-cycle V1 acquisition]
\label{dec:v11v100-next}
For a local polynomial interaction on the imaginary Gaussian block,
derive a verifiable sectorial real-part bound, a zero-free
denominator estimate, and an adjacent martingale-difference
criterion.  Keep reflection factorisation separate and require the
full re-treeing density identified by the V95 and V100 audits.
\end{decision}

\begin{assessment}[Progress after V100]
\label{ass:v11v100-status}
The eleventh cycle closes with the real positive-gravity obstruction,
the alternative complex-contour architecture, and an exact projective
quadratic conformal block all recorded.  The first task of the new
cycle is no longer to choose a formal contour, but to control a
specific interacting perturbation of that block.
\end{assessment}

\subsection{Parent record}

This V100 note was formed by appending the preceding material to the
validated V99 source, whose record was
\[
\begin{array}{c|c}
\text{lines}&29531\\
\text{bytes}&963057\\
\text{SHA--256}&
\texttt{729C2BACB4AF86FFA13DC5C53861A897EB84773E6F05B6D397D1111BA562BC3F}.
\end{array}                                                    \tag{100.3}
\]

% =============================================================================
% END APPENDED ELEVENTH-SERIES V100 MATERIAL
% =============================================================================
\end{document}
